\documentclass[11pt,a4paper,openany,ja=standard]{ltjsbook}

% --------------------------------------------------
% 数学・科学・図版パッケージ
% --------------------------------------------------
\usepackage{amsmath,amssymb,amsfonts,amsthm}
\usepackage{mathtools}
\usepackage{bm}
\usepackage{graphicx}
\usepackage{wrapfig}
\usepackage{xcolor}
\usepackage{tikz}
\usetikzlibrary{calc,angles,quotes,intersections,patterns,arrows.meta}
\usepackage{tcolorbox}
\tcbuselibrary{skins,breakable}
\usepackage{geometry}
\geometry{top=25mm,bottom=25mm,left=25mm,right=25mm}
\usepackage{hyperref}
\hypersetup{
    colorlinks=true,
    linkcolor=blue!80!black,
    citecolor=green!60!black,
    urlcolor=blue!90!black,
    pdftitle={自然哲学の数学的諸原理（プリンキピア）全訳版},
    pdfauthor={アイザック・ニュートン / Project Gutenberg No. 28233}
}

% --------------------------------------------------
% 原著対応マクロ
% --------------------------------------------------
\newcount \Zw
\newcount \Zh
\newcounter{wrapwidth}

\usepackage{adjustbox}
\newcommand\pngright[4]{%
\par\begin{center}%
\IfFileExists{images/#1}{\adjustbox{max width=0.85\linewidth,keepaspectratio}{\includegraphics{images/#1}}}{\fbox{#1}}%
\end{center}\par}

\newcommand\pngcent[3]{%
\par\begin{center}%
\IfFileExists{images/#1}{\adjustbox{max width=0.85\linewidth,keepaspectratio}{\includegraphics{images/#1}}}{\fbox{#1}}%
\end{center}\par}

\newcommand\pngrightsc[5]{%
}

\renewcommand\minute{^{\prime}}
\newcommand\second{^{\prime\prime}}
\newcommand\third{^{\prime\prime\prime}}
\newcommand\fourth{^{\mathrm{iv}}}
\newcommand\fifth{^{\mathrm{v}}}
\newcommand\fsp{\hspace{0.3em}}
\newcommand\wsp{\hspace{0.5em}}
\newcommand\longs{\textrm{\raisebox{1pt}{$\scriptstyle \int$}\hspace{-1.5pt}}}
\newcommand\gesperrt[2]{#1}
\newcommand\spreadout[1]{#1}
\newcommand\condpagelarge[1]{\vspace{0.8em}\noindent\textbf{\large #1}\par\vspace{0.4em}}
\newcommand\tblong{\vspace{0.5em}}
\newcommand\QED{\hfill$\blacksquare$}
\newcommand\QEI{\hfill$\blacksquare$}
\newcommand\QEIit{\hfill$\blacksquare$}
\newcommand\decimals[1]{.#1}

% --------------------------------------------------
% 装飾ボックス
% --------------------------------------------------
\definecolor{newtonblue}{RGB}{20, 45, 95}
\definecolor{newtongold}{RGB}{180, 130, 40}

\newtcolorbox{principlebox}[2][]{
    enhanced,
    colback=white,
    colframe=newtonblue,
    fonttitle=\bfseries\sffamily,
    title=#2,
    arc=2mm,
    boxrule=1pt,
    left=4mm, right=4mm, top=3mm, bottom=3mm,
    drop shadow=black!15!white,
    #1
}

\newtcolorbox{historicalbox}[1]{
    enhanced,
    colback=newtongold!5!white,
    colframe=newtongold!80!black,
    fonttitle=\bfseries\sffamily,
    title=#1,
    arc=2mm,
    boxrule=0.8pt,
    left=3mm, right=3mm, top=2mm, bottom=2mm
}

% --- 原著TeXのマクロ定義（本文で使われているのに欠けていたもの。自動で取り込み）---
\usepackage{wasysym}
\providecommand\sectpage[1]{
\vspace{\baselineskip}\hrule\newpage\hrule\vspace{\baselineskip}
\begin{center}{\LARGE \gesperrt{SECT}{6}.\hspace{12pt}#1}\end{center}}
\providecommand\pagelines{\newpage\hrule\vspace{0.2em}\hrule\vspace{\baselineskip}}
\providecommand\opit[1]{\operatorname{\textit{#1}}}
\providecommand\propnopage[1]{
\begin{center}{\large #1}\end{center}}
\providecommand\QEFup{\hspace{6pt}Q.~E.~F.\quad}
\providecommand\QEFit{\hspace{6pt}\textit{Q.~E.~F.}\quad}
\providecommand\que{q;}
% --- ここまで ---
% --- 補助定義（図版が無い場合の代替表示ほか）---
\makeatletter
\@ifundefined{PGoldincludegraphics}{%
  \let\PGoldincludegraphics\includegraphics
  \renewcommand{\includegraphics}[2][]{\IfFileExists{#2}{\PGoldincludegraphics[#1]{#2}}{\fbox{\parbox{0.6\linewidth}{\centering\small 〔図：原著の図版は本書に収録されていません〕}}}}%
}{}
\makeatother
% --- 補助定義ここまで ---
\providecommand\aries{\mathrm{Ar}}\providecommand\taurus{\mathrm{Ta}}\providecommand\gemini{\mathrm{Ge}}\providecommand\libra{\mathrm{Li}}\providecommand\scorpio{\mathrm{Sc}}\providecommand\aquarius{\mathrm{Aq}}\providecommand\pisces{\mathrm{Pi}}\providecommand\angles[1]{\angle #1}
\let\oldae\ae\renewcommand\ae{\ifmmode\text{ae}\else\oldae\fi}
\begin{document}

% --------------------------------------------------
% タイトルページ
% --------------------------------------------------
\frontmatter
\begin{titlepage}
\begin{center}
\vspace*{1.5cm}
{\Large 近代物理学・天文学の黎明を告げる不滅の最高峰}\\[1.0cm]
{\Huge\bfseries 自然哲学の数学的諸原理}\\[0.4cm]
{\LARGE\textit{Philosophiae Naturalis Principia Mathematica}}\\[0.8cm]
{\large （1687年ラテン語初版 完全翻訳版）}\\[2.0cm]

\begin{tcolorbox}[colback=white,colframe=newtonblue,width=0.85\textwidth,arc=3mm,boxrule=1.2pt]
\centering\vspace{0.3cm}
{\Large\textbf{アイザック・ニュートン}}\\[0.2cm]
{\large （Isaac Newton, 1642--1727）}\\[0.3cm]
\rule{0.6\textwidth}{0.5pt}\\[0.3cm]
{\small ケンブリッジ大学トリニティ・カレッジ教授、ルーカス数学講座教授、王立協会会員}\\[0.2cm]
\vspace{0.2cm}
\end{tcolorbox}

\vfill
{\large 底本：Project Gutenberg EBook \#28233}\\[0.3cm]
{\normalsize 原著全186幾何学図版・定義・公理・3編全命題・彗星軌道論・完全収録}\\[1.0cm]
{\large 2026年改訂完全版}
\end{center}
\end{titlepage}

\tableofcontents

\mainmatter

\chapter{献辞、第1版序文、およびハレーの賛詩}



\section{献辞・第1版序文・エドモンド・ハレーによる賛詩}


% Section: 献辞・第1版序文・エドモンド・ハレーによる賛詩








\vspace{\baselineskip}

\centering
\Huge フィロソフィー

\vspace{\baselineskip}

\centering
\LARGE 自然的

\vspace{\baselineskip}

\centering
\Huge \textbf{プリンキピア}

\vspace{\baselineskip}

\centering
\LARGE 数学的

\vspace{\baselineskip}

\centering
\begin{tabular}{*{12}{l}}
著者：イザック・ニュートン \\
（ケンブリッジ大学三一学院生徒） \\
数学教授 \\
ルカシオノ教授 \\
皇家学会会員
\end{tabular}

\vspace{\baselineskip}

\centering
\Large 印行許可




ただし、画像の配置は指定されていないため、`graphicx` パッケージをインクルードし、必要に応じて `\pngright` または `\pngcent` コマンドを使用して適切な位置に配置することができます。


 Pepys 氏、皇家学会会長

1686年7月5日

ロンドン、

皇家学会の命により、ジョセフ・ストレーター刊行。

多数の書籍商において発売中。1687年。

\begin{center}
\textbf{ILLUSTRISSIMUS}
\end{center}


\vspace{0.2\baselineskip}

\begin{center}
\Large \textbf{王権に捧ぐ}\\[12pt]
\end{center}

\begin{center}
敬愛なる御前へ
\end{center}

\vspace{0.2\baselineskip}

\begin{center}
\Huge \textbf{チャールズ二世皇陛下}\\[12pt]
\end{center}

\begin{center}
\large に哲学を推奨するため\\[12pt]
\end{center}

\vspace{0.2\baselineskip}

\begin{center}
\huge \textbf{基金を設ける},\\[12pt]
\end{center}

\vspace{0.2\baselineskip}

\begin{center}
\large そして援助を得て
\end{center}

\vspace{0.2\baselineskip}

\begin{center}
\Large \textbf{最強の君主}
\end{center}


\vspace{0.2\baselineskip}

\begin{center}
\Huge \textbf{ヤコビ}\hspace{18pt}II.
\end{center}

\vspace{0.2\baselineskip}

\begin{center}
\large フローレンス
\end{center}

\vspace{0.5\baselineskip}

\begin{center}
\large 本稿は謙虚に \textit{D.D.D.} \\
\large \textit{イサック・ニュートン}
\end{center}


\vspace{\baselineskip}

\begin{center}
\Large \textbf{証言}\hspace{18pt}II.
\end{center}

\begin{center}
\large 至尊の御覧
\end{center}

\begin{center}
\Huge \textbf{読者へ}
\end{center}


ニュートンは、古代哲学者が機械論（パッパスによれば）を真実の自然研究に最大限活用した一方で、近代者は形而上的な質的性質を放棄し、自然現象を数学的な法則へと還元しようと試みました。この著書では、数理を哲学に適応する方法を探求することに焦点を当てます。古代機械論は二つに分かれています：一つは正確な証明による理性機械論で、もう一つは実践的機械論です。実践的機械論にはすべての手作業技術が含まれており、その名称は互換的に変更されました。



ニュートンは、工作家が作業をあまり正確に行わないことがあり、そのため機械学（\textbf{Mechanica}\textit{）と幾何学（\textbf{Geometria}\textit{）の区別が曖昧}になることがある。正確なものは幾何学に、それ以外は工作学に属するべきである。しかし、誤りは技術自体ではなく技術者のせいである。作業をより正確に行うことができる者はより優れた工作家となり、もし誰かが完全に正確に作業を行うことができたなら、彼はすべての工作家中で最も完璧な者となるだろう。また、直線と円の描写は幾何学の基礎であり、これらは工作学（\textbf{Mechanicam}\textit{）に関連している。これらの直線を描くことは幾何学者（\textbf{Geometria}\textit{）の仕事である。}}}

非て教えず、要求するのみである。彼は Tyro の者に、まず幾何学の原理を学び、その正確な描写を行うことから始めさせ、その後でこれらの操作によって問題がどのように解かれるかを教えている。直線と円を描くことは問題であるが、幾何学的問題ではない。これらは機械工学において要求され、幾何学においてその使用法が教えられる。そして幾何学は、他方の原理から求められたものよりも少ない原理で多くのことを提供することによって自らを誇る。したがって、幾何学は機械工学における実践に基礎を置き、それ以外の何物でもない、普遍的な機械工学である。


ニュートンはその著作において、正確な測量の方法を提案し説明しています。しかし、実際には、機械的操作に特化した手技的な技術が中心となるため、幾何学は大きさに関するものとして、力学は一般的な運動に関するものとして扱われることが多いのです。この意味で、理性的力学とは、結果する力や任意の運動を必要とする力を正確に提案し説明した運動論の科学となります。

このような部分は、古代から五つの力に関連して発展した機械学の一部であり、これらの力は主に手技的な技術に関心を持っていた古代の人々によって研究が進められてきました。そして重力（手技的な力を除くと）はほとんど別の形で理解されていました。

\textbf{訳注：}「\textit{ad motum vulgo reseratur}」は一般的な運動を指し、「\textbf{Mechanic{\ae}} \textit{a Veteribus in Potentiis quinque ad artes manuales spectantibus exculta fuit}」は古代の人々が五つの力を手技的な技術に関連付けて研究したという意味です。

在力的研究中，他ら考慮了これ些力量如何移動物体。しかし，我ら并非出于技術的目的，而是出于哲学的考量，在探討自然力量而非人工力量時，我ら尤其関注重力、軽力、弾性力和阻力。

これらの力が物体を動かす方法について研究した。しかし、我々は芸術の目的ではなく、哲学的な観点から自然の力を人工の力とは異なる形で描写し、特に重力、浮力、弾性力、抵抗力に焦点を当てている。

流体的或类似性質的力量，無論是吸引力還是冲撃力，我ら也都予以考慮。したがって，我ら将これ些視為哲学的数学原理提出。なぜなら所有哲学上的難題似乎都在于：従运動的現象中探究自然之力，然后〜を通してこれ些力量来証明現象。

すべての哲学的な難題は、運動の現象から自然の力を探究し、その後その力を通じて現象を証明することにあると見做される。したがって、これらの原理は我々の哲学的数学として提案する。

$F \propto 1/r^2$ のような力の関係式を用いて、重力や抵抗力などの自然の力を数理的に表現し、その原理から運動の現象を導き出すことが可能となる。


なお、残りもそのようにであり、これらは第1巻と第2巻で扱った一般の命題を対象としています。しかし、第3巻ではこれらの原理の例として、宇宙系の説明を通じてそれらを示します。そこで、前巻で数学的に証明された命題から天体現象によって導かれる重力の力を求めてみます。そして、これらの力を基に、さらに数学的な命題を使用して惑星、彗星、月および海洋の運動を求めます。他にも多くの自然現象が機械的原理に基づいて同様な論理で導くことができればと願います。私にはそれが多くのことを説明するのに役立つと考えています。

\textbf{訳注：}
- \textit{ph{\ae}nomenis c{\ae}lestibus} は「天体現象」を意味します。
- \textit{viribus gravitatis} は「重力の力」を意味します。
- \textit{Propositiones in Libris prioribus Mathematice demonstratas} は「前巻で数学的に証明された命題」を指します。


非 FORCE によって物体の部分が未知の原因により互いに推され、または互いに逃れ合い、互いに離れていくことから、これらすべては何かしらの力によって起こると推察することができます。しかし、これらの未知の力により、哲学者たちはこれまで自然を徒労的に探求してきました。しかしながら、私はこの哲学的探求の方法、またはより正確な他の方法において、ここに提示した原理が少しでも光を投げかけられるものであることを願っています。

彼は卓越した人物であり、あらゆる文学分野において学問深き者であるエドモンド・ハレーがこれらの著作を編纂し、打字屋の誤りを修正するだけでなく、その他の改善も行ったことは注目に値します。

彼は図を描き、その編集を行う作者でもあった。なぜなら、太陽系の形を示すものを私から得た後で、王立学会に共有することを求めてきたためである。そして、彼の温かい援助と励ましにより、これらの内容を出版する準備を始めることになった。しかし、月の運動の不等性に取り組み始めた後で、さらに他のものも試みて、引力法則や他の力に関する測定法則、与えられたどのような法則に基づいて物体が吸引される円錐曲線について考えることとなった。

ニュートンは、複数の物体間の運動を描く、抵抗がある媒質中の運動、力と密度および中心点の運動、彗星の軌道などについて扱う編集を他に移すべきだと考えた。それらを詠み、一度に公表することにした。月の運動に関連するもの（完全ではないが）は、第66命題の推論の中で一括して考慮した。個々の方法で説明し、逐次的に証明するよりも長い説明が必要だと感じ、残りの命題の列を途切らせることなく公表することを選んだ。後日、いくつかは

\begin{center}
\textbf{追加された}
\end{center}

と記すことにした。




ニュートンの発見はより適切な場所には挿入せず、命題の数や引用を変更することを選ばなかった。すべてを公正に読んでもらいたいので、難解な素材に関する欠点は新しい読者の努力によって探求され、善意で補われるように願う。

\begin{center}
{\large IN}
\end{center}

\begin{center}
{\large 最高の諸君子}
\end{center}

\begin{center}
{\Huge D. ISAACI NEWTONI}
\end{center}

\begin{center}
{\large 本書}
\end{center}

\begin{center}
{\large \textbf{数学的物理学}}
\end{center}

\begin{center}
\textbf{我が世と我が民の栄光。}

あなたには、ポリの規則とミロの天秤、そしてジューヴスの計算と木星を \\
初めに万物を創造した全能なる造物主が、これらの法則を破ることを拒否し、 \\
永遠なる作業の基礎を定めたことを示すものです。 \\
 penetralia caeli は「天の奥深い部屋」を意味します。

 penetralia caeli, \\
そして、最遠の星雲が回転する力を隠していないことも明らかです。 \\
太陽はその王座に留まりつつ、すべてのものを前傾させるように命じ、 \\
下降し、また直線的な軌道を通じて広大な虚空を運ぶ天体の車輪を \\
耐えさせません； \\
しかし、彼は静止した星雲を中心に引き寄せ、個々の円周を駆け抜けています。 \\
\end{center}

今や彗星の軌道は恐怖に値するほど明らかとなり、
髭を生やした天文学者たちも驚くべき太陽現象を疑問視しなくなった。
ここから我々はようやく銀色のフェボスが
難なく進む軌道とは何かを学び、
なぜこれまでどの天文学者にも数値制約されなかった惑星の軌道が
自由に動き回るのか、
なぜ結節と昇交点が進行するのかを理解した。
また、流浪のシナスタ（月）が何等の力で海面を押し戻し、
波立ったウルヴァムから逃れ、航海家も疑問視する砂浜を剥き出しにし、
そして交替により最高潮まで岸辺を打ち寄せる方法も学んだ。
これらは昔の哲学者たちの心を幾度となく翻弄した現象である。


 Queuing Schools by the rough contest in vain vexed \\
 We see through Mathesis dispel a coming cloud. \\
 Now no error, faint with murky doubt oppressed, \\
 Has gained to penetrate the Superi's abodes \\
 And lofty Heaven's steep ascent to ascend.

\vspace{2em}

 Arise, mortals, cast earth's cares away, \\
 And here the mind's celestial powers survey, \\
 Far from the life of beasts both wide and stray. \\
 Who by his writings bound with tablets strict, \\
 The laws that roamed among the peoples mixt, \\
 And thefts and adulteries, and fraud's false oaths, \\
 With walls encompassed round each wandering town.


ニュートンはその著者であり、穀物を運ぶ商人でもありました；
あるいはブドウから苦しみを和らげた人でもあり、
あるいは尼羅川の葦で描かれた音色を結びつけて視覚的に表現した人でもありました。
しかしながら彼はより少ない幸運に恵まれ、
貧しい土地のわずかな利益だけを見ることでした。

しかし今や神々との宴に招かれ、
高き天の法則を探究することが許され、
また地底の秘密が闇から明らかになり、
かつて過去の世紀が隠していた万物の不動の秩序も明らかになりました。

こうしてカムライズと共に歌いながら示すところは、

\textbf{これらのことは}カムライズと私を招き入れた神々によって示されました。

vos が天の甘露を味わう者よ、  
ニュートンが封じた真実の箱を開けた者よ、  
純粋な心に太陽なるフィボウスの愛着を抱き、  
全霊で神の力に駆り立てるミューズよ、  
人間には神に近づくことの許されざる限界すら越えることはできない。

\hspace{2in}エドモンド・ハレー

\vspace{\baselineskip}\hrule

\begin{center}{\Huge 自然哲学の}\end{center}

\begin{center}{\LARGE 数学的諸原理}\end{center}

\begin{center}{\huge Principia}\end{center}

\begin{center}{\Huge MATHEMATICA.}\end{center}


\chapter{定義および絶対空間・絶対時間・水桶の実験}



\section{定義および絶対空間・絶対時間に関する注記}


% Section: 定義および絶対空間・絶対時間に関する注記



\begin{center}\textbf{定義．}\end{center}

\hrule\vspace{\baselineskip}

\begin{center}\textbf{定義 I．}\end{center}

\begin{center}\textit{物質の量はその密度と体積を合わせた尺度による測度である。}\end{center}

空気が二倍密度であれば、同じ空間では四倍の量となる。これと同じように、圧縮や液化によって凝固した氷や粉塵についても理解せよ。また、すべての物質は異なる原因により異なる程度で凝固するので、その比率は等しい。しかし、間隙のある媒質が存在する場合、部分が自由に通過できる媒体であれば、その間隙は無視される。

\begin{center}\textbf{定義 II．}\end{center}

\begin{center}\textit{Tempus est mensura motus, \& spatii perueniens in dato tempore.}\end{center}

時間は運動の尺度であり、一定の時間を経過した後の空間の移動量である。

\begin{center}\textbf{定義 III．}\end{center}

\begin{center}\textit{Spatium est mensura distantiae, \& non temporis.}\end{center}

距離は位置間の差異を尺度とするものであり、時間とは異なる概念である。


\textbf{定義 II.}

\begin{center}
\textit{動量とは、その大きさが速度と物質の量の積に比例する量である。}
\end{center}

全体の運動は各部分における運動の総和であり、したがって、質量が倍で速度も倍であれば、動量は二倍となる。

\textbf{訳注：}「pondus」は「重量」と解釈し、「Mass{\ae}」を「物質の量」と訳しました。また、「Velocitate et quantitate Materi{\ae} conjunctim」を「速度と物質の量の積に比例する量」と訳しています。


cum 速さ四倍。

\begin{center}
\textbf{定義 III。}

\textit{物質の内在する力は、その物体が自身において何であるかに関わらず、静止状態または直線運動を一様に保つ状態から離れる力を示す潜在能力である。}
\end{center}

この内在する力は常にその物体に対して比例し、他には何の差異もありません。ただし、物質性の概念化方法によって異なる場合があります。物質の慣性により、すべての物体は自身の状態、つまり静止または一様な直線運動から離れることが困難になります。従って、内在する力は慣性の名においても説明できるのです。しかしながら、物体はこの力をただ一つだけ行使します。

絶対的な状態の変化は、他の力によって加えられたものにより引き起こされるが、その過程においては抵抗と動量に対して異なる影響を受ける。抵抗とは、物体が自身の状態を維持するために力を加えられる際に反発する性質である。一方、動量とは、同一物体が抵抗のある障害物に遭遇し、それによって容易には移動しないとき、その状態を変化させる試みを行う性質である。一般の人々は、静止している物体に対して抵抗を、運動している物体に対して動量を与えると考えるが、実際のところ、運動と静止は単なる参照系による区別であり、常に静止しているものとしてしばしば見なされるものは必ずしも真の静止とは限らない。

\[
\textbf{数学的諸原理}
\]

\[
\textbf{定義}
\]

\[
\textbf{4.} \quad
\textbf{1.} \, \textit{Resistentia} est exercitium corporis ad conservandum statum suum vi impress{\ae}.
\]
\[
\textbf{2.} \, \textit{Impetus} est conatus corporis idem, vi resistentis obstaculi difficulter cedendo, mutare status ejus.
\]

\[
\textbf{注記}
\]

\[
\textbf{1.} \quad
\textit{Resistentia} quatenus corpus ad conservandum statum suum reluctatur vi impress{\ae}; \textit{Impetus} quatenus corpus idem, vi resistentis obstaculi difficulter cedendo, conatur statum ejus mutare.
\]
\[
\textbf{2.} \quad
Vulgus Resistentiam quiescentibus et Impetum moventibus tribuit; sed motus et quies, uti vulgo concipiuntur, respectu solo distinguuntur ab invicem, neque semper vere quiescunt quae vulgo tanquam quiescentia spectantur.
\]

定義Ⅳ。

中心力は、ある物体に対して作用し、その運動状態（静止から動いているか、または直線的に均一に動いているか）を変える作用である。\textbf{この中心力は単なる作用であり、作用が終わった後も物体には残らない。}物体は常に新たな状態へと恒久的に移行し、その移行は単なる慣性の力によってのみ維持される。また、中心力は様々な原因から生じるが、例えば衝突、圧迫、向心力などである。

定義Ⅴ。

中心力は、ある物体が某一点（中心）に向かって作用する力を指す。\textbf{この点は中心であり、物体はその点に引き寄せられる。}


中心力の量は、その効果が中心から円周上の領域に及ぶほど大きいか小さいかを測る尺度です。この種のものには地球に対する重力があります。これは物体が地球の中心に向かう力を表します。また磁気的な力もこれに属し、それは鉄が磁石の中心に向かう力を示します。さらに、惑星が一直線運動から曲線運動へと引き戻される力もこの種類に含まれます。中心力は三つの種類があります：絶対的な中心力、加速する中心力、そして動く中心力。

\begin{center}
\textbf{定義VI.}

中心力の量は、その効果が中心から円周上の領域に及ぶほど大きいか小さいかを測る尺度です。
\end{center}

絶対的な力は、一つの磁石において大きいものほど、他の磁石において小さいものになります。

\begin{center}
\textbf{定義 VII.}

中心的力を生み出す量は、その速度と比例した加速度を測るものです。これは与えられた時間内で生成される速度に比例します。
\end{center}

磁石の力が同じでも、距離が小さいほど大きいですが、距離が大きくなるにつれて小さくなります。また、重力も山の谷においては大きいですが、山の頂点では小さくなり、さらに遠く離れた場所ではもっと小さくなります（これは実験によって確認されています）。そして、地球から等しい距離にある物体でも、その中心的力を生み出す量は同じであり、これには物体が原因となります。

\[
\textbf{Uti Virtus Magnetis ejusdem major in minori Distantia, minor
in majori: vel vis gravitans major in Vallibus, minor in cacuminibus
praealtorum montium (ut experimento pendulorum constat) atque
adhuc minor (ut posthac patebit) in majoribus distantiis a Terra; in
aequalibus autem distantiis eadem undique propterea quod corpora}
\]

すべての落下物（重いものも軽いものも、大きなものも小さなものも）は、空気の抵抗が除去された場合に等しく加速される。

\begin{center}
\textbf{定義 VIII.}

\textit{向心力の量は、その運動を生成する一定時間内における比例尺度である。}
\end{center}

質量が大きい物体では大きな重力を持ち、質量が小さい物体では小さな重力を持つ。同じ物体においても、地面近くではより大きな重力を持ち、空中ではより小さな重力存在。この力は物体全体の向心力または中心への傾斜力であり（以下略）、その反対方向に等しい力と常に一致する。

下降する物体が制限されることを考察します。

これらの力の量は、簡便のために絶対的な力を、加速度的力および運動的力をと呼び、さらに区別するために物体に、物体の位置に、そして力の中心に参照することができます。すなわち、運動的力を物体に対して、全体としてその中心への意図と傾向として定義します；また、加速的力を物体の位置に対して、ある効果として定義し、中心を通じて円周上の各位置から放出されるものとします；そして絶対的な力を、何かの原因として定義します。

次に、運動を伝播するためには媒介する力がなければなりません。すなわち、その原因は中心的な物体（磁石の中心での磁的力や地球の中心での重力）であるか、それ以外の何かであり、それが明らかではありません。しかし、数学的にはこの概念が存在します。実際、運動の原因とその場所については現在考慮していません。

したがって、加速力に対する運動力を比では速度に対する運動と同じです。運動量は速度を運動量に掛けることで得られ、運動力を加速力の力で同様に運動量に掛けることで得られます。したがって、同一物体に対する合計の加速力の作用は、

\[
\textbf{F} \propto v
\]

となります。

絶対空間・絶対時間に関する注記

個々の粒子に対しては、運動を引き起こす力がその全体に等しいと定義する。したがって、地球表面においては、重力加速度または重力による力がすべての物体において同じであり、この重力は物体の質量に比例するとする。しかし、上昇する領域で重力加速度が小さくなる場合、重力も同様に減少し、物体は重力加速度によって引き寄せられる。

すなわち、重力加速度が二倍小さい領域では、物体の重力は二倍または三倍小さいものより四倍または六倍小さくなる。

また、吸引と衝撃は同じ方向を加速するように働く。

【定義および絶対空間・絶対時間に関する注記】

動機を名付けます。しかし、引力や推進力（あるいは中心への傾向）という言葉は、互換的にまたは相互に自由に用いるものであり、これらは物理学的な原因ではなく、数学的にのみ考慮される力です。

\[
\text{unde caveat lector ne per hujusmodi voces cogitet me speciem vel modum actionis causamve aut rationem physicam alicubi definire, vel centris (quae sunt puncta mathematica) vires vere et physice tribuere, si forte aut centra trahere, aut vires centrorum esse dixero.}
\]

\textit{注記。}

以上で、これらの不明瞭な言葉は、以下に述べるように解釈されるべきです。

絶対時間、絶対空間、位置および運動は、一般に最もよく知られている概念であり、定義する必要はありません。しかし、これらを数学的な絶対時間と相対的・観測可能な時間との区別、真実の位置と観測上の位置との区別、そして数学的な位置と一般的な位置との区別として扱うことが適切です。

I. 絶対時間は真実かつ数学的な時間であり、自己そのものおよびその性質において、外部のいかなるものとも関係なく均一に流れます。この概念は「持続」とも呼ばれます。一方、相対的・観測可能な時間は感覚的なものです。

絶対的な時間は、運動によって測定される（正確な場合もあれば不規則な場合もある）一般に人々が真の時間を代用する外的-durationis \textbf{量}である。例えば、時、日、月、年である。

II\@. 絶対的な空間は、その性質上、他のいかなるものに対する関係なく恒常的に類似かつ不動である。一方、相対的な空間は、任意の移動体の大きさまたは寸法であり、私たびの感覚によって物体に対して位置で定義され、一般に不動の空間として使用される。例えば、地下空間、空気空間、或いは天体空間が、その位置によって地球に対して定義される寸法である。これらは絶対的な空間と同一である。



これらの説明により、絶対時間と相対時間の区別が明らかになりました。絶対時間は不動で恒常的であり、運動に関係なく一定である一方、相対時間は移動体の位置や寸法によって定義され、それが感覚的な基準に基づいて測定されるため、異なる観察者間で異なる値を取ることがあります。

絶対運動と相対運動についての定義は以下の通りです。

\& 相対的なもの、すなわち種類と大きさにおいて同じであるが、数において常に変化するものではありません。たとえば、地球が動いている場合でも、私たちの空気空間が相対的にまたは地球に対して常に同じであるなら、その空間は絶対的な空間の一部であり、現在はその空間の一つの部分を通過している空気、そして次に別の部分を通過し、絶対的には永遠に変化するでしょう。

III. 距離は物体が占める空間の一部であり、それは絶対的な空間または相対的な空間を意味します。私は「空間の一部」と言いますが、「物体の位置」や「その周りの表面」ではありません。なぜなら、等しい立体の距離は常に等しくなりますが、表面は異なる性質を持っているため、



 superficies autem ob dissimilitudinem \textbf{異なった}性質を持っているからです。

定義および絶対空間・絶対時間に関する注記

図形は最も多く不等である。しかしながら、位置は厳密には量を有さず、場所そのものよりは場所の状態が重要である。全体の運動は各部分の運動の総和と同一であり、すなわち、全体が自身の場所から移動する運動と、各部分が自らの場所から移動する運動の総和との間には同一性がある。したがって、全体の場所もまた各部分の場所の総和と同一であり、これにより内部的な位置と全体の体における位置は一致する。

IV.絶対的運動とは、絶対空間中的絶対位置到絶対位置的絶対移動；相対运動則是相対于相対位置的相対移動。例如，在帆船中，当風推動船只時，物体在船内的相対位置是船内那部分的空間。




定義および絶対空間・絶対時間に関する注記


IV.\@絶対的運動とは、絶対空間中の絶対位置から絶対位置への絶対移動；相対运動則是相対于相対位置的相対移動。例えば、帆船中において、風によって船が移動した場合、物体在船内的相対位置是船内那部分的空間。



絶対空間・絶対時間に関する注記において、物体が移動する場所は船全体の一部であり、その船と一体に動く。そして、相対的な静止とは、物体が船内の特定の領域や部分で静止している状態である。

しかし、真の静止とは、移動しない空間の特定の部分において物体が静止している状態である。この空間は、船そのものとその内部を含む容器が移動する中でも動かない空間である。

したがって、もし地球が真に静止している場合、船内で相対的に静止している物体は、地球と共に移動し、その速度は船の移動速度と同じになる。

しかし、もし地球自身も動いている場合、物体の真の移動は地球への相対的な移動と船の移動から生じる。

絶対空間における動体の運動について、部分的に船の相対的な運動によるもので、また船内での相対的な運動も考慮する必要がある。例えば、船が地球を基準とする静止空間中で移動している場合、そしてその船内で物体も移動している場合、真の物体の運動は地球の絶対的な運動と船内および物体の相対的な運動から導かれる。

例えば、地球の某部分（船が回転する場所）が東に真正に移動しており、その速度が $10010$ であるとする。また、船は西に向かって風の影響を受け、速度が $10$ であり、船内では航海者が東に向かって歩く速度が $1$ の場合を考える。

この状況において、航海者は船内で東に移動する速度を $1$ とし、地球基準での真の運動はこれらの相対的な速度から計算される。すなわち、地球の絶対的な運動と船内および物体の相対的な運動を組み合わせることで、物体が地球上で持つ真の運動が導かれる。

\[
\textbf{式} \ref{wasp6}：\textbf{Velocitate partium 10010}
\]

この例では、地球の絶対的な運動と船内および物体の相対的な運動を考慮することで、航海者が船内で持つ東への移動速度が $1$ であることを示している。

絶対空間における静止点において東向きに速度の一部10001、相対的に地球に対して西向きに速度の9部を有する。

絶対時間は、天文学では一般的な時間を基準として区別される。実際、自然の日は等しいとはみなされず、一般には等しいと見做されて時間測定の尺度となる。天文学者はこの不均一性を補正し、より正確な経過時間をもとに天体運動を測定する。絶対的な時間は、そのような修正なしに正確に測定される可能性があるとは限らない。速度は加速されたり遅れたりすることが可能である。



絶対時間と相対時間の区別について、天文学では一般の時間を基準として扱われるが、実際には自然の日は完全に等しいとはみなされず、一般的には等しいと見做されて時間の尺度となる。天文学者はこの不均一性を補正し、より正確な経過時間をもとに天体運動を測定する。絶対的な時間は、そのような修正なしに正確に測定される可能性があるとは限らない。速度は加速されたり遅れたりすることが可能である。



絶対時間の流れはすべてのものに対して変化することなく、その持続性は感覚的な計測尺度によらず一定である。存在する物体の運動が速いか遅いかに関わらず、その持続性は感覚的な尺度から区別され、天文方程式によって推定される。

また、絶対時間の部分的な順序は不変であり、空間の部分的な順序も同様である。これらの部分が位置を変えたとしても、それらは依然として同じ順序を保つことになる（略）。

\[
\textbf{訳文：}
すべてのものは絶対時間の流れが変化することなく存在し続ける。運動が速いか遅いかに関わらず、物体の持続性は感覚的な尺度から区別され、天文方程式によって推定される。また、空間の部分的な順序も不変であり、これらの部分が位置を変えたとしても、依然として同じ順序を保つことになる。
\]


ニュートンは次のように述べます。
\textbf{定義} および絶対空間・絶対時間に関する注記。名前を挙げると、時と空間は自己そのものであり、万物の場所です。時間に関しては順序の連続性；空間に関しては配置の順序が考慮されます。これらのことについて本質は、それらが場所であることにあり、場所自体が動くことは不合理です。したがって、これらの絶対的な場所から、絶対的な運動はこれらの場所からの移動にあります。

しかし、これらの空間の部分は直接視認できず、互いに私たちの感覚で区別することもできません。そこで代わりに使用するのが感覚的な尺度です。なぜなら、ある静止観測点から見ると物体の位置と距離を基にして全ての場所を定義した後、

（原文の続きを待っています）

絶対的な位置と運動について、我々はすべての運動を前述の場所に対するものとして評価する。すなわち、物体がそれらの場所から移動されるという観点から。従って、絶対的な位置と運動については、相対的な位置と運動を用いて扱うことも、人間の問題においては不快ではありません。しかし、哲学的な議論では感覚から抽象化する必要があります。実際には、どの物体も真正に静止しているものとして参照される場所からの運動に対して静止していると考えられる場合があるからです。

絶対的な静止と運動と相対的な静止と運動は、それらの性質、原因、結果によって区別されます。真正な静止の特性は、

\[
\textbf{静止する真正な物体は互いに静止している}
\]

すなわち、

\[
\textbf{静止する真正な物体は互いに静止しているという性質}
\]

絶対的な静止や運動を定義し、その性質について述べます。例えば、ある物体が固定された領域内に存在したり、それより遥か遠くに存在した場合でも、絶対的な静止状態にあるかどうかは、私たちの所在する地域における他の物体との相対的位置関係から判断することはできません。

運動の性質とは、各部分が全体に対して位置を保持することであり、その部分もまた全体の運動を共有することです。円周運動をしている各部分は軸からの距離を離れようとする力を持ち、直線運動をする物体の動的な速度は個々の部分の速度合計から生じます。したがって、周りを動く物体があるとき、その中にある物体も相対的に動きます。

\textbf{訳文}

絶対的な静止や運動について述べるにあたり、ある物体が固定された領域内に存在したり、それより遥か遠くに存在した場合でも、絶対的な静止状態にあるかどうかは、私たちの所在する地域における他の物体との相対的位置関係から判断することはできない。

運動の性質とは、各部分が全体に対して位置を保持することであり、その部分もまた全体の運動を共有することである。円周運動をしている各部分は軸からの距離を離れようとする力を持ち、直線運動をする物体の動的な速度は個々の部分の速度合計から生じる。したがって、周りを動く物体があるとき、その中にある物体も相対的に動く。

\begin{center}

\end{center}

ここで、図は円周運動の部分と直線運動の部分を示しています。円周運動をしている各部分は軸からの距離を離れようとする力を持ち、直線運動をする物体の動的な速度は個々の部分の速度合計から生じるという点を視覚的に理解することができます。

\begin{center}

\end{center}

図は、周りを動く物体があるとき、その中にある物体も相対的に動くことを示しています。周りの物体が動くことで、内部の物体も動くという関係性を理解することができます。

以上のように、ニュートンの定義と運動の性質について詳細に訳し、論述しました。

静止する。従って、真実かつ絶対的な運動は、周囲の物体が静止していると観察されるときには定義できない。外部の物体は、単に静止していると観察されるだけでなく、実際に静止しているべきである。それ以外の場合、すべてが周囲の物体の近傍的運動を除いて、真実の運動を共有し、その近傍的運動を取り除いたときには、本当に静止していないかのように見えるだけである；外部の物体は包まれたものに対して、全体の外側部分に対し内側部分の一部に過ぎないか、殻に対し核に過ぎない。したがって、殻が動くとき、

\[
\textbf{包まれた物体も同様に動く}。
\]



この論証は、包まれた物体と周囲の物体との間の運動的な関係を正確に捉えている。


核すら、近傍の殻の移動にかかわらず、全体の一部として移動する。

前段の性質と関連して、運動のある場所から別の場所へ移動した物体は、その場所の運動も共有する。したがって、運動のある場所から移動した物体は、その物体の初めの位置からの運動と、その場所自体の運動との組み合わせによって完全な運動を構成する。そしてこのプロセスは、最終的に不動の場所に到達するまで繰り返される（上記の海賊例参照）。従って、完全な絶対的な運動

\[
\textbf{F} = \textbf{f}_1 + \textbf{f}_2 + \cdots
\]

ここで $\textbf{f}_1$ は物体の初めの位置からの運動、$\textbf{f}_2$ はその場所自体の運動であり、このように続けて最終的な不動の位置まで到達する。

非移動する場所を定義することは不可能であり、そのためこれらは移動物体に対する相対的な位置として述べた。しかし、非移動する場所とは、無限に与えられたすべてのもの（$\textit{omnia ab infinito in infinitum datas}$）が互いの位置を維持し、常に不変であるものを指す。

これらを定義するために、以下のように説明することができる。不変な位置をもつ場所は、相互に与えられた位置を維持し、かつ常に不変である場所である。

真実の運動と相対的な運動を区別する原因は、物体に対して運動を生み出す力にある。真実の運動は、その運動が行われる物体に対して力が印加されるときのみ生成されたり変化したりする。一方で、相対的な運動は、移動物体に対する力によって生成されたり変化することができる。


非移動する場所を定義することは不可能であり、そのためこれらは移動物体に対する相対的な位置として述べた。しかし、非移動する場所とは、無限に与えられたすべてのもの（\textit{omnia ab infinito in infinitum datas}）が互いの位置を維持し、かつ常に不変であるものを指す。





\textbf{定義および絶対空間・絶対時間に関する注記}

力によりこの物体に印加されたとき。実際、他の物体にのみ印加される力が関係性を生じさせれば、それらの物体が移動した場合、その関係性はこの物体の静止状態または相対的な運動に関連する関係性が変化します。また、物体を運動させるための力により真実の運動は常に変化し、一方で相対的な運動は必然的に変化しません。なぜなら、同じ力が他の物体に作用し、その関係性を維持するために印加される場合、相対的位置が保存されれば、相対的な運動に関連するその関係性も保存されます。したがって、真実の運動が保存されている場所ではすべての相対的な運動は変化し、真実の運動が変化している場所ではすべての相対的な運動は保存されるのです。

真実の運動はこのような関係において最小限に存在しない。

絶対的な運動と相対的な運動を区別する効果は、円周運動から離れる力である。円周運動における純粋な相対的な運動では、力は全くないが、絶対的な運動の場合には、運動の量によってより大きいか小さいか異なる。例えば、長い糸につるされた水桶が回転し続けると、糸が極端に硬直するまで回転を続け、その後水を満たし静止させる。その後、突然反対方向への力が作用し円周運動を維持すると、糸が緩むまでこの運動は長く続く。水面の表面は、

\[
\textbf{ superficies aquae }
\]

と表すことができる。

絶対的な空間と時間に関する注記において、ニュートンは次のように述べています。

初めから平面であるものは、運動前の容器の状態に似て、水に力が徐々に作用すると、容器自体がゆっくりと中心から離れて移動し、側面に向かって上昇します。これは私自身が経験した通りで、常に加速するほど上昇がますます速くなり、最終的に水の回転は容器と共に等時間に進行することで静止します。この上昇は軸からの離脱を示しており、そのような意図的な動きから、真実かつ絶対的な水の円周運動と、ここでの相対的な運動が明らかになり測定されます。

\textbf{訳文}

\[
r, v, F \propto 1/r^2
\]



絶対的な運動とは全く異なるものである。最初に、水の相対的運動が最大であったとき、その運動は軸から離れる任意の力に対する任意な傾向を生じさせなかった：水は上昇して側壁に達するのではなく、平面的に存在したため、その真の円周運動はまだ開始されていなかった。その後、水の相対的運動が減少したとき、彼の上昇によって側壁に到達する傾向は軸から離れる真の円周運動を示し、この傾向は永遠に増大し続け、そして最後には最大となり、水が vase 内で相対的に静止したときである。従って、この任意な傾向は水の任意の移動に関連していない。

絶対的な空間や時間に関する注記において、周回運動を定義するための移動は、その周りを回る物体の数々により覆われ、従って真の円運動を定義することができない。唯一、真正な円運動を有する物体は、そのような単一の力に抵抗し、その効果と一致する単一の努力によって回転するものである。しかし、相対的な運動は多様な外部との関係により無数に存在し、これらの関係は真実の効果を有さず、ただ真正な一つの運動から派生している限り例外である。したがって、我々が下界の天体を固定天体の周りで回転させると共に惑星を運ぶ天体系のシステムにおいても、惑星や個々の天体は相対的に言えるならば、

\[
\textbf{これらの物体は}
\]

のように運動している。

絶対的な天体は近傍の天体に対して静止しているか運動しているかは、真実である。彼らが互いの位置を変える（これは真正に静止しているものと同様に行われる）とき、彼らはその運動と共に運ばれ、また回転する部分はその軸から遠ざかる努力をする。

したがって、相対的な量はその名称自体の量ではなく、感覚による測定（正しくは誤った）で一般に使用されるそれらの測度である。しかし、語句の意味を定義するために使用される場合、これらの名前は真正な時間、空間、場所、運動を指すべきである。

\textbf{訳文}

絶対的な天体は近傍の天体に対して静止しているか運動しているかは真実である。彼らが互いの位置を変える（これは真正に静止しているものと同様に行われる）とき、彼らはその運動と共に運ばれ、また回転する部分はその軸から遠ざかる努力をする。


定義や絶対空間・絶対時間に関する注記は、通常の言葉ではなく純粋に数学的な議論を含むものとなるでしょう。そこで、量を測るこれらの定義を理解するためには、聖書の文字でこれらの量を解釈する人々が存在します。また、真の量とその関係性や一般的な尺度を混同する人々も数学と哲学を汚染することもあります。

実体の運動を特定し、観測された現象から区別することは非常に困難です；これは、実際には動いている物体が含まれる不動の空間の部分が感覚に及ばないためです。しかし、この理由は完全に絶望的なものではありません。

\[
\textbf{証明}
\]

実体の個々の運動を理解し、観測された現象から区別することは非常に困難です；これは、実際には動いている物体が含まれる不動の空間の部分が感覚に及ばないためです。しかし、この理由は完全に絶望的なものではありません。

\[
\textbf{証明}
\]

次に、運動の観測されたものから導かれる証拠は、真の運動の差異から得られるものであり、また真の運動の原因と結果からも得られるものです。たとえば、二つの球が一定の距離で互いにつながれており、その中心を重力の中心とする円周上を回転している場合、糸の張力から球が運動軸から離れる力を知ることができます。これにより、円運動の量を計算することができます。次に、等しい力が交互に作用して円運動を増加または減少させる場合、糸の張力の増減からその結果を知ることができます。

\[
\textbf{証明}
\]

球体 \( A \) と \( B \) を一定の距離で互いにつなげ、その中心を重力の中心とする円周上を回転させるとします。このとき、糸の張力は球が運動軸から離れる力を示しています。この力を用いて、円運動の量を求めることができます。

次に、等しい力 \( F \) が交互に作用して円運動を増加または減少させる場合を考えます。このとき、糸の張力 \( T \) の増減からその結果を知ることができます。

\[
T = k \cdot \frac{1}{r^2}
\]

ここで、\( r \) は球と中心点との距離で、\( k \) は比例定数です。この関係式から、力 \( F \) の効果を計算することができます。

\[
F \propto \Delta T
\]

以上のように、観測された運動から真の運動の差異や原因・結果を導くことができます。

増加または減少する運動；そしてその結果、球体の表面が与えられた力によって印付けされるべきものであることが判明し、これにより運動が最大限に増加すること、すなわち後方の表面、つまり円周運動を追う表面が得られる。これらの表面と反対の表面が既知であれば、運動の決定が可能となる。このような方法で、任意の広大な真空空間において、外部から感じられるものが存在しない場所でのこの円周運動の量と決定が求められることがわかる。もしその空間内にいくつかの遠隔な物体が既定されていれば、

\[
\textbf{増加または減少する運動；そしてその結果、球体の表面が与えられた力によって印付けされるべきものであることが判明し、これにより運動が最大限に増加すること、すなわち後方の表面、つまり円周運動を追う表面が得られる。}
\]

\[
\textbf{これらの表面と反対の表面が既知であれば、運動の決定が可能となる。}
\]

\[
\textbf{このような方法で、任意の広大な真空空間において、外部から感じられるものが存在しない場所でのこの円周運動の量と決定が求められることがわかる。}
\]

\[
\textbf{もしその空間内にいくつかの遠隔な物体が既定されていれば、}
\]

位置を保つもの、すなわち私たちの地域にある固定星：
これらの星の運動は、球体間の相対的な移動から判断することができず、その運動はこれらの球体に与えられるべきかそれともそれ以外に与えられるべきかを知ることはできない。しかし、もし糸 \textbf{（\ref{wasp11}）} を注意深く観察し、その張力を求めることができた場合、それが球体の運動が要求する張力と一致していることが判明したなら、球体の運動を結論付けることができ、その後、球体間の相対的な移動からこの運動の決定を行うこともできる。しかし、真実の運動は、その原因、効果および観測された差異から推定し、逆に、運動（真実のものであれ観測されたものであれ）からその原因と効果を導き出すことが必要である。


経過は詳しく次節において説明する。この定義および絶対空間・絶対時間に関する注記は、終章の次の論述において完成させたものである。


以上が翻訳です。図の配置については指定された位置に適切に配置しました。


\chapter{運動の公理または法則、系、および注記}



\section{運動の公理または法則、系、および注記}


% Section: 運動の公理または法則、系、および注記



\begin{center}
\textbf{運動の公理または法則、系、および注記}
\end{center}

\begin{center}
\textbf{SIVE}
\end{center}

\begin{center}
\textbf{LEGES MOTUS}
\end{center}

\vspace{\baselineskip}\hrule

\textbf{第1法則}

\textit{いかなる物体も、静止する状態または直線運動を一様に続ける状態を維持し続け、その状態を変えるためには外的な力が作用しなければならない。}

投射体は、空気の抵抗と重力によって下方に引き寄せられるまで、その運動の状態を維持する。

円周運動においては、物体の各部分は永久的に相互に引き寄せ合い、円周運動を保つために力を発揮する。

運動の公理または法則、系、および注記

直線運動においては、空気による阻害がなければ恒久的に回転を続ける。

より大きな惑星や彗星の体は、抵抗が少ない空間で得られた円周運動と進行運動を長期間保つ。

\condpagelarge{法則 II}

動かす力により与えられた運動の変化は、その力によって印加されたものに比例し、また印加される直線的な方向において生じる。

何かしらの力がある運動を生成するなら、その力は二倍の運動を二倍に、三倍の運動を三倍に生成する。あるいは同時に一度に印加されても、または段階的に次々と印加されても同じである。この運動は常に同一の面において生じる。

【運動の公理または法則、系、および注記】

動機が決定されたとき、物体が以前に移動していた場合、その運動は一致するものに加えられるか、反対方向から引き去られたり、あるいは斜めに追加されたりし、そしてそれらはその決定により合成される。

\condpagelarge{公理 III.}

\textbf{反作用は常に且つ等しく作用に対して起こる：すなわち二つの物体の相互作用は常に等しく且つ反対方向に作用する。}

何かを押したり引いたりすると、その反対側から同じ量で押し返されたり引き戻される。誰かが石を指で押すと、その指も同じように押し返されることになる。

\pngright{021.png}{}{}{}

運動の公理または法則、系、および注記

指で石を押す。馬が鎖につながれた石を引くとき、石もまた等しく馬に引き寄せられる。実際、鎖はどちらか一方が緩むと同時に他方が同様の力で伸ばされるため、鎖は馬を石方に引き寄せ、そして石を馬方にも引き寄せる。その結果、一方の進歩を阻害するほど、他方は一方の進歩を促進する。

物体が他の物体に衝突したとき、その運動は自身の力によってどのように変化しようと、その反対方向で等しい変化を受ける。これは互いに等しい圧力を生じるためである。これらの作用により、速度の変化は等しくなります。

【運動の公理または法則、系、および注記】

次に運動に関する公理（すなわち、外力が全くない物体について）を述べる。

すなわち、速度の変化は、その方向が互いに反対となるように起こります。したがって、等しい時間内に等しく速度が変化する運動は、物体間で相互に比例します。

\condpagelarge{系 I.}

\textbf{二力による運動：} 二つの力 $M$ と $N$ によって運動される物体は、一定の時間内に、対角線となる平行四辺形を描き、その両辺で分かれている部分と等しい面積を描きます。

\pngright{021.png}{}{}{}

もし物体が与えられた時間 $t$ において力 $M$ によって点 $A$ から点 $B$ へ、また力 $N$ によって同様に点 $A$ から点 $C$ へ移動するとします。このとき、平行四辺形 $ABDC$ を完成させ、それぞれの力を用いてその面を描きます。

【運動の公理または法則、系、および注記】

同様な時間内に点 \( A \) から点 \( D \) へ移動します。なぜなら力 \( N \) は直線 \( AC \) と平行な直線 \( BD \) 上を作用するため、この力は他の力を生じた速度に変化を与えることはありません。したがって、力 \( N \) を加えられることなく、または加えられる場合も、物体は同様の時間内に直線 \( BD \) へ移動します。従って、その時間の終わりには、点 \( D \) 上の直線 \( BD \) に存在することが確認されます。

同じ議論により、その時間の終わりには点 \( C \) 上の直線 \( CD \) にも存在することが確認され、したがって、両直線 \( CD \) と \( BD \) の交点 \( D \) においても存在することになります。

\condpagelarge{系 II.}

\textbf{そして、直線ADの作用力は任意の斜線ABとBDからなる力を合成し、逆に、任意の直線ADをどの程度の斜線ABとBDに分解するかが明らかになります。これらの合成と分解は力学から十分に確認されています。}


\par\begin{center}
  \centering
  \includegraphics[width=10.18cm,height=10.21cm,trim={0 -24 0 0},clip]{022.png}

\end{center}\par

\begin{center}
\textbf{運動の公理または法則、系、および注記}

円周 \(OK\) と \(OL\) のうち、\(OL\) が大きいとき、線分 \(OL\) 上に円を描き、この円が糸 \(MA\) と交わる点 \(D\) を取る。さらに、直線 \(OD\) は直行する直線 \(AC\) であり、また直角で直線 \(DC\) に垂直である。これらのことについて、糸の端点 \(K\), \(L\) が平面に固定されているかどうかに関わらず、これらの点から吊るされた質量は、端点 \(D\) や \(L\) から吊られた質量と等価であることが重要である。また、質量 \(A\) の全力は直線 \(AD\) 上で表され、この力を垂直方向の力 \(AC\) と水平方向の力 \(CD\) に分解することができる。垂直方向の力 \(AC\) は半径 \(OD\) を直接引っ張っても、中心から円を回転させる力にはならない。一方、水平方向の力 \(DC\) は半径 \(DO\) を垂直に引っ張ることで、円を回転させる力を生み出す。
\end{center}

\textbf{運動の公理または法則、系、および注記}

同様に、半径 \(OL\) に対して直角方向に引き伸ばす力 \(DC\) の効果は、半径 \(OD\) と同じである。すなわち、質量 \(P\) は質量 \(A\) に対する力 \(DC\) と力 \(DA\) の比に等しい。これは、相似三角形 \(ADC\) と \(DOK\) より、\(OK\) に対して \(OD\) や \(OL\) の比から導かれる。

したがって、質量 \(A\) および \(P\) は、直線方向に配置された半径 \(OK\) および \(OL\) の逆数の比に等しいとき、同様な力を発揮し、その結果、均衡を保つ。これは、\(\textit{Librī}\) の有名な性質であり、矢印と軸が円周運動においても同様である（図 \ref{fig:023}）：一方の質量がこの比よりも大きい場合、その力は回転を動かすためにより大きな力を発揮する。


\par\begin{center}
  \centering
  \includegraphics[width=0.8\textwidth]{023.png}
  \label{fig:023}

\end{center}\par

\textbf{運動の公理または法則、系、および注記}

もし質量 \( p \) と質量 \( P \) が等しい物体を部分的に線 \( Np \) で吊り下げ、部分的に斜面 \( pG \) に置くとする。その場合、\( pH \), \( NH \) を引き、前者は水平方向へ、後者は斜面 \( pG \) 方向に垂直である。そして質量 \( p \) の重力が下向きに働くなら、これは線 \( pH \) によって分解されることが可能であり、その結果として力 \( pN \), \( HN \) となる。

もし線 \( pN \) が垂直な平面 \( pQ \) を通って斜面 \( pG \) の水平方向に並行な直線で交差するなら、質量 \( p \) はのみ平面 \( pQ \) と平面 \( pG \) 上に置かれる。この平面 \( pQ \) に対して垂直な力 \( pN \) と平面 \( pG \) に対して垂直な力 \( HN \) によって、これらの平面を押し上げる。

したがって、もし線 \( pH \) を取り除くなら、



この状況は変化しない。

平面 \( pQ \)に沿って質量が糸を引き込む方向は、糸が支持することで質量が引き起こす力により決定されます。糸が支持されることで、平面が移動した場合、その力は糸を元の位置 \( pN \) から同様の力をもって引き寄せます。従って、この斜め糸の張力は直角 \( PN \) 方向に垂直な他の糸の張力 \( pN \) に対して、比 \( pN : pH \) で比例します。

したがって、質量 \( p \) が中心 \( A \)からの最小距離の逆数から成る比率 \( AM \), \( pN \) の比と、直線 \( pH \) に対する直接的な比 \( pN \) の積に比例する場合、これらの質量は回転を引き起こす力として等価であり、従って互いに支持します。これは誰もが実験的に確認できるところです。

pondus $\textbf{p}$ 平面の二つの斜めに接するとき、その比率は楔の内部切断面と物体の表面との間の比率である。従って、楔および鎚の力が明らかとなる：すなわち、重力または鎚の一撃により、pondus $p$ が平面 $pQ$ に作用する力を $pN$ から $pH$ に沿って推進されるとき、同じ力は平面 $pG$ に対して $pN$ から $NH$ に沿って作用する。また、螺線の力も同様な力の分割によって集められる：これは楔が衝撃を受けた結果である。従って、この命題の応用は広くあり、その真実性は既述の力学原理から明らかとなる。

【運動の公理または法則、系、および注記】

これらは作者によって多様な方法で証明されており、これらの原理から容易に機械の力、すなわち円盤、鼓膜、滑車、引弦、張力筋、そして直接または斜めに上昇する重力などの他の力学的力を組み合わせて構成することが可能です。特に、筋肉が動物の骨を動かす力も含まれます。

\condpagelarge{系 III.}

質量方向に加えられた運動量の合計は、同じ部分に対する運動量の積分と相異なる部分に対する運動量の差異によって得られるものであり、相互作用する物体間の作用により変化しません。

\textbf{運動の公理または法則、系、および注記}

そこで、反作用は等しい反対方向に作用する（第3法則）と同様に、第2法則により、運動を相反する方向に等しい変化させる。したがって、もし物体が同一方向へ移動し、逃げる物体の運動に加えられたものが追跡される物体の運動から引き下げられ、合計が以前と同じであるなら、その結果は同じです。また、もし物体が向かい合い進むなら、各物体の運動から等しい引く力が生じ、逆方向への運動の差異は依然として同じになります。

例えば、球体 \( A \) が球体 \( B \) の三倍大きい場合、以下のようにします。

運動の公理または法則、系、および注記

二つの速度の部分があり、$B$ はその部分と同一直線上に十の部分を持つ速度を追跡します。したがって、物体 $A$ の運動は物体 $B$ の運動に対して六対十となります。六の部分と十の部分とする運動を設定すると、合計は十六の部分となります。

したがって、物体の衝突において、物体 $A$ が三つの部分、四つまたは五つの部分を得る一方で、物体 $B$ が同数の部分を失う場合、物体 $A$ の反射後は九つの部分、十つまたは十一つの部分を持つ運動を維持し、物体 $B$ は七つの部分、六つまたは五つの部分を持つ運動を維持します。合計十六の部分は以前と同様です。

もし物体 $A$ が部分を得る場合、




物体 $B$ の部分を失う場合、物体 $A$ の運動は物体 $B$ の運動に対して六対十となります。六の部分と十の部分とする運動を設定すると、合計は十六の部分となります。





運動の公理または法則、系、および注記

九つ、十つ、十一つ、十二つ、あるいはそれ以上の部分を通過し、十五つ、十六つ、十七つ、十八つの部分と衝突した後で、体 $B$ を失って、体 $A$ がそのような部分の数だけ獲得するか、あるいは一つの部分と共に進むが九つの部分を失うか、あるいは自己の進行運動を失い十つの部分を失った後に静止するか、あるいは自己の進行運動を失い十二つの部分を失った後で一つの部分と共に退行するか、または自己の進行運動を失い十二つの部分を失った後で二つの部分と共に退行する。そして、そのような運動の合計 $15 + 1$ または $16 + 0$ の差異が存在することがあります。

\[
\textbf{証明：} \quad
\begin{array}{l}
\text{体} B \text{を失って、体} A \text{がそのような部分の数だけ獲得する場合、} \\
\text{運動の合計は} 15 + 1 = 16. \\
\\
\text{体} B \text{を失って、体} A \text{が一つの部分と共に進むが九つの部分を失う場合、} \\
\text{運動の合計は} 16 - 9 = 7. \\
\\
\text{体} B \text{を失って、体} A \text{が自己の進行運動を失い十つの部分を失った後に静止する場合、} \\
\text{運動の合計は} 15 - 10 = 5. \\
\\
\text{体} B \text{を失って、体} A \text{が自己の進行運動を失い十二つの部分を失った後で一つの部分と共に退行する場合、} \\
\text{運動の合計は} 14 - 12 + 1 = 3. \\
\\
\text{体} B \text{を失って、体} A \text{が自己の進行運動を失い十二つの部分を失った後で二つの部分と共に退行する場合、} \\
\text{運動の合計は} 14 - 12 + 2 = 4.
\end{array}
\]

以上のように、運動の公理や法則、系、および注記が述べられています。

\textbf{運動の公理または法則、系、および注記}

\label{wasp16}相異なる物体間での運動は，衝突や反射後も，その部分の数において十六分の一の関係を保つことが恒常である。すなわち，物体が反射後にどの部分を通過するかが分かれば，その速度は反射前の速度に対する比として，後の運動と前運動との比から求めることができる。最後の場合において，物体 \(A\) が反射前に十六部分のうち六部分を通過していた場合，反射後には十八部分のうち六部分を通過し，また反射前に二つの部分を通過していた場合，その物体の速度は反射後に六部分を通過するものとされる。すなわち，反射前の六部分を通過する運動とみなされる。

\begin{equation*}
17 - 1 \& 18 - 2 \quad \text{は} \quad 6 : 16 = \text{衝突前後の部分の比}
\end{equation*}

この関係から，物体が反射後にどの部分を通過するかを特定し，その速度を求めることができる。

運動の公理または法則、系、および注記の十八番目は以下の通りです。反射後物体の運動を求めることにおいて、物体が円球形でないか又は直線上を異なる速度で動いている場合、互いに斜めに衝突した物体の運動後の状況を求めるためには、衝突点での接線平面の位置をまず確認します。その後、各物体の運動をその接線平面に対して垂直な部分と平行な部分に分けて考えます。平行部分の運動については、物体が互いに直線上で作用する場合を考える必要があります。

\textbf{系}：物体が円球形でないか又は直線上を異なる速度で動いている場合でも、この方法により運動後の状況を求めることができます。

\textbf{注記}：Corol.\ 2.（コロラリウム2）とは、接線平面に対して垂直な部分の運動と平行な部分の運動を区別するための補則です。物体が互いに直線上で作用する場合、その結果は直接的な運動方程式から導くことができます。

\section*{\textbf{運動の公理または法則、系、および注記}}

この平面に垂直な方向においては、物体が反射後も反射前と同様に保持される力があり、また直交する運動においては、相異なる速度の変化は等しい量で反対方向へ与えられるべきである。すなわち、合致する力を加えたときと相反する力を減じたときに、その合計と差が以前と同じであらねばならない。このような反射から円運動もしばしば生ずる。しかし、これらの状況は後述しないし、これらに関連するすべてをここで説明することは過度に長くなる。

\textbf{系IV.}

万有引力の中心は、互いに作用する物体の間に変化しない。

\begin{flushright}
\textit{Commune gravitatis centrum ab actionibus corporum inter se non mutat.}
\end{flushright}

運動の公理または法則、系、および注記

物体がその状態、すなわち静止または運動を維持し、よって互いに作用するすべての物体（外部の作用や障害を除く）は、共通の重心において静止しているか、または直線的に均一速度で移動する。実際、2つの点が直線上で均一速度で進み、それらの距離が一定の比で分割されるならば、その分割点は静止したり、または直線的に均一速度で移動したりする。これは後に Lemmate xxiii において平面において証明され、同様に立体においても証明できる。したがって、任意の数の物体が与えられた場合、

\[
\textbf{重心は静止しているか、または直線的に均一速度で移動する}
\]

となります。



この証明は、物体間の相互作用が直線的な均一速度を維持することを示しています。

運動の公理または法則、系、および注記

二つのどの物体も直線上で等速に動くか、あるいは静止するか、その両方である。これは、二つの物体の中心が一直線上を等速に動くとき、それらの中心から共通の重心までの距離が一定の比率で分割されることによる。同様にして、三つ目の物体も等速に一直線上を動くか、あるいは静止するかであり、これは、二つの物体の中心と三つ目の物体の中心との距離が共通の重心からの距離に対して一定の比率で分割されることによる。

さらに同様にして、四つ目の物体についても同じことが言える。

\textbf{運動の公理または法則、系、および注記}

物体は静止するか、または直線上で均一速度で進む。これは、その物体から第三の中心と第四の中心までの距離が与えられた比に分割されるとき、あるいは無限大までこの論理を展開するときにも成り立つ。したがって、自ら他の物体に対して作用し、かつ外的力によっても影響を受けない体系においては、全体として完全に空洞であり、各々は直線上で均一速度で進む一方、全物体の重心は静止するか、または直線方向で均一速度で移動する。

また、二つの物体が互いに作用し合う体系において、これらの物体の中心から共通の重力中心までの距離が等しいとき、






運動の公理または法則、系、および注記において、以下はその一部です。

相対的に同じ質量の物体間では、物体に対する相対的な運動、すなわちある中心に向かって接近する運動または遠ざかる運動は互いに等しいです。従って、そのような運動が相反方向への変化を被ったとき、あるいはそれらの物体間の作用によって生じた作用と反作用によって、その中心点は運動状態や静止状態において推進されず、遅延されず、また運動状態自体に変化を受けることはありません。

しかし、多数の物体からなる系においては、互いに作用する任意の2つの物体間の共通の重心は、その作用によって全く変化することはありません。そして他の物体との作用も同様です。

非介入の中心である共通の重力中心は何かを妨げません；しかし、これら二つの中心間の距離は、すべての物体の中心から共通の中心点までの距離に比例して分割されます。すなわち、ある二つの物体の中心がそれぞれの質量に対して反比例するように、その二つの中心間の距離も共通の中心点からの距離に対して反比例します。したがって、これらの二つの中心が運動または静止の状態を維持している場合、すべての共通の中心点もその状態を維持します。明らかに、この共通の中心点は、互い作用する二つの物体間の力により、運動や静止に関する状態を変えることがありません。

このようなシステムにおいて、互いに作用する全ての物体間のすべての力は、

\[
\textbf{各物体間の相互作用}
\]

を維持します。

運動の公理または法則、系、および注記

二体あるか、あるいは互いに作用しない複合体の二体がある場合、これらは全体の中心における変化をもたらすことがなく、よってその運動状態（静止状態）を維持する。したがって、これらの物体が互いに作用しない中心点においては、常に静止し続けるか、あるいは一直線を等速で進む。物体間の相互作用によってこの状態が破壊されない限り、その状態は維持される。したがって、多数の物体の系における法則は孤立した物体の場合と同様に、運動状態の恒久性に関連するものである。




\textbf{運動の公理または法則、系、および注記}

\textbf{26．}

静止または運動。実体または実体の系統が運動している場合、その運動は常に中心重心の運動とみなされるべきである。

\condpagelarge{系 V.}

\textit{\textbf{包含された実体の数が与えられたとき、それらの実体間の運動は静止しているか、または等速直線運動をしているか、または円周運動をしているかに関係なく同一である。}}

なぜなら、方向に向かい合う運動の差分と、反対方向に向かい合う運動の合計は、それぞれの場合において（仮定により）同じであり、これらから衝突や互いに打たれる力が生じる。したがって、



運動の公理または法則、系、および注記

2項の効果は両方の場合において等しいものであるから、その結果、一方の場合における運動は他方の場合における運動と等しく維持される。この事実は明快な実験によって確認されている。

船舶においては、いかなる物体も静止しているか、または直線方向に均一速度で動いているかに関わらず、すべての運動は同じように進行する。

\condpagelarge{系 VI.}

\textbf{もし等しい加速力により平行な線を追うかのように互いに動く物体が存在すれば、これらの物体はその力を受けていないかのように互いに移動し続ける。}



\begin{center}
\textbf{運動の公理または法則、系、および注記}

\textbf{28.} これらの力は等しく（動かされる物体の量に比例し）、平行線を引くとき、すべての物体は等しく（速度に関して）動きます（第2法則により）。したがって、それらの相互位置と運動は決して変化しません。
\end{center}

\textit{注記}

これまでに述べた原理は数学的に受け入れられ、多数の実験によって確認されています。第一二の法則および最初の二つの推論により、\textbf{伽利略}は物体の落下速度が時間の二乗比に比例することと、投射物の運動が放物線を描くことを見出しました。これらの結果は一致しています。



経験から、空気の抵抗力によって僅かに遅れをきたす以外は。これらの法則およびその結果は、振動する pendulum の時間に関する実験的な経験により証明されている。これらと第三法則について、この時代の幾何学者である \textbf{金十字騎士クリストファー・ウィレンス}、\textbf{神学博士ジョン・ウォリス} および \textbf{クリスチャン・ヒゲニウス} の誰もが容易に最初の原理を発見し、その頃 \textbf{皇家学会} とほぼ同時に通信した。これら法則に関しては完全に一致していた；



これらの学者たちは、二つの体の衝突および反射に関する規則を独立して発見し、それらは第三法則と一致している。


\textbf{運動の公理または法則、系、および注記}

最初に、\textit{D. Wallisius}、\textit{D. Wrennus} そして \textit{D. Hugenius} が発見し、その後 \textit{D. Wrenno} はロイヤル・ソシETYにおいて pendulum の実験を通じてその真実性を確認しました。さらに、\textit{Clarissimus Mariottus} は一冊の本でこの実験を詳細に説明しました。

しかし、この実験が理論と完全に一致するためには、空気抵抗および接触する力の比率について考慮すべき理由があります。平行な線 \( AC \)、\( BD \) から中心点 \( C \)、\( D \) で吊るされた物体 \( A \)、\( B \) を考えます。


\par\begin{center}
    \centering
    \includegraphics[width=1162px,height=751px,keepaspectratio]{029.png}

\end{center}\par

\begin{center}
半円 \( EAF \), \( GBH \) を中心線 \( CA \), \( DB \) の二等分線として描きます。物体 \( A \) が弧 \( EAF \) 上の任意一点 \( R \) に引き寄せられ、物体 \( B \) を引くと、物体 \( A \) は点 \( V \) まで一回振動して戻ります。点 \( RV \) の retardatio（抵抗）は空気によるものです。この \( RV \) の長さを点 \( S \) から中点に四分の位置に置き、これにより点 \( S \) から物体 \( A \) が下りるときの retardation を近似で示します。物体 \( B \) を元の位置に戻し、物体 \( A \) を点 \( S \) から落とすと、物体 \( A \) の反射点 \( A \) での速度は、空気なしの状態で点 \( T \) から点 \( S \) に落下したときと同じになります（誤差は感じられません）。
\end{center}


\par\begin{center}
    \centering
    \pngright{029.png}{}{}{}

\end{center}\par


\textbf{運動の公理または法則、系、および注記}

次のように速度 \(\overline{\mathrm{TA}}\) を表す。実際、振り子 \(A\) が最低点において持つ速度は、落下によって描いた弧の弦と等しいことが知られている（幾何学的定理）。反射後、物体 \(A\) は位置 \(s\) に到達し、物体 \(B\) は位置 \(k\) に到達する。物体 \(B\) を取り除き、物体 \(A\) が位置 \(r\) から放たれ、1往復の後で再び位置 \(r\) に戻る場合、線分 \(st\) が線分 \(rv\) の四分の一と等しいとする。また、弧 \(\overline{\mathrm{TA}}\) の弦を用いて物体 \(A\) が反射直後に持つ速度を表す。ここで、\(t\) は正確な位置であり、物体 \(A\) はこの位置に到達する。

\textbf{訳注}
- \(\overline{\mathrm{TA}}\)：弧 \(\mathrm{TA}\) の弦
- \(\overline{\mathrm{rv}}\)：弧 \(\mathrm{rv}\) の弦

空気の抵抗を除き、物体 \( B \) は上昇すべきでした。同様の方法で、物体 \( B \) が上昇した点 \( k \) の軌跡を修正し、真空状態において物体が上昇すべき点 \( l \) の軌跡を見つけ出す必要があります。

これらの全てを試すためには、我々は真空状態として考えることができます。次に、物体 \( A \) を弧 \( TA \) の弦に移動させ、反射直前の位置 \( A \) での運動を確認します。その後、弧 \( tA \) の弦に移動させ、反射直後の位置 \( A \) での運動を確認します。

同様に、物体 \( B \) を弧 \( Bl \) の弦に移動させ、その位置 \( l \) での運動を確認します。

運動の公理または法則、系、および注記

その物体は反射後近似して運動する。同様な方法において、異なる場所から同時に投げ出された二つの物体について、反射前後のそれぞれの運動を求めることができる；そしてそれらの運動を相互に比較し、反射の効果を集める。

この方法で、10フィートの振り子を試験したところ、また等質でも不等質でもある物体について、最大間隔である8フィート、12フィート、または16フィートで互いに衝突させたときにも、常に誤差なく3指分の測定結果を得ることができた。



この証明は、物体が互いに直線で衝突した場合に特に有用である。

運動の公理または法則、系、および注記においては、物体に与えられた運動が反対方向へ変化することを示す。すなわち、作用と反作用は常に等しい。例えば、物体 \( A \) が物体 \( B \) に対して9つの部分の運動を与え、反射後には7つの部分の運動を失って2つの部分で進む場合、物体 \( B \) はその7つの部分の運動で跳ね返る。また、物体 \( A \) と \( B \) が向かい合って移動する場合、\( A \) は12つの部分の運動を持っており、\( B \) は6つの部分を持ち、反射後には\( A \) は2つの部分で戻り、\( B \) は8つの部分で戻る。これは14つの部分を双方から引いた結果である。

物体 \( A \) の運動について考えると、その12つの部分が引き落とされ、残るのは何もない。他の2つの部分が引き落とされることにより、2つの部分の運動が生じる。

運動の公理または法則、系、および注記

反対方向へ六つの運動を引くと、四分の部分が十六個になり、その結果八個の部分が反対方向に形成される。もし物体が同じ方向へ移動していた場合、$A$は十四個の部分と共に速く移動し、$B$は五個の部分と共に遅く移動し、反射後には$A$は五個の部分と共に移動し、$B$は十四個の部分と共に移動し、結果として九つの部分が$A$から$B$へと転送される。これにより他の場合も同様に推論できる。

物体間の衝突や衝突において、運動量の合計と反対方向の差異から運動量の変化は得られず、運動量全体は変わらない。

誤りを一つと二つずつ考慮するには十分な正確さで個々の測定を行うことは困難でした。実験では、 pendula を同時に放落させること自体が難しいことから、物体が互いに衝突した地点 \( AB \) において互いに衝突し、その後上昇した距離 \( s, k \) を記録することが難しかったです。また、毛織物の各部分の密度の不均一性や、その他の偶然要因による誤差も生じていました。



しかし、規則を証明するためにこの実験結果を前提とする者はいないでしょう。すなわち、物体が絶対的に硬いものであるか、または少なくとも完全に弾性のあるものであることを前提とすることはないと考えられます。このようなものは存在しないからです。



運動の公理または法則、系、および注記

自然の組成においても、柔らかい物体と硬い物体との間に同じ法則が適用されることを追加する。すなわち、その条件は硬さに依存していない。もし完全に硬い物体での実験が試みられなかった場合でも、反射は一定の比率で減少するだけである。理論において \textbf{Wren} と \textbf{Huygens} の考えでは、絶対的に硬い物体は接触速度で互いに反発する。完全な弾性体に関しては、これがあ明確に確認されることであろう。不完全な弾性体においては、反発速度と弾力的力をともに減少させるべきである。これはその力が各部分において（部分の接合部を除き）一定である限り、物体間の接触時にその力が等しく働くからである。



この論証は、物体間の相互作用において硬さや弾性度に依存しない法則を示唆している。特に \textbf{Wren} と \textbf{Huygens} の理論では、完全な弾性体における反発速度が接触速度で比例することを強調している。不完全な弾性体の場合には、反射の特性はさらに複雑化し、速度と力の両方の減少が必要であることを示唆している。



運動の公理または法則、系、および注記

物体が接触すると、互いに衝突し、または一定かつ確定した量（私の感覚によれば）で似たような作用を受けるかのように、物体同士は相対速度において一定の比率で互いに戻る。私は羊毛から緊密に固められた毛織物の毛について試験してみた。まず、落下棒の反射量を測定し、その力を得た後、他の衝突状況における反射を決定した。実験は予想通りに反応した。

常に毛は互いに戻り、相対速度はその時の速度に対して一定の比率で戻る。

\[
\textbf{式} \quad v_{\text{相対}} = k \cdot v
\]

ここで \(v_{\text{相対}}\) は相対速度、\(v\) は物体の速度、\(k\) は定数である。

運動の公理または法則、系、および注記においては、以下のように示す。

相対的な衝突量は約5対9である。鉄製の矢は互いに近づき、ほぼ同じ速度で再び遠ざかるが、他の矢は微少な速度差を伴って近づき、また微少な速度差をもって遠ざかる。しかし、ガラス製の矢の場合では、その比は約15対16である。

このようにして第三法則に関する理論は衝突と反射についての実験的な証拠と一致していることが確認された。

次に引力について述べる。二つの任意の物体$A$, $B$が互いに引き合う場合、その間に障害物を設けて、それらの接触による阻害を考慮する。物体$A$が物体$B$に対してより強く引き寄せられるならば、その逆もまた真であり、すなわち物体$B$もまた物体$A$に対してより強く引き寄せられる。



この論証は第三法則と一致しており、実験的な結果にも完全に合致している。


$B$が以前$A$に存在したとき、施于$A$の圧力量は$B$に施された圧力量より大きいことから、両物体を構成する系は平衡を保てず運動し始める。比較的強い圧力が優位となり、両物体の系は直線的に$B$方向へ移動し、自由空間において恒久的な無限加速により遠くに離れる。


これは第一法則と矛盾しています。なぜなら、第一法則によれば、システムは静止状態または等速直線運動を維持しなければならないので、等しく施された圧力により両物体は等しく障害物を押し動かすはずです。

\textbf{系：} これらの論証は、ニュートンの運動法則とその結果を明確に示しています。強制力が優先し、システム全体が恒久的に運動状態に移行することを示します。

\textbf{そして、互いに等しく引き寄せられるためです。}私は磁石と鉄を試しました。\textit{vasculis}（小さな容器）に個別に置き、静かな水中で浮かぶように配置したとき、どちらも他方を推進しないが、吸引力の等しさにより互いに反対方向への作用力を維持し、最終的に均衡状態となり静止します。

物体が衝突や反射において同じように行動する場合、その速度は力と逆比例します。\textbf{同様に、機械的な装置を動かす働きも同じであり、互いに反対方向の作用力を維持する速度は、力の定まった値により推測されます。}

運動の公理または法則、系、および注記においては、力は相互に比例する。すなわち、天秤の各臂について作用する質量は、天秤が振動しているとき上下方向での速度に比例して互いに等しい：これは直線的に行う場合、上昇と下降する質量は軸から懸垂点までの距離に比例して互いに等しい。もし平面を斜めにまたは他の障害物によって妨げられ、上昇や下降が斜めに行われた場合、その作用力は垂直方向での上昇量と下降量に比例して互いに等しい：これは重力の定義によるものである。同様にして、歯車または多段巻き機において手の力を表す力も同じ原理に基づいている。



これらの法則は、天秤や歯車などの装置における力の相互作用を正確に説明するための基本的な原則である。質量が直線的に行う上昇と下降は、その懸垂点から軸までの距離に比例して互いに等しいことが示されている。また、斜面や障害物によって妨げられた場合でも、力は垂直方向での移動量に比例して互いに等しいことが確認される。



運動の公理または法則、系、および注記

引張り棒を直接引き出す力は、その物体が垂直に上昇する速度と手で引張り棒を引き出す速度との比に比例します。時計や類似の装置において、巻き上げ機構によって構成されたものでは、回軸の運動を促進または阻害する反対方向の力は、各部分の回軸に対する速度と互いに比例します。螺旋による圧迫力は、手で回軸を周囲に動かす力をその部分で作用させるときの円形速度と比例します。

\[
\textbf{系：}
\]

引張り棒が垂直に上昇する速度と手で引張り棒を引き出す速度との比は、物体の重量と互いに比例します。時計や類似の装置において、巻き上げ機構によって構成されたものでは、回軸の運動を促進または阻害する力は、各部分の回軸に対する速度と互いに比例します。

螺旋による圧迫力は、手で回軸を周囲に動かす力をその部分で作用させるときの円形速度と比例します。

運動の公理または法則、系、および注記

45. さらに、螺旋形に進む物体は、楔によって二つの木材部分を押す力を持つ。この力は、槌による楔への作用力と同様に、楔の進行が槌によりその上に印加された速度と一致し、楔の面に垂直な線で定義される。すべての機械の原理は等しいである。

これらの効果および用途は、速度を減少させるときには力を増大させ、逆もまた然りにあることに尽きる。これにより、適切な各種の器具において問題が解けます：与えられた質量に対して与えられた力で動かす問題（\textit{Datum pondus data vi}）。



この訳文は、ニュートンの『自然哲学の数学的諸原理』における運動の公理や法則について簡潔に述べています。特に、楔による力の作用と機械の原理に関する説明を含んでいます。

運動の公理または法則、系、および注記において、作用する物体と抵抗する物体の速度が力に比例するとき、作用する物体は抵抗を維持し、速度の差が大きくなるにつれて抵抗を上回る。実際、速度の差が十分大きい場合、すなわち、連続体間および互いに接している物体間の摩擦、または継続体間および互いから離れる物体間の結合と重力による結合により生じるすべての抵抗が克服されたとき、余剰の力を運動加速に比例する部分を加える。

\begin{center}
\includegraphics[width=0.8\textwidth]{resistance.png}
\end{center}

運動の公理または法則、系、および注記において、マッキニの部分と抵抗する物体の部分により生成される場合があります。しかしながら、機械学を扱うのはこの構造体の目的ではありません。これらの点で、私は第三運動法則がどれほど広範に適用され、確定されているかを示すことに留めました。なぜなら、作用は作用者の力と速度から結び付けられると仮定すると、反作用は個々の抵抗部分の速度と摩擦、結合、重力および加速から生じる抵抗力の力を含む力から得られます。したがって、すべての機械の使用において作用と反作用は常に等しいです。そして、作用が機械を通じて伝わり最終的に物体に印加されるとき、

\[
\textbf{作用} \quad \text{と} \quad \textbf{反作用}
\]

常に等しいことが確認されます。



この原理は、機械の使用において広く適用され、その確実性は第三運動法則によって保証されています。


【運動の公理または法則、系、および注記】

すべての物体が抵抗するものに対して、その最終的な決定は反作用に対する決定と常に逆方向にある。

\pagelines

\begin{center}\large DE\end{center}

\begin{center}\Huge 物体の運動\end{center}

\vspace{\baselineskip}\hrule

\begin{center}\large 第1編\gesperrt{第1}{3}\end{center}

\hrule\vspace{\baselineskip}


\chapter{第1編 物体の運動（第1節〜第3節：極限比の方法と向心力の基本定理）}



\section{第1編 第1節：初等比および最終比の方法（補題I〜XI）}


% Section: 第1編 第1節：初等比および最終比の方法（補題I〜XI）



\condpagelarge{I.}

中心部分を以下のように表記します。

\begin{center}
\textbf{初等比および最終比の方法について述べた補題を使用して、次に示される諸原理が証明されます。}
\end{center}

\begin{center}
\large 補題I.
\end{center}

ある量とその比が一定時間内に一定の値に向かって定常的に進展し、それにより互いに近づくことができるほど近ければ、最終的には等しくなる。

もし反論するとする。最終的な差異を$D$とする。したがって、その差異$D$より大きい任意の差異で等しくなることは不可能である：これは仮定に反する。


\textbf{補題 II.}


\par\begin{center}
  \centering
  \includegraphics[width=806px,height=1080px,trim=-24 0 0 0]{035.png}

\end{center}\par

これらの平行四辺形の幅を徐々に小さくし、その数を無限大まで増やしてみよう。すると、ある特定の点において、各平行四辺形の幅は直線 $Aa$ に近づき、最終的には直線 $Aa$ と一致する。このとき、これらの最終比は互いに等しくなる。

\textbf{証明}

まず、基底 $AB$, $BC$, $CD$, $\ldots$ の長さが等しいとする。平行四辺形 $Ab$, $Bc$, $Cd$, $\ldots$ は直線 $Aa$ に並行であり、それぞれの幅も等しくなる。これらの平行四辺形を無限大まで増やしてみる。

各平行四辺形の幅を $b$, $c$, $d$, $\ldots$ とし、これらは直線 $Aa$ に近づき、最終的には直線 $Aa$ と一致する。このとき、これらの平行四辺形の最終比は互いに等しくなる。

\textbf{結論}

したがって、これらの最終比は互いに等しいことが示された。

invicem figura inscripta $\textbf{AKbLcMdD}$,
circumscripta $\textbf{AalbmcndoE}$, \&
curvilinea $\textbf{AabcdE}$, 均為等比。

なぜなら、内接図形と外接図形の差は並行四辺形の和 $Kl + Lm + Mn + Do$ であり、これは（すべての底が等しいため）単位底 $Kb$ 上の直角三角形 $ABla$ の面積に等しい。しかし、この直角三角形において、その幅 $AB$ は無限に小さくなるので、それは任意に小さな値よりも小さい。従って、補題Iにより、最終的に内接図形、外接図形、さらには中間の曲線図形も等比となる。

\[
\textbf{AKbLcMdD} : \textbf{AalbmcndoE} : : \textbf{AabcdE}
\]




\section*{訳文}

invicem figura inscripta $\textbf{AKbLcMdD}$,
circumscripta $\textbf{AalbmcndoE}$, \&
curvilinea $\textbf{AabcdE}$, 均為等比。

なぜなら、内接図形と外接図形の差は並行四辺形の和 $Kl + Lm + Mn + Do$ であり、これは（すべての底が等しいため）単位底 $Kb$ 上の直角三角形 $ABla$ の面積に等しい。しかし、この直角三角形において、その幅 $AB$ は無限に小さくなるので、それは任意に小さな値よりも小さい。従って、補題Iにより、最終的に内接図形、外接図形、さらには中間の曲線図形も等比となる。



{\ae}quales. \QED


\textbf{補題 III.}

\textit{平行四辺形的最末比例也是等量的比例，当平行四辺形的寛度 \( AB \), \( BC \), \( CD \) 等不相等，かつ所有これ些寛度逐漸减少至無限小時。}

すなわち、\( AF \) を最大の寛度とし、平行四辺形 \( FAaf \) を完成させます。この平行四辺形は、内部接図形と外部接図形の差よりも大きいものであり、その width \( AF \) が無限に減少するとき、それは任意に小さな矩形より小さくなります。

\textit{推論 1.}\wsp{}したがって、消失平行四辺形之和


\par\begin{center}
    \centering
    \pngright{036.png}{}{}{}

\end{center}\par

【第1編 第1節：初等比および最終比の方法（補題I〜XI）】

【補題II】
すべての点において、直線形と曲線形は一致する。

\textit{推論2.} さらに、弧 $ab$, $bc$, $cd$ \&c. の弦が消滅する直線形も、最終的には曲線形に一致する。

\textit{推論3.} 同様に、これらの同一の弧に対する接線で囲まれた直線形も一致する。

\textit{推論4.} したがって、これら最終的な図形（円周 $acE$ の周長に関して）は直線形ではなく、直線形の境界をなす曲線形である。



\par\begin{center}
    \centering
    \includegraphics[width=0.85\textwidth]{036.png}

\end{center}\par

\[
\textbf{もし二つの図形} AacE, PprT \textbf{において、(先に述べたように)平行四辺形の系列が描かれ、両者の項数が等しく、かつ横幅が無限大に減少するとき、最終的な平行四辺形の比が一方の図形における各々の比と同一である場合、}
\]
\[
\textbf{このとき、図形} AacE, PprT \textbf{は互いに同じ比を有します。}
\]

\[
\textbf{なぜなら、個々の平行四辺形の比が等しいとき、(補題 III により)全体の和の比も等しくなります。したがって、一方の図形に対する全体の和の比は他方の図形に対する全体の和の比と同一であり、}
\]
\[
\textbf{また、一方の図形自体の比も他方の図形自体の比と同一である。}
\]

\begin{center}
\includegraphics[width=0.8\linewidth]{037.png}
\end{center}

\[
\text{初めに、前段の和図形と後段の和図形は等比関係にある。}
\]

\textbf{補論}\quad したがって、ある種の量が同じ分割数で分割され、その部分が分割数が無限大になると大きさが無限小となるとき、これらの部分は最初の部分に対して第二の部分を、他の部分もそれぞれ順序通りに持つ一定の比で得られる。全体も互いに対しても同じ一定の比を持つことになる。なぜなら、この補題の図形において、各段階の平行四辺形が互いに分割されるとき、各段階の和が常に等しいからである。

【第1編 第1節：初等比および最終比の方法（補題I〜XI）】


$\text{平行四辺形総和の比は} : \text{部分の数が無限大に増加し, 面積が無限小になるとき, ある平行四辺形に対する最終比は, 即ち (仮定により) ある部分に対する最終比は}.$

\condpagelarge{補題V.}

\textbf{相似な図形のすべての辺, 互いに対応するものであるものは比例的である, 曲線のものも直線のものも同様であり, 面積は辺の二重比に比例する。}

\condpagelarge{補題VI.}

\pngright{037.png}{}{}{}
%Illustration

\textbf{任意の位置で与えられた弧 $AB$ が弦 $AB$ を引くとき,}

punctum aliquo $\textbf{A}$, in medio curvatur$\textbf{ae}$ continu$\textbf{ae}$, tangatur a recta utrumque producta $\textbf{AD}$; dein puncta $\textbf{A}$, $\textbf{B}$ ad invicem accendant \& coeant; dico quod angulus $\textbf{BAD}$ sub chorda \& tangente contentus minuetur in infinitum \& ultimo evanescet.

Nam producatur $\textbf{AB}$ ad $b$ \& $\textbf{AD}$ ad $d$, \& punctis $\textbf{A}$, $\textbf{B}$ coeuntibus, nulla\que adeo ipsius $Ab$ parte $AB$ jacente amplius intra curvam, manifestum esi quod hae recta $Ab$,

\pngright{038.png}{}{}{}

vel coincidet cum tangente $Ad$, vel ducetur inter tangentem \& curvam. Sed casus posterior est contra naturam Curvatur$\textbf{ae}$, ergo


% 画像の配置

\par\begin{center}
    \centering
    \includegraphics[width=4cm,height=3cm]{038.png}

\end{center}\par

\textbf{補題 VII.}

同様に置いたとき、私は弧、弦および接線の最終比が等しい比であると主張します。詳細は図 Lem.\ 6 および Lem.\ 8 vi 参照してください。

実際、$AB$ および $AD$ を $b$ および $d$ に引き延ばし、$BD$ の平行線として $bd$ を引きます。弧 $Ab$ が弧 $AB$ に相似であるとします。また、点 $A$ および $B$ が重なるとき、角 $dAb$（補題より）は消えます；したがって直線 $Ab$ および $Ad$ は中間の弧 $Ab$ と一致し、従って等しいです。よって、これらは常に比例する直線となります。

\textbf{補題I.} 直線 \( AB \), \( AD \)，および弧の中間部分 \( AB \) は等比関係を有する。

\begin{center}
\includegraphics[width=1057px,height=424px,viewport=-24 0 1057 424]{038.png}
\end{center}

\textit{補題Iの推論1.} よって、直線 \( AB \) 上の点 \( B \) を通って任意に引かれた接線平行線 \( BF \) が直線 \( AF \) を点 \( A \) から通過して無限に弧 \( AB \) と交わるとき、この直線 \( BF \) は弧 \( AB \) の極限となる弧に対して等比関係を有する。これは、平行四辺形 \( AFBD \) を完成させると、常に直線 \( BF \) が直線 \( AD \) に対する等比関係を保つからである。

\textit{補題Iの推論2.} また、点 \( B \) と \( A \) を通って他の直線 \( BE \), \( BD \), \( AF \), \( AG \) を引き、これらの直線が接線 \( AD \) およびその平行線 \( BF \) と交わるとき、極限の比は等比関係を有する。


\textbf{第1編 第1節：初等比および最終比の方法（補題I〜XI）}

すべての点 \(AD\), \(AE\), \(BF\), \(BG\) の横軸 \(AD\), \(AE\), \(BF\), \(BG\) および弧 \(AB\) 間の比は等しい比となる。

\textit{補題3.} よって、これらすべての線分は最終的な比において互いに置き換えることができる。

\textbf{補題8.}

\textit{もし与えられた直線 \(AR\), \(BR\) と弧 \(AB\), 弦 \(AB\) および接線 \(AD\) が与えられたとき、三角形 \(ARB\), \(ARB\), \(ARD\) を構成する。そして点 \(A, B\) が互いに近づく場合、私は主張する：最終的な三角形の形は相似性となり、最終的な比も等しい比となる。}


---
$AB$, $AD$, $AR$ をそれぞれ $b$, $d$, $r$ とし、$RD$ に対して平行な線分 $rbd$ を引き、弧 $AB$ に似た弧 $Ab$ を引きます。点 $A$, $B$ が重なるとき、角度 $bAd$ は消滅します。したがって、三角形 $rAb$, $rAb$, $rAd$ は一致し、同名で等しいとみなされます。従って、最終的にこれらの三角形 $RAB$, $RAB$, $RAD$ は互いに似て等しくなります。証明終了。

\textit{補則}：したがって、これらの三角形はすべての最終比での議論において自身として扱えます。
---

\textbf{互換性がある。}

\textit{証明 IX.}


\par\begin{center}
    \centering
    \includegraphics[width=1050px,height=1171px,trim=-24 0 0 0,clip]{039.png}

\end{center}\par

\textit{直線 $AE$ と曲線 $AC$ が与えられた位置で互換的に接する角度 $\angle A$ を持つとき、$AE$ 上の点 $B$, $C$ はそれぞれ直線 $BD$, $EC$ に与えられた角度で順次配置される。$BD$, $EC$ は曲線 $AC$ と交わる点 $B$, $C$ を通る。その後、点 $B$, $C$ が点 $A$ に近づくとき、三角形 $ADB$ および $AEC$ の面積比は最後の直線 $AD$ 上の点 $A$ における辺の比の二倍となる。}

証明するには、$AD$ 上から $Ad$, $Ae$ を $AD$, $AE$ に比例して引き取り、それぞれを直線として引く。



$db, ec$を直線$DB, EC$とし、これらは平行かつ比例します。$AC$を点$c$まで延長し、曲線$Abc$を三角形$ABC$に相似な曲線として描きます。また、直線$Ag$を曲線上の点$A$で接線とします。次に、$F, G, f, g$を$DB, EC$上に順次置きます。

このとき、点$B, C$が点$A$に向かって収束し、角度$cAg$が消滅するにつれて、曲線弧$Abd, Ace$は直線弧$Afd, Age$と一致します。従って、補題Vの論証により、これらの弧は辺$Ad, Ae$との二重比にあります。

しかし、これら弧に対応する面積$ABD, ACE$は常に辺$AD, AE$に対する比例した面積$ABD, ACE$であり、それらの辺$AD, AE$も同様です。

\[
\therefore AB:AC = AD^2 : AE^2.
\]
\textbf{証明終了}

\textbf{補題X}

\textit{任意の規則的な力によって物体が描く空間は、その運動が始まったときの時間比に二倍比をとったものである。}

時間を直線 $AD$, $AE$ で表し、速度を直線 $DB$, $EC$ の縦方向成分で表すとする。この速度によって描かれた空間は、これらの縦方向成分によって描かれた面積 $ABD$, $ACE$ と等しい。したがって、運動が始まったときの時間比 $AD:AE$ に二倍比を取ると、描かれた空間の比 $ABD:ACE$ となる（補題IXより）。 \textbf{証明終了}

\textit{推論1.} したがって、類似な物体は、類似な時間比において相似な空間を描く。


figurarum partes temporibus proportionalibus describentium
errores, qui viribus aequalibus in partibus istis ad corpora similiter
applicatis generantur, \& mensurantur a locis figurarum, ad quae
corpora temporibus idem proportionalibus absque viribus istis pervenirent,
sunt ut quadrata temporum in quibus generantur quam proxime.

\textbf{補題2.} これらの誤差は、等しい力で等しい部分に類似の物体に適用されるとき、そしてこれら物体が等しい時間比において同じ力なしに到達する位置から測定されるとき、生成された時間の平方と比例します。

\condpagelarge{補題XI.}

\textbf{補題11.} 接触角の弧は最後の瞬間における比率は二倍です。

【第1編 第1節：初等比および最終比の方法（補題I〜XI）】

\textbf{補題 I}

接線と接点から中心までの弧を含む角に対する弦の投射について。

\pngright{041.png}{}{}{}

\textbf{案例 1}
弧 \( AB \) が与えられ、その接点 \( A \) から接線 \( AD \) が引かれます。この接線 \( AD \) に対する垂線 \( BD \) と弧 \( AB \) の投射をそれぞれ \( AG \) および \( BG \) として、これらは点 \( G \) で交わります。次に、点 \( D \), \( B \), \( G \) をそれぞれ点 \( d \), \( b \), \( g \) へと近づけ、最終的に直線 \( BG \) および \( AG \) の交点を点 \( J \) として、点 \( D \), \( B \) が点 \( A \) に近づくときを考えます。明らかに、点 \( D \), \( B \) が点 \( A \) に近づくにつれて、距離

\[
\textbf{AG} - \textbf{BG}
\]

は一定の比で減少します。

この証明は図を参照してください。

\textbf{GJは任意に小さな長さをもつことができる。} これは（円周上の点$ABG$, $Abg$を通る）円の性質から、$AB^2 = AG \times BD$ かつ $ab^2 = ag \times bd$であるので、比$AB^2$に対する$Ab^2$の比は、比$AG$に対する$Ag$の比と比$BD$に対する$bd$の比から構成される。しかし、$JG$を任意に小さな長さとすることができるため、比$AG$に対する$Ag$の差が任意に与えられた差より小さいようにすることが可能であり、これにより比$AB^2$に対する$Ab^2$の差も任意に与えられた差より小さいようにすることが可能である。

\[
BD : bd = \text{任意に与えられた差}
\]
従って、補題Iにより、最終比 \( AB^2 \) は \( Ab^2 \) に等しい比 \( BD : bd \) である。□

\textit{案2.} 現在、\( BD \) が任意の与えられた角度で \( AD \) に傾斜しているとしても、最終比 \( BD : bd \) は以前と同様であり、従って \( AB^2 : Ab^2 \) も同じである。□

\textit{案3.} そして、角 \( D \) の値が与えられていなくても、角 \( D \) と \( d \) は常に等しくなり、任意に与えられた差ほど近づくので、最終的には等しくなる。

\textbf{補題I.I}：線分 \( BD \) と \( bd \) は同じ比で互いに接する弧 \( AB \) と \( Ab \)、およびそれらの正弦 \( BC \) と \( bc \) も最終的に等しい。従って、これらの正方形も最終的に互いに等しい。

\textbf{証明}：補題Iにより、線分 \( BD \) と \( bd \) は弧 \( AB \) と \( Ab \) の比で互いに接する。また、接する弧の正弦も最終的に等しくなるため、これらの正方形も最終的に互いに等しい。

\textbf{補題I.II}：三角形 \( ADB \) と \( Adb \) は最終的に三倍比で辺 \( AD \) と \( Ad \) の比であり、辺 \( DB \) と \( db \) の比では二倍半分比である。これは、辺 \( AD \) と \( DB \) および \( Ad \) と \( db \) の組合せ比によるものである。

同様に、三角形 \( ABC \) と \( Abc \) も最終的に三倍比である。

【第1編 第1節：初等比および最終比の方法（補題I〜XI）】

\textbf{補題 3}

直線 \( BC \) と \( bc \) の比。

\textbf{命題 3.} また、もし \( DB \) と \( db \) が最終的に平行であり、それぞれ \( AD \) と \( Ad \) の二倍比であるならば、曲線弧 \( ADB \) と \( Adb \) （パラボラの性質によれば）は直線三角形 \( ADB \) と \( Adb \) の三等分点であり、また弦 \( AB \) と \( Ab \) も同じ三角形の三等分点である。従ってこれらの曲線弧およびこれらの弦もそれぞれ \( AD \) と \( Ad \) の接線の比、そして弦と弧 \( AB \) と \( Ab \) の比において三倍比となる。

\textbf{補説。} 

それ以外において、これら全てについて無限に大きい接角より小さい接角を、つまり円が自らの接線と含む接角よりも小さい接角を、または等しくない接角を、また無限に小さい接角より大きい接角を仮定する。これは曲率半径 $A$ における点 $A$ の曲率を無限に大きくも小さくもないことを意味し、すなわち有限長の区間 $AJ$ を有することである。実際、$DB$ を $AD^3$ とすることができる。この場合、円は点 $A$ において接線 $AD$ と曲線 $AB$ の間に引くことができず、したがって接角は無限に小さい円弧の接角よりも小さくなる。同様の議論により、もし $DB$ を次々と $AD^4$, $AD^5$, $AD^6$, $AD^7$, などとする場合も同様である。

\[
\textbf{証明}
\]

点 $A$ の曲率を無限に大きくも小さくないことを示すために、$DB$ を $AD^3$ と仮定する。このとき、円は点 $A$ において接線 $AD$ と曲線 $AB$ の間に引くことができず、したがって接角 $\angle DAB$ は無限に小さい円弧の接角よりも小さくなる。

同様に、もし $DB$ を次々と $AD^4$, $AD^5$, $AD^6$, $AD^7$, などとする場合も同様である。


habebitur series angulorum contactus pergens in infinitum, quorum quilibet posterior est infinite minor priore. Et si fiat $DB$ successive ut $AD^{2}$, $AD^{\frac{3}{2}}$, $AD^{\frac{4}{3}}$, $AD^{\frac{5}{4}}$, $AD^{\frac{6}{5}}$, $AD^{\frac{7}{6}}$, \&c.\ habebitur alia series infinita angulorum contactus, quorum primus est ejusdem generis cum circularibus, secundus infinite major, \& quilibet posterior infinite major priore. Sed \& inter duos quosvis ex his angulis potest series utrumque in infinitum pergens angulorum intermediorum inseri, quorum quilibet posterior erit infinite major

\begin{center}
\textbf{訳文}

無限に及ぶ接線角の級数が存在し、その各項は前項に対して無限に小さい。もし $DB$ が次のように進化するとき，すなわち $AD^{2}$, $AD^{\frac{3}{2}}$, $AD^{\frac{4}{3}}$, $AD^{\frac{5}{4}}$, $AD^{\frac{6}{5}}$, $AD^{\frac{7}{6}}$, \&c. と進化するとき、別の無限の接線角の級数が存在し、その最初の項は円弧と同じ種類であり、次の項は無限に大きい、そして各項は前項に対して無限に大きい。しかし、これらの角度のどの2つを選んでも、それらの間に入りうる中間角度の級数があり、その各項は前項に対して無限に大きい。
\end{center}

priore. すなわち、$AD^{2}$ と $AD^3$ の間に入れる数列として、
$$
AD^{\frac{13}{6}}, AD^{\frac{11}{5}}, AD^{\frac{9}{4}}, AD^{\frac{7}{3}}, AD^{\frac{5}{2}},
AD^{\frac{8}{3}}, AD^{\frac{11}{4}}, AD^{\frac{14}{5}}, AD^{\frac{17}{6}},
\&c.
$$
を挿入することができます。また、この数列の任意の二つの角度の間には、無限個の微分した角度の間に入れる新たな数列が可能であり、

$$
\textbf{新しい角度の間の無限な差異}。
$$

新規性の自然な境界を知るまでは、これらは不明です。

曲線の線と曲面に関する証明が示された場合、

$$
\textbf{これらの原理は曲面を有する立体の表面にも容易に適用されます}。
$$


内容は、初等比および最終比の方法（補題I〜XI）に関する前文です。実際には、この Lemmata は、古の幾何学者の手法を模して、反論法（reductio ad absurdum）を使って複雑な証明から逃れることを目的としています。無限小法による方法は、証明をより簡潔にしますが、無限小仮説は堅固ではなく、その方法は幾何学的に見劣りするため、私はその後の事項の証明を最終的な量が消失する和と比、最初に現れる量の比、すなわち、和と比の極限まで導くことにしました。そのため、これらの極限の証明はできるだけ簡潔に省略することにしました。

\[
\textbf{訳文：}
\]

ニュートンは、無限小法を用いる方法と同様のことを提供します。既に証明された原理に基づいて無限小法を使用することはより安全です。したがって、以下では、量を個々の部分からなるものとして考慮する場合や、直線を曲線とする場合でも、無限小ではなく消滅する有限部分を考えることに留めます。定まった部分の和または比ではなく、その和または比の極限を常に理解し、これらの証明法は以前の補題の方法に戻されることを常に念頭に置いてください。

反論がありますが、消滅する量には最終的な比はありません。なぜなら、それらが消滅する直前ではまだ最終的ではないからです。


$$\begin{aligned}
かつて存在しなかったものは、存在しない。しかし同様の論理と議論により、ある物体が特定の場所に到達するまでの最終速度も存在しないことを主張できる。なぜなら、その物体が最終的な位置に達する直前におよびその後の速度は存在しないからであり、その物体が最終的な位置に達した瞬間の速度こそが最終速度であるからだ。この答えは容易である。最終速度とは、物体が特定の位置に到達し、運動を停止する直前の速度と直後に速度を持たないとき、そしてその物体が最終的な位置に到達し、運動を停止したときに持つ速度を指すものである。同様に、消失する量の最終比も理解されるべきである。
\end{aligned}$$

rationem quantitatum non antequam evanescunt, non postea,
sed quacum evanescunt. Pariter \& ratio prima nascentium est
ratio quacum nascuntur. Et summa prima \& ultima est quacum
esse (vel augeri \& minui) incipiunt \& cessant. Extat limes quem
velocitas in fine motus attingere potest, non autem transgredi.
\textbf{これが消滅する直前に、または後ではなく、その過程で存在する量の比です。同様に、初生する量の最初の比はその生成時に存在します。また、最初と最後の総和は存在（増加や減少）が始まるときと終わりるときに存在します。} 最終的な速度が運動の終了時に達する限界がありますが、それを超えることはありません。

\textbf{これが最終速度です。そして、量または比の限界はそれぞれ開始と終了時に存在する最初と最後のものと等しい比です。ここでの限界は確定的で定義されていますので、その値を正確に決定する問題は真の幾何学的な問題となります。} しかし、幾何学的なすべてのことは他の幾何学的な要素において表現されます。




定められ、証明されるためには、極限する量の最終比が与えられた場合、最終量も与えられることが可能である。従って、すべての量は無限小を構成することになるが、これはユークリッドが『要素』第10巻で示した反対の主張に矛盾する。

しかし、この反論は仮定に基づいており、偽である。極限する量の最終比は、実際には最終量の比ではなく、減少する量の無制限に近づく比の限界であり、

\[
\textbf{すなわち}
\]

無限小を構成する量の限界である。

\textbf{より小さい比を求めることが可能で、任意に与えられた差よりも小さく接続されますが、決してそれを超えることはありません。また、量が無限小になる前に到達することもありません。これは無限大の量において明確になります。}

もし二つの量の差が与えられていれば、これらの量を無限に増加させることで、その最終的な比が得られます。つまり、等しい比が得られるのです。しかし、それゆえに最終的な量または最大の量が得られるわけではありません。したがって、以下では、想像が容易になる場合、最小の量、あるいは消滅する量、あるいは最終的な量について述べるときには、その量を大きさで定義されたものとして理解してほしいと注意します。

次のように考えよ。無限小の減少を繰り返すことを。

\begin{center}
\includegraphics[width=0.8\textwidth]{045.png}
\end{center}


\section{第1編 第2節：向心力の決定（命題I〜X）}


% Section: 第1編 第2節：向心力の決定（命題I〜X）



\textbf{第2節：向心力の決定}

\textit{向心力を決定するための面積に関する定理。}

\textbf{命題I．定理I．}

\textit{物体が不動中心を向いた円運動で描く面積は、その時間に比例し、平面内に存在する。}


\par\begin{center}
    \centering
    \includegraphics[width=1642px,height=1513px,trim=-24 0 0 0,clip]{045.png}

\end{center}\par

時間が等しい部分に分割され、最初の時間部分では物体が固有の力により直線 $AB$ を描くとする。次に同じ時間部分では、もし何か障害がないなら、物体は直線 $c$ へと進むことになる（第1法則による）。

\[
\text{線分 } Bc \text{ が } AB \text{ と等しいので、中心から引かれた半径 } AS, BS, cS \text{ に作用したとき、}
\]
\[
\text{弧 } ASB, BSc \text{ も等しくなる。しかし、物体が } B \text{ に達すると、向心力は一瞬の衝撃を伴って大きな力を発揮し、}
\]
\[
\text{物体を直線 } Bc \text{ から離れるようにして、直線 } BC \text{ へと進ませる。平行な直線 } BS \text{ 上に引かれた直線 } cC
\]
\[
\text{が } BC \text{ と交わり、} C \text{ で物体は次の時間の半分を経て、法則Iの補題により } C \text{ に存在する。このとき、}
\]
\[
\text{平面 } ASB \text{ と同じ平面にあることが確認できる。} SC \text{ を結び、三角形 } SBC \text{ を考えると、平行線 } SB, cC
\]
\[
\text{により、三角形 } SBC \text{ は三角形 } SBc \text{ と等しい。したがって、また三角形 } SAB \text{ も等しい。同様の議論により}
\]

\textbf{第1編 第2節：向心力の決定（命題I〜X）}

% -----File: 046.png---
向心力が点$C$, $D$, $E$, \ldots{}に順次作用し、その結果として物体は各時間部分において直線$CD$, $DE$, $EF$, \ldots{}を描く。これら直線は同一平面に存在し、また三角形$SCD$は三角形$SBC$と$SDE$と等しいであり、同様に三角形$SEF$も等しいである。したがって等しい時間部分において等しい面積が描かれる：そして合成すると、任意の面積和$SADS$と$SAFS$はそれぞれ描写時間比に等しい。

既に点数を増加させ、三角形の幅を無限大まで小さくする。これらの

\textit{注}：ここでニュートンは、三角形の幅を無限小量とする極限過程について言及していますが、具体的な証明や論理展開は次の段落に続きます。

最終的な周界 $ADF$ は（第3 LemmatisのCorollarium quartumに従って）曲線である。したがって、物体がこの曲線上の接点から永久に引き寄せられる向心力は継続的に働き続け、円錐曲線 $SADS$, $SAFS$ の任意の面積が時間とともに常に比例するとき、これらの面積も同じ時間において比例する。

\textbf{命題 I.} 証明終了 \QED

\textit{Corol.\,1.} 中心媒質が抵抗しない場合、任意の面積が時間に比例しないとき、力は半径に対して収束しない。

\textit{Corol.\,2.} 全ての中心媒質において、面積の描写が加速されるとき、



これらの証明は、物体が曲線上から引き寄せられる向心力を継続的に作用させるための論理を示しています。また、面積と時間の比例性についての命題も述べられています。

向心力は放射線の交点に指向せず、それらから離れて円錐曲線の軌道上を下降する。

命題II. 問題II.

任意の物体が曲線線上を動くとき、その物体が一直線運動を一定速度で進行するために直線方向に引かれた放射線を通じて固定点または一直線運動で進行する点へ円錐曲線を描く面積について、時間比に比例する力は、その物体に対して向心力をもたらす。

1. つまり、任意の物体が曲線上を動くとき、その物体は一定速度で一直線運動を離れて曲線方向に歪められる。これは第一法則（Leg. I）により明らかである。そして、その歪められる力は物体が一直線運動から逸脱する力である。



この図において、$O$ は固定点であり、$PQ$ は一直線運動の軌道を表し、物体が曲線 $PAQ$ 上を動くとき、放射線 $OP, OQ, OR, OS, \ldots$ によって描かれる面積が時間比に比例する力は、その物体に対して向心力をもたらす。

\& 三角形 \( SAB \), \( SBC \), \( SCD \) 等を固定点 \( S \) 囲み，等しい時間内に等しい面積を描くとする。
これらは直線 \( cC \) と平行な直線 \( BS \) 上の点 \( B \) で，また直線 \( dD \) と平行な直線 \( CS \) 上の点 \( C \) で動く（第1編 第40命題及び要素論第二定理による）。
すなわち，これらは固定点 \( S \) に向かって一直線に作用する。

\textbf{証明：} これらの三角形は等しい時間内に等しい面積を描くので，直線 \( BS \) 上の点 \( B \) と直線 \( CS \) 上の点 \( C \) は固定点 \( S \) に向かって一直線上に作用する。

\textbf{案2：} また，第五定理の推論により，物体が曲線形を描く平面において静止している場合でも同様である。

\textbf{命題I.} すなわち、向心力が複数の力から成る場合、物体はその力を含む同一方向に移動し、点$S$を通る直線上を均一に運動する。

\textit{\textbf{補説。}} 向心力の合成力によって物体が押されるとき、この命題の意味は、これらの力の合計が点$S$に向かう方向にあるということである。一方で、ある力が描かれた平面に垂直な線を追跡する場合、その力は物体が運動平面から偏るよう作用し、しかし描かれた面積の量には影響を与えないため、合成力の計算において無視される。

\begin{center}
\includegraphics[width=0.8\textwidth]{figura_descripta_perpendicularem.png}{幅}{高さ}{オフセット}
\end{center}

\textbf{命題 III. 理論 III.}

\textit{いかなる物体も、他の物体に対して中心から放射された軌道上を動かされると、その中心周りに一定の時間比で面積を描く。この物体は、向心力によって他の物体に引き寄せられることと、加速力を受けることにより、これら二つの力を合せた力によって動かされる。}

なぜなら（法則の補則第 6）\textbf{、}もし新しい力が物体に対してその他の物体に引き寄せられる力と等しく且つ反対方向にある場合、両方の物体は平行線を描く軌道上を動かされると、最初の物体は以前と同じ面積を描く。しかし、物体に対して作用する力は、向心力によって他の物体に引き寄せられることと、加速力を受けることにより、これら二つの力を合せた力である。


---
あるものは、すでに自己に等しく反対の力を受けて崩壊するので、その他の物体は（第1法則により）静止するか一直線上に均一速度で移動し、第一の物体は異なる力の差によって中心周囲に時間比に比例した面積を描く。従って、その他の物体に対する力の差は（理論2により）中心と同一である。（証明終）

\textit{命題1.} ここで、一方の物体が他方の物体に対して半径ベクトルを通じて時間比に比例した面積を描くなら、全体の力（単純な力還是複合的な力であっても）は
---

\textbf{証明終}

\textit{\textbf{命題1.}} 一方の物体が他方の物体に対して半径ベクトルを通じて時間比に比例した面積を描くなら、全体の力（単純な力還是複合的な力であっても）は


この訳文では、証明終了の部分を「\textbf{証明終}」とし、命題1のタイトルを「\textit{\textbf{命題1.}}」としています。また、日本語で自然に表現するために、一部の用語を適切に変更しています。

【第1編 第2節：向心力の決定（命題I〜X）】

\textbf{命題II}

物体が以前に引き寄せられた力を多数の力によって合成した場合、その合計の加速力は物体が他の物体を引き寄せたときに与えた力と等しいとする法則のコロラリウムにより、元の引く力は物体に対して向心力として作用します。

\textbf{コロラリウム 2.} そして、もしこれらの面積が時間においてほぼ比例している場合、残りの力は他の物体に対してほぼ比例して作用します。

\textbf{コロラリウム 3.} 逆に、もし残りの力が他の物体に対してほぼ比例して作用する場合、これらの面積は時間においてほぼ比例しています。

\textbf{命題IV．推論4}

もし、ある物体が一直線に引かれた半径によって別の物体を描く面積の比が非常に異なるとき、そしてその他の物体は静止しているか直線運動をしているとき、向心力がその他の物体に向かう作用は存在しないか、または非常に微弱な他の力と混合され、合成される。これらのすべての力を合計した全作用は、面積を均一に描く中心に対して指向するものとなる。

同様のことが成り立つ。

$\text{もし他の物体が向心力の作用により移動するとするとき、全作用力を引いた残りの力が物体に働く場合、}\ \textbf{その面積速度は中心点を指し、物体はこの中心に向かって作用する力によって軌道上に保持されます。また、すべての円運動は、物体が直線運動から引き戻され、軌道に留まるものとして正しく定義されます。}$

\textit{解説}

面積速度は、物体が最も強く影響を受ける力に対して中心点を指す中心点であり、物体はこの中心に向かって作用する力を通過して軌道上に保持されるため、次に述べる際には、面積速度をその指数として使用することとします。

\textbf{第1編 第2節：向心力の決定（命題I〜X）}

\textit{自由空間で等速円運動を行う物体が異なる円周上にある場合、これらの物体に対する向心力をどのように決定するか？}

\textbf{命題IV. 定理IV.}


\par\begin{center}
    \centering
    \includegraphics[width=1000px,height=1254px,trim=-24 0 0 0]{049.png}

\end{center}\par

異なる円周 \( BD \) と \( bd \) 上を等速で回転する物体 \( B \) と \( b \)，これらが弧 \( BD \) および \( bd \) を描くとき，物体はそれぞれ接線 \( BC \) および \( bc \) を描き，これらの接線は弧 \( BD \) および \( bd \) に等しい長さを持つ。これは明らかである。

\[
\textbf{証明}
\]

物体 \( B \) と \( b \) が円周 \( BD \) および \( bd \) 上を等速で回転し，それぞれ弧 \( BD \) および \( bd \) を描くとき，これらの物体は接線 \( BC \) および \( bc \) を描き，これら接線の長さは弧 \( BD \) および \( bd \) の長さと等しい。これは直感的に明らかである。

\[
\textbf{結論}
\]

異なる円周上を等速で回転する物体に対する向心力は，これらの物体が描く弧の平方と円の半径との比に比例し，すなわち

\[
F \propto \left(\frac{\text{弧の平方}}{\text{円の半径}}\right)
\]

となる。

向心力は、円周上の接線から物体を引き戻すものであり、これらの力は円の中心に対して最初に生成される距離 \(CD\) と \(cd\) の比にあります。実際には、これらは円の半径で描かれた面積に対応する時間の比であるため、理論IIにより、円周 \(BD\) に対する円弧 \(bt\) の比として存在します。また、論理Vにより、生成される線分 \(tk\) は生じる線分 \(dc\) に対して \(bt^2\) と \(bd^2\) の比であります。これらは等しいことから、

\[
\textbf{線分 } CD \textbf{ は線分 } kt \textbf{ に対して} \frac{\textit{弧 } BD}{\textit{弧 } bt}
\]

また、

\[
\textbf{生成される線分 } tk \textbf{ は生成される線分 } dc \textbf{ に対して} \frac{bt^2}{bd^2}
\]

\textbf{命題 I. Corol. 1.} 向心力は円運動の速度の平方を沿った半径に比例する。

\textbf{命題 I. Corol. 2.} 復元すると、周期的な時間の平方を沿った半径に比例する力は互いに反比例する。すなわち（幾何学者によれば）：

$$
\frac{BD \times bt}{Sb} : \frac{bd^{\text{quad.}}}{Sb} = \frac{BD^{\text{quad.}}}{SB} : \frac{bd^{\text{quad.}}}{Sb}
$$

これは、等しい比 $\frac{bt}{Sb}$ と $\frac{BD}{SB}$ の結果である。

\textbf{命題 I. Corol. 1.} 向心力は円運動の速度の平方を沿った半径に比例する。

\textbf{命題 I. Corol. 2.} 復元すると、周期的な時間の平方を沿った半径に比例する力は互いに反比例する。すなわち（幾何学者によれば）：

$$
\frac{BD \times bt}{Sb} : \frac{bd^{\text{quad.}}}{Sb} = \frac{BD^{\text{quad.}}}{SB} : \frac{bd^{\text{quad.}}}{Sb}
$$


\begin{center}
\includegraphics[width=0.8\textwidth]{050.png}
\end{center}

\textbf{命題 I. Corol.\,3.} ここで、周期的な時間が測定されたとします。すると、向心力と速度は半径に比例し、逆に、半径と速度は向心力に比例します。

\textbf{命題 II. Corol.\,4.} 周期的な時間の平方が半径に比例するとき、向心力は等しくなり、速度は半径の半分の比で変化します。逆もまた然りです。

\textbf{命題Ⅴ．}

\textit{もし周期的時間の平方が半径の平方と比例するならば、向心力は半径と反比例し、且つ等しい速度を持つ物体も同様である。逆に、}

\textbf{命題Ⅵ．}

\textit{もし周期的時間の平方が半径の立方と比例するならば、向心力は半径の平方と反比例し、速度は半径の半分の比で比例する。逆に、}

\textbf{命題Ⅶ．}

\textit{同様にして、相似な図形において中心が似た位置にある相似な部分を描く物体について、時間、速度、及び向心力に関するすべては、}



先の証明を、これらの場合に適用する。

\textit{注記：}

天体において（私たちはWrennus, Hookius \& Halleusも含め）第六補則の結果が得られることが確認されています。よって、向心力が中心からの距離の二乗に反比例する力を考慮するためには、後述しますので詳しく説明しません。

また、前節の証明から得られる利点により，任意の力（例えば重力）に対する向心力の比も収集できます。すなわち，その力は時間に沿って物体が通過するとき、

\[
F \propto 1/r^2
\]

となることを示します。

向心力が物体を一定時間内に特定の弧$BC$を推進し、その軌跡は点$D$までの距離$CD$を覆う。この距離$CD$は運動開始時に等しい弧$BD$の平方と円の直径との比である。また、物体は同じ力で一定時間継続して運動すると、特定の弧を描くときと同じ時間内に直線運動で等しい時間の平方根に比例する距離を描く。すなわち、その向心力を与えられた時間内で任意の弧を描くときに物体が直線運動で描く距離は、重力によって同じ時間内に落下したときの距離と等しくなる。これらの命題により、

\[
\textbf{arcum } BC, \textbf{ impellat ipsum per spatium } CD, quod ipso motus initio aequale est quadrato arcus illius BD ad circuli diametrum applicato; \& corpus omne vi eadem in eandem semper plagam continuata, describat spatia in duplicata ratione temporum: Vis
\]
\[
\textbf{illa, quo tempore corpus revolvens arcum quemvis datum describit, efficiet ut corpus idem recta progrediens describat spatium quadrato arcus illius ad circuli diametrum applicato aequale; adeoque est ad vim gravitatis ut spatium illud ad spatium quod grave cadendo eodem tempore describit.}
\]
\[
\textbf{Et hujusmodi Propositionibus}
\]

\textbf{ヒゲノイウス}は、優れた著書『振動時計に関する論文』において、重力の力と円運動する物体に対する遠心力との関係を明らかにしました。

さらに、以下の命題を示すことができます。任意の円周上を描く多角形を考えます。そして、物体が多角形の各辺上にある一定の速度で動いている場合、その物体が各頂点から円周へ反射されるとき、各反射時に円周に衝突する力はその物体の速度と等しくなります。したがって、与えられた時間における総力はその速度と反射回数の積に比例します。すなわち（多角形が特定の種類の場合）、

\[
\text{総力} \propto \text{速度} \times \text{反射回数}
\]

となります。



長い線分が与えられた時間に描かれる長さと、同じ長さが円の半径に引かれたとき、すなわちその長さの平方が円の半径に引かれたとき、もし無限に小さくなった多角形が円に一致し、与えられた時間に描かれた弧の平方が円の半径に引かれるとき、これらは等しいです。これは、物体が共通の中心に向かって力によって描く図形を描く速度で物体を引きつける力を表します。これとは反対に、円は物体を中心から引き離す力を与えます。

\textbf{命題Ⅴ．問題Ⅰ}

任意の点での速度が与えられたとき、物体が力によって共通の中心に向かって一定の図形を描く中心


\textbf{これを発見すること。}


\par\begin{center}
    \centering
    \includegraphics[width=1150px,height=934px,keepaspectratio]{052a.png}

\end{center}\par

図に示したように、直線 \(PT\), \(TQV\), \(VR\) が点 \(P\), \(Q\), \(R\) で接し、これらは点 \(T\) と \(V\) で交わる。これらの接線上からそれぞれ垂直線 \(PA\), \(QB\), \(RC\) を引き、各点 \(P\), \(Q\), \(R\) で体が有する速度に比例する長さとする；すなわち、\(PA\) は \(QB\) に、そして \(QB\) は \(RC\) にそれぞれその比をもつようにし、これらはそれぞれ点 \(P\), \(Q\) での速度と点 \(Q\), \(R\) での速度の比に対応する。垂直線の端点 \(A\), \(B\), \(C\) を直角に見立てる。

\textbf{命題I}：点 \( S \), \( D \), および \( T \) は一直線上にある。

証明：引く線 \( AD \), \( DBE \), および \( EC \) が点 \( D \) と \( E \) で交わり、中心 \( S \) を通って線 \( TD \) および \( VE \) が交わる。

次に、物体が点 \( P \) および \( Q \) から中心に向かって放射したときに描く面積が時間の比に比例するとき、点 \( P \) および \( Q \) で描かれた面積はそれぞれ点 \( S \) から中心へ向かう垂線 \( PT \) および \( QT \) と直角に引いた速度の比に等しい。したがって、垂線 \( AP \), \( BQ \) は直接比例するものであり、つまり、点 \( D \) から中心への垂線を引いたときの垂線 \( AD \) および \( BD \) の比に等しい。

以上より容易にわかるように、点 \( S \), \( D \), および \( T \) は一直線上にある。

\[
\textbf{そして似たような論証により、点} \ S, \ E, \ V \ \textbf{は一直線上にあり、ゆえに中心} \ S \ \textbf{は直線} \ TD, \ VE \ \textbf{の交点に向かって回転する。}
\]
\QED

\condpagelarge{命題VI. 問題V.}

\pngright{052b.png}{}{}{}

\textbf{もし物体} \ P \ \textbf{が中心} \ S \ \textbf{を中心に回転し、曲線} \ APQ \ \textbf{を描き、かつ直線} \ ZPR \ \textbf{とこの曲線を点} \ P \ \textbf{で接する。また、曲線} \ Q \ \textbf{から点} \ P \ \textbf{を通る平行線} \ QR \ \textbf{を引いて、直線} \ SP \ \textbf{と垂直な直線} \ QT \ \textbf{を下ろす。このとき私は主張する：向心力は} \ SP \ \textbf{の立方体に反比例する。}
\[
\textit{Si corpus } P \textbf{ revolvendo circa centrum } S, \textbf{ describat lineam quamvis curvam } APQ, \textbf{ tangat vero recta } ZPR \textbf{ curvam illam in puncto quovis } P, \& \textbf{ ad tangentem ab alio quovis curvae } Q \textbf{ agatur } QR \textbf{ distantiae } SP \textbf{ parallela, ac demittatur } QT \textbf{ perpendicularis ad distantiam } SP: \textbf{ Dico quod vis centripeta sit reciproce ut solidum }
\]

\[
\frac{SP \textbf{quad.} \times QT \textbf{quad.}}{QR}
\]
もし、その固体の量が常に最後に等しいとすると、点 \(P\) および \(Q\) が合流するときです。

なぜなら、図形を無限に伸長させた場合、線分 \(QRPT\) の微小な部分 \(QR\) は与えられた時間において、向心力（第2法則により）と比例し、また与えられた速度の平方（第1補題により）と比例します。従って、どちらか一方が与えられているとき、向心力と速度の平方は連続的に比例します。したがって、向心力は線分 \(QR\) に直接比例し、時間の平方に逆比例します。

また、時間は面積 \(SPQ\) に等しく、その二倍 \(SP \times QT\) に等しいです。

\[
\textbf{すなわち、} SP \& QT の積は、向心力  F \textbf{ が QR に比例し、}
\]
\[
\textbf{一方で} SP^2 \textbf{ と } QT^2 \textbf{ の比の逆数に比例する。すなわち、}
\]
\[
F \propto \frac{SP^2 \times QT^2}{QR}.
\]

\textit{\textbf{補則}}：したがって、任意の図形が与えられ、その中に向心力を受ける点が存在する場合、その図形の円周上を回る物体に対する向心力の法則を求めることができる。すなわち、計算すべきは
\[
\frac{SP^2 \times QT^2}{QR}
\]
がこの力に対して反比例することである。これについては、次節の問題において例を示す。



\textbf{命題 VII. 問題 II.}

\textit{円周上を回る物体に対して、与えられた円周上の一定点に向かう向心力の法則を求めよ。}


\par\begin{center}
    \centering
    \includegraphics[width=1043px,height=662px,trim=-24 0 0 0,clip]{053.png}

\end{center}\par

設、円周 $SQPA$ の中心を $S$、円周上の点 $P$ を物体とする。また、点 $Q$ は物体が最も近づくと予想される位置とする。直径 $SA$ および直線 $SP$ に対して、点 $P$ から垂直に下ろした垂線 $PK$ と、点 $Q$ から平行に引いた直線 $LR$（円周 $SQPA$ 上の点 $L$ で交わる）とを引き、さらに直線 $TQ$ および $PR$ の交点を $Z$ とする。

\textbf{証明}

まず、直角三角形 $SPK$ において、$\angle SPK = 90^\circ$ である。また、直線 $LR$ は円周上を接するので、$\angle LRP = \angle SLP$（円の内接角定理）。

次に、直線 $TQ$ と $PR$ の交点 $Z$ を用いて、向心力に関する法則を求めよう。直線 $SP$ 上で平行な直線 $LR$ が引かれていることから、$\angle ZRP = \angle QSP$（対応角定理）。

ここで、円周上の任意の点を考慮すると、物体は一定の速度で円周上を移動するため、向心力は常に物体を中心に引き戻す。したがって、向心力を求めるには、物体が円周上を回る際に直角三角形や接線との関係から得られる幾何学的な法則を利用することが必要となる。

以上の証明により、与えられた円周上の一定点 $Q$ に向かう向心力の法則を求めることができる。

向心力の決定（命題I〜X）【第2節】

三角形 \( ZQR \), \( ZTP \), および \( SPA \) の相似性から、
\[
RP^2 \quad (\text{すなわち} \; QRL)
\]
は
\[
QT^2
\]
に比例し、また
\[
SA^2
\]
は
\[
SP^2
\]
に比例します。

したがって、
\[
\frac{QRL \times SP^2}{SA^2}
\]
は
\[
QT^2
\]
これらに $\frac{SP^2}{QR}$ を掛け、点 $P$, $Q$ が合致するときに $RL$ のかわりに $SP$ を書けば、$\frac{SP^5}{SA^2}$ は $\frac{QT^2 \times SP^2}{QR}$ に等しくなる。ゆえに（定理Vの系により）向心力は $\frac{SP^5}{SA^2}$ に反比例する。すなわち（$SA^2$ が与えられているから）距離 $SP$ の平方立方（5乗）に反比例する。これが求めるものであった。\QED

\condpagelarge{命題 VIII. 問題 III.}

\pngright{054.png}{}{}{}

\textit{物体が円 $PQA$ の上を動くとする。この運動のために必要な、向心力の法則を求めよ。ただしその力は、あまりに遠い点に向かうので、その点へ引いた線 $PS$, $RS$ がすべて平行とみなせるものとする。}

円の中心 $C$ から半径 $CA$ を、これらの平行線を垂直に切って $M$, $N$ で交わるように引き、$CP$ を結ぶ。相似な三角形 $CPM$ と $TPZ$、あるいは（補助定理VIIIにより）$TPQ$ から、$CP^2$ 対 $PM^2$ は $PQ^2$、あるいは（補助定理VIIにより）$PR^2$ 対 $QT^2$ に等しい。また、円の性質から、長方形 $QR \times (RN + QN)$ は $PR^2$ に等しい。点 $P$, $Q$ が合致するとき、$RN + QN$ は $2PM$ に等しくなる。ゆえに $CP^2$ 対 $PM^2$ は $QR \times 2PM$ 対 $QT^2$ に等しく、したがって $\frac{QT^2}{QR}$ は $\frac{2PM^3}{CP^2}$ に等しく、$\frac{QT^2 \times SP^2}{QR}$ は $\frac{2PM^3 \times SP^2}{CP^2}$ に等しい。ゆえに（定理Vの系により）向心力は $\frac{2PM^3 \times SP^2}{CP^2}$ に反比例し、すなわち（定まった比 $\frac{2SP^2}{CP^2}$ を無視して）$PM^3$ に反比例する。\QED

\condpagelarge{注.}

同様の論法で、物体は、縦線の3乗に反比例する向心力で、きわめて遠い力の中心に向かうならば、楕円の上も、さらには双曲線や放物線の上も動く。

\condpagelarge{命題 IX. 問題 IV.}

\pngright{055.png}{}{}{}

\textit{物体がうず巻線 $PQS$ の上を回り、この線がすべての半径 $SP$, $SQ$ などと与えられた角で交わるとする。うず巻線の中心に向かう向心力の法則を求めよ。}

限りなく小さい角 $PSQ$ を与える。すべての角が与えられているから、図形 $SPQRT$ は形が与えられる。ゆえに比 $\frac{QT}{RQ}$ が与えられ、$\frac{QT^2}{QR}$ は $QT$ に、すなわち $SP$ に比例する。いま角 $PSQ$ を任意に変えると、接角 $QPR$ に張る直線 $QR$ は（補助定理XIにより）$PR$ すなわち $QT$ の二乗の比で変わる。ゆえに $\frac{QT^2}{QR}$ は前と同じまま、すなわち $SP$ に比例する。したがって $\frac{QT^2 \times SP^2}{QR}$ は $SP^3$ に比例し、すなわち（定理Vの系により）向心力は距離 $SP$ の3乗に比例する〔原文には reciproce（反比例）の語がない。定理Vの系からは反比例となる〕。\QED

\condpagelarge{補助定理 XII.}

\textit{与えられた楕円のまわりに描かれた平行四辺形は、すべて互いに等しい。双曲線でも、その直径のまわりに描かれた平行四辺形について同様である。}

どちらも円錐曲線の性質から明らかである。


\condpagelarge{命題 X. 問題 V.}

\textit{物体が楕円の上を回るとする。楕円の中心に向かう向心力の法則を求めよ。}

$CA$, $CB$ を楕円の半軸、$GP$, $DK$ を共役な直径、$PF$, $Qt$ を直径への垂線、$Qv$ を直径 $GP$ への縦線とする。平行四辺形 $QvRP$ を作れば、（円錐曲線の性質により）$PvG$ 対 $Qv^2$ は $PC^2$ 対 $CD^2$ であり、（相似な三角形 $Qvt$, $PCF$ により）$Qv^2$ 対 $Qt^2$ は $PC^2$ 対 $PF^2$ である。比を合わせれば、$PvG$ 対 $Qt^2$ は、$PC^2$ 対 $CD^2$ と $PC^2$ 対 $PF^2$ との合成であり、すなわち $vG$ 対 $\frac{Qt^2}{Pv}$ は $PC^2$ 対 $\frac{CD^2 \times PF^2}{PC^2}$ である。$Pv$ のかわりに $QR$ を、（補助定理XIIにより）$CD \times PF$ のかわりに $BC \times CA$ を、また（点 $P$, $Q$ が合致するとき）$vG$ のかわりに $2PC$ を書き、外項と内項を掛け合わせれば、$\frac{QT^2 \times PC^2}{QR}$ は $\frac{2BC^2 \times CA^2}{PC}$ に等しくなる。ゆえに（定理Vの系により）向心力は $\frac{2BC^2 \times CA^2}{PC}$ に反比例し、すなわち（$2BC^2 \times CA^2$ が与えられているから）$\frac{1}{PC}$ に反比例し、いいかえれば距離 $PC$ に正比例する。\QED

\pngright{056.png}{}{}{}

\textit{系1.}\wsp{}したがって、逆に、もし力が距離に比例するなら、物体は、力の中心に中心をもつ楕円の上を、あるいは楕円が変わりうる円の上を動く。

\textit{系2.}\wsp{}また、同じ中心のまわりのすべての図形での回転の周期は、等しい。実際、それらの時間は、相似な楕円では（命題IVの系3, 7により）等しい。長軸を共有する楕円では、それらは互いに、楕円全体の面積に正比例し、同時に描かれる面積の微小部分に反比例する。すなわち、短軸に正比例し、主頂点での物体の速さに反比例し、いいかえれば、それらの軸に正比例し、他方の軸への縦線に反比例する。したがって（正比と反比が等しいから）等しさの比になる。

\condpagelarge{注.}

もし楕円が、中心が無限に遠ざかるとき放物線に変わるなら、物体はこの放物線の上を動き、いま無限に遠い中心に向かう力は一様になる。これがガリレオの定理である。さらに、放物線である円錐曲線が、切られた円錐に対する平面の傾きを変えて双曲線に変わるなら、物体はその周の上を動き、向心力は遠心力に変わる。


\section{第1編 第3節：偏心円錐曲線における物体の運動（命題XI〜XVII）}


% Section: 第1編 第3節：偏心円錐曲線における物体の運動（命題XI〜XVII）



\begin{center}\textit{離心的な円錐曲線における物体の運動について}\end{center}

\condpagelarge{命題XI. 問題VI.}

\textit{物体が楕円の上を回るとする。楕円の焦点に向かう向心力の法則を求めること。}

\pngright{056.png}{}{}{}

上の楕円の焦点を$S$とする。$SP$を引き、楕円の直径$DK$と$E$で、縦線$Qv$と$x$で交わらせ、平行四辺形$QxPR$を完成させる。$EP$が長半軸$AC$に等しいことは明らかである。なぜなら、楕円のもう一つの焦点$H$から$EC$に平行に線$HI$を引けば、（$CS$と$CH$が等しいので）$ES$と$EI$が等しく、したがって$EP$は$PS$と$PI$の半和、すなわち（$HI$と$PR$が平行で、角$IPR$と$HPZ$が等しいから）$PS$と$PH$の半和になり、$PS$と$PH$の和は全体の軸$2AC$に等しいからである。$SP$へ垂線$QT$を下ろし、楕円の主たる通径（すなわち$\frac{2BC^2}{AC}$）を$L$と呼ぶと、$L \times QR : L \times Pv = QR : Pv$、すなわち$PE$（$AC$）$: PC$であり、$L \times Pv : Gv \times Pv = L : Gv$であり、$Gv \times Pv : Qv^2 = CP^2 : CD^2$であり、（補助定理VIIIにより）点$Q$と$P$が一致するとき$Qv^2 : Qx^2$は等しいという比にある。そして$Qx^2$、すなわち$Qv^2$は$QT^2$に対して、$EP^2 : PF^2$、すなわち$CA^2 : PF^2$、あるいは（補助定理XIIにより）$CD^2 : CB^2$の比にある。これらすべての比を合成すれば、$L \times QR : QT^2$は、$AC : PC$、$L : Gv$、$CP^2 : CD^2$、$CD^2 : CB^2$の合成、すなわち$AC \times L$（すなわち$2CB^2$）$\times CP^2 : PC \times Gv \times CB^2$、あるいは$2PC : Gv$になる。ところが点$Q$と$P$が一致するとき、$2PC$と$Gv$は等しい。したがってそれらに比例する$L \times QR$と$QT^2$も等しい。これらの等しいものに$\frac{SP^2}{QR}$を掛ければ、$L \times SP^2$は$\frac{SP^2 \times QT^2}{QR}$に等しくなる。ゆえに（定理Vの系により）向心力は$L \times SP^2$に反比例し、すなわち距離$SP$の二乗に反比例する。\QED

問題Vを放物線と双曲線に移したのと同じ簡潔さで、ここでも同じことをしてもよい。しかしこの問題の重要さと、後に使われることを考えて、他の場合も証明で確かめるのを惜しまない。

\condpagelarge{命題XII. 問題VII.}

\textit{物体が双曲線の上を動くとする。図形の焦点に向かう向心力の法則を求めること。}

\pngright{060.png}{}{}{}

$CA$、$CB$を双曲線の半軸、$PG$、$KD$を共役な直径、$PF$、$Qt$を直径への垂線、$Qv$を直径$GP$に対する縦線とする。$SP$を引き、直径$DK$と$E$で、縦線$Qv$と$x$で交わらせ、平行四辺形$QRPx$を完成させる。$EP$が、横半軸$AC$に等しいことは明らかである。なぜなら、双曲線のもう一つの焦点$H$から$EC$に平行に線$HI$を引けば、$CS$と$CH$が等しいので$ES$と$EI$が等しくなり、したがって$EP$は$PS$と$PI$の半差、すなわち（$HI$と$PR$が平行で、角$IPR$と$HPZ$が等しいから）$PI$と$PH$の半差になり、それらの差は全体の軸$2AC$に等しいからである。$SP$へ垂線$QT$を下ろす。そして双曲線の主たる通径（すなわち$\frac{2BC^2}{AC}$）を$L$と呼ぶと、$L \times QR : L \times Pv = QR : Pv$、すなわち$PE$（$AC$）$: PC$であり、$L \times Pv : Gv \times Pv = L : Gv$であり、$Gv \times Pv : Qv^2 = CP^2 : CD^2$であり、（補助定理VIIIにより）$Qv^2 : Qx^2$は、点$Q$と$P$が一致するとき、等しいという比になる。そして$Qx^2$、すなわち$Qv^2$は$QT^2$に対して、$EP^2 : PF^2$、すなわち$CA^2 : PF^2$、あるいは（補助定理XIIにより）$CD^2 : CB^2$の比にある。これらすべての比を合成すれば、$L \times QR : QT^2$は、$AC : PC$、$L : Gv$、$CP^2 : CD^2$、$CD^2 : CB^2$の合成、すなわち$AC \times L$（すなわち$2CB^2$）$\times PC^2 : PC \times Gv \times CB^2$、あるいは$2PC : Gv$になる。ところが点$Q$と$P$が一致するとき$2PC$と$Gv$は等しい。したがってそれらに比例する$L \times QR$と$QT^2$も等しい。これらの等しいものに$\frac{SP^2}{QR}$を掛ければ、$L \times SP^2$は$\frac{SP^2 \times QT^2}{QR}$に等しくなる。ゆえに（定理Vの系により）向心力は$L \times SP^2$に反比例し、すなわち距離$SP$の二乗に反比例する。\QED

同じようにして、この向心力が遠心力に変わるとき、物体は共役な双曲線の上を動くことが示される。

\condpagelarge{補助定理XIII.}

\textit{放物線の、任意の頂点に属する通径は、その頂点から図形の焦点までの距離の4倍である。}これは円錐曲線論から明らかである。

\condpagelarge{補助定理XIV.}

\textit{放物線の焦点から、その接線へ下ろした垂線は、焦点から接点までの距離と、焦点から図形の主頂点までの距離との比例中項である。}

\pngright{062.png}{}{}{}

$APQ$を放物線、$S$をその焦点、$A$を主頂点、$P$を接点、$PO$を主直径に対する縦線、$PM$を主直径と$M$で交わる接線、$SN$を焦点から接線への垂線とする。$AN$を結ぶ。$MS$と$SP$、$MN$と$NP$、$MA$と$AO$が等しいから、直線$AN$と$OP$は平行であり、したがって三角形$SAN$は$A$で直角であり、等しい三角形$SMN$、$SPN$に相似である。ゆえに$PS : SN = SN : SA$である。\QED

\textit{系1.}\wsp{}$PS^2 : SN^2 = PS : SA$である。

\textit{系2.}\wsp{}$SA$が与えられているから、$SN^2$は$PS$に比例する。

\textit{系3.}\wsp{}任意の接線$PM$と、焦点からそれに下ろした垂線$SN$との交点は、放物線に主頂点で接する直線$AN$の上に落ちる。

\condpagelarge{命題XIII. 問題VIII.}

\textit{物体が放物線の周上を動くとする。この図形の焦点に向かう向心力の法則を求めること。}

\pngright{062.png}{}{}{}

補助定理の作図をそのままにし、$P$を放物線の周上の物体とし、物体が次に動く位置$Q$から、$SP$に平行に$QR$を、垂直に$QT$を引き、さらに接線に平行な$Qv$を引いて、直径$YPG$と$v$で、距離$SP$と$x$で交わらせる。三角形$Pxv$と$MSP$が相似で、一方の辺$SM$、$SP$が等しいから、他方の辺$Px$、すなわち$QR$と$Pv$は等しい。ところが円錐曲線論から、縦線$Qv$の平方は、通径と直径の切片$Pv$との矩形に等しく、すなわち（補助定理XIIIにより）矩形$4PS \times Pv$、すなわち$4PS \times QR$に等しい。そして点$P$と$Q$が一致するとき、比$Qv : Qx$は（補助定理8により）等しいという比になる。ゆえにこの場合、$Qx^2$は矩形$4PS \times QR$に等しい。ところで（角$QxT$、$MPS$、$PMO$が等しいから）$Qx^2 : QT^2 = PS^2 : SN^2$であり、すなわち（補助定理XIVの系1により）$PS : AS$、すなわち$4PS \times QR : 4AS \times QR$である。したがって（『原論』第V巻命題9により）$QT^2$と$4AS \times QR$は等しい。これらの等しいものに$\frac{SP^2}{QR}$を掛ければ、$\frac{SP^2 \times QT^2}{QR}$は$SP^2 \times 4AS$に等しくなる。したがって（定理Vの系により）向心力は$SP^2 \times 4AS$に反比例し、すなわち、$4AS$が与えられているから、距離$SP$の二乗に反比例する。\QED

\textit{系I.}\wsp{}最後の三つの命題から、次のことが導かれる。もし任意の物体$P$が、位置$P$から任意の直線$PR$に沿って任意の速さで出発し、同時に、中心からの距離の二乗に反比例する向心力に動かされるならば、この物体は、力の中心に焦点をもつ、ある円錐曲線の上を動く。逆もまた成り立つ。

\textit{系II.}\wsp{}また、物体が位置$P$から出発する速さが、ごく短い時間のうちに微小な線$PR$を描くことができるような速さであり、向心力が同じ時間に同じ物体を距離$QR$だけ動かすことができるようなものであるならば、この物体は、通径が量$\frac{QT^2}{QR}$（微小な線$PR$、$QR$を限りなく小さくしたときの最終の値）である、ある円錐曲線の上を動く。これらの系では、円を楕円に含め、物体が直線的に中心へ落ちる場合は除外している。

\condpagelarge{命題XIV. 定理VI.}

\textit{もし多くの物体が共通の中心のまわりを回り、向心力が中心からの距離の二乗に反比例して減少するならば、私は言う。軌道の通径は、物体が中心へ引いた動径によって同じ時間に描く面積の二乗の比にある。}

実際、問題VIIIの系IIにより、通径$L$は、量$\frac{QT^2}{QR}$、すなわち点$P$と$Q$が一致するときの最終の値に等しい。ところで、与えられた時間における極小の線$QR$は、それを生じさせる向心力に比例し、すなわち（仮定により）$SP^2$に反比例する。したがって$\frac{QT^2}{QR}$は$QT^2 \times SP^2$に比例し、すなわち通径$L$は面積$QT \times SP$の二乗に比例する。\QED

\textit{系.}\wsp{}したがって、楕円の全面積と、それに比例する軸の矩形とは、通径の平方根と周期の時間との比を合成した比にある。

\condpagelarge{命題XV. 定理VII.}

\textit{同じ仮定のもとで、私は言う。楕円における周期の時間は、主軸の3/2乗の比にある。}

\pngright{064.png}{}{}{}

実際、短軸は、長軸（私はこれを主軸と呼ぶ）と通径との比例中項であり、したがって軸の矩形は、通径の平方根と主軸の3/2乗の比とを合成した比にある。ところがこの矩形は、定理VIの系により、通径の平方根と周期の時間そのものとを合成した比にある。両辺から通径の平方根の比を除けば、主軸の3/2乗の比が周期の時間の比に等しいことが残る。\QED

\textit{系.}\wsp{}したがって、楕円における周期の時間は、直径が楕円の長軸に等しい円における周期の時間と同じである。
\condpagelarge{命題XVI. 定理VIII.}

\textit{同じ仮定のもとで、物体のところで軌道に接する直線を引き、共通の焦点からそれらの接線へ垂線を下ろすとき、物体の速さは、垂線に反比例し、通径の平方根に正比例する比に比例する。}（命題X・XIの図を見よ。）

焦点$S$から接線$PR$へ垂線$SY$を下ろす。すると物体$P$の速さは、量$\frac{SY^2}{L}$の平方根に反比例する。なぜなら、その速さは、与えられた微小時間に描かれる極小の弧$PQ$に比例し、したがって（補助定理VIIにより）接線$PR$に比例し、さらに（$PR : QT = SP : SY$であるから）$\frac{SP \times QT}{SY}$に、すなわち$SY$に反比例し$SP \times QT$に正比例する量に比例するからである。そして$SP \times QT$は、与えられた時間に描かれる面積に比例し、それは（定理VIにより）通径の平方根に比例する。\QED

\textit{系1.}\wsp{}通径は、垂線の二乗比と、速さの二乗比とを合成した比になる。

\textit{系2.}\wsp{}共通の焦点から最も遠い位置と最も近い位置での物体の速さは、距離に反比例し、通径の平方根に正比例する比を合成した比になる。このとき垂線は、そのまま距離になるからである。

\textit{系3.}\wsp{}したがって円錐曲線上の、焦点からの距離が最小の位置での速さは、同じ距離を中心から取った円周上の速さに対して、通径とその距離の2倍との比の平方根の比をなす。

\textit{系4.}\wsp{}楕円上を回る物体の、共通の焦点からの平均の距離における速さは、同じ距離の円周上を回る物体の速さと等しい。すなわち（定理IVの系6により）距離の平方根に反比例する。なぜなら、このとき垂線は短半軸であり、短半軸は距離と通径との比例中項だからである。この比に、通径の平方根に正比例する比を逆に合成すれば、距離の平方根に反比例する比ができる。

\textit{系5.}\wsp{}同じ図形、あるいは通径の等しい別の図形では、物体の速さは、焦点から接線へ下ろした垂線に反比例する。

\textit{系6.}\wsp{}放物線では、速さは物体の焦点からの距離の平方根に反比例する。楕円ではこの比より小さく、双曲線では大きい。なぜなら（補助定理XIVの系2により）、焦点から放物線の接線へ下ろした垂線は、距離の平方根に比例するからである。

\textit{系7.}\wsp{}放物線では、どこでも速さは、同じ距離を円周上で回る物体の速さに対して、2と1との比の平方根の比をなす。楕円ではこの比より小さく、双曲線では大きい。なぜなら、この定理の系2により、放物線の頂点での速さはこの比になり、この定理の系6と定理IVにより、すべての距離で同じ比が保たれるからである。したがって放物線では、どこでも速さは、距離の半分の円周上を回る物体の速さに等しく、楕円ではそれより小さく、双曲線では大きい。

\textit{系8.}\wsp{}どんな円錐曲線でも、回る物体の速さは、その曲線の通径の半分の距離にある円周上を回る物体の速さに対して、その距離が、焦点から曲線の接線へ下ろした垂線に対するのと同じ比をなす。これは系5から明らかである。

\textit{系9.}\wsp{}ところで（定理IVの系6により）この円周上を回る物体の速さは、他の任意の円周上を回る物体の速さに対して、距離の平方根に反比例する比をなすから、等比によって、円錐曲線上を回る物体の速さは、同じ距離の円周上を回る物体の速さに対して、共通の距離と曲線の通径の半分との比例中項が、共通の焦点から接線へ下ろした垂線に対するのと同じ比をなす。

\condpagelarge{命題XVII. 問題IX.}

\textit{向心力が中心からの距離の2乗に反比例し、その力の絶対量が知られているとする。与えられた位置から、与えられた速さで、与えられた直線に沿って出発する物体が描く曲線を求めよ。}

\pngright{067.png}{}{}{}

点$S$に向かう向心力が、与えられた任意の軌道$pq$上で物体$p$を回らせるものであり、その物体の位置$p$での速さが知られているとする。位置$P$から直線$PR$に沿って、与えられた速さで物体$P$が出発し、やがて向心力に押されて円錐曲線$PQ$の方へ曲がるとする。すると直線$PR$はこの曲線に点$P$で接する。同様に、ある直線$pr$が軌道$pq$に点$p$で接するとし、$S$からそれらの接線へ垂線が下ろされたと考えれば、（定理VIIIの系1により）求める円錐曲線の通径と、与えられた軌道の通径との比は、垂線の二乗比と速さの二乗比とを合成した比になり、したがって与えられる。それを$L$とする。さらに円錐曲線の焦点$S$も与えられている。

角$RPS$の二直角に対する補角を角$RPH$とすれば、もう一つの焦点$H$がのっている直線$PH$が、位置で与えられる。$PH$へ垂線$SK$を下ろし、共役半軸$BC$を立てると、（『原論』第II巻命題13により）
\[
SP^2 - 2\,KP \times PH + PH^2 = SH^2 = 4CH^2 = 4BH^2 - 4BC^2
= (SP + PH)^2 - L \times (SP + PH)
\]
\[
= SP^2 + 2\,SP \times PH + PH^2 - L \times (SP + PH)
\]
となる。両辺に$2\,KP \times PH + L \times (SP + PH) - SP^2 - PH^2$を加えると、
\[
L \times (SP + PH) = 2\,SP \times PH + 2\,KP \times PH
\]
となる。すなわち$SP + PH$と$PH$との比は、$2SP + 2KP$と$L$との比に等しい。これから$PH$は、長さも位置も与えられる。

すなわち、$P$での物体の速さが、通径$L$が$2SP + 2KP$より小さくなるようなものであれば、$PH$は接線$PR$に関して直線$PS$と同じ側にあり、したがって図形は楕円で、与えられた焦点$S$、$H$と、主軸$SP + PH$から決まる。もし速さがちょうど、通径$L$が$2SP + 2KP$に等しくなるようなものであれば、$PH$の長さは無限大になり、図形は直線$PK$に平行な軸$SH$をもつ放物線となって、それから決まる。さらに物体がそれより大きな速さで位置$P$を出発するならば、$PH$は接線の反対側に取らなければならない。そのとき接線は二つの焦点の間を通るので、図形は、$SP$と$PH$との差に等しい主軸をもつ双曲線となって、それから決まる。\QED

\textit{系1.}\wsp{}したがってどんな円錐曲線でも、主頂点$D$、通径$L$、一方の焦点$S$が与えられれば、$DH$と$DS$との比を、通径と、通径から$4DS$を引いた差との比にとることによって、もう一方の焦点$H$が与えられる。なぜなら、この系の場合、$SP + PH$と$PH$との比が$2SP$と$L$との比になることから、$DS + DH$と$DH$との比が$4DS$と$L$との比になり、分割して、$DS$と$DH$との比が$4DS - L$と$L$との比になるからである。

\textit{系2.}\wsp{}したがって主頂点$D$での物体の速さが与えられれば、軌道は簡単に求められる。すなわち、その通径を、距離$DS$の2倍に対して、この与えられた速さと、距離$DS$の円周上を回る物体の速さとの比の二乗の比にとり（定理VIIIの系3による）、次に$DH$と$DS$との比を、通径と、通径から$4DS$を引いた差との比にとればよい。

\textit{系3.}\wsp{}また、物体が任意の円錐曲線上を動いていて、何かの衝撃によって軌道からはじき出されたとき、そのあとで物体が進む軌道を知ることができる。物体自身の運動と、その衝撃だけが生む運動とを合成すれば、衝撃を受けた位置から、位置で与えられた直線に沿って物体が出てゆく運動が得られるからである。

\textit{系4.}\wsp{}さらに物体が外から加えられる何らかの力によって絶えず乱されるときも、その力がいくつかの点で生じさせる変化を集め、その並びの規則にならって中間の位置での連続的な変化を見積もれば、その進路はほぼ分かる。

\chapter{第1編 物体の運動（第4節〜第7節：円錐曲線の決定と直線昇降運動）}



\section{第1編 第4節：焦点が与えられたときの楕円・放物線・双曲線軌道の決定（命題XVIII〜XXI）}


% Section: 第1編 第4節：焦点が与えられたときの楕円・放物線・双曲線軌道の決定（命題XVIII〜XXI）



\begin{center}\textit{与えられた焦点から、楕円・放物線・双曲線の軌道を求めること}\end{center}

\begin{center}\large \textit{補助定理XV.}\end{center}

\pngright{069.png}{}{}{}

\textit{楕円または双曲線の二つの焦点$S$、$H$から、任意の第三の点$V$へ二本の直線$SV$、$HV$を折れ線のように引く。このうち$HV$は図形の主軸に等しく、もう一方の$SV$は、そこへ下ろした垂線$TR$によって点$T$で二等分されるとする。このとき垂線$TR$は、その円錐曲線のどこかに接する。逆に、接するならば、$VH$は図形の主軸に等しい。}

実際、$VH$が円錐曲線と点$R$で交わるとし、$SR$を結ぶ。$TS$と$TV$が等しいから、角$TRS$と角$TRV$は等しい。したがって$RT$は角$VRS$を二等分し、ゆえに図形に接する。逆もまた同じである。\QED

\condpagelarge{命題XVIII. 問題X.}

\textit{焦点と主軸の長さが与えられたとき、与えられた点を通り、位置の与えられた直線に接する、楕円または双曲線の軌道を描くこと。}

\pngright{070a.png}{}{}{}

$S$をこれらの図形に共通の焦点、$AB$を任意の軌道の主軸の長さ、$P$を軌道が通るべき点、$TR$を軌道が接すべき直線とする。$P$を中心に、軌道が楕円なら半径$AB - SP$で、双曲線なら半径$AB + SP$で、円$HG$を描く。接線$TR$へ垂線$ST$を下ろし、それを$V$まで延ばして$TV$が$ST$に等しくなるようにする。$V$を中心に半径$AB$で円$FH$を描く。この方法によって、二点$P$、$p$が与えられても、二本の接線$TR$、$tr$が与えられても、点$P$と接線$TR$が与えられても、二つの円を描くことになる。それらの共通の交点を$H$とし、焦点$S$、$H$と、与えられた主軸とによって軌道を描く。

これでよいと私は言う。実際、描かれた軌道は（楕円では$PH + SP$が、双曲線では$PH - SP$が主軸に等しいから）点$P$を通り、（先の補助定理によって）直線$TR$に接する。同じ論法で、この軌道は二点$P$、$p$を通るか、二本の接線$TR$、$tr$に接する。\QED

\condpagelarge{命題XIX. 問題XI.}

\textit{与えられた焦点のまわりに、与えられた点を通り、位置の与えられた直線に接する放物線の軌道を描くこと。}

\pngright{070b.png}{}{}{}

$S$を焦点、$P$を点、$TR$を描こうとする軌道の接線とする。$P$を中心に半径$PS$で円$FG$を描く。焦点から接線へ垂線$ST$を下ろし、それを$V$まで延ばして$TV$が$ST$に等しくなるようにする。もう一つの点$p$が与えられているなら、同様にもう一つの円$fg$を描かなければならない。もう一つの接線$tr$が与えられているなら、もう一つの点$v$を求めなければならない。次に、二点$P$、$p$が与えられているときは二つの円$FG$、$fg$に接する直線$IF$を引き、二本の接線$TR$、$tr$が与えられているときは二点$V$、$v$を通る直線を引き、点$P$と接線$TR$が与えられているときは円$FG$に接して点$V$を通る直線を引く。

$FI$へ垂線$SI$を下ろし、それを$K$で二等分して、軸を$SK$、主頂点を$K$として放物線を描く。これでよいと私は言う。実際、$SK$と$IK$が等しく、$SP$と$FP$が等しいから、放物線は点$P$を通る。また（補助定理XIVの系3により）$ST$と$TV$が等しく、角$STR$が直角であるから、直線$TR$に接する。\QED

\condpagelarge{命題XX. 問題XII.}

\textit{与えられた焦点のまわりに、種類（形）の与えられた任意の軌道で、与えられた点を通り、位置の与えられた直線に接するものを描くこと。}

\pngright{071.png}{}{}{}

\textit{場合1.}\wsp{}焦点$S$が与えられ、二点$B$、$C$を通る軌道$ABC$を描くとする。軌道の形が与えられているから、主軸と焦点間の距離との比が与えられる。その比で$KB$と$BS$、$LC$と$CS$をとる。$B$、$C$を中心に半径$BK$、$CL$で二つの円を描き、これらに$K$、$L$で接する直線$KL$へ垂線$SG$を下ろして、それを点$A$、$a$で、$SA$と$AG$の比、$Sa$と$aG$の比が、どちらも$SB$と$BK$の比に等しくなるように分ける。そして軸を$Aa$、頂点を$A$、$a$として軌道を描く。

これでよいと私は言う。実際、描いた図形のもう一つの焦点を$H$とする。$SA : AG = Sa : aG$であるから、分割して、$Sa - SA$すなわち$SH$と、$aG - AG$すなわち$Aa$との比は同じ比になり、したがって描くべき図形の主軸と焦点間の距離との比に等しい。ゆえに描いた図形は、描くべきものと同じ種類である。また$KB : BS$と$LC : CS$が同じ比であるから、この図形は円錐曲線の性質から明らかなように、点$B$、$C$を通る。\QED

\pngright{072.png}{}{}{}

\textit{場合2.}\wsp{}焦点$S$が与えられ、二本の直線$TR$、$tr$にどこかで接する軌道を描くとする。焦点から接線へ垂線$ST$、$St$を下ろし、それらを$V$、$v$まで延ばして、$TV$、$tv$が$TS$、$ts$に等しくなるようにする。$Vv$を点$O$で二等分し、垂線$OH$を限りなく立てる。直線$VS$を限りなく延ばして、$VK$と$KS$の比、$Vk$と$kS$の比が、描くべき軌道の主軸と焦点間の距離との比に等しくなるように、点$K$、$k$で分ける。$Kk$を直径とする円を描き、これが直線$OH$と交わる点を$H$とする。焦点$S$、$H$と、$VH$に等しい主軸とによって軌道を描く。

これでよいと私は言う。実際、$Kk$を$X$で二等分し、$HX$、$HS$、$HV$、$Hv$を結ぶ。$VK : KS = Vk : kS$であるから、合比により$VK + Vk : KS + kS$、分割により$Vk - VK : kS - KS$、すなわち$2VX : 2KX$であり、$2KX : 2SX$でもある。したがって$VX : HX = HX : SX$となって、三角形$VXH$と$HXS$は相似である。ゆえに$VH : SH = VX : XH$、したがって$= VK : KS$となる。よって描いた軌道の主軸$VH$は、その焦点間の距離$SH$に対して、描くべき軌道の主軸がその焦点間の距離に対してもつのと同じ比をもつ。したがって同じ種類である。さらに$VH$、$vH$が主軸に等しく、$VS$、$vS$が直線$TR$、$tr$によって垂直に二等分されるから、補助定理XVにより、それらの直線が描いた軌道に接することは明らかである。\QED

\pngright{073.png}{}{}{}

\textit{場合3.}\wsp{}焦点$S$が与えられ、与えられた点$R$で直線$TR$に接する軌道を描くとする。直線$TR$へ垂線$ST$を下ろし、それを$V$まで延ばして$TV$が$ST$に等しくなるようにする。$VR$を結び、直線$VS$を限りなく延ばして、$VK : SK$と$Vk : Sk$が、描くべき楕円の主軸と焦点間の距離との比に等しくなるように、点$K$、$k$で分ける。$Kk$を直径とする円を描いて、延長した直線$VR$がこれと交わる点を$H$とし、焦点$S$、$H$と、直線$HV$に等しい主軸とによって軌道を描く。

これでよいと私は言う。実際、$VH : SH = VK : SK$であり、したがって軌道の主軸と焦点間の距離との比が、描くべき軌道のそれに等しいことは、場合2で示したことから明らかである。ゆえに描いた軌道は描くべきものと同じ種類である。また角$VRS$を二等分する直線$TR$が、点$R$で軌道に接することは、円錐曲線論から明らかである。\QED

\pngright{074.png}{}{}{}

\textit{場合4.}\wsp{}さて焦点$S$のまわりに軌道$APB$を描くとする。この軌道は直線$TR$に接し、接線の外の与えられた任意の点$P$を通り、かつ、主軸$ab$と焦点$s$、$h$をもつ図形$apb$に相似であるとする。接線$TR$へ垂線$ST$を下ろし、それを$V$まで延ばして、$TV$が$ST$に等しくなるようにする。角$VSP$、$SVP$に、それぞれ等しい角$hsq$、$shq$をつくる。$q$を中心に、半径$ab$に対して$SP$が$VS$に対するのと同じ比になる長さで円を描き、図形$apb$と点$p$で交わらせる。$sp$を結び、$SH$を、$sh$に対して$SP$が$sp$に対するのと同じ比になるように引き、角$PSH$を角$psh$に、角$VSH$を角$psq$に等しくなるようにする。最後に、焦点$S$、$H$と、距離$VH$に等しい主軸とによって円錐曲線を描く。

これでよいと私は言う。実際、$sv$を引いて、$sp$に対して$sh$が$sq$に対するのと同じ比になり、角$vsp$を角$hsq$に、角$vsh$を角$psq$に等しくするとき、三角形$svh$と$spq$は相似になる。したがって$vh : pq = sh : sq$、すなわち（三角形$VSP$と$hsq$が相似であるから）$VS : SP$、つまり$ab : pq$となる。ゆえに$vh$と$ab$は等しい。さらに三角形$VSH$と$vsh$が相似であるから、$VH : SH = vh : sh$、すなわち今描いた円錐曲線の主軸と焦点間隔との比は、主軸$ab$と焦点間隔$sh$との比に等しい。したがって今描いた図形は図形$apb$に相似である。この図形は、三角形$PSH$が三角形$psh$に相似であるから、点$P$を通る。また$VH$がその主軸に等しく、$VS$が直線$TR$によって垂直に二等分されるから、この図形は直線$TR$に接する。\QED

\condpagelarge{補助定理XVI.}

\textit{与えられた三つの点から、与えられていない第四の点へ、三本の直線を引く。その差が与えられているか、または零である（等しい）とき、その点を求めること。}

\pngright{075.png}{}{}{}

\textit{場合1.}\wsp{}与えられた三点を$A$、$B$、$C$とし、求める第四の点を$Z$とする。直線$AZ$と$BZ$の差が与えられているから、点$Z$は、$A$、$B$を焦点とし、その与えられた差を主軸とする双曲線の上にある。その主軸を$MN$とする。$MA$に対して$MN$が$AB$に対するのと同じ比になるように$PM$をとり、$AB$に垂直に$PR$を立て、$PR$へ垂線$ZR$を下ろす。するとこの双曲線の性質から、$ZR : AZ = MN : AB$となる。同じ論法で、点$Z$は、$A$、$C$を焦点とし、$AZ$と$CZ$の差を主軸とするもう一つの双曲線の上にもある。$AC$に垂直な$QS$を引くことができ、この双曲線の任意の点$Z$から$QS$へ垂線$ZS$を下ろすと、それは$AZ$に対して、$AZ$と$CZ$の差が$AC$に対するのと同じ比になる。こうして$ZR$と$ZS$が$AZ$に対してもつ比は与えられ、したがって$ZR$と$ZS$の比も与えられる。ゆえに直線$RP$、$SQ$が点$T$で交わるとき、点$Z$は位置の与えられた直線$TZ$の上にある。同じ方法で、$B$、$C$を焦点とし、$BZ$と$CZ$の差を主軸とする第三の双曲線によって、点$Z$がのっている別の直線も見つけることができる。こうして二本の直線の軌跡が得られれば、求める点$Z$は、それらの交点として得られる。\QED

\textit{場合2.}\wsp{}三本の線のうち二本、たとえば$AZ$と$BZ$が等しいときは、点$Z$は距離$AB$を二等分する垂線の上にあり、もう一つの直線の軌跡は上と同様に見つかる。\QED

\textit{場合3.}\wsp{}三本とも等しいときは、点$Z$は、点$A$、$B$、$C$を通る円の中心にある。\QED

この補助定理は、ヴィエタが復元したアポロニウスの『接触論』によっても、問題として解かれる。

\condpagelarge{命題XXI. 問題XIII.}

\textit{与えられた焦点のまわりに、与えられた点を通り、位置の与えられた直線に接する軌道を描くこと。}

\pngright{076.png}{}{}{}

焦点$S$、点$P$、接線$TR$が与えられ、もう一つの焦点$H$を求めるとする。接線へ垂線$ST$を下ろし、それを$Y$まで延ばして$TY$が$ST$に等しくなるようにすると、$YH$は主軸に等しい。$SP$、$HP$を結べば、$SP$は$HP$と主軸との差になる。この方法で、より多くの接線$TR$、あるいはより多くの点$P$が与えられても、つねに、これらの点$Y$または$P$から焦点$H$へ引いた同じ数の線$YH$または$PH$に行き着く。それらは主軸に等しいか、与えられた長さ$SP$だけ主軸と異なる。したがって互いに等しいか、与えられた差をもつ。こうして先の補助定理により、そのもう一つの焦点$H$が与えられる。焦点が主軸の長さ（それは$YH$であるか、軌道が楕円なら$PH + SP$、双曲線なら$PH - SP$である）とともに得られれば、軌道が得られる。\QED

\condpagelarge{\textit{注.}}

\pngright{077.png}{}{}{}

三つの点が与えられた場合は、次のようにすればより手早く解ける。与えられた点を$B$、$C$、$D$とする。結んだ$BC$、$CD$を$E$、$F$まで延ばし、$EB : EC = SB : SC$、$FC : FD = SC : SD$となるようにする。引いて延ばした$EF$へ垂線$SG$、$BH$を下ろし、限りなく延ばした$GS$の上に、$GA : AS$と$Ga : aS$がどちらも$HB : BS$に等しくなるように$A$、$a$をとる。すると$A$が頂点で、$Aa$が軌道の主軸になる。この軌道は、$GA$が$AS$より小さいか、等しいか、大きいかに応じて、楕円、放物線、双曲線になる。点$a$は、最初の場合には点$A$と直線$GK$の同じ側に落ち、第二の場合には無限遠へ去り、第三の場合には直線$GK$の反対側に落ちる。

実際、$GF$へ垂線$CI$、$DK$を下ろせば、$IC : HB = EC : EB$、すなわち$SC : SB$となる。入れ替えて$IC : SC = HB : SB$、すなわち$GA : SA$となる。同様の論法で、$KD : SD$も同じ比であることが示される。したがって点$B$、$C$、$D$は、焦点$S$のまわりに次のように描いた円錐曲線の上にある。すなわち、焦点$S$から曲線の各点へ引いたすべての直線は、同じ点から直線$GK$へ下ろした垂線に対して、その与えられた比をもつ。

これとあまり違わない方法でこの問題の解法を示しているのは、著名な幾何学者ド・ラ・イール氏で、その『円錐曲線論』第VIII巻命題XXVにある。
\section{第1編 第5節：焦点がいずれも与えられないときの軌道の決定（命題XXII〜XXXI）}


% Section: 第1編 第5節：焦点がいずれも与えられないときの軌道の決定（命題XXII〜XXXI）





\begin{center}\textit{焦点がどちらも与えられないときの軌道を求めること}\end{center}

\begin{center}{\large \textit{補助定理XVII.}}\end{center}

\textit{与えられた円錐曲線上の任意の点$P$から、その円錐曲線に内接する四辺形$ABCD$の四つの辺$AB$、$CD$、$AC$、$DB$（限りなく延長したもの）へ、それぞれ与えられた角で四本の直線$PQ$、$PR$、$PS$、$PT$を、各辺に一本ずつ引く。このとき、向かい合う二辺へ引いた二本の積$PQ \times PR$は、残りの向かい合う二辺へ引いた二本の積$PS \times PT$に対して、一定の比をもつ。}

\pngright{078.png}{}{}{}

\textit{場合1.}\wsp{}まず、向かい合う辺へ引いた線が、残りの辺のどちらかに平行である場合を考える。すなわち$PQ$と$PR$は辺$AC$に平行、$PS$と$PT$は辺$AB$に平行とする。さらに、向かい合う辺のうち二辺、たとえば$AC$と$BD$が平行であるとする。そのとき、これらの平行な辺を二等分する直線は円錐曲線の直径の一つであり、$RQ$をも二等分する。$RQ$が二等分される点を$O$とすれば、$PO$はその直径に対する縦線（ordinatim applicata）である。$PO$を$K$まで延ばして$OK$を$PO$に等しくすれば、$OK$は直径の反対側の縦線になる。したがって点$A$、$B$、$P$、$K$は円錐曲線上にあり、$PR$が$AB$と与えられた角で交わるから、（アポロニウスの第III巻命題17・18により）矩形$PQK$と矩形$AQB$とは一定の比をもつ。ところが$QK$と$PR$は等しい。なぜなら$OK$と$OP$が等しく、$OQ$と$OR$が等しいから、それらの差が等しいからである。したがって矩形$PQK$と矩形$PQ \times PR$は等しい。ゆえに矩形$PQ \times PR$は、矩形$AQB$に対して、すなわち矩形$PS \times PT$に対して、一定の比をもつ。\QED

\pngright{079a.png}{}{}{}

\textit{場合2.}\wsp{}次に、四辺形の向かい合う辺$AC$と$BD$が平行でないとする。$AC$に平行な$Bd$を引き、これが直線$ST$と$t$で、円錐曲線と$d$で交わるようにする。$Cd$を結んで$PQ$と$r$で交わらせ、$PQ$に平行な$DM$を引いて、$Cd$と$M$で、$AB$と$N$で交わらせる。三角形$BTt$と$DBN$が相似であるから、$Bt$すなわち$PQ$と$Tt$との比は、$DN$と$NB$との比に等しい。同様に$Rr$と$AQ$すなわち$PS$との比は、$DM$と$AN$との比に等しい。そこで前項どうし、後項どうしを掛けると、矩形$PQ \times Rr$と矩形$Tt \times PS$との比は、矩形$NDM$と矩形$ANB$との比に等しい。そして（場合1により）矩形$QPr$と矩形$SPt$との比もこれに等しく、分割すれば、矩形$QPR$と矩形$PS \times PT$との比もこれに等しい。\QED

\newpage

\pngright{079b.png}{}{}{}

\textit{場合3.}\wsp{}最後に、四本の線$PQ$、$PR$、$PS$、$PT$が、辺$AC$、$AB$に平行でなく、それらに対して任意に傾いているとする。そのときは、$AC$に平行な$Pq$、$Pr$を引き、$AB$に平行な$Ps$、$Pt$を引く。三角形$PQq$、$PRr$、$PSs$、$PTt$の角が与えられているから、比$PQ : Pq$、$PR : Pr$、$PS : Ps$、$PT : Pt$が与えられ、したがって合成した比$PQ \times PR : Pq \times Pr$と$PS \times PT : Ps \times Pt$も与えられる。ところが先に示したことにより、比$Pq \times Pr : Ps \times Pt$は与えられている。ゆえに比$PQ \times PR : PS \times PT$も与えられる。\QED

\condpagelarge{\textit{補助定理XVIII.}}

\textit{同じ仮定のもとで、四辺形の向かい合う二辺へ引いた線の積$PQ \times PR$が、残りの二辺へ引いた線の積$PS \times PT$に対して一定の比をもつならば、それらの線を引く点$P$は、四辺形のまわりに描いた円錐曲線上にある。}

\pngright{080.png}{}{}{}

点$A$、$B$、$C$、$D$と、無数の点$P$のうちの一つ、たとえば$p$を通る円錐曲線が描かれたと考える。私は、点$P$がつねにこの曲線上にあると言う。もしそうでないとし、$AP$を結んで、これがこの円錐曲線と$P$とは別の点、たとえば$b$で交わるとしよう。すると、これらの点$p$、$b$から四辺形の辺へ、与えられた角で直線$pq$、$pr$、$ps$、$pt$と$bk$、$b\textrm{r}$、$b\longs$、$bd$を引くとき、（補助定理XVIIにより）$bk \times b\textrm{r} : bd \times b\longs$は$pq \times pr : ps \times pt$に等しく、（仮定により）それは$PQ \times PR : PS \times PT$に等しい。また、四辺形$bkA\longs$と$PQAS$が相似であるから、$bk : b\longs = PQ : PS$である。したがって前の命題の項を、この命題の対応する項にあてはめれば、$b\textrm{r} : bd = PR : PT$となる。ゆえに等角の四辺形$D\textrm{r}bd$と$DRPT$は相似であり、したがってそれらの対角線$Db$、$DP$は一致する。そこで$b$は直線$AP$、$DP$の交点にあり、したがって点$P$と一致する。ゆえに点$P$は、どこにとっても、指定した円錐曲線上にある。\QED

\textit{系.}\wsp{}したがって、共通の点$P$から、位置の与えられた三本の直線$AB$、$CD$、$AC$へ、それぞれ与えられた角で三本の直線$PQ$、$PR$、$PS$を引き、そのうち二本の積$PQ \times PR$が第三のものの平方$PS^2$に対して一定の比をもつならば、これらの線を引く点$P$は、直線$AB$、$CD$に点$A$、$C$で接する円錐曲線上にあり、逆もまた成り立つ。実際、三本の直線$AB$、$CD$、$AC$の位置を保ったまま、直線$BD$を直線$AC$に重ね、さらに直線$PT$を直線$PS$に重ねる。すると矩形$PS \times PT$は$PS^2$になる。また、曲線を点$A$と$B$、$C$と$D$で切っていた直線$AB$、$CD$は、それらの点が一致したので、もはや曲線を切ることはできず、接するだけになる。

\condpagelarge{\textit{注.}}

この補助定理では円錐曲線という名を広い意味で用いており、円錐の頂点を通る直線の切断も、底面に平行な円の切断も含まれる。実際、点$p$が、四点$A$、$B$、$C$、$D$のうち任意の二点を結ぶ直線の上に落ちるとき、円錐曲線は一対の直線に変わり、その一方は点$p$が落ちたその直線、もう一方は他の二点を結ぶ直線になる。

四辺形の向かい合う二つの角の和が二直角に等しく、その辺へ四本の線$PQ$、$PR$、$PS$、$PT$を垂直に、あるいは互いに等しい任意の角で引き、引いた二本の積$PQ \times PR$が他の二本の積$PS \times PT$に等しいならば、円錐曲線は円になる。同じことは、四本の線を任意の角で引いても起こる。すなわち、引いた二本の積$PQ \times PR$が他の二本の積$PS \times PT$に対して、後の二本$PS$、$PT$を引く角$S$、$T$の正弦の積が、先の二本$PQ$、$PR$を引く角$Q$、$R$の正弦の積に対するのと同じ比になる場合である。その他の場合には、点$P$の軌跡は、ふつう円錐曲線と呼ばれる三つの図形のどれかになる。

また四辺形$ABCD$の代わりに、向かい合う二辺が対角線のように互いに交わる四角形を用いることもできる。さらに四点$A$、$B$、$C$、$D$のうち一つまたは二つが無限遠に去ることもあり、そのときには、それらの点へ集まる図形の辺が平行になる。その場合、円錐曲線は残りの点を通り、平行線の向きに無限遠へ去る。

\condpagelarge{補助定理XIX.}

\textit{点$P$を求めること。すなわち、そこから位置の与えられた他の四本の直線$AB$、$CD$、$AC$、$BD$へ、それぞれ与えられた角で四本の直線$PQ$、$PR$、$PS$、$PT$を引くとき、引いた二本の積$PQ \times PR$が、他の二本の積$PS \times PT$に対して一定の比をもつような点を。}

\pngright{082.png}{}{}{}

積の一方をつくる二本$PQ$、$PR$を引く直線$AB$、$CD$が、位置の与えられた他の二直線と、点$A$、$B$、$C$、$D$で交わるとする。そのうち一つ$A$から、点$P$をそこに見いだしたい任意の直線$AH$を引く。これが向かい合う直線$BD$、$CD$と交わる点を、$BD$と$H$、$CD$と$I$とする。図形の角がすべて与えられているから、比$PQ : PA$と$PA : PS$が与えられ、したがって比$PQ : PS$が与えられる。これを、与えられた比$PQ \times PR : PS \times PT$から取り除けば、比$PR : PT$が与えられる。さらに、与えられた比$PI : PR$と$PT : PH$を加えれば、比$PI : PH$が与えられ、したがって点$P$が与えられる。\QED

\textit{系1.}\wsp{}したがって、無数の点$P$の軌跡の任意の点$D$へ接線を引くこともできる。実際、点$P$と$D$が一致するとき、すなわち$AH$が点$D$を通るとき、弦$PD$は接線になる。この場合、消えてゆく$IP$と$PH$の最終の比は、上のようにして求められる。そこで$AD$に平行な$CF$を引いて、$BD$と$F$で交わらせ、これをその最終の比で$E$において分ければ、$DE$が接線になる。なぜなら、$CF$と、消えてゆく$IH$とは平行であり、$E$と$P$において相似に分けられているからである。

\pngright{083.png}{}{}{}

\textit{系2.}\wsp{}したがって、すべての点$P$の軌跡を決めることもできる。点$A$、$B$、$C$、$D$のうち任意の一つ、たとえば$A$を通って軌跡の接線$AE$を引き、別の任意の点$B$を通って接線に平行な$BF$を引き、軌跡と$F$で交わらせる。点$F$は先の補助定理によって見つかる。$BF$を$G$で二等分すれば、引いた$AG$は直径となり、これに対して$BG$と$FG$が縦線になる。この$AG$が軌跡と$H$で交わるとすれば、$AH$は横辺（latus transversum）であり、これに対する通径（latus rectum）は、$BG^2 : AG \times GH$の比で決まる。もし$AG$が軌跡とどこでも交わらず、線$AH$が無限であるならば、軌跡は放物線で、その通径は$\frac{BG^2}{AG}$である。もしどこかで交わるならば、点$A$と$H$が$G$の同じ側にあるとき軌跡は双曲線であり、$G$がその間にあるとき楕円である。ただし、もし角$AGB$が直角であって、しかも$BG^2$が矩形$AGH$に等しい場合には、円が得られる。

このようにして、ユークリッドが始め、アポロニウスが続けた、四本の線についての古人の問題は、計算ではなく、古人が求めたような幾何学的な作図として、この系で示される。

\condpagelarge{補助定理XX.}

\textit{平行四辺形$ASPQ$の向かい合う二つの角$A$、$P$が、任意の円錐曲線に点$A$、$P$で接し、それらの角の一方の辺$AQ$、$AS$を限りなく延ばしたものが、同じ円錐曲線と$B$、$C$で交わるとする。交わった点$B$、$C$から、円錐曲線上の任意の第五の点$D$へ二本の直線$BD$、$CD$を引き、これらが平行四辺形の他の二辺$PS$、$PQ$（限りなく延ばしたもの）と$T$、$R$で交わるとする。このとき、切りとられた辺の部分$PR$と$PT$は、つねに互いに一定の比をもつ。逆に、それらの切りとられた部分が互いに一定の比をもつならば、点$D$は、四点$A$、$B$、$P$、$C$を通る円錐曲線上にある。}

\pngright{084.png}{}{}{}

\textit{場合1.}\wsp{}$BP$、$CP$を結び、点$D$から二本の直線$DG$、$DE$を引く。そのうち$DG$は$AB$に平行で、$PB$、$PQ$、$CA$と$H$、$I$、$G$で交わるとし、$DE$は$AC$に平行で、$PC$、$PS$、$AB$と$F$、$K$、$E$で交わるとする。すると（補助定理XVIIにより）矩形$DE \times DF$と矩形$DG \times DH$との比は一定である。ところが$PQ : DE$すなわち$IQ$は$PB : HB$に等しく、したがって$PT : DH$に等しい。入れ替えて$PQ : PT = DE : DH$である。また$PR : DF = RC : DC$であり、したがって$IG$または$PS$と$DG$の比に等しい。入れ替えて$PR : PS = DF : DG$である。これらの比を合成すれば、矩形$PQ \times PR$と矩形$PS \times PT$との比は、矩形$DE \times DF$と矩形$DG \times DH$との比に等しくなり、したがって一定の比である。ところが$PQ$と$PS$は与えられている。ゆえに比$PR : PT$が与えられる。\QED

\textit{場合2.}\wsp{}逆に、$PR$と$PT$が互いに一定の比にあると仮定するなら、同様の推論をさかのぼってゆくことで、矩形$DE \times DF$と矩形$DG \times DH$とが一定の比にあることになり、したがって点$D$は（補助定理XVIIIにより）点$A$、$B$、$P$、$C$を通る円錐曲線上にあることになる。\QED

\textit{系1.}\wsp{}したがって、$BC$を引いて$PQ$と$r$で交わらせ、$PT$の上に、$Pr$に対して$PT$が$PR$に対するのと同じ比になる$Pt$をとれば、$Bt$は円錐曲線の点$B$での接線になる。実際、点$D$が点$B$に近づいて一致し、弦$BD$が消えるとき、$BT$が接線になり、$CD$と$BT$は$CB$と$Bt$に一致する、と考えればよい。

\textit{系2.}\wsp{}逆に、$Bt$が接線であって、円錐曲線の任意の点$D$へ$BD$、$CD$が集まるならば、$PR : PT = Pr : Pt$となる。そして逆に、$PR : PT = Pr : Pt$ならば、$BD$、$CD$は円錐曲線のどこかの点$D$に集まる。

\textit{系3.}\wsp{}一つの円錐曲線が別の円錐曲線を切る点は、四つより多くない。なぜなら、もし可能だとして、二つの円錐曲線が五つの点$A$、$B$、$C$、$D$、$P$を通るとし、直線$BD$がそれらを点$D$、$d$で切り、直線$Cd$が$PQ$を$\mathrm{r}$で切るとする。すると$PR : PT = P\mathrm{r} : PT$となり、すなわち$PR$と$P\mathrm{r}$が等しくなって、仮定に反するからである。

\condpagelarge{補助定理XXI.}

\textit{動く無限の二直線$BM$、$CM$が、与えられた点$B$、$C$を極として引かれ、それらの交点$M$が、位置の与えられた第三の直線$MN$を描くとする。また別の動く無限の二直線$BD$、$CD$が、先の二直線と、与えられた点$B$、$C$において与えられた角$MBD$、$MCD$をなすように引かれるとする。私は言う。これら二直線$BD$、$CD$は、その交点$D$によって円錐曲線を描く。逆に、直線$BD$、$CD$が、その交点$D$によって、点$B$、$C$、$A$を通る円錐曲線を描き、それらの交点が曲線のある点$A$に落ちるとき、他の二直線$BM$、$CM$が直線$BC$に重なるならば、点$M$は位置の与えられた直線に接する。}

\pngright{086.png}{}{}{}

実際、直線$MN$上に点$N$を与えておき、動点$M$が定点$N$に落ちるとき、動点$D$が定点$P$に落ちるとする。$CN$、$BN$、$CP$、$BP$を結び、点$P$から直線$PT$、$PR$を引いて、$BD$、$CD$と$T$、$R$で交わらせ、角$BPT$を角$BNM$に、角$CPR$を角$CNM$に等しくする。仮定により角$MBD$と角$NBP$が等しく、角$MCD$と角$NCP$も等しいから、共通部分$NBD$と$MCP$を除けば、角$NBM$と$PBT$が等しく、角$NCM$と$PCR$が等しい。したがって三角形$NBM$と$PBT$は相似であり、三角形$NCM$と$PCR$も相似である。ゆえに$PT : NM = PB : NB$であり、$PR : NM = PC : NC$である。よって$PT$と$PR$は$NM$に対して一定の比をもち、したがって互いに一定の比をもつ。ゆえに補助定理XXにより、点$P$（動く直線$BT$と$CR$の絶えず変わる交点）は円錐曲線に接する。\QED

逆に、点$D$が点$B$、$C$、$A$を通る円錐曲線に接し、直線$BM$、$CM$が直線$BC$に重なるときに点$D$が曲線のある点$A$に落ちるとする。そして点$D$が曲線の他の任意の二点$p$、$P$に順に落ちるとき、動点$M$が定点$n$、$N$に順に落ちるとする。その$n$、$N$を通って直線$nN$を引けば、これが動点$M$の永続的な軌跡になる。実際、もし可能だとして、点$M$がある曲線の上を動くとしよう。すると、点$M$が絶えず曲線に接するとき、点$D$は五つの点$C$、$p$、$P$、$B$、$A$を通る円錐曲線に接する。しかしすでに示したことから、点$M$が絶えず直線に接するときにも、点$D$は同じ五つの点$C$、$p$、$P$、$B$、$A$を通る円錐曲線に接する。したがって二つの円錐曲線が同じ五つの点を通ることになり、補助定理XXの系3に反する。ゆえに点$M$が曲線上を動くとするのは不合理である。\QED

\condpagelarge{命題XXII. 問題XIV.}

\textit{与えられた五つの点を通る軌道を描くこと。}

\pngright{087.png}{}{}{}

五つの点$A$、$B$、$C$、$D$、$P$が与えられたとする。そのうちの一つ$A$から、他の任意の二点$B$、$C$（これらを極と呼ぶ）へ直線$AB$、$AC$を引き、第四の点$P$を通ってこれらに平行な$TPS$、$PRQ$を引く。次に二つの極$B$、$C$から、第五の点$D$を通る無限の二直線$BDT$、$CRD$を引き、これらが先に引いた$TPS$、$PRQ$（前の線は前の線と、後の線は後の線と）と$T$、$R$で交わるようにする。最後に、直線$PT$、$PR$から、$TR$に平行な直線$tr$を引くことによって、$PT$、$PR$に比例する任意の$Pt$、$Pr$を切り取る。それらの端$t$、$r$と極$B$、$C$を通って引いた$Bt$、$Cr$が点$d$で交わるならば、その点$d$は求める軌道上にある。

実際、その点$d$は（補助定理XXにより）四つの点$A$、$B$、$P$、$C$を通る円錐曲線上にあり、線$Rr$、$Tt$が消えるとき、点$d$は点$D$に一致する。したがってこの円錐曲線は五つの点$A$、$B$、$C$、$D$、$P$を通る。\QED

\condpagelarge{\textit{別法.}}

与えられた点から任意の三つ$A$、$B$、$C$を結び、そのうち二つ$B$、$C$を極として、大きさの与えられた角$ABC$、$ACB$を回転させる。その脚$BA$、$CA$を、最初は点$D$に、次に点$P$に当て、他の脚$BL$、$CL$が、どちらの場合にも交わる点$M$、$N$を記す。無限の直線$MN$を引く。そして、それらの動く角を極$B$、$C$のまわりに、脚$BL$、$CL$（または$BM$、$CM$）の交点、いまこれを$m$とすれば、それが常にその無限の直線$MN$に落ちるという規則で回転させる。すると脚$BA$、$CA$（または$BD$、$CD$）の交点、いまこれを$d$とすれば、それが求める軌道$PADdB$を描く。

\pngright{086.png}{}{}{}

実際、点$d$は補助定理XXIにより、点$B$、$C$を通る円錐曲線に接し、点$m$が点$L$、$M$、$N$に近づくとき、点$d$は（作図により）点$A$、$D$、$P$に近づく。したがって五つの点$A$、$B$、$C$、$D$、$P$を通る円錐曲線が描かれる。\QED

\textit{系1.}\wsp{}これによって、軌道に与えられた任意の点$B$、$C$で接する直線を、すばやく引くことができる。どちらの場合でも、点$d$を点$C$に近づければ、直線$Cd$が求める接線になる。

\textit{系2.}\wsp{}これによって、軌道の中心、直径、通径も、補助定理XIXの系2におけるように求められる。

\condpagelarge{\textit{注.}}

前の場合の作図は、$BP$を結び、必要なら延ばして、$Bp$を$BP$に対して$PR$が$PT$に対するのと同じ比にとり、$p$を通って$SPT$に平行な無限の直線$p\mathrm{D}$を引き、その上に$p\mathrm{D}$を常に$Pr$に等しくとり、直線$B\mathrm{D}$、$Cr$を引いて$d$で交わらせることにすれば、少し簡単になる。実際、$Pr : Pt$、$PR : PT$、$pB : PB$、$p\mathrm{D} : Pt$は同じ比であるから、$p\mathrm{D}$と$Pr$はつねに等しい。この方法で軌道の点は最も手早く見つかる。ただし、二番目の方法のように、曲線を機械的に描くほうがよければ、そうしてもよい。

\condpagelarge{命題XXIII. 問題XV.}

\pngright{090a.png}{}{}{}

\textit{与えられた四つの点を通り、位置の与えられた直線に接する軌道を描くこと。}

\textit{場合1.}\wsp{}接線$HB$、接点$B$、他の三点$C$、$D$、$P$が与えられたとする。$BC$を結び、$BH$に平行に$PS$を、$BC$に平行に$PQ$を引いて、平行四辺形$BSPQ$を完成させる。$BD$を引いて$SP$と$T$で交わらせ、$CD$を引いて$PQ$と$R$で交わらせる。最後に、$TR$に平行な任意の$tr$を引くことによって、$PQ$、$PS$から、それぞれ$PR$、$PT$に比例する$Pr$、$Pt$を切り取る。すると、引いた$Cr$、$Bt$の交点$d$は（補助定理XXの系2により）つねに描くべき軌道上に落ちる。

\condpagelarge{\textit{別法.}}

\pngright{090b.png}{}{}{}

大きさの与えられた角$CBH$を極$B$のまわりに回転させ、同時に、両方向に延ばした任意の直線の半径$DC$を極$C$のまわりに回転させる。角の脚$BC$がその半径と交わる点$M$、$N$を記す。このとき、もう一方の脚$BH$がその同じ半径と点$D$、$P$で交わるものとする。次に、引いた無限の直線$MN$に対して、その半径$CP$または$CD$と角の脚$CB$が絶えず交わるようにする。すると、もう一方の脚$BH$とその半径との交点が、求める軌道を描く。

実際、上の問題の作図において点$A$を点$B$に近づければ、線$CA$と$CB$は一致し、線$AB$はその最後の位置で接線$BH$になる。したがって、そこで置いた作図は、ここに述べた作図と同じになる。ゆえに脚$BH$と半径との交点は、点$C$、$D$、$P$を通り、点$B$で直線$BH$に接する円錐曲線を描く。\QED

\pngright{091.png}{}{}{}

\textit{場合2.}\wsp{}接線$HI$の外にある四つの点$B$、$C$、$D$、$P$が与えられたとする。二組の$BD$、$CP$を結んで、$G$で交わらせ、接線とは$H$、$I$で交わらせる。接線を点$A$で、$HA : AI$が、「$BH$と$HD$の比例中項」と「$CG$と$GP$の比例中項」とがつくる矩形の、「$PI$と$IC$の比例中項」と「$DG$と$GB$の比例中項」とがつくる矩形に対する比に等しくなるように切る。すると$A$が接点である。実際、$PI$に平行な$HX$が軌道を任意の点$X$、$Y$で切るとき、（円錐曲線論により）$HA^2 : AI^2$は、矩形$XHY$と矩形$BHD$（すなわち矩形$CGP$と矩形$DGB$）との比に、矩形$BHD$と矩形$PIC$との比を合成したものに等しい。接点$A$が見つかれば、場合1のように軌道を描く。\QED

点$A$は、点$H$と$I$の間にとることも、外にとることもできる。それに応じて軌道は二通りに描かれる。

\condpagelarge{命題XXIV. 問題XVI.}

\textit{与えられた三つの点を通り、位置の与えられた二直線に接する軌道を描くこと。}

\pngright{092.png}{}{}{}

接線$HI$、$KL$と点$B$、$C$、$D$が与えられたとする。$BD$を引いて接線と点$H$、$K$で交わらせ、$CD$を引いて接線と点$I$、$L$で交わらせる。引いたこれらの線を、$HR : KR$が「$BH$と$HD$の比例中項」と「$BK$と$KD$の比例中項」の比に等しくなるように点$R$で、$IS : LS$が「$CI$と$ID$の比例中項」と「$CL$と$LD$の比例中項」の比に等しくなるように点$S$で、切る。$RS$を引いて接線と$A$、$P$で交わらせれば、$A$、$P$が接点である。

実際、$A$と$P$が接線上のどこかにある接点であるとし、点$H$、$I$、$K$、$L$のうちの任意の一つ$I$を通って、接線$KL$に平行な直線$IY$を引き、曲線と$X$、$Y$で交わらせ、その上に$IX$と$IY$の比例中項$IZ$をとる。すると円錐曲線論により、矩形$XIY$（すなわち$IZ^2$）と$LP^2$との比は、矩形$CID$と矩形$CLD$との比に等しい。すなわち（作図により）$SI^2 : SL^2$に等しく、したがって$IZ : LP = SI : SL$である。ゆえに点$S$、$P$、$Z$は一直線上にある。さらに接線が$G$で交わるとき、（円錐曲線論により）矩形$XIY$（すなわち$IZ^2$）と$IA^2$との比は、$GP^2 : GA^2$に等しく、したがって$IZ : IA = GP : GA$である。ゆえに点$P$、$Z$、$A$は一直線上にあり、したがって点$S$、$P$、$A$は一直線上にある。同じ論法によって、点$R$、$P$、$A$が一直線上にあることも示される。ゆえに接点$A$、$P$は直線$SR$上にある。これらが見つかれば、軌道は、上の問題の場合1のように描かれる。\QED

\condpagelarge{補助定理XXII.}

\textit{図形を、同じ種類の別の図形に変えること。}

\pngright{093.png}{}{}{}

任意の図形$HGI$を変えたいとする。任意に平行な二直線$AO$、$BL$を引いて、位置の与えられた第三の直線$AB$と$A$、$B$で交わらせ、図形の任意の点$G$から直線$AB$へ、$OA$に平行に$GD$を引く。次に、直線$OA$上の与えられた点$O$から点$D$へ直線$OD$を引き、$BL$と$d$で交わらせる。その交点から、直線$BL$と任意の与えられた角をなす直線$gd$を立て、その長さを、$Od$に対して$GD$が$OD$に対するのと同じ比になるようにとる。すると$g$は、新しい図形$hgi$の中で点$G$に対応する点になる。同じ方法で、最初の図形の一つ一つの点は、新しい図形の同じ数だけの点を与える。そこで、点$G$が連続の運動で最初の図形のすべての点を巡ると考えれば、点$g$も同様に連続の運動で新しい図形のすべての点を巡り、その図形を描く。

区別のために、$DG$を最初の縦線、$dg$を新しい縦線、$BD$を最初の横線（abscissa）、$Bd$を新しい横線、$O$を極、$OD$を切り取る半径、$OA$を最初の縦の半径、$Oa$（これによって平行四辺形$OABa$が完成する）を新しい縦の半径と呼ぶことにしよう。

さて私は言う。もし点$G$が位置の与えられた直線に接するならば、点$g$も位置の与えられた直線に接する。もし点$G$が円錐曲線に接するならば、点$g$も円錐曲線に接する。ここでは円錐曲線の中に円も数えている。さらに、もし点$G$が解析的な三次の線に接するならば、点$g$も同じく三次の線に接し、より高い次数の曲線についても同様である。点$G$、$g$が接する二つの線は、つねに同じ解析的次数のものである。

実際、$ad : OA = Od : OD = dg : DG = AB : AD$であるから、$AD$は$\frac{OA \times AB}{ad}$に等しく、$DG$は$\frac{OA \times dg}{ad}$に等しい。いま、点$D$が直線に接するならば、横線$AD$と縦線$DG$の関係を与える任意の方程式のなかで、不定量$AD$と$DG$は一次元にしか上がらない。そこでこの方程式のなかで、$AD$の代わりに$\frac{OA \times AB}{ad}$を、$DG$の代わりに$\frac{OA \times dg}{ad}$を書けば、新しい方程式が生じ、そのなかで新しい横線$ad$と新しい縦線$dg$は一次元にしか上がらず、したがって直線を表す。もし最初の方程式のなかで$AD$と$DG$（またはそのどちらか）が二次元に上がっていたなら、第二の方程式のなかで$ad$と$dg$も同じく二次元に上がる。三次元以上についても同様である。第二の方程式の不定量$ad$、$dg$と、最初の方程式の$AD$、$DG$は、つねに同じ次元の数まで上がり、したがって点$G$、$g$が接する線は、同じ解析的次数のものである。

さらに私は言う。もし最初の図形のなかで、ある直線が曲線に接するならば、この直線を移したものは、新しい図形のなかで曲線に接し、逆もまた成り立つ。なぜなら、曲線の任意の二点が互いに近づいて、最初の図形のなかで一致するならば、移された同じ点も新しい図形のなかで一致し、したがってそれらの点を結ぶ直線は、どちらの図形でも曲線の接線になるからである。これらの主張の証明は、もっと幾何学的な流儀で構成することもできるが、簡潔さのために省く。

したがって、直線で囲まれた図形を別の図形に変えるには、直線の交点を移し、それらを通る直線を新しい図形のなかに引けば足りる。曲線で囲まれた図形を変える必要があるなら、点と接線、およびその他の、曲線を定めるのに役立つ直線を移さなければならない。この補助定理は、与えられた図形をより簡単な図形に変えることによって、より難しい問題の解法に役立つ。実際、集まる任意の直線は、最初の縦の半径$AO$として、集まる直線の交点を通る任意の直線を用いることで、平行線に変わる。このようにすると、その交点は無限遠に去り、無限に遠い点へ向かう線は平行だからである。そして問題が新しい図形のなかで解かれたのち、逆の操作によってこの図形を最初の図形に変えれば、求める解が得られる。

この補助定理は、立体的な問題の解法にも役立つ。なぜなら、二つの円錐曲線が現れ、その交点によって問題が解かれる場合には、そのうち一つを円に変えてよいからである。平面的な問題の作図でも、直線と円錐曲線は直線と円に変わる。

\condpagelarge{命題XXV. 問題XVII.}

\textit{与えられた二つの点を通り、位置の与えられた三直線に接する軌道を描くこと。}

任意の二つの接線どうしの交点と、第三の接線が二つの与えられた点を通る直線と交わる点を通って、無限の直線を引く。それを最初の縦の半径として用い、先の補助定理によって図形を新しい図形に変える。この図形のなかで、それら二つの接線は平行になり、第三の接線は、二つの点を通る直線に平行になる。

\pngright{096.png}{}{}{}

$hi$、$kl$を二つの平行な接線、$ik$を第三の接線、$hl$をこれに平行で、この新しい図形のなかで円錐曲線が通るべき点$a$、$b$を通る直線とし、平行四辺形$hikl$を完成させる。直線$hi$、$ik$、$kl$を点$c$、$d$、$e$で、$hc$が矩形$ahb$の平方根に対して、また$ic : id$、$ke : kd$が、$hi$と$kl$の和が、「直線$ik$」と「矩形$ahb$の平方根」と「矩形$alb$の平方根」の三つの線の和に対するのと同じ比になるように切る。すると$c$、$d$、$e$が接点である。

実際、円錐曲線論により、$hc^2$と矩形$ahb$、$ic^2$と$id^2$、$ke^2$と$kd^2$、$el^2$と矩形$alb$は、同じ比にある。したがって$hc$と矩形$ahb$の平方根、$ic$と$id$、$ke$と$kd$、$el$と矩形$alb$の平方根は、その比の平方根の比にあり、合成して、すべての前項$hi$と$kl$が、すべての後項（矩形$ahb$の平方根と、直線$ik$と、矩形$alb$の平方根）に対してもつ与えられた比にある。こうしてその与えられた比から、新しい図形のなかの接点$c$、$d$、$e$が得られる。直前の補助定理の逆の操作によってこれらの点を最初の図形に移し、そこで問題XIVの場合1によって軌道を描く。\QED

なお、点$a$、$b$が点$h$、$l$の間にあるか外にあるかに応じて、点$c$、$d$、$e$も点$h$、$i$、$k$、$l$の間か外にとらなければならない。もし点$a$、$b$の一方が点$h$、$l$の間に落ち、他方が外に落ちるなら、この問題は不可能である。

\condpagelarge{命題XXVI. 問題XVIII.}

\textit{与えられた一つの点を通り、位置の与えられた四直線に接する軌道を描くこと。}

任意の二つの接線の共通の交点から、残りの二つの接線の共通の交点へ、無限の直線を引き、これを最初の縦の半径として用いて、（補助定理XXIIにより）図形を新しい図形に変える。すると、縦の半径に集まっていた二組の接線は、いまや平行になる。それらを、平行四辺形$hikl$をつくる$hi$と$kl$、$ik$と$hl$とする。$p$を、この新しい図形のなかで、最初の図形で与えられた点に対応する点とする。図形の中心$O$を通って$pq$を引き、$Oq$を$Op$に等しくすれば、$q$が、この新しい図形のなかで円錐曲線が通るべきもう一つの点になる。補助定理XXIIの逆の操作によってこの点を最初の図形に移せば、そこに軌道を描くべき二つの点が得られる。これらを通る軌道は、問題XVIIによって描くことができる。\QED

\condpagelarge{補助定理XXIII.}

\textit{位置の与えられた二直線$AC$、$BD$が、与えられた点$A$、$B$で終わり、互いに一定の比をもち、不定の点$C$、$D$を結ぶ直線$CD$が、与えられた比で$K$において分けられるとする。私は言う。点$K$は、位置の与えられた直線上にある。}

\pngright{097.png}{}{}{}

実際、直線$AC$、$BD$が$E$で交わるとし、$BE$の上に、$BG : AE = BD : AC$となるように$BG$をとり、$FD$を$EG$に等しくする。すると$EC : GD$、すなわち$EC : EF$は$AC : BD$に等しく、したがって与えられた比であり、ゆえに三角形$EFC$は形が与えられる。$CF$を、$CK : CD$の比で$L$において分ければ、三角形$EFL$も形が与えられ、したがって点$L$は位置の与えられた直線$EL$上にある。$LK$を結ぶと、$FD$が与えられ、$LK : FD$が与えられた比であるから、$LK$が与えられる。これに等しく$EH$をとれば、$ELKH$は平行四辺形になる。ゆえに点$K$は、平行四辺形の位置の与えられた辺$HK$上にある。\QED

\condpagelarge{補助定理XXIV.}

\textit{三本の直線が任意の円錐曲線に接し、そのうちの二本は平行で位置が与えられているとする。私は言う。これら二本に平行な、曲線の半径（半径方向の線分）は、接点と第三の接線とにはさまれたそれら二本の線分の比例中項である。}

\pngright{098.png}{}{}{}

$AF$、$GB$を、円錐曲線$ADB$に点$A$、$B$で接する二本の平行線とする。$EF$を、曲線に点$I$で接する第三の直線で、先の二本の接線と$F$、$G$で交わるものとする。$CD$を、接線に平行な図形の半径とする。私は言う。$AF$、$CD$、$BG$は連比例する。

実際、共役な直径$AB$、$DM$が、接線$FG$と$E$、$H$で交わり、互いに$C$で交わるとし、平行四辺形$IKCL$を完成させる。円錐曲線の性質から、$EC : CA = CA : LC$であり、したがって分割して$EC - CA : CA - CL$、すなわち$EA : AL$である。合成して$EA : EA + AL$、すなわち$EL$に対して、$EC : EC + CA$、すなわち$EB$に等しい。したがって（三角形$EAF$、$ELI$、$ECH$、$EBG$が相似であるから）$AF : LI = CH : BG$となる。また円錐曲線の性質から$LI$、すなわち$CK$と$CD$との比は、$CD : CH$に等しい。したがって乱れ比の等比により$AF : CD = CD : BG$である。\QED

\textit{系1.}\wsp{}したがって、二つの接線$FG$、$PQ$が、平行な接線$AF$、$BG$と、$F$、$G$および$P$、$Q$で交わり、互いに$O$で交わるならば、（乱れ比の等比により）$AF : BQ = AP : BG$、分割して$FP : GQ$、したがって$FO : OG$となる。

\textit{系2.}\wsp{}これによって、点$P$と$G$、$F$と$Q$を通って引いた二直線$PG$、$FQ$は、図形の中心と接点$A$、$B$を通る直線$ACB$の上で交わる。

\condpagelarge{補助定理XXV.}

\textit{平行四辺形の四つの辺（限りなく延長したもの）が任意の円錐曲線に接し、任意の第五の接線によって切り取られるとする。切り取られた部分は、平行四辺形の向かい合う角に向かってとる。私は言う。一つの辺の切り取られた部分の、その辺全体に対する比は、隣り合う辺の、接点と第三の辺とにはさまれた部分の、この隣り合う辺の切り取られた部分に対する比に等しい。}

\pngright{100.png}{}{}{}

平行四辺形$MIKL$の四つの辺$ML$、$IK$、$KL$、$MI$が、円錐曲線に点$A$、$B$、$C$、$D$で接し、第五の接線$FQ$がこれらの辺を$F$、$Q$、$H$、$E$で切るとする。私は言う。$ME : MI = BK : KQ$、$KH : KL = AM : MF$である。実際、前の補助定理の系により、$ME : EI = AM$すなわち$BK : BQ$であり、合成して$ME : MI = BK : KQ$である。\QED　また$KH : HL = BK$すなわち$AM : AF$であり、分割して$KH : KL = AM : MF$である。\QED

\textit{系1.}\wsp{}したがって、平行四辺形$IKLM$が与えられれば、矩形$KQ \times ME$が与えられ、またこれに等しい矩形$KH \times MF$も与えられる。実際、それらの矩形は、三角形$KQH$と$MFE$が相似であるから等しい。

\textit{系2.}\wsp{}さらに、第六の接線$eq$が引かれて、$KI$、$MI$と$e$、$q$で交わるならば、矩形$KQ \times ME$は矩形$Kq \times Me$に等しく、$KQ : Me = Kq : ME$であり、分割して$Qq : Ee$となる。

\textit{系3.}\wsp{}したがって、$Eq$、$eQ$を結んで二等分し、二等分点を通る直線を引けば、それは円錐曲線の中心を通る。実際、$Qq : Ee = KQ : Me$であるから、その直線は、（補助定理XXIIIにより）$Eq$、$eQ$、$MK$のすべての中点を通り、直線$MK$の中点は曲線の中心だからである。

\condpagelarge{命題XXVII. 問題XIX.}

\textit{位置の与えられた五本の直線に接する軌道を描くこと。}

\pngright{101.png}{}{}{}

位置の与えられた接線$ABG$、$BCF$、$GCD$、$FDE$、$EA$が与えられたとする。四本の任意の接線に囲まれた四辺形$ABFE$の対角線$AF$、$BE$を二等分すれば、（補助定理XXVの系3により）二等分点を通る直線は、軌道の中心を通る。さらに、他の任意の四本の接線に囲まれた四辺形$BGDF$の対角線（と言ってよければ）$BD$、$GF$を二等分すれば、二等分点を通る直線は曲線の中心を通る。したがって中心は、これら二つの二等分線の交点として得られる。それを$O$とする。任意の接線$BC$に平行に$KL$を、中心$O$が平行線の間のちょうど真ん中にくるような距離に引けば、引いた$KL$は描くべき軌道に接する。これが他の任意の二つの接線$CD$、$FDE$と$L$、$K$で交わるとする。平行でない接線$CL$、$FK$と平行線$CF$、$KL$との交点$C$と$K$、$F$と$L$を通って$CK$、$FL$を引き、$R$で交わらせる。そして直線$OR$を引き延ばせば、それが平行な接線$CF$、$KL$を接点で切る。これは補助定理XXIVの系2から明らかである。同じ方法で他の接点も見つけることができ、そのうえで、問題XIVの場合1によって軌道を描く。\QED

\condpagelarge{\textit{注.}}

軌道の中心、あるいは漸近線が与えられる問題は、先の問題に含まれる。実際、点と接線が中心とともに与えられれば、中心から反対側に等しい距離にある、同じ数の他の点と他の接線が得られる。漸近線は接線とみなさなければならず、その無限に遠い端（そう言ってよければ）は接点とみなさなければならない。任意の接線の接点が無限に遠ざかると考えれば、接線は漸近線に変わり、問題XVおよび問題XIVの場合1の作図は、漸近線が与えられる問題の作図に変わる。

\pngright{102.png}{}{}{}

軌道が描かれたあと、その軸と焦点は次の方法で求められる。補助定理XXIの作図と図において、動く角$PBN$、$PCN$の脚$BP$、$CP$（その交点によって軌道が描かれたもの）が、互いに平行になり、その位置を保ちながら、その図の極$B$、$C$のまわりに回転するようにする。その間に、これらの角のもう一方の脚$CN$、$BN$が、その交点$K$または$k$によって円$IBKGC$を描くようにする。この円の中心を$O$とする。この中心から、軌道が描かれていたあいだ、もう一方の脚$CN$、$BN$が集まっていた定規$MN$へ垂線$OH$を下ろし、円と$K$、$L$で交わらせる。そして、もう一方の脚$CK$、$BK$が、定規に近いほうの点$K$で交わるとき、最初の脚$CP$、$BP$は、長軸に平行で短軸に垂直になる。同じ脚が遠いほうの点$L$で交わるときは、逆のことが起こる。したがって軌道の中心が与えられれば、軸が与えられる。これらが与えられれば、焦点はすぐに見つかる。

\vspace{\baselineskip}

さらに軸の平方は互いに$KH : LH$の比にあり、これによって、形の与えられた軌道を、与えられた四つの点を通って描くことは容易である。実際、与えられた点のうち二つを極$C$、$B$とすれば、第三の点が動く角$PCK$、$PBK$を与える。次に、軌道の形が与えられているから、比$OH : OK$が与えられ、中心$O$、半径$OH$で円を描き、第四の点を通ってその円に接する直線を引けば、軌道を描くための定規$MN$が与えられる。これによって逆に、形の与えられた台形は（ある不可能な場合を除いて）、与えられた任意の円錐曲線に内接させることができる。

\vspace{\baselineskip}

ほかにも、形の与えられた軌道を、与えられた点と接線から描くのに役立つ補助定理がある。次のようなものがその種類である。位置の与えられた任意の点を通って直線を引き、これが与えられた円錐曲線を二点で切るとし、その交点の間隔を二等分するとき、二等分点は、先の曲線と同じ種類で、先の曲線の軸に平行な軸をもつ別の円錐曲線に接する。しかし私は、より役に立つことへ急ぐ。

\condpagelarge{補助定理XXVI.}

\textit{形と大きさの与えられた三角形の三つの角を、すべてが平行ではない、位置の与えられた同じ数の直線の上に、一つずつ置くこと。}

\pngright{104a.png}{}{}{}

位置の与えられた三本の無限の直線$AB$、$AC$、$BC$が与えられ、三角形$DEF$を、その角$D$が直線$AB$に、角$E$が直線$AC$に、角$F$が直線$BC$に接するように置かなければならないとする。$DE$、$DF$、$EF$の上に、それぞれ角$BAC$、$ABC$、$ACB$に等しい角をふくむ、三つの円の弓形$DRE$、$DGF$、$EMF$を描く。これらの弓形は、文字の列$DRED$が文字の列$BACB$と、文字の列$DGFD$が文字の列$ABCA$と、文字の列$EMFE$が文字の列$ACBA$と、同じ順序で一周するような側に、線$DE$、$DF$、$EF$に対して描く。そのあとで、これらの弓形を円に完成させる。最初の二つの円が互いに$G$で交わるとし、それらの中心を$P$と$Q$とする。$GP$、$PQ$を結び、$AB$に対して$GP$が$PQ$に対するのと同じ比になるように$Ga$をとり、中心$G$、半径$Ga$で円を描いて、最初の円$DGE$と$a$で交わらせる。次に$aD$を結んで第二の円$DFG$と$b$で交わらせ、$aE$を結んで第三の円$GEc$と$c$で交わらせる。そして図形$abcDEF$に相似で等しい図形$ABCdef$を完成させる。私は、これで作図できたと言う。

\pngright{104b.png}{}{}{}

実際、$aD$と$n$で交わるように$Fc$を引く。$aG$、$bG$、$PD$、$QD$を結び、$PQ$を$R$まで延ばす。作図により角$EaD$は角$CAB$に等しく、角$EcF$は角$ACB$に等しいから、三角形$anc$は三角形$ABC$と等角である。したがって角$anc$、すなわち$FnD$は角$ABC$に等しく、したがって角$FbD$に等しい。ゆえに点$n$は点$b$に一致する。さらに、中心の角$GPD$の半分である角$GPQ$は、円周角$GaD$に等しい。また、中心の角$GQD$の補角の半分である角$GQR$は、円周角$GbD$に等しい。したがってそれらの補角$PQG$、$abG$は等しく、ゆえに三角形$GPQ$と$Gab$は相似で、$Ga : ab = GP : PQ$、すなわち（作図により）$Ga : AB$である。よって$ab$と$AB$は等しく、したがって、いま相似であることを示した三角形$abc$と$ABC$は等しくもある。さらに、三角形$DEF$の角$D$、$E$、$F$が、それぞれ三角形$abc$の辺$ab$、$ac$、$bc$に接するから、図形$abcDEF$に相似で等しい図形$ABCdef$を完成させることができ、それを完成させることによって問題が解かれる。\QED

\textit{系.}\wsp{}これによって、長さの与えられた部分が、位置の与えられた三本の直線の間にはさまれるような直線を引くことができる。三角形$DEF$を考え、点$D$が辺$EF$に近づき、辺$DE$、$DF$が一直線上に置かれて、三角形が一本の直線に変わるとする。その直線のうち、与えられた部分$DE$は位置の与えられた直線$AB$、$AC$の間に、与えられた部分$DF$は位置の与えられた直線$AB$、$BC$の間に、はさまれなければならない。前の作図をこの場合にあてはめれば、問題は解かれる。

\condpagelarge{命題XXVIII. 問題XX.}

\textit{形と大きさの与えられた軌道で、その与えられた部分が、位置の与えられた三直線の間にはさまれるものを描くこと。}

\pngright{106.png}{}{}{}

曲線$DEF$に相似で等しい軌道を描くとする。その軌道は、位置の与えられた三本の直線$AB$、$AC$、$BC$によって、この曲線の与えられた部分$DE$、$EF$に相似で等しい部分に切られるものとする。直線$DE$、$EF$、$DF$を引き、この三角形$DEF$の角$D$、$E$、$F$を、位置の与えられたそれらの直線の上に置く（補助定理XXVIによる）。次に、その三角形のまわりに、曲線$DEF$に相似で等しい軌道を描く。\QED

\condpagelarge{補助定理XXVII.}

\textit{形の与えられた四辺形で、その角が、位置の与えられた四本の直線（すべてが平行でもなく、一点に集まりもしない）の上に、一つずつ乗るものを描くこと。}

\pngrightsc{108.png}{}{}{}{}

位置の与えられた四本の直線$ABC$、$AD$、$BD$、$CE$が与えられ、そのうち第一の直線が、第二の直線と$A$で、第三の直線と$B$で、第四の直線と$C$で交わるとする。そして、四辺形$FGHI$に相似な四辺形$fghi$を描くとする。その角$f$は、与えられた角$F$に等しくて直線$ABC$に接し、他の角$g$、$h$、$i$は、他の与えられた角$G$、$H$、$I$に等しくて、他の直線$AD$、$BD$、$CE$にそれぞれ接するものとする。

$FH$を結び、$FG$、$FH$、$FI$の上に、同じ数の円の弓形$FSG$、$FTH$、$FVI$を描く。最初の弓形$FSG$は角$BAD$に等しい角を、第二の弓形$FTH$は角$CBE$に等しい角を、第三の弓形$FVI$は角$ACE$に等しい角を、ふくむものとする。これらの弓形は、線$FG$、$FH$、$FI$に対して、文字の列$FSGF$が文字の列$BADB$と同じ円の順序をなし、文字の列$FTHF$が文字の列$CBEC$と同じ順序で、文字の列$FVIF$が文字の列$ACEA$と同じ順序で、一周するような側に描かなければならない。弓形を円に完成させ、$P$を最初の円$FSG$の中心、$Q$を第二の円$FTH$の中心とする。$PQ$を結んで両方向に延ばし、その上に$QR$を、$PQ$に対して、$BC$が$AB$に対するのと同じ比になるようにとる。$QR$は、文字$P$、$Q$、$R$の円の順序が、文字$A$、$B$、$C$の順序と同じになるように、点$Q$に対してそちら側にとる。中心$R$、半径$RF$で第四の円$FNc$を描き、第三の円$FVI$と$c$で交わらせる。$Fc$を結んで、最初の円と$a$で、第二の円と$b$で交わらせる。$aG$、$bH$、$cI$を引き、図形$abcFGHI$に相似な図形$ABCfghi$を作る。すると四辺形$fghi$が、作るべきその四辺形になる。

実際、最初の二つの円$FSG$、$FTH$が互いに$K$で交わるとする。$PK$、$QK$、$RK$、$aK$、$bK$、$cK$を結び、$QP$を$L$まで延ばす。円周角$FaK$、$FbK$、$FcK$は、中心角$FPK$、$FQK$、$FRK$の半分であり、したがって、それらの半分$LPK$、$LQK$、$LRK$に等しい。ゆえに図形$PQRK$は図形$abcK$と等角で相似であり、したがって$ab : bc = PQ : QR$、すなわち$AB : BC$である。さらに作図により、角$FaG$、$FbH$、$FcI$は、角$fAg$、$fBh$、$fCi$に等しい。ゆえに図形$abcFGHI$に相似な図形$ABCfghi$を完成させることができる。そうすれば四辺形$fghi$は、四辺形$FGHI$に相似になり、その角$f$、$g$、$h$、$i$で直線$AB$、$AD$、$BD$、$CE$に接する。\QED

\textit{系.}\wsp{}これによって、位置の与えられた四本の直線の間に、決まった順序ではさまれた部分が、互いに与えられた比をもつような直線を引くことができる。角$FGH$、$GHI$を、直線$FG$、$GH$、$HI$が一直線に並ぶまで大きくし、この場合に問題を作図すれば、位置の与えられた四本の直線$AB$と$AD$、$AD$と$BD$、$BD$と$CE$の間にはさまれる部分$fg$、$gh$、$hi$が、線$FG$、$GH$、$HI$と同じ比をもち、同じ順序を保つ、直線$fghi$が引かれる。同じことは次のようにすれば、より手早くできる。

\pngright{109.png}{}{}{}

$AB$を$K$まで、$BD$を$L$まで、$BK : AB = HI : GH$、$DL : BD = GI : FG$となるように延ばし、$KL$を結んで直線$CE$と$i$で交わらせる。$iL$を$M$まで、$LM : iL = GH : HI$となるように延ばし、$LB$に平行に$MQ$を引いて直線$AD$と$g$で交わらせ、$gi$を引いて$AB$、$BD$と$f$、$h$で交わらせる。私は、これで作図できたと言う。

実際、$Mg$が直線$AB$を$Q$で、$AD$が直線$KL$を$S$で切るとし、$BD$に平行な$AP$を引いて$iL$と$P$で交わらせる。すると$Mg : Lh$（$Mi : Li$、$gi : hi$、$AK : BK$）と$AP : BL$は同じ比になる。$DL : RL$が同じ比になるように$DL$を$R$で切る。$gS : gM$、$AS : AP$、$DS : DL$が比例するから、等比により、$gS : Lh = AS : BL = DS : RL$となる。混合して、$BL - RL : Lh - BL = AS - DS : gS - AS$となる。すなわち$BR : Bh = AD : Ag$、したがって$BD : gQ$である。入れ替えて$BR : BD = Bh : gQ$、すなわち$fh : fg$となる。ところが作図により$BR : BD = FH : FG$である。ゆえに$fh : fg = FH : FG$となる。さらに$ig : ih = Mi : Li$、すなわち$IG : IH$であるから、線$FI$、$fi$が$g$、$h$と$G$、$H$において相似に分けられていることは明らかである。\QED

この系の作図において、$LK$を引いて$CE$と$i$で交わらせたあと、$iE$を$V$まで、$EV : iE = FH : HI$となるように延ばし、$BD$に平行に$Vf$を引いてもよい。中心$i$、半径$IH$で円を描いて$BD$と$X$で交わらせ、$iX$を$Y$まで、$iY$が$IF$に等しくなるように延ばし、$BD$に平行に$Yf$を引くことにしても、同じことになる。

\condpagelarge{命題XXIX. 問題XIX.}

\textit{形の与えられた軌道で、位置の与えられた四本の直線によって、順序と形と比の与えられた部分に切られるものを描くこと。}

\pngright{110.png}{}{}{}

曲線$FGHI$に相似な軌道$fghi$を描くとする。その部分$fg$、$gh$、$hi$は、曲線の部分$FG$、$GH$、$HI$に相似で比例し、位置の与えられた直線$AB$と$AD$、$AD$と$BD$、$BD$と$EC$の間に、第一のものは第一のものの間に、第二のものは第二のものの間に、第三のものは第三のものの間に、はさまれるものとする。直線$FG$、$GH$、$HI$、$FI$を引き、四辺形$FGHI$に相似な四辺形$fghi$を描く。その角$f$、$g$、$h$、$i$が、位置の与えられた直線$AB$、$AD$、$BD$、$CE$に、一つずつ、述べた順序で接するものとする。次に（補助定理XXVIIにより）この四辺形のまわりに、曲線$FGHI$に似た軌道を描く。

\condpagelarge{\textit{注.}}

この問題は次のようにも作図できる。$FG$、$GH$、$HI$、$FI$を結び、$GF$を$V$まで延ばし、$FH$、$IG$を結んで、角$FGH$、$VFH$に等しい角$CAK$、$DAL$を作る。$AK$、$AL$が直線$BD$と$K$、$L$で交わるとし、そこから$KM$、$LN$を引く。そのうち$KM$は角$AKM$を角$GHI$に等しくし、$AK$に対して$HI$が$GH$に対するのと同じ比にする。$LN$は角$ALN$を角$FHI$に等しくし、$AL$に対して$HI$が$FH$に対するのと同じ比にする。$AK$、$KM$、$AL$、$LN$は、文字の列$CAKMC$、$ALK$、$DALND$が、文字の列$FGHIF$と同じ順序で一周するような側に、線$AD$、$AK$、$AL$に対して引く。引いた$MN$が直線$CE$と$i$で交わるとする。角$iEP$を角$IGF$に等しくし、$PE : Ei = FG : GI$とする。$P$を通って$QPf$を引き、これが直線$AED$とつくる角$PQE$を角$FIG$に等しくし、直線$AB$と$f$で交わらせて、$fi$を結ぶ。$PE$、$PQ$は、文字の列$PEiP$、$PEQP$が、文字の列$FGHIF$と同じ円の順序をなすような側に、線$CE$、$PE$に対して引く。そして線$fi$の上に、同じ文字の順序で四辺形$fghi$を、四辺形$FGHI$に相似に作り、与えられた形の軌道をそのまわりに描けば、問題は解かれる。

\pngcent{111.png}{}{}

ここまでは軌道を求めることについて述べた。物体が、求めた軌道の上をどのように動くかを決めることが、残っている。

\section{第1編 第6節：与えられた軌道における運動の決定（命題XXX〜XXXI）}


% Section: 第1編 第6節：与えられた軌道における運動の決定（命題XXX〜XXXI）



\begin{center}\textit{与えられた軌道における運動を求めること}\end{center}

\condpagelarge{命題XXX. 問題XXII.}

\textit{与えられた放物線の軌道を動く物体について、指定された時刻における位置を求めること。}

\pngright{112.png}{}{}{}

$S$を放物線の焦点、$A$を主頂点とし、$4AS \times M$を、動径$SP$によって描かれた（物体が頂点を通りすぎたあとに描かれた、あるいは頂点に着く前にこれから描かれる）放物線の面積$APS$とする。この面積は、それに比例する時間から分かる。$AS$を$G$で二等分し、垂線$GH$を$3M$に等しく立てる。すると、中心$H$、半径$HS$で描いた円は、放物線を求める位置$P$で切る。

実際、軸へ垂線$PO$を下ろすと、$HG^2 + GS^2$（$= HS^2 = HP^2 = GO^2 + (PO - HG)^2$）$= GO^2 + HG^2 - 2HG \times PO + PO^2$である。両辺から$HG^2$を除けば、$GS^2 = GO^2 - 2HG \times PO + PO^2$、すなわち$2HG \times PO$（$= GO^2 + PO^2 - GS^2 = AO^2 - 2GA \times AO + PO^2$）$= AO^2 + \frac{3}{4}PO^2$となる。$AO^2$の代わりに$AO \times \frac{PO^2}{4AS}$を書き、すべての項を$3PO$で割って$2AS$を掛けると、$\frac{4}{3}GH \times AS$（$= \frac{1}{6}AO \times PO + \frac{1}{2}AS \times PO = \frac{AO + 3AS}{6} \times PO = \frac{4AO - 3SO}{6} \times PO$ $=$ 面積$APO -$面積$SPO$）$=$ 面積$APS$となる。ところが$GH$は$3M$であったから、$\frac{4}{3}HG \times AS$は$4AS \times M$である。ゆえに面積$APS$は$4AS \times M$に等しい。\QED

\textit{系1.}\wsp{}したがって$GH$は$AS$に対して、物体が弧$AP$を描く時間が、物体が頂点$A$と、焦点$S$から軸に立てた垂線との間の弧を描く時間に対するのと同じ比をなす。

\textit{系2.}\wsp{}また、動く物体を通る円$ASP$がつねにそれを通っているとき、点$H$の速さは、物体が頂点$A$でもっていた速さに対して、3と8の比をなす。したがって線$GH$は、物体が$A$から$P$まで動くあいだに、頂点$A$でもっていた速さで描くことのできる直線に対して、同じ比をなす。

\textit{系3.}\wsp{}これによって逆に、任意に指定した弧$AP$を物体が描いた時間も求められる。$AP$を結び、その中点で垂線を立てて、直線$GH$と$H$で交わらせればよい。

\condpagelarge{補助定理XXVIII.}

\textit{任意に直線で切り取られた面積が、項の数と次元が有限の方程式によって一般に求められるような卵形（ovalis）の図形は存在しない。}

卵形の内に任意の点を与え、その点を極として、一本の直線を絶えず回転させる。その間に、その直線上で、極から動点を、その直線の卵形内の長さに比例する速さで、絶えず出発させる。この運動によって、その点は無限の巻きをもつ螺旋を描く。いま、もし卵形の面積が有限の方程式によって求められるなら、同じ方程式によって、この面積に比例する、極から点までの距離も求められ、したがって螺旋のすべての点が有限の方程式によって求められ、ゆえに、位置の与えられた任意の直線と螺旋との交点も、有限の方程式によって求められることになる。

ところが、限りなく延ばした任意の直線は、螺旋を無数の点で切り、二つの線の交点の一つを求める方程式は、その交点のすべてを同じ数の根として示すから、交点の数と同じだけの次元にのぼる。二つの円は互いに二点で切り合うから、一方の交点は、他方の交点も同時に求められる二次元の方程式によってしか求められない。二つの円錐曲線は四つの交点をもちうるから、そのどれかを一般に求めるには、すべてが同時に求められる四次元の方程式によるほかはない。なぜなら、それらの交点を別々に求めようとしても、すべての交点の法則と条件は同じだから、どの場合も同じ計算になり、したがって同じ結論になる。それゆえ、その結論はすべての交点を一度に含んで区別なく示さなければならない。このため、円錐曲線と三次曲線の交点も、六つありうるので、六次元の方程式によって一度に得られ、二つの三次曲線の交点は、九つありうるので、九次元の方程式によって一度に得られる。もしそうでなければ、すべての立体的な問題を平面的な問題に、立体より高い問題を立体的な問題に、還元できてしまうだろう。同じ理由で、二直線と円錐曲線の交点は、つねに二次元の方程式で、三直線と三次曲線の交点は三次元の方程式で、四直線と四次曲線の交点は四次元の方程式で得られ、以下限りなくそうである。したがって、直線の無数の交点は、すべての交点の法則も計算も同じであるため、次元と根の数が無限の方程式を必要とし、それによってすべてが同時に示される。

極から、その切る直線へ垂線を下ろし、その垂線を切る直線とともに極のまわりに回転させるとき、螺旋の交点は互いに移りかわり、最初のもの、すなわち最も近いものは、一回転ののちには二番目に、二回転ののちには三番目に、以下同様になる。その間、方程式は、切る直線の位置を決める量の大きさが変わるにつれて変わるだけである。したがって、それらの量が一回転ごとに最初の大きさにもどるとき、方程式も最初の形にもどり、同一の方程式がすべての交点を示すことになって、その根は無数にある。ゆえに、直線と螺旋の交点は、有限の方程式によって一般には求められない。したがって、直線で切り取られた面積が、そのような方程式によって一般に示される卵形は存在しない。

同じ論法で、螺旋を描く極と点の間隔を、卵形の切り取られた周の長さに比例してとるならば、周の長さも有限の方程式によっては一般に示されないことが示される。

\condpagelarge{\textit{系.}}

したがって、焦点から動く物体へ引いた動径によって描かれる楕円の面積は、与えられた時間から有限の方程式によっては得られず、それゆえ、幾何学的に有理な曲線の描画によっては決められない。私が幾何学的に有理な曲線と呼ぶのは、そのすべての点が、方程式で定められる長さ、すなわち長さの複雑な比によって決められる曲線であり、その他のもの（螺旋、円積曲線、トロコイドなど）を、幾何学的に無理な曲線と呼ぶ。実際、数と数との比であるか否かという長さは（『原論』第十巻のように）算術的に有理または無理な長さである。そこで私は、時間に比例する楕円の面積を、次のように、幾何学的に無理な曲線によって切り取る。

\condpagelarge{命題XXXI. 問題XXIII.}

\textit{与えられた楕円の軌道を動く物体について、指定された時刻における位置を求めること。}

楕円$APB$について、$A$を主頂点、$S$を焦点、$O$を中心とし、$P$を、求めたい物体の位置とする。$OA$を$G$まで、$OG : OA = OA : OS$となるように延ばす。垂線$GH$を立て、中心$O$、半径$OG$で円$EFG$を描き、定規$GH$を底として、輪$GEF$を、自分の軸のまわりに回転させながら進ませる。その間に、その輪の点$A$がトロコイド$ALI$を描くようにする。そうしたうえで、$GK$を、輪の円周$GEFG$に対して、物体が$A$から進んで弧$AP$を描いた時間が楕円上の一周期に対するのと同じ比になるようにとる。垂線$KL$を立ててトロコイドと$L$で交わらせ、$KG$に平行に引いた$LP$が、楕円と、求める物体の位置$P$で交わる。

\pngright{116.png}{}{}{}

実際、中心$O$、半径$OA$で半円$AQB$を描き、延長した$LP$が弧$AQ$と$Q$で交わるとし、$SQ$、$OQ$を結ぶ。$OQ$が弧$EFG$と$F$で交わるとし、$OQ$へ垂線$SR$を下ろす。面積$APS$は面積$AQS$に比例し、すなわち扇形$OQA$と三角形$OQS$との差に比例する。すなわち矩形$\frac{1}{2}OQ \times AQ$と$\frac{1}{2}OQ \times SR$との差に比例し、$\frac{1}{2}OQ$が与えられているから、弧$AQ$と直線$SR$との差に比例する。したがって（比$SR : $弧$AQ$の正弦、$OS : OA$、$OA : OG$、$AQ : GF$が等しく、分割して$AQ - SR : GF - $弧$AQ$の正弦であることから）、$GK$、すなわち弧$GF$と弧$AQ$の正弦との差に比例する。\QED

\condpagelarge{\textit{注.}}

\pngright{117.png}{}{}{}

ところで、この曲線を描くことの難しさのために、実際の機械的な作業では、近似の作図を用いるほうがよい。任意の楕円$APB$について、$AB$を長軸、$O$を中心、$S$を焦点、$OD$を短半軸、$AK$を通径の半分とする。$AS$を$G$で、$AG : AS = BO : BS$となるように切り、長さ$L$を、$\frac{1}{2}GK$に対して$AO^2$が矩形$AS \times OD$に対するのと同じ比になるように求める。$OG$を$C$で二等分し、中心$C$、半径$CG$で半円$GFO$を描く。最後に、角$GCF$を、四直角に対して、物体が求める弧$AP$を描いた与えられた時間が楕円の周期、すなわち一周の時間に対するのと同じ比になるようにとる。$AO$へ垂線$FE$を下ろし、それを$F$の方へ$N$まで、$EN$が長さ$L$に対して、その角の正弦$EF$が半径$CF$に対するのと同じ比になるように延ばす。中心$N$、半径$AN$で描いた円は、楕円を、物体の求める位置$P$で、ごく近似的に切る。

\pngright{118.png}{}{}{}

実際、周期の半分が終わると、物体$P$はつねに遠日点（最高の点）$B$にあり、残りの半分が終わると、必要なとおり近日点に戻っている。そして近点（遠日点・近日点）の近くにあるとき、生じつつある扇形$ASP$、$GCF$の最初の比と、消えてゆく扇形$BSP$、$OCF$の最後の比は、楕円全体と円全体との比に等しい。なぜなら、点$P$、$F$、$N$が、軸$AB$にきわめて近い位置$p$、$f$、$n$に落ちるとき、$An$と$pn$が等しいから、弧$Ap$に垂直で、したがって軸と点$K$で交わる直線$nq$は、弧$Ap$を二等分する。ゆえに$\frac{1}{2}Ap : Gn = AK : GK$、したがって$Ap : Gn = 2AK : GK$である。また$Gn : Gf = EN : EF$、すなわち$L : CF$、すなわち$\frac{GK \times AO^2}{2AS \times OD} : CF$、あるいは$GK \times AO^2 : 2AS \times OD \times CF$である。そして等比により$Ap : Gf = 2AK : GK$と、$GK \times AO^2 : 2AS \times OD \times CF$との合成、すなわち$AK \times AO^2 : AS \times OD \times CF$、すなわち（$AK \times AO$が$OD^2$に等しいから）$AO \times OD : AS \times CF$となる。したがって$Ap \times \frac{1}{2}AS : Gf \times \frac{1}{2}GC = AO \times OD \times AS : AS \times CF \times GC$、すなわち$AO \times OD : CG^2$、つまり生じつつある扇形$ASp$と、生じつつある扇形$GCf$との比は$AO \times OD : CG^2$であり、したがって楕円全体の面積と円全体の面積との比に等しい。\QED　消えてゆく扇形$BSP$、$OCF$の最後の比も、もっと長い論法で示すことができる。したがって点$P$の位置は、近点の近くでは十分に正確に求められている。求積点（quadratura、軌道の1/4の位置）では、誤差は楕円全体の面積の約500分の1か、それより少し大きくなるのがふつうである。しかしこの誤差は、次の、さらに進んだ作図によって、ほとんど消える。

点$G$、$O$を通って、適当な大きさの円弧$GTO$を引く。次に$EF$を両側に$T$、$N$まで、$EN : FT = \frac{1}{2}L : CF$となるように延ばす。中心$N$、半径$AN$で円を描いて、上のように楕円と$P$で交わらせる。弧$GTO$は、その点$T$を一つ求めることで決まる。それがその場合の作図を正確にする。

楕円の長軸が通径よりずっと大きく、頂点の近くの物体の運動を知りたいとき（彗星の理論でこの場合が起こる）、点$G$から軸$AB$に垂直な直線$GI$を引き、$GK$に対して、面積$AVPS$が矩形$AK \times AS$に対するのと同じ比をなす長さを$GI$ととり、中心$I$、半径$AI$で円を描くことができる。この円は、楕円を、物体の求める位置$P$で、ごく近似的に切るからである。同じ作図で（必要な変更を加えて）この問題は双曲線でも解ける。これらの作図は上と同じように証明され、もし図形が（遠い頂点$B$が無限遠に去って）放物線に変わるなら、問題XXIIの正確な作図に移る。

\pngright{119.png}{}{}{}

もしその位置$P$をもっと正確に決めなければならないときは、次のようにする。まず、角$B$を、半径に等しい弧が張る角$57\decimals{29578}$度に対して、焦点間の距離$SH$が楕円の直径$AB$に対するのと同じ比になるように求める。さらに長さ$L$を、半径に対して同じ比を逆にしたものになるように求める。これらが一度求められれば、問題は、次の解析によってしだいに解ける。

上の作図によって（あるいは何らかの推測によって）、物体の位置$P$を近似的に知る。軸へ縦線$PR$を下ろせば、楕円の直径の比から、外接円$AQB$の縦線$RQ$が分かる。それは、$AC$を半径として、角$ACQ$の正弦である。その角を粗い計算で近い数として求めれば足りる。また、時間に比例する角、すなわち、物体が弧$AP$を描いた時間が楕円上の一周期に対するのと同じ比で四直角に対する角を求める。その角を$N$とする。次に、角$D$を、角$B$に対して、角$ACQ$の正弦が半径に対するのと同じ比になるようにとり、角$E$を、角$N - ACQ + D$に対して、長さ$L$が、同じ長さ$L$から（角$ACQ + \frac{1}{2}D$が直角より小さいときは）その余弦を引いたもの、（大きいときは）その余弦を加えたものに対するのと同じ比になるようにとる。その後、角$F$を、角$B$に対して、角$ACQ + E$の正弦が半径に対するのと同じ比になるようにとり、角$G$を、角$N - ACQ - E + F$に対して、長さ$L$が、同じ長さ$L$から（角$ACQ + E + \frac{1}{2}F$が直角より小さいときは）その余弦を引いたもの、（大きいときは）加えたものに対するのと同じ比になるようにとる。第三回には、角$H$を、角$B$に対して、角$ACQ + E + G$の正弦が半径に対するのと同じ比になるようにとり、角$I$を、角$N - ACQ - E - G + H$に対して、長さ$L$が、同じ長さ$L$から（角$ACQ + E + G + \frac{1}{2}H$が直角より小さいときは）その余弦を引いたもの、（大きいときは）加えたものに対するのと同じ比になるようにとる。こうして限りなく続けることができる。最後に、角$ACq$を、角$ACQ + E + G + I + \ldots$に等しくとり、その余弦$Cr$と、正弦$qr$に対して楕円の短軸が長軸に対するのと同じ比をなす縦線$pr$とから、修正された物体の位置$p$が得られる。もし角$N - ACQ + D$が負になるときは、$E$の符号$+$をつねに$-$に、$-$を$+$に変えなければならない。同じことは$G$と$I$の符号についても、角$N - ACQ - E + F$や$N - ACQ - E - G + H$が負になるときに成り立つ。なお、無限級数$ACQ + E + G + I + \ldots$はきわめて速く収束するので、第二項$E$より先へ進む必要がほとんどない。この計算は、面積$APS$が、弧$AQ$と、焦点$S$から半径$CQ$へ垂直に下ろした直線との差に比例する、という定理にもとづいている。

\pngright{121.png}{}{}{}

同様の計算で、この問題は双曲線でも解ける。その中心を$C$、頂点を$A$、焦点を$S$、漸近線を$CK$とする。時間に比例する面積$APS$の大きさが分かっているとし、それを$A$とする。その面積をおよそ切り取る直線$SP$の位置を推測する。$CP$を結び、$A$と$P$から、漸近線に、他方の漸近線に平行に$AI$、$PK$を引き、対数表によって面積$AIKP$が、したがってこれに等しい面積$CPA$が得られる。これを三角形$CPS$から引けば、面積$APS$が残る。面積$A$と$APS$の半差$\frac{1}{2}APS - \frac{1}{2}A$（または$\frac{1}{2}A - \frac{1}{2}APS$）を、焦点$S$から接線$PT$へ下ろした垂線$SN$で割ると、長さ$PQ$が生じる。$PQ$は、面積$APS$が面積$A$より大きいときは$A$と$P$の間にとり、そうでなければ点$P$の反対側にとる。すると点$Q$が、物体のより正確な位置になる。計算を繰り返せば、いっそう正確な位置がつねに得られる。

これらの計算によって、問題は一般に解析的に解かれる。ただし天文学の用途には、次の個別の計算のほうが適している。楕円の半軸を$AO$、$OB$、$OD$、通径を$L$とするとき（先の、点$A$、$O$、$B$、$D$を示した図を見よ）、まず角$Y$を、その正接が半径に対して、半軸の差$AO - OD$が和$AO + OD$に対するのと同じ比になるように求め、次に角$Z$を、その正接が半径に対して、焦点間の距離$SH$と半軸の差$AO - OD$との矩形が、短半軸$OQ$と、長半軸と通径の4分の1との差$AO - \frac{1}{4}L$との矩形の3倍に対するのと同じ比になるように求める。これらの角が一度求められれば、物体の位置は次のように決まる。弧$BP$が描かれた時間に比例する角$T$、すなわち平均運動に等しい角をとる。そして角$V$（平均運動の第一の均差）を、角$Y$（最大の第一の均差）に対して、角$T$の2倍の正弦が半径に対するのと同じ比になるようにとる。さらに角$X$（第二の均差）を、角$Z$（最大の第二の均差）に対して、角$T$の2倍の正矢が半径の2倍に対するのと同じ比、あるいは（同じことだが）角$T$の正弦の平方が半径の平方に対するのと同じ比になるようにとる。角$T$、$V$、$X$の和$T + X + V$（角$T$が直角より小さいとき）、あるいは差$T + X - V$（角$T$が直角より大きく二直角より小さいとき）に等しい角$BHP$（均差を加えた平均運動）をとる。すると、$HP$が楕円と$P$で交わるとき、引いた$SP$は、時間に比例する面積$BSP$を近似的に切り取る。この方法は、極めて小さい角$V$、$X$（望むなら秒で表す）の最初の二、三桁を求めれば足りるので、十分に簡単に思われる。均差を加えた平均運動の角$BHP$が求められれば、真の運動の角$HSP$と距離$SP$は、私が深く敬愛するソールズベリー司教セス・ウォード博士の、よく知られた方法で、すぐに得られる。

ここまでは、曲線上の物体の運動について述べた。しかし、運動体が直線に沿って降りたり昇ったりすることもあるので、そうした運動にかかわることを、これから述べることにする。

\section{第1編 第7節：物体の直線的な上昇および下降（命題XXXII〜XXXIX）}


% Section: 第1編 第7節：物体の直線的な上昇および下降（命題XXXII〜XXXIX）



\begin{center}\textit{物体の直線的な上昇と下降について}\end{center}

\condpagelarge{命題XXXII. 問題XXIV.}

\textit{向心力が、場所の中心からの距離の2乗に反比例すると仮定して、物体が直線的に落下して、与えられた時間に描く空間（距離）を定めること。}

\pngright{123.png}{}{}{}

\textit{場合1.}\wsp{}もし物体が垂直に落ちないならば、物体は、下の焦点が中心と一致する、ある円錐曲線を描く。これは命題XI、XII、XIIIとそれらの系から明らかである。その円錐曲線を$ARPB$、下の焦点を$S$とする。まず、もしその図形が楕円ならば、その長軸$AB$の上に半円$ADB$を描き、落下する物体を通って、軸に垂直な直線$DPC$を引く。$DS$、$PS$を引けば、面積$ASD$は面積$ASP$に比例し、したがって時間にも比例する。軸$AB$を保ったまま、楕円の幅を絶えず小さくしてゆけば、面積$ASD$はつねに時間に比例したままである。その幅を限りなく小さくすれば、軌道$APB$は軸$AB$に、焦点$S$は軸の端$B$に一致し、物体は直線$AC$の上を落下し、面積$ABD$は時間に比例するようになる。したがって、物体が位置$A$から垂直に落下して与えられた時間に描く距離$AC$は、時間に比例して面積$ABD$をとり、点$D$から直線$AB$へ垂線$DC$を下ろせば、求められる。\QED

\pngright{126.png}{}{}{}

\textit{場合2.}\wsp{}もし上の図形$RPB$が双曲線ならば、同じ主直径$AB$の上に直角双曲線$BD$を描く。面積$CSP$、$CBfP$、$SPfB$は、面積$CSD$、$CBED$、$SDEB$に対して、それぞれ、高さ$CP$、$CD$の与えられた比にあり、また面積$SPfB$は、物体$P$が弧$PB$を動く時間に比例するから、面積$SDEB$もその同じ時間に比例する。双曲線$RPB$の通径を、横辺を保ったまま限りなく小さくすれば、弧$PB$は直線$CB$に、焦点$S$は頂点$B$に、直線$SD$は直線$BD$に一致する。したがって面積$BDEB$は、物体$C$が直線的に落下して線$CB$を描く時間に比例する。\QED

\textit{場合3.}\wsp{}同様の論法で、もし図形$RPB$が放物線ならば、同じ主頂点$B$のもとに、つねに与えられたままの別の放物線$BED$を描く。その間に、物体$P$が周上を動く先の放物線の通径を小さくして零にすれば、それは線$CB$と一致し、放物線の弓形$BDEB$は、その物体$P$または$C$が中心$B$へ落下する時間に比例する。\QED

\condpagelarge{命題XXXIII. 定理IX.}

\textit{これまでに見いだしたことを仮定して、私は言う。落下する物体の任意の位置$C$での速さは、中心$B$、半径$BC$の円周上を動く物体の速さに対して、物体の、円または双曲線の遠い頂点$A$からの距離$CA$が、図形の主半径$\frac{1}{2}AB$に対してもつ比の平方根の比をなす。}

実際、$CD$と$CP$が比例するから、直線$AB$は両方の図形$RPB$、$DEB$に共通の直径である。それを$O$で二等分し、図形$RPB$に点$P$で接して、その共通の直径$AB$（必要なら延長したもの）と$T$で交わる直線$PT$を引く。$SY$をこの直線へ、$BQ$をこの直径へ、垂線とし、図形$RPB$の通径を$L$とおく。定理VIIIの系9により、線$RPB$上を中心$S$のまわりに動く物体の、任意の位置$P$での速さは、同じ中心$S$のまわりに距離$SP$の円周上を動く物体の速さに対して、矩形$\frac{1}{2}L \times SP$と$SY^2$との比の平方根の比をなす。ところで円錐曲線の性質から、$AC \times CB : CP^2 = 2AO : L$であり、したがって$\frac{2CP^2 \times AO}{AC \times CB}$が$L$に等しい。ゆえにそれらの速さは互いに、$\frac{CP^2 \times AO \times SP}{AC \times CB}$と$SY^2$との比の平方根の比にある。さらに円錐曲線の性質から$CO : BO = BO : TO$であり、合比または分割により$CB : BT$に等しい。したがって分割または合成により$BO \mp CO : BO = CT : BT$、すなわち$AC : AO = CP : BQ$となる。そして$\frac{CP^2 \times AO \times SP}{AC \times CB}$は$\frac{BQ^2 \times AC \times SP}{AO \times BC}$に等しい。さて、図形$RPB$の幅$CP$を限りなく小さくして、点$P$が点$C$に、点$S$が点$B$に、線$SP$が線$BC$に、線$SY$が線$BQ$に一致するようにする。すると、いま線$CB$上を直線的に落下する物体の速さは、中心$B$、半径$BC$の円周上を動く物体の速さに対して、$\frac{BQ^2 \times AC \times SP}{AO \times BC}$と$SY^2$との比の平方根の比になり、すなわち（等しいという比$SP : BC$と$BQ^2 : SY^2$を無視して）、$AC : AO$の平方根の比になる。\QED

\pngright{126.png}{}{}{}

\textit{系.}\wsp{}点$B$と$S$が一致するとき、$TC : ST = AC : AO$となる。

\condpagelarge{命題XXXIV. 定理X.}

\textit{もし図形$BED$が放物線ならば、私は言う。落下する物体の任意の位置$C$での速さは、物体が中心$B$のまわりに、その距離$BC$の半分を半径として等速に円を描くことのできる速さに等しい。}

実際、放物線$RPB$を中心$S$のまわりに描く物体の、任意の位置$S$での速さは、（定理VIIIの系7により）、同じ$S$のまわりに距離$SP$の半分を半径として等速に円を描く物体の速さに等しい。放物線の幅$CP$を、放物線の弧$PfB$が直線$CB$に、中心$S$が頂点$B$に、間隔$SP$が間隔$BP$に一致するまで限りなく小さくすれば、命題は明らかになる。\QED

\condpagelarge{命題XXXV. 定理XI.}

\textit{同じ仮定のもとで、私は言う。不定の動径$SD$によって描かれる図形$DES$の面積は、物体が、図形$DES$の通径の半分に等しい半径で中心$S$のまわりを等速に回って、同じ時間に描くことのできる面積に等しい。}

\pngright{127.png}{}{}{}

実際、物体$C$が落下して、きわめて短い時間に微小な線$Cc$を描き、その間に別の物体$K$が、中心$S$のまわりに円$OKk$上を等速に回って弧$Kk$を描くと考える。垂線$CD$、$cd$を立てて、図形$DES$と$D$、$d$で交わらせる。$SD$、$SK$、$Sk$を結び、$Dd$を引いて軸$AS$と$T$で交わらせ、それへ垂線$SY$を下ろす。

\textit{場合1.}\wsp{}さて、もし図形$DES$が円または双曲線ならば、その主直径$AS$を$O$で二等分すれば、$SO$が通径の半分になる。そして$TC : TD = Cc : Dd$、$TD : TS = CD : SY$であるから、等比により$TC : TS = CD \times Cc : SY \times Dd$となる。ところが命題33の系により$TC : ST = AC : AO$である（点$D$、$d$が一致するときの線の最終の比をとるのである）。ゆえに$AC : AO$、すなわち$AC : SK$は、$CD \times Cc : SY \times Dd$に等しい。さらに、落下する物体の$C$での速さは、中心$S$のまわりに距離$SC$の円周上を動く物体の速さに対して、（定理IXにより）$AC : AO$、すなわち$SK$の平方根の比をなす。そしてこの速さは、円$OKk$を描く物体の速さに対して、（定理IVの系6により）$SK : SC$の平方根の比をなす。等比により、最初の速さと最後の速さ、すなわち線$Cc$と弧$Kk$との比は、$AC : SC$の平方根の比、すなわち$AC : CD$の比になる。したがって$CD \times Cc$は$AC \times Kk$に等しく、ゆえに$AC : SK = AC \times Kk : SY \times Dd$、したがって$SK \times Kk$は$SY \times Dd$に等しく、$\frac{1}{2}SK \times Kk$は$\frac{1}{2}SY \times Dd$に等しい。すなわち面積$KSk$は面積$SDd$に等しい。こうして、時間の一つ一つの微小な部分ごとに、二つの面積の微小な部分$KSk$、$SDd$が生じる。それらの大きさを小さくして数を限りなく増すとき、等しいという比に達するから、（補助定理IVの系により）同時に生じた全体の面積は、つねに等しい。\QED

\pngright{128a.png}{}{}{}

\textit{場合2.}\wsp{}もし図形$DES$が放物線ならば、上と同様に、$CD \times Cc : SY \times Dd = TC : ST$、すなわち2対1であることが見いだされ、したがって$\frac{1}{4}CD \times Cc$は$\frac{1}{2}SY \times Dd$に等しい。ところで落下する物体の$C$での速さは、（定理Xにより）距離$\frac{1}{2}SC$の円を等速に描くことができる速さに等しい。そしてこの速さは、半径$SK$の円を描くことができる速さ、すなわち線$Cc$と弧$Kk$との比で、（定理IVの系6により）$SK : \frac{1}{2}SC$の平方根の比、すなわち$SK : \frac{1}{2}CD$の比をなす。したがって$\frac{1}{2}SK \times Kk$は$\frac{1}{4}CD \times Cc$に等しく、ゆえに$\frac{1}{2}SY \times Dd$に等しい。すなわち面積$KSk$は面積$SDd$に等しく、上のとおりである。証明すべきことは以上である。

\condpagelarge{命題XXXVI. 問題XXV.}

\pngright{128b.png}{}{}{}

\textit{与えられた位置$A$から落下する物体の、落下の時間を定めること。}

直径$AS$（初めの、物体の中心からの距離）の上に半円$ADS$を描き、またそれに等しい半円$OKH$を中心$S$のまわりに描く。物体の任意の位置$C$から縦線$CD$を立てる。$SD$を結び、面積$ASD$に等しい扇形$OSK$をつくる。定理XIによって明らかなように、物体は落下して距離$AC$を、別の物体が中心$S$のまわりを等速に回って弧$OK$を描くことのできる、まさにその同じ時間で描く。求めるべきことは以上である。

\condpagelarge{命題XXXVII. 問題XXVI.}

\textit{与えられた位置から上向きまたは下向きに投げ出された物体の、上昇または下降の時間を定めること。}

\pngrightsc{129a.png}{}{}{}{}

物体が、与えられた位置$G$から、直線$ASG$に沿って、任意の速さで出発するとする。この速さと、物体が与えられた距離$SG$で中心$S$のまわりを回りうる円周上の等速の速さとの比の二乗の比で、$CA$を$\frac{1}{2}AS$に対してとる。もしその比が2対1ならば、点$A$は無限遠に落ち、この場合には頂点$S$、軸$SC$、任意の通径をもつ放物線を描かなければならない。これは定理Xから明らかである。もしその比が2対1より小さいか大きいかならば、前の場合には円を、後の場合には直径$SA$の上に直角双曲線を描かなければならない。これは定理IXから明らかである。次に、中心$S$、通径の半分に等しい半径で円$HKk$を描き、上昇または下降する物体の任意の二つの位置$G$、$C$に垂線$GI$、$CD$を立てて、円錐曲線または円と$I$、$D$で交わらせる。それから$SI$、$SD$を結び、弓形$SEIS$、$SEDS$に等しい扇形$HSK$、$HSk$をつくる。すると定理XIによって、物体$G$は、物体$K$が弧$Kk$を描くことのできる時間と同じ時間に、距離$GC$を描く。\QED

\condpagelarge{命題XXXVIII. 定理XII.}

\textit{向心力が、場所の中心からの高さ（距離）に比例すると仮定して、私は言う。落下する物体の時間、速さ、描く空間（距離）は、それぞれ、弧、弧の正矢、弧の正弦に比例する。}

\pngright{129b.png}{}{}{}

物体が、任意の位置$A$から直線$AS$に沿って落下するとし、力の中心$S$、半径$AS$で四分円$AE$を描き、$CD$を任意の弧$AD$の正弦とする。すると物体$A$は、時間$AD$で距離$AC$を落下して描き、位置$C$で速さ$CD$を得る。これは、命題XXXIIが命題XIから示されたのと同じように、命題Xから示される。\QED

\textit{系1.}\wsp{}したがって、一つの物体が位置$A$から落下して中心$S$に至る時間と、別の物体が回りながら四分円の弧$ADE$を描く時間は、等しい。

\textit{系2.}\wsp{}ゆえに、物体が任意の位置から中心まで落下する時間はすべて等しい。なぜなら、回る物体の周期はすべて、（命題IVの系3により）等しいからである。

\condpagelarge{命題XXXIX. 問題XXVII.}

\textit{どんな種類の向心力が仮定されても、また曲線図形の求積が許されるならば、直線的に上昇または下降する物体の、それぞれの位置での速さと、物体が任意の位置に達するまでの時間とを求めること。またその逆を求めること。}

\pngright{130.png}{}{}{}

任意の位置$A$から直線$ADEC$上を物体$E$が落下するとし、その位置$E$に、その場所で中心$C$に向かう向心力に比例する垂線$EG$をつねに立てる。$BFG$を、点$G$が絶えず接する曲線とする。運動の初めには$EG$が垂線$AB$に一致するものとする。すると任意の位置$E$での物体の速さは、曲線の面積$ABGE$の平方根に比例する。\QED　$EG$の上に、面積$ABGE$の平方根に反比例して$EM$をとり、$ALM$を、点$M$が絶えず接する曲線とすれば、物体が落下して線$AE$を描く時間は、曲線の面積$ALME$に比例する。求めるべきことは以上である。

実際、直線$AE$の上に、与えられた長さのごく小さい線$DE$をとり、$DLF$を、物体が$D$にあったときの線$EMG$の位置とする。もし向心力が、面積$ABGE$の平方根が落下する物体の速さに比例するようなものであれば、その面積自体は速さの二乗に比例する。すなわち、$D$と$E$での速さの代わりに$V$と$V + I$を書けば、面積$ABFD$は$V^2$に、面積$ABGE$は$V^2 + 2VI + I^2$に比例し、分割して面積$DFGE$は$2VI + I^2$に比例する。したがって$\frac{DFGE}{DE}$は$\frac{2I \times V + I^2}{DE}$に比例し、すなわち、生じつつある量の最初の比をとるならば、長さ$DF$は量$\frac{2I \times V}{DE}$に、したがってその半分の$\frac{I \times V}{DE}$にも比例する。ところで、物体が落下して微小な線$DE$を描く時間は、その線に正比例し、速さ$V$に反比例する。また力は、速さの増加$I$に正比例し、時間に反比例する。したがって、生じつつある量の最初の比をとるならば、$\frac{I \times V}{DE}$に、すなわち長さ$DF$に比例する。ゆえに$DF$または$EG$に比例する力は、物体を、面積$ABGE$の平方根に比例する速さで落下させる。\QED

さらに、与えられた長さの任意の微小な線$DE$を描く時間は、速さに比例し、したがって面積$ABFD$の平方根に反比例する。$DL$も、したがって生じつつある面積$DLME$も、同じ平方根に反比例する。ゆえに時間は面積$DLME$に比例し、すべての時間の和は、すべての面積の和に比例する。すなわち（補助定理IVの系により）線$AE$を描く全時間は、全面積$AME$に比例する。\QED

\textit{系1.}\wsp{}もし$P$が、物体が落下しなければならない位置であって、ある一様な、知られた向心力（ふつう仮定される重力のようなもの）に押されて、位置$D$で、別の物体が何らかの力で落下して同じ位置$D$で得た速さに等しい速さを得るものとする。そして垂線$DF$の上に、$DR$を、その一様な力が位置$D$での他方の力に対するのと同じ比で$DF$に対するようにとり、矩形$PDRQ$を完成させて、これに等しい面積$ABFD$を切り取れば、$A$が、別の物体が落下してきた位置になる。実際、矩形$EDRS$を完成させるとき、面積$ABFD$と面積$DFGE$との比は$V^2 : 2V \times I$、したがって$\frac{1}{2}V : I$、すなわち、全体の速さの半分と、不均一な力で落下する物体の速さの増加との比である。同様に、面積$PQRD$と面積$DRSE$との比は、全体の速さの半分と、一様な力で落下する物体の速さの増加との比である。これらの増加は、（生じつつある時間が等しいから）生じさせる力、すなわち縦線$DF$、$DR$に比例し、したがって生じつつある面積$DFGE$、$DRSE$に比例する。ゆえに（等比により）全体の面積$ABFD$、$PQRD$は、互いに全体の速さの半分の比にあり、したがって（速さが等しいから）等しい。

\textit{系2.}\wsp{}したがって、任意の物体が、任意の位置$D$から、与えられた速さで上向きまたは下向きに投げ出され、向心力の法則が与えられるならば、その物体の他の任意の位置$e$での速さは、縦線$eg$を立てて、その速さを、位置$D$での速さに対して、矩形$PQRD$に曲線の面積$DFge$を加えたもの（位置$e$が位置$D$より下にあるとき）、あるいは引いたもの（上にあるとき）の平方根が、矩形$PQRD$だけの平方根に対するのと同じ比にとることによって求められる。すなわち$\sqrt{PQRD \pm DFge} : \sqrt{PQRD}$である。

\textit{系3.}\wsp{}時間も、縦線$em$を、$PQRD \pm DFge$の平方根に反比例してたて、物体が線$De$を描いた時間を、別の物体が一様な力で$P$から落ちて、落ちて$D$に達した時間に対して、曲線の面積$DLme$が矩形$2PD \times DL$に対するのと同じ比にとることによって、分かる。実際、一様な力で下降する物体が線$PD$を描いた時間は、同じ物体が線$PE$を描いた時間に対して、$PD : PE$の平方根の比をなす。すなわち（微小な線$DE$がいままさに生じつつあるとき）$PD : PD + \frac{1}{2}DE$、あるいは$2PD : 2PD + DE$の比をなし、分割して、その同じ物体が微小な線$DE$を描いた時間に対して$2PD : DE$、したがって矩形$2PE \times DL$が面積$DLME$に対するのと同じ比をなす。そして、一方の物体が微小な線$DE$を描いた時間は、他方の物体が不均一な運動で線$De$を描いた時間に対して、面積$DLME$が面積$DLme$に対するのと同じ比をなす。等比により、最初の時間と最後の時間の比は、矩形$2PD \times DL$が面積$DLme$に対するのと同じ比になる。

\chapter{第1編 物体の運動（第8節〜第10節：任意中心力・近点移動・拘束面振子）}



\section{第1編 第8節：任意の向心力によって巡る物体の軌道の決定（命題XL〜XLII）}


% Section: 第1編 第8節：任意の向心力によって巡る物体の軌道の決定（命題XL〜XLII）



\begin{center}\textit{任意の向心力に動かされて物体が回る軌道を求めること}\end{center}

\condpagelarge{命題XL. 定理XIII.}

\textit{物体が、任意の向心力に押されて任意の仕方で動き、別の物体が直線的に上昇または下降するとする。そして、等しい高さのどこか一つの位置で二つの物体の速さが等しいならば、すべての等しい高さにおいて二つの速さは等しい。}

\pngcent{136.png}{}{}

ある物体が$A$から$D$、$E$を通って中心$C$へ落下し、別の物体が$V$から曲線$VIKk$に沿って動くとする。中心$C$、任意の半径で同心円$DI$、$EK$を描き、これらが直線$AC$と$D$、$E$で、曲線$VIK$と$I$、$K$で交わるとする。$IC$を結んで$KE$と$N$で交わらせ、$IK$へ垂線$NT$を下ろす。円周の間隔$DE$、すなわち$IN$をきわめて小さくし、物体は$D$と$I$で等しい速さをもつとする。距離$CD$、$CI$が等しいから、$D$と$I$での向心力は等しい。これらの力を、等しい微小な線$DE$、$IN$で表す。力$IN$を（運動の法則の系2により）二つの力$NT$と$IT$に分解すると、力$NT$は、物体の進路$ITK$に垂直な線$NT$の方向にはたらくので、その進路での物体の速さを少しも変えず、物体を直線の進路から引きもどすだけで、物体を軌道の接線から絶えずそらせて、曲線の道$ITKk$の上を進ませる。この結果を生むことにこの力はすべて使われる。もう一方の力$IT$は、物体の進路に沿ってはたらくので、すべてが物体を加速し、与えられたごく短い時間に、それ自身に比例する加速を生む。したがって、$D$と$I$での物体の加速は、等しい時間では（生じつつある線$DE$、$IN$、$IK$、$IT$、$NT$の最初の比をとるならば）線$DE$、$IT$に比例し、等しくない時間では、それらの線と時間とを合成したものに比例する。時間は、速さが等しいから、描かれる道$DE$、$IK$に比例し、したがって線$DE$と$IK$による物体の進行の間の加速は、$DE$と$IT$、$DE$と$IK$を合成した比、すなわち$DE^2$と矩形$IT \times IK$の比になる。ところが矩形$IT \times IK$は$IN^2$に等しく、すなわち$DE^2$に等しい。したがって、物体が$D$、$I$から$E$、$K$へ移るあいだに生じる加速は等しい。ゆえに$E$と$K$での物体の速さは等しく、同じ論法で、それに続く等しい距離でも、つねに等しいことが見いだされる。\QED　また同じ論法で、等しい速さで、中心から等しい距離にある物体は、等しい距離まで上昇するあいだ、等しく減速される。\QED

\textit{系1.}\wsp{}したがって、もし物体や糸でつるした振子が振動するか、きわめてなめらかで完全に摩擦のない何かの妨げによって曲線上を動くように強いられ、別の物体が直線的に上昇または下降し、それらの速さが任意の同じ高さで等しいならば、それらの速さは、他のあらゆる等しい高さでも等しい。なぜなら、絶対的に滑らかな容器の妨げは、横向きの力$NT$と同じ結果をもたらすからである。物体はそれによって減速も加速もされず、ただ直線の進路から離れるように強いられるだけである。

\textit{系2.}\wsp{}さらに、もし量$P$が、振動する物体、または任意の軌道を回る物体が、軌道の任意の点で、そこでもつ速さで上向きに投げ出されたとき、到達しうる最大の中心からの距離であるとし、量$A$が、軌道の他の任意の点での物体の中心からの距離であって、向心力がつねに$A$の任意の冪$A^{n-1}$（その指数$n-1$は、任意の数$n$から1を引いたもの）に比例するならば、あらゆる高さ$A$での物体の速さは、$\sqrt{nP^n - nA^n}$に比例し、したがって与えられる。なぜなら、上昇または下降する物体の速さは、（命題XXXIXにより）まさにこの比にあるからである。

\condpagelarge{命題XLI. 問題XXVIII.}

\textit{どんな種類の向心力が仮定されても、また曲線図形の求積が許されるならば、物体が動く軌道と、求めた軌道における運動の時間とを求めること。}

\pngcent{136.png}{}{}

任意の力が中心$C$に向かってはたらくとし、軌道$VITKk$を求めるとする。中心$C$、任意の半径$CV$で円$VXY$を描き、同じ中心で、軌道と$I$、$K$で、直線$CV$と$D$、$E$で交わる他の任意の円$ID$、$KE$を描く。次に、円$KE$、$VY$と$N$、$X$で交わる直線$CNIX$を引き、円$VXY$と$Y$で交わる直線$CKY$を引く。点$I$、$K$は互いにきわめて近いものとし、物体は$V$から$I$、$T$、$K$を通って$k$へ進むとする。$A$を、別の物体が落下して、位置$D$で、前の物体の$I$での速さに等しい速さを得るべき高さとする。命題XXXIXの仮定のもとで、与えられたごく短い時間に描かれる微小な線$IK$は、速さに、したがって面積$ABFD$の平方根に比例し、また時間に比例する三角形$ICK$が与えられるから、$KN$は高さ$IC$に反比例する。すなわち、ある量$Q$が与えられ、高さ$IC$を$A$とすれば、$\frac{Q}{A}$に比例する。これを$Z$とおく。$Q$の大きさを、ある場合に$\sqrt{ABFD}$が$Z$に対して$IK : KN$になるようなものとすれば、つねに$\sqrt{ABFD} : Z = IK : KN$、$ABFD : Z^2 = IK^2 : KN^2$、分割して$ABFD - Z^2 : Z^2 = IN^2 : KN^2$であり、したがって$\sqrt{ABFD - Z^2} : Z = IN : KN$となり、ゆえに$A \times KN$は$\frac{Q \times IN}{\sqrt{ABFD - Z^2}}$に等しい。

そこで、$YX \times XC$が$A \times KN$に対して、$YC : KC$の二乗比にあるから、矩形$YX \times XC$は$\frac{Q \times IN \times CX^2}{A^2 \sqrt{ABFD - Z^2}}$に等しい。したがって、垂線$DF$の上に、$Db$、$Dc$を、それぞれ$\frac{Q}{2\sqrt{ABFD - Z^2}}$と$\frac{Q \times CX^2}{2A^2 \sqrt{ABFD - Z^2}}$に等しくつねにとり、点$b$、$c$が絶えず接する曲線$ab$、$cd$を描く。さらに点$V$から直線$AC$へ垂線$Vad$を立て、曲線の面積$VDba$、$VDdc$を切り取り、縦線$Ez$、$Ex$も立てる。矩形$Db \times IN$、すなわち$DbzE$は、矩形$A \times KN$の半分、すなわち三角形$ICK$に等しく、矩形$Dc \times IN$、すなわち$Dc \times E$は、矩形$YX \times CX$の半分、すなわち三角形$XCY$に等しい。すなわち、面積$VDba$、$VIC$の生じつつある微小な部分$DbzE$、$ICK$はつねに等しく、面積$VDcd$、$VCX$の生じつつある微小な部分$DExc$、$XCY$はつねに等しいから、生じた面積$VDba$は生じた面積$VIC$に等しく、したがって時間に比例する。また生じた面積$VDdc$は、生じた扇形$VCX$に等しい。ゆえに、物体が位置$V$を離れてからの任意の時間が与えられれば、それに比例する面積$VDba$が与えられ、それによって物体の高さ$CD$、すなわち$CI$が与えられる。また面積$VDcd$と、それに等しい扇形$VCX$が、その角$VCI$とともに与えられる。角$VCI$と高さ$CI$が与えられれば、その時間が経ったとき物体がいる位置$I$が与えられる。\QED

\textit{系1.}\wsp{}これによって、物体の最大・最小の高さ、すなわち軌道の遠近点（アプシス）を、手早く求めることができる。遠近点は、中心を通って引いた直線$IC$が軌道$VIK$に垂直に当たる点に落ちる。これは、直線$IK$と$NK$が等しいとき、したがって面積$ABFD$が$Z^2$に等しいときに起こる。

\textit{系2.}\wsp{}また、軌道が直線$IC$を別の点で切る角$KIN$は、物体の与えられた高さ$IC$から手早く求められる。すなわち、その正弦を、半径に対して$KN : IK$、すなわち$Z$と面積$ABFD$の平方根との比にとればよい。

\textit{系3.}\wsp{}中心を$C$、主頂点を$V$として、任意の円錐曲線$VRS$を描き、その任意の点$R$から接線$RT$を引いて、限りなく延ばした軸$CV$と点$T$で交わらせる。次に$CR$を結んで、直線$CP$を、切り取られた$CT$に等しくなるように、また角$VCP$を扇形$VCR$に比例するように引く。そして、中心$C$に向かって、場所の中心からの距離の3乗に反比例する向心力がはたらき、物体が位置$V$から、直線$CV$に垂直な方向に、適切な速さで出発するとする。すると、その物体は、点$P$が絶えず接する軌道の上を進む。したがって、もし円錐曲線$CVRS$が双曲線ならば、物体は中心へ落ちてゆく。もしそれが楕円ならば、物体は絶えず上昇して無限遠へ去る。逆に、物体が位置$V$から任意の速さで出発し、出発したとおりに、斜めに中心へ降りるか、中心から斜めに昇るかするとき、図形$CVRS$が双曲線か楕円かは、角$VCP$を与えられたある比で増減することで、その軌道を求められる。さらに、向心力が遠心力に変わるならば、物体は、角$VCP$を楕円の扇形$CVRC$に比例させ、長さ$CP$を長さ$CT$に等しくして求められる軌道$VPQ$の上を、斜めに上昇する。これらすべては、前の命題から、ある曲線の求積によって導かれる。その求積の求め方は、かなり容易なので、簡潔さのために省く。

\condpagelarge{命題XLII. 問題XXIX.}

\textit{向心力の法則が与えられたとき、与えられた位置から、与えられた速さで、与えられた直線に沿って出発した物体の運動を求めること。}

前の三つの命題と同じ仮定のもとで、物体が位置$I$から微小な線$IT$に沿って、別の物体が一様な向心力によって位置$P$から落下して$D$で得ることのできる速さと同じ速さで出発するとする。この一様な力と、最初の物体が$I$で押される力との比は、$DR : DF$であるとする。物体は$k$へ進み、中心$C$、半径$Ck$で円$ke$を描いて、直線$PD$と$e$で交わらせ、曲線$ALMm$、$BFGg$、$abzv$、$dcxw$の縦線$em$、$eg$、$ev$、$ew$を立てる。与えられた矩形$PDRQ$と、最初の物体を動かす向心力の与えられた法則とから、曲線$BFGg$、$ALMm$が、問題XXVIIIの作図とその系1によって与えられる。次に、与えられた角$CIT$から、生じつつある$IK$、$KN$の比が与えられ、そこから、問題XXVIIIの作図によって、量$Q$が、曲線$abzv$、$dcxw$とともに与えられる。したがって、任意の時間$Dbve$が経過したとき、物体の高さ$Ce$、すなわち$Ck$と、面積$Dcwe$、これに等しい扇形$XCy$と、角$XCy$と、物体がそのときいる位置$k$とが与えられる。\QED

これらの命題で私たちは、向心力が中心から遠ざかるにつれて、人が想像しうるどんな法則にしたがって変化してもよいが、中心から等しい距離では、どこでも同じであると仮定している。ここまで私たちは、動かない軌道における物体の運動を考えてきた。力の中心のまわりを回る軌道における物体の運動について、少し付け加えることが残っている。
\section{第1編 第9節：運動する軌道における物体の運動と近点移動（命題XLIII〜XLV）}


% Section: 第1編 第9節：運動する軌道における物体の運動と近点移動（命題XLIII〜XLV）





\begin{center}\textit{動く軌道における物体の運動と、遠近点（アプシス）の運動について}\end{center}

\condpagelarge{命題XLIII. 問題XXX.}

\textit{力の中心のまわりを回る任意の軌道において、物体を、同じ軌道が静止しているときの別の物体と同じように動かすことができるようにすること。}

\pngright{140.png}{}{}{}

位置の与えられた軌道$VPK$の上を、物体$P$が$V$から$K$へ向かって進みながら回るとする。中心$C$から、つねに$CP$に等しい$Cp$を引き、角$VCp$を角$VCP$に比例させる。すると、線$Cp$が描く面積と、線$CP$が描く面積$VCP$との比は、線$Cp$を描くものの速さと線$CP$を描くものの速さとの比に等しい。すなわち角$VCp$と角$VCP$との比に等しく、したがって与えられた比であり、ゆえに時間に比例する。面積が、線$Cp$が不動の平面上に描く面積が時間に比例するのだから、物体は、適切な大きさの向心力に押されて、点$p$とともに、その同じ点$p$が上に述べた仕方で不動の平面上に描く曲線の上を、回ることができるのは明らかである。角$VCv$を角$PCp$に等しくし、線$Cv$を線$CV$に等しくして、図形$vCp$を図形$VCP$に等しくする。すると、つねに$p$にある物体は、回転する図形$vCp$の周上を動き、静止した図形$VPK$の上で別の物体$P$が、それに相似で等しい弧$VP$を描くのと同じ時間に、その弧$vp$を描く。そこで、命題VIの系によって、点$p$が不動の平面上に描くその曲線の上を物体が回るための向心力を求めれば、問題は解かれる。\QED

\condpagelarge{命題XLIV. 定理XIV.}

\textit{静止した軌道を動く物体と、同じ軌道を回転しながら動く別の物体とを、同じように動かすことのできる二つの力の差は、共通の高さの3乗に反比例する。}

静止した軌道の部分$VP$、$PK$に、回転する軌道の相似で等しい部分$vp$、$pk$をとる。点$k$から直線$pC$へ垂線$kr$を下ろし、それを$m$まで延ばして、$mr : kr$が角$VCp$と角$VCP$との比に等しくなるようにする。物体の高さ$PC$と$pC$、$KC$と$kC$は、つねに等しいから、位置$P$、$p$にある物体のそれぞれの運動を（運動の法則の系2により）二つに分けるとき（その一つは中心へ、すなわち線$PC$、$pC$に沿い、もう一つはそれに垂直な、線$PC$、$pC$に垂直な方向に決まる）、中心へ向かう運動は等しく、物体$p$の横向きの運動は、物体$P$の横向きの運動に対して、線$pC$の角運動と線$PC$の角運動との比、すなわち角$VCp$と角$VCP$との比をなすことは明らかである。したがって、物体$P$がその運動によって点$K$に着く時間に、物体$p$は、中心へ等しい運動によって、$P$から$C$の方へ等しく動かされ、したがってその時間が経ったとき、点$k$を通って直線$pC$に垂直な線$mkr$上のどこかにある。そして横向きの運動によって、直線$pC$からの距離を、もう一方の物体が直線$PC$から得る距離に対して、この運動の横向きの運動の、もう一方の横向きの運動に対する比にあるだけ得る。ところで$kr$は、もう一方の物体が直線$pC$から得る距離に等しく、$mr : kr$は角$VCp$と角$VCP$との比、すなわち物体$p$の横向きの運動と物体$P$の横向きの運動との比に等しいから、物体$p$は、その時間が経ったとき位置$m$にあることは明らかである。これは、物体$P$と$p$が線$pC$、$PC$に沿って等しく動くとき、したがってそれらの線に沿って等しい力で押されるときに成り立つ。

そこで、角$pCn$を角$pCk$に対して、角$VCp$が角$VCP$に対するのと同じ比になるようにとり、$nC$を$kC$に等しくすれば、物体$p$は、その時間が経ったとき実際には$n$にある。したがって、もし角$mCp$が角$kCp$より大きいなら、すなわち軌道$Vpk$が順行（進む向き）に動くなら、より大きな力で押され、小さいなら、すなわち軌道が逆行するなら、より小さい力で押される。そして力の差は、位置の間隔$mn$に比例する。物体$p$は、その力のはたらきによって、与えられたその時間のあいだにこの間隔だけ移されなければならないのである。中心$C$、半径$Cn$または$Ck$の円を描いたと考え、これが延長した線$mr$、$mn$と$s$、$t$で交わるとすると、矩形$mn \times mt$は矩形$mk \times ms$に等しく、したがって$mn$は$\frac{mk \times ms}{mt}$に等しい。さて三角形$pCk$、$pCn$は大きさが与えられているから、$kr$と$mr$、およびそれらの差$mk$と和$ms$は、高さ$pC$に反比例し、したがって矩形$mk \times ms$は高さ$pC$の平方に反比例する。また$mt$は$\frac{1}{2}mt$に、すなわち高さ$pC$に正比例する。これらは生じつつある線の最初の比であり、したがって$\frac{mk \times ms}{mt}$、すなわち生じつつある微小な線$mn$と、それに比例する力の差は、高さ$pC$の3乗に反比例する。\QED

\textit{系1.}\wsp{}したがって、位置$P$と$p$、あるいは$K$と$k$での力の差は、物体$P$が不動の軌道で弧$PK$を描くのと同じ時間に、物体が$r$から$k$まで円運動で回ることのできる力に対して、$mk \times ms$が$rk^2$に対するのと同じ比をなす。すなわち、角$VCP$と角$VCp$の比にしたがって、与えられた量$F$、$G$をとるならば、$G^2 - F^2$と$F^2$の比をなす。したがって、中心$C$、任意の半径$CP$または$Cp$で、物体$P$が不動の軌道で回りながら、ある時間に中心へ引いた動径によって描いた全面積$VPC$に等しい円の扇形を描くとき、不動の軌道で物体$P$を、動く軌道で物体$p$を回らせる力の差は、ある物体が、中心へ引いた動径によって、面積$VPC$が描かれたのと同じ時間にその扇形を一様に描くことのできる向心力に対して、$G^2 - F^2$と$F^2$の比をなす。なぜなら、その扇形と面積$pCk$は、それらが描かれる時間に比例するからである。

\textit{系2.}\wsp{}もし軌道$VPK$が、焦点を$C$、遠い方の遠近点を$V$とする楕円であって、それに相似で等しい楕円$vpk$が、つねに$pC$が$PC$に等しく、角$VCp$と角$VCP$の比が与えられた$G : F$になるように置かれるならば、また高さ$PC$または$pC$の代わりに$A$を、楕円の通径の代わりに$2R$を書くならば、動く楕円の上を物体が回ることのできる力は、$\frac{F^2}{A^2} + \frac{RG^2 - RF^2}{A^3}$に比例し、逆もまた成り立つ。実際、不動の楕円の上を物体が回る力を量$\frac{F^2}{A^2}$で表せば、$V$での力は$\frac{F^2}{CV^2}$になる。ところで、距離$CV$の円の上を、楕円を回る物体が$V$でもつ速さで物体が回ることのできる力は、楕円を回る物体が遠近点$V$で押される力に対して、楕円の通径の半分が円の半径$CV$に対するのと同じ比をなし、したがって$\frac{RF^2}{CV^3}$に等しい。そしてこれに対して$G^2 - F^2$が$F^2$に対する比をなす力は、$\frac{RG^2 - RF^2}{CV^3}$に等しい。この力は（この系1により）、不動の楕円$VPK$の上の物体$P$と、動く楕円$vpk$の上の物体$p$とを回らせる力の差である。ゆえに（この命題により）その差は、他の任意の高さ$A$では、高さ$CV$でのそれ自身に対して$\frac{1}{A^3} : \frac{1}{CV^3}$の比にあるから、あらゆる高さ$A$で同じ差は$\frac{RG^2 - RF^2}{A^3}$に等しい。したがって、不動の楕円$VPK$の上で物体を回すことのできる力$\frac{F^2}{A^2}$に、余りの$\frac{RG^2 - RF^2}{A^3}$を加えれば、動く楕円$vpk$の上で同じ時間に物体を回すことのできる全体の力$\frac{F^2}{A^2} + \frac{RG^2 - RF^2}{A^3}$ができる。

\textit{系3.}\wsp{}同じ仕方で、もし不動の軌道$VPK$が、中心を力の中心$C$にもつ楕円であって、それに相似で等しい同心の動く楕円$vpk$が置かれ、この楕円の通径を$2R$、横辺（長軸）を$2T$とし、角$VCp$がつねに角$VCP$に対して$G : F$であるならば、不動の楕円と動く楕円の上で、等しい時間に物体を回すことのできる力は、それぞれ$\frac{F^2 A}{T^3}$と$\frac{F^2 A}{T^3} + \frac{RG^2 - RF^2}{A^3}$に比例する、ということが導かれる。

\textit{系4.}\wsp{}一般に、もし物体の最大の高さ$CV$を$T$とし、軌道$VPK$が$V$でもつ曲率の半径、すなわち同じ曲がり方の円の半径を$R$とし、任意の不動の軌道$VPK$の上で物体を回すことのできる向心力を、位置$V$で$\frac{F^2}{T^2}$と呼び、他の位置$P$では不定に$X$と呼び、高さ$CP$を$A$と呼び、$G$を$F$に対して角$VCp$と角$VCP$の与えられた比にとるならば、同じ物体が、円運動をする同じ軌道$vpk$の上で、同じ時間に同じ運動を行うことのできる向心力は、力の和$X + \frac{VR G^2 - VR F^2}{A^3}$に比例する。

\textit{系5.}\wsp{}したがって、任意の不動の軌道における物体の運動が与えられれば、力の中心のまわりでのその角運動を、与えられた比で増減することができ、そこから、物体が新しい向心力によって回る新しい不動の軌道を求めることができる。

\pngright{145.png}{}{}{}

\textit{系6.}\wsp{}したがって、位置の与えられた直線$CV$に、不定の長さの垂線$VP$を立て、$PC$を結んで、それに等しい$Cp$を引き、角$VCp$を角$VCP$に対して与えられた比にするならば、点$p$が絶えず接するその曲線$Vpk$の上を物体が回りうる力は、高さ$Cp$の3乗に反比例する。なぜなら、物体$P$は、慣性の力によって、ほかのいかなる力にも押されなければ、直線$VP$上を一様に進むことができるからである。高さ$CP$、$Cp$の3乗に反比例する中心$C$への力を加えれば、（すでに示したことによって）その直線の運動は、曲線$Vpk$へ曲げられる。この曲線$Vpk$は、命題XLIの系3で見いだされた曲線$VPQ$と同じであり、そこで、このような力に引かれる物体は、斜めに上昇する、と私たちが述べた曲線である。

\condpagelarge{命題XLV. 問題XXXI.}

\textit{円に最も近い軌道の、遠近点の運動を求めること。}

この問題は、算術的に、次のようにして解かれる。すなわち、上の命題の系2または系3のように、動く楕円の上を回る物体が不動の平面上に描く軌道を、遠近点を求めたい軌道の形に近づけ、その物体が不動の平面上に描く軌道の遠近点を求めるのである。ところで軌道が同じ形になるのは、それらを描く向心力が、互いに比べて、等しい高さで比例するようになるときである。点$V$を遠い方の遠近点とし、最大の高さ$CV$の代わりに$T$を、他の任意の高さ$CP$または$Cp$の代わりに$A$を、高さの差$CV - CP$の代わりに$X$を書く。楕円の焦点$C$のまわりを回る楕円（系2のように）に動かされる物体を動かす力は、系2では$\frac{F^2}{A^2} + \frac{RG^2 - RF^2}{A^3}$に比例した。これはすなわち$\frac{F^2 A + RG^2 - RF^2}{A^3}$に比例する。$A$の代わりに$T - X$を代入すると、$\frac{RG^2 - RF^2 + TF^2 - F^2 X}{A^3}$に比例する。同様に、他の任意の向心力も、分母が$A^3$であるような分数に直さなければならず、分子は、対応する項を比べて、比例するように定めなければならない。例によって明らかになる。

\textit{例1.}\wsp{}向心力が一様である、すなわち$\frac{A^3}{A^3}$に比例するとしよう。（分子で$A$の代わりに$T - X$を書けば）$\frac{T^3 - 3T^2 X + 3TX^2 - X^3}{A^3}$に比例する。分子の対応する項、すなわち与えられた項と与えられた項、与えられていない項と与えられていない項を比べれば、$RG^2 - RF^2 + TF^2 : T^3 = -F^2 X : -3T^2 X + 3TX^2 - X^3$、すなわち$-F^2 : -3T^2 + 3TX - X^2$となる。いま軌道は円に最も近いとするから、軌道が円と一致するとし、$R$と$T$を等しくおき、$X$を限りなく小さくすると、最終の比は$RG^2 : T^3 = -F^2 : -3T^2$、すなわち$G^2 : T^2 = F^2 : 3T^2$、入れ替えて$G^2 : F^2 = T^2 : 3T^2$、すなわち$1 : 3$となる。したがって$G : F$、すなわち角$VCp$と角$VCP$の比は$1 : \sqrt{3}$である。ゆえに、不動の楕円の上の物体が、遠い方の遠近点から近い方の遠近点へ降りるあいだに、角$VCP$（と言ってよければ）180度をなすとき、動く楕円の上の別の物体、したがってここで問題にしている不動の軌道の上の物体は、遠い方の遠近点から近い方の遠近点へ降りるあいだに、角$VCp$を$\frac{180}{\sqrt{3}}$度なす。これは、一様な向心力のもとで物体が描くこの軌道と、回る楕円のなかで回転を繰り返す物体が静止した平面上に描くあの軌道とが相似であることによる。上の項の比較によって、これらの軌道は、一般にではなく、円形にできるだけ近づくときに、相似になる。したがって、一様な向心力のもとでほぼ円形の軌道を回る物体は、遠い方の遠近点と近い方の遠近点の間で、つねに、中心に対して角$\frac{180}{\sqrt{3}}$度、すなわち103度55分をなす。遠い方の遠近点から近い方の遠近点へ着いて、一度この角をなしたあと、そこから遠い方の遠近点に戻って、もう一度同じ角をなし、以下限りなくこれを繰り返す。

\textit{例2.}\wsp{}向心力が、高さ$A$の任意の冪$A^{n-3}$、すなわち$\frac{A^n}{A^3}$に比例するとしよう。ここで$n-3$と$n$は、整数でも分数でも、有理数でも無理数でも、正でも負でもよい任意の冪の指数を表す。その分子$A^n$、すなわち$(T - X)^n$を、私たちの収束級数の方法によって不定の級数に展開すれば、$T^n - nXT^{n-1} + \frac{n^2 - n}{2} X^2 T^{n-2} - \ldots$となる。この項を、もう一方の分子$RG^2 - RF^2 + TF^2 - F^2 X$の項と比べれば、$RG^2 - RF^2 + TF^2 : T^n = -F^2 : -nT^{n-1} + \frac{n^2 - n}{2} XT^{n-2} - \ldots$となる。そして、軌道が円形に近づくときの最終の比をとれば、$RG^2 : T^n = -F^2 : -nT^{n-1}$、すなわち$G^2 : T^{n-1} = F^2 : nT^{n-1}$、入れ替えて$G^2 : F^2 = T^{n-1} : nT^{n-1}$、すなわち$1 : n$となる。したがって$G : F$、すなわち角$VCp$と角$VCP$の比は$1 : \sqrt{n}$である。それゆえ、楕円で物体が遠い方の遠近点から近い方の遠近点へ降りるあいだになす角$VCP$は180度であるから、冪$A^{n-3}$に比例する向心力で任意の物体が描くほぼ円形の軌道において、物体が遠い方の遠近点から近い方の遠近点へ降りるあいだになす角$VCp$は、$\frac{180}{\sqrt{n}}$度に等しい。この角を繰り返して、物体は近い方の遠近点から遠い方の遠近点へ戻り、以下限りなくこれが続く。たとえば、向心力が物体の中心からの距離に比例する、すなわち$A$、あるいは$\frac{A^4}{A^3}$に比例するならば、$n$は4に等しく、$\sqrt{4}$は2に等しい。したがって、遠い方の遠近点と近い方の遠近点との間の角は$\frac{180}{2}$度、すなわち90度に等しい。したがって、一周の4分の1が終わると物体は近い方の遠近点に着き、さらに4分の1が終わると遠い方の遠近点に着き、以下かわるがわる限りなく続く。これは命題Xからも明らかである。なぜなら、この向心力に押される物体は、中心が力の中心にある不動の楕円の上を回るからである。もし向心力が距離に反比例する、すなわち$\frac{1}{A}$、あるいは$\frac{A^2}{A^3}$に比例するならば、$n = 2$であり、したがって遠い方と近い方の遠近点の間の角は$\frac{180}{\sqrt{2}}$度、すなわち127度17分である。したがってそのような力で回る物体は、この角を絶えず繰り返しながら、かわるがわる、遠い方の遠近点から近い方へ、近い方から遠い方へと、永遠に着く。さらに、もし向心力が、高さの11乗の4乗根（11/4乗）に反比例する、すなわち$A^{\frac{11}{4}}$に反比例し、したがって$\frac{1}{A^{11/4}}$、あるいは$\frac{A^{1/4}}{A^3}$に比例するならば、$n$は$\frac{1}{4}$に等しく、$\frac{180}{\sqrt{n}}$度は360度に等しい。したがって、遠い方の遠近点を離れて、そののち絶えず降りてゆく物体は、一回転を完全に終えたときに近い方の遠近点に着き、それからまた絶えず昇りながらもう一回転を終えて、遠い方の遠近点に戻る。こうして、かわるがわる永遠に続く。

\textit{例3.}\wsp{}高さの冪の任意の指数として$m$と$n$を、任意の与えられた数として$b$、$c$をとり、向心力が$\frac{bA^m + cA^n}{A^3}$、すなわち$\frac{b(T - X)^m + c(T - X)^n}{A^3}$に比例するとしよう。これを（同じ収束級数の方法によって）展開すれば、
\[
\frac{bT^m - mbXT^{m-1} + \frac{m^2 - m}{2} bX^2 T^{m-2} + cT^n - ncXT^{n-1} + \frac{n^2 - n}{2} cX^2 T^{n-2} - \ldots}{A^3}
\]
に比例する。分子の項を比べれば、$RG^2 - RF^2 + TF^2 : bT^m + cT^n = -F^2 : -mbT^{m-1} - ncT^{n-1} + \frac{m^2 - m}{2} bXT^{m-2} + \frac{n^2 - n}{2} cXT^{n-2} - \ldots$となる。軌道が円形に近づくときに出てくる最終の比をとれば、$G^2 : bT^{m-1} + cT^{n-1} = F^2 : mbT^{m-1} + ncT^{n-1}$、入れ替えて$G^2 : F^2 = bT^{m-1} + cT^{n-1} : mbT^{m-1} + ncT^{n-1}$となる。この比は、最大の高さ$CV$、すなわち$T$を算術的に1で表せば、$G^2 : F^2 = b + c : mb + nc$、したがって$1 : \frac{mb + nc}{b + c}$になる。これから$G : F$、すなわち角$VCp$と角$VCP$の比は、$1 : \sqrt{\frac{mb + nc}{b + c}}$である。したがって、不動の楕円における、遠い方の遠近点と近い方の遠近点の間の角$VCP$は180度であるから、量$\frac{bA^m + cA^n}{A^3}$に比例する向心力で物体が描く軌道における、同じ二つの遠近点の間の角$VCp$は、$180\sqrt{\frac{b + c}{mb + nc}}$度に等しい。そして同じ論法で、もし向心力が$\frac{bA^m - cA^n}{A^3}$に比例するならば、遠近点の間の角は$180\sqrt{\frac{b - c}{mb - nc}}$度であることが分かる。より難しい場合でも、同じように問題は解ける。向心力が比例する量は、つねに、分母が$A^3$である収束級数に展開しなければならない。そのうえで、この分子$RG^2 - RF^2 + TF^2 - F^2 X$の与えられた部分と与えられていない部分を、同じ比にあるとおかなければならない。そして余分な量を消し、$T$の代わりに1を書けば、$G : F$の比が得られる。

\textit{系1.}\wsp{}したがって、向心力が高さの何らかの冪に比例するならば、その冪を遠近点の運動から求めることができ、逆もまた成り立つ。すなわち、物体が同じ遠近点に戻るまでの全角運動が、一回転、すなわち360度の角運動に対して、ある数$m$が別の数$n$に対するのと同じ比にあり、高さを$A$と呼ぶならば、力は、高さの冪$A^{\frac{n^2}{m^2} - 3}$に比例し、その指数は$\frac{n^2}{m^2} - 3$である。これは例2から明らかである。ここから、その力は、高さの3乗の比より大きな比では減少しえないことが明らかになる。そのような力で回る物体が、遠近点を離れて降り始めるならば、近い方の遠近点、すなわち最小の高さには決して着かず、中心まで降りてゆき、命題XLIの系3で扱ったあの曲線を描く。しかしもし、遠近点を離れて、わずかでも昇り始めるならば、無限に昇ってゆき、遠い方の遠近点には決して着かない。実際、同じ系と命題XLIVの系6で扱われたあの曲線を描くからである。同様に、中心から遠ざかるにつれて力が高さの3乗より大きな比で減少するとき、遠近点を離れる物体は、降り始めるか昇り始めるかにしたがって、中心まで降りるか、無限に昇ってゆく。しかし、中心から遠ざかるにつれて力が、高さの3乗より小さな比で減少するか、あるいは高さの任意の比で増加するならば、物体は決して中心まで降りることはなく、いつか近い方の遠近点に着く。逆に、物体が遠近点から遠近点へかわるがわる降りたり昇ったりして、決して中心に近づかないならば、力は中心から遠ざかるにつれて、増加するか、高さの3乗より小さい比で減少する。そして、物体が遠近点から遠近点へ早く戻れば戻るほど、力の比はその3乗の比から遠く離れる。たとえば、物体が、8回、4回、2回、あるいは$1\frac{1}{2}$回の回転で、遠い方の遠近点から遠い方の遠近点へ、かわるがわる降りたり昇ったりして戻るならば、すなわち$m : n$が$8$、$4$、$2$、$1\frac{1}{2}$に対して$1$であるならば、$\frac{n^2}{m^2} - 3$は$\frac{1}{64} - 3$、$\frac{1}{16} - 3$、$\frac{1}{4} - 3$、$\frac{4}{9} - 3$に等しく、力は$A^{\frac{1}{64} - 3}$、$A^{\frac{1}{16} - 3}$、$A^{\frac{1}{4} - 3}$、$A^{\frac{4}{9} - 3}$に比例し、すなわち$A^{3 - \frac{1}{64}}$、$A^{3 - \frac{1}{16}}$、$A^{3 - \frac{1}{4}}$、$A^{3 - \frac{4}{9}}$に反比例する。物体が一回転ごとに同じ遠近点（動かない）に戻るならば、$m : n$は$1 : 1$であり、したがって$A^{\frac{n^2}{m^2} - 3}$は$A^{-2}$、すなわち$\frac{1}{A^2}$に等しく、したがって力の減少は高さの二乗の比であり、これは先に示したとおりである。もし物体が、一回転の4分の3、3分の2、3分の1、あるいは4分の1で同じ遠近点に戻るならば、$m : n$は$\frac{3}{4}$、$\frac{2}{3}$、$\frac{1}{3}$、$\frac{1}{4}$に対して$1$であり、したがって$A^{\frac{n^2}{m^2} - 3}$は$A^{\frac{16}{9} - 3}$、$A^{\frac{9}{4} - 3}$、$A^{9 - 3}$、$A^{16 - 3}$に等しく、したがって力は、$A^{\frac{11}{9}}$、$A^{\frac{3}{4}}$に反比例するか、あるいは$A^6$、$A^{13}$に比例する。最後に、物体が遠い方の遠近点から遠い方の遠近点へ進むあいだに、一回転を完全に終え、そのうえさらに3度をなすならば、したがってその遠近点が物体の一回転ごとに進行方向に3度進むならば、$m : n$は363度$: 360$度であり、したがって$A^{\frac{n^2}{m^2} - 3}$は$A^{-\frac{265707}{131769}}$に等しく、したがって向心力は$A^{\frac{265707}{131769}}$、すなわち$A^{2\frac{4}{243}}$に反比例する。したがって向心力は、二乗の比よりわずかに大きい比で減少し、その比は二乗の比に、三乗の比よりも$60\frac{3}{4}$倍近い。

\textit{系2.}\wsp{}さらに、もし物体が、高さの二乗に反比例する向心力によって、力の中心に焦点をもつ楕円を回り、この向心力に別の任意の外来の力が加えられるか除かれるならば、その外来の力から生じる遠近点の運動を（例3により）知ることができ、逆もまた成り立つ。たとえば、楕円の上で物体を回す力が$\frac{1}{A^2}$に比例し、除かれた外来の力が$cA$に比例し、したがって残りの力が$\frac{A - cA^4}{A^3}$に比例するならば、（例3において）$A$は1に等しく、$n$は4に等しく、したがって遠近点の間の回転の角は$180\sqrt{\frac{1 - c}{1 - 4c}}$度に等しい。その外来の力が、物体を楕円で回すもう一方の力より$357\decimals{45}$倍小さい、すなわち$c$が$\frac{100}{35745}$であるとすれば、$180\sqrt{\frac{1 - c}{1 - 4c}}$は$180\sqrt{\frac{35645}{35345}}$、すなわち$180\decimals{7602}$、すなわち180度45分37秒になる。したがって、遠い方の遠近点を離れる物体は、角運動180度45分37秒で近い方の遠近点に着き、この運動を2倍にして遠い方の遠近点に戻る。したがって遠い方の遠近点は、一回転ごとに進みながら、1度31分14秒をなす。

ここまでは、その平面が力の中心を通る軌道における物体の運動について述べた。物体の運動を、中心を通らない平面においても決めることが残っている。重い物体の運動を扱う著者たちは、任意の与えられた平面における、斜めの上昇と下降も、垂直な上昇と下降も、ともに考えるのがふつうである。同様の権利で、任意の力で中心に向かい、中心を通らない平面にもたれる物体の運動も、ここで考えなければならない。平面は、物体を妨げないように、きわめてなめらかで完全に摩擦のないものと仮定する。それどころか、これらの証明では、物体がのっていて、そこにのることで接する平面の代わりに、これらに平行で、そのなかで物体の中心が動き、動くことによって軌跡を描く平面を用いる。そして同じ法則で、曲面の上でなされる物体の運動もそのうちに決める。

\section{第1編 第10節：与えられた曲面上の運動および振子の運動（命題XLVI〜LVI）}


% Section: 第1編 第10節：与えられた曲面上の運動および振子の運動（命題XLVI〜LVI）



\begin{center}\textit{与えられた曲面上の物体の運動と、糸でつるした振子の往復運動について}\end{center}

\condpagelarge{命題XLVI. 問題XXXII.}

\textit{どんな種類の向心力が仮定されても、力の中心と、物体が回る任意の平面とが与えられ、曲線図形の求積が許されるならば、与えられた位置から、与えられた速さで、その平面上の与えられた直線に沿って出発した物体の運動を求めること。}

\pngright{153.png}{}{}{}

$S$を力の中心、$SC$をこの中心から与えられた平面までの最短の距離、$P$を位置$P$から直線$PZ$に沿って出発する物体、$Q$をその同じ物体が自分の軌道上を回っているときの位置、$PQR$を与えられた平面上に描かれたその軌道で、求めるべきものとする。$CQ$、$QS$を結び、$QS$上に、物体が中心$S$に向かって引かれる向心力に比例して$SV$をとり、$CQ$に平行に$VT$を引いて、$SC$と$T$で交わらせる。すると力$SV$は（運動の法則の系2により）力$ST$、$TV$に分解される。そのうち$ST$は、物体を平面に垂直な直線に沿って引くので、この平面上の物体の運動をまったく変えない。もう一方の力$TV$は、平面の方向にはたらいて、物体を平面上の与えられた点$C$に向かってまっすぐに引くので、物体を、あたかも力$ST$が取り除かれて、物体が力$TV$だけによって、自由な空間で中心$C$のまわりを回るかのように、この平面上で動かす。さて、物体$Q$が自由な空間で与えられた中心$C$のまわりを回る向心力$TV$が与えられれば、命題XLIIによって、物体が描く軌道$PQR$も、与えられた任意の時刻に物体がいる位置$Q$も、さらにその位置$Q$での物体の速さも、与えられる。逆もまた成り立つ。\QED

\condpagelarge{命題XLVII. 定理XV.}

\textit{向心力が物体の中心からの距離に比例すると仮定すれば、任意の平面上を回るすべての物体は楕円を描き、等しい時間で回転を終える。また、直線上を行ったり来たりして動く物体は、行って戻る一往復を、同じ時間で終える。}

実際、上の命題と同じ仮定のもとで、任意の平面$PQR$上を回る物体$Q$が中心$S$に向かって引かれる力$SV$は、距離$SQ$に比例する。したがって$SV : SQ = TV : CQ$であるから、物体が軌道の平面上の与えられた点$C$に向かって引かれる力$TV$は、距離$CQ$に比例する。ゆえに、平面$PQR$上にある物体が点$C$に向かって引かれる力は、距離に応じて、それぞれの物体が中心$S$に向かって引かれる力に等しい。したがって物体は、任意の平面$PQR$の上で点$C$のまわりに、自由な空間で中心$S$のまわりに動くのと、同じ時間で同じ図形のなかを動く。ゆえに（命題Xの系2と命題XXXVIIIの系2により）つねに等しい時間で、その平面上で中心$C$のまわりに楕円を描くか、その平面上で中心$C$を通って引いた直線上を行ったり来たりする周期を完了する。\QED

\pngcent{153.png}{}{}

\condpagelarge{\textit{注.}}

これらに近いのは、曲面上での物体の上昇と下降である。平面上に曲線が描かれ、それが力の中心を通る任意の与えられた軸のまわりに回転して、その回転によって曲面を描くとし、物体が、その中心がつねにこれらの曲面上にあるように動くとする。もしそれらの物体が、斜めに昇ったり降りたりしながら行ったり来たりするなら、その運動は、軸を通る平面の中で、したがってその回転によってそれらの曲面が生じた曲線の上でなされる。したがってそれらの場合には、これらの曲線上の運動を考えれば足りる。

\condpagelarge{命題XLVIII. 定理XVI.}

\textit{輪が球の外側に直角に接し、輪が回るように回転しながら大円の上を進むとき、輪の周上の任意に与えられた点が、球に触れてから描いた曲線の道のりの長さは、その間に球に触れた弧の半分の正矢の2倍に対して、球と輪の直径の和が球の半径に対するのと同じ比をなす。}

\condpagelarge{命題XLIX. 定理XVII.}

\textit{輪が凹んだ球の内側に直角に接し、回転しながら大円の上を進むとき、輪の周上の任意に与えられた点が、球に触れてから描いた曲線の道のりの長さは、その間に球に触れた弧の半分の正矢の2倍に対して、球と輪の直径の差が球の半径に対するのと同じ比をなす。}

\pngcent{156.png}{}{}

$ABL$を球、$C$をその中心、$BPV$をそれにのった輪、$E$を輪の中心、$B$を接点、$P$を輪の周上の与えられた点とする。この輪が大円$ABL$上を$A$から$B$を通って$L$の方へ進み、その間に弧$AB$、$PB$がつねに互いに等しくなるように回転し、輪の周上の与えられた点$P$が、その間に曲線の道$AP$を描くと考える。$AP$を、輪が$A$で球に触れてから描かれた曲線の道の全体とすれば、この道の長さ$AP$と、弧$\frac{1}{2}PB$の正矢の2倍との比は、$2CE$と$CB$との比に等しい。

実際、直線$CE$（必要なら延長したもの）が輪と$V$で交わるとし、$CP$、$BP$、$EP$、$VP$を結び、延長した$CP$へ垂線$VF$を下ろす。$PH$、$VH$が円に$P$、$V$で接して$H$で交わり、$PH$が$VF$と$G$で交わるとし、$VP$へ垂線$GI$、$HK$を下ろす。また中心$C$、任意の半径で円$nom$を描いて、直線$CP$と$n$で、輪の周$Bp$と$o$で、曲線の道$AP$と$m$で交わらせ、中心$V$、半径$Vo$で円を描いて、延長した$VP$と$q$で交わらせる。

輪は、進むあいだつねに接点$B$のまわりに回転しているから、直線$BP$が、輪の点$P$の描くその曲線$AP$に垂直であること、したがって直線$VP$がこの曲線に点$P$で接することは明らかである。円$nom$の半径を徐々に増やして、ついに距離$CP$に等しくするとき、消えてゆく図形$Pnomq$と図形$PFGVI$が相似であるから、消えてゆく微小な線$Pm$、$Pn$、$Po$、$Pq$の最終の比、すなわち曲線$AP$、直線$CP$、円弧$BP$の瞬間の増加分と、直線$VP$の減少分の比は、線$PV$、$PF$、$PG$、$PI$の比に等しい。ところで$VF$は$CF$に、$VH$は$CV$に垂直であり、したがって角$HVG$と$VCF$は等しい。また角$VHP$は（四辺形$HVEP$の$V$と$P$での角が直角だから）角$VEP$と合わせて二直角になり、したがって角$CEP$に等しいから、三角形$VHG$と$CEP$は相似である。これから$EP : CE = HG : HV$、すなわち$HP$、すなわち$KI : KP$となり、分割して$CB : CE = PI : PK$、後項を2倍して$CB : 2CE = PI : PV$となる。したがって線$VP$の減少分、すなわち線$BV - VP$の増加分と、曲線$AP$の増加分との比は、与えられた比$CB : 2CE$であり、ゆえに（補助定理IVの系により）それらの増加分によって生じた長さ$BV - VP$と$AP$は、同じ比にある。ところで$BV$を半径とするとき、$VP$は角$VPB$、すなわち$\frac{1}{2}BEP$の余弦であり、したがって$BV - VP$はその同じ角の正矢であり、したがって半径$\frac{1}{2}BV$のこの輪では、$BV - VP$は弧$\frac{1}{2}BP$の正矢の2倍になる。ゆえに$AP$と、弧$\frac{1}{2}BP$の正矢の2倍との比は、$2CE : CB$である。\QED

線$AP$を、区別のために、前の命題では球の外側のサイクロイド、後の命題では球の内側のサイクロイドと呼ぶことにする。

\textit{系1.}\wsp{}したがって、サイクロイドの全体$ASL$を描き、それを$S$で二等分すれば、部分$PS$の長さと$VP$の長さ（これは$EB$を半径とするとき、角$VBP$の正弦の2倍である）との比は、$2CE : CB$で、したがって与えられた比である。

\textit{系2.}\wsp{}また、サイクロイドの半周$AS$の長さは、輪の直径$BV$に対して$2CE$が$CB$に対するのと同じ比をもつ直線に等しい。

\textit{系3.}\wsp{}したがって、球の半径が与えられるならば、その長さは矩形$BEC$に比例する。

\condpagelarge{命題L. 問題XXXIII.}

\pngrightsc{159.png}{}{}{}{}

\textit{つるした物体を、与えられたサイクロイドの上で振動させること。}

中心$C$で描いた球$QVS$の内側に、点$R$で二等分された、サイクロイド$QRS$が与えられ、その両端$Q$、$S$で球の表面と両側で交わるとする。$CR$を引いて弧$QS$を$O$で二等分し、それを$A$まで延ばして、$CA : CO = CO : CR$にする。中心$C$、半径$CA$で外側の球$ABD$を描き、この球の内側に、直径$AO$の輪によって、内側の球に$Q$、$S$で接し、外側の球と$A$で交わる二つの半サイクロイド$AQ$、$AS$を描く。その点$A$から、長さ$AR$に等しい糸$APT$によって物体$T$をつるし、半サイクロイド$AQ$、$AS$の間を振動させる。すなわち、振子が垂線$AR$から離れるたびに、糸が、その上の部分$AP$で、動きが向かっている側の半サイクロイド$APS$に沿い、それを障害物としてそのまわりに曲がり、残りの部分$PT$（これにはまだ半サイクロイドが当たっていない）が直線に伸びるようにする。すると重り$T$は、与えられたサイクロイド$QRS$の上で振動する。\QED

実際、糸$PT$が、サイクロイド$QRS$と$T$で、円$QOS$と$V$で交わるとし、$CV$を引いて円$ABD$と$B$で交わらせる。糸の直線の部分$PT$に対して、両端$P$、$T$から垂線$PB$、$TW$を立てて、直線$CV$と$B$、$W$で交わらせる。サイクロイドの生成から、それらの垂線$PB$、$TW$が$CV$から切り取る長さ$VB$、$VW$が、輪の直径$OA$、$OR$に等しいこと、したがって点$B$が円$ABD$に落ちることは明らかである。ゆえに$TP : VP$（$\frac{1}{2}BV$を半径とするとき、角$VBP$の正弦の2倍）は、$BW : BV$、すなわち$AO + OR : AO$、すなわち（$CA : CO$、$CO : CR$が比例し、分割して$AO : OR$も比例するから）$CA + CO$、すなわち$2CE : CA$に等しい。したがって命題XLIXの系1により、長さ$PT$はサイクロイドの弧$PS$に等しく、糸全体$APT$はサイクロイドの半弧$APS$に等しい。すなわち（命題XLIXの系2により）長さ$AR$に等しい。したがって逆に、もし糸がつねに長さ$AR$に等しいままなら、点$T$はサイクロイド$QRS$の上を動く。\QED

\textit{系.}\wsp{}糸$AR$は、サイクロイドの半弧$APS$に等しい。

\condpagelarge{命題LI. 定理XVIII.}

\textit{球の中心$C$へ、あらゆる場所から向かう向心力が、それぞれの場所の中心からの距離に比例し、この力だけがはたらいて、物体$T$が（いま述べた仕方で）サイクロイド$QRS$の周上を振動するならば、私は言う。どんなに等しくない振れ幅の振動でも、その時間は等しい。}

実際、サイクロイドの無限に延長した接線$TW$に垂線$CX$を下ろし、$CT$を結ぶ。物体$T$が$C$に向かって押される向心力は距離$CT$に比例するから、（運動の法則の系2により）それは部分$CX$、$TX$に分けられる。そのうち$CX$は、物体を$P$からまっすぐ遠ざける向きに押して、糸$PT$を張り、その抵抗によってすべて打ち消され、ほかのいかなる効果も生まない。もう一方の部分$TX$は、物体を横向きに、すなわち$X$の方へ押して、サイクロイド上の運動をまっすぐに加速する。したがって物体の加速は、この加速する力に比例して、各瞬間に長さ$TX$に比例する。すなわち、$CV$、$WV$が与えられ、$TX$、$TW$がそれらに比例するから、長さ$TW$に比例し、すなわち（命題XLIXの系1により）サイクロイドの弧$TR$の長さに比例する。そこで、二つの振子$APT$、$Apt$を垂線$AR$から異なる角度に引いて同時に放せば、それらの加速は、つねに、描かれるべき弧$TR$、$tR$に比例する。ところで、初めに描かれた部分は加速に比例し、すなわち初めに描かれるべき全体に比例するから、描かれずに残っている部分と、それに続く加速も、これらの部分に比例するので、やはり全体に比例する。以下同様である。したがって加速、したがって生じた速さ、これらの速さによって描かれた部分と描かれるべき部分は、つねに全体に比例する。ゆえに、描かれるべき部分は、互いに与えられた比を保ちながら同時に消え、すなわち、振動する二つの物体は同時に垂線$AR$に達する。そして逆に、最低の位置$R$から垂線が、同じトロコイドの弧に沿って逆向きに昇るとき、下降の運動を加速したのと同じ力によって、各位置で減速されるから、同じ弧を昇る運動と降る運動の速さは等しく、したがって等しい時間になされることが明らかである。したがって、垂線の両側にあるサイクロイドの二つの部分$RS$と$RQ$は相似で等しいから、二つの振子は、それらの振動を、全体でも半分でも、つねに同じ時間で行う。\QED

\condpagelarge{命題LII. 問題XXXIV.}

\textit{振子の、それぞれの位置での速さと、振動の全体と、振動の個々の部分とが、それぞれ費やす時間とを定めること。}

\pngright{161.png}{}{}{}

任意の中心$G$、サイクロイドの弧$RS$に等しい半径$GH$で半円$HKMG$を描き、半径$GK$で二等分する。そして、中心$G$に向かって、場所の中心からの距離に比例する向心力がはたらき、周$HIK$上でのその力が、球$QOS$の周上での、その中心に向かう向心力に等しい（命題L・LIの図を見よ）とする。そして振子$T$が最も高い位置$S$から放たれるのと同時に、ある物体$L$が$H$から$G$へ落下するとする。物体が押される力は、初めには等しく、描かれるべき距離$TR$、$GL$にはつねに比例し、したがって$TR$と$LG$が等しいとき位置$T$、$L$でも等しいから、その二つの物体が、初めには等しい距離$ST$、$HL$を描き、したがってそののちも等しく押されつづけ、等しい距離を描くことは明らかである。ゆえに命題XXXVIIIにより、物体が弧$ST$を描く時間と、一回の振動の時間との比は、弧$HI$（物体$H$が$L$に達する時間）と半円$HKM$（物体$H$が$M$に達する時間）との比に等しい。そして、振子の物体の位置$T$での速さと、最低の位置$R$での速さとの比は、（すなわち物体$H$の位置$L$での速さと、位置$G$での速さとの比、すなわち弧$HI$、$HK$が一様な流れで増加するときの線$HL$の瞬間の増加分と線$HG$の瞬間の増加分との比は）縦線$LI$と半径$GK$との比、すなわち$\sqrt{SR^2 - TR^2}$と$SR$との比に等しい。したがって、等しくない振動でも、等しい時間に、振動の全弧に比例する弧が描かれるから、与えられた時間から、すべての振動における速さと描かれた弧が得られる。これが初めに求めるべきことであった。

\pngright{162.png}{}{}{}

次に、糸でつるした二つの物体が、等しくないサイクロイドの上で振動するとし、それらの半弧に等しい直線$GH$、$gh$をとり、中心$G$、$g$、半径$GH$、$gh$で半円$HZKM$、$hzkm$を描く。それらの直径$HM$、$hm$の上に、等しい微小な線$HY$、$hy$をとり、垂線$YZ$、$yz$を立てて、円周と$Z$、$z$で交わらせる。振子の物体は運動の初めには球$QOS$の周上にあり、したがって中心に向かって等しい力で押され、まっすぐに中心へ向かって動き始めるから、同時に進む距離は初めには等しい。そこで、物体$H$、$h$が、$H$と$h$で同じ力に押されるとし、$HY$、$hy$が運動のまさに初めに描かれる等しい距離であるとすれば、弧$HZ$、$hz$が等しい時間を表す。これらの生じつつある弧の最初の比の二乗は、矩形$GHY$、$ghy$の比、すなわち線$GH$、$gh$の比に等しく、したがって半径の比の平方根の比にとった弧が、等しい時間を表す。ゆえに、円$HKM$の全体の時間、すなわち一つのサイクロイドでの一回の振動に対応する時間と、円$hkm$の全体の時間、すなわち別のサイクロイドでの一回の振動に対応する時間との比は、半周$HKM$と、この半周ともう一つの円$hkm$の半周との比例中項との比、すなわち直径$HM$と直径$hm$との比の平方根の比、すなわち先のサイクロイドの周と別のサイクロイドの周との比の平方根の比に等しい。したがってその時間は、任意のサイクロイドでは（命題XLIXの系3により）、サイクロイドを描いた輪の半径と、その半径と球の半径との差とがつくる矩形$BEC$の平方根に比例する。\QED　また同じ時間は（命題Lの系により）糸の長さ$AR$の平方根に比例する。\QED

さらに、同心の球に相似なサイクロイドを描くとする。それらの周は球の半径に比例し、周の対応する位置での力は、場所の球の共通の中心からの距離に比例し、すなわち球の半径に比例し、したがってサイクロイドの周、および周の相似な部分に比例するから、周の相似な部分が相似な振動で描かれる時間は等しく、したがってすべての振動は等時的（isochron）である。ゆえに、与えられた球での振動の時間は、長さ$AR$の平方根に比例し、したがって（$AC$が与えられているから）数$\frac{AR}{AC}$の平方根に比例し、すなわち数$\sqrt{\frac{AR}{AC}}$に（そのままの比で）比例する。そしてこの数$\sqrt{\frac{AR}{AC}}$は、比$AR : AC$が保たれる限り（相似なサイクロイドでそうであるように）つねに同じままで、したがって異なる球でもサイクロイドが相似であるときは時間に比例する。したがって、他の任意の与えられた球での、したがってすべての同心の球での振動の時間は、数$\sqrt{\frac{AR}{AC}}$に比例し、すなわち、糸の長さ$AR$の平方根に正比例し、球の半径$AC$の平方根に反比例する比を合成した比にあることは明らかである。\QED

最後に、異なる球の絶対的な力が等しくないとすれば、等しい時間になされる加速は力に比例する。したがって、時間を、力の平方根に反比例してとるならば、速さは同じ平方根に正比例し、したがってこれらの時間に描かれる距離は等しい。ゆえに、任意の絶対的な力をもつ、すべての球とサイクロイドでの振動は、振子の長さの平方根に正比例し、振子の中心と球の中心との距離の平方根に反比例し、絶対的な力の平方根にも反比例する比を合成した比にあり、すなわちその力を$V$と呼ぶなら、数$\sqrt{\frac{AR}{AC \times V}}$の比にある。\QED

\textit{系1.}\wsp{}したがって、振動する物体、落下する物体、回る物体の時間も、互いに比べることができる。実際、サイクロイドが球の内側に描かれるときの輪の直径を、球の半径に等しくすれば、サイクロイドは球の中心を通る直線になり、振動は、この直線上の下降と、それに続く上昇になる。これから、任意の位置から中心へ下る時間と、これに等しい時間、すなわち物体が球の中心のまわりを任意の距離で回りながら四分円の弧を描く時間が得られる。実際、この時間は（第二の場合により）、任意のトロコイド$APS$での半振動の時間に対して、$\frac{1}{2}BC : \sqrt{BEC}$の比をなす。

\textit{系2.}\wsp{}これによって、レン氏（クリストファー・レン）とホイヘンス氏が通常のサイクロイドについて見いだしたことも導かれる。実際、球の直径を無限に大きくすれば、その球面は平面に変わり、向心力は、この平面に垂直な線に沿って一様にはたらき、私たちのサイクロイドは通常のサイクロイドになる。この場合、その平面と、描く点との間のサイクロイドの弧の長さは、その平面と描く点との間の輪の弧の半分の正矢の4倍に等しくなる。これはレン氏が見いだしたとおりである。そして二つのそのようなサイクロイドの間の振子は、相似で等しいサイクロイドの上で等しい時間に振動する。これはホイヘンスが示したとおりである。さらに、重い物体の、一回の振動の時間に落ちる距離は、ホイヘンスが示したものになる。

私たちが示した命題は、地球の実際のしくみにあてはまる。すなわち、輪が地球の大円の上を進むとき、釘の運動によって球の外にサイクロイドを描く。また鉱山や洞窟の下にぶらさげた振子は、すべての振動が等時的になるように、球の内側のサイクロイドの上で振動しなければならない。なぜなら、重力は（第三巻で教えるように）、地球の表面から上へ進むにつれて中心からの距離の二乗の比で、下へ進むにつれて単純な比で、減少するからである。

\condpagelarge{命題LIII. 問題XXXV.}

\textit{曲線図形の求積が許されるならば、与えられた曲線上の物体が、つねに等時的な振動をするような力を求めること。}

\pngright{165.png}{}{}{}

物体$T$が、任意の曲線$STRQ$の上を振動するとし、その軸$OR$は力の中心$C$を通るとする。その曲線に物体の任意の位置$T$で接する直線$TX$を引き、この接線$TX$の上に、弧$TR$に等しい$TY$をとる。その弧の長さは、図形の求積によって通常の方法で知られるからである。点$Y$から、接線に垂直な直線$YZ$を立てる。$CT$を引いてその垂線と$Z$で交わらせれば、向心力は直線$TZ$に比例する。\QED

実際、物体が$T$から$C$に向かって引かれる力を、それに比例してとった直線$TZ$で表せば、それは力$TY$、$YZ$に分解される。そのうち$YZ$は、物体を糸$PT$の方向に引くので、その運動を少しも変えず、もう一方の力$TY$は、曲線$STRQ$上の物体の運動をまっすぐに加速するか、まっすぐに減速する。したがって、これが描かれるべき道$TR$に比例するとき、物体の加速や減速は、二つの振動（大きいものと小さいもの）の、描かれるべき比例した部分において、つねにそれらの部分に比例し、したがってそれらの部分を同時に描かせる。ところで、全体にいつも比例する部分を同時に描く物体は、全体も同時に描く。\QED

\textit{系1.}\wsp{}したがって、もし物体$T$が、中心$A$から直線の糸$AT$でつるされて、円弧$STRQ$を描き、その間に平行な線に沿って下向きに、一様な重力に対して弧$TR$とその正弦$TN$との比にある何らかの力で押されるならば、個々の振動の時間は等しい。実際、$TZ$と$AR$が平行であるから、三角形$ANT$と$TYZ$は相似であり、したがって$TZ : AT = TY : TN$である。すなわち、一様な重力の力を与えられた長さ$AT$で表せば、振動を等時的にする力$TZ$は、重力$AT$に対して、$TY$に等しい弧$TR$と、その弧の正弦$TN$との比をなす。

\textit{系2.}\wsp{}したがって時計において、運動を保つために機械が振子に加える力が、重力と合成されて、全体の下向きの力がつねに、弧$TR$と半径$AR$の矩形を正弦$TN$で割ってできる線に比例するようにできるならば、すべての振動は等時的になる。

\condpagelarge{命題LIV. 問題XXXVI.}

\textit{曲線図形の求積が許されるならば、物体が、任意の向心力のもとで、力の中心を通る平面上に描かれた任意の曲線の上を、降りたり昇ったりする時間を求めること。}

\pngright{167.png}{}{}{}

物体が、任意の位置$S$から、力の中心$C$を通る平面上に与えられた任意の曲線$STtR$に沿って降りるとする。$CS$を結び、それを数えきれないほど多くの等しい部分に分け、$Dd$をそれらの部分の一つとする。中心$C$、半径$CD$、$Cd$で円$DT$、$dt$を描き、曲線$STtR$と$T$、$t$で交わらせる。向心力の法則と、物体が落ちてきた高さ$CS$が与えられるから、他の任意の高さ$CT$での物体の速さが、命題XXXIXによって与えられる。ところで物体が微小な線$Tt$を描く時間は、この微小な線の長さ（すなわち角$tTC$の割線）に正比例し、速さに反比例する。この時間に比例して、点$D$を通って直線$CS$に垂直な縦線$DN$をとれば、$Dd$が与えられているから、矩形$Dd \times DN$、すなわち面積$DNnd$が、同じ時間に比例する。したがって$SNn$を、点$N$が絶えず接するその曲線とすれば、面積$SNDS$は、物体が降りて線$ST$を描いた時間に比例する。ゆえにそのようにして求めた面積から、時間が得られる。\QED

\condpagelarge{命題LV. 定理XIX.}

\textit{もし物体が、軸が力の中心を通る任意の曲面上を動き、物体から軸へ垂線が下ろされ、それに平行で等しい線が軸の任意の点から引かれるならば、私は言う。その平行線は、時間に比例する面積を描く。}

\pngright{168.png}{}{}{}

$BSKL$を曲面、$T$をその上を回る物体、$STtR$を物体がその上に描く軌道、$S$を軌道の始点、$OMNK$を曲面の軸、$TN$を物体から軸への垂線、$OP$を、これに平行で等しく、軸上の与えられた点$O$から引いた線、$AP$を、回る線$OP$の平面$AOP$上に点$P$が描く軌道の跡、$A$を、点$S$に対応する跡の始点、$TC$を物体から中心へ引いた線、$TG$をその一部で、物体が中心$C$に向かって押される向心力に比例するもの、$TM$を曲面に垂直な線、$TI$をその一部で、物体が曲面を押す圧力の力に、また逆に物体が曲面から$M$の方へ押される力に比例するもの、$PHTF$を物体を通る軸に平行な線、$GF$、$IH$を点$G$、$I$からその平行線$PHTF$へ垂直に下ろした線とする。

さて私は言う。動き始めてから動径$OP$が描いた面積$AOP$は、時間に比例する。実際、力$TG$は（運動の法則の系2により）力$TF$、$FG$に分解され、力$TI$は力$TH$、$HI$に分解される。ところで力$TF$、$TH$は、平面$AOP$に垂直な線$PF$の方向にはたらくので、この平面に垂直な方向の物体の運動を変えるだけである。したがって、平面の方向になされる運動、すなわち、この平面上に軌道の跡$AP$を描く点$P$の運動は、あたかも力$TF$、$TH$が取り除かれて、物体が力$FG$、$HI$だけによって動かされるかのような運動と同じである。すなわち、物体が平面$AOP$上で、中心$O$に向かい、力$FG$と$HI$の和に等しい向心力によって曲線$AP$を描く場合と同じである。ところがそのような力では、面積$AOP$は（命題Iにより）時間に比例して描かれる。\QED

\textit{系.}\wsp{}同じ論法で、もし物体が、同じ与えられた任意の直線$CO$上の二つ以上の中心に向かう力に動かされて、自由な空間で任意の曲線$ST$を描くならば、面積$AOP$はつねに時間に比例する。

\condpagelarge{命題LVI. 問題XXXVII.}

\textit{曲線図形の求積が許され、与えられた中心に向かう向心力の法則と、軸がその中心を通る曲面とが与えられるならば、その曲面上で、与えられた位置から、与えられた速さで、その曲面上の与えられた方向へ出発した物体が描く軌道を求めること。}

前の命題で作った構成のまま、物体が位置$S$から、求めるべき軌道$STtR$へ出発するとし、高さ$SC$での与えられた速さから、他の任意の高さ$TC$での速さが与えられる。物体はその速さで、与えられたごく短い時間に、その軌道の微小な部分$Tt$を描くとし、$Pp$を、平面$AOP$上に描かれたその跡とする。$Op$を結び、中心$T$、半径$Tt$で曲面上に描いた小円の跡を、同じ平面$OAPp$上に描かれた楕円$PpQ$とする。小円が大きさと位置で与えられているから、その楕円$PpQ$が与えられる。そして面積$POp$は時間に比例し、したがって与えられた時間から与えられるから、$Op$が位置で与えられ、そこから、それと楕円との共通の交点$p$が、軌道の跡$APp$が直線$OP$を切る角$OPp$とともに与えられる。そしてそこから、命題XLIで曲線$VIKk$が同様のデータから求められたのと同じ方法で、軌道の跡$APp$が求められる。次に、跡の各点$P$から平面$AOP$へ垂線$PT$を立て、曲面と$T$で交わらせれば、軌道の各点$T$が得られる。\QED

\chapter{第1編 物体の運動（第11節〜第14節：多体引力・球殻引力定理・光学屈折）}



\section{第1編 第11節：相互に向心力で引き合う球体の運動（二体・多体問題、命題LVII〜LXIX）}


% Section: 第1編 第11節：相互に向心力で引き合う球体の運動（二体・多体問題、命題LVII〜LXIX）



\sectpage{XI.}

\begin{center}\textbf{相互に向心力で引き合う球体の運動}\end{center}

ここまで私は、動かない中心に引かれる物体の運動を述べてきたが、そのような中心は、自然界にはほとんど存在しない。実際、引力はふつう物体に向かって生じ、引く物体と引かれる物体の作用は、法則IIIにより、つねに相互的で等しい。したがって、物体が二つある場合には、引くほうも引かれるほうも静止することはできず、両方とも（法則の系4により）、あたかも相互の引力によるかのように、共通の重心のまわりを回る。また、物体がもっと多くあるとき（それらが一つの物体に引かれるにせよ、すべて互いに引き合うにせよ）、それらは、共通の重心が静止するか、あるいは直線上を一様に動くように、互いに動かなければならない。こうした理由から、私はこれから、互いに引き合う物体の運動を述べることにし、向心力を引力として考える。物理的に言うならば、おそらく、衝撃と呼ぶほうがより正しいであろうが。というのは、いま私たちは数学のなかにいるのだから、物理的な議論は脇において、数学の読者がより容易に理解できるように、なじみのある言葉を用いるのである。

\condpagelarge{命題 LVII. 定理 XX.}

\textit{互いに引き合う二つの物体は、共通の重心のまわりにも、互いのまわりにも、相似な図形を描く。}

実際、共通の重心からの距離は、物体に反比例し、したがって互いに与えられた比をなし、合成すれば、物体間の全距離に対しても与えられた比をなす。ところで、これらの距離は、その端点のまわりを共通の角運動で運ばれる。それは、つねに一直線上にあって、互いの傾きを変えないからである。そして、互いに与えられた比をなし、等しい角運動で端点のまわりを運ばれる直線は、それらの端点のまわりに（これらの端点とともに静止するか、あるいは角運動でない何らかの運動で動く平面の中で）、まったく相似な図形を描く。したがって、これらの距離が回されることによって描かれる図形は相似である。\QED

\condpagelarge{命題 LVIII. 定理 XXI.}

\textit{二つの物体が、任意の力で互いに引き合いながら、同時に共通の重心のまわりを回るとする。このとき、物体がそのように動いて互いのまわりに描く図形に対して、同じ力のもとで、どちらか一方の物体を静止させて、そのまわりに、相似で等しい図形を描くことができる、と私は言う。}

物体 $S$, $P$ が共通の重心 $C$ のまわりを回り、$S$ から $T$ へ、$P$ から $Q$ へ進むとする。与えられた点 $s$ から、$SP$, $TQ$ に等しくて平行な線 $sp$, $sq$ をつねに引く。すると、点 $p$ が動かない点 $s$ のまわりを回って描く曲線 $pqv$ は、物体 $S$, $P$ が互いのまわりに描く曲線に相似で等しい。したがって（定理XXにより）、同じ物体が共通の重心 $C$ のまわりに描く曲線 $ST$ と $PQV$ に相似である。それは、線 $SC$, $CP$ と $SP$（または $sp$）の互いの比が与えられているからである。

\pngright{172.png}{}{}{}

\textit{場合1.}\wsp{}その共通の重心 $C$ は、法則の系4により、静止するか、直線上を一様に動く。まず静止していると仮定し、$s$ と $p$ に二つの物体、$s$ に動かないもの、$p$ に動くものを置く。それらは物体 $S$, $P$ に相似で等しいとする。ついで、直線 $PR$ と $pr$ が曲線 $PQ$ と $pq$ に $P$ と $p$ で接するとし、$CQ$ と $sq$ を $R$ と $r$ まで延長する。図形 $CPRQ$ と $sprq$ が相似だから、$RQ$ 対 $rq$ は $CP$ 対 $sp$ に等しく、したがって与えられた比である。そこで、物体 $P$ が物体 $S$ に、したがって中間の中心 $C$ に引かれる力が、物体 $p$ が中心 $s$ に引かれる力に対して、その与えられた比をなすとすれば、これらの力は、等しい時間に、物体を接線 $PR$, $pr$ から弧 $PQ$, $pq$ へ、それらに比例する間隔 $RQ$, $rq$ だけ、つねに引くであろう。したがって、後の力は、物体 $p$ を、曲線 $pqv$ の上で回らせるであろう。この曲線は、前の力が物体 $P$ を回らせる曲線 $PQV$ に相似で、回転は同じ時間で完了するであろう。ところが、これらの力は互いに $CP$ 対 $sp$ の比をなさず、（物体 $S$ と $s$、$P$ と $p$ が相似で等しく、距離 $SP$, $sp$ が等しいから）互いに等しい。したがって、物体は等しい時間に接線から等しく引かれる。それゆえ、後の物体 $p$ がより大きな間隔 $rq$ だけ引かれるには、より長い時間が必要で、それは間隔の平方根の比になる。それは、補助定理Xにより、運動のはじめに描かれる距離が時間の二乗に比例するからである。そこで、物体 $p$ の速さを物体 $P$ の速さに対して、距離 $sp$ 対距離 $CP$ の比の平方根にとる。そうすれば、同じ平方根の比をなす時間に、全体の比をなす弧 $PQ$, $pq$ が描かれる。そして物体 $P$, $p$ は、つねに等しい力で引かれて、静止した中心 $C$ と $s$ のまわりに、相似な図形 $PQV$, $pqv$ を描く。その後者 $pqv$ は、物体 $P$ が動く物体 $S$ のまわりに描く図形に、相似で等しい。\QED

\textit{場合2.}\wsp{}いま、共通の重心が、物体が互いに動く空間とともに、直線上を一様に進むと仮定する。すると、法則の系6により、この空間内のすべての運動は前と同じように行われ、したがって物体は、互いのまわりに前と同じ図形を描く。ゆえに図形 $pqv$ は相似で等しい。\QED

\textit{系1.}\wsp{}したがって、自分たちの距離に比例する力で互いに引き合う二つの物体は、（命題Xにより）共通の重心のまわりにも、互いのまわりにも、同心の楕円を描く。逆に、そのような図形が描かれるなら、力は距離に比例する。

\textit{系2.}\wsp{}また、自分たちの距離の平方に反比例する力で引き合う二つの物体は、（命題XI, XII, XIIIにより）共通の重心のまわりにも、互いのまわりにも、図形が描かれる中心に焦点をもつ円錐曲線を描く。逆に、そのような図形が描かれるなら、向心力は距離の平方に反比例する。

\textit{系3.}\wsp{}共通の重心のまわりを回る任意の二つの物体は、その中心へも互いへも引いた半径で、時間に比例する面積を描く。

\condpagelarge{命題 LIX. 定理 XXII.}

\textit{共通の重心 $C$ のまわりを回る二つの物体 $S$ と $P$ の周期は、一方の物体 $P$ が、他方を静止させたまま $S$ のまわりを回り、物体が互いのまわりに描く図形に相似で等しい図形を描く場合の周期に対して、他方の物体 $S$ の、和 $S + P$ に対する比の平方根の比をなす。}

実際、上の命題の証明から、任意の相似な弧 $PQ$ と $pq$ が描かれる時間は、距離 $CP$ と $SP$（または $sp$）の比の平方根の比、すなわち物体 $S$ の、和 $S + P$ に対する比の平方根の比をなす。そして合成すれば、すべての相似な弧 $PQ$, $pq$ が描かれる時間の和、すなわち全体の相似な図形が描かれる全時間は、同じ平方根の比をなす。\QED

\condpagelarge{命題 LX. 定理 XXIII.}

\textit{二つの物体 $S$ と $P$ が、自分たちの距離の平方に反比例する力で互いに引き合いながら、共通の重心のまわりを回っているとする。このとき、一方の物体 $P$ がこの運動によってもう一方の $S$ のまわりに描く楕円の主軸は、同じ物体 $P$ が、もう一方の $S$ を静止させたまま同じ周期で描くことのできる楕円の主軸に対して、二つの物体の和 $S + P$ が、この和ともう一方の物体 $S$ との間にある二つの比例中項のうちの第一のものに対するのと同じ比をなす。}

もし描かれた二つの楕円が互いに等しいならば、周期は、先の定理により、物体 $S$ と和 $S + P$ との比の平方根の比になるはずである。そこで、後の楕円（$S$ が静止している場合）の周期をこの比で小さくすれば、二つの周期は等しくなる。このとき楕円の主軸は、第7の定理により、この比が3/2乗になるような比で小さくなる。すなわち、$S$ 対 $S + P$ の比がその3乗になるような比で小さくなる。したがって、もう一方の楕円の主軸に対しては、$S + P$ と $S$ との間にある二つの比例中項の第一のもの対 $S + P$ の比になる。逆にいえば、動いている物体のまわりに描かれる楕円の主軸は、静止している物体のまわりに描かれる楕円の主軸に対して、$S + P$ 対、$S + P$ と $S$ との間にある二つの比例中項の第一のもの、という比をなす。\QED

\condpagelarge{命題 LXI. 定理 XXIV.}

\textit{二つの物体が、任意の力で互いに引き合い、ほかには動かされも妨げられもしないとき、どのように動くにしても、それらの運動は、互いに引き合うのではなく、それぞれが、共通の重心に置かれた第三の物体から同じ力で引かれるとした場合と同じになる。そして、引く力の法則は、物体の共通の中心からの距離についても、物体間の全距離についても、同じである。}

実際、物体が互いに引き合う力は、物体に向かうことによって、中間にある共通の重心に向かい、したがって、中間の物体から生じるのと同じである。\QED

そして、どちらの物体の共通の中心からの距離の、同じ物体の他方の物体からの距離に対する比は与えられているから、一方の距離の任意のべきの、他方の距離の同じべきに対する比が与えられる。また、一つの距離と与えられた量とから何らかの仕方で導かれる量の、もう一つの距離と、それと同数の与えられた量（距離の与えられた比を前の量に対してもつもの）から同様に導かれる別の量に対する比も、与えられる。したがって、もし一方の物体が他方に引かれる力が、物体の互いの距離に正比例または反比例し、あるいはその距離の任意のべきに比例し、あるいはまたこの距離と与えられた量から何らかの仕方で導かれる任意の量に比例するならば、同じ物体が共通の重心に引かれる力も、同様に、引かれる物体の共通の中心からの距離に、あるいはその距離の同じべきに、あるいはこの距離と類似の与えられた量から同様に導かれる量に、正比例または反比例する。すなわち、引く力の法則は、どちらの距離についても同じである。\QED

\condpagelarge{命題 LXII. 問題 XXXVIII.}

\textit{自分たちの距離の平方に反比例する力で互いに引き合い、与えられた場所から落とされる二つの物体の運動を定めること。}

物体は、直前の定理により、あたかも共通の重心に置かれた第三の物体に引かれるかのように動く。そしてその中心は、運動のはじめには（仮定により）静止しているから、（法則の系4により）つねに静止する。したがって、物体の運動は、（問題XXVにより）あたかもその中心に向かう力で押されるかのように定めなければならず、そうすれば、互いに引き合う物体の運動が得られる。\QED

\condpagelarge{命題 LXIII. 問題 XXXIX.}

\textit{自分たちの距離の平方に反比例する力で互いに引き合い、与えられた場所から、与えられた直線に沿って、与えられた速さで出発する二つの物体の運動を定めること。}

物体のはじめの運動が与えられているから、共通の重心の一様な運動が与えられ、またこの中心とともに直線上を一様に動く空間の運動、さらにこの空間に対する物体のはじめの運動も与えられる。ところで、そのあとの運動は、（法則の系5と直前の定理により）この空間の中で、あたかも空間そのものがその共通の重心とともに静止していて、物体が互いに引き合わず、その中心にある第三の物体に引かれるかのように行われる。したがって、この動く空間の中で、与えられた場所から、与えられた直線に沿って、与えられた速さで出発し、その中心に向かう向心力にとらえられる一方の物体の運動は、問題IXと問題XXVIによって定めなければならず、同時にもう一方の物体の、向かい側の運動が得られる。この運動に、上で見いだした、空間と、その中で回る物体の系の一様な前進運動を合成すれば、動かない空間での物体の絶対運動が得られる。\QED

\condpagelarge{命題 LXIV. 問題 XL.}

\textit{物体が互いに引き合う力が、中心からの距離に単純な比で増すとき、多くの物体の互いの運動を求めること。}

まず、共通の重心 $D$ をもつ二つの物体 $T$ と $L$ を仮定する。これらは、定理XXIの系1により、$D$ に中心をもつ楕円を描き、その大きさは問題Vから知られる。

\pngright{177.png}{}{}{}

いま、第三の物体 $S$ が、先の二つの物体 $T$ と $L$ を加速力 $ST$, $SL$ で引き、また逆にそれらから引かれるとする。力 $ST$ は、法則の系2により、力 $SD$, $DT$ に分解され、力 $SL$ は力 $SD$, $DL$ に分解される。ところで、力 $DT$, $DL$ は、それらの和 $TL$ に比例し、したがって物体 $T$ と $L$ が互いに引き合う加速力に比例する。これらを物体 $T$ と $L$ の力に、前者は前者に、後者は後者に加えると、前と同じように距離 $DT$ と $DL$ に比例する力を合成するが、前の力より大きい。したがって（命題Xの系1と、命題IVの系1, 7により）これらの物体は、前と同じように楕円を描くが、運動はより速くなる。残りの加速力 $SD$ と $SD$ は、物体に比例する起動力 $SD \times T$ と $SD \times L$ によって、これらの物体を、$DS$ に平行な線 $TI$, $LK$ に沿って等しく引くから、物体の互いの位置を変えず、ただ、それらを線 $IK$ に等しく近づける。この線は、物体 $S$ の中点を通り、線 $DS$ に垂直に引かれたものと考えよ。ところで、この線 $IK$ への近づきは、物体 $T$ と $L$ の系を一方の側に、物体 $S$ を他方の側に置いて、適当な速さで、共通の重心 $C$ のまわりを回らせることによって妨げられる。このような運動で、物体 $S$ は（距離 $CS$ に比例する起動力 $SD \times T$ と $SD \times L$ の和によって中心 $C$ のほうへ引かれるから）同じ $C$ のまわりに楕円を描き、点 $D$ は、$CS$ と $CD$ が比例するから、向かい側に、これと相似な楕円を描く。ところで、物体 $T$ と $L$ は、起動力 $SD \times T$ と $SD \times L$（前者は前者に、後者は後者に）によって、（述べたように）平行線 $TI$ と $LK$ に沿って等しく引かれながら、（法則の系5と6により）動く中心 $D$ のまわりに、前と同じように自分たちの楕円を描きつづける。\QED

いま第四の物体 $V$ を加えれば、同様の論法により、この物体と点 $C$ は、すべての物体の共通の重心 $B$ のまわりに楕円を描くと結論される。そのとき、先の物体 $T$, $L$, $S$ の中心 $D$, $C$ のまわりの運動は保たれるが、いくらか速くなる。そして同じ方法で、さらに多くの物体を加えることができる。\QED

これらは、物体 $T$ と $L$ が互いに引き合う加速力が、ほかの物体が引く力より、距離の比に応じて大きいか小さいかの場合である。いま、すべての物体の相互の加速的な引力が、互いに、距離に、引く物体を掛けたものに比例するとすれば、前のことから容易に導かれるように、すべての物体は、等しい周期で、さまざまな楕円を、すべての物体の共通の重心 $B$ のまわりに、動かない平面の中で描く。\QED

\condpagelarge{命題 LXV. 定理 XXV.}

\textit{力が、それらの中心からの距離の平方に反比例して減少する多くの物体は、互いに楕円の上を動き、焦点へ引いた半径が、ほぼ時間に比例する面積を描く、ということが可能である。}

前の命題では、多くの運動が正確に楕円の上で行われる場合を証明した。力の法則が、そこで仮定した法則から離れるほど、物体は互いの運動をいっそう乱し、ここで仮定した法則にしたがって互いに引き合う物体が、互いの距離のある比を保たない限り、正確に楕円の上を動くことはできない。しかし、つぎの場合には、楕円からあまり外れない。

\textit{場合1.}\wsp{}いくつかの小さな物体が、ある最大の物体のまわりを、それからさまざまな距離で回り、それぞれに向かう絶対力が、その物体に比例するとする。そして、すべての物体の共通の重心は、（法則の系4により）静止するか、直線上を一様に動くから、小さな物体は、最大の物体がこの重心から感じられるほどには決して離れないほど小さいとする。すると、その最大の物体は、感じられる誤差なしに、静止するか、直線上を一様に動き、小さな物体は、この最大の物体のまわりを楕円の上で回り、それへ引いた半径で、時間に比例する面積を描く。ただし、最大の物体のその共通の重心からの誤差や、小さな物体どうしの作用によって、誤差が生じる分を除く。ところで、小さな物体を、この誤差と相互の作用が、与えられた任意の量より小さくなるまで小さくしていくことができ、したがって、軌道が楕円に合い、面積が時間に対応するまで、与えられた任意の量より小さくない誤差なしに、小さくできる。\QED

\textit{場合2.}\wsp{}いま、いま述べた仕方で最大の物体のまわりを回る小さな物体の系、あるいは互いのまわりを回る二つの物体のほかの任意の系が、直線上を一様に進み、その間に、はるか遠くにある別の、きわめて大きな物体の力で、側方に押されるとする。そして、平行線に沿って物体を押す等しい加速力は、物体の互いの位置を変えず、系全体を、部分の互いの運動を保ったまま、一緒に運ぶから、最大の物体への引力から、引かれる物体の互いの運動にはまったく変化は生じないことは明らかである。ただし、加速的な引力の不等や、引力が働く線の互いの傾きから生じる分を除く。そこで、最大の物体への加速的な引力がすべて、互いに距離の平方に反比例するとし、最大の物体の距離を、それからほかへ引いた直線の差と、互いの傾きが、与えられた任意のものより小さくなるまで増やすならば、系の部分の互いの運動は、与えられた任意の量より小さくない誤差なしに、保たれる。そして、部分の互いの距離がきわめて小さいため、系全体は、あたかも一つの物体のように引かれるから、この引力によって、一つの物体のように動く。すなわち、その重心で、最大の物体のまわりに、ある円錐曲線（引力が弱ければ双曲線か放物線、強ければ楕円）を描き、最大の物体へ引いた半径で、時間に比例する面積を掃く。ただし、部分の距離（まことにきわめて小さく、思いのままに小さくできる）がもたらしうる誤差を除く。\QED

同様の論法で、さらに複雑な場合へ、限りなく進むことができる。

\textit{系1.}\wsp{}場合2で、最大の物体が二つまたは多くの物体の系に近づくほど、系の部分の互いの運動はいっそう乱される。それは、最大の物体からこれらへ引いた直線の互いの傾きがより大きくなり、比の不等もより大きくなるからである。

\textit{系2.}\wsp{}ところが、系の部分の最大の物体への加速的な引力が、その最大の物体からの距離の平方に、互いに反比例しないと仮定すれば、最も大きく乱される。とくに、この比の不等が、最大の物体からの距離の比の不等より大きいならば、そうである。というのは、加速力が、平行線に沿って等しく働けば、互いの運動を乱さないのだから、乱れは作用の不等から生じなければならず、不等が大きいか小さいかに応じて、大きくも小さくもなるからである。より大きな衝撃の超過が、いくつかの物体に働き、ほかの物体には働かないと、必然的にそれらの互いの位置を変える。そしてこの乱れは、線の傾きと不等から生じる乱れに加わって、全体の乱れをより大きくする。

\textit{系3.}\wsp{}したがって、もしこの系の部分が、著しい乱れなしに、楕円や円の上を動くならば、それらは、ほかの物体に向かう加速力に押されないか、押されてもごくわずかであるか、あるいはほぼ等しく平行線に沿って押されていることは明らかである。

\condpagelarge{命題 LXVI. 定理 XXVI.}

\textit{三つの物体が、力が距離の平方に反比例して減るように互いに引き合い、どの二つの物体の第三の物体への加速的な引力も、互いに距離の平方に反比例するとする。そして小さいほうの二つが、共通の平面で最大の物体のまわりを回るとする。このとき、最内の最大の物体のまわりを回る内側の物体は、それへ引いた半径で、最大の物体が、これらの引力によって動かされる場合のほうが、最大の物体が小さい物体に引かれずに静止していたり、はるかに小さくまたははるかに大きく引かれたり、はるかに小さくまたははるかに大きく動かされたりする場合よりも、時間にいっそう比例する面積を描き、半径が出会う点を焦点にもつ楕円の形に、いっそう近づく図形を描く、と私は言う。}

\pngright{185.png}{}{}{}

これは、前の命題の系2の証明からほぼ明らかであるが、より明確で、より強い論法で、次のように証明される。

\textit{場合1.}\wsp{}小さな物体 $P$ と $Q$ が同じ平面で最大の物体 $S$ のまわりを回り、そのうち $P$ は内側の軌道 $PAB$ を、$Q$ は外側の軌道 $QE$ を描くとする。$QK$ を物体 $P$ と $Q$ の平均の距離とし、その平均の距離での物体 $P$ の $Q$ に向かう加速的な引力を、同じ線で表す。$QL$ を、$QL$ 対 $QK$ が $QK$ 対 $QP$ の二乗の比になるようにとれば、$QL$ は、任意の距離 $QP$ での物体 $P$ の $Q$ に向かう加速的な引力である。$PS$ を結び、それに平行に $LM$ を引いて $QS$ と $M$ で交わらせれば、引力 $QL$ は（法則の系2により）引力 $QM$, $LM$ に分解される。こうして物体 $P$ は、三重の加速力で押される。その一つは $S$ に向かい、物体 $S$ と $P$ の相互の引力から生じる。この力だけなら、物体 $P$ は、物体 $S$ が動かないにせよ、この引力で動かされるにせよ、そのまわりに、半径 $PS$ が時間に比例する面積を描き、物体 $S$ の中心に焦点をもつ楕円を描くはずである。これは問題VIと定理XXIの系から明らかである。第二の力は引力 $LM$ で、これは $P$ から $S$ に向かうから、前の力に加われば、それと一致し、こうして、定理XXIの系3により、面積はなお時間に比例して描かれる。ところが、これは距離 $PS$ の平方に反比例しないから、前の力と合成して、この比から外れる力をつくり、それは、他が同じなら、この力の前の力に対する比が大きいほど、いっそう外れる。したがって、（問題VIIIの系1と定理XXIの系2により）焦点 $S$ のまわりに楕円を描く力は、その焦点に向かい、距離 $PS$ の平方に反比例しなければならないから、この合成された力は、この比から外れて、軌道 $PAB$ を、$S$ に焦点をもつ楕円の形から外れさせる。それは、この比からの外れが大きいほど、いっそう外れ、したがってまた、他が同じなら、第二の力 $LM$ の第一の力に対する比が大きいほど、いっそう外れる。さらに、第三の力 $QM$ は、物体 $P$ を、$QS$ に平行な線に沿って引くから、前の力と合成して、もはや $P$ から $S$ に向かわない力をつくる。この力は、他が同じなら、この第三の力の前の力に対する比が大きいほど、この向きからいっそう外れる。したがって、物体 $P$ は、半径 $SP$ で、もはや時間に比例する面積を描かず、この比例からの外れは、この第三の力のほかの力に対する比が大きいほど、大きくなる。そして軌道 $PAB$ の、前に述べた楕円の形からの外れを、この第三の力は二つの理由で増す。一つは、$P$ から $S$ に向かわないこと、もう一つは、距離 $PS$ の平方に比例しないことである。これらがわかれば、面積が時間に最もよく比例するのは、ほかの力が変わらないまま、第三の力が最小になるときであり、また軌道 $PAB$ が前述の楕円の形に最も近づくのは、第一の力が変わらないまま、第二の力も第三の力も、とりわけ第三の力が最小になるときである。

物体 $S$ の $Q$ に向かう加速的な引力を、線 $QN$ で表す。そして、もし加速的な引力 $QM$ と $QN$ が等しければ、これらは物体 $S$ と $P$ を、等しく、平行線に沿って引くから、それらの互いの位置を少しも変えない。そのとき、これらの物体の互いの運動は（法則の系6により）、これらの引力がなくなった場合と同じになる。同じ理由で、もし引力 $QN$ が引力 $QM$ より小さければ、それは引力 $QM$ の部分 $QN$ を取り去り、部分 $MN$ だけが残って、これが時間と面積の比例、および軌道の楕円の形を乱す。そして同様に、もし引力 $QN$ が引力 $QM$ より大きければ、差 $MN$ だけから、比例と軌道の乱れが生じる。こうして、引力 $QN$ によって、上の第三の引力 $QM$ はつねに引力 $MN$ に還元され、第一と第二の引力はまったく変わらない。したがって、面積と時間が比例に、軌道 $PAB$ が前述の楕円の形に最も近づくのは、引力 $MN$ がないか、できるだけ小さいときである。すなわち、物体 $Q$ に向かう物体 $P$ と $S$ の加速的な引力が、できるだけ等しさに近づくときである。いいかえれば、引力 $QN$ が、ゼロでも、すべての $QM$ の最小のものより小さくもなく、すべての引力 $QM$ の最大と最小のいわば中間、すなわち引力 $QK$ よりもはるかに大きくもはるかに小さくもないときである。\QED

\textit{場合2.}\wsp{}いま、小さな物体 $P$, $Q$ が、異なる平面で最大の物体 $S$ のまわりを回るとする。すると、軌道 $PAB$ の平面にある線 $PS$ に沿って働く力 $LM$ は、前と同じ効果をもち、物体 $P$ をその軌道の平面から外さない。ところが、もう一つの力 $NM$ は、$QS$ に平行な線に沿って働くから（したがって物体 $Q$ が交点線の外にあるときには、軌道 $PAB$ の平面に傾く）、前に述べた経度の運動の乱れのほかに、緯度の運動の乱れを引き起こし、物体 $P$ をその軌道の平面から引き離す。そしてこの乱れは、物体 $P$ と $S$ の互いの任意の位置では、それを生じる力 $MN$ に比例し、したがって、$MN$ が最小のとき、すなわち（すでに述べたように）引力 $QN$ が引力 $QK$ よりはるかに大きくもはるかに小さくもないとき、最小になる。\QED

\textit{系1.}\wsp{}ここから、容易に推論されるように、もし多くの小さな物体 $P$, $Q$, $R$ などが最大の物体 $S$ のまわりを回るなら、最内の物体 $P$ の運動は、外側の物体の引力によって、最も乱されない。それは、最大の物体 $S$ が、同様に、ほかの物体から、加速力の比に応じて引かれ動かされ、ほかの物体も互いにそうされるときである。

\textit{系2.}\wsp{}三つの物体 $S$, $P$, $Q$ の系で、どの二つの物体の第三の物体への加速的な引力も、互いに距離の平方に反比例するならば、物体 $P$ は、物体 $S$ のまわりに半径 $PS$ で、合 $A$ と衝 $B$ の近くでは、矩 $C$, $D$ の近くより速く面積を描く。実際、物体 $P$ を押し、物体 $S$ を押さず、線 $PS$ に沿って働かないすべての力は、面積の描かれ方を、先行する向きか後続の向きかに向けられるのに応じて、加速するか減速する。そのような力が $NM$ である。この力は、物体 $P$ が $C$ から $A$ へ進むあいだは先行の向きに向かい、運動を加速し、ついで $D$ まで後続の向きに向かい、運動を減速し、ついで $B$ まで先行の向き、最後に $B$ から $C$ へ移るときには後続の向きとなる。

\textit{系3.}\wsp{}同じ論法で、物体 $P$ は、他が同じなら、合と衝では、矩におけるより速く動くことが明らかである。

\textit{系4.}\wsp{}物体 $P$ の軌道は、他が同じなら、矩では、合と衝におけるより曲がりが大きい。というのは、速い物体は、直線の道からそれにくいからである。さらに、力 $NM$ は、合と衝では、物体 $S$ が物体 $P$ を引く力と逆向きで、したがってその力を減らし、物体 $P$ は、物体 $S$ のほうへ押される力が小さいほど、直線の道からそれにくい。

\textit{系5.}\wsp{}したがって、物体 $P$ は、他が同じなら、矩では、合と衝よりも、物体 $S$ から遠ざかる。これらは、離心率の運動を除いたときのことである。というのは、もし物体 $P$ の軌道が離心的なら、その離心率は（この命題の系9で、すぐあとで示されるように）アプシスが朔望にあるとき最大になり、そこから、物体 $P$ が遠いアプシスに達するとき、朔望では矩におけるより、物体 $S$ から遠く離れることがありうるからである。

\textit{系6.}\wsp{}中心の物体 $S$ の向心力は、それによって物体 $P$ が軌道上にとどめられるが、矩では力 $LM$ が加わって増し、朔望では力 $KL$ が引かれて減る。そして力 $KL$ の大きさのために、増すより減るほうが大きい。ところで、この向心力は（命題IVの系2により）半径 $SP$ に単純な比で正比例し、周期の二乗に反比例する、という合成された比をなす。したがって、この合成された比は、力 $KL$ の作用によって小さくなり、したがって、軌道の半径 $SP$ が変わらなければ、周期は、その向心力が小さくなる比の平方根の比で、増えることが明らかである。そして、この半径を増しても減らしても、周期は、命題IVの系6により、この半径の3/2乗の比より多く増すか少なく減る。もしその中心の物体の力が少しずつ弱まるなら、物体 $P$ は、つねに引かれる力が小さくなり、つねに中心 $S$ からより遠くへ遠ざかるであろう。逆に、もしその力が増せば、より近づくであろう。ゆえに、その力を減らす遠い物体 $Q$ の作用が、かわるがわる増減するなら、半径 $SP$ も同時にかわるがわる増減し、周期は、半径の3/2乗の比と、遠い物体 $Q$ の作用の増減によって中心の物体 $S$ の向心力が減るか増す比の平方根の比とから合成された比で、増減する。

\textit{系7.}\wsp{}前に述べたことから、物体 $P$ が描く楕円の軸、すなわちアプシス線は、角運動に関して、かわるがわる進み、また退くが、しかし進むほうが大きく、物体が一回転するごとに、進みの超過分だけ、後続の向きに運ばれることも導かれる。というのは、物体 $P$ が矩で物体 $S$ のほうへ押される力は、力 $MN$ が消えているから、力 $LM$ と、物体 $S$ が物体 $P$ を引く向心力とから合成される。前の力 $LM$ は、距離 $PS$ が増すと、この距離とほぼ同じ比で増し、後の力はその二乗の比で減るから、これらの力の和は、距離 $PS$ の二乗より小さい比で減り、したがって、命題XLVの系1により、遠地点、すなわち遠いアプシスを後退させる。ところが、合と衝では、物体 $P$ が物体 $S$ のほうへ押される力は、物体 $S$ が物体 $P$ を引く力と力 $KL$ との差であり、この差は、力 $KL$ がほぼ距離 $PS$ の比で増すために、距離 $PS$ の二乗より大きい比で減るから、命題XLVの系1により、遠地点を進ませる。朔望と矩の間の位置では、遠地点の運動は、どちらの原因にも合わせて依存し、一方か他方の超過に応じて、進むか退く。そこで、朔望での力 $KL$ が、矩での力 $LM$ のほぼ2倍であるから、全回転での超過は力 $KL$ のほうにあり、毎回転ごとに遠地点を後続の向きに運ぶ。この系と前の系の真実は、二つの物体 $S$, $P$ の系が、軌道 $QE$ にある多くの物体 $Q$, $Q$, $Q$ などに、四方から取り巻かれていると考えれば、いっそう容易に理解される。実際、これらの作用によって、$S$ の作用は四方で小さくなり、距離の二乗より大きい比で減少する。

\textit{系8.}\wsp{}ところで、アプシスの進みや後退は、物体が近いアプシスから遠いアプシスへ進むときに、向心力の減少が、距離 $SP$ の二乗より大きい比か小さい比で起こること、また、近いアプシスへもどるときの同様の増加に依存し、したがって、遠いアプシスでの力の近いアプシスでの力に対する比が、距離の逆二乗の比から最も離れるとき、最大になるのだから、アプシスは、朔望では、減らす力 $KL$、すなわち $NM - LM$ によって、いっそう速く進み、矩では、加える力 $LM$ によって、よりゆっくり退くことは明らかである。ところで、進みの速さや退きの遅さが続く時間が長いので、この不等は、はるかに大きくなる。

\textit{系9.}\wsp{}もし、ある物体が、自分の中心からの距離の平方に反比例する力で、この中心のまわりを楕円の上で回り、まもなく、遠いアプシス、すなわち遠地点から近いアプシスへ下るときに、その力が、新しい力が絶えず加わることによって、減少した距離の二乗より大きい比で増すとしよう。すると、この新しい力の絶えざる加わりによって、つねに中心へ押された物体は、減少した距離の二乗の比で増す力だけで押される場合より、いっそうこの中心のほうへ向かい、したがって楕円の軌道より内側の軌道を描いて、近いアプシスでは前より中心に近づくことは明らかである。それゆえ軌道は、この新しい力が加わることによって、いっそう離心的になる。いま、物体が近いアプシスから遠いアプシスへ遠ざかるとき、その力が、前に増したのと同じ度合いで減るなら、物体は前の距離にもどるであろう。したがって、もし力がより大きい比で減るなら、物体は、いまやより小さく引かれて、より大きな距離まで上り、こうして軌道の離心率はさらに増す。ゆえに、向心力の増減の比が、一回転ごとに増すなら、離心率はつねに増し、逆にその比が減るなら、離心率は減る。ところで、物体 $S$, $P$, $Q$ の系では、軌道 $PAB$ のアプシスが矩にあるとき、この増減の比は最小で、アプシスが朔望にあるとき最大になる。もしアプシスが矩に置かれるなら、アプシスの近くでは比が距離の二乗より小さく、朔望の近くでは大きく、その大きいほうの比から、すでに述べたように、遠地点の最も速い運動が生じる。しかし、アプシスの間の進行における全体の増減の比を考えるなら、これは距離の二乗より小さい。近いアプシスでの力の遠いアプシスでの力に対する比は、遠いアプシスの楕円の焦点からの距離の、近いアプシスのその焦点からの距離に対する比の二乗より、小さい。逆に、アプシスが朔望に置かれるとき、近いアプシスでの力の遠いアプシスでの力に対する比は、距離の二乗より大きい。というのは、矩で物体 $S$ の力に加えられる力 $LM$ は、より小さい比の力を合成し、朔望で物体 $S$ の力から引かれる力 $KL$ は、より大きい比の力を残すからである。ゆえに、アプシスの間を通るときの増減の全体の比は、矩で最小、朔望で最大である。したがって、アプシスが矩から朔望へ移るときには、絶えず増し、楕円の離心率を増し、朔望から矩へ移るときには、絶えず減り、離心率を減らす。

\textit{系10.}\wsp{}緯度の誤差の比を求めるために、軌道 $QES$ の平面が動かないままであると仮定する。誤差の原因の上述の説明から、その原因のすべてである力 $NM$, $ML$ のうち、力 $ML$ は、つねに軌道 $PAB$ の平面に沿って働くから、緯度の運動を決して乱さないこと、また力 $NM$ は、交点が朔望にあるとき、やはり同じ軌道の平面に沿って働くから、これらの運動を乱さないが、矩にあるときには最も乱し、物体 $P$ を軌道の平面から絶えず引き離して、物体が矩から朔望へ移るあいだは平面の傾きを小さくし、逆に朔望から矩へ移るあいだはその傾きを大きくすることは、明らかである。したがって、物体が朔望にあるとき、傾きはすべての中で最小になり、物体が次の交点に近づくとき、ほぼもとの大きさにもどる。ところが、もし交点が矩のあとの八分点、すなわち $C$ と $A$、$D$ と $B$ の間に置かれるならば、いま述べたことから、物体 $P$ がどちらかの交点から、そこから90度の位置まで進む間、平面の傾きは絶えず小さくなり、ついで次の矩までの45度を通る間は傾きが増し、そののち、次の交点までのさらに45度を通る間は、ふたたび減ることがわかる。したがって傾きは、増えるより減るほうが大きく、そのため、あとの交点では、つねに前の交点より小さい。同様の推論によって、交点が他方の八分点、すなわち $A$ と $D$、$B$ と $C$ の間にあるときは、傾きは減るより増すほうが大きい。ゆえに、傾きは交点が朔望にあるとき、すべての中で最大である。交点が朔望から矩へ移るあいだは、物体が交点に達するごとに減り、交点が矩にあり物体が朔望にあるとき、すべての中で最小になる。ついで、前に減ったのと同じ度合いで増し、交点が次の朔望に達すれば、はじめの大きさにもどる。

\textit{系11.}\wsp{}交点が矩にあるとき、物体 $P$ は、自分の軌道の平面から絶えず引き離される。それも、交点 $C$ から合 $A$ を通って交点 $D$ へ進む間は $Q$ のほうへ、交点 $D$ から衝 $B$ を通って交点 $C$ へ進む間は反対の側へ引かれる。したがって、交点 $C$ からの運動で、物体は、次の交点に達するまで、自分の軌道のはじめの平面 $CD$ から絶えず遠ざかる。したがってこの交点では、そのはじめの平面 $CD$ から最も遠く離れて、軌道 $QES$ の平面を通る。その場所は、その平面のもう一方の交点 $D$ ではなく、そこから物体 $Q$ の側へ寄った点である。したがって、これが新しい交点の位置で、前の方へ向かう。同様の論法で、物体がこの交点から次の交点へ移るあいだにも、交点は退きつづける。ゆえに、矩に置かれた交点は絶えず退き、朔望（そこでは緯度の運動はまったく乱されない）では静止し、中間の位置では、両方の条件に関与して、よりゆっくり退く。したがって、毎回転ごとに、つねに逆行するか静止するかのどちらかで、先行する向きに運ばれる。

\textit{系12.}\wsp{}これらの系で述べたすべての誤差は、物体 $P$, $Q$ の合のときのほうが、それらの衝のときより、やや大きい。それは、生じさせる力 $NM$, $ML$ がより大きいからである。

\textit{系13.}\wsp{}これらの系の論拠は、物体 $Q$ の大きさに依存しないから、前のすべては、物体 $Q$ の大きさが、そのまわりを二つの物体 $S$, $P$ の系が回るほど大きく定められる場合にも成り立つ。そして物体 $Q$ を大きくすると、その向心力も増し、物体 $P$ の誤差はそれから生じるから、これらの誤差は、すべて（距離が等しければ）、物体 $Q$ が物体 $P$, $S$ の系のまわりを回る他方の場合より、この場合のほうが大きくなる。

\textit{系14.}\wsp{}ところで、物体 $Q$ が遠いとき、力 $NM$, $ML$ は、ほぼ、力 $QK$ と、比 $PS$ 対 $QS$ との積に比例する。すなわち、距離 $PS$ と物体 $Q$ の絶対力がともに与えられていれば、$QS^3$ に反比例する。そして、これらの力 $NM$, $ML$ は、これまでの系で扱った誤差とすべての効果の原因であるから、効果はすべて、物体 $S$, $P$ の系が変わらないとき、ほぼ、物体 $Q$ の絶対力に正比例し、距離 $QS$ の3乗に反比例する、という合成された比をなすことは明らかである。したがって、もし物体 $S$ と $P$ の系が、遠い物体 $Q$ のまわりを回るならば、これらの力 $NM$, $ML$ とそれらの効果は、（命題IVの系2と6により）周期の二乗に反比例する。そこから、もし物体 $Q$ の大きさがその絶対力に比例するなら、これらの力 $NM$, $ML$ とそれらの効果は、物体 $S$ から見た遠い物体 $Q$ の見かけの直径の3乗に正比例し、逆も同様である。実際、これらの比は、上の合成された比と同じだからである。

\textit{系15.}\wsp{}そして、軌道 $QE$ と $PAB$ の形、互いの比と傾きが変わらないまま、それらの大きさが変わり、物体 $Q$ と $S$ の力が変わらないか、あるいは任意の与えられた比で変わるとき、これらの力（すなわち、物体 $P$ を直線の道から軌道 $PAB$ へそらす物体 $S$ の力と、同じ物体 $P$ をその軌道からそれさせる物体 $Q$ の力）は、つねに同じ仕方、同じ比で働く。そこで、すべての効果は相似で比例し、効果の時間も比例しなければならない。すなわち、すべての線の誤差は軌道の直径に比例し、角の誤差は前と同じで、相似な線の誤差や等しい角の誤差の時間は、軌道の周期に比例する。

\textit{系16.}\wsp{}したがって、もし軌道の形と互いの傾きが与えられ、物体の大きさ、力、距離が任意に変わるなら、ある一つの場合の与えられた誤差と誤差の時間から、ほかの任意の場合の誤差と誤差の時間を、ほぼ求めることができる。ただし、次の方法のほうが簡単である。力 $NM$, $ML$ は、他が同じなら半径 $SP$ に比例し、これらの周期的な効果は（補助定理Xの系2により）、力と物体 $P$ の周期の二乗との積に比例する。これらが物体 $P$ の線の誤差であり、したがって、中心 $S$ から見た角の誤差（すなわち、遠地点と交点の運動、および経度と緯度の見かけの誤差のすべて）は、物体 $P$ の毎回転ごとに、ほぼ回転の時間の二乗に比例する。これらの比を系14の比と合わせれば、物体 $S$, $P$, $Q$ の任意の系で、$P$ が近くの $S$ のまわりを、$S$ が遠くの $Q$ のまわりを回るとき、物体 $P$ の、中心 $S$ から見た角の誤差は、その物体 $P$ の毎回転ごとに、物体 $P$ の周期の二乗に正比例し、物体 $S$ の周期の二乗に反比例する。そしてそこから、遠地点の平均運動は、交点の平均運動に対して与えられた比をなし、どちらの運動も、物体 $P$ の周期に正比例し、物体 $S$ の周期の二乗に反比例する。軌道 $PAB$ の離心率と傾きを増減しても、遠地点と交点の運動は、それらがあまり大きくない限り、感じられるほどには変わらない。

\textit{系17.}\wsp{}ところで、線 $LM$ は、ときに半径 $PS$ より大きく、ときに小さいから、平均の力 $LM$ を、その半径 $PS$ で表す。すると、これは平均の力 $QK$ または $QN$（これは $QS$ で表せる）に対して、長さ $PS$ 対長さ $QS$ の比をなす。ところで、物体を $Q$ のまわりの軌道にとどめる平均の力 $QN$ または $QS$ は、物体 $P$ を $S$ のまわりの軌道にとどめる力に対して、半径 $QS$ 対半径 $PS$ の比と、物体 $P$ の $S$ のまわりの周期の、物体 $S$ の $Q$ のまわりの周期に対する比の二乗との合成された比をなす。そして等比により、平均の力 $LM$ は、物体 $P$ を $S$ のまわりの軌道にとどめる力（あるいは、同じ物体 $P$ が、同じ周期で、任意の動かない点 $S$ のまわりを距離 $PS$ で回りうる力）に対して、この周期の比の二乗の比をなす。ゆえに、周期が距離 $PS$ とともに与えられれば、平均の力 $LM$ が与えられ、それが与えられれば、線 $PS$, $MN$ の比例によって、力 $MN$ もほぼ与えられる。

\textit{系18.}\wsp{}物体 $P$ が物体 $S$ のまわりを回るのと同じ法則で、多くの流体の物体が、同じ $S$ のまわりを、それから等しい距離で動くとし、ついでこれらを接触させて、物体 $S$ と同心の、丸い流体の輪をつくると考える。すると、輪の各部分は、すべての運動を物体 $P$ の法則にしたがって行い、物体 $S$ にいっそう近づき、それら自身と物体 $Q$ の合と衝では、矩におけるより速く動く。そしてこの輪の交点、すなわち物体 $Q$ または $S$ の軌道の平面との交わりは、朔望では静止し、朔望以外では先行する向きに動き、とくに矩では最も速く、ほかの位置ではより遅い。輪の傾きも変化し、その軸は毎回転ごとに振動し、回転が完了すれば、交点の歳差によって運び去られる分を除いて、もとの位置にもどる。

\textit{系19.}\wsp{}いま、物体 $S$ の球が、流体でない物質から成っていて、この輪にまで広がり伸びて、まわりにえぐった溝に水を含み、同じ周期で、自分の軸のまわりを一様に回るとする。この液体は、（前の補助定理でのように）かわるがわる加速され減速されて、朔望では球の表面より速く、矩では遅くなり、こうして溝の中を、海のように満ち引きする。もし物体 $Q$ の引力が取り去られれば、水は、球の静止した中心のまわりを回っても、満ち引きの運動をまったく得ない。球が直線上を一様に進みながら、同時に自分の中心のまわりを回る場合も、同じ理由である。

そこで物体 $Q$ を近づければ、その不等な引力によって、まもなく水は乱される。実際、近い水への引力は大きく、遠い水への引力は小さい。そして力 $LM$ は、矩では水を下へ引いて、朔望まで下降させ、力 $KL$ は、朔望では同じ水を上へ引き、その下降を止めて、矩まで上昇させる。

\textit{系20.}\wsp{}もし輪が固くなり、球が小さくなるなら、流れて引く運動はやむが、傾きの振動運動と交点の歳差は残る。球が輪と同じ軸をもち、同じ時間で回転を完了し、自分の表面で輪の内側に接して、それにくっついているとし、その運動に加わるなら、両者の全体は振動し、交点は後退する。というのは、球は、すぐあとで述べるように、あらゆる印象を受け取ることについて、無差別だからである。球を取り去った輪の最大の傾きの角は、交点が朔望にあるときである。そこから交点が矩へ進むにつれて、輪は傾きを小さくしようとし、この努力によって、球全体に運動を与える。球は、輪が逆の努力によってこの運動を取り去り、反対の側に新しい運動を与えるまで、与えられた運動を保つ。こうして、減少する傾きの最大の運動は、交点の矩で起こり、傾きの角が最小になるのは、矩のあとの八分点であり、ついで、傾きがもどる最大の運動は朔望で、最大の角は次の八分点で起こる。輪を取り去った球でも、赤道の地方が極の近くよりやや高いか、あるいはやや密な物質でできているなら、理由は同じである。実際、赤道の地方での物質の超過分が、輪の代わりをするからである。そして、この球の向心力を任意に増して、その部分がすべて、地球の重い部分のように下へ向かうと仮定しても、この系と前の系の現象は、そのためにほとんど変わらない。

\textit{系21.}\wsp{}球の赤道近くの余分な物質が交点を後退させるのと同じ理由で、この増加によってこの後退は増し、減少によって減り、取り去ることによってなくなる。もし余分以上の物質が取り去られるなら、すなわち、球が赤道近くで極の近くよりへこんでいるか、まばらになるなら、交点は後続の向きに動く。

\textit{系22.}\wsp{}そこから逆に、交点の運動から、球の構造がわかる。すなわち、もし球が極をいつも同じに保ち、運動が先行の向きなら、赤道の近くの物質は余分である。後続の向きなら、不足している。均質で完全に円い球が、自由な空間にはじめ静止しているとし、ついで任意の衝撃が表面に斜めに加えられて押され、そこから、一部は円の、一部は直線の運動を得るとする。この球は、中心を通るすべての軸に対して無差別であり、ある軸、またはある軸の位置に、他の軸よりも傾きやすいということがないから、自分の力では、自分の軸と軸の傾きを決して変えないことは明らかである。いま球が、前と同じ表面の部分に、任意の新しい衝撃で斜めに押されるとする。衝撃が早いか遅いかは効果を少しも変えないから、これら二つの衝撃が続けて与えられても、同時に与えられた場合と同じ運動を生じることは明らかである。すなわち、球が両方から（法則の系2により）合成された単一の力で押された場合と同じ運動、したがって、傾きが与えられた軸のまわりの単一の運動である。同じ理由は、第一の運動の赤道上の任意の別の場所に、第二の衝撃が加えられる場合、また、第二の衝撃だけが生じる運動の赤道上の任意の場所に第一の衝撃が加えられる場合、したがって、任意の場所に加えられた二つの衝撃の場合にも成り立つ。これらは、それぞれ別々に生じる運動の赤道の交わる場所に、同時に一度に与えられた場合と同じ円運動を生じる。ゆえに、均質で完全な球は、いくつもの別々の運動を保たず、与えられたすべてを合成して一つにまとめ、自分としては、傾きがつねに不変に与えられた単一の軸のまわりを、つねに単一の一様な運動で回る。しかし、向心力も、軸の傾きや回転の速さを変えることはできない。球が、自分の中心と、力が向けられる中心とを通る任意の平面で、二つの半球に分けられると考えるなら、その力は、つねに両方の半球を等しく押すから、球の回転運動を、どちらの側にも傾けない。ところが、極と赤道の間のどこかに、山の形に積み上げた新しい物質を加えれば、これは、自分の運動の中心から離れようとする絶えざる努力によって、球の運動を乱し、極をその表面の上でさまよわせて、自分と、自分の向かい側の点とのまわりに、絶えず円を描かせるであろう。このさまよいの異常は、その山をどちらかの極に置くか（この場合には、系21により、赤道の交点は進む）、あるいは赤道に置くか（この場合には、系20により、交点は後退する）、あるいは最後に、軸の反対側に新しい物質を加えて、山が動くあいだつりあうようにしなければ、正されない。そしてこのとき、交点は、山とこの新しい物質が極と赤道のどちらに近いかに応じて、進むか退く。
\condpagelarge{命題 LXVII. 定理 XXVII.}

\textit{同じ引力の法則を仮定するとき、外側の物体 $Q$ が、内側の物体 $P$, $S$ の共通の重心 $C$ のまわりに、その中心へ引いた半径で描く面積は、最も内側で最大の物体 $S$ のまわりに、それへ引いた半径で描きうる面積よりも、時間にいっそう比例する。また、その軌道は、同じ中心に焦点をもつ楕円の形に、いっそう近づく。}

実際、物体 $Q$ の $S$ と $P$ への引力は、その絶対的な引力を合成する。この引力は、最大の物体 $S$ よりも、物体 $S$ と $P$ の共通の重心 $C$ のほうへいっそう向けられ、また距離 $QS$ の平方よりも、距離 $QC$ の平方にいっそう反比例する。よく考えてみれば、このことは容易にわかる。

\condpagelarge{命題 LXVIII. 定理 XXVIII.}

\textit{同じ引力の法則を仮定するとき、外側の物体 $Q$ が、内側の物体 $P$, $S$ の共通の重心 $C$ のまわりに、その中心へ引いた半径で描く面積は、最も内側で最大の物体が、これらの引力を、ほかの物体と同じように受けて動かされる場合のほうが、それが引かれずに静止したり、はるかに強くまたははるかに弱く引かれたり、はるかに強くまたははるかに弱く動かされたりする場合よりも、時間にいっそう比例する。また、その軌道は、同じ中心に焦点をもつ楕円の形に、いっそう近づく。}

\pngright{185.png}{}{}{}

証明は、命題LXVIとほぼ同じ仕方でできるが、論述がより長くなるので、省く。次のように見積もれば足りるであろう。直前の命題の証明から、物体 $Q$ が合わさった力で押される中心は、その二つの物体の共通の重心に最も近いことが明らかである。もしこの中心がその共通の重心に一致し、三つの物体の共通の重心が静止しているならば、物体 $Q$ は一方の側に、ほかの二つの物体の共通の重心は他方の側に、すべての物体の静止した共通の中心のまわりに、正確な楕円を描くであろう。これは、命題LVIIIの系2を、命題LXIVおよびLXVで証明したことと照らし合わせれば明らかである。この楕円運動は、二つの物体の重心と、第三の物体 $Q$ が引かれる中心との距離によって、いくらか乱される。さらに三つの物体の共通の重心に運動を与えれば、乱れは増す。したがって、乱れは、三つの物体の共通の重心が静止しているとき、すなわち、最も内側で最大の物体 $S$ が、ほかの物体と同じ法則で引かれるときに最小になる。そして、その三つの物体の共通の重心が、物体 $S$ の運動を減らすことによって動き始め、それからしだいにいっそう動かされるようになるとき、乱れはつねにより大きくなる。

\textit{系.}\wsp{}ここから、もし多くの小さな物体が最大の物体のまわりを回るなら、次のように推論できる。すべての物体が、絶対力に正比例し距離の平方に反比例する加速力で互いに引き合い動かし合い、また各軌道の焦点が、それより内側のすべての物体の共通の重心に置かれる場合には（すなわち、第1の最も内側の軌道の焦点は最大で最も内側の物体の重心に、第2の軌道のそれは二つの最も内側の物体の共通の重心に、第3の軌道のそれは三つの内側の物体の共通の重心に、以下同様に）、最も内側の物体が静止して、すべての軌道に共通の焦点とされる場合よりも、描かれる軌道は楕円にいっそう近づき、面積の描かれ方はいっそう一様になる。

\condpagelarge{命題 LXIX. 定理 XXIX.}

\textit{多くの物体 $A$, $B$, $C$, $D$ などの系で、ある物体 $A$ が、ほかのすべての物体 $B$, $C$, $D$ などを、引く物体からの距離の平方に反比例する加速力で引き、別の物体 $B$ も、ほかのすべての物体 $A$, $C$, $D$ などを、引く物体からの距離の平方に反比例する力で引くならば、引く物体 $A$, $B$ の絶対力は、互いに、それらの力をもつ物体 $A$, $B$ そのものに比例する。}

実際、物体 $B$, $C$, $D$ すべての $A$ への加速的な引力は、等しい距離では、仮定により互いに等しく、同様に、すべての物体の $B$ への加速的な引力も、等しい距離では互いに等しい。ところで、物体 $A$ の絶対的な引力の、物体 $B$ の絶対的な引力に対する比は、等しい距離での、すべての物体の $A$ への加速的な引力の、すべての物体の $B$ への加速的な引力に対する比に等しく、これは、物体 $B$ の $A$ への加速的な引力の、物体 $A$ の $B$ への加速的な引力に対する比に等しい。ところが、物体 $B$ の $A$ への加速的な引力の、物体 $A$ の $B$ への加速的な引力に対する比は、物体 $A$ の質量の物体 $B$ の質量に対する比に等しい。それは、（定義II, VII, VIIIにより）加速力に、引かれる物体を掛けて生じる起動力が、（運動の法則IIIにより）互いに等しいからである。ゆえに、物体 $A$ の絶対的な引力の、物体 $B$ の絶対的な引力に対する比は、物体 $A$ の質量の物体 $B$ の質量に対する比である。\QED

\textit{系1.}\wsp{}したがって、もし系の個々の物体 $A$, $B$, $C$, $D$ などが、別々に見て、ほかのすべてを、引く物体からの距離の平方に反比例する加速力で引くならば、それらすべての物体の絶対力は、互いに、物体そのものに比例する。

\textit{系2.}\wsp{}同じ論法で、もし系の個々の物体 $A$, $B$, $C$, $D$ などが、別々に見て、ほかのすべてを、引く物体からの距離の任意のべきに、反比例または正比例する加速力で、あるいは、各引く物体からの距離によって、任意の共通の法則で定められる加速力で引くならば、それらの物体の絶対力は、物体に比例することは明らかである。

\textit{系3.}\wsp{}力が距離の二乗の比で減少する物体の系で、もし小さい物体が、最大の物体のまわりを、その最大の物体の中心に共通の焦点をもつ、できるだけ正確な楕円の上を回り、最大の物体へ引いた半径で、できるだけ時間に比例する面積を描くならば、それらの物体の絶対力は、互いに、正確にか、ほぼ、物体の比をなす。逆も成り立つ。これは、命題LXVIIIの系を、この命題の系1と照らし合わせれば明らかである。

\condpagelarge{注.}

これらの命題によって、私たちは、向心力と、その力が向けられる中心の物体との類比へ導かれる。実際、物体に向けられる力が、その物体の本性と量に依存することは、磁気の場合のように、道理にかなっている。このような場合には、物体の引力は、その各粒子にそれぞれの力を割り当て、力の総和を求めることによって、見積もらなければならない。私はここで「引力」という語を、一般に、物体が互いに近づこうとする任意の努力の意味に使う。その努力が、物体が互いに求め合う作用から生じるにせよ、放出された精気によって互いに動かし合うことから生じるにせよ、あるいは、エーテルや空気、または物体的であれ非物体的であれ何らかの媒質の作用が、その中に浮かぶ物体を互いに押すことから生じるにせよ、である。同じ一般的な意味で、私は「衝撃」という語も用いる。この論文では、力の種類や物理的な性質ではなく、数学的な量と比を考えるからであり、それは定義で説明したとおりである。数学では、どんな条件を仮定しても、そこから導かれる力の量と比を調べなければならない。ついで物理学に下るときには、これらの比を現象と比べて、どのような力の条件が、引力をもつ物体の各種類にふさわしいかを知らなければならない。そしてそのとき初めて、力の種類、原因、物理的な理由について、より安全に論じることができるであろう。そこで、いま述べた仕方で引力をもつ粒子から成る球形の物体が、互いにどのような力で作用し合うはずか、またそこからどのような運動が生じるかを、見ることにしよう。
\section{第1編 第12節：球体の引力（殻定理、命題LXX〜LXXXIV）}


% Section: 第1編 第12節：球体の引力（殻定理、命題LXX〜LXXXIV）



\sectpage{第12節}

\textbf{球体表面の点に対する向心力}

\begin{center}
\textbf{球形表面の物体に対する引力について}
\end{center}

\propnopage{命題LXX. 理論XXX.}

\textbf{もし球形表面の各点に指向する力が等しく、その強度は半径から各点までの距離の二乗反比例であるとすると、内部にある物体はこれらの力により表面のどの部分にも引き寄せられない。}

\pngright{201.png}{}{}{}
%Illustration

球体 $HIKL$ の表面を考えてみましょう。その中点に物体 $P$ を置きます。物体 $P$ から球面上の点 $H, K$ への線分 $HK$ が作用します。


% 略


\textbf{第12節：球体の引力（殻定理、命題LXX〜LXXXIV）}

\textit{IL} は \textit{HI} と \textit{KL} の最小接する弧であり；また三角形 \textit{HPI} と \textit{LPK} （補則 VII の第3節による）が相似であるから、これらの弧は距離 \textit{HP} と \textit{LP} に比例し、そして球面の任意部分 \textit{S}，これらを直線 \textit{HI} と \textit{KL} で通る点 \textit{P} で終了するものは、その二重比においても終焉します。したがって、これらの部分が体 \textit{P} 上に行使される力は互いに等しいです。これらは直接の部分と距離の平方の逆数に比例し、そして両方の比率が合成すると

\[
\textbf{viarum} \propto \frac{\textbf{particul{\ae}}}{(\textbf{distantiarum})^2}
\]

となります。

\textbf{等位性。}したがって、引力は対向する部分に対して等位的に作用し、互いに相殺します。同様の論証によれば、球面全体上の任意の点から発生するすべての引力は、対向する引力によって完全に相殺されます。従って、物体 \( P \) はこれらの引力により何ら推進されることはありません。\QED

\condpagelarge{命題LXXI. 定理XXXI.}

\textit{同様の仮定のもとで、私は次のように主張します。球体外に置かれた物体は、その中心から一定距離にある点に対する引力は、その距離の平方に反比例します。}

\pngright{202.png}{}{}{}

\[
\text{設} \, AHKB, \, ahkb \, \text{は等しい球面の面積であり、中心点} \, S, \, s \, \text{から直径} \, AB, \, ab \, \text{を描いたもので、} \, P, \, p \, \text{はこれらの直径延長線上に位置する粒子です。}
\]

\[
\text{粒子} \, P, \, p \, \text{から線} \, PHK, \, PIL, \, phk, \, pil \, \text{を引くと、最大円} \, AHB, \, ahb \, \text{の等しい弧} \, HK, \, hk \, \& \, HL, \, hl \, \text{が減少します。}
\]

\[
\text{これらの線} \, HK, \, hk \, \& \, HL, \, hl \, \text{に対して、垂線} \, SD, \, sd \, \text{を引いて} \, PL, \, pl \, \text{と交点} \, F, \, f \, \text{を取る。}
\]

\[
\text{さらに、直径方向に垂線} \, IQ, \, iq \, \text{を引きます。そして} \, DS = ds \, \text{かつ} \, ES = es \, \text{であるから、}
\]

球体の引力（殻定理）

次に、点 \( E \) と \( e \) の微小な角 \( DPE \) および \( dpe \)，線分 \( PE, PF \) および \( pe, pf \)，さらに線素 \( DF, df \) を等しいものとして考えます。これらの最終的な比は、微小な角 \( DPE \) および \( dpe \) が消滅するときの等しいものです。

この設定に基づいて、線分 \( PI \) と \( PF \) の比は \( RI \) と \( DF \) の比に等しく，また \( pf \) と \( pi \) の比は \( DF \) または \( df \) と \( ri \) の比に等しいです。従って、
\[ PI \times pf : PF \times pi = RI : ri, \]
これは（補題 VII の第 3 コロラリウムにより）弧 \( IH \) と弧 \( ih \) の比に等しい。

また、線分 \( PI \) と \( PS \) の比は \( IQ \) と \( SE \) の比に等しく，さらに \( ps \) と \( pi \) の比は \( SE \) または \( se \) と \( iq \) の比に等しいです。従って、
\[ PI \times ps : PS \times pi = IQ : iq. \]

\[
\textbf{そして、比}\quad PI^2 \times pf \times ps \textbf{を}\quad pi^2 \times PF \times PS \textbf{に}
\]
\[
\textbf{結びつけると、} IH \times IQ \textbf{が} ih \times iq \textbf{に比するようになります。}
\]
\[
\textbf{すなわち、半円} AKB \textbf{を直径} AB \textbf{を中心に弧} IH \textbf{によって描く円の面積と、}
\]
\[
\textbf{半円} akb \textbf{を直径} ab \textbf{を中心に弧} ih \textbf{によって描く円の面積との比です。そして、}
\]
\[
\textbf{これらの円面上で体積} P \& p \textbf{に垂直な線方向に引き寄せられる力は（仮定により）、}
\]
\[
\textbf{それぞれの円面がその距離からなる正方形に比するように、すなわち、}
\]



\[
\textbf{以下のように比します。}
\]

\[
pf \times ps \text{ ad } PF \times PS.
\]
これらの力は、法則のコロラリ2により分解された結果力を沿って中心に向かう直線 \( PS, ps \) に作用します。すなわち、三角形 \( PIQ \) と \( PSF \)、および \( piq \) と \( psf \) が相似であるため、
\[
PS \text{ ad } PF \quad \& \quad ps \text{ ad } pf.
\]
したがって、粒子 \( P \) の引力 \( S \) 方向に対する引力と粒子 \( p \) の引力 \( s \) 方向に対する引力は、
\[
\frac{PF \times pf \times ps}{PS} \text{ ad } \frac{pf \times PF \times PS}{ps}
\]
すなわち、
\[
ps^2 \text{ ad } PS^2.
\]
同様の議論により、


球体の表面 $KL$, $kl$ によって描かれる粒子を引きつける力は、$ps^2$ に比例し、$PS^2$ に比例します；これらの比は同じです。また、すべての円形表面の力を区別するためには、常に $sd = SD$ および $se = SE$ を採用できます。そして、組合せにより、すべての球形表面の力が粒子に対して同一の比で働きます。$\blacksquare$

\condpagelarge{命題LXXII. 定理XXXII.}

円錐曲線 $KL$, $kl$ によって描かれる粒子を引きつける力は、$ps^2$ に比例し、$PS^2$ に比例します；これらの比は同じです。また、すべての球面の表面が粒子に対して作用する力を区別するためには、常に $sd = SD$ および $se = SE$ を採用できます。そして、組合せにより、すべての球面の表面が粒子に対して同一の比で作用します。$\blacksquare$


\textbf{球体の引力（殻定理）}

球体からの吸引力は、その中心から一定距離にある粒子に対して、距離の二乗に反比例し、また与えられた球体の直径とその粒子までの距離との比が与えられる。このとき、粒子が球体の中心に向かって引き寄せられる力は、球体の半径に比例する。

まず、異なる2つの球体からそれぞれ独立して粒子を引き寄せると仮定し、各球体の中心からの距離が直径に比例することも仮定する。さらに、球体は粒子に等しい部分で構成されていると想定する。これより、一つの粒子に対する一方の球体からの吸引力は、その球体の同じ部分に対するもう一方の粒子への吸引力に対して比例する。

\textbf{証明}

球体 \( S \) から一定距離 \( r \) にある粒子 \( P \) に働く引力を考慮する。この引力は距離 \( r \) の二乗に反比例し、また球体 \( S \) の直径と粒子 \( P \) からの距離との比が与えられる。

球体 \( S \) を粒子 \( P \) に対して等しい部分で構成すると仮定する。このとき、球体 \( S \) 内の任意の点 \( A \) から粒子 \( P \) への引力は、球体 \( S \) の中心 \( O \) からの距離 \( OA \) の二乗に反比例し、また粒子 \( P \) からの距離 \( PA \) との比が与えられる。

一方、球体 \( S \) 内の点 \( A \) から粒子 \( P \) への引力は、球体 \( S \) の半径 \( R \) に比例する。したがって、粒子 \( P \) に対する球体 \( S \) の吸引力は、球体 \( S \) の半径 \( R \) に比例することが証明される。

\[
\textbf{結論}
\]

粒子 \( P \) に対する球体 \( S \) の吸引力は、球体 \( S \) の半径 \( R \) に比例する。

\textbf{第12節：球体の引力（殻定理、命題LXX〜LXXXIV）}

\textbf{【第10部分】}

球体に対する粒子の数は、各粒子の比と半径の平方との逆比から成る合算比に等しい。すなわち、ある二つの球体に対する粒子の数は、各粒子の比と半径の平方の逆比から成る合算比に比例する。また、粒子は球体と同一であるため、粒子の比は直径の三重比に相当し、距離は直径と比例する。したがって、前記の比は直径に対する後者の比を二倍して逆数としたものに等しい。\QED

\textit{命題1.} すなわち、もし球体周囲の円軌道上を同一材料から成る粒子が引力によって引き寄せられるとするとき、これらの粒子の距離はそれぞれの球体の直径と比例し、その場合における周期は等しい。\QED

\textbf{証明}

まず、粒子の数について述べる。ある二つの球体に対する粒子の数は、各粒子の比と半径の平方との逆比から成る合算比に比例する。すなわち、

\[
\text{粒子の数} = \frac{\sum p_i}{r_1^2 + r_2^2}
\]

ここで \(p_i\) は各粒子の比、\(r_1, r_2\) はそれぞれの球体の半径である。

次に、粒子の比と距離について述べる。粒子は球体と同一であり、したがって粒子の比は直径の三重比に相当し、距離は直径と比例する。すなわち、

\[
p_i \propto d^3, \quad d = r
\]

ここで \(d\) は直径である。

以上の関係から、前記の比は直径に対する後者の比を二倍して逆数としたものに等しい。すなわち、

\[
\text{粒子の数} \propto \frac{1}{r^2}
\]

\textbf{命題1の証明}

もし球体周囲の円軌道上を同一材料から成る粒子が引力によって引き寄せられるとするとき、これらの粒子の距離はそれぞれの球体の直径と比例し、その場合における周期は等しい。\QED

\textbf{補遺1}


\textbf{命題LXXIII.}

球体内部の粒子は、与えられた点から等しい中心力が二倍の距離に比例して減少する場合、その粒子は球体内を周期的に移動し、これらの周期は等しく、距離と直径との比は等しいです。

\textbf{推論2.}

逆もまた然りである。すなわち、周期が等しいとき、距離は直径に比例します。この二つは定理IVの推論3によります。

\begin{center}
\includegraphics[width=763pt,height=981pt,viewport=-24 0 763 981]{204.png}
\end{center}

\textit{証明.}

もし与えられた球体上の点から等しい中心力が二倍の距離に比例して減少する場合、粒子は球体内を周期的に移動し、これらの周期は等しく、距離と直径との比も等しいことを示す。

まず、粒子が球体内部を円錐曲線軌道上を移動すると仮定します。この軌道は中心力の性質により、円または楕円となります。ここで、中心力は点から二倍の距離に比例して減少することを考慮に入れます。

次に、粒子が球体内部を移動する周期を考察します。粒子が一定の速度で球体内を移動すると仮定し、その速度は軌道上の位置によって変化します。中心力が二倍の距離に比例して減少するという条件により、粒子の運動は円錐曲線軌道上での運動となります。

周期は粒子が軌道を1周する時間であり、この周期は粒子の速度と軌道の形状に関連しています。粒子が球体内部を移動するとき、その速度は位置に依存しますが、中心力が二倍の距離に比例して減少するという条件により、粒子の運動は一定の規則性を持つことが示されます。

したがって、周期は等しくなり、距離と直径との比も等しいことが証明されました。

球体 \(ABCD\) 内に中心 \(S\) をもつ球面が描かれており、その内側に点 \(P\) があるとする。また、中心 \(S\) から点 \(P\) までの距離 \(SP\) で内接する球面 \(PEQF\) を描く。

このとき、第3章の理論によれば、同心球面 \(AEBF\) の差 \(AEBF\) は引力間の反対作用により相互に打ち消されるため、点 \(P\) に対して全く作用しないことが明らかである。残る唯一の引力は内側の球面 \(PEQF\) によるものである。

さらに、第32章の理論によれば、この引力は距離 \(PS\) の逆平方率に比例する。

\textbf{証明終了}

\condpagelarge{\textit{補説。}}

Superficies ex quibus solida componuntur, hic non sunt pure \textbf{数学的}, すなわち純粋に幾何学的なものではなく、非常に薄い円錐面で、その厚さは実質的に虚零とみなされる。すなわち、最終的には球体を構成する円錐面の数が増加し、その厚さが無限大まで減少する円錐面からなるものであり、これは最初に一般論において説明した方法による。同様に、線分、面積、立体が組み立てられる点群についても、これらは等しい大きさの部分を表すものとして理解されるべきである。




\textbf{Iisdem positis, dico quod球体外の粒子はその中心から求めた距離の平方反比例となる力で引き寄せられる。}

まず、球体を無数の同心円柱面に分割し、粒子が各円柱面上から受ける引力はその粒子から中心までの距離の平方反比例であるとします（命題XXXIによる）。これらを合成すると、全体の球体からの合計引力は同じ比となる。

\textbf{Corol.\ 1.} \quad 次に、均質な中心からの異なる距離での引力について述べる。

この証明は以下のように構成されます。

【第1編 第12節：球体の引力（殻定理、命題LXX〜LXXXIV）】

\textbf{命題LXXVII. 定理2.}

球体に対する引力は球体に比例する。すなわち、球体同士の距離が直径に対して反比例するとすれば、その引力は直径に比例する。大きな距離が減少したとき、現在の距離が等しいとすれば、引力は二倍となる。したがって、他の引力に対して三倍となる、すなわち球体に比例する。

\textbf{補題2.}

任意の距離における引力は、その距離の平方に比例する。

\textbf{補題3.}

同質の球体外にある粒子について、粒子に対する引力は球体の半径に比例する。

【証明】

まず、球体 \( S \) と \( T \) の直径がそれぞれ \( d_S \) および \( d_T \) であり、その間の距離が \( r \) であるとする。補題2により、引力は距離 \( r \) の平方に比例する。したがって、引力 \( F \) は以下のように表される：

\[ F \propto \frac{1}{r^2} \]

次に、直径 \( d_S \) と \( d_T \) の関係を考慮すると、引力 \( F \) は直径 \( d_S \) または \( d_T \) に比例する。すなわち：

\[ F \propto d_S \quad \text{または} \quad F \propto d_T \]

したがって、球体の引力は球体の直径に比例し、さらに距離 \( r \) の逆平方率に従う。

\[
F \propto \frac{d_S}{r^2}
\]

ここで、球体の半径を \( R \) とすると：

\[ F \propto \frac{R}{r^2} \]

したがって、粒子に対する引力は球体の半径に比例する。

\[
F \propto R
\]

\textbf{第12節：球体の引力（殻定理、命題LXX〜LXXXIV）}

\textbf{命題 LXXV. 理論 XXXV.}

もし与えられた球面上の各点がその他の点から等しい中心吸引力を放つとし、この力は各点からの距離の二倍比に比例して減少するとき、任意の相似な球体は中心間の距離の平方反比に比例した力で引き寄せられる。

証明：各点の吸引はその点から中心までの距離の平方に反比例しているとします。すると、与えられた球面上の各点が他の点から等しい中心吸引力を放つとき、任意の相似な球体は中心間の距離の平方反比に比例した力で引き寄せられることが示されます。

\textbf{証明：}

各点の吸引はその点から中心までの距離の平方に反比例しているとします。すなわち、点 \( P \) が中心 \( O \) からの距離 \( r \) のとき、吸引力 \( F \) は
\[
F \propto \frac{1}{r^2}
\]
と表されます。

与えられた球面上の各点が他の点から等しい中心吸引力を放つとします。この場合、任意の相似な球体は中心間の距離 \( R \) の平方反比に比例した力で引き寄せられることが示されます。

証明は以上です。

球体からの引力について、半径 \(\textbf{S}\) から中心点まで距離 \(r\) の場所にある粒子の引力は（理論 XXXIにより）、その引力が一粒子である中心点に集中した場合と等しい。すなわち、この引力は中心点に存在する一粒子が各球体の部分から等しく引き寄せられる力と等しい。

また、その粒子の引力は（理論 XXXIVにより）半径 \(r\) の平方に反比例し、これは球体からの引力とも等しい。

\[
F \propto \frac{1}{r^2}
\]

このようにして、球体からの引力は中心点に集中した場合と同様の性質を持つことが示されました。

\textbf{定理}\textit{1.} 球体の引力は、他の同種球体に対する比において同一である。すなわち、引く球体が中心から球体への距離の平方に比例する。

\textbf{証明}\textit{QED}

\textbf{補則}\textit{1.} 同じ球体を引く球体に対する引力は、引かれる球体と引く球体との中心間の距離の二乗に対して比例する。すなわち、引く球体が引かれる球体に及ぼす引力は、引かれる球体が引く球体に及ぼす引力と同じであり、両者の引力の大きさは互換的に等しい。

\textbf{補則}\textit{2.} さらに、引かれる球体もまた引力を発生する場合でも同様である。各点が他の点を引き、その力は互いに逆方向で等しく働くため、すべての引力において（法則3により）、引く点と引かれる点に対する引力の大きさは二倍となるが、その比は保たれる。

\QED

\textbf{補則 3.} 上で述べたすべてのことは、球体が中心にあり、その外側を移動する物体についても同様に成り立つ。すなわち、上記で円錐曲線周囲の物体の運動について証明したすべては、球体が中心にあり、物体が球体内を移動する場合にも適用される。

\textbf{補則 4.} 上記で円錐曲線中心周囲の物体の運動について証明したものは、物体が球体内を移動する場合にも同様に成り立つ。

\pngright{207.png}{}{}{}

\textbf{命題 LXXVI.} \textit{理論 XXXVI.}

球体の中心にある引力場において、物体が球体内を移動する場合の運動について述べる。

\begin{center}
\includegraphics[width=1761px,height=888px,viewport=-24 0 1785 912]{207.png}
\end{center}

\textbf{密度と引力：}異なる種類の物体は中心から与えられた距離に沿って円周方向に移動するとき、その引力は相似であるが、同じ種類の物体は中心から与えられた距離にあるとき、引力は減少し、これは体への引力がその距離の二乗反比例であることを示す。つまり、このような球体一つが他の球体を引き付ける力全体は、中心間の距離の平方に反比例すると主張する。

\begin{center}
\includegraphics[width=0.8\textwidth]{207.png}
\end{center}

同心相似な球殻 \(AB\), \(CD\), \(EF\) などがあり、これらは内部から外部まで materiam を合成している。

\textbf{第12節：球体の引力（殻定理、命題LXX〜LXXXIV）}

\textbf{【ラテン語テキスト】}
densiorem versus centrum, vel subduct{\ae} relinquant tenuiorem;
\& h{\ae}, per Theor.\ XXXV, trahent Sph{\ae}ras alias quotcun{\que} concentricas $GH$, $IK$, $LM$, \&c.\ singul{\ae} singulas, viribus
reciproce proportionalibus quadrato distanti{\ae} $SP$. Et componendo
vel dividendo, summa virium illarum omnium, vel excessus
aliquarum supra alias, hoc est, vis qua Sph{\ae}ra tota ex concentricis
quibuscun{\que} vel concentricarum differentiis composita
$AB$, trahit totam ex concentricis quibuscun{\que} vel concentricarum
differentiis

\textbf{【日本語LaTeX形式】}
中心に密度の大きいものより小さいものを引き下げるか、または微小量を引いたもの；これらは理論 XXXV によれば、同心的な他の球体 $GH$, $IK$, $LM$，など、それぞれ単独で、距離 $SP$ の平方に反比例する力によって引き寄せられる。そして加算法や減算法により、これらのすべての力を合計または差し引くと、全体の球体 $AB$ は、同心的な部分またはその差異からなる球体 $AB$ が作る引力で、他の同心的な部分またはその差異からなる球体を引き寄せることになる。

\textbf{訳注}
理論 XXXV は、球体間の引力が距離の平方に反比例することを述べています。


球体の複合的な引力 \( GH \) は、同じ比にあります。
球殻の数を無限大に増加させます。すなわち、密度と引力の力が円周から中心に向かって進行するとき、その比がどの程度増加しても減少しても、任意で密度が欠けた部分に非引力性の物質を追加して球殻の形態を任意の形にすることができます。そして、これらの球殻の一つが他の球殻を引きつける力は（先の議論によれば）依然としてその距離の平方に反比例する比であります。\(\blacksquare\)

\textit{補則1.} したがって、このような球体が互いに接触していれば、それぞれの引力は依然として同じ距離の平方に反比例する比であります。


\textbf{すべての物体は互いに引き寄せ合う。}相似する球体間の加速引力は、中心間の任意の距離において、その距離から定まる二次式に比例し、吸引される球体同士である。

\textit{\textbf{補則2.}}異なる中心間の距離において、相似する球体が互いに引き寄せ合う力は、中心間の距離の平方に比例して、接線方向に作用する。

\textit{\textbf{補則3.}}動機引力、すなわち任意の球体の質量は、その球体が他の球体に吸引されるとき、中心間の距離が等しい場合においても、吸引される球体と同時に作用し、それが球体内の任意の点における質量の総和に比例する。

\textbf{4.} 異なる距離における引力は，中心間の距離の平方に比例した面積によって乗算されます。

\textbf{5.} 同じ球体からの引力が互いに作用する場合，その引力は等しいです。なぜなら，双方から与えられる力により引力が倍増し，その比は保たれます。

\textbf{6.} あるいくつかの球体が他の球体周りで静止している場合，それぞれが個別に回転しており，中心間の距離が異なる場合でも，これらの球体間の引力は等しいです。


球体の回転と静止する物体の静止時直径との比が等しい；周期的な時間は等しくなる。
\textit{補則7.} そして逆もまた然りである。もし周期的な時間が等しいならば、距離は直径に比例する。

\textit{補則8.} 上記で円錐曲線囲繞体の運動について述べたすべてが、球体を引き寄せ、その形状と条件が既に説明されている場合においても、中心に位置しているときにも成り立つ。

\textit{補則9.} また、回転する球体であるならば、引き付けられる球体でもある。

条件のどの部分も既に述べられています。

\condpagelarge{命題 LXXVII. 定理 XXXVII.}

\pngright{210.png}{}{}{}

もし球体上の各点に向心力が、その点から物体までの距離に比例するもので作用するとします。このとき、二つの球体が互いに引き合う合力は、中心間の距離に比例することを主張します。

\textit{証明 1.} 考えられる球体 \( ACBD \) の中心を \( S \)、物体 \( P \) を吸引する点として \( PASB \) を通る軸を考えます。平面 \( EF \) と \( ef \) は球体を垂直に交差します。

\textbf{証明}：各点に向心力が距離に比例すると仮定した場合、物体 \( P \) から中心 \( S \) までの距離が短くなるほど、その吸引力は強くなります。二つの球体間の合力は、それぞれの吸引力を加算することで得られます。この合力は中心間の距離に比例することが示されます。

\textbf{証明}：各点に向心力が距離に比例すると仮定した場合、物体 \( P \) から中心 \( S \) までの距離が短くなるほど、その吸引力は強くなります。二つの球体間の合力は、それぞれの吸引力を加算することで得られます。この合力は中心間の距離に比例することが示されます。

\textbf{証明}：各点に向心力が距離に比例すると仮定した場合、物体 \( P \) から中心 \( S \) までの距離が短くなるほど、その吸引力は強くなります。二つの球体間の合力は、それぞれの吸引力を加算することで得られます。この合力は中心間の距離に比例することが示されます。

\textbf{証明}：各点に向心力が距離に比例すると仮定した場合、物体 \( P \) から中心 \( S \) までの距離が短くなるほど、その吸引力は強くなります。二つの球体間の合力は、それぞれの吸引力を加算することで得られます。この合力は中心間の距離に比例することが示されます。

\textbf{証明}：各点に向心力が距離に比例すると仮定した場合、物体 \( P \) から中心 \( S \) までの距離が短くなるほど、その吸引力は強くなります。二つの球体間の合力は、それぞれの吸引力を加算することで得られます。この合力は中心間の距離に比例することが示されます。

\textbf{証明}：各点に向心力が距離に比例すると仮定した場合、物体 \( P \) から中心 \( S \) までの距離が短くなるほど、その吸引力は強くなります。二つの球体間の合力は、それぞれの吸引力を加算することで得られます。この合力は中心間の距離に比例することが示されます。

\textbf{証明}：各点に向心力が距離に比例すると仮定した場合、物体 \( P \) から中心 \( S \) までの距離が短くなるほど、その吸引力は強くなります。二つの球体間の合力は、それぞれの吸引力を加算することで得られます。この合力は中心間の距離に比例することが示されます。

\textbf{証明}：各点に向心力が距離に比例すると仮定した場合、物体 \( P \) から中心 \( S \) までの距離が短くなるほど、その吸引力は強くなります。二つの球体間の合力は、それぞれの吸引力を加算することで得られます。この合力は中心間の距離に比例することが示されます。

\textbf{証明}：各点に向心力が距離に比例すると仮定した場合、物体 \( P \) から中心 \( S \) までの距離が短くなるほど、その吸引力は強くなります。二つの球体間の合力は、それぞれの吸引力を加算することで得られます。この合力は中心間の距離に比例することが示されます。

\textbf{証明}：各点に向心力が距離に比例すると仮定した場合、物体 \( P \) から中心 \( S \) までの距離が短くなるほど、その吸引力は強くなります。二つの球体間の合力は、それぞれの吸引力を加算することで得られます。この合力は中心間の距離に比例することが示されます。

\textbf{証明}：各点に向心力が距離に比例すると仮定した場合、物体 \( P \) から中心 \( S \) までの距離が短くなるほど、その吸引力は強くなります。二つの球体間の合力は、それぞれの吸引力を加算することで得られます。この合力は中心間の距離に比例することが示されます。

\textbf{証明}：各点に向心力が距離に比例すると仮定した場合、物体 \( P \) から中心 \( S \) までの距離が短くなるほど、その吸引力は強くなります。二つの球体間の合力は、それぞれの吸引力を加算することで得られます。この合力は中心間の距離に比例することが示されます。

\textbf{証明}：各点に向心力が距離に比例すると仮定した場合、物体 \( P \) から中心 \( S \) までの距離が短くなるほど、その吸引力は強くなります。二つの球体間の合力は、それぞれの吸引力を加算することで得られます。この合力は中心間の距離に比例することが示されます。

\textbf{証明}：各点に向心力が距離に比例すると仮定した場合、物体 \( P \) から中心 \( S \) までの距離が短くなるほど、その吸引力は強くなります。二つの球体間の合力は、それぞれの吸引力を加算することで得られます。この合力は中心間の距離に比例することが示されます。

\textbf{証明}：各点に向心力が距離に比例すると仮定した場合、物体 \( P \) から中心 \( S \) までの距離が短くなるほど、その吸引力は強くなります。二つの球体間の合力は、それぞれの吸引力を加算することで得られます。この合力は中心間の距離に比例することが示されます。

\textbf{証明}：各点に向心力が距離に比例すると仮定した場合、物体 \( P \) から中心 \( S \) までの距離が短くなるほど、その吸引力は強くなります。二つの球体間の合力は、

\textbf{そして、球体から中心までの距離が等しい場合、}点 \( Gg \) は平面と軸の交点であり、点 \( H \) は平面 \( EF \) 上の任意の点です。


\par\begin{center}
    \centering
    \includegraphics[width=0.8\textwidth]{209.png}

\end{center}\par

点 \( H \) から球体中心 \( S \) 方向への向心力は、距離 \( PH \) に比例し、従って線分 \( PG \) または中心 \( S \) 方向に指向する長さ \( PG \) に比例します。したがって、平面 \( EF \) 上のすべての点（すなわち全体の力）が球体中心 \( S \) 方向に引き寄せられる力を表すものは、各点を距離 \( PG \) で掛けると得られます。つまり、平面 \( EF \) 下およびその距離 \( PG \) の積です。

同様に、平面 \( ef \) の力も同じ論理により推定されます。

球体の中心 $S$ に向かって粒子 $P$ が引き寄せられる力は、その距離 $Pg$ に比例する平面 $Pg$ を引いた平面の面積と等しい。すなわち、平面 $EF$ が距離 $Pg$ に引かれたときの面積と等しくなり、両方の平面の総力を考慮すると、平面 $EF$ が中心点 $PS$ の二倍の距離 $2PS$ に引かれたときの面積と等しい。また、同様の論証により、球体全体に対するすべての平面の力は、各点から中心までの距離が異なるため、これらの力を一方向から他方向へ均等に加えると、その総和は以下のようになる。

すなわち、球体全体に対するすべての平面の力は、各点から中心までの距離 $PS$ の二倍である点 $2PS$ における平面 $EF$ の面積と比例する。また、これは等価な平面 $EF$ および $ef$ の総和が同じ距離 $PS$ に引かれたときの面積と等しい。

以上を類似した論証により、球体全体に対するすべての平面の力は、各点から中心までの距離が異なるため、これらの力を一方向から他方向へ均等に加えると、その総和は以下のようになる。すなわち、

\[
\textbf{球体全体に対するすべての平面の力} = \textbf{平面} EF + ef \textbf{が距離} PS \textbf{に引かれたときの面積}
\]

となります。

\[
\textbf{平面を $PS$ に引いた総和は，球体全体が自らの中心 $S$ から物体 $P$ まで引かれた距離として表されます。}
\]
証明終了。\QED

\textit{\textbf{案2.}} 物体 $P$ を球体 $ACBD$ 内に引き込むとします。同様の議論により，その球体が引き込まれる力は距離 $PS$ に比例することが証明されます。\QED

\textit{\textbf{案3.}} 現在，球体を無数の物体 $P$ で構成するものとします；そして，各物体が引き込まれる力は，その物体から最初の球体の中心まで引かれた距離に比例することが知られています。したがって，これは全体として物体から発生することと同じです。\QED


球体の引力（殻定理）

唯一の中心点 \( \textbf{S}_1 \) に存在し、各粒子が球体 \( \textbf{S}_2 \) 内部全体に及ぼす引力は、その球体 \( \textbf{S}_2 \) 全体に対する引力と等しい。すなわち、球体 \( \textbf{S}_2 \) が各粒子から受ける引力は、中心点 \( \textbf{S}_1 \) からの距離に比例するものである。証明終了。\QED

\textit{4. 案例} 球体が互いに引き合うとき、その引力は比例定数を保つ。

\textit{場合5.}\wsp{}次に、微小体 $p$ を球 $ACBD$ の内部に置く。平面 $ef$ が微小体に及ぼす力は、その平面と距離 $pg$ との積に比例し、これと逆向きに働く平面 $EF$ の力は、その平面と距離 $pG$ との積に比例する。したがって両者を合わせた力は、これらの積の差に比例する。すなわち、等しい平面の和に、距離の差の半分を掛けたものに比例し、いいかえれば、その和に、微小体から球の中心までの距離 $pS$ を掛けたものに比例する。同様の論法により、球全体におけるすべての平面 $EF$, $ef$ の引力、すなわち球全体の引力は、すべての平面の和、つまり球全体に、微小体から球の中心までの距離 $pS$ を掛けたものに比例する。\QED

\textit{場合6.}\wsp{}また、数えきれないほどの微小体 $p$ から成る新しい球が、前の球 $ACBD$ の内部にあるとしても、前と同じように証明できる。すなわち、一方の球が他方の球に及ぼす引力も、両者が互いに及ぼし合う引力も、中心間の距離 $pS$ に比例する。\QED


\textbf{命題 LXXVIII.} \textbf{定理 XXXVIII.}

\textit{\textbf{もし球体が中心から円周へ移動する過程において異質で不均一であるとし、その円周を通過する過程においては特定の}}\wsp{}値に達すると仮定されるとき、

証明終了。\QED

omnem a centro distantiam sint undi\que similares; \& vis attractiva puncti cujus\que sit ut distantia corporis attracti: dico quod vis tota qua hujusmodi Sphaer\ae duae se mutuo trahunt sit proportionalis distanti\ae inter centra Sphaerarum.

前の命題から、命題LXXVが命題LXXVから示されたのと同じ仕方で示される。

\textit{Corol.} Quae superius in Propositionibus X. \& LXIV. de motu corporum circa centra Conicarum Sectionum demonstrata sunt,

\textbf{第1編 第12節：球体の引力（殻定理、命題LXX〜LXXXIV）}


すべての引力が球体の物体間で発生し、その条件がすでに述べられた場合、これらの引力は同じ条件を持つ球体を吸引する。

\textit{\textbf{補遺。}}

これまでに二つの重要な引力の例を説明した。すなわち、中心指向力が距離の平方反比律で減少するか、または単純な距離比例で増加する場合である。これらの引力は球体が円錐曲線内を回転させることを可能にし、また球体の中心から遠ざかるとき、あるいは近づくときにその力が同じ法則により減少または増加することを意味する。

\textbf{訳注：}

- \textit{Corporum Sph{\ae}ricorum} は「球体の物体」を指します。
- \textit{Conicis Sectionibus} は「円錐曲線」を指します。


球体の引力に関する他の場合については、その結論がより優雅ではないため、逐一述べることは長く、また無駄である。私は一般的な方法ですべてを一度に把握し、決定することを選ぶべきである。以下のように進める。

\textbf{命題 XXIX.}

中心 \( S \) から円 \( AEB \) を描き、中心 \( P \) から円 \( EF \), \( ef \) を描き、これらの円は前記の円 \( AEB \) と交わる点 \( E \), \( e \) で交わり、直線 \( PS \) と交わる点 \( F \), \( f \) で交わる。さらに、垂線 \( ED \), \( ed \) を直線 \( PS \) 上から下ろす。私は言います：もし弧の距離 \( EF \), \( ef \) の差が無限に小さくなると理解されるなら、消滅する線分 \( Dd \) に対する最終的な比は

\[
\lim_{EF, ef \to 0} \frac{Dd}{\text{一定の線}}
\]

となります。

\[
\text{もし線分}\, Pe\,\text{が弧}\, EF\,\text{を点}\, q\,\text{で交わり、直線}\, Ee\,\text{が弧}\, Ee\,\text{が消滅するときと一致し、}
\]
\[
\text{直線}\, Pe\,\text{の延長が直線}\, PS\,\text{と点}\, T\,\text{で交わり、点}\, S\,\text{から}\, PE\,\text{に垂線}\, SG\,\text{を引き、}
\]
\[
\triangle EDT, \, edt, \, EDS\,\text{が相似であるので、}\, Dd : Ee = DT : ET\, \text{または}\, DE : ES.
\]
\[
\text{\( \triangle Eqe, \, ESG\) が相似（補題 VIII 及び Lem. VII の補則 3）なので、}
\]
\[
Ee : qe = Ff : SG,
\]
\[
\text{よって}\, Dd : Ff = DE : SG.
\]
\[
\text{これは（相似三角形の性質により）}
\]



\textbf{命題LXXIX.} \textbf{定理XXXIX.}

\textit{もし平面$EFfe$が無限に薄く、すなわち$EFe$の幅が消失するにつれて円錐回転体を描くと、この体は球面状の凸凹形状$S$を作り出し、その粒子は等しい力を中心的に引き寄せ合うとする。このとき、体$DEq \times Ff$に比例し、かつ粒子が位置$Ff$において同じ物体を引き付ける力に比例する複合的な比によって、物体が点$P$に引かれる力は定義される。}


\par\begin{center}
  \centering
  \includegraphics[width=1820px,height=1150px,keepaspectratio]{215.png}

\end{center}\par

\[
\text{次に、球面} FE \text{の表面を考えます。この表面は弧} FE \text{によって生成され、任意の点} d \text{で線} de \text{が} r \text{に切断されます。}
\]
\[
\text{すると、弧} rE \text{によって生成された円環部分が得られ、これは} Dd \text{のような小線分を保持します。（この点は} PE \text{である球の半径で、アキメデスが} \textit{De Sphaera et Cylindro} \text{の中で証明しました。)}
\]
\[
\text{そして、この力は} PE \text{または} Pr \text{上の線に沿って} \textbf{球面の表面} \text{上での円錐面上にある各点で作用し、これは同じ円環部分と等価です。つまり、} Dd \text{そのものや、これと同等な} PE \text{の半径} \textbf{球体} \text{上の与えられた直角三角形}
\]

\textbf{その線 $Dd$: そして、直線 $PS$ に中心 $S$ に向かって伸びる線分 $Dd$ の比は $PD$ に対する $PE$ の比と等しくなるので}

\[
\frac{Dd}{PS} = \frac{PD}{PE}
\]

\textbf{ここで、線分 $DF$ を無数の等しい部分に分割し、各部分を $Dd$ と名付けます。また、面積 $FE$ も同様に無数の等しい円帯状部分に分割します。これらの力はすべての $PD \times Dd$ の和と等しくなります。}

\textbf{すなわち、すべての線分 $Dd$ が互いに等しいと仮定し、与えられた値をもとに各 $PD \times Dd$ の合計を計算できるので}

\[
\frac{1}{2}PFq - \frac{1}{2}PDq = \frac{1}{2}PEq - \frac{1}{2}PDq = \frac{1}{2}DEq
\]

ここで、$P$, $D$, $E$, $F$ は点を表し、$q$ は無限小量を示します。

在 \(Dd\)；すなわち、もし \(\frac{1}{2}Dd\) を無視すれば、\(DE\) は平方となります。現在、面 \(FE\) を高さ \(Ff\) まで引きます；そして、体 \(EFfe\) の作用する力は点 \(P\) の粒子に対して \(DEq. \times Ff\) となるでしょう。仮定します：あるデータの粒子 \(Ff\) 在位置 \(PF\) から点 \(P\) に作用する力を与えます。しかし、その力が与えられない場合、体 \(EFfe\) の作用する力は体 \(DEq. \times Ff\) と同様であり、その無指定の力とは別に存在します。\QED

\condpagelarge{命題 LXXX. 定理 XL.}

\textbf{もし}球面 \(AEB\) を中心 \(S\) で描くとき、等しい粒子を各点に置くと、







ニュートンは次のように述べています：

$$\text{向心力が等しい力をもって} \, AB \, \text{軸上に位置する任意の一点} \, P \, \text{から} \, D \, \text{点を基点として} \, DE \, \text{という直線を引くと、この直線は球体} \, S \, \text{との交点} \, E \, \text{で垂直に交差します。}$$
$$\text{さらに、} \, D \, \text{から} \, N \, \text{までの距離} \, DN \, \text{を求めるために、} \, DE \times PS \, \text{と} \, PE \, \text{の積} \, \frac{DEq. \times PS}{PE} \, \text{が} \, S \, \text{上の一点} \, P \, \text{から} \, PE \, \text{の距離} \, PE \, \text{に及ぼす力と比例するという条件を用います。}$$
$\text{このとき、点} \, P \, \text{が} \, S \, \text{に対する引力によって引き寄せられる総力を} \, AB \, \text{軸と曲線} \, ANB \, \text{の下部に囲まれた面積} \, \textbf{ANB} \, \text{と比例すると主張します。}$

ここで、\textbf{ANB}はニュートンが指す曲線部分を強調表示しています。

\textbf{すなわち、命題と最後の推論において構築された状態を保つとして考えます。円球体の軸$AB$を無数に等しい部分$Dd$に分割し、円球全体を同数の凹凸面を持つ円錐曲線板$EFfe$で分割します。そして直線$dn$を垂直に引きます。}

\textbf{先述の命題により、板$EFfe$が粒子$P$を引き込む力は$DEq \times Ff$と等しく、またこれは距離$PE$または$PF$における各部分の力を同時に作用させたものと等しいです。}

\textbf{さらに、最後の推論により、$Dd$は$Ff$に対して$PE$は$PS$に対して等しい比を成し、従って$Ff$は等しくなります。}

\[
\frac{PS \times Dd}{PE};
\] そして \( DEq. \times Ff \) は \( Dd \) に等しくなり、すなわち
\[
DEq. \times PS
\]
を \( PE \) で割った値と等しい。したがって、板 \( EFfe \) の引力は
\[
Dd \times \frac{DEq. \times PS}{PE}
\] に比例し、また個々の粒子に対して \( PF \) の距離に沿って作用する引力の総和は（仮定より）\( DN \times Dd \)，すなわち消滅する面積 \( DNnd \) に比例します。したがって、すべての板の引力は物体 \( P \) に対して
\[
DNnd
\] の総和と等しくなります。すなわち球体の全体的な引力は面積 \( ABNA \) に比例します。\(\blacksquare\)

\textit{Corol.\,1.} したがって、中心への引力が各粒子に対して常に一定である場合、
\[
DN
\] は一定になります。

\[
\frac{DEq. \times PS}{PE}
\]
これは球体から粒子を引きつける全引力と等しい。

\textbf{補則2.} 特定の粒子間の引力は粒子同士の距離に反比例するものとする。このとき、
\[
DN = \frac{DEq. \times PS}{PEq.}
\]
により得られる引力は粒子 \(P\) が球体全体から引きつけられる引力と等しい。

\textbf{補則3.} 特定の粒子間の引力は粒子同士の距離の立方に反比例するものとする。このとき、
\[
DN = \frac{DEq. \times PS}{PE^{\text{qq.}}}
\]
により得られる引力は粒子が球体全体から引きつけられる引力と等しい。

\textbf{命題 LXXXI. 問題 XLI.}

既に前提を設定したならば、面積 \( ABNA \) を測定しなければならない。


\par\begin{center}
    \centering
    \includegraphics[width=1820px,height=1150px,trim=-24 0 0 0,clip]{215.png}

\end{center}\par

点 \( P \) を通って球面に接する直線 \( PH \) を引き、軸 \( PAB \) に対して正規線 \( HI \) を下ろし、\( PI \) の中点を \( L \) とする。すると、

\[
\textbf{証明：} \quad \text{点 } H \text{ から球面に接する直線 } PH \text{ は, 点 } P \text{ から球の中心までの距離と等しい.}
\]

したがって、\( PI = 2PL \) であり、

\[
\frac{DEq. \times PS}{PE \times V} = DN
\]

ここで \( DEq. \) は点 \( D \) から弧 \( EQ \) までの距離を表し、\( PS \) は点 \( P \) から点 \( S \) の距離、\( PE \) は点 \( P \) から点 \( E \) の距離、\( V \) は引力の比例定数である。よって、

\[
DN = \frac{DEq. \times PS}{PE \times V}
\]

ここで \( DN \) は面積 \( ABNA \) を表す。したがって、球体全体に対する粒子の吸引力は面積 \( ABNA \) に比例する。

以上より、


\textbf{結論：} \quad 粒子が球体全体から引き寄せられる力は, 面積  ABNA \text{ に比例する。
}

\textbf{訳文}

\textit{（第12命題、第2編、要素論）} \(PE^2\) は \(PS^2 + SE^2 + 2PS \cdot SD\) に等しい。また、\(SE^2\) または \(SH^2\) （三角形 \(SPH\) と \(SHI\) の相似性から）は直角三角形 \(PSI\) の面積に等しい。したがって、\(PE^2\) は \(PS\) および \(PS + SI + 2SD\) の下部の面積に等しく、これは \(PS\) および \(2LS + 2SD\) の下部の面積に等しく、さらにこれは \(PS\) および \(2LD\) の下部の面積に等しい。

また、\(DE^2\) は \(SE^2 - SD^2\) または \(SE^2 - SL^2 + 2SL \cdot LD - LD^2\) に等しく、これは \(SLD - LD^2 - ALB\) に等しい。なぜなら、\(LS^2 - SE^2\) または \(LS^2 - SA^2\) （第2編の第6命題による）は直角三角形 \(ALB\) の面積に等しいからである。

したがって、
\[2SLD - LD^2 - ALB\]
と表すことにする。

\textbf{DEq.} および量
\[
\frac{DEq. \times PS}{PE \times V}
\]
は、命題前段第4条の補遺により、長さ \( DN \) を順次適用すると、以下の3つの部分に分解される：
\[
\frac{2SLD \times PS}{PE \times V} - \frac{LDq. \times PS}{PE \times V} - \frac{ALB \times PS}{PE \times V}
\]
ここで、もし \( V \) を逆比とし、\( PE \) を \( PS \) および \( 2LD \) の中比例とするとき、これらの3つの部分は順次適用される曲線の長さに相当し、その面積は一般的な方法で明らかとなる。証明終了。\(\QEFup\)


球体の引力（殻定理）

1. 例。もし中心力が各球体粒子に対して互換的に半径距離に比例するものとするとき、$V$ を半径距離 $PE$ と置き換え、$2PS \times LD$ を $PEq.$ と置き換え、すると $DN$ は
\[
SL - \frac{1}{2}LD - \frac{ALB}{2LD}
\]
となる。

1. ここで $DN$ を $2SL - LD - \frac{ALB}{LD}$ と等しく設定し、$2SL$ の部分を長さ $AB$ 方向に引くと直角形の面積 $2SL \times AB$ が描かれる。また、$LD$ の無限小部分を正規方向に引き、運動継続の法則に従って
\[
\textbf{動径} \cdot \textbf{速度}
\]
と増加または減少するとき、面積 $2SL \times AB$ は一定となる。

この証明は、運動の連続性と中心力の性質を基にしています。

\[
\text{令}\, LD \, \text{の長さに対する面積}\, \frac{LBq. - LAq.}{2}，すなわち面積
\]

\begin{center}
\includegraphics[width=823px,height=1021px,offset=-12pt]{216.png}
\end{center}

\[
SL \times AB;
\]
これを前の面積から引くと、残る面積は
\[
SL \times AB.
\]

また、第三の部分 $\frac{ALB}{LD}$ を局所的な運動により同方向に引き延ばすと、双曲線の面積を描くことになります。この面積を前の面積 $SL \times AB$ から引くと、求めたい面積 $ABNA$ が得られます。

したがって、このような構成は問題の解法を導きます。点 $L$, $A$, $B$ に対して直線 $Ll$, $Aa$, $Bb$ を垂直に引き、これらの直線はそれぞれ点 $L$, $A$, $B$ から面積の計算に使用されます。

\textbf{第1編 第12節：球体の引力（殻定理、命題LXX〜LXXXIV）}

\textbf{【50】} ここで、点 \(a\) と \(b\) を通る双曲線 \(ab\) を直線 \(Ll\), \(LB\) の漸近線を用いて描く。そして、弦 \(ba\) によって領域 \(aba\) の面積が求められるものとする。

\textbf{命題 LXXV}：点 \(A\) と \(B\) の領域 \(ABNA\) の面積は、求めるべき面積である。


\par\begin{center}
    \centering
    \includegraphics[width=8.86cm,height=9.51cm,trim={0 -24 0 0},clip]{217.png}

\end{center}\par

\textit{例題 2}：もし中心力が各球体部分に向かう力を、半径の立方反比例とし、あるいは（等価であるため）任意の平面に接した直線から半径までの距離の立方反比例とするとき、\( \frac{PE^{\text{立方}}}{2AS^2} \) を \( V \) とする。次に、\( 2PS \times LD \) を \( PE^2 \) とすると、

\[
DN \propto
\]

\[
\frac{SL \times ASq.}{PS \times LD} - \frac{ASq.}{2PS} - \frac{ALB \times ASq.}{2PS \times LDq.}
\]
これは、比例する量 $PS$, $AS$, $SI$ について
\[
\frac{LSI}{LD} - \frac{1}{2}SI - \frac{ALB \times SI}{2LDq.}
\]
と等しい。

これらの部分を引くと、

\[
3 \textbf{ 部分は} AB の長さに相当し、
\]

\[
\textbf{第1部分} \frac{LSI}{LD} \textbf{は双曲線の面積} を生成します；
\]


\textbf{第2部分} $\frac{1}{2}SI$ \textbf{は半径} $AB \times SI$ の半分の面積を生成します；



\[
\textbf{第3部分} \frac{ALB \times SI}{2LDq.} \textbf{は} ALB \times SI から 2LA と 2LB の差を除いた面積 を生成します、
\]

\[
つまり \frac{1}{2}AB \times SI。
\]

\textbf{第1部分から第3部分を引くと}


\par\begin{center}
    \centering
    \pngright{217.png}{}{}{}

\end{center}\par

\[
\textbf{第12節：球体の引力（殻定理、命題LXX〜LXXXIV）}
\]


\[
\textbf{命題 LXXVII}
\]
球体の引力を求めるための面積を計算する。

次に、図形 \(ABNA\) の面積を求めます。これは二次曲線と三次曲線の面積の和から得られます。
\[
\textbf{証明}
\]
二次曲線と三次曲線の面積の差が残る面積 \(ABNA\) となります。

次に、問題の構築方法を示します。点 \(L, A, S, B\) を通じて直線 \(Ll, Aa, Ss, Bb\) を引き、これらの中から \(Ss\) が直線 \(SI\) と等しい長さを持つようにします。次に、点 \(s\) を通り、直線 \(Ll, LB\) の漸近線を描き、直線 \(Aa, Bb\) との交点 \(a\) および \(b\) を求めることで、双曲線 \(asb\) を得ます。この双曲線の面積から直角三角形 \(2ASI\) を引くと、求める面積 \(ABNA\) が残ります。

\textit{例3}
もし中心指向力が各球体粒子に及ぶ力を、その粒子からの距離の四倍反比例するとする場合、

\[
\textbf{証明終了}
\]

$$\frac{PE^4}{2AS^3}$$
を$V$とし、$\surd 2PS \times LD$ を$PE$とする。従って$$DN = \frac{SL \times SI^{\frac{3}{2}}}{\surd 2 \times LD^{\frac{3}{2}}} - \frac{SI^{\frac{3}{2}}}{2\surd 2 \times LD^{\frac{1}{2}}} - \frac{ALB \times SI^{\frac{3}{2}}}{2\surd 2 \times LD^{\frac{5}{2}}}$$となる。これらの三つの部分を$AB$の長さに引くと、面積が等しい四つあることがわかる。すなわち、

\[
\frac{\surd 2 \times SL \times SI^{\frac{3}{2}}}{LA^{\frac{1}{2}}} - \frac{\surd 2 \times SL \times SI^{\frac{3}{2}}}{LB^{\frac{1}{2}}},
\]

\[
\frac{LB^{\frac{1}{2}} \times SI^{\frac{3}{2}} - LA^{\frac{1}{2}} - SI^{\frac{3}{2}}}{\surd 2}
\]

\[
\frac{ALB \times SI^{\frac{3}{2}}}{3\surd 2 \times LA^{\frac{3}{2}}} - \frac{ALB \times SI^{\frac{3}{2}}}{3\surd 2 \times LB^{\frac{3}{2}}}.
\]
そして、後者を前者から引いたとき、

\[
\frac{8SI \textbf{立方}}{3LI}
\]

を得る。したがって、球体の中心に引き寄せられる粒子 \( P \) の全引力は

\[
\frac{SI \textbf{立方}}{PI},
\]

すなわち \( PS \textbf{立方} \times PI \) と反比例する。

同様の方法で、球体内にある粒子に対する吸引力も決定できるが、次に示す定理によりより簡便に行うことができる。

\textbf{命題 LXXXII. 定理 XLI.}

\textit{中心 $S$、半径 $SA$ で描いた球において、$SI$, $SA$, $SP$ を連比例にとる。このとき、球の内部の任意の位置 $I$ にある微小体の引力は、球の外部の位置 $P$ にある微小体の引力に対して、中心からの距離 $IS$ と $PS$ の比の平方根と、それらの位置 $I$, $P$ で中心に向かう向心力の比の平方根とから合成された比をなす。}

\par\begin{center}
    \centering
    \includegraphics[width=1820px,height=1150px,keepaspectratio]{215.png}

\end{center}\par
たとえば、球の粒子の向心力が、引かれる微小体からの距離に反比例するとしよう。このとき、$I$ にある微小体が球全体から引かれる力は、$P$ で引かれる力に対して、距離 $SI$ 対 $SP$ の比の平方根と、$I$ における向心力（中心にある一つの粒子から生じるもの）対 $P$ における向心力（同じ中心の粒子から生じるもの）の比の平方根とから合成された比をなす。後の比の平方根は、距離 $SI$, $SP$ の逆比の平方根である。この二つの平方根の比を合成すると、等しさの比になる。したがって、球全体が $I$ と $P$ に及ぼす引力は等しい。

同様の計算により、球の粒子の力が距離の二乗に反比例するなら、$I$ での引力は $P$ での引力に対して、距離 $SP$ 対球の半径 $SA$ の比になることが導かれる。もしそれらの力が距離の三乗に反比例するなら、$I$ と $P$ での引力は互いに $SP^2$ 対 $SA^2$ の比になる。四乗に反比例するなら、$SP^3$ 対 $SA^3$ の比になる。したがって、この最後の場合に $P$ での引力が $PS^3 \times PI$ に反比例することが見いだされていたから、$I$ での引力は $SA^3 \times PI$ に反比例する。すなわち（$SA^3$ は一定だから）$PI$ に反比例する。以下、同様に限りなく進む。

定理は次のように証明される。前に作図したままの状態で、物体が任意の位置 $P$ にあるとき、縦線 $DN$ は
\[ \frac{DE^2 \times PS}{PE \times V} \]
に比例することが見いだされた。したがって、$IE$ を引けば、もう一つの位置 $I$ に対する縦線は、文字を置きかえて
\[ \frac{DE^2 \times IS}{IE \times V} \]
に比例する。

球の任意の点 $E$ から生じる向心力が、距離 $IE$, $PE$ において互いに $PE^n$ 対 $IE^n$ の比をなすとする（$n$ は $PE$, $IE$ のべき指数を表す）。すると、これらの縦線は
\[ \frac{DE^2 \times PS}{PE \times PE^n} \quad\text{と}\quad \frac{DE^2 \times IS}{IE \times IE^n} \]
に比例し、その比は $PS \times IE \times IE^n$ 対 $IS \times PE \times PE^n$ である。相似な三角形 $SPE$, $SEI$ から、$IE$ 対 $PE$ は $IS$ 対 $SE$（すなわち $SA$）に等しい。そこで $IE$ 対 $PE$ の比のかわりに $IS$ 対 $SA$ の比を書くと、縦線の比は $PS \times IE^n$ 対 $SA \times PE^n$ となる。ところが、$PS$ 対 $SA$ は距離 $PS$, $SI$ の比の平方根であり、$IE^n$ 対 $PE^n$ は距離 $PS$, $IS$ における力の比の平方根である。ゆえに縦線、したがってそれらが描く面積、そしてそれらに比例する引力は、これらの二つの平方根の比から合成された比をなす。\QED
\textbf{第12節：球体の引力（殻定理、命題LXX〜LXXXIV）}

\textbf{命題 LXXXIII. 問題 XLII.}

中心に位置する粒子が球体のどの部分を引き寄せられるかの力を見つける。

球体 \(P\) の中心に置かれた粒子があり、この球体の平面 \(RDS\) によって切り取られた部分 \(RBSD\) がある。この平面 \(RDS\) は球面 \(RBS\) を含み、また中心 \(P\) から球面 \(EFG\) を描く。

\textbf{証明}

粒子が球体の表面 \(RBS\) 上の任意の点 \(S\) に引き寄せられる力は、その点 \(S\) からの直線 \(PS\) と平面 \(RDS\) の交点 \(I\) との距離 \(SI\) に比例する。したがって、粒子が球体の表面 \(RBS\) 上の任意の点 \(S\) に引き寄せられる力は、その点 \(S\) からの直線 \(PS\) と平面 \(RDS\) の交点 \(I\) との距離 \(SI\) に比例する。これらの距離 \(SI\) は、各点 \(S\) から中心 \(P\) までの半径の半分と、平面 \(RDS\) に対する角度の正弦との比に比例する。

したがって、粒子が球体の表面 \(RBS\) 上の任意の点 \(S\) に引き寄せられる力は、各点 \(S\) から中心 \(P\) までの半径の半分と、平面 \(RDS\) に対する角度の正弦との比に比例する。これらの比は合成比となる。

\(\blacksquare\)

\begin{center}
\includegraphics[width=902px,height=1280px,offset=-12pt]{220.png}
\end{center}

次に、洞窟 \(DB\) を点 \(F\) において引き、領域を部分 \(BREFGS\)，\(FEDG\) と分ける。

しかしながら、その表面は純粋な数学的なものではなく、物理学的なものとして考えられ、最小限の深さを持つとする。この深さを \(O\) と名付け、アキメデスによって証明された通り、表面 \(PF \times DF \times O\) となる。

さらに、各球体の引力は、その距離の指数 \(n\) の位相に反比例すると仮定する。すなわち、点 \(P\) を引きつける力 \(FE\) の引力は
\[
\frac{DF \times O}{PF^{n - 1}}
\]
となるとする。

\[
\textbf{第12節：球体の引力（殻定理、命題LXX〜LXXXIV）}
\]

\textbf{命題 LXXXIII}

球体 \( S \) の中心から外側に位置する点 \( P \) に作用する力は、その点 \( P \) が所在する円錐曲線の殻 \( RBSD \) を引き出す力と等しい。このとき、垂線 \( FN \) が点 \( O \) まで引かれて、点 \( P \) から垂線 \( FN \) に沿って動径 \( DB \) の長さを順次適用した曲線面積 \( BDLIB \) は、球体 \( S \) の中心 \( S \) から点 \( P \) まで引かれた力 \( v \) と比例する。すなわち、
\[
\textbf{命題 LXXXIV}
\]

\textit{任意の円錐曲線殻 \( EBK \) の軸上に位置する点 \( P \) から外側に引かれた力を求めること。}

球体 \( S \) の中心 \( S \) から外側に位置する点 \( P \) に作用する力は、円錐曲線殻 \( RBSD \) を引き出す力と等しい。このとき、垂線 \( FN \) が点 \( O \) まで引かれて、点 \( P \) から垂線 \( FN \) に沿って動径 \( DB \) の長さを順次適用した曲線面積 \( BDLIB \) は、球体 \( S \) の中心 \( S \) から点 \( P \) まで引かれた力 \( v \) と比例する。すなわち，
\[
v : \text{面積 } BDLIB = \text{定数}
\]

\textit{証明終了}

\begin{center}
\textbf{命題LXXIII}

球体の表面 $EFK$ が円弧 $EBKF$ と円弧 $EFKD$ に分割されるとき、部分 $EBKF$ の引力を命題LXXXIで求め、部分 $EFKD$ の引力を命題LXXXIIIで求め、その和が全体の段体 $EBKD$ の引力となる。
\end{center}

\textit{証明。} 球体の表面 $EFK$ が円弧 $EBKF$ と円弧 $EFKD$ に分割されるとき、部分 $EBKF$ の引力は命題LXXXIで求め、部分 $EFKD$ の引力は命題LXXXIIIで求め、その和が全体の段体 $EBKD$ の引力となる。

\textit{補遺。}

球形物体の引力について説明した後、他の物体の引力法則に進むべき時が来た。しかし個々の吸引成分を逐次扱うことは本稿の目的には適さない。必要なのは以下の命題である。



球体の引力についてのより一般的な法則や、それらから生じる運動に関するものについては、哲学的な問題において一定の用に立つため、追加することとする。

 % 例として画像を配置します。実際の使用では適切な画像ファイル名を使用してください。


\section{第1編 第13節：非球形物体の引力（命題LXXXV〜XCIII）}


% Section: 第1編 第13節：非球形物体の引力（命題LXXXV〜XCIII）



\sectpage{XIII.}

\begin{center}\textbf{非球形物体の引力について。}\end{center}

\propnopage{命題LXXXV. 定理XLII.}

\textbf{もし引き起こされる物体において、引き起こす物体と接触している場合に比べて、引き起こされる物体と引き起こす物体が僅かな距離で離れているときの引力がはるかに強いとすれば、引き起こす粒子による引く力は、物体が引き寄せられるときの距離から比例するよりはるかに速い速度で減少する。}

なぜなら、もし引く力が粒子からの距離の二乗比で減少するとすれば、球形物体に対する引力は、

\begin{center}\textit{その理由は}\end{center}

\textbf{（命題LXXXV）}球体の中心から吸引される物体との距離の平方に反比例する引力は、接触による影響で著しく増加することはありません。また、吸引が物体の移動方向に向かって減少する比率よりも小さい比率で減少する場合、接触によりさらに少ない増加です。したがって、球体を吸引する物質に関する命題が明らかになります。そして、内側に物体が構成された球状の天体は、外側の物体を引きつける力についても同じ比率を持ちます。さらに、内部に物体が存在する球状の天体は、その空洞を通じて吸引される場合、より大きな引力を持つことが確認されます。

\begin{center}

\end{center}

\textbf{命題LXXXVI（理論XLIII）}

異なる物体間の引力は（命題LXXで示したように）相殺されるため、接触点においては引力が存在しない。もし球体とこれらの楕円体との接触部分からそれぞれ球的部局が取り除かれて新たに楕円的部局が追加されると、これらの引力物体の形状は自由度により変化するが、これらが接触点から遠ざかった部分については、接触による引力過剰が著しく増大することが可能である。したがって、すべての形状を持つ物体に関する命題が確定される。

\textbf{証明}

異なる物体間の引力は相殺される（命題LXX参照）。よって、接触点においては引力が存在しない。

次に、球体と楕円体との接触部分からそれぞれ球的部局が取り除かれ、新たに楕円的部局が追加された場合を考える。これらの物体の形状は自由度により変化するが、これらが接触点から遠ざかった部分については、接触による引力過剰が著しく増大することが可能である。

したがって、すべての形状を持つ物体に関する命題が確定される。$\blacksquare$


\par\begin{center}
  \centering
  \includegraphics[width=0.8\textwidth]{fig13-3.png}
  \label{fig:contact}

\end{center}\par

もし，物体を構成する各部分からなる引力体において，引かれる物体に対する力が物体からの距離に従って三倍以上速い速度で減少する場合、接触時に引力は非常に強くなります。引く側と引かれる側の間の距離が最小であるときより強くなるのです。

なぜなら，ラテックス球に対し近づく粒子の引力が無限大になることを示すのは，第41問題の例2および3による解釈から明らかです。また，これらの例と第41定理を比較すると，これは容易に理解できます。


非球形物体に対する引力について、円凸凹体に対して物体が吸引される場合や、物体がその内部に配置される場合を論じます。また、これらの球体や円錐に対し、接触点以外の任意の場所に吸引可能な物質を加えたり取り除いたりすることで、吸引する物体が任意の形状を持つ場合でも、命題全体はすべての物体について成り立つことが確認されます。\QED

\textbf{命題 LXXXVII. 定理 XLIV.}

\textit{もし二つの物体が互いに似ていて、かつ材料が等しく吸引可能な性質を持つとし、それぞれ自身に対して比例する粒子を引き合うとする。}

\textbf{そして、類似の位置に置かれた場合：物体全体に対する粒子の加速度引力は、該当する粒子部分に対する粒子の加速度引力と同様に、全体に対して比例し、同じように配置されます。}

なぜなら、物体を粒子とみなすと、これらの粒子は全体に対して比例しており、また同じように配置されるので、任意の一つの物体における粒子に対する引力は、他の物体における該当する粒子に対する引力と同等であり、さらに個々の粒子部分に対する引力は、第一物体の各粒子部分に対する引力と同等であるから、加算すると、第一物体全体に対する引力は\ldots{}

\textbf{命題 LXXXV. 1.}

したがって、引き力の個々の力を引く物体とその距離を増加させるとき、その力は反比例する。 \QED

\textit{Corol.\,1.} より明らかに、引き力の個々の力を引く物体とその距離を増加させるとき、全体の物体に対する加速引力は、各物体の質量とそれぞれの距離の逆数の積に比例する。すなわち、引き力の個々の力を引く物体との距離が二乗で減少するとき、全体の物体に対する加速引力は各物体の質量とそれぞれの距離の二乗の逆数の積に比例する。例えば、引く物体との距離が二乗で減少する個々の物体の力があるとき、全体の物体は $A \text{立方}$ と $B \text{立方}$ の比にあり、また、全体の物体の辺の長さと引き力の個々の力を引く物体に対する各物体の距離の二乗の逆数との積も同様である。

非球形物体の引力について、距離から $A$ と $B$ の関係が与えられたとき、加速引力はそれぞれの物体に対して以下のようになります。

$A \& B : \text{加速引力} = \frac{A \, \text{立方}}{A \, \text{平方}} \quad \& \quad \frac{B \, \text{立方}}{B \, \text{平方}}$

すなわち、物体の側面である $A$ と $B$ の立方比です。もし個々の力が吸引される粒子からの距離の立方に比例して減少する場合、全体的な加速引力は以下のようになります。

$A \& B : \text{加速引力} = \frac{A \, \text{立方}}{A \, \text{立方}} \quad \& \quad \frac{B \, \text{立方}}{B \, \text{立方}}$

すなわち、等しいです。もし個々の力が吸引される粒子からの距離の四倍に比例して減少する場合、全体的な引力は以下のようになります。

$A \& B : \text{引力} = \frac{A \, \text{立方}}{A \, \text{平方平方}} \quad \& \quad \frac{B \, \text{立方}}{B \, \text{平方平方}}$

すなわち、物体の側面である $A$ と $B$ の平方比の逆数です。

\textbf{命題LXXXVIII.} \textbf{理論XLV.}

\textit{2.\,コロラリウム.} したがって、逆に、類似の物体が類似の粒子に対して引き力を持つとき、粒子の位置から物体までの距離の冪乗比に引き力の減少率を結びつけることができる。もしその減少率が各々の距離の冪乗比であるならば、

\[
\text{物体全体の引く力はその中心に向かいます。}
\]

ここで、\(\textbf{A}\) と \(\textbf{B}\) は立方体を表します。そして、この論証は他の場合にも適用されます。

\[
\textbf{Corol.\,2.} \quad \text{したがって、逆に、類似の物体が類似の粒子に対して引き力を持つとき、粒子の位置から物体までの距離の冪乗比に引き力の減少率を結びつけることができる。もしその減少率が各々の距離の冪乗比であるならば}
\]

\[
v \propto 1/r^n
\]

ここで \(v\) は引く力を表し、\(r\) は粒子から物体までの距離を表します。指数 \(n\) は具体的な関数により定義されます。

\[
\text{物体全体の引く力はその中心に向かいます。}
\]

中心の重力点は；これは、球体から類似かつ等しい定数密度を持つ物体の引力と同様に、その重心の中心にある中心である。


\par\begin{center}
  \centering
  \includegraphics[width=1043px,height=990px,trim=-24]{224.png}

\end{center}\par

物体 \( RSTV \) の部分点 \( A \), \( B \) が任意の粒子 \( Z \) を力 \( AZ \), \( BZ \) で引き寄せるとする。部分点が互いに等しいと仮定すれば、これらは距離 \( AZ \), \( BZ \) に比例します。また、部分点が不等であると仮定すれば、これらの比はそれぞれ距離 \( AZ \), \( BZ \) の導出された値に比例します。そして、力はそれぞれの部分点から粒子への距離に部分点の質量を乗じた積で表されます。したがって、
\[
A \times AZ, \quad B \times BZ
\]
と表すことができます。これらの力を結びつける線分を引きます。

\textbf{第13節：非球形物体の引力（命題LXXXV〜XCIII）}

\textit{第11部分}

円弧 \( AB \) に沿って点 \( A \) から点 \( B \) まで移動し、点 \( G \) を通るように接続します。このとき、\( AG \) に対する \( BG \) の比が点 \( B \) に対する点 \( A \) の比と等しくなるようにします。すると、点 \( G \) は点 \( A \) および点 \( B \) の重心となることになります。

力 \( A \times AZ \) を法則の補則2に従って分解すると、力 \( A \times GZ \) と \( A \times AG \) に分解されます。同様に、力 \( B \times BZ \) は力 \( B \times GZ \) および \( B \times BG \) に分解されます。

しかし、力 \( A \times AG \) および \( B \times BG \) は、点 \( A \) に対する点 \( B \) の比と円弧 \( BG \) に対する円弧 \( AG \) の比が等しいことから、これらの力は等しくなります。さらに、これらは反対方向に作用するため、互いに相殺されます。

残る力は \( A \times GZ \) および \( B \times GZ \) です。これらは点 \( Z \) から重心 \( G \) 方向に作用し、合力は \( A + B \times GZ \) となります。


非球形物体の引力（命題LXXXV〜XCIII）の一部

これらの部分物質 \( A \) と \( B \) を組み合わせると、これは、その共通の重心 \( G \) にある点での重力による力を表します。すなわち、球体を形成します。

同様の議論により、第三の部分物質 \( C \) を追加すると、この力と部分物質 \( A + B \times GZ \) が重心 \( G \) に向かう力 \( A + B \times GZ \) の合成を組み合わせます。これによって生成された力は、その他の部分物質 \( C \) および重心 \( G \) に向けられます。すなわち、三つの部分物質 \( A \), \( B \), \( C \) の共通の重心に向かいます。また、これは部分物質 \( C \) と球体がその共通の重心にある場合と同じ力となります。

\[
\textbf{訳注：} GZ は重心から部分物質までの距離を表します。
\]




\section*{第1編 第13節：非球形物体の引力（命題LXXXV〜XCIII）}


\textbf{命題 LXXXV.} 物体 $RSTV$ の各部分から物体 $Z$ に及ぼす力を分解すると、その力は無限大まで延びていく。したがって、物体 $RSTV$ 全体の合力は、中心の重力点を保つ条件下で、球体と同等である。

\textbf{証明} \textit{Eadem est igitur vis tota particularum omnium corporis cujusque $RSTV$ ac si corpus illud, servato gravitatis centro, figuram globi indueret.}

\textbf{命題 LXXXVI.} 由此，被吸引物体 \( Z \) 的運動與假設吸引物體 \( RSTV \) 為球形時之運動相同。従って，若吸引物體處於靜止或等速直線運動狀態，則被吸引物體將在以吸引物體的重力中心為焦點的橢圓上運動。

修正後：
\textbf{命題 LXXXVI.} 由此，被吸引物体 \( Z \) 的运動与假設吸引物体 \( RSTV \) 為球形時之运動相同。故若吸引物体處於靜止或等速直線運動狀態，則被吸引物體將在以吸引物体的重力中心為焦點的橢圓上運動。

修正：

\textbf{命題 LXXXVI.} 由此，被吸引物体 \( Z \) 的运動与假設吸引物体 \( RSTV \) 為球形時之运動相同。故若吸引物体處於靜止或等速直線運動狀態，則被吸引物體將在以吸引物体的重力中心為焦點的橢圓上運動。


\textbf{命理 XLVI.} 若有多個物体由等質部分組成，各部分的力量





\textbf{訳文}

非球形物体の引力は、各点から中心点までの距離に比例します。引力はすべての力の合計であり、その結果、任意の粒子が引き寄せられる方向は、引く物体の共通の重心に向かいます。また、この引力は、引く物体を共通の重心で球体と仮定した場合と同じです。

同様に証明し、前節の命題も確認できます。

\textit{補遺} より：従って、引き寄せられた物体の運動は、引く物体を共通の重心で球体と仮定した場合と同じになります。したがって、引く物体の共通の重心が静止しているか、または直線で一様に進行しているなら、物体の運動も同様です。

\textbf{命題 XC. 問題 XLIV.}

\textit{円周上の各点に中心力として減少する力を向ける場合、物体が直線を沿ってどの位置でも引力を受けたときにその直線は円の平面における中心から垂直となる力を求める。}

中心 \( A \) から任意の距離 \( AD \) を引いたとき、直線 \( AP \) が平面に垂直である。


\par\begin{center}
    \centering
    \includegraphics[width=1025px,height=1125px,keepaspectratio]{226.png}

\end{center}\par

\textbf{円を描き、任意の点} \( E \) から物体 \( P \) を引き寄せる力について説明する。直線 \( PE \) は点 \( E \) から物体 \( P \) への方向を示す。直線 \( PA \) 上に点 \( F \) をとるとき、\( PF \) は \( PE \) と等しい長さとする。また、正規線 \( FK \) を引き、この正規線は点 \( E \) から物体 \( P \) を引き寄せている力の方向を表すものとする。さらに、曲線 \( IKL \) を考え、点 \( K \) が常に接する直線とし、円の平面が \( L \) で交わる。直線 \( PA \) 上に点 \( H \) をとるとき、\( PH \) は \( PD \) と等しい長さとする。また、曲線 \( IKL \) の直線 \( HI \) との交点を \( I \) とする。すると、物体 \( P \) は\(\textbf{この条件} \)のもとで運動する。

\textbf{第13節：非球形物体の引力（命題LXXXV〜XCIII）}

\textbf{命題 LXXXV. 定理 85.}

円周上に引かれた面積 $AHIL$ を、点 $P$ から垂直方向に引き下ろした高さ $AP$ に沿って円周へ向けての引力を表すと、この引力は $\textbf{AHIL} \cdot AP$ となる。$\blacksquare$

\textbf{証明.}

まず、線分 $AE$ 上で最小値を持つ点 $Ee$ を選びます。点 $P$ から $Pe$ を結び、さらに $PA$ 上で $Pf$ を選び、$Pe$ と等しい長さにします。引力は、円環上の任意の点 $E$ が自身の中心 $A$ へ引き寄せられる力を表し、この力は $FK$ の比となると仮定します。したがって、点 $E$ から物体 $P$ への引力は $\frac{AP \times FK}{PE}$ となります。

また、円環全体が物体 $P$ を点 $A$ 方向へ引き寄せるとするとき、その引力は各点の引力を加算したものです。この円環の引力は、円環の面積と等しい $\frac{AP \times FK}{PE}$ となります。

しかし、円環は円弧 $AE$ の長さと線分 $Ee$ の幅の積に比例します。よって、この矩形の面積は $(\textbf{AE} \cdot Ee)$ であり、これは円環全体の引力を表します。

以上より、非球形物体の引力は円周上に引かれた面積と高さの積に比例することが示されました。


\textbf{比例}$PE \& AE, Ee \& cE$は直角三角形の面積$PE \times cE$または$PE \times Ff$に等しいです。したがって、環が物体$P$を点$A$に向けて引き力を持つとき、その引力は比例します：

$PE \times Ff \quad$ \text{および} $\quad \frac{AP \times FK}{PE}$

の積に等しくなります。すなわち、
\[
Ff \times AP \times FK, 
\]
または円錐形領域$FKkf$の面積を$AP$で掛けると同様です。したがって、円周上にあるすべての環が中心点$A$および半径$AD$からなる円を描くとき、物体$P$を点$A$に向けて引き力を持つ総合力は以下の面積に比例します：
\[
\textbf{領域}AHIKL \quad \text{を} \quad AP \quad で掛けると得られる全体の面積。
\]
\QED

\textit{\textbf{補題1。}}\wsp{}したがって、引き力が点に対して距離の二乗に反比例する場合、




\textbf{命題LXXXV.} このとき、すなわち \( FK \) が \( \frac{1}{PF^2} \) に比例する場合、面積 \( AHIKL \) は \( \frac{1}{PA} - \frac{1}{PH} \) に比例します。したがって、点 \( P \) の粒子が円周に対しての吸引力は \( 1 - \frac{PA}{PH} \) に比例し、すなわち \( \frac{AH}{PH} \) に比例します。

\textit{\textbf{補題2.}} より普遍的に、点 \( P \) の粒子に対する吸引力は、距離 \( D \) に対して力が \( D^n \) の冪に反比例するとき、すなわち \( FK \) が \( \frac{1}{D^n} \) に比例する場合、面積 \( AHIKL \) は \( \frac{1}{PA^{n-1}} - \frac{1}{PH^{n-1}} \) に比例します。したがって、点 \( P \) の粒子に対する吸引力は \( \frac{1}{PA^{n-2}} - \frac{PA}{PH^{n-1}} \) に比例します。


\textbf{命題 XC I. 問題 XL V.}

円の直径を無限大に増加させ、数 \( n \) を 1 より大きいとするとき、点 \( P \) の粒子が全体的な平面に及ぼす引力は、\( PA^{n-2} \) と反比例する。これはなぜなら另一极限
\[
\frac{PA}{PH^{n-1}}
\]
趨于零。

【図】

円錐体 \( ADEFG \) 内に粒子 \( P \) を引き込む。各点への中心力は距離の逆数で減少していると仮定する。


\par\begin{center}
  \centering
  \includegraphics[width=0.8\textwidth]{228.png}

\end{center}\par

軸 \( AB \)。任意の円 \( RFS \) がこの軸に対して垂直に截断され、その固体を切断し、直径 \( FS \) 上の平面 \( PALKB \) において、軸を通過する点 \( P \) 処から固体への引力の長さ \( FK \) を (命題 XC による) 求めます。この点 \( K \) が曲線 \( LKI \) と平面 \( AL \) および \( BI \) の交点で、これらは軸 \( A \) および \( B \) に交わるとき、物体 \( P \) に対する引力は固体に対して円錐面 \( LABI \) の面積に比例します。証明終了。

\textit{命題1.} もし固体が円柱である場合、軸 \( AB \) 回転して平行四辺形 \( ADEB \) から形成される平面 \( ADEB \) であれば、

\[
\textbf{証明：}
\]

円柱の底面は円 \( AL \) および \( BI \) と一致します。したがって、引力 \( FK \) の長さは円錐面 \( LABI \) の面積に比例します。証明終了。

\[
\textbf{補足：}
\]

円柱の場合、底面の円 \( AL \) および \( BI \) が平面 \( ADEB \) 上に存在することから、引力 \( FK \) の長さは円錐面 \( LABI \) の面積に比例します。

非球形物体の引力（命題LXXXV〜XCIII）における以下は、第13節の一部です。

\begin{center}
\textbf{命題 LXXXIX}

物体 \( P \) が円柱体に引き寄せられる力を描写し、各点へ向かう中心力は、その点から円柱までの距離の平方と反比例するとする。すると、物体 \( P \) がこの円柱体 \( C \) に対して受ける引力は \( BA - PE + PD \) となる。

命題 XC

直線 \( FK \) の順次適用（Corol. 1, Prop. XC）により、上記の式は \( 1 - \frac{PF}{PR} \) と等しい。

この式の第1項を長さ \( AB \) に引くと、面積 \( 1 \times AB \) を描写する。また、第2項を長さ \( PB \) に引くと、面積 \( 1 \times (PE - AD) \) を描写する（これは曲線 \( LKI \) の平方計算から容易に示される）。
\end{center}


\par\begin{center}
\centering
\includegraphics[width=0.8\textwidth]{figura.png}
\label{fig:principia}

\end{center}\par

このように、各項を長さに引くことで、それぞれの面積が描写されます。

\textbf{命題LXXXVII}

直線 \( PA \)は、点 \( P \)から点 \( D \)までの距離 \( PD - AD \)に引かれた面積1を描き、直線 \( PB \), \( PA \)の差 \( AB \)は、直線 \( PE - PD \)に引かれた面積の差1を描く。

まず、面積 \( 1 \times AB \)から最後の面積 \( 1 \in PE - PD \)を引くと、残る面積 \( LABI \)は \( 1 \in (AB - PE + PD) \)に等しい。したがって、この面積に対する力は \( AB - PE + PD \)に比例する。


\par\begin{center}
    \centering
    \includegraphics[width=1344px,height=1081px,trim={24 0 0 0},clip]{229.png}

\end{center}\par

\textbf{補則2}

したがって、球体 \( AGBCD \)から外側に位置する任意の物体 \( P \)に対する力は、軸 \( AB \)上にあるとき、その大きさは \( AB - PE + PD \)に比例する。

\textit{証明終了}

非球形物体の引力（命題LXXXV〜XCIII）【第13節：第24段落】

円錐曲線の軸上に垂直に引かれた線分 \( ER \) は、点 \( P \) から垂線 \( PE \) に引かれるとき、常に線分 \( PD \) と等しくなります。ここで、線分 \( PD \) は球楕円体を交わる点 \( D \) まで引き出されます。球楕円体の頂点 \( A, B \) から軸 \( AB \) に垂直な線分 \( AK, BM \) が引かれ、それぞれ線分 \( AP, BP \) と等しくなります。したがって、これらの円錐曲線は点 \( K \) および \( M \) で交わり、線分 \( KM \) を結び、同じく \( KMRK \) の部分を除きます。

球楕円体の中心を \( S \)、最大半径を \( SC \) とします。すると、球楕円体から物体 \( P \) を引き込む力は、直径 \( AB \) 上に描かれた球から同様の物体を引き込む力に対して比例します。

\[
\textbf{証明}
\]

まず、線分 \( ER \) が点 \( E \) で垂直に交わる円錐曲線 \( NKRM \) の性質を考えます。この円錐曲線は球楕円体を交わる点 \( D \) まで垂線 \( PD \) と等しい長さを持ちます。

次に、球楕円体の頂点 \( A, B \) から軸 \( AB \) に対して垂直な線分 \( AK, BM \) を引きます。これらの線分はそれぞれ \( AP, BP \) と等しくなります。したがって、これらは円錐曲線 \( NKRM \) の交点 \( K \) および \( M \) で交わります。

線分 \( KM \) を結び、同じく \( KMRK \) の部分を除きます。球楕円体の中心 \( S \) と最大半径 \( SC \) を定義します。このとき、球楕円体から物体 \( P \) を引き込む力は、直径 \( AB \) 上に描かれた球から同様の物体を引き込む力に対して比例します。

\[
\textbf{証明終了}
\]

\textbf{同様に、}
\[
\frac{AS \times CS^2 - PS \times KMRK}{PS^2 + CS^2 - AS^2}
\]
は
\[
\frac{AS^3}{3PS^4}
\]
に等しい。

そして同じ計算に基づいて、球楕円体の弧に対する引力も求めることができる。


\par\begin{center}
  \centering
  \includegraphics[width=898px,height=881px,keepaspectratio]{230.png}

\end{center}\par

\textbf{補題 3.} 仮定するように、球楕円体内の任意直径上に粒子を配置した場合、その粒子に対する引力は中心からの距離に比例する。

この事実は次のように証明できる。$AGOF$ を球楕円体が引き寄せ、$S$ をその中心、$P$ を引張られた粒子とする。粒子 $P$ に対して働く力は

\[
\frac{AS \times CS^2 - PS \times KMRK}{PS^2 + CS^2 - AS^2}
\]
であり、

\[
\frac{AS^3}{3PS^4}
\]

に等しい。

\textbf{さらに、半径}\(\textbf{SPA}\)，そして任意の直径\(\textbf{DE, FG}\)が楕円体に交わる点として\(\textbf{D, E, F, G}\)を定める。これらの楕円体は内部同心で相似であり、\(\textbf{PCM, HLN}\)はこれら2つの楕円体の表面である。前記楕円体が体\(\textbf{P}\)を通るとき、直線\(\textbf{DE, FG}\)をそれぞれ点\(\textbf{B, C}\)で交わす平面\(\textbf{PCM}\)がある。後者の楕円体は直線\(\textbf{DE, FG}\)をそれぞれ点\(\textbf{H, I, K, L}\)で交わす。これらの楕円体は共通の軸を持つものとする。したがって、直線\(\textbf{DE, PB, HI}\)が互いに等しい部分\(\textbf{DP, BE, FP, CG, DH, IE, FK, LG}\)をそれぞれ截取する。


\par\begin{center}
  \centering
  \includegraphics[width=0.8\textwidth]{230.png}

\end{center}\par

\textbf{同一点で二等分}され、直線 $FG$, $PC$ および $KL$ となるようにします。ここで、$DPF$ と $EPG$ を対立する円錐を表し、これらは垂直角 $DPF$ と $EPG$ の非常に小さな角度で描かれます。また、線分 $DH$, $EI$ も非常に短いものとします。そして、球形楕円面上の部分線分 $DHKF$ および $GLIE$ は、線分 $DH$ および $EI$ の長さが等しいことから、これらは互いに

\[
\textbf{その距離の平方比} \quad \textbf{すなわち} \quad \textbf{点} \, P \, \textbf{からの距離の平方比}
\]

となります。したがって、点 $P$ からこの部分線分に及ぼす引力は等しくなります。

また、同様の理由により、無数の同心球形楕円面上の部分線分も互いに等しい引力を及ぼします。

\& これらが共有軸を持つ場合、空間 \(DPF\) と \(EGCB\) は粒子に分割され、すべての粒子は等しく物体 \(P\) を反対方向へ引き寄せます。したがって、円錐セグメント \(EGCB\) と円錐線分 \(DPF\) の力は等しく、互いに相殺します。また、物質全体が内球形体 \(PCBM\) 内部にある場合の力を考慮すると、物体 \(P\) は単なる内球形体 \(PCBM\) から引き寄せられます。したがって（命題 LXXII の補則 3）、その引力は物体 \(A\) が全体の球形体 \(AGOD\) から引き寄せられる力に比例し、距離 \(PS\) に対する距離

\[
\frac{PS}{SA}
\]

となります。

\textbf{第1編 第13節：非球形物体の引力（命題LXXXII〜XCIII）}

\textit{第46問. 命題 XCII. 与えられた吸引体について、その点々に指向する吸引力の減少率を求める方法を示す。}

与えられた吸引体として球体や円柱体などの規則的な形状を考える（図231参照）。これらの物体に対する引力法則は、任意の距離での吸引力の減少率に対応している（命題LXXX, LXXXI, XCI 参照）。したがって、実験によって各距離での吸引力を求めて、その結果から引力法則全体を明らかにすると、求めなければならない個々の吸引力の減少率も得られる。


\par\begin{center}
    \centering
    \pngright{231.png}{}{}{}

\end{center}\par

\textbf{命題 XCIII. 理論 XLVII.}

\textit{もし一方を平面、他方を無限大なる部分から成る球形でない物体が、等しい力を持ち、その力を物体から離れるほど力が減少し、距離のべき乗より大きいべき乗率で減少する粒子から構成されているとすると、その物体全体が平面の一方に配置された微小な物体を引き寄せられる：この場合、物体の引力は平面からの面積への垂線距離に対して、べき乗率で減少し、そのべき乗指数は3より小さいべき乗距離の指数よりも小さくなります。}

\textbf{証明：} 物体が無限大な部分から成る球形でないものと仮定します。この物体は等しい力を持ち、その力を物体から離れるほど力が減少し、距離のべき乗より大きいべき乗率で減少する粒子から構成されています。

微小な物体を平面の一方に配置します。物体の引力は平面からの面積への垂線距離に対して、べき乗率で減少します。このべき乗指数は3より小さいべき乗距離の指数よりも小さくなります。

\textbf{証明終了}

\[
\textbf{1.} \quad \text{案。} \, LGGl \, \text{が平面で、その平面に終止する固体を考えてみましょう。} \, I \, \text{に向かって} \, solidum \, \text{が} \, J \, \text{に落ちているとします。} \, mHM, \, nIN \, \&c. \, \text{は平面} \, GL \, \text{に対して平行な無数の直線です。} \, C \, \text{が固体外にある物体である場合、} \, CGHI \, \text{をこれらの無数の平面に垂直に引きます。} \, \text{そして、固体上の各点から} \, C \, \text{への引力は距離の} \, n \, \text{次方根比例で減少します。} \, (n \, \text{は3より小さい自然数}) \, \text{であるとします。}
\]

\[
\textbf{したがって、(命題 XC の補則 3) 平面} \, mHM \, \text{によって引き伸ばされる任意の点} \, P \, \text{に対する引力は} \, C \, \text{から} \, P \, \text{までの距離} \, r \, \text{の} \, n \, \text{次方根比例で減少します。}
\]

\[
F \propto \frac{1}{r^n}
\]

\textbf{C}は\(\frac{1}{CH^{2-n}}\)と反比例します。平面\(mHM\)において、長さ\(HM\)を\(CH^{n-2}\)に対する反比例するものとして定めると、その力は\(HM\)に比例します。同様に、平面\(lGL\), \(nIN\), \(oKO\)等において、長さ\(GL\), \(IN\), \(KO\)等をそれぞれ\(CG^{n-2}\), \(CI^{n-2}\), \(CK^{n-2}\)に対する反比例するものとして定めると、各平面の力はこれらの長さに比例し、したがって、全体的な立体の力を表す総合力は無限大方向への領域\(GLOK\)を無限大まで積分した面積に比例します。しかし、その領域は既知の平方曲線によって求めることができます。


\par\begin{center}
  \centering
  \includegraphics[width=0.8\textwidth]{232.png}

\end{center}\par

\textbf{方法は} \( CG^{n-3} \) に反比例するため、全体の物体に対する力も \( CG^{n-3} \) に反比例します。\QED

\textit{案2.}\wsp{}平面 \( lGL \) の内側に物体 \( C \) を配置し、距離 \( CK \) を \( CG \) と等しくします。部分的な物体 \( LGloKO \)，平行な平面 \( lGL, oKO \) で終了する，物体 \( C \) は平面 \( OK \) の中間にある任意の部分を引き寄せません。反対方向の点の作用が等しいとすれば、相互に相殺します。したがって、物体 \( C \) は平面 \( OK \) を超えて物体に対して引き寄せられます。これで、

\[
\textbf{証明終了}
\]


\textbf{しかし、力（最初の場合）は} \textbf{$CK^{n-3}$ に反比例し}、これは（等しい $CG$, $CK$ について）$CG^{n-3}$ に反比例することを意味します。 \QED

\textit{\textbf{補題1.}}\wsp{}したがって、平面 $LG$, $IN$ に平行で無限大である固体 $LGIN$ の各端点で終了する場合、その引力は、全体の無限体 $LGKO$ の引力から後部部分 $NIKO$ の引力を引いたものとして明らかになります。これは $KO$ 方向へ無限大に伸長します。

\textit{\textbf{補題2.}}\wsp{}もしこの無限体の後部部分が、その吸引力と前部部分の吸引力との合算がほとんどない場合、それはほぼ零です。

\textbf{動量が失われ、より近い部分の引力は距離を増加させるにつれてほぼ $CG^{n-3}$ の比例で減少します。}

\textit{\textbf{補則 3.}}\wsp{}したがって、中間平面から一方に有限な物体が平面内の粒子を引き込む場合、粒子と平面との距離と引力施す物体の寸法との比が非常に小さいとき、引力を施す物体は同種の部分から成り立つもので、その吸引力は任意より大きい次元の距離の比例で減少するものとする。

引力は物体全体に対して、近似してその距離の冪乗に反比例し、その指数は非常に小さな距離を含む項で三つ目より小さいものである。恒星由来の粒子からなる物体において、粒子間の引力が距離の立方に反比例するという主張は成り立たない。なぜなら、この場合、無限遠方の部分の引力は二条の推論において常に有限遠方の部分の引力より無限大であるからである。

\textit{補遺。}

ある物体が与えられた平面に対して垂直に引き寄せられるとき、

$$
\textbf{vis attractiva corporis totius decrescet quamproxime in ratione potestatis, cujus latus sit distantia illa perexigua, \& Index ternario minor quam Index potestatis prioris.}
$$

$$
\textbf{De corpore ex particulis constante, quarum vires attractivae decrescunt in ratione potestatis triplicat\ae distantiarum, assertio non valet, propterea quod, in hoc casu, attractio partis illius ulterioris corporis infiniti in Corollario secundo, semper est infinite major quam attractio partis citerioris.}
$$

$$
\textbf{Si corpus aliquod perpendiculariter versus planum datum trahatur,}
$$


この問題は、（命題XXVIIにより）この平面へまっすぐ落ちる物体の運動を求め、（運動法則の系2により）その運動を、同じ平面に平行な線に沿った一様な運動と合成することで解ける。逆に、平面に向かって垂線の方向に働く引力の法則を求める場合は、引かれる物体が与えられた任意の曲線上を動くという条件のもとで、問題3の例にならって計算すれば解ける。

ただし、こうした計算は、縦線を収束する級数に展開して簡単にするのがふつうである。たとえば、底辺 $A$ に、任意に与えた角で縦線の長さ $B$ が立てられ、それが底辺の任意のべき $A^{\frac{m}{n}}$ に比例するとしよう。そして、縦線の向きに底辺へ引かれる（あるいは底辺から押しのけられる）物体が、縦線の上端がいつもたどる曲線の上を動くようにする力を求めるとする。私は底辺が最小の部分 $O$ だけ増すと考え、縦線 $\overline{A + O}^{\frac{m}{n}}$ を無限級数
\[ A^{\frac{m}{n}} + \frac{m}{n} OA^{\frac{m - n}{n}} + \frac{mm - mn}{2nn} O^2 A^{\frac{m - 2n}{n}} + \text{など} \]
に展開する。そして、この級数のうち $O$ が2次の項、すなわち項 $\frac{mm - mn}{2nn} O^2 A^{\frac{m - 2n}{n}}$ に、力が比例するものとする。したがって、求める力は $\frac{mm - mn}{nn} A^{\frac{m - 2n}{n}}$ に、同じことだが $\frac{mm - mn}{nn} B^{\frac{m - 2n}{m}}$ に比例する。

たとえば、縦線が放物線をたどるなら、$m = 2$, $n = 1$ だから、力は一定の量 $2B^0$ に比例し、したがって力は一定である。ゆえに一定の力のもとで、物体はガリレオが示したとおり放物線上を動く。また、縦線が双曲線をたどるなら、$m = -1$, $n = 1$ だから、力は $2B^3$ に比例する。ゆえに縦線の3乗に比例する力のもとで、物体は双曲線上を動く。しかし、この種の命題はこのくらいにして、まだ触れていない運動についての別の命題に進もう。


\pngright{235.png}{}{}{}



\section{第1編 第14節：向心力に向かって引き寄せられる極微小体の運動（光学・屈折への適用、命題XCIV〜XCVIII）}


% Section: 第1編 第14節：向心力に向かって引き寄せられる極微小体の運動（光学・屈折への適用、命題XCIV〜XCVIII）



\textbf{第14節}

\begin{center}
\textbf{向心力に向かって引き寄せられる極微小体の運動（光学・屈折への適用、命題XCIV〜XCVIII）}
\end{center}

\textbf{命題 XCIV. 理論 XLVIII.}

もし相似な二つの媒質が平面に並行な間隔で区切られ、かつ物体がその空間を垂直方向から一方の媒質に対して引き寄せられたり押し付けられる場合、他の力は何らかの形で作用したり阻害されないとして、

\textbf{ただし}引張りは、同一部分に対し、それぞれ異なる距離から平面に及ぶとき、

\[
F \propto 1/r^2
\]

ここで \( r \) は物体と平面との距離を表します。


向心力に向かって引き寄せられる極微小体の運動（光学・屈折への適用）

命題 XCIV. 定理 4.

もし同じ条件が得られれば、入射角の正弦と反射角の正弦との比は一定である。


\par\begin{center}
    \centering
    \includegraphics[width=1308pt,height=1498pt,keepaspectratio]{235.png}

\end{center}\par

\textit{例 1.} 平面 $Aa$ と平面 $Bb$ を平行な二平面とする。物体が平面 $Aa$ 上の直線 $GH$ に垂直入射し、中間領域を通過して平面 $Bb$ から出射する。この作用により曲線 $HI$ を描き、直線 $IK$ 沿って出射する。平面 $Bb$ の出射方向に対して直角に引いた垂線 $IM$ が存在し、






向心力に向かって引き寄せられる極微小体の運動（光学・屈折への適用、命題XCIV〜XCVIII）

線 \( GH \) が点 \( M \) で直交し、平面 \( Aa \) が点 \( R \) で直交し、射出線 \( KI \) が \( HM \) の点 \( L \) で交わる。中心を \( L \) として、半径 \( LI \) を描く円は、\( HM \) 上の点 \( P \) と \( Q \) および延長線 \( MI \) 上の点 \( N \) で交わる。まず、引き寄せまたは衝撃が一様であるとするとき（ガリレオによって証明された）、曲線 \( HI \) は放物線であり、その性質は与えられた直角三角形の底辺と点 \( I \) から点 \( M \) までの距離が等しい直角三角形の斜辺の平方に等しいというものです。したがって、点 \( H \) から点 \( M \) の距離は二等分される点 \( L \) で二等分されます。よって、点 \( M \) から点 \( I \) への線分を下ろすと、

\[
\textbf{証明終了}
\]


垂線 \( LO \) は \( MO \), \( OR \) と等しい。さらに等しい量 \( IO \), \( ON \) を加えると, \( MN \), \( IR \) も等しくなる。したがって, \( IR \) が与えられたとき, 同様に \( MN \) も与えられ, 直角三角形 \( NMI \) の直角の対辺 \( NM \) は, 假想線 \( HMq \) と等しい。すなわち, \( IM \), つまり \( HMq \) が与えられた比で与えられる。

また, 直角三角形 \( NMI \) は直角三角形 \( PMQ \) に等しく, これは平方差 \( ML^2 - PL^2 \) や \( LI^2 \) に等しい。さらに, \( HMq \) は \( LMq \) の四分之一と比を成す。したがって, \( ML^2 - LI^2 \) に対する \( ML^2 \), そして部分的に, \( LI^2 \) に対する \( ML^2 \), さらに半分の比も与えられる。

\[
LI \textbf{ ad } ML.
\]
次に、任意の三角形 \( LMI \)において、各角の正弦は対辺と比例します。従って、入射角 \( LMR \) の正弦に対する反射角 \( LIR \) の正弦の比が与えられます。
\QED

\textbf{案2.} 今、極微小体が平行な平面区間 \( AabB \), \( BbcC \) 等を順に通過し、各平面に対する力はそれぞれ一定でありながら異なるとします。そして、既に証明した通り、平面 \( Aa \) における入射角の正弦は平面 \( Bb \) からの反射角の正弦に対して一定の比で与えられます。この正弦は、

\[
\sin(\text{入射角}) : \sin(\text{反射角})
\]

と表されます。

向心力に向かって引き寄せられる極微小体の運動（光学・屈折への適用、命題XCIV〜XCVIII）

この極微小体が平面 \( Bb \) に至る入射角の正弦を表す \(\textbf{sinus incidenti}\ae\) は、第3平面 \( Cc \) から出射角の正弦 \(\textbf{sinus emergenti}\ae\) への比で与えられます。そして、この比は任意に定めることができます；同様に、第4平面 \( Dd \) から出射角の正弦も任意の比で与えられます。これを無限に繰り返すことで、最初の平面から最後の平面までの出射角の正弦が与えられ、平面間の距離は減少し、その数は無限大となります。これにより、向心力または推進力の作用は任意に定められた法則に従って継続します。


\par\begin{center}
  \centering
  \includegraphics[width=1284px,height=1121px,trim=-12 0 0 0]{237.png}

\end{center}\par


向心力に向かって引き寄せられる極微小体の運動（光学・屈折への適用、命題XCIV〜XCVIII）【第14節】

\textbf{命題 XCIV.} ある角度で平面に衝突する極微小体の速度比は、その角度から平面を離れて再び飛ぶときの速度比と恒常的に等しい。

証明：同様に置くとする。すなわち、衝突前の速度と衝突後の速度は、出射角の正弦と入射角の正弦との比である。

\textbf{証明終了} \QED


\textbf{命題 XCV.} 同じ条件において、極微小体の速度比は、その角度から平面を離れて再び飛ぶときの速度比と恒常的に等しい。

同様に置くとする。すなわち、衝突前の速度と衝突後の速度は、出射角の正弦と入射角の正弦との比である。

\textbf{証明終了} \QED


\condpagelarge{命題 XCVI.}


同様にして置くべく、すなわち \(AH\) と \(Id\) を等しくし、直線 \(AG\) および \(dK\) を入射線 \(GH\) および \(IK\) に対して垂直に引き上げる。次に \(GH\) 上に点 \(T\) をとり、その長さを \(IK\) に等しいものとする。さらに平面 \(Aa\) から垂線を下ろす。


\[
AH = Id, \quad AG \perp GH, \quad dK \perp IK
\]
\[
GH において TH = IK
\]

\[
AG, dK は直線 GH, IK に対して垂直に引き上げられた。
\]

\[
Tv.
\]
通常、物体の運動は二つに分けることができる。一つは平面 \( Aa \), \( Bb \), \( Cc \) などに対して垂直な運動であり、もう一つはこれらの平面に対して平行な運動である。引力または衝撃が直線を垂直方向に作用させた場合、その力は平行方向の運動を変化させないため、物体は等しい時間内に等しい距離を平行方向で移動する。すなわち、点 \( H \) から線分 \( AG \) までの距離と、点 \( I \) から線分 \( dK \) までの距離が等しいとき、物体は等しい時間内に線分 \( GH \) および \( IK \) を描く。したがって、入射速度は反射後の速度に対して等しい。

\pngright{238.png}{}{}{}
\[
\textbf{図}
\]

\textbf{命題 XCIV.}

出射角は入射角に等しい、すなわち、$GH$ が $IK$ や $TH$ に比例し、これは $AH$ や $Id$ が $vH$ に比例することを意味します。すなわち（放射線 $TH$ や $IK$ の観点から）出射角の正弦は入射角の正弦に等しい。\QED

\textbf{命題 XCIV.}

\textit{同様にして、かつ前後の運動が入射前に速いものであると仮定するとき、私は極微小体が放射線を描く直線を傾け、反射されるという主張をする。そして反射角は入射角に等しい。}

\begin{center}
\includegraphics[width=1293pt,height=640pt,trim=-24 0 0 0]{238.png}
\end{center}

次のように考えます。平面 $Aa$、$Bb$、$Cc$ \&c. 上を描く極微小体がパラボラ形の弧 $HP$、$PQ$ を描くとします。\pngright{238.png}{}{}{}

向心力に向かって引き寄せられる極微小体の運動（光学・屈折への適用、命題XCIV〜XCVIII）

$QR, \&c.$ そして、入射線 $GH$ が平面最初の $Aa$ に対してなす角を $\textbf{AQ}$ とする。すなわち、入射角は円の半径に対する正弦と同一であり、この正弦が入射正弦と出射正弦（平面 $Dd$）との比に等しい空間 $DdeE$ の比である。そして、既に出射正弦を半径とし等しくしたとき、出射角は直角となり、従って出射線は平面 $Dd$ と一致する。

物体がこの平面 $Dd$ 上の点 $R$ に達すると、出射線が同様の平面と一致することから明らかに、物体は平面 $Ee$ 方向には進めないことがわかる。しかし、




物体が平面 \(Dd\) 上の点 \(R\) に達したとき、出射線がその平面と一致することから明らかに、物体は平面 \(Ee\) 方向には進めないことがわかる。しかし、



\textbf{これには、極微小体が一直線 $Rd$ から進むことができません。}これは、常に直進方向から媒質に引き寄せられたり押し出されたりするためです。その結果、平面 $Cc$, $Dd$ の間で放物線弧 $QRq$ を描きながら反転します。この放物線の頂点（Galilaeiによる証明によれば）は主に点 $R$ にあります；平面 $Cc$ において同じ角度で点 $q$ に、かつ先に点 $Q$ に切ります。その後、弧 $qp$, $ph$ \&c. を通って、前の弧 $QP$, $PH$ と相似かつ等しい角度で他の平面を切ります；そして先に点 $P$, $H$ \&c. に切ります。最後には直線 $Rd$ から再び進み出ます。

\textbf{同様に、極微小体は} $h$ において入射角 \(\theta\) で平面 \(Aa\), \(Bb\), \(Cc\), \(Dd\), \(Ee\) に沿って無限に小さくするとき、これらの間隔を増加させ、引力または推進力が与えられた法則に従って継続的になるようにします。そして、出射角は常に入射角と等しくなります。

\textit{補遺}。


\par\begin{center}
  \centering
  \includegraphics[width=7.75in,height=5.89in,trim={0cm 24cm 0cm 0cm},clip]{240.png}

\end{center}\par

これらの引力は、与えられた接線比に従って生じる光の反射と屈折とほとんど異なりません。

\textbf{翻訳}

ut invenit \textbf{Snellius}, \& per consequens secundum datam Sinuum rationem, ut exposuit \textbf{Cartesius}. Nam Lucem successive propagari \& spatio quasi decem minutorum primorum a Sole ad Terram venire, jam constat per Phaenomena Satellitum \textbf{Jovis}, Observationibus diversorum Astronomorum confirmata. Radii autem in aere existentes (uti dudum \textbf{Grimaldus}, luce per foramen in tenebrosum cubiculum admissa, invenit, \& ipse quoque expertus sum) in transitu suo prope corporum vel opacorum vel perspicuorum angulos (quales sunt nummorum

\textbf{LaTeX形式}


ut \textbf{Snellius} invenit, \& per consequens secundum datam Sinuum rationem, ut \textbf{Cartesius} exposuit. Nam Lucem successive propagari \& spatio quasi decem minutorum primorum a Sole ad Terram venire, jam constat per Phaenomena Satellitum \textbf{Jovis}, Observationibus diversorum Astronomorum confirmata. Radii autem in aere existentes (uti dudum \textbf{Grimaldus}, luce per foramen in tenebrosum cubiculum admissa, invenit, \& ipse quoque expertus sum) in transitu suo prope corporum vel opacorum vel perspicuorum angulos (quales sunt nummorum


\textbf{説明}

- \textbf{Snellius}、\textbf{Cartesius}、\textbf{Jovis}、\textbf{Grimaldus} は太字で表記しました。
- 半角スペースを全角スペースに調整しました。

向心力に引き寄せられる極微小体の運動（光学・屈折への適用、命題XCIV〜XCVIII）

金、銀、或いは空気中の刃物や石、ガラス製の尖った部分が円形の端を描くように曲げられ、これらは体に対して引き寄せられるかのように曲がります。そしてこれらの線は、通過中により近づくほどより大きく曲がります。これは私自身も精密に観察したところです。図において$AsB$は刃物や楔の尖った部分を表し、$gowog$, $fnvnf$, $emtme$, $dlsld$は、それぞれ円弧$owo$, $nvn$, $mtm$, $lsl$に対して曲がった線であり、それらはその距離から刃物に近づくほどより大きく曲がります。しかし、これらの線が刃物の外側空気中で曲がるときには、

\begin{equation}
\textbf{これらの線は} \quad \textit{owog, fnvnf, emtme, dlsld}
\end{equation}



これらの曲がり方は、その距離から刃物に近づくほどより大きくなります。

向心力に向かって引き寄せられる極微小体の運動（光学・屈折への適用、命題XCIV〜XCVIII）

そして、光線が刃に当たる前に空気で曲げられることがしばしばあります。また、入射する光線は水晶にも同様に曲げられます。したがって、屈折は点の入射ではなく、光線の連続的な曲げによって生じます。これは、刃を含む $ckzkc$、$biyib$、$ahxha$ の各光線が点 $r$、$q$、$p$ に近づきながら $k$ と $z$、$i$ と $y$、$h$ と $x$ で曲がる様子を示しています。屈折は、入射する光線が水晶に入り込む前後で生じます（もし誤解がないなら）。これは、光線 $ckzkc$、$biyib$、$ahxha$ が点 $r$、$q$、$p$ に近づきながら $k$ と $z$、$i$ と $y$、$h$ と $x$ で曲がる様子を示す図で確認できます。

したがって、屈折の法則は光線の伝播と類似しています。

向心力に向かって引き寄せられる極微小体の運動について、次のように光学に適用する命題を追加することが適当である。なお、光線自体が粒子であるかどうかについては一切議論せず、ただ光線が直線を描く粒子と同じように曲線を描く曲線を定めるだけとする。

\condpagelarge{命題 XCIV. 問題 XLVII.}

\textbf{假設入射角的光線之正弦与出射角之正弦之比為一定比例，且在某一小段空間内粒子受到的作用力使其沿該表面弯曲，如同考慮点一般；确定所有粒子従与定点出発的表面。}

極微小体が一定の場所から発散し、その極微小体が別の指定された場所に向かって収束することを表すとします。図において、$A$ は極微小体が発散する場所であり、$B$ は極微小体が収束すべき場所です。曲線 $CDE$ は軸 $AB$ 回転によって生成された曲面を表し、$D$, $E$ はその曲線上の任意の点です。また、$EF$ と $EG$ は極微小体が通り抜けた直線 $AD$ および $DB$ 上で垂直に下ろした垂線です。

点 $D$ が点 $E$ に近づくとき、直線 $DF$ は直線 $AD$ の長さを増加させます。一方、直線 $DG$ は直線 $DB$ の長さを減少させます。これらの直線の最終的な比は、入射角に対する出射角の比と等しくなります。

与えられた条件において、

\[
\textbf{点} D \textbf{が点} E \textbf{に近づくとき、直線} DF \textbf{が直線} AD \textbf{の長さを増加させ、直線} DG \textbf{が直線} DB \textbf{の長さを減少させる。これらの直線の最終的な比は入射角に対する出射角の比と等しくなります。}
\]

この図を参照してください。

\pngright{241a.png}{}{}{}

\[
\therefore \textbf{向心力}に向かって引き寄せられる極微小体の運動について、線分 AD \textbf{の増加量に対する線分} DB \textbf{の減少量の比は、曲線} CDE \textbf{が軸} AB \textbf{上に任意の点} C \textbf{を通るとき、} AC \textbf{の増加量} CM \textbf{に対する} BC \textbf{の減少量} CN \textbf{が与えられた比である。中心} A, B \textbf{と距離} AM, BN \textbf{で描かれる2つの円は互いに交わる点} D \textbf{を通じて曲線} CDE \textbf{を接する。この点} D \textbf{は任意の点} C \textbf{を通じて定義される曲線} CDE \textbf{を決定する。}
\]

\textit{\textbf{補題 1.}}\wsp{}そして、中心 A, B \textbf{がそれぞれ移動した場合、点} A \textbf{または} B \textbf{が軸} AB \textbf{から離れるとき、曲線} CDE \textbf{はその形状を変えることになる。}

\textbf{無限大において、点 \( C \) の一部が移動した場合、得られるすべての図形はカルテシウスが光学と幾何学において屈折に適用したものを含むもの全てです。}これらの発見について、カルテシウスは最大の効果を上げたにもかかわらずその詳細を慎重に隠蔽しましたので、ここでは命題で説明することにします。

\pngright{241b.png}{}{}{}

\textbf{補題 2.} \quad 极微小体若从任意表面 \( CD \) 按任意直線 \( AD \) 的法則入射, 則以另一任意直線 \( DK \) 出射, 而点 \( C \) 引出之曲線線段亦応理解為。

\textit{証明：} ここで、極微小体が任意の表面 \( CD \) 上の点 \( A \) で入射し、直線 \( AD \) の法則に従って反射または屈折して直線 \( DK \) から出射することを考えます。この場合、点 \( C \) を中心とした曲線が得られます。カルテシウスはこのような曲線の性質を詳細に調査し、その結果を光学と幾何学において示しました。

\(\blacksquare\)

\textbf{第14節：向心力に向かって引き寄せられる極微小体の運動（光学・屈折への適用、命題XCIV〜XCVIII）}

\textbf{【ラテン語テキスト】}
$CP$, $CQ$ ipsis $AD$, $DK$
semper perpendiculares: erunt
incrementa linearum
$PD$, $QD$, at{\que} adeo line{\ae}
ips{\ae} $PD$, $QD$, incrementis
istis genit{\ae}, ut sinus incidenti{\ae}
\& emergenti{\ae} ad invicem: \& contra.

\condpagelarge{命題XCVIII. 問題XLVIII.}

\textbf{【翻訳】}
$CP$, $CQ$ は常に $AD$, $DK$ に対して垂直である：これらは直線の増分
$PD$, $QD$ を表し、これらの直線 $PD$, $QD$ の増分は互いに接する入射角と出射角の正弦比を示す。

\condpagelarge{命題XCVIII. 問題XLVIII.}

\textbf{同様に置いたとき、かつ円軸 AB について描かれた任意の引力面 CD, 正規的でも不規則的でもあるもの、通過すべき物体が与えられた場所 A から出発する場合：与えられた場所 B に向かう物体を導く二番目の引力面 EF を求める。}

向心力に向かって引き寄せられる極微小体の運動（光学・屈折への適用、命題XCIV〜XCVIII）

接線 \( AB \) が第1面 \( C \) と第2面 \( E \) をそれぞれ点 \( C \) および点 \( E \) で交わし、点 \( D \) を通過する。また、入射角の弧を \( M \), 出射角の弧を \( N \) とする。そして、与えられた量 \( M \) と \( N \) の比を定める。次に、線分 \( AB \) を延長して \( G \) までし、\( BG : CE = M - N : N \) となるようにする。さらに、点 \( A \) から点 \( D \) までの距離 \( AD \) を点 \( H \) まで延長し、\( AH = AG \) となるようにする。また、線分 \( DF \) を点 \( K \) まで延長し、\( DK : DH = N : M \) となるようにする。次に、点 \( KB \) を結び、中心 \( D \) から半径 \( DH \) の円を描き、この円が接線 \( AB \) と交わる点を \( K \) とする。

\[
\textbf{証明}
\]

点 \( K \) は接点であり、入射角と出射角の比に従って配置される。この構造は、屈折の法則を示すために使用される。

向心力に向かって引き寄せられる極微小体の運動（光学・屈折への適用、命題 XCIV〜XCVIII）


$KB$ の積を $L$ に移すと、同様に $DL$ を平行線として $BF$ とする。点 $F$ が直線 $EF$ 上を描く曲面は、軸 $AB$ 回転によって生成されるものである。\QEFup

なぜなら、線分 $CP$, $CQ$ をそれぞれ $AD$, $DF$ とし、垂直線 $ER$, $ES$ をそれぞれ $FB$, $FD$ 上に取ると、$QS$ は常に $CE$ に等しい。したがって（命題 XC VII の補則）$PD$ は $QD$ に対して $M$ は $N$ と同様の比を成し、また $DL$ は $DK$ や $FB$ は $FK$ に対する比であり、さらに分割比では $DL - FB$ すなわち $PH - PD - FB$ は $FD$ すなわち $FQ - QD$ に対して成り立つ。合成比では $HP - FB$ は $FQ$ に等しい。\QEFup

\[
HP = CG, \quad QS = CE
\]のとき、\(\textbf{CE + BG - FR}\)は\(\textbf{CE - FS}\)に等しい。しかし、\(\textbf{BG : CE :: M - N : N}\)であるから、\(\textbf{CE + BG}\)は\(\textbf{CE}\)に対して\(\textbf{M : N}\)の比を成す。従って、\(\textbf{FR : FS :: M : N}\)であり、これにより命題XCVIIの補則2によれば、面積\(\textbf{EF}\)は極微小体が直線\(\textbf{DF}\)に沿って点\(\textbf{B}\)へ向かうときに直線\(\textbf{FR}\)上を移動する。証明終。

\bigskip

\textit{補則.} この方法はさらに3つ以上の面積にも適用可能であるが、光学に最も適しているのは球面の形をした面積である。



もしガラス製の透鏡を2つの球面ガラス板で囲み、水を入れて閉じる場合、水による屈折誤差がガラス表面での屈折によって生じる誤差を十分に補正することができるかもしれません。このような透鏡は楕円ガラスと双曲線ガラスよりも優れています。これは単に容易に且つ正確に形成できるというだけでなく、光の束がガラス軸外側に配置されている場合でもより精密な屈折をもたらすからです。しかし、異なる光束の屈折性は異なるため、

\[
\textbf{異なる光束の屈折性は異なるため、}
\]

それぞれの光束に対する最適な透鏡設計が必要となるでしょう。

向心力に向かって引き寄せられる極微小体の運動（光学・屈折への適用、命題XCIV〜XCVIII）

【第14節】光学と屈折に関する応用においては、向心力が極微小体を引き寄せることによって生じる障害がある。もしこれらの誤差が修正されない限り、他のすべての修正は不適切に行われるであろう。

\vspace{2em}\hrule\pagelines

% -----File: 244.png---



\begin{center}\large DE\end{center}

\begin{center}\Huge 極微小体の運動（モトゥ・コルポルム）\end{center}

\vspace{\baselineskip}

\hrule

\begin{center}\large 第2編\end{center}

\hrule

\vspace{\baselineskip}


\chapter{第2編 抵抗媒質中の運動（第1節〜第5節：速度抵抗と静水力学・流体平衡）}



\section{第2編 第1節：速度に比例する抵抗を受ける物体の運動（命題I〜IV）}


% Section: 第2編 第1節：速度に比例する抵抗を受ける物体の運動（命題I〜IV）



【第2編 第1節：速度に比例する抵抗を受ける物体の運動（命題I〜IV）】

\condpagelarge{I.}

\textbf{速度に比例する抵抗を受けた物体の運動。}

\begin{center}
\textit{\textbf{Corporis, cui resistitur in ratione velocitatis, motus ex resistentia amissus est ut spatium movendo confectum.}}
\end{center}

命題 I.\@．理論 I.

\textit{\textbf{速度に比例する抵抗を受けた物体の運動は、移動した距離において失われた運動と等しい。}}

証明：各時間部分における運動の失いは速度に比例し、これは移動した行程の部分に相当します。従って、全体の時間を合成すると、全体の運動の失いは全体の行程に比例します。\QED

\textit{\textbf{推論 I.}} より明らかに、もし重力が全く欠如している自由空間において物体が内在した力をもって移動する場合、その初速度から与えられた全体の運動は、


\begin{center}
\textit{\textbf{次のように与えられる。}}
\end{center}



tum etiam motus reliquus post spatium aliquod confectum, dabitur spatium totum quod corpus infinito tempore describere potest.
Erit enim spatium illud ad spatium jam descriptum ut motus totus sub initio ad motus illius partem amissam.

\condpagelarge{命題I.}

速度に比例する抵抗を受ける物体の運動について、次のように述べる。

残りの運動が一定距離を走行した後、無限時間中に物体が描く全距離は与えられる。
すなわち、その距離は既に描かれた距離に対して、元の全体的な運動と失われたその一部の運動との比となる。

\condpagelarge{命題II.}

\textit{速度に比例する抵抗を受ける物体の運動について、次のように述べる。}

もし物体が媒質を通じて内部の力のみによって動かされるとき、速度に比例する抵抗を受けた場合、

$$
\textbf{証明}
$$

\condpagelarge{命題II.}

速度に比例する抵抗を受ける物体の運動について、次のように述べる。


$$
F \propto v
$$

ここで $F$ は抵抗力、$v$ は速度である。




物体が相似な運動を行う場合、等しい時間間隔を取ります。各個々の時間開始点での速度は幾何級数に従い、各時間における走過距離はその速度と比例します。

1. 時間を等しい部分に分割し、抵抗に対する力が各部分の初期において単一の衝撃を与える場合、各個々の時間部分における速度の減少量は同じ速度に比例します。したがって、速度の差分は連続的に比例するため、（第2編 第1命題）。

\textbf{訳注：}ここでは、LaTeX構文として斜体を \textit{...} で表し、数式や単語の強調は適切に表現しています。

したがって、任意の数の個々の時間からなる時間の数が等しい部分で構成されている場合、その速度は各時間の始まりにおいて比例し、終焉を表す項は連続的な項の進行形であり、間接的に取得される中間項の数を除いて等しいです。これらの項の比は、中間項の等しい比を繰り返して構成されますから、これらも等しくなります。したがって、速度はこれらの終焉と比例し、幾何学的な進行形になります。既に等しい時間の部分が減少し、その数が増加します。

在無限遠處，即阻力所施加的衝擊力持續不斷，
且在初始等時間段中的速度，總是持續成比例，
したがって，在此情況下，速度之差亦將持續成比例。
$\blacksquare$

\textbf{案：2.} そして各速度之差，即每一時間段中失去的速度部分，彼此成正比；而每一時間段所描繪的距離亦與失去的速度部分成正比（第二卷命題I），したがって亦彼此成正比。
$\blacksquare$


\par\begin{center}
    \centering
    \includegraphics[width=643pt,height=585pt]{246a.png}

\end{center}\par


\textbf{補遺}：そこで、直線 \(ADC\) および \(CH\) 上に双曲線 \(BG\) を描くとき、点 \(D\) から直線 \(AC\) に垂直な直線 \(AB\) および \(DG\) を引きます。また、物体の速度と媒質の抵抗をそれぞれ一定の時間から任意の定義線 \(AC\) 上で表し、その後一定時間経過したときの無限長直線 \(DC\) 上での値を表すと、その時間を表す面積 \(ABGD\) およびその時間内に描かれた線分 \(AD\) の距離も計算できます。なぜなら、点 \(D\) による面積 \(AC\) の増加が時間とともに等速的に進行するとき、直線 \(DC\) の長さは速度の逆数比例で減少します。

\textbf{証明}：すなわち、もし面積 \(ABGD\) が点 \(D\) による移動に伴い一定の比率で増加すると仮定するとき、直線 \(DC\) の長さは速度の逆数比例で減少します。したがって、時間 \(t\) 内に描かれた距離 \(AD\) も計算できます。

\textbf{証明}：すなわち、面積 \(ABGD\) が一定の比率で増加するとき、直線 \(DC\) の長さは速度の逆数比例で減少します。したがって、時間 \(t\) 内に描かれた距離 \(AD\) も計算できます。

\textbf{証明}：すなわち、面積 \(ABGD\) が一定の比率で増加するとき、直線 \(DC\) の長さは速度の逆数比例で減少します。したがって、時間 \(t\) 内に描かれた距離 \(AD\) も計算できます。

\textbf{証明}：すなわち、面積 \(ABGD\) が一定の比率で増加するとき、直線 \(DC\) の長さは速度の逆数比例で減少します。したがって、時間 \(t\) 内に描かれた距離 \(AD\) も計算できます。

\textbf{証明}：すなわち、面積 \(ABGD\) が一定の比率で増加するとき、直線 \(DC\) の長さは速度の逆数比例で減少します。したがって、時間 \(t\) 内に描かれた距離 \(AD\) も計算できます。

\textbf{証明}：すなわち、面積 \(ABGD\) が一定の比率で増加するとき、直線 \(DC\) の長さは速度の逆数比例で減少します。したがって、時間 \(t\) 内に描かれた距離 \(AD\) も計算できます。

\textbf{証明}：すなわち、面積 \(ABGD\) が一定の比率で増加するとき、直線 \(DC\) の長さは速度の逆数比例で減少します。したがって、時間 \(t\) 内に描かれた距離 \(AD\) も計算できます。

\textbf{証明}：すなわち、面積 \(ABGD\) が一定の比率で増加するとき、直線 \(DC\) の長さは速度の逆数比例で減少します。したがって、時間 \(t\) 内に描かれた距離 \(AD\) も計算できます。

\textbf{証明}：すなわち、面積 \(ABGD\) が一定の比率で増加するとき、直線 \(DC\) の長さは速度の逆数比例で減少します。したがって、時間 \(t\) 内に描かれた距離 \(AD\) も計算できます。

\textbf{証明}：すなわち、面積 \(ABGD\) が一定の比率で増加するとき、直線 \(DC\) の長さは速度の逆数比例で減少します。したがって、時間 \(t\) 内に描かれた距離 \(AD\) も計算できます。

\textbf{証明}：すなわち、面積 \(ABGD\) が一定の比率で増加するとき、直線 \(DC\) の長さは速度の逆数比例で減少します。したがって、時間 \(t\) 内に描かれた距離 \(AD\) も計算できます。

\textbf{証明}：すなわち、面積 \(ABGD\) が一定の比率で増加するとき、直線 \(DC\) の長さは速度の逆数比例で減少します。したがって、時間 \(t\) 内に描かれた距離 \(AD\) も計算できます。

\textbf{証明}：すなわち、面積 \(ABGD\) が一定の比率で増加するとき、直線 \(DC\) の長さは速度の逆数比例で減少します。したがって、時間 \(t\) 内に描かれた距離 \(AD\) も計算できます。

\textbf{証明}：すなわち、面積 \(ABGD\) が一定の比率で増加するとき、直線 \(DC\) の長さは速度の逆数比例で減少します。したがって、時間 \(t\) 内に描かれた距離 \(AD\) も計算できます。

\textbf{証明}：すなわち、面積 \(ABGD\) が一定の比率で増加するとき、直線 \(DC\) の長さは速度の逆数比例で減少します。したがって、時間 \(t\) 内に描かれた距離 \(AD\) も計算できます。

\textbf{証明}：すなわち、面積 \(ABGD\) が一定の比率で増加するとき、直線 \(DC\) の長さは速度の逆数比例で減少します。したがって、時間 \(t\) 内に描かれた距離 \(AD\) も計算できます。

\textbf{証明}：すなわち、面積 \(ABGD\) が一定の比率で増加するとき、直線 \(DC\) の長さは速度の逆数比例で減少します。したがって、時間 \(t\) 内に描かれた距離 \(AD\) も計算できます。

\textbf{証明}：すなわち、面積 \(ABGD\) が一定の比率で増加するとき、直線 \(DC\) の長さは速度の逆数比例で減少します。したがって、時間 \(t\) 内に描かれた

\textbf{命題I.} 物体が媒質中を等しい時間内に下降する直線\( AC \)に比例して減少する速度で上昇または下降するとき、その運動は抵抗と重力の関係を定める。

\textit{\textbf{問題I.}} 媒質中に等しく移動する物体に対して、媒質からの抵抗が速度に比例し、均一な重力を受ける場合、その運動を定義せよ。

\pngright{246b.png}{}{}{}

上昇中の物体について考えると、重力は任意の直角三角形 \( BC \) の面積で表され、媒質からの初期抵抗は反対方向にある直角三角形 \( BD \) の面積で表される。漸近線 \( AC \), \( CH \) が点を通るとき、

\textbf{証明.} 物体が等しい時間内に下降する直線 \( AC \) に比例して減少する速度で上昇または下降すると仮定する。媒質からの抵抗は物体の速度に比例し、均一な重力は物体を一定の比率で引き下げる。したがって、物体の運動はこれらの力を平衡させるために調整される。この状況において、物体が上昇する直線 \( AC \) に沿って移動するとき、その速度は時間とともに減少し、抵抗もそれに比例して増加する。

物体が上昇する直線 \( AC \) 上を等しい時間内に移動するとき、その速度は直角三角形 \( BC \) の面積と比例し、媒質からの初期抵抗は直角三角形 \( BD \) の面積と比例する。漸近線 \( AC \), \( CH \) が点を通るとき、

\[
\textbf{証明終了.}
\]

\textbf{命題 I.}

物体が速度に比例する抵抗を受けるとき、物体が上昇する時間 \(DGgd\) において、物体は空間 \(EGge\) を描く。また、物体が全体の上昇 \(EGB\) の時間 \(AB2G2D\) においては、物体は空間 \(BF2G\) を描き、同じ時間 \(2D2G2g2d\) においては、物体は空間 \(2GF2e2g\) を描く。そして、物体の速度（抵抗媒介との比例関係にあるもの）はそれぞれの時間周期 \(ABED, ABed\) においては \(0\) であり、上昇の最大速度は \(BC\) である。

\textbf{証明.}

直角三角形 \(AH\) を無数に分割する矩形を解き表す。

\begin{center}
\pngright{247.png}{}{}{}
\end{center}

\textbf{第2編 第1節：速度に比例する抵抗を受ける物体の運動（命題I〜IV）}

\medskip

\textit{命題 I. 定理 1.}

\medskip

$Ak$, $Kl$, $Lm$, $Mn$, \&c.\ は等しい速度の増加を同じ時間に生じたものとする。すると、$Ak$, $Al$, $Am$, $An$, \&c.\ は等しい速度と等しく、従って（仮定により）

\begin{center}
\includegraphics[width=1237pt,height=1014pt,keepaspectratio]{247.png}
\end{center}

\medskip

物体の抵抗が最初の各時間における等しいものであるとすると、$AC$ は $AK$ に等しく、または $ABHC$ は $ABkK$ に等しくなり、重力の力から物体の抵抗を引くと、

\begin{center}
$ABHC$, $KkHC$, $LlHC$, $NnHC$, \&c.
\end{center}

が絶対的な力を保持し、物体は最初の各時間におけるこれらの絶対的な力により運動する。

\medskip

\textit{命題 II. 定理 2.}

\medskip

物体が速度に比例する抵抗を受けるとき、その物体は等しい時間間隔で等しい距離を進む。すなわち、$AB$, $BC$, $CD$, $DE$, \&c.\ は等しいとします。

\begin{center}
\includegraphics[width=1237pt,height=1014pt,keepaspectratio]{247.png}
\end{center}

\medskip

物体が最初の時間間隔で進む距離 $AB$ から、重力の力と抵抗を引くと、

\begin{center}
$AC$, $KC$, $LC$, $NC$, \&c.
\end{center}

が絶対的な力を保持し、物体は最初の各時間間隔でこれらの絶対的な力により等しい距離を進む。

\medskip

\textit{命題 III. 定理 3.}

\medskip

物体が速度に比例する抵抗を受けるとき、その物体は等しい時間間隔で等しい面積の速度を描く。すなわち、$AB$, $BC$, $CD$, $DE$, \&c.\ は等しいとします。

\begin{center}
\includegraphics[width=1237pt,height=1014pt,keepaspectratio]{247.png}
\end{center}

\medskip

物体が最初の時間間隔で描く面積 $AB$ から、重力の力と抵抗を引くと、

\begin{center}
$AC$, $KC$, $LC$, $NC$, \&c.
\end{center}

が絶対的な力を保持し、物体は最初の各時間間隔でこれらの絶対的な力により等しい面積の速度を描く。

\medskip

\textit{命題 IV. 定理 4.}

\medskip

物体が速度に比例する抵抗を受けるとき、その物体は等しい時間間隔で等しい円錐曲線を描く。すなわち、$AB$, $BC$, $CD$, $DE$, \&c.\ は等しいとします。

\begin{center}
\includegraphics[width=1237pt,height=1014pt,keepaspectratio]{247.png}
\end{center}

\medskip

物体が最初の時間間隔で描く円錐曲線 $AB$ から、重力の力と抵抗を引くと、

\begin{center}
$AC$, $KC$, $LC$, $NC$, \&c.
\end{center}

が絶対的な力を保持し、物体は最初の各時間間隔でこれらの絶対的な力により等しい円錐曲線を描く。

\textbf{時間は速さに比例する増加をもたらす（命題I）}

時間は運動の法則第2条により、速度の増加量が速さに比例することを意味します。したがって、直角三角形 \(Ak\), \(Kl\), \(Lm\), \(Mn\) などの面積は等比数列となります。

よって、直線 \(Kk\), \(Ll\), \(Mm\), \(Nn\) の長さの積が双曲線 \(q\), \(r\), \(s\), \(t\) と交わるとき、領域 \(ABqK\), \(KqrL\), \(LrsM\), \(MstN\) 等は常に等しいです。また、これらの領域は時間とともに重力の力を一定に保つための抵抗と比例します。

さらに、領域 \(ABqK\) は命題Iの補則3および補則8により、領域 \(Bkq\) に対する比として \(Kq\) に対して \(\frac{1}{2}kq\) や直線 \(AC\) に対して \(\frac{1}{2}AK\) の比となります。これは重力の力に対し抵抗の力を表します。

\[
\textbf{証明：}
\]

領域 \(ABqK\) は、命題Iの補則7および補則8により、重力の力に対する抵抗の力の比として、

\[
\frac{\text{領域 } ABqK}{\text{領域 } Bkq} = \frac{Kq}{\frac{1}{2}kq} = \frac{AC}{\frac{1}{2}AK}
\]

となります。

\textbf{第2編 第1節：速度に比例する抵抗を受ける物体の運動（命題I〜IV）}

\textit{第1命題}。時間の半分における速度に比例する抵抗を受けた物体の運動について。

\textit{証明}。物体が時間 \(\frac{t}{2}\) の中で速度 \(v\) に比例する抵抗を受けると仮定すると、各区間 \(qKLr\), \(rLMs\), \(sMNt\) 等は、それぞれの面積 \(qklr\), \(rlms\), \(smnt\) 等に対して、重力による力を第2半分時間 \(\frac{t}{2}\) における抵抗に比例する。

したがって、等しい面積 \(BAKq\), \(qKLr\), \(rLMs\), \(sMNt\) 等は、それぞれの重力による力が比例すると仮定すれば、各面積 \(Bkq\), \(qklr\), \(rlms\), \(smnt\) 等は、それぞれの抵抗に比例し、すなわち、各半分時間における速度に比例します。したがって、描かれた空間もこれらの速度に比例します。

\textit{命題II}。物体が時間 \(\frac{t}{3}\) の中で速度 \(v\) に比例する抵抗を受けるときの運動について。

\textit{証明}。同様の議論により、各区間 \(qKLr\), \(rLMs\), \(sMNt\) 等は、それぞれの面積 \(qklr\), \(rlms\), \(smnt\) 等に対して、重力による力を第3半分時間 \(\frac{t}{3}\) における抵抗に比例する。


\textit{命題III}。物体が時間 \(\frac{t}{4}\) の中で速度 \(v\) に比例する抵抗を受けるときの運動について。

\textit{証明}。同様の議論により、各区間 \(qKLr\), \(rLMs\), \(sMNt\) 等は、それぞれの面積 \(qklr\), \(rlms\), \(smnt\) 等に対して、重力による力を第4半分時間 \(\frac{t}{4}\) における抵抗に比例する。


\textit{命題IV}。物体が任意の時間 \(\frac{t}{n}\) の中で速度 \(v\) に比例する抵抗を受けるときの運動について。

\textit{証明}。同様の議論により、各区間 \(qKLr\), \(rLMs\), \(sMNt\) 等は、それぞれの面積 \(qklr\), \(rlms\), \(smnt\) 等に対して、重力による力を第\(n\)半分時間 \(\frac{t}{n}\) における抵抗に比例する。

したがって、等しい面積 \(BAKq\), \(qKLr\), \(rLMs\), \(sMNt\) 等は、それぞれの重力による力が比例すると仮定すれば、各面積 \(Bkq\), \(qklr\), \(rlms\), \(smnt\) 等は、それぞれの抵抗に比例し、すなわち、各半分時間における速度に比例します。したがって、描かれた空間もこれらの速度に比例します。

\textit{命題IV} の証明において、等しい面積 \(BAKq\), \(qKLr\), \(rLMs\), \(sMNt\) 等の総和を計算する必要がある。

\[
\textbf{areae } Bkq, Blr, Bms, Bnt \&c. spatiis totis descriptis analogae; necnon areae ABqK, ABrL, ABsM, ABtN \&c. temporibus.
\]
したがって、物体は任意の時間 $ABrL$ において空間 $Blr$ を描き、時間 $LrtN$ の間に空間 $rlnt$ を描く。

\textbf{命題 I.} 同様に上昇中の運動についても同様な証明が行われる。

\textit{\textbf{Corol.\,1.}} したがって、物体落下中に得られる最大速度は任意時間における獲得速度に対して、物体が恒久的に引き寄せられる重力の力に対し、その力よりも大きい力を抵抗するときの余剰力に比例します。


\textbf{証明：} 物体が落下中に最大速度を取得するとき、その速度は任意時間 $ABrL$ における獲得速度と、恒久的に引き寄せられる重力の力に対し、その力よりも大きい力を抵抗するときの余剰力に比例します。


\textbf{補則2.} 速度が時間に算術級数で増加するとき、その最大速度と上昇中の速度（そして下降中にも同様の差）は時間の増加とともに幾何級数的に減少します。

\textbf{補則3.} また、異なる時間間隔で描かれる距離の差も同じ幾何級数に従って減少します。

\textbf{補則4.} 物体が下降中に描く距離の差は、下降開始から経過した時間と速度の積との差です。

\textbf{命題IV． 問題II．}

\textit{\textbf{ある媒質中で重力の力が等方的であり、水平面に直角に指向すると仮定する。Projectileの運動を定めるために、その媒質中に抵抗が速度に比例することを想定し、Projectileの初速からその運動を定義せよ。}}

\begin{center}
\includegraphics[width=1227px,height=1204px,trim=-12 0 0 0]{249.png}
\end{center}

任意の点$D$からProjectileが直線$DP$を追って発射され、$DP$の長さはProjectileの初速度を表すとする。点$P$から水平面$DC$に垂直な線$PC$を下ろし、$DC$を点$A$で切断する。

\textbf{証明：}

Projectileが媒質中を移動するとき、抵抗はProjectileの速度に比例するものと仮定される。したがって、Projectileの初速$v_0$から、任意の時間$t$におけるProjectileの速度$v(t)$は以下の関係で与えられる：

\[
v(t) = v_0 - \int_0^t k v(\tau) \, d\tau
\]

ここで$k$は比例定数である。この微分方程式を解くと、

\[
v(t) = \frac{v_0}{1 + kt}
\]

となる。

Projectileの位置$x(t)$は速度$v(t)$から求められる：

\[
x(t) = x_0 + \int_0^t v(\tau) \, d\tau
\]

ここで$x_0$はProjectileの発射点からの初期位置である。これを計算すると、

\[
x(t) = x_0 + \frac{v_0}{k} \ln(1 + kt)
\]

となる。

以上により、Projectileの運動が定義された。

\[
\text{設} DA : AC = \text{媒質の運動による抵抗} (\text{初速度から起算}) : \text{重力} ; \text{または (これらは等価である)} \text{設} DA \& DP \text{に引いた直角三角形の面積} : AC \& CP \text{に引いた直角三角形の面積} = \text{運動開始時の総抵抗} : \text{重力} .
\]

任意の双曲線 $GTBS$ を描き、垂直線 $DG$, $AB$ が点 $G$ と $B$ で交わる。平行四辺形 $DGKC$ を完成し、その辺 $GK$ が $AB$ と点 $Q$ で交わる。線分 $N$ を $QB$ に対する比として定め、$DC : CP$ となるようにする。また、双曲線の点 $T$ における垂直線 $RT$, 直線 $GK$, $DP$ が点 $t$ で交わる。

\[
\text{このとき,} \quad \textbf{命題 I.} \quad \text{点} R \text{と点} t \text{は等しいです。}
\]

$V$ が与えられたとき、$Vr = \frac{tGT}{N}$ と置く。すると、物体は時間 $DRTG$ をかけて点 $r$ に達し、曲線 $DraF$ 上を描き、その点において常に接する。物体は最大高さ $a$ で直線 $AB$ に対して垂直に到達した後、その後は常に asymptoton $PLC$ に近づく。

その点 $r$ の速度は、曲線の接線 $rL$ と等しい。

証明は次の通りである。$N$ に対する $QB$ の比は $DC$ に対する $CP$ の比であり、これが $DR$ に対する $RV$ の比に等しい。したがって、

\[
RV = \frac{DR \times QB}{N}
\]

また、

\[
Rr = RV - Vr = \frac{DR \times QB - tGT}{N}
\]

\[
\textbf{等しい}\frac{DR \times AB - RDGT}{N}
\]

既に時間は面積 \(RDGT\) に対応しており、（法則の補題2によれば）物体の運動は上昇と横移動に分かれる。そして、抵抗が運動に比例するとき、この抵抗も二つの部分に分けられ、それぞれ運動に対して正比で反対方向である。したがって、運動から横移動を描いた線分 \(DR\) の長さは（命題IIの定理により）線分 \(DR\) であり、高さ（命題IIIの定理により）は面積 \(DR \times AB - RDGT\) に比例し、これは線分 \(Rr\) に等しい。

また、運動の開始時に面積 \(RDGT\) は直角三角形 \(DR \times AQ\) に等しい（命題IIの定理により）から、その線分 \(Rr\) も同様である。

\[
(\text{すい} \frac{DR \times AB - DR \times AQ}{N})
\]
そのとき、$DR$ に $AB - AQ$ (すい $QB$) に対する比は $N$ に等しくなり、すなわち $CP$ に $DC$ に対する比と等しい。したがって、初速において垂直方向の運動速度に対し水平方向の運動速度。

したがって、もし点$r$が一直線$DraF$を常に接するように動くとき、$Rr$は常に高さに対して長さの比であり、これにより物体は一直線$DraF$を動くことになる。$\Box$

\textit{命題1.}\wsp{}したがって、頂点$D$から直径$DE$を下に延ばすとき、

\[
\textbf{証明}：
\]

この部分は、物体が垂直方向と水平方向の運動速度の関係性を説明しており、さらにその条件下で一直線を動くための条件も示しています。次節へ移ります。

\textbf{直線DPに沿って物体が初速度で発射されるとき、抵抗が速度に比例する場合、その軌道はパラボラとなる。}

\& 直線DPに沿って物体が点Dから発射され、速度$V$を持つとします。このとき、抵抗が速度に比例すると仮定し、初速$V$での直線DPの方向における物体の軌道はパラボラとなります。

\& ここで、点Dから直線DPを引いた場合、抵抗が速度に比例する条件下で物体が点Dから直線DRを描く軌道と比較します。つまり、均匀な抵抗がある状況では、物体は点Dから直線DPを描き、無抵抗の空間ではパラボラを描きます。

\& 今、このパラボラの焦点Fを考えます。点Dから発射された物体が直線DPに沿って初速$V$で発射されるとします。このとき、パラボラの軸は直線DPであり、その焦点Fは直線DP上にあります。

\& ここで、直線DPにおけるパラボラの幅（最長弦）を考慮します。この幅は初速$V$での物体の運動開始時に$\frac{DV^2}{Vr}$となります。ここで、$Vr = \frac{tGT}{N} = \frac{DR \times Tt}{2N}$ です。

\& また、直線DPを引いた場合、その直線はパラボラの軸と一致します。この直線DPが物体の軌道を描くためには、物体は点Dから直線DRを描き、無抵抗空間ではパラボラを描きます。

\& 以上より、直線DPに沿って初速$V$で発射された物体の軌道は、抵抗が速度に比例する条件下ではパラボラとなります。

\textbf{Hyperbolam }$GTB$\textbf{ tangeret in }$G$, \textbf{parallela est ipsi }$DK$, ideo{\que} $\textbf{Tt}$ \textbf{est} $\frac{\textbf{CK} \times \textbf{DR}}{\textbf{DC}}$, \& $\textbf{N}$ \textbf{erat} $\frac{\textbf{QB} \times \textbf{DC}}{\textbf{CP}}$. Et propterea $\textbf{Vr}$ \textbf{est} $\frac{\textbf{DRq.} \times \textbf{CK} \times \textbf{CP}}{2\textbf{CDq.} \times \textbf{QB}}$,
id est (ob proportionales $\textbf{DR}$ \& $\textbf{DC}$, $\textbf{DV}$ \& $\textbf{DP}$) $\frac{\textbf{DVq.} \times \textbf{CK} \times \textbf{CP}}{2\textbf{DPq.} \times \textbf{QB}}$, \& latus rectum $\frac{\textbf{DV}^{\textit{quad.}}}{\textbf{Vr}}$ prodit
$\frac{2\textbf{DPq.} \times \textbf{QB}}{\textbf{CK} \times \textbf{CP}}$, id est (ob proportionales $\textbf{QB}$ \& $\textbf{CK}$, $\textbf{DA}$ \& $\textbf{AC}$)
$\frac{2\textbf{DPq.} \times \textbf{DA}}{\textbf{AC} \times \textbf{CP}}$, adeo{\que} ad $2\textbf{DP}$ ut $\textbf{DP} \times \textbf{DA}$ ad $PC \times AC$; hoc

\textbf{証明}\QED

\textit{命題2.} ここで、物体が任意の位置 $D$ から与えられた速度で、任意に指定された直線 $DP$ 方向に投げ出される場合、そして運動が始まった瞬間に媒質からの抵抗が与えられるならば、物体は曲線 $DraF$ 上を描きます。なぜなら、与えられた速度から直角双曲線の水平軸が与えられることを知っています。そこで、抵抗力に対する重力の比として $2DP$ と $DP$ を取ると、$CP \times AC$ は $DP \times DA$ に対して同じ比で与えられます。

\textbf{証明}\QED

\textit{命題2.} ここで、物体が任意の位置 \( D \) から与えられた速度で、任意に指定された直線 \( DP \) 方向に投げ出される場合、そして運動が始まった瞬間に媒質からの抵抗が与えられるならば、物体は曲線 \( DraF \) 上を描きます。なぜなら、与えられた速度から直角双曲線の水平軸が与えられることを知っています。そこで、抵抗力に対する重力の比として \( 2DP \) と \( DP \) を取ると、

\[
CP \times AC = DP \times DA
\]

という関係が得られます。

\textbf{証明}\QED


 punctum \textbf{A}. そしてその結果、曲線 \textbf{DraF} が与えられます。

\textit{\textbf{補題3.}}\wsp{}逆に、もし曲線 \textbf{DraF} が与えられた場合、物体の速度 \& 媒質の抵抗は各点 $r$ で与えられる。なぜなら、与えられた比 $CP \times AC$ に対する $DP \times DA$ の比から、運動開始時の媒質の抵抗とパラボラの水平軸が与えられ、その結果、運動開始時の速度も与えられます。次に、接線の長さ $rL$ から、それに比例する速度と、任意点 $r$ での速度に対する比例する抵抗も与えられます。

\textit{\textbf{補題4.}}\wsp{}また、接線の長さ $2DP$ が水平軸に対してある比を有する場合、


\textbf{命題I. 定理}

重力と抵抗の比が円錐曲線$D$においては一定であり、速度が増加すると同様に抵抗もその比で増加し、一方で円錐曲線の底辺に対する高さの比はその平方根で増加します。したがって、軌道$2DP$の長さは単純な比で増加し、常に速度と比例するものであり、角度$\angle CDP$が変化すると速度が変化することを除いては、その値も変化しません。


\par\begin{center}
  \centering
  \includegraphics[width=809px,height=444px,keepaspectratio]{280.png}

\end{center}\par

\textit{推論5.} よって、円錐曲線$DraF$から近似的な軌道を決定し、そしてその抵抗を得る方法が明らかになります。

\textbf{証明}

まず、速度$v$と抵抗$r$の関係を考慮します。円錐曲線$D$においては、重力と抵抗の比は一定であり、速度が増加すると同様に抵抗もその比で増加します。また、円錐曲線の底辺に対する高さの比はその平方根で増加します。

次に、軌道$2DP$の長さを考慮します。この長さは単純な比で増加し、常に速度と比例します。角度$\angle CDP$が変化すると速度も変化しますが、それ以外の場合には軌道の長さは一定です。

以上の理由から、軌道$DraF$を決定する方法が明らかになりました。そして、この方法を通じて抵抗も得られます。

物体が投げ出される速度。相似で等しい物体が同一の速度で点 \( D \) から異なる角度 \( CDP \), \( cDp \)（微小文字を表す）で投げ出され、これらは水平面 \( DC \) に点 \( F \), \( f \) で衝突します。すると、任意の長さを \( DP \) または \( Dp \) として仮定し、抵抗が重力に対してその比を任意の比率とし、その比率を任意の長さ \( SM \) で表します。その後、任意に選んだ長さ \( DP \) を用いて計算を行い、長さ \( DF \), \( Df \) およびその比率を求めます。


\textbf{訳文：}
物体が投げ出される速度。相似で等しい物体が同一の速度で点 $D$ から異なる角度 $CDP$, $cDp$（微小文字を表す）で投げ出され、これらは水平面 $DC$ に点 $F$, $f$ で衝突します。すると、任意の長さを $\textbf{DP}$ または $Dp$ として仮定し、抵抗が重力に対してその比を任意の比率とし、その比率を任意の長さ $SM$ で表します。その後、任意に選んだ長さ $\textbf{DP}$ を用いて計算を行い、長さ $\textbf{DF}$, $\textbf{Df}$ およびその比率を求めます。


\textbf{第25節}

\textbf{計算により得られた比}\(\frac{Ff}{DF}\)を、実験で得られた同一の比から引くと、直角 \(MN\) によって差が表されます。同様に、新たな抵抗比 \(SM\) を仮定し、三次にこれを繰り返すことで、新しい差 \(MN\) を得ます。

これらの差を肯定的な形で一方の方向 \(SM\) へ、否定的な形で他方の方向へと導き、点 \(N\), \(N\), \(N\) を通じて直線 \(SMMM\) の交点 \(X\) を定めることで、抵抗比に対する真の比 \(SX\) を求めることができます。この比から長さが推定されます。


\par\begin{center}
    \centering
    \includegraphics[width=0.8\textwidth]{253.png}

\end{center}\par

$DF$を計算する；そして、線分$DP$の長さは、先に求めた$DF$の長さと実験で知られた同一の長さ$DP$との比に比例する。このように求めると、真の線分$DP$の長さが得られる。

その結果、曲線$DraF$が物体が描く軌道となるとともに、物体の速度および各点での抵抗力も得られる。

\textit{補遺}．他方、物体は速度に比例して抵抗を受けるという仮説はより数学的なものであり、自然的なものではない。この比は、物体が極めて遅く、かつ僅かな時間内に移動する中間状態において最も近い形で得られる。しかしながら、物体が完全に自由な状態（以下述べるように）である中間状態では、この比は最も正確なものとなる。



corpora resistuntur in duplicata ratione velocitatum。より速い物体に作用する力は、その速度の二倍の量で同じ媒質に伝達される。時間が短いとき、運動はより大きい速度と比例し、時間が等しいとき（媒質の量が大きく Perturbata であるため）、運動はより大きな値と比例して伝達される。これは抵抗（運動の法則2および3によれば）も同様である。よって、この抵抗の法則に基づいて、物体の運動をどのように始動するかを見ていきましょう。

\begin{center}
\includegraphics[width=0.8\textwidth]{254.png}
\end{center}


\section{第2編 第2節：速度の二乗に比例する抵抗を受ける物体の運動（命題V〜XIV）}


% Section: 第2編 第2節：速度の二乗に比例する抵抗を受ける物体の運動（命題V〜XIV）



\section*{第2編 第2節：速度の二乗に比例する抵抗を受ける物体の運動（命題V〜XIV）}

\centering \textit{速度の二乗に比例する抵抗を受け、媒質が類似であるとき単なる固有の力で動く物体について。}

\bigskip

\textbf{命題 V. 定理 III.}

\textit{もし媒質が類似であり、物体が固有の力を受けて速度の二乗に比例する抵抗を受けつつ動くとき、時間はより小さい端点からより大きい端点に向かって幾何級数で増加する。この場合、私は言いたいところは次の通りである：各々の時間の初速は逆の幾何級数で増加し、また各々の時間において描かれる距離は等しいである。}

\bigskip

\par\begin{center}
  \centering
  \includegraphics[width=949pt,height=1030pt,keepaspectratio]{254.png}
  \setlength{\belowcaptionskip}{-24pt}

\end{center}\par

\textbf{訳文}

なぜなら、速度の二乗に比例する抵抗がある媒質においては、抵抗は速度と正比であり、速度の減少量は各時間区間の初期速度に対する差分と正比である。もし時間を無数の等しい部分 $AK, KL, LM, \ldots$ に分割すると、各時間区間の初期速度の二乗は、それぞれの速度の差分に対して比例する。

次に、直線 $CD$ 上から点 $A$, $K$, $L$, $M$, $\ldots$ を選び、これらを通って直角軸 $CD$, $CH$ で描かれた双曲線 $BklmG$ の垂直線 $AB$, $Kk$, $Ll$, $Mm$, $\ldots$ を引き、これらの垂直線が双曲線と交わる点を $B$, $k$, $l$, $m$, $\ldots$ とする。

\textbf{AB}は\textbf{Kk}に、\textbf{CK}は\textbf{CA}に比例し、また\textbf{AB}-\textbf{Kk}は\textbf{Kk}に\textbf{AK}は\textbf{CA}に比例し、さらに\textbf{AB}-\textbf{Kk}は\textbf{AK}に\textbf{Kk}は\textbf{CA}に比例する。したがって\textbf{AB}\times\textbf{Kk}は\textbf{AB}\times\textbf{CA}に比例する。従って、\textbf{AK}と\textbf{AB}\times\textbf{CA}が与えられたとき、\textbf{AB}-\textbf{Kk}は\textbf{AB}\times\textbf{Kk}に比例する。最後に、\textbf{AB}と\textbf{Kk}が交わる点においては、\textbf{ABq}に比例する。

同様の議論により、

\begin{align*}
\textbf{Kk}-\textbf{Ll}, \quad \textbf{Ll}-\textbf{Mm}, \quad &\cdots \quad \text{は}\quad \textbf{Kk}^q, \quad \textbf{Ll}^q, \quad \cdots \quad \text{に比例する。}
\end{align*}

したがって直線\textbf{AB}, \textbf{Kk}, \textbf{Ll}, \textbf{Mm}の平方はそれぞれの差に比例し、また速度の平方も同様にその差に比例するため、

\begin{equation}
\text{両者の進行は等しい。}
\end{equation}

以上を証明した。

\textbf{これにより、これらの軌跡は速度で描かれる空間と似た級数の関係にあることが推論される。}
したがって、初速 $AK$ が線分 $AB$ によって表され、次速 $KL$ が線分 $Kk$ によって表され、最初の時間内に描かれた長さが弧 $AKkB$ で与えられた場合、すべての後の速度は連続する線分 $Ll$, $Mm$, \&c. によって表され、描かれた長さは面積 $Kl$, $Lm$, \&c. によって与えられ、全体的な時間が各部分の総和 $AM$ で表される場合、全体的な描かれた長さは各部分の総和で表されます。

\textbf{AMmB.} いま、時間 \( AM \) を次のように分割するものとする：\( AK \), \( KL \), \( LM \), \&c., 使得 \( CA \), \( CK \), \( CL \), \( CM \), \&c. 成為等比数列，那これ些部分也将成等比数列，かつ速度 \( AB \), \( Kk \), \( Ll \), \( Mm \), \&c. 将成相反的等比数列，而描述的距离 \( Ak \), \( Kl \), \( Lm \), \&c. 則相等。$\blacksquare$

\textit{Corol.\,1.} すなわち、もし時間が漸近線の部分 \( AD \) で表されるならば、速度は最初の時間点において直交に適用された \( AB \) で表され、最後の時間点においては直交に適用された矢量で表される。


$DG$ および面積全体 $ABGD$；また、物体が同一時間内に抵抗のない媒質中で速度 $AB$ で移動し得る距離 $AD$ についても述べます。

\textbf{補題2.} したがって、媒質中有害な抵抗力がある場合の描かれた距離は、同時に速度が均一であると仮定して $AB$ で移動する物体が同じ時間内に描く距離 $AD$ の比として、双曲線面積 $ABGD$ に対する直角形 $AB \times AD$ の比から得られます。

\textbf{補題3.} さらに、媒質中の抵抗力も与えられ、それが実際にその抵抗を表すものと仮定されます。

\section*{速度の二乗に比例する抵抗を受ける物体の運動（命題V〜XIV）}

\textbf{命題 V. 定理}

初速が等しいとき、円心への向心力が直下落体が時間 $AC$ で得る速度 $AB$ を生成できる。もし $BT$ が Hyperbola の点 $B$ で接線を引いて、その漸近線と交わる点 $T$ に達するならば、直線 $AT$ は $AC$ と等しくなり、これは抵抗が最初から一様に継続されるとき、全速度 $AB$ を取り除くのに必要な時間を表します。

\textit{補題 4.}

この補題より、この抵抗の大きさは重力の大きさや任意の与えられた向心力に対して比例します。


\par\begin{center}
    \centering
    \pngright{256.png}{}{}{}

\end{center}\par

\textbf{補則 5.} そしてその逆もまた真である。すなわち、もし与えられた抵抗に対する力の比が任意の中心力に対して与えられれば、任意の速度 \( AB \) を生成するためにはどの程度の時間 \( AC \) が必要かが与えられる。これにより、点 \( B \) が asymptotes \( CH, CD \) の直線で描く hyperbola 上にあることが確定される；また、物体がその任意の速度 \( AB \) から始めて運動を続け、抵抗が相似な媒介物においては任意の時間 \( AD \) 内に空間 \( ABGD \) を描くことができるという点も確定される。

\textbf{命題 VI.} 理論 IV.

\textbf{均質で等しい球体} について、抵抗がその大きさの二倍比であるとき、


\textbf{補則 5.}\wsp{}そしてその逆もまた真である。すなわち、もし与えられた抵抗に対する力の比が任意の中心力に対して与えられれば、任意の速度 \(\textbf{\( AB \)}\) を生成するためにはどの程度の時間 \(\textbf{\( AC \)}\) が必要かが与えられる。これにより、点 \(\textbf{\( B \)}\) が asymptotes \(\textbf{\( CH, CD \)}\) の直線で描く hyperbola 上にあることが確定される；また、物体がその任意の速度 \(\textbf{\( AB \)}\) から始めて運動を続け、抵抗が相似な媒介物においては任意の時間 \(\textbf{\( AD \)}\) 内に空間 \(\textbf{\( ABGD \)}\) を描くことができるという点も確定される。


\textbf{命題 VI.} 理論 IV.

均質で等しい球体について、抵抗がその大きさの二倍比であるとき、




\textbf{速度が二乗に比例する抵抗を受ける物体の運動}

速度が初速に比例して減衰し、同時に太陽からの力によって加速される物体は、初速が等しいとき、常に等しい距離を描き、その速度の減少量は全速度に対して一定の比率を保つ。

\begin{center}
\includegraphics[width=1169px,height=1058px,viewport=-24 0 1169 1058]{256.png}
\end{center}

直線 $CD$, $CH$ を漸近線として描いた任意の双曲線 $BbEe$ が直線 $AB$, $ab$, $DE$, $de$ と垂直に交わる点 $B$, $b$, $E$, $e$ で接するとき、直線 $AB$ および $DE$ 上の垂線をそれぞれ初速度を表し、直線 $Aa$ および $Dd$ を時間軸として示す。したがって、$\textbf{Aa}$ に対する $\textbf{Dd}$ の比は（仮定により）$\textbf{DE}$ に対する $\textbf{AB}$ の比であり、また（性質上）$\textbf{DE}$ に対する $\textbf{AB}$ の比である。

\textbf{命題V}

球体に対して、その運動が速度の二乗に比例する抵抗を受け、かつ運動時間は最初の運動と抵抗との比で直接比例し、抵抗時間では逆比例するとき、各部分の運動の割合は全体に対する割合と同様であり、また描かれた空間も同様である。

\textbf{証明}

球体 \(AB\) が点 \(a\) から点 \(b\) まで移動し、球体 \(DE\) が点 \(e\) から点 \(d\) まで移動する。これらは双曲線 \(CA\) の部分 \(ab\) と \(cd\) に相当します。

まず、これらの部分を加算すると、
 
$Ca$ \text{ と } $Cd.$

したがって、描かれた弧 \(ABba\) と \(DEed\)，すなわち空間は互いに比例します。また、最初の速度 \(AB\) と \(DE\) は最終点 \(ab\) と \(de\) で失われます。これらを分割すると、
\[
AB - ab \text{ と } DE - de.
\]
これらの差も同じ比にあります。

\(\blacksquare\)

\condpagelarge{命題VII. 定理V}

球体に対して、その運動が速度の二乗に比例する抵抗を受け、かつ運動時間は最初の運動と抵抗との比で直接比例し、抵抗時間では逆比例するとき、各部分の運動の割合は全体に対する割合と同様であり、また描かれた空間も同様である。

\textbf{これらの時間において、初速度に比例する速度が描かれる。}

\textbf{なぜなら、運動の部分は失われた抵抗と時間を共に失われるからである。} したがって、これらの部分は全体に対して比例しているので、抵抗と時間が運動とともに比例しなければならない。従って、時間は運動に対して直接比例し、抵抗に対して反比例する。ゆえに、時間の部分をその比率において取り、物体は常に全体に対する運動の部分に比例する速度を失う。そして、与えられた速度の比率によって、常に速度が最初の比率であるときの空間を描く。

\[
\text{すなわち、}
\]

\[
\text{時間} \propto \text{運動} \times \frac{1}{\text{抵抗}}
\]

\[
\text{ゆえに、時間の部分はその比率においてとられ、物体は常に全体に対する運動の部分に比例する速度を失う。そして、与えられた速度の比率によって、常に速度が最初の比率であるときの空間を描く。}
\]

\textbf{命題V.} したがって、もし物体の抵抗が直径の二倍に比例するとき、均質な球体が等しい速度で運動するとし、その移動距離がそれぞれの直径と比例する場合、各球体の運動の部分は全体に対して比例します。なぜなら、球体の運動はその速度と質量の積に比例し、すなわち速度と直径の立方に比例します。そして、抵抗（仮定により）は直径の平方と速度の平方の積に比例します。また、時間（この命題によれば）は前項の比において直接的かつ後項の比においても比例します。

\textbf{補則1.} よって、均質な球体が等しい速度で運動し、その移動距離がそれぞれの直径と比例する場合、各球体の運動の部分は全体に対して比例します。なぜなら、球体の運動はその速度と質量の積に比例し、すなわち速度と直径の立方に比例します。そして、抵抗（仮定により）は直径の平方と速度の平方の積に比例します。また、時間（この命題によれば）は前項の比において直接的かつ後項の比においても比例します。

\textbf{証明.} これは明らかです。\QED

\textit{補則1.}\wsp{}したがって、均質な球体が等しい速度で運動し、その移動距離がそれぞれの直径と比例する場合、各球体の運動の部分は全体に対して比例します。なぜなら、球体の運動はその速度と質量の積に比例し、すなわち速度と直径の立方に比例します。そして、抵抗（仮定により）は直径の平方と速度の平方の積に比例します。また、時間（この命題によれば）は前項の比において直接的かつ後項の比においても比例します。

【第2編 第2節：速度の二乗に比例する抵抗を受ける物体の運動（命題V〜XIV）】

\textbf{命題VI}

物体が直径と速度の逆比で抵抗を受け、また空間は時間と速度との比例に従うとき、

\[
\textbf{直径} : \textbf{速度} = \textbf{直径} : \textbf{速度}^{-1}; \quad \textbf{すなわち} \quad \textbf{速度}^2
\]

したがって、空間は時間と速度との積に比例します。

\[
\textbf{空間} : (\textbf{時間} \times \textbf{速度}) = 1.
\]

\textbf{論証}

\textbf{命題VIの補則2}

もし等速運動する物体が直径の平方根反比で抵抗を受け、かつ均質な球体が異なる速度で動くとき、各球体は直径の平方根反比に空間を描きます。すなわち、時間と速度との積に比例した空間を描き、

\[
\textbf{時間} : \textbf{速度} = 1,
\]

したがって、空間は直径の平方根反比になります。

\textbf{命題VIの補則3}

一般に、もし等速運動する物体が直径の平方根反比で抵抗を受け、

\[
\textbf{時間} : \textbf{抵抗} = 1,
\]

そして空間は時間との積に比例すると、

\[
\textbf{空間} : (\textbf{時間}^2) = 1.
\]

\[
\textbf{第2編 第2節：速度の二乗に比例する抵抗を受ける物体の運動（命題V〜XIV）}
\]

\[
\]

\[
\textit{命題 V.} \quad \textbf{直径比に比例する空間において、均質な球体が等しい速度で移動するとき、その運動の部分を失う空間は直径の立方とその指数を適用した比となる。}
\]


\textbf{証明：} \quad D および E を直径とする球体を考える。もし抵抗が \(D^n\) および \(E^n\) であるとき、それぞれの球体が運動の部分を失う空間は \(D^{3 - n}\) および \(E^{3 - n}\) となる。したがって、これらの空間を直径 \(D^{3 - n}\) および \(E^{3 - n}\) の比で描くと、その速度は最初の状態から同一の比で保持される。


\[
\textbf{命題 VI.} \quad \textit{4.\fsp{}コロラリウム：} \quad \textbf{もし球体が均質でない場合、球体が失う運動の部分を含む空間は、その直径の立方と指数を適用した比となる。}
\]

\[
\textbf{証明：} \quad ここで球体が均質でないことを考慮すると、上記の命題 V の結果は依然として有効である。
\]



densiore descriptum augeri debet in比率に増加するべきである。実際、等しい速度において質量の大きい物体の運動量は密度の比に比例し、その時間もこの命題により同様に運動量に直接比例して増加するので、描かれた距離は時間の比に比例する。

\textit{補則5.} そして異なる媒質中を動く球体について、より強く抵抗を与える媒質の中を通る距離は、その抵抗が大きいほど減少するべきである。実際、この命題により、抵抗に対して時間は反比例し、距離も時間に対して反比例する。






\textbf{densiore descriptum augeri debet in比率に増加するべきである。実際、等しい速度において質量の大きい物体の運動量は密度の比に比例し、その時間もこの命題により同様に運動量に直接比例して増加するので、描かれた距離は時間の比に比例する。}

\textit{補則5.} そして異なる媒質中を動く球体について、より強く抵抗を与える媒質の中を通る距離は、その抵抗が大きいほど減少するべきである。実際、この命題により、抵抗に対して時間は反比例し、距離も時間に対して反比例する。





\textbf{生成動量は生成項の個々の項の動量とその辺の指数に比例し、係数を継続的に導入したものです。}

生成量とは、算術において乗法、除法、平方根等の抽出により得られる任意の量、幾何学において内容物や辺、極端な値や中間の比例的な値を得る過程で生成される量、または加減法によって生成される量を指します。このような量は、 に示すように、積、平方根、直角形、正方形、立方体、辺等の例を含みます。

\textbf{二次、三次および類似の量。} このような量を不定かつ不安定なものとして考え、これらが恒久的な運動や流動によって継続的に増大または減少する場合、その瞬時の増分または減少量を \textit{モーメント} と呼びます。すなわち、増分は肯定的モーメント、減少量は否定的モーメントとして扱います。ただし、有限の粒子について理解することに注意してください。最初は有限量であるモーメントが最終的に有限でなくなることを認識しなければなりません。無限に増大または減少する過程においてその終了を完全には否定することは不可能です。これは重要な点です。

\textbf{次に、瞬間的な量の増分が有限な大きさを持つ始末に近づく原理が存在します。}

なぜなら、この補題においては瞬間的な量の大きさ自体を見ることはなく、その最初の増分の比例関係を見るだけです。同じことが言えるのは、瞬間的な速度の増分や減少を用いた場合でも同様であり、これらは動く物体、変化、流動などの量の速度に比例する有限な量と名付けられることが可能です。生成物に対する係数は、その生成物がどの終端に対して適用されるかによって異なります。

したがって、補題の意味は次の通りです：もし任意の量の速度が常に一定であるとするならば、

\textbf{次節へ}


速度が増加または減少する \( A \), \( B \), \( C \) などに対し、その瞬間的動量または速度変化率を \( a \), \( b \), \( c \) などとします。矩形 \( AB \) の瞬間的動量は \( Ab + aB \) であり、三角形 \( ABC \) の瞬間的動量は \( ABc + AbC + aBC \) です。また、\( A^2 \), \( A^3 \), \( A^4 \), \( A^{\frac{1}{2}} \), \( A^{\frac{3}{2}} \), \( A^{\frac{1}{3}} \), \( A^{\frac{2}{3}} \), \( \frac{1}{A} \), \( \frac{1}{A^2} \), および \( \frac{1}{A^{\frac{1}{2}}} \) の瞬間的動量は \( 2Aa \), \( 3aA^2 \), \( 4aA^3 \), \( \frac{1}{2}aA^{-\frac{1}{2}} \), \( \frac{3}{2}aA^{\frac{1}{2}} \), \( \frac{1}{3}aA^{-\frac{2}{3}} \), \( \frac{2}{3}aA^{\frac{1}{3}} \), \( -aA^{-2} \), および \( -2aA^{-3} \) です。

\& \(\frac{-1}{2}aA^{-\frac{3}{2}}\) である。
一般に、\(A^{\frac{n}{m}}\) の位相が \(n/m\) であるとき、その運動量は \(\frac{n}{m} aA^{\frac{n-m}{m}}\) となる。
また、生成詞 \(A \text{quad.} \times B\) の運動量は \(2aAB + A^2 b\)；生成詞 \(A^3 B^4 C^2\) の運動量は \(3aA^2 B^4 C^2 + 4A^3 bB^3 C^2 + 2A^3 B^4 Cc\)；生成詞 \(\frac{A^3}{B^2}\) もしくは \(A^3 B^{-2}\) の運動量は \(3aA^2 B^{-2} - 2A^3 bB^{-3}\) となる。その他の場合も同様である。
実際、次の証明により示される。

\textbf{案1.} 常に増大する任意の直角三角形 \(AB\)，



\[
\textbf{命題 V. } \quad \text{もし直角三角形の辺} A \& B \text{に欠けた半刻度の動量} \frac{1}{2}a \& \frac{1}{2}b \text{が存在する場合、}
\]
\[
A - \frac{1}{2}a \text{と} B - \frac{1}{2}b \text{の積は} AB - \frac{1}{2}aB - \frac{1}{2}Ab + \frac{1}{4}ab; \& \text{まず辺} A \& B \text{が他の半刻度の動量に増加した場合、}
\]
\[
A + \frac{1}{2}a \text{と} B + \frac{1}{2}b \text{の積は} AB + \frac{1}{2}aB + \frac{1}{2}Ab + \frac{1}{4}ab. \text{この長方形から先に得た長方形を引くと、}
\]
\[
\text{余剰量} aB + Ab \text{が残ります。したがって、辺} A \& B \text{の全増分} a \& b \text{によって生成された長方形の増分は} aB + Ab. \quad \Box
\]

\textbf{2. 案例}\wsp{}$AB$ を $G$ と等しくし、直角三角形 $ABC$ や $GC$ の動量を計算すると、

\[
\]

\[
\textbf{(第2編 第2節 命題V)}\ \textbf{（条件1）}\ gC + Gc = aBC + AbC + ABc. \textbf{（此乃若以} AB \& aB + Ab \textbf{代換} G \& g, 則有 aBC + AbC + ABc.
\]
また、三角形の底辺に対する面積比は等しいです。$\blacksquare$

\textit{\textbf{(条件3)}}\wsp{}令 $A$, $B$, $C$ 皆等しくし, \& 同じ $A^2$, 即即矩形 $AB$ 的動量 $aB + Ab = 2aA$. そのうえ, 同様地, $A^3$, 即体積 $ABC$ 的動量 $aBC + AbC + ABc = 3aA^2$. 按此論理輯, 動量的等級 $A^n$ 為 $naA^{n - 1}$. $\blacksquare$

\textit{\textbf{(条件4)}}\wsp{}由此可知，当 $\frac{1}{A}$ 在 $A$ 中時動量為 1. 故其自身的動量乘積為 \textbf{（此処原文未完）}.

\textbf{第5命題}

速度の二乗に比例する抵抗を受ける物体の運動において、質量が \( A \) の物体の速度が \( a \) であるとき、その動量は \( 1 \)，すなわち \( 0 \) となる。従って、質量が \( \frac{1}{A} \) の物体の動量は \( -\frac{a}{A^2} \) となる。

一般に、質量が \( A^{-n} \) の物体の速度が \( a \) であるとき、その動量は \( n \cdot a \cdot A^{-(n+1)} \) となる。証明終了。\QED

\textit{5. 案例}\wsp{}次に、質量が \( A^{\frac{1}{2}} \) の物体の速度が \( a \) であるとき、その動量は \( 2A^{\frac{1}{2}} \cdot a = 2a\sqrt{A} \) となる。これは第3命題により明らかであるから、従って質量が \( A^{\frac{1}{2}} \) の物体の動量は \( 2a\sqrt{A} \) となる。

\[
\textbf{証明終了}
\]

\[
\frac{a}{2A^{\frac{1}{2}}} \textbf{sive} 2aA^{-\frac{1}{2}}.
\]
一般に、もし \( A^{\frac{m}{n}} \) を \( B \) と仮定すると、\( A^m \) は \( B^n \) と等しくなります。したがって、\( maA^{m - 1} \) は \( nbB^{n - 1} \) と等しくなり、また \( maA^{-1} \) は \( nbB^{-1} \) や \(\frac{nb}{A^{\frac{m}{n}}}\) と等しくなります。さらに、\(\frac{m}{n}aA^{\frac{m-n}{n}}\) は \( b \) と等しい、すなわち \( A^{\frac{m}{n}} \) の動量と等しい。

\textbf{案6.} よって、任意の \( A^m B^n \) の動量は、\( A^m \) を \( B^n \) に掛けるものと \( B^n \) 自身の動量を合わせたものと等しくなります。


\QED


\[
\textbf{導入} \, A^m \text{ に至る道は } m A^{m-1} + n B^{n-1}; \text{ すなわち、指標の指数 } m \& n \text{ が整数または分数であり、正または負である場合も同様です。また、複数の指標の比は等しいです。}
\]
\[
\textbf{命題} \, \textit{1.}\wsp{} \text{したがって、連続する比例項において、与えられた項を除く他の項のモーメントは、その他の項を与えられた間隔数で乗じたものと等しい。}
\]
\[
\textbf{例} \, A, B, C, D, E, F \text{ が連続する比例であるとする；与えられた } C \text{ を除く他の項のモーメントは、}
\]

\textbf{2.} そして，四つの比例之中項が与えられたとき、極端の動量は同じ極端に等しい。

\textit{Corol.\,\textbf{2.}}\hspace*{\fontdimen2\the\textfont2} そのうえ，もし在四個比例中插入両個中項，則両端的動量将与これ些端点相等。これ一点同様适用于与定直角三角形的辺。

\textbf{3.} および，もし与出両平方数之和或差，則両辺的動量将与其辺成反比。

\textit{Corol.\,\textbf{3.}}\hspace*{\fontdimen2\the\textfont2} もし与出両個平方数之和或差，則これ両辺的動量将与其辺成反比。

\condpagelarge{注.}

十年前に私が、最もすぐれた幾何学者 G. G. ライプニッツと交わした手紙の中で、私は、極大と極小を決め、接線を引き、そのほか同様のことを、無理数の項でも有理数の項と同じようにやりとげる方法を、自分が会得していることを知らせ、そして、次の文を含む文字を並べかえて、これを隠しておいた。〔任意の数の流量を含む方程式が与えられたとき、流率を求め、また逆に求めること〕。すると、この高名な人物は、自分もそのような方法に行き当たったと返事を寄こし、自分の方法を知らせてきたが、それは、言葉や記号の形式を除けば、私のものとほとんど違わなかった。どちらの方法の基礎も、この補助定理に含まれている。
\textbf{Prop.\ VIII. Theor.\ VI.}

\textbf{もし物体が均一な重力作用を受け、媒質から均一な抵抗を受けて垂直に上昇または下降する場合、全体の空間が等しい部分に区切られたとき、各部分の初点における絶対的な力を合算すると（上昇する場合は重力に抵抗を加え、下降する場合は重力を減ずる）、その絶対的な力は幾何級数に属することが言える。}

すなわち、物体が下降するときの絶対的な力を表す線分 $AC$ と、媒質からの抵抗を表す不確定な線分 $AK$ を与え、物体が下降している場合の絶対的な力は次のようになる。

\textbf{命題V.} 物体が直線 \( AP \) 上を動くとき、その速度は \( AK \) と \( AC \) の算術平均比に等しくなり、また抵抗の半分の比率に比例します。この条件下で、与えられた時間における抵抗の増加量は直線 \( KL \) と等しいです。

\begin{center}
\includegraphics[width=2323px,height=2211px]{263.png}
\end{center}

物体が直線 \( AP \) 上を移動し、その速度の増加量は直線 \( PQ \) に沿って生じるとき、中心 \( C \) を通って直角双曲線 \( BNS \) を描き、この双曲線は直線 \( CA, CH \) と直交する。ここで、\( AK \) は \( AP \) の比に等しいので、抵抗の動量 \( KL \) もその動量と同じ比率で増加します。

\textbf{証明}： \( AK \) が \( AP \) の比であるとすると、抵抗の動量 \( KL \) も同じ比率で増加する。

\textbf{2APQ}，すなわち \(\textbf{AP}\) に \(\textbf{KC}\) の比です。速度の増分 \(\textbf{PQ}\) は運動法則第2条により，生成する力 \(\textbf{KC}\) に比例します。\(\textbf{KL}\) を \(\textbf{KN}\) と合成すると，\(\textbf{KL} \times \textbf{KN}\) は \(\textbf{AP} \times \textbf{KC} \times \textbf{KN}\) の比となります；すなわち，与えられた \(\textbf{KC} \times \textbf{KN}\) に対する \(\textbf{AP}\) の比です。しかしながら，双曲線面積 \pngright{263.png}{}{}{} の \(\textbf{KNOL}\) が \(\textbf{KL} \times \textbf{KN}\) に対する最終的な比は，点 \(\textbf{K}\) と \(\textbf{L}\) が重なるとき，等しいです。したがって，双曲線面積 \(\textbf{ABOL}\) の消失は \(\textbf{AP}\) に比例します。よって，速度 \(\textbf{AP}\) に対する恒常的な比をもつ各部分 \(\textbf{KNOL}\) を合成すると，全体の双曲線面積 \(\textbf{ABOL}\) が得られます。

\textbf{よって、その速度で描かれた距離と比例する面積は存在します。}既に述べた面積$ABMI$、$IMNK$、$KNOL$等を等分割し、力の絶対値$AC$、$IC$、$KC$、$LC$等は幾何級数に属するものとなります。\QED 同様の議論により、物体の上昇において、点$A$に対する反対側の面積$ABmi$、$imnk$、$knol$等が存在し、力の絶対値$AC$、$iC$、$kC$、$lC$等は連続的に比例します。したがって、上昇と下降においてすべての距離が等しい場合、力の絶対値$lC$、$kC$、$iC$、$AC$、$IC$、$KC$、$LC$等も連続的に比例するものとなります。\(\blacksquare\)

\textbf{命題V．推論1}

したがって、描かれた距離を双曲線面積 \(ABNK\) で表す場合、物体の重力と速度、そして媒質の抵抗はそれぞれ直線 \(AC\), \(AP\), および \(AK\) によって表されることが可能です。逆もまた然りです。

\textbf{推論2}

また、無限に下降する物体が得ることができる最大速度を表すのは直線 \(AC\) です。

\textbf{推論3}

したがって、媒質の抵抗が特定された場合、その抵抗に対抗して最大速度を得るためには、それを直線 \(AC\) と比較することになります。


\textbf{速度}が重力の力と知られている抵抗との比で半分であることを与えます。

\textit{\textbf{補題 4.}}\wsp{}しかし、時間の粒子$\,NKLO$は下降時に描く空間の粒子は最小限であり、それは長方形$KN \times PQ$に等しいです。なぜなら、空間$NKLO$は速度と時間の粒子との比で与えられるからで、したがって時間の粒子はその空間を速度に適用したもの、すなわち最小限の長方形$KN \times KL$を速度$AP$に適用するものに等しいです。かつ先に述べたように$KL$は$AP \times PQ$に比例します。従って、時間の粒子は$KN \times PQ$に比例し、


\textbf{すなわち}


この証明を完了するために、以下に続きます。


\textbf{時間の粒子}は$\,KN \times PQ$に比例します。


 
\textbf{同様に、}$\frac{PQ}{CK} \textbf{}$。 $\qedsymbol$


\textbf{命題 5.} 同じ議論により、物体が上昇中に描く空間の粒子時間 \( nklo \) は \(\frac{pq}{Ck}\) に比例する。




\section*{\textbf{命題 IX. 定理 VII.}}

既に証明したことを基にして、私は次のように主張します。円弧の扇形角と双曲線扇形角の接線が速度に比例する場合、且つ適切な半径が存在するとき、上昇全体の時間は円弧の扇形に比例し、下降全体の時間は双曲線の扇形に比例します。

直線 \( AC \)，重力の方向を示す垂直で等しい直線



\[
ADを引く。中心Dから半径ADで四分円AtEと双曲線AVZを描く。
\]

% -----File: 265.png---
軸はAX、主頂点はA、漸近線はDCを持つ。
\textbf{接線} \(Ap\)と\(AP\)を結ぶ。

\begin{center}
\includegraphics[width=2323px,height=2211px]{263.png}
\end{center}

\textit{案 1.} まず、四分円ADtと三角形ADpの部分 \(Dvq\) を切って、それぞれの瞬間的な要素または最小微粒子を引く。




$tDv$ と $pDq$. 片方の粒子が角度 $\angle D$ の共通角を持つもう一方の粒子と等しいとき、これらの粒子は辺の二重比に比例します。したがって、粒子 $tDv$ は $\frac{qDp}{pD \textbf{quad.}}$ に比例します。

ここで、$pD \textbf{quad.}$ は $AD \textbf{quad.} + Ap \textbf{quad.}$ であり、これは $AD \textbf{quad.} + Ak \times AD$ や $AD \times Ck$ に等しい。また、$qDp$ は $\frac{1}{2}AD \times pq$ です。したがって、セクターの粒子 $vDt$ は $\frac{pq}{Ck}$ に比例し、これはコロラリウム5（命題VIII）により時間の粒子と等しいことが分かります。

各粒子を合計すると、セクター $ADt$ 内のすべての粒子 $tDv$ の和は、各速度が減少する粒子 $Ap$ が失われた時間の総和に比例します。


\textbf{答}：これらの反応がその速度を完全に失うまで；これは、全体のセクター \( ADt \) は全体の未来時間に対する昇降比であることを意味します。
\(\blacksquare\)

\textit{\textbf{例2}}：線分 \( DQV \) を引くと、セクター \( DAV \) の部分を \( TDV \) と \( PDQ \) に分割し、これらの部分は互いに

$\frac{DTq.}{DPq.} = \frac{DXq.}{DAq.} \quad$ \text{または} $\quad \frac{TXq.}{APq.}$

であり、さらに
 
$\frac{DXq. - TXq.}{ADq. - APq.}.$


しかし、双曲線 \( DXq. - TXq. \) は \( ADq. \)，そして仮定により \( APq. \) である。

\[
est \ AD \times AK.
\]
従って、各部分は互いに次の比で比例します：
\[
\frac{ADq.}{AD \times AK}; 
\]
すなわち、
\[
\frac{AD}{AD - AK} \textbf{ または } \frac{AC}{CK}.
\]
したがって、セクターの部分 \( TDV \)は
\[
\frac{PDQ \times AC}{CK},
\]
であり、従って、与えられた \( AC \) および \( AD \) の値から
\[
\frac{PQ}{CK};
\]
したがって、コロラリウム5.命題VIII.第2編において、部分時間の総和は速度 \( AP \) の各部分 \( PQ \) によって生成される全速度の増加に比例します。
最後に、部分時間を合成すると、すべての速度 \( AP \) の部分 \( PQ \) が生成する総時間はセクター \( ADT \) の総部分であると分かります。

\textbf{証明終了。}\QED

\textit{命題 1.}\wsp{}次に、もし \( AB \) が \( AC \) の四分之一と等しいとき、物体が落下中に時間 \( ATD \) 内に描く面積 \( ABRP \) は、同じ時間内に最大速度 \( AC \) の半分で均一に進むことで描くことができる面積 \( ATD \) に対する比として表されます。なぜなら、\( AC \) と \( AP \) の比が \( AP \) と \( AK \) の比と同じであるとき、面積 \( 2APQ \) は \( AC \times KL \) に等しい（命題 1. 論理学 II の補題）から、\( KL \) と \( PQ \) の比が \( 2AP \) と \( AC \) の比と同じであるため、面積 \( LKN \) は \( PQ \times \frac{1}{2}AD \) や \( DPQ \) に等しく、さらに \( 2AP \times KN \) と \( \frac{1}{2}AC \times AD \) の比と同じであるからです。

\[
\text{しかし、} DPQ \text{は} DTV \text{に} CK \text{は} AC \text{に比例します。}
\]
\[
\text{従って、等しい比から} LKN \text{は} DTV \text{に} 2AP \times KN \times CK \text{は} \frac{1}{2}AC^3 \text{に比例します；}
\]
\[
\text{すなわち、} CKN \& \frac{1}{4}AC^3 \text{が等しいとき、} AP \text{は} AC \text{に比例します；}
\]
\[
\text{これは落ちる物体の速度} \textbf{と} \text{最大速度} \textbf{との比} \text{に相当します。}
\]
\[
\text{したがって、図形} ABKN \& AVD \text{の面積} LKN \& DTV \text{は速度に比例するので、}
\]
\[
\text{これらの面積全体は同時に生成された空間全体と等しい比を持ちます；}
\]
\[
\text{したがって、最初から生成された面積} ABKN \& AVD \text{は最初から描かれた全距離と等しい比を持ちます。} \qquad \blacksquare
\]

\textbf{補則2.} また、上昇中に描かれる距離についても同様の結論が成り立ちます。すなわち、そのすべての距離は、速度 \( AC \) で等しい時間に描かれた距離 \( ABnk \) に対して、面積 \( ABnk \) はセクター \( ADt \) に比例します。

\textbf{補則3.} 落下体 \( ATD \) の速度は、同一の時間内に抵抗がない媒質で得られる速度 \( APD \) の三角形に対するセクタ \( ATD \) に比例します。実際、媒質が抵抗がない場合の速度は時間 \( ATD \) に比例し、抵抗がある場合には \( AP \) に比例し、これは三角形 \( APD \) に比例することを意味します。

\textbf{そして、物体が下降初速で互いに比を成す速度は、領域 \( ATD \) と \( APD \) の比に等しい。}

\textit{\textbf{補題4.}}\hspace*{1em}同様の議論により、上昇中の物体の速度は、同一時間内に抵抗のない空間で全ての上昇運動を失う速度 \( AP \) に対して、三角形 \( ApD \) に対する円錐曲線セクター \( AtD \) の比、すなわち直線 \( Ap \) に対する弧 \( At \) の比に等しい。

\textit{\textbf{補題5.}}\hspace*{1em}したがって、物体が抵抗のある媒質で落下し、速度 \( AP \) を得るのに要する時間は、同一時間内に抵抗のない空間で最大速度 \( AC \) を得るのに要する時間に対して比を成す。

\[
\text{ut Sector } ADT \text{ ad triangulum } ADC: \& \text{ tempus, quo velocitatem } Ap \text{ in Medio resistente ascendendo possit amittere, ad tempus } quo velocitatem eandem in spatio non resistente ascendendo posset amittere, ut arcus  At \text{ ad ejus Tangentem } Ap.
\]

\textbf{命題VI.} 速度の二乗に比例する抵抗を受ける物体の運動において、物体が抵抗のある空間で上昇する時間を与えられた時間とすると、その比は扇形 \( ADT \) に対する三角形 \( ADC \) の比と同じである。また、物体が同様な速度で抵抗がない空間で上昇する時間を与えたとき、その比も同じである。

\textit{命題VIの推論6.} よって、与えられた時間から移動距離は求められる。無限に下降する物体について、推論2と3により最大速度が与えられるので、そこから与えられた時間において半分の速度で移動できる距離も求められ、その時間も求められる。

\[
\text{Corol.\fsp{}6.}\wsp{}Hinc ex dato tempore datur spatium ascensu vel descensu descriptum. Nam corporis in infinitum descendentis datur velocitas maxima, per Corol.\ 2.\ \& 3.\ Theor.\ VI, Lib.\ II. inde{\que} datur \& spatium quod semisse velocitatis illius dato tempore describi potest, \& tempus quo corpus velocitatem illam in
\]

\[
\text{落下する物体が非抵抗性の空間で速度の二乗に比例する速度を獲得できる。そして、セクター } ADT \text{ や } ADt \text{ を三角形 } ADC \text{ に対する時間比として取り扱うと、}
\]
\[
\text{そのときの速度は} AP \text{ や } Ap, \text{ 面積は} ABKN \text{ や } ABkn \text{ となる。これらはセクターの比に等しく、移動した距離の比に等しい。}
\]

\textit{\textbf{補則7.}}\wsp{}そして逆推して、与えられた上昇または下降の距離 \( ABnk \) や \( ABNK \) から、時間 \( ADt \) や \( ADT \) が得られる。

\condpagelarge{\textbf{命題X.} \textbf{証明III.}}

\textit{\textbf{Tendat uniformis vis gravitatis directe ad planum Horizontis, sitque resistentia ut medii densitas \& quadratum velocitatis conjunctim:}}
\[
\text{均質な重力の力が水平面に直接指向し、抵抗は媒質の密度と速度の二乗の積に等しいとする：}
\]

求められるのは、媒質の密度が各点でどのようなものかであり、その結果、物体が任意の曲線 \(AK\) 上を動く速度も同様に求められます。


\par\begin{center}
  \centering
  \includegraphics[width=1190px,height=966px,keepaspectratio]{271.png}

\end{center}\par

平面 \(AK\) が図示の平面に垂直であり、曲線 \(ACK\) 上で物体 \(C\) が動きます。また、直線 \(FCf\) は点 \(C\) でこの曲線を接しています。

さらに、物体 \(C\) は現在 \(A\) から \(K\) へと曲線 \(ACK\) 上を進み、その後同じ曲線 \(ACK\) 上を戻ります。また、進むときには媒質により阻害され、戻るときには同様に推進されます。したがって、各点において物体の速度は等しくなります。

\[
\textbf{証明}
\]

物体 \(C\) が曲線 \(ACK\) 上を動く際、媒質による抵抗と推進力によって物体の運動が制限されることが示されます。進むときには阻害され、戻るときには推進されるため、各点での速度は等しくなります。

\[
\textbf{命題V}
\]

曲線 \(ACK\) 上で物体 \(C\) が動くとき、媒質の密度が各点でどのように変化するかを求める方法について述べます。物体 \(C\) の運動は媒質の密度によって制限され、その結果、物体の速度も各点で異なる値を持つことがわかります。

\[
\textbf{命題VI}
\]

曲線 \(ACK\) 上での物体 \(C\) の速度を求める方法について述べます。物体 \(C\) が進むときと戻るときの速度は等しくなり、媒質による抵抗と推進力によって制限されます。

\[
\textbf{命題VII}
\]

曲線 \(ACK\) 上での物体 \(C\) の運動を詳しく分析します。物体 \(C\) が各点でどのような速度を持つか、そしてその速度は媒質の密度によってどのように影響されるかについて説明します。

\[
\textbf{命題VIII}
\]

曲線 \(ACK\) 上での物体 \(C\) の運動をさらに詳しく分析します。物体 \(C\) が各点でどのような速度を持つか、そしてその速度は媒質の密度によってどのように影響されるかについて更に詳細に説明します。

\[
\textbf{命題IX}
\]


\[
\textbf{命題X}
\]


\[
\textbf{命題XI}
\]

曲線 \(ACK\) 上での物体 \(C\) の運動をさらに詳しく分析します。物体 \(C\) が各点でどのような速度を持つか、そしてその速度は媒質の密度によってどのように影響されるかについて更に詳細に説明します。

\[
\textbf{命題XII}
\]


\[
\textbf{命題XIII}
\]


\[
\textbf{命題XIV}
\]

曲線 \(ACK\) 上での物体 \(C\) の運動をさらに詳しく分析します。物体 \(C\) が各点でどのような速度を持つか、そしてその速度は媒質の密度によってどのように影響されるかについて更に詳細に説明します。

以上の証明と命題により、曲線 \(ACK\) 上での物体 \(C\) の運動に関する詳細な理解を得ることができます。

常に物体が進行し、退行するときの速度は存在します。

しかし、等しい時間内に、進行中の物体が最小限の弧$CG$を描き、退行中の物体が弧$Cg$を描くとします。そして、これらの物体が点$C$, $G$, $g$からそれぞれ直線$CH$, $Ch$の長さで移動し、その時間内に媒介力および重力の作用により平面水平線$AK$へ垂直な線$CB$, $GD$, $gd$を引くとします。これらの線$GD$と$gd$はそれぞれ弧$CG$と$Cg$の接線上で点$F$と点$f$に到達します。媒介力により、進行中の物体の速度は長さ$CH$の逆比例する関係にあるとします。

\textbf{CF}の長さを描き、物体が重力の力により$F$から$G$へ移動する。そして、抵抗の力によって線分$HF$が生成され、同時に重力の力によって線分$FG$が生成される。したがって（第1巻第10 Lem.より）、線分$FG$は重力の力と時間の平方の積に比例し、また（与えられた重力により）時間の平方に比例する。一方で、線分$HF$は抵抗と時間の平方に比例し、これは抵抗と線分$FG$に比例することを意味する。従って、生成される抵抗は線分$HF$と直接、線分$FG$と逆に比例し、すなわち$\frac{HF}{FG}$に比例する。これらは生成される線分の関係において同じである。また、有限長の線分に対するこれらの比は

\begin{equation*}
\textbf{CF} \quad \text{を描く}
\end{equation*}

物体が重力の力により$F$から$G$へ移動する。そして、抵抗の力によって線分$HF$が生成され、同時に重力の力によって線分$FG$が生成される。したがって（第1巻第10 Lem.より）、線分$FG$は重力の力と時間の平方の積に比例し、また（与えられた重力により）時間の平方に比例する。一方で、線分$HF$は抵抗と時間の平方に比例し、これは抵抗と線分$FG$に比例することを意味する。従って、生成される抵抗は線分$HF$と直接、線分$FG$と逆に比例し、すなわち$\frac{HF}{FG}$に比例する。これらは生成される線分の関係において同じである。また、有限長の線分に対するこれらの比は

\begin{equation*}
\textbf{CF} \quad \text{を描く}
\end{equation*}

同様に、有限な長さを持つ線分におけるこれらの比も同様です。

非正確である。

同様の論理により、$fg$ は時間の平方に比例し、等しい時間間隔において $FG$ と等しくなる。また、物体が退行するときの衝撃量は $\frac{hf}{fg}$ に比例する。そして、物体が退行するときの衝撃量と進行するときの抵抗量は運動開始時に等しい比で測定され、したがって、$\frac{hf}{fg}$ および $\frac{HF}{FG}$ も同じ比で測定される。従って、$fg$ および $FG$ が等しいとき、$hf$ および $HF$ も等しくなる。これにより、$CF$ と $CH$（または $Ch$）および $Cf$ は算術級数に属し、したがって $HF$ は半差となる。

\[
\textbf{これらの} Cf \& CF; そして、先に述べた抵抗は
\]
\[
\frac{HF}{FG} \textbf{と等しく、つまり} \frac{Cf - CF}{FG} \textbf{と等しい。}
\]

\[
\textbf{この抵抗は中間体の密度と速度の二乗に比例する。}
\]
\[
\textbf{速度は描かれた長さ CF と時間 } \sqrt{FG} \textbf{との比に直接、そして逆に比例する、つまり} \frac{CF}{\sqrt{FG}} \textbf{と等しく、したがって速度の二乗は} \frac{CF^2}{FG} \textbf{と等しい。}
\]
\[
\therefore \textbf{これらの抵抗} \frac{Cf - CF}{FG} \textbf{は中間体の密度と} \frac{CF^2}{FG} \textbf{との積に比例する；そして、したがって中間体の密度は} \frac{Cf - CF}{FG} \textbf{と} \frac{CF^2}{FG} \textbf{との比に直接、そして逆に比例する、つまり} \frac{Cf - CF}{CF^2}
\]

\textbf{補遺}

\textit{1.} したがって、$Cf$において$Ck$を$CF$に等しい長さとして取り、平面水平線$AK$に対して直角$ki$を引くと、$ki$は曲線$ACK$の点$l$で交わる。そのとき、中間密度は
\[
\frac{FG - kl}{CF \times (FG + kl)}
\]
となる。

これは$FC$と$KC$が$\sqrt{fg}$または$\sqrt{FG}$と$\sqrt{kl}$に比例することから、さらに分母分子で割ることによって得られる。すなわち
\[
\frac{fk}{kC} = \frac{\sqrt{FG} - \sqrt{kl}}{\sqrt{kl}},
\]
これはまた
\[
\frac{CF - Cf}{CF} = \frac{\sqrt{FG} - \sqrt{kl}}{\sqrt{kl}}
\]
と等しい。したがって、$\sqrt{FG}$と$\sqrt{kl}$の和または差を用いて項を引くと
\[
\frac{FG - kl}{kl + \sqrt{FG} \times kl},
\]
すなわち
\[
\frac{FG - kl}{FG + kl}
\]
となる。初等比は



\[
kl + \sqrt{FG} \times kl \& FG + kl \textbf{は等しい。} これを以下のように書き換える。
\]
\[
\frac{FK - Kl}{FK + Kl} \textbf{を} \frac{Cf - CF}{CF} \textbf{とする; }
\]
\[
\textbf{そして, 中心の密度は} \frac{Cf - CF}{CF^2} \textbf{であったので,}
\]

\textbf{結果として} $\frac{FG - kl}{CF \times FG + kl}$ \textbf{となる。}

\textit{命題 2.}\wsp{}従って, $2HF$ \& $Cf - CF$ が等しいとき, そして $FG$ \& $kl$ (等比の関係にある)を $2FG$ と組み合わせると,

2HF \textbf{は} CF \textbf{に対して} FG - kl \textbf{は} 2FG \textbf{に対して等しい;} 

 
\textbf{従って, } $HF$ \textbf{は} $FG$ \textbf{に対して, 即ち抵抗は重力に対して}

\[
\textbf{直角三角形} CF \textbf{と} FG - kl \textbf{の積に対する} 4FG^2 \textbf{に対して等しい。}
\]

\textbf{命題 3.}

次に、曲線 \( AB \) を基底または軸 \( AB \) と垂直に引かれた高さ \( BC \) の関係によって定義する場合（これは通常の方法です）で、高さ \( BC \) の値を級数収束型の級数に分解し、その最初の項を使用して問題を素早く解決します。これについては、以下に示す例をご覧ください。

\textbf{例 1.}

半円 \( ACK \) を直径 \( AK \) 上で描き、この曲線上で物体が運動するための中間密度を求めます。

直径 \( AK \) の中点 \( O \) を取ります；そして、\( OK = n \), \( OB = a \) とします。

\textbf{BC} に \textit{e}, または \textit{BD} や \textit{Bi} が接するとき，\textbf{DGq.} または \textbf{OGq. - ODq.} は等しいとします。すなわち，
\[ nn - aa - 2ao - oo = ee - 2ao - oo \]
または
\[ ee - 2aa - oo. \]
これらの式を \textit{nostrum} の方法で平方根を抽出すると，得られる結果は
\[ DG = e - \frac{ao}{e} - \frac{oo}{2e^2} - \frac{aaoo}{2e^4} - \frac{ao^3}{2e^4} - \frac{a^3o^3}{2e^6} + \cdots. \]
ここで，\textit{nn} を \textit{ee + aa} と置き換えると，
\[ DG = e - \frac{ao}{e} - \frac{nnoo}{2e^4} - \frac{anno^3}{2e^6} + \cdots. \]

このような級数を次のように項に区分します。最初の項は，無限に小さい量 \textit{o} が存在しない項で定義します。次に，そのような量 \textit{o} が存在する項を定義します。



\textbf{次に、速度の二乗に比例する抵抗を受ける物体の運動について述べます。}

\textbf{命題V（定理5）}：速度の二乗に比例する抵抗を受ける物体が直線$BC$上から点$B$に向かって移動し、その位置を表す直線$e$と平行な直線$DF$との差を示す直線$\frac{ao}{e}$で定義されるとき、この直線は常に直線$BC$の位置を表します。また、直線$IF$は直線$ID$と直線$BC - ID$の差を表し、これが完成された平行四辺形$BC - ID$から引かれた線分です。したがって、接線$CF$の位置も常に定義されます。この場合、直線$IF$に対する直線$IC$の比は$\frac{ao}{e}$に対する比$o$または$a$に対する比$e$と等しいとします。

\textbf{証明}：まず、直線$BC$が点$B$から始まり、直線$DF$との差を表す直線$\frac{ao}{e}$で定義されるとき、直線$IF$は直線$ID$と直線$BC - ID$の差を表します。これは平行四辺形$BC - ID$から引かれた線分です。したがって、接線$CF$の位置も常に定義されます。

ここで、直線$IF$に対する直線$IC$の比は$\frac{ao}{e}$に対する比$o$または$a$に対する比$e$と等しいとします。この関係を用いて、物体の運動を解析します。

\textbf{命題VI（定理6）}：速度の二乗に比例する抵抗を受ける物体が直線$BC$上から点$B$に向かって移動し、その位置を表す直線$e$と平行な直線$DF$との差を示す直線$\frac{ao}{e}$で定義されるとき、この直線は常に直線$BC$の位置を表します。また、直線$IF$は直線$ID$と直線$BC - ID$の差を表し、これが完成された平行四辺形$BC - ID$から引かれた線分です。したがって、接線$CF$の位置も常に定義されます。

\textbf{証明}：同様に、直線$IF$に対する直線$IC$の比は$\frac{ao}{e}$に対する比$o$または$a$に対する比$e$と等しいとします。この関係を用いて、物体の運動を解析します。

\textbf{命題VII（定理7）}：速度の二乗に比例する抵抗を受ける物体が直線$BC$上から点$B$に向かって移動し、その位置を表す直線$e$と平行な直線$DF$との差を示す直線$\frac{ao}{e}$で定義されるとき、この直線は常に直線$BC$の位置を表します。また、直線$IF$は直線$ID$と直線$BC - ID$の差を表し、これが完成された平行四辺形$BC - ID$から引かれた線分です。したがって、接線$CF$の位置も常に定義されます。

\textbf{証明}：同様に、直線$IF$に対する直線$IC$の比は$\frac{ao}{e}$に対する比$o$または$a$に対する比$e$と等しいとします。この関係を用いて、物体の運動を解析します。

\textbf{命題VIII（定理8）}：速度の二乗に比例する抵抗を受ける物体が直線$BC$上から点$B$に向かって移動し、その位置を表す直線$e$と平行な直線$DF$との差を示す直線$\frac{ao}{e}$で定義されるとき、この直線は常に直線$BC$の位置を表します。また、直線$IF$は直線$ID$と直線$BC - ID$の差を表し、これが完成された平行四辺形$BC - ID$から引かれた線分です。したがって、接線$CF$の位置も常に定義されます。

\textbf{証明}：同様に、直線$IF$に対する直線$IC$の比は$\frac{ao}{e}$に対する比$o$または$a$に対する比$e$と等しいとします。この関係を用いて、物体の運動を解析します。

\textbf{命題IX（定理9）}：速度の二乗に比例する抵抗を受ける物体が直線$BC$上から点$B$に向かって移動し、その位置を表す直線$e$と平行な直線$DF$との差を示す直線$\frac{ao}{e}$で定義されるとき、この直線は常に直線$BC$の位置を表します。また、直線$IF$は直線$ID$と直線$BC - ID$の差を表し、これが完成された平行四辺形$BC - ID$から引かれた線分です。したがって、接線$CF$の位置も常に定義されます。

\textbf{証明}：同様に、直線$IF$に対する直線$IC$の比は$\frac{ao}{e}$に対する比$o$または$a$に対する比$e$と等しいとします。この関係を用いて、物体の運動を解析します。

\textbf{命題X（定理10）}：速度の二乗に比例する抵抗を受ける物体が直線$BC$上から点$B$に向かって移動し、その位置を表す直線$e$と平行な直線$DF$との差を示す直線$\frac{ao}{e}$で定義されるとき、この直線は常に直線$BC$の位置を表します。また、直線$IF$は直線$ID$と直線$BC - ID$の差を表し、これが完成された平行四辺形$BC - ID$から引かれた線分です。したがって、接線$CF$の位置も常に定義されます。

\textbf{証明}：同様に、直線

第三項は、$\frac{nnoo}{2e^3}$ で表され、接線と曲線の間にある線素 $FG$ を示し、点 $C$ における接触角 $FCG$ や、曲線が $C$ 点において持つ曲率を定める。もし線素 $FG$ が有限な大きさを持つ場合、それは無限に続く次項とともに第三項とみなされる。しかし、線素 $FG$ が無限小になる場合、その後の項はさらに無限に小さくなるため、これらは無視できる。第四項は $\frac{anno^3}{2e^5}$ で表され、曲率の変化を示す。第五項は曲率の変化の変化を示す。

\textbf{変化量は、さらなる変化も同様に。従って、これらの級数は曲線に関する問題の解法において軽視できない有用性を示す。}

また、$CF$ は $CIq$ と $IFq$ の正方形からなり、これは $BDq$ と第二項の平方からなる。さらに、$FG + kl$ は第三項の二倍に等しく、$FG - kl$ は第四項の二倍に等しい。なぜなら、$DG$ の値は $il$ の値に変換され、$FG$ の値は $kl$ の値に変換されるからである。ここで、$Bi$ を $BD$ と書き換え、$-o$ を $+o$ とする。

\[
FG = -\frac{nnoo}{2e^3} - \frac{anno^3}{2e^5} + \cdots
\]

\[
kl = -\frac{nnoo}{2e^3} + \frac{anno^3}{2e^5}
\]
そして、これら項の合計は \(\displaystyle -\frac{nnoo}{e^3}\), 差は \(\displaystyle -\frac{anno^3}{e^5}\) である。

ここでは第五項以降を省略し、これらの項がこの問題において考慮されるよりもずっと小さいものとして扱う。よって、もし系列を普遍的に次のように表すとすると
\[
\text{Series} = \mp Qo - Roo - So^3 + \cdots
\]
\(\displaystyle CF\) は
\[
\sqrt{\overline{oo + QQoo}}
\]
\(\displaystyle FG + kl\) は \(\displaystyle 2Roo\)
\(\displaystyle FG - kl\) は \(\displaystyle 2So^3\) である。

これらの値を表すために、\(\displaystyle CF\), \(\displaystyle FG + kl\), \(\displaystyle FG - kl\) をそれぞれ
\[
\text{これらの値, そしてその中間密度}
\]
と書き換える。 

\[
\frac{\textbf{FG} - \textbf{kl}}{\textbf{CF} \operatorname{in} \textbf{FG} + \textbf{kl}}
\]
は

\[
\frac{\textbf{S}}{\textbf{R} \surd \overline{1 + QQ}}
\]
と等しくなります。

したがって、系列に収束する問題を導き出します。ここで $Q$, $\textbf{R}$, および $\textbf{S}$ をそれぞれ系列の項に対応するものとして書き換えることと、さらに抵抗を任意の点 $\textbf{G}$ で体が受けるものとするとき、その抵抗は重力に対する比として

\[
\frac{\textbf{S} \surd \overline{1 + QQ}}{2\textbf{RR}}
\]

であり、

\[
\textbf{S}
\]
は物体が直線 $\textbf{CF}$ に従って点 $\textbf{C}$ から出発し、円錐曲線の直径 $\textbf{CB}$ と水平軸 $\frac{\textbf{1} + \textbf{QQ}}{\textbf{R}}$ を持つ放物線上を動くときの速度となります。

この条件で問題は解けます。

\textbf{そのようにして問題を解いた場合、}\(\sqrt{\overline{1 + \frac{aa}{ee}}}\) や \(\frac{n}{e}\) を \(\sqrt{\overline{1 + QQ}}\) として書き換えると、\(\frac{nn}{2e^3}\) が \(R\)、そして \(\frac{ann}{2e^3}\) が \(S\) となる。すると、中間密度は \(\frac{a}{ne}\) に比例し、これは与えられた \(n\) の値によって \(\frac{a}{e}\) や \(\frac{OB}{BC}\) に比例し、すなわち接線の長さ \(CT\) が半径 \(OL\) に対して垂直に接する点 \(K\) まで達する長さに比例します。抵抗は重力に対する比として \(a\) と \(n\) の比、すなわち円周半径 \(OK\) に対する \(OB\) の比となります。また、速度は \(\sqrt{2BC}\) に比例します。したがって、物体 \(C\)

\textbf{確実に一定の速度で、直線} $\mathrm{OK}$ と平行な方向から点 $\mathrm{L}$ 離開し、中間密度 $C$ は各点において弧 $\mathrm{TCT}$ の接線長さと比例し、ある特定地点 $C$ における抵抗もまた重力の力 $OB$ に対する比率として $OK$ に比例する。その物体は円弧 $\mathrm{LCK}$ を描く。\QEI

しかし、同じ物体が点 $\mathrm{A}$ から直線 $\mathrm{AK}$ と垂直な方向へと出発した場合、$OB$ もしくは $a$ は中心 $\mathrm{O}$ の両側に反対方向へ移動する必要があり、従ってその符号も変更され、$-a$ が +$a$ として書き換わる。これにより中間密度は

\begin{equation}
\textbf{如上所述，もし同一物体従点 \( A \) 沿着与直線 \( AK \) 垂直的方向离开，則 \( OB \) 或 \( a \) 需要向中心 \( O \) 的相反方向移動，したがって其符号也需要改変為 \(-a\)。由此できる得出中密度：}
\end{equation}

\textbf{如上所述，もし同一物体従点 \( A \) 沿着与直線 \( AK \) 垂直的方向离开，則 \( OB \) 或 \( a \) 需要向中心 \( O \) 的相反方向移動，したがって其符号也需要改変為 \(-a\)。由此できる得出中密度：}

\[
C = -\frac{OB}{OK}
\]

\textbf{すなわち、中間密度は重力の力 \( OB \) に対する比率として \(-\frac{OB}{OK}\) となる。}

\[
-\frac{a}{e}.
\]
しかし、負の密度（すなわち、運動を加速するもの）は自然には存在しないとされています。したがって、物体が点 \( A \) から上昇して円弧 \( AL \) を描くことは自然に起こり得ません。この結果を得るためには、物体が中間の力によって加速されるべきであり、抵抗によって妨げられるべきではありません。

\textit{例2.}\wsp{}線 $ALCK$ を放物線とし、その軸 $OL$ は水平線 $AK$ に垂直であるとする。そして、投射体がその上を動くようにする、媒質の密度を求めよ。

放物線の性質から、長方形 $ADK$ は、縦線 $DG$ と、ある与えられた直線とがつくる長方形に等しい。すなわち、その直線を $b$、$AB$ を $a$、$AK$ を $c$、$BC$ を $e$、$BD$ を $o$ と呼べば、長方形 $(a + o)(c - a - o)$、すなわち $ac - aa - 2ao + co - oo$ は、$b$ と $DG$ の長方形に等しく、したがって $DG$ は
\[ \frac{ac - aa}{b} + \frac{c - 2a}{b}o - \frac{oo}{b} \]
に等しい。ここで、この級数の第2項 $\frac{c - 2a}{b}o$ を $Qo$ に、その係数 $\frac{c - 2a}{b}$ を $Q$ に書き直し、また第3項 $\frac{oo}{b}$ を $Roo$ に、その係数 $\frac{1}{b}$ を $R$ に書き直さなければならない。ところが、項はこれ以上ないから、第4項 $So^3$ の係数 $S$ は消えなければならず、したがって、媒質の密度が比例する量 $\frac{S}{R\sqrt{1 + QQ}}$ はゼロになる。ゆえに、どんな媒質の密度もなくても、投射体は、かつてガリレオが証明したとおり、放物線の上を動く。\QED

\textit{例3.}\wsp{}線 $AGK$ を双曲線とし、その漸近線 $NX$ は水平面 $AK$ に垂直であるとする。そして、投射体がこの線の上を動くようにする、媒質の密度を求めよ。

\pngright{275.png}{}{}{}

もう一つの漸近線を $MX$ とし、延長した縦線 $DG$ と $V$ で交わらせる。双曲線の性質から、長方形 $XV \times VG$ は与えられる。ところが $DN$ 対 $VX$ の比は与えられているから、長方形 $DN \times VG$ も与えられる。それを $bb$ とし、平行四辺形 $DNXZ$ を作って、$BN$ を $a$、$BD$ を $o$、$NX$ を $c$ と呼び、与えられた比 $VZ$ 対 $ZX$、すなわち $DN$ を $\frac{m}{n}$ とする。すると、$DN$ は $a - o$ に等しく、$VG$ は $\frac{bb}{a - o}$、$VZ$ は $\frac{m}{n}(a - o)$ に等しく、$GD$ すなわち $NX - VZ - VG$ は
\[ c - \frac{m}{n}a + \frac{m}{n}o - \frac{bb}{a - o} \]
に等しい。項 $\frac{bb}{a - o}$ を収束級数
\[ \frac{bb}{a} + \frac{bb}{aa}o + \frac{bb}{a^3}o^2 + \frac{bb}{a^4}o^3 + \text{など} \]
に展開すれば〔原文は最後の項を $\frac{bb3}{a^4}o^3$ と印刷しているが、$bb$ の誤りと見た〕、$GD$ は
\[ c - \frac{m}{n}a - \frac{bb}{a} + \frac{m}{n}o - \frac{bb}{aa}o - \frac{bb}{a^3}o^2 - \frac{bb}{a^4}o^3 - \text{など} \]
に等しくなる。この級数の第2項 $\frac{m}{n}o - \frac{bb}{aa}o$ は $Qo$ のかわりに、第3項は符号を変えて $\frac{bb}{a^3}o^2$ を $Ro^2$ のかわりに、第4項も符号を変えて $\frac{bb}{a^4}o^3$ を $So^3$ のかわりに用い、それらの係数 $\frac{m}{n} - \frac{bb}{aa}$, $\frac{bb}{a^3}$, $\frac{bb}{a^4}$ を、上の規則の $Q$, $R$, $S$ のかわりに書かなければならない。こうすれば、媒質の密度は
\[ \frac{\frac{bb}{a^4}}{\frac{bb}{a^3}\sqrt{1 + \frac{mm}{nn} - \frac{2mbb}{naa} + \frac{b^4}{a^4}}} \quad\text{すなわち}\quad \frac{1}{\sqrt{aa + \frac{mm}{nn}aa - \frac{2mbb}{n} + \frac{b^4}{aa}}} \]
となる〔原文は根号の中を $1 - \frac{mm}{nn}$ と印刷しているが、次の式から $+$ とした〕。$VZ$ 上に $VY$ を $VG$ に等しくとれば、これは $\frac{1}{XY}$ に比例する。実際、$aa$ と $\frac{mm}{nn}aa - \frac{2mbb}{n} + \frac{b^4}{aa}$ は、$XZ$ と $ZY$ の平方だからである。そして抵抗は、重力に対して $XY$ 対 $YG$ の比をなすことがわかり、速さは、物体が、頂点 $G$、直径 $DG$、通径 $\frac{XY^2}{VG}$ の放物線の上を進むときのそれである。そこで、各点 $G$ での媒質の密度が距離 $XY$ に反比例し、ある点 $G$ での抵抗の重力に対する比が $XY$ 対 $YG$ であるとすれば、点 $A$ から適当な速さで放たれた物体は、その双曲線 $AGK$ を描く。\QED
\textbf{例4.} 軌道 \( AGK \) が無限に延びる直線として定められる。この Hyperbola は中心 \( X \)、漸近線 \( MX \), \( NX \) で描かれる。


\textbf{命題V．}曲線 \(XZDN\) の直角三角形 \(XZDN\) の辺 \(ZD\) が双曲線を交わり、その漸近線と交わる点をそれぞれ \(G\) および \(V\) とするとき、\(VG\) は任意の指数 \(n\) で \(ZX\) や \(DN\) の何らかのべき \(ND^n\) に反比例し、求めるのはこの曲線上での物体の運動速度の二乗に比例する抵抗を受ける場合の中間密度。

まず、\(DN\), \(BD\), \(NX\) をそれぞれ \(A\), \(O\), \(C\) と書き、さらに \(VZ\) は \(ZX\) や \(DN\) の比 \(d : e\) で与えられるとする。また、
\[ VG = \frac{b^2}{D N^n} \]
であるから、
\[ D N = A - O, \quad V G = \frac{b^2}{A - O^n}, \quad V Z = \frac{d}{e} (A - O), \quad G D \text{ または } N X - V Z - V G = C - \frac{d}{e} A + \frac{d}{e} O - \frac{b^2}{A - O^n}. \]

Resolvatur terminus ille \(\frac{bb}{A - O^n}\) in seriam infinitam
\[
\frac{bb}{A^n} + \frac{nbbO}{A^{n+1}} + \frac{nn + n}{2A^{n+2}} bbO^2 + \frac{n^3 + 3nn + 2n}{6A^{n+3}} bbO^3 \&c.
\]
ac fiet \(GD\) \(\textbf{等しく}\)
\[
C - \frac{d}{e} A - \frac{bb}{A^n} + \frac{d}{e} O - \frac{nbb}{A^{n+1}} O - \frac{nn + n}{2A^{n+2}} bbO^2 - \frac{n^3 + 3nn + 2n}{6A^{n+3}} bbO^3 \&c.
\]
Hujus seriei terminus secundus \(\frac{d}{e} O - \frac{nbb}{A^{n+1}} O\) usurpandus est pro \(Qo\),
tertius \(\frac{nn + n}{2A^{n+2}} bbO^2\) pro \(Ro^2\),
quartus \(\frac{n^3 + 3nn + 2n}{6A^{n+3}} bbO^3\) pro

\textbf{第70節}

次に、中間体の密度は任意の点 \( G \)において
\[
\frac{S}{R \times \sqrt{\overline{1 + QQ}}}
\]
であり、ここで \( Q = r \) とします。この式を展開すると、
\[
\frac{n + 2}{3\sqrt{A^2 + \frac{dd}{ee} A^2 - \frac{2dnbb}{eA^n} A + \frac{nnb^4}{A^{2n}}}}
\]
となります。

ここで、点 \( VZ \)における \( VY \)を定義します。\( VY = n \times VG \)とすると、これは \( XY \)に反比例します。

なぜなら、
\[
A^2 + \frac{dd}{ee} A^2 - \frac{2dnbb}{eA^n}
\]
は \( XZ \)および \( ZY \)の平方であり、また
\[
A^2 + \frac{dd}{ee} A^2 - \frac{2dnbb}{eA^n} + \frac{nnb^4}{A^{2n}}
\]
は \( XY \)の平方に等しいからです。

したがって、抵抗は点 \( G \)において重力に対して
\[
S : 2RR = \frac{XY}{A} : \frac{3nn + 3n}{n + 2} VG
\]
と比例します。すなわち、
\[
\text{抵抗} : \text{重力} = XY : \frac{3nn + 3n}{n + 2} VG.
\]

したがって、投射された物体の速度は
\[
XY : \frac{3nn + 3n}{n + 2} VG
\]
となります。

\textbf{命題V．}

物体が放たれ、軌道が双曲線であるとき、頂点を \( G \)、直径を \( GD \) と \( L \) 横軸 \(\frac{1 + QQ}{R}\) や \(\frac{2XY \text{（平方）}}{\overline{nn + n} \operatorname{in} VG}\) を持つ場合、物体は双曲線を描きます。証明終了。\QEIit

\condpagelarge{\textit{補遺．}}

物体が抵抗のない媒質中で双曲線軌道を描くことはありませんが、ここに示したように抵抗のある媒質中では双曲線軌道を描きます。明らかに、抵抗が均一な媒質中で放たれた物体が描く線分は、この双曲線よりも近づいています。その線分は実際には双曲線の種類ですが、頂点からアスンプトゥスに近い部分ではより遠く離れており、

\(\boxed{\text{ここに続く証明や論述が続きます}}\)

\textbf{第2編 第2節：速度の二乗に比例する抵抗を受ける物体の運動（命題V〜XIV）}

部分 \( a \) から頂点より遠い部分へ向かうとき、これらの部分は頂点に近づく速度の比が双曲線 \( y = \frac{k}{x} \) のような双曲線と等しい比で進む。しかし、これら二つの双曲線間の差異は実際の物理現象において問題にならないほど小さいため、実践的な観点からはこれらの部分を双曲線として近似することができます。そして、より有用なのはこの近似双曲線よりも正確だが複雑な双曲線である可能性が高い。しかしながら、これらはその用途のために導き出される。

平行四辺形 \( XYGT \) を完成させます。そして、これらの双曲線の性質から容易に得られるように、直線 \( GT \) が双曲線を点 \( G \) で接していることがわかります。したがって、媒質中の密度 \( \rho \) 在 \( G \) の値は接線 \( GT \) の逆数と等しいです。

\[
\rho_G = \frac{1}{GT}
\]

\& 速度は \(\surd \frac{GTq}{GV}\) に比例し、抵抗は重力の力を \(GT\) に対して \(\frac{3nn + 3n}{n + 2} GV\) に比例します。

したがって、物体が点 \(A\) から直線 \(AH\) 上に投げ出され、双曲線 \(AGK\) を描き、\(AH\) の延長が漸近線 \(NX\) と交わる点 \(H\)、そして \(AI\) の延長がもう一つの漸近線 \(MX\) と交わる点 \(I\) に達する場合、

点 \(A\) の密度は \(AH\) に対して反比例し、物体の速度は \(\surd \frac{AHq}{AI}\) に比例し、そして同地点での抵抗は重力に対して \(AH\) に対して \(\frac{3nn + 3n}{n + 2} AI\) に比例します。従って、次のような規則が導けます。

\textbf{規則1.}\wsp{}もし点 \(A\) の密度を保ちつつ角度が変更される場合、

\(\textit{Reg.\,1.}\) もし点 \(A\) の密度を保持しつつ角度が変更される場合、

\[
NAH, \text{manebunt longitudines } AH, AI, HX.
\]
ideo si longitūdēs illāe in aliquo casu inveniantur, hyperbola deinceps ex dato quovis angulo $NAH$ expedite determinari potest.

\textbf{規則2.} \textit{Si servetur tum angulus } $NAH$, \textit{tum medium densitas in } $A$, \& mutetur velocitās quacum corpus projicitur; \textbf{manebit} longitūdō $AH$, \textbf{et} mutabitur $AI$ in duplicata ratione velocitatis reciproce.

\textbf{規則3.} \textit{Si tam angulus } $NAH$ \textit{quam corporis velocitās in } $A$, gravitas \textbf{et} acceleratrix servētur, \& proportionēs resistentiae in $A$ ad

\[
\textbf{規則4.} \textit{mutabitur ut eamdem sit cum ratione velocitatis reciproce.}
\]


引力が任意の比率で増加すると、比例 \( AH \) に対する \( AI \) の比も同じ比率で増加し、円錐曲線の斜辺が直角保たれるとき、かつその長さに比例する \(\frac{AH}{AI}\) も同様に増加します。したがって、\( AH \) は同じ比率で減少し、\( AI \) もその二倍比率で減少します。

\textbf{真実の定理4.} 一方で、抵抗に対する比も質量に対して増加します。これは、特定の密度が小さくなるか、媒質の密度が大きくなるか、または質量が減少するとき、抵抗はより少ない比率で減少します。

\[
\textbf{証明：} 媒質の密度が円錐曲線の頂点近くで等しい量になるとき、
\]



この証明は、媒質の密度が円錐曲線の頂点近くで等しい量になる場合を考慮します。

\textbf{第2編 第2節：速度の二乗に比例する抵抗を受ける物体の運動（命題V〜XIV）}


質量が中庸である点 \( A \)より大きい量 \( GT \)と直線 \( AH \)上の接線 \( AH \)との接線比 \( GT \)を求めるべきであり、また点 \( A \)における密度は第3則により、接線 \( AH \)の半分よりも僅かに小さい比率で減少するべきである。

\textit{命題5.} 与えられた長さ \( AH \), \( AI \)があり、図形 \( AGK \)を描く必要がある。直線 \( HN \)を延長して \( X \)まで引き、\( HX \)を \( n + 1 \)乗した値と \( AI \)で等しくする。点 \( A \)を通って中心 \( X \)と漸近線 \( MX \), \( NX \)を用いて双曲線を描き、その定義は任意の接線 \( VG \)に対する比 \( AI \)が \( XV^n \)に対する \( XI^n \)の比であるとする。

\textbf{命題VI. 定理6.}

速度の二乗に比例する抵抗を受ける物体の運動において、物体が点 \( A \) から上昇するとき、即ち点 \( G \) に向かって下落するとき、\( n \) の数値が大きいほど、双曲線はより正確に近似し、小さいほど双曲線はより不正確に近似します。逆もまた然りです。円錐双曲線は中庸の比を保ち、他の双曲線よりも単純です。したがって、もしこの種の双曲線が存在する場合、物体が直線 \( AN \) によって表される任意の直線上の点 \( K \) に到達し、\( A \) を通る直線を通過したとすると、点 \( M \) および \( N \) で直線 \( AN \) が漸近線 \( MX \) および \( NX \) と交わることを求め、線分 \( NK \) は線分 \( AM \) に等しいとします。


\par\begin{center}
    \centering
    \includegraphics[width=809pt,height=444pt,keepaspectratio]{280.png}

\end{center}\par

\textbf{訳文}

7．次に、これらの現象からこの双曲線を決定する便利な方法が明らかになります。相似で等しい物体が同じ速度で異なる角度 \(HAK\) および \(hAk\) から水平面に投げ出され、それぞれ点 \(K\) と \(k\) で当たるとき、\(AK\) に対する \(Ak\) の比を表す数 \(d\) と \(e\) を定めます。次に、任意の長さの直線 \(AI\) 上に垂直に点 \(H\) および \(h\) を取ります。これらの点から \(K\) および \(k\) までの距離 \(AK\) および \(Ak\) の長さを図的に求めるには、第6命題の方法を使用します。もし \(AK\) に対する \(Ak\) の比が \(d\) に対する \(e\) の比と等しければ、点 \(H\) から \(A\) までの距離 \(AH\) は正しく求められます。

\textbf{図}





\begin{center}
\textbf{命題 V（定理 5）}

直線 \( SM \) の無限長さを仮定し、もし \( SM \) の長さが仮定された \( AH \) に等しいとすると、任意の直線 \( MN \) を垂直に引き、直線 \( AK \) と \( Ak \) の比 \( \frac{AK}{Ak} - \frac{d}{e} \) で与えられた任意の直線の長さに比例する長さ \( MN \) を得る。同様の方法により、多数の点 \( N \) を求めることができる：そしてその後、曲線 \( NNXN \) を通ってすべてを描くと、この曲線は求める \( AH \) の長さに等しい \( SX \) を切り離す。機械工学のために、各角度 \( HAK \) において \( AH \) と \( AI \) の長さが同じであることが十分である。
\end{center}


\par\begin{center}
\centering
\includegraphics[width=0.8\textwidth]{280.png}
\label{fig:280}

\end{center}\par

\textbf{保持する．}図形を用いて抵抗を求めることにした場合、中間的な長さ$AH$，$HX$が既に求められているので、直線$AH$の位置を求めるには、命題Ⅷにより、点$A$と$K$から垂直に水平線に対して立たせた直線$AC$，$KF$の位置を求め、$AC$は下向きに伸び、かつ等しい長さで$AI$と一致するようにします．

\newpage

\pngright{275.png}{}{}{}

\textit{\textbf{命題Ⅷ．}}已知速度発射的物体在与定点$K$落下的位置，若已求得長度$AH$和$HX$；〜のために确定直線$AH$的位置，〜によれば図示，従点$A$和$K$分别作水平線的垂線$AC$和$KF$，其中$AC$向下延伸，かつ其長度与$AI$相等．

\[
\text{すなわち}\quad \frac{1}{2}HX.
\]

漸近線として
\[
AK, \quad KF
\]
を描き、接点が点 \( C \) である双曲線の共役軸を描きます。中心は \( A \)，半径は \( AH \) の円は点 \( H \) で双曲線と交わります。そして，射出された物体は直線 \( AH \) 上から発射され，点 \( K \) に到達します。\QEIit

なぜなら，点 \( H \)，\( AH \) の長さが与えられているので，円周上の任意の位置で定められます。線分 \( CH \) を引くと，この線分は \( AK \) および \( KF \) と交わり，それぞれ点 \( C \)，\( F \) に到達します。また，平行線 \( CH \)， \( MX \) および等しい長さの \( AC \)， \( AI \) から， \( AE \) も \( AM \) と同じ長さとなります。\(\blacksquare\)

\textbf{命題V}

物体が速度の二乗に比例する抵抗を受けるとき、その運動について次のようなことが言える。

\textit{ KN} と等しい $CE$ および $FH$ は比において互いに等しい。すなわち、$CE$ は $AE$ に対して $FH$ は $KN$ に対して同じ比をもつ。したがって、$CE$ および $FH$ の大きさは等しい。

点 $H$ は、直線 $AKF$ 上の双曲線の漸近線 $AK$, $KF$ との交点となる。この双曲線の共役軸は点 $C$ を通る。そして、これらの双曲線と円が共有する交点を求めるために、上記の条件から得られる交点 $H$ が存在することが確認される。

\textit{注意すべきは}、この操作は直線 $AKN$ が水平面に平行である場合でも、または任意角度で傾斜している場合でも同様に行われる。これは、二つの交点 $H$, $H$ から生じる角 $\angle NAH$, $\angle NAH$ のうち小さい方を選べばよいことからも明らかである。そして、機械工学においては、

\textbf{命題VI}

これらの結果は、物体が速度の二乗に比例する抵抗を受ける場合でも同様に適用される。

\textit{注意すべきは}、この操作は直線 $AKN$ が水平面に平行である場合でも、または任意角度で傾斜している場合でも同様に行われる。これは、二つの交点 $H$, $H$ から生じる角 $\angle NAH$, $\angle NAH$ のうち小さい方を選べばよいことからも明らかである。

\textbf{命題VII}

物体が速度の二乗に比例する抵抗を受ける場合でも、これらの結果は同様に適用される。

\textbf{命題 V.}

円を一度描き、その後無限長の規則 \(CH\) を点 \(C\) に適用し、その部分 \(FH\) と円弧 \(FK\) 及び直線 \(HK\) の間にある部分が点 \(C\) における部分 \(CE\) と等しくなるようにする。

\begin{center}
\includegraphics[width=960px,height=1246px,viewport=-24 0 960 1246]{281.png}
\end{center}

\textbf{命題 VI.}

上記について、双曲線に関するものも容易に適用できる。すなわち、直線 \(XV\) の頂点 \(X\) における接線を表す平行四辺形 \(XAGK\) を考える。各軸の部分 \(IA\), \(VG\) を任意の \(abscissarum XI, XV\) のべき \(XI^n\), \(XV^n\) とし、直線 \(XT\), \(TG\), \(HA\) を引く。

\textbf{命題 VII.}

同様に、放物線についても適用できる。すなわち、直線 \(XV\) の頂点 \(X\) における接線を表す平行四辺形 \(XAGK\) を考える。各軸の部分 \(IA\), \(VG\) を任意の \(abscissarum XI, XV\) のべき \(XI^n\), \(XV^n\) とし、直線 \(XT\), \(TG\), \(HA\) を引く。

\textbf{命題 VIII.}


\textbf{命題 IX.}


\textbf{命題 X.}


\textbf{命題 XI.}

同様に、放物線についても適用できる。すなわち、直線 \(XV\) の頂点 \(X\) における接線を表す平行四辺形 \(XAGK\) を考える。各軸の部分 \(IA\), \(VG\) を任意の \(abscissarum XI, XV\) のべき \(XI^n\), \(XV^n\) とし、直線 \(XT\), \(TG\), \(HA\) を引く。

\textbf{命題 XII.}


\textbf{命題 XIII.}


\textbf{命題 XIV.}


$XT$ が $VG$ に平行であり、$TG$, $HA$ がそれぞれ点 $G$ と $A$ で放物線を接している場合、物体が任意の位置 $A$ から直線 $AH$ の延長線上に沿って正当な速度で投げ出されると、その軌跡はこの放物線を描きます。ただし、各点 $G$ における媒質の密度が接線 $GT$ に対して逆比例するものとします。そして、点 $G$ での速度は、抵抗のない空間において物体が通過する速度と同じになります。

この場合、点 $G$ での放物線の頂点は $VG$ の延長線上にある $G$ 点であり、その焦点からの距離（直径）は $VG$ です。また、放物線の実軸は
\[
\surd \frac{2TGq.}{\overline{nn - n} X VG}
\]
となります。そして、点 $G$ での抵抗は重力に対する比として $TG$ に比例します。

\[
\frac{3n n - 3n}{n - 2} VG \,.\] ここで、もし水平線 \(NAK\) を示すとし、質量密度 \(M\) が点 \(A\) で一定であり、物体を投げ出す速度が変化したとき、角度 \(\textbf{NAH}\) のうちいずれかが変更されるとしても、長さ \(AH\), \(AI\), \(HX\) は不変となります。従って、放物線の頂点 \(X\) と直線 \(XI\) の位置が与えられます。さらに、\(VG\) に対する \(IA\) の比を \(XV^n\) に対する \(XI^n\) の比として定めることにより、すべての放物線の点 \(G\) が得られます。これにより、投射体が通過する全ての点が特定されます。


\section{第2編 第3節：速度とその二乗の和に比例する抵抗を受ける運動（命題XV〜XVI）}


% Section: 第2編 第3節：速度とその二乗の和に比例する抵抗を受ける運動（命題XV〜XVI）



第三編 第三節

\begin{center}
\textbf{速度とその二乗の和に比例する抵抗を受ける運動（命題XV〜XVI）}
\end{center}

\bigskip

\textbf{\underline{命題XI}}．理論 VIII．

\medskip

\textit{
もし物体が速度に比例し、またその二乗に比例して抵抗を受け、そして媒介物において同一の内在的力を受けて運動する場合、}

進歩時間が算術級数で与えられたと仮定すると、

速度に反比例する量が一定量を増加させば、

これらの量は幾何級数で与えられる。


\bigskip

\par\begin{center}
  \centering
  \includegraphics[width=10.12cm,height=10.61cm,trim={0 -24 0 0},clip]{283.png}

\end{center}\par

\begin{center}
\textbf{訳文}

中心 \( C \) と漸近線 \( CADd \) および \( CH \) 上に、双曲線 \( BEeS \) を描きます。また、漸近線 \( CH \) に平行な直線 \( AB \) および \( DE \) を引きます。直線 \( CD \) 上に点 \( A \) および \( G \) を取ります。そして、時間が双曲線 \( ABED \) の面積を一様に増加させるとき、速度は長さ \( DF \) で表すことができ、その逆数 \( GD \) と与えられた直線 \( CG \) の和は直線 \( CD \) 上の等比級数増加する長さ \( CD \) を構成します。

\end{center}

\begin{center}
\textbf{訳文（続）}

すなわち、時間 \( t \) における面積 \( DEed \) の最小増分を \( dA \) とすると、\( Dd \) は \( DE \) の逆数に比例し、従って直線 \( CD \) 上の点 \( D \) から点 \( d \) までの距離 \( Dd \) は \( DE \) に直接比例します。

\end{center}


\par\begin{center}
\centering
\includegraphics[width=0.8\textwidth]{283.png}

\end{center}\par

\[
CD. \textbf{ipsius} \frac{1}{GD} \textbf{decrementum}, quod (per hujus Lem.\ II.) est \frac{Dd}{GDq.}, erit ut \frac{CD}{GDq.} seu \frac{CG + GD}{GDq.}, id est, ut \frac{1}{GD} + \frac{CG}{GDq.}.
\]
したがって、時間 \( ABED \) において、各部分 \( EDde \) を加えると、速度は一様に増加し、一方で \( \frac{1}{GD} \) は同じ比率で減少します。

実際、速度の減少量は抵抗に比例しており、これは仮定により、二つの量の合計に比例します。その一つは速度に比例し、もう一つは速度の平方に比例するからです。したがって、\( \frac{1}{GD} \) の減少量もこれらの量の合計に比例します。

\[
\frac{1}{GD} \& \frac{CG}{GDq.}, \textbf{quarum prior est ipsa } \frac{1}{GD}, \textbf{\& posterior } \frac{CG}{GDq.} \textbf{ est ut } \frac{1}{GDq.}. \textbf{ Proinde } \frac{1}{GD}, \textbf{ ob analogum decrementum, } \textbf{est ut velocitas. Et si quantitas } GD \textbf{ ipsi } \frac{1}{GD} \textbf{ reciproce proportionalis quantitate data } CG \textbf{ augeatur, summa } CD, \textbf{ tempore } ABED \textbf{ uniformiter crescente, crescet in progressione Geometrica.}
\]

\textit{\textbf{Corol.\fsp{}1.}}\wsp{}したがって、与えられた点 $A$, $G$ において、時間 $ABED$ を双曲線形面積 $ABED$ で表すならば、速度は $GD$ の逆数 $\frac{1}{GD}$ で表されることが可能となる。

\textbf{補則2．}しかし、速度の逆数 \( GA \) と \( GD \) の比を初速における速度の逆数として取ると、任意の時間区間 \( ABED \) 終了時の速度の逆数が点 \( G \) となることが分かります。この点 \( G \) を得た後では、任意の他の時間においてはその速度も求めることができます。

\textbf{命題XII．定理IX．}

\textbf{同様に置き立てば、私は次のように主張します：描かれた距離が等差数列をなすとき、ある一定量によって増加した速度は等比数列をなすでしょう。}

直線 \( CD \) の漸近線上に点 \( R \) を取ります。さらに、垂線  を引き、以下に述べる論証を進めます。

\textbf{RS}が双曲線$S$に遭遇するとき、描かれた距離は双曲形面積$RSED$と一致し、速度は直線$GD$の長さと比例し、与えられた直線$CG$との合成で減少する幾何級数を形成します。ここで、空間距離$RSED$が算術的に増加すると仮定します。

実際には、与えられた空間距離の増分$EDde$に対応する線段$Dd$は、直線$GD$の減少量であり、$ED$と反比例し、また$CD$と直接比例します。すなわち、速度の減少量は、その時間に反比例して与えられた空間部分$DdeE$が描かれるとき、以下の関係で表されます：

\[
\text{速度の減少量} \propto \frac{1}{ED}
\]

これは、直線$GD$と与えられた直線$CG$の和に比例します。

抵抗と時間は連続的に比例し、すなわち速度の和に直接比例し、速度の二乗に反比例します。また、速度自体とも反比例します。したがって、速度の減少量および線分 \( GD \) の減少量は、与えられた量と減少する量の積に直接比例します。そして、相似な減少量は常に相似な減少量を生じます：すなわち、速度と線分 \( GD \) です。\(\blacksquare\)

\textit{Corol.\,1.} よって、速度を長さ \( GD \) で表した場合、



抵抗 \& 時間 conjunctim, id est directe ut summa duarum
quantitatum, quarum una est velocitas, altera ut velocitatis
quadratum, \& inverse ut velocitas; adeoque directe ut summa dearum
quantitatum, quarum una datur, altera est ut velocitas.
Igitur decrementum tam velocitatis quam line{\ae} $GD$, est ut quantitas
data \& quantitas decrescens conjunctim, \& propter analoga
decrementa, analog{\ae} semper erunt quantitates decrescentes: nimirum
velocitas \& linea $GD$. $\blacksquare$

\textit{Corol.\,1.}\wsp{}Igitur si velocitas exponatur per longitudinem $GD$,


\textbf{命題XV}：描かれた距離は双曲線面積 \(DESR\) に比例する。

\textit{\textbf{推論2}}：任意の点 \(R\) を仮定した場合、点 \(G\) を求めることができます。ここで、\(GD : GR\) は初速度に対する任意の距離 \(ABED\) の後の速度の比とします。点 \(G\) を求めることが出来れば、与えられた速度から距離を決定し、逆に距離から速度を決定することも可能です。\pngright{285.png}{}{}{}

\textit{\textbf{推論3}}：したがって、命題XIにより与えられた時間から速度を求めることができ、この命題により与えられた速度から距離を求めることができます。従って、与えられた時間から距離を求めることも可能です。\pngright{285.png}{}{}{}

\textbf{命題XVI}\@. 理論X。

\[
\begin{array}{c}
\textbf{図}285
\end{array}
\]

次に、重力によって直線的に上昇または下降する物体に対して、抵抗は速度とその二乗の比例に作用すると仮定します。このとき、私は次のように主張します：円と双曲線の直径方向に並行な直線を共役直径の端点を通じて引くと、これらの直線からある定点まで引かれた区間の長さが速度と比例する場合、時間は各セクターの面積と比例します。逆もまた然りです。

\textbf{例1.}まず、物体が上昇していると仮定し、中心を \( D \) とするとき、\ldots

\& 任意の半径 \( DB \) 上に四分円 \( BETF \) を描き、半径 \( DB \) の端点 \( B \) から無限遠方に至る平行線 \( BAP \) を引き、半径 \( DF \) と平行な位置に点 \( A \) を定めます。また、速度の比例因子としてセグメント \( AP \) を選びます。
そして、抵抗の一部が速度とその二乗の比であるとするとき、全体の抵抗は \( P \) 点において \( AP^2 + 2PAB \) の比で与えられます。
半径 \( DA \) と半径 \( DP \) が円を交わる点 \( E \) および \( T \) を結び、重力 \( DA^2 \) の比を定義し、重力を抵抗 \( P \) 点における比 \( AP^2 + 2PAB \) に比例するものとするとき、

\[
\textbf{重力} : \textbf{抵抗} = DA^2 : AP^2 + 2PAB
\]

\textbf{昇降時間の全時間は、円弧 \(EDTE\) の扇形と等しい。}


\par\begin{center}
  \centering
  \includegraphics[width=1519px,height=1396px,trim=-24 0 0 0,clip]{286.png}

\end{center}\par

次に、速度 \(AP\) の動量 \(PQ\) を切断し、扇形 \(DET\) の動量 \(DTV\) が与えられたとき、時間の動量に対する応答として、速度の減少 \(PQ\) は重力の力を表す \(DBq\) と抵抗の力を表す \(APq + 2BAP\) の和に比例します。これは \textbf{第2編 第12命題}により、つまり \(DP^2\) に比例します。したがって、面積 \(DPQ\) は \(PQ\) に対して比例し、面積 \(DTV\)（これは \(DPQ\) の面積に対する比が \(DTq\) に対して \(DPq\) の比である）は与えられた \(DTq\) に比例します。


速度とその二乗の和に比例する抵抗を受ける運動において、領域 $EDT$ は時間の進行とともに等速的に減少し、これは具体的な部分領域 $DTV$ の差分によって説明されます。従って、昇天時間に対する比例関係が成立します。\QED

\textit{案2.} 体の昇降における速度を $AP$ の長さに比例させ、抵抗を $APq + 2BAP$ とするとき、重力の強度が $DAq$ で表すことができない場合、$BD$ を $ABq - BDq$ の長さとし、これは重力に対する比例関係を持つものとします。また、頂点 $F$ を中心として、直線 $DF$ を垂直に引きます。



速度とその二乗の和に比例する抵抗を受ける運動において、領域 $\textbf{EDT}$ は時間の進行とともに等速的に減少し、これは具体的な部分領域 $\textbf{DTV}$ の差分によって説明されます。従って、昇天時間に対する比例関係が成立します.\QED

\textit{案2.} 体の昇降における速度を $\textbf{AP}$ の長さに比例させ、抵抗を $\textbf{APq + 2BAP}$ とするとき、重力の強度が $\textbf{DAq}$ で表すことができない場合、$\textbf{BD}$ を $\textbf{ABq - BDq}$ の長さとし、これは重力に対する比例関係を持つものとします。また、頂点 $\textbf{F}$ を中心として、直線 $\textbf{DF}$ を垂直に引きます。


双曲線 \(FTVE\) の半径軸の共役直径を \(DB\) と \(DF\) とする。これらは直線 \(DA\) を点 \(E\), \(DP\), \(DQ\) で交わし、それぞれ点 \(T\) と \(V\) に達する。すると、未来の上昇時間は双曲線セクター \(TDE\) の面積と比例する。

なぜなら、速度の減少量 \(PQ\) は与えられた時間内に粒子が経験した抵抗の総和 \(APq. + 2ABP\) および重力 \(ABq. - BDq.\) に比例し、これは \(BPq. - BDq.\) である。また、面積 \(DTV\) は面積 \(DPQ\) に対して \(DTq.\) は \(DPq.\) に対して比例するから、直線 \(DA\) 上の点 \(G\) から垂直に下ろした垂線 \(GT\) を引き、この垂線 \(GT\) の長さ \(GTq.\) や \(GDq. - DFq.\) は \(BDq.\) に対して比例し、さらに \(GDq.\) は \(PBq.\) に対して比例し、また部分的に \(DFq.\) は \(BPq. - DBq\) に対して比例する。

\[
\text{なぜなら}、面積  DPQ \text{ は} PQ \text{ に比例し、すなわち} BPq - BDq \text{ に比例しますから、}
\]
\[
\text{面積 } DTV \text{ は与えられた量 } DFq \text{ に比例します。したがって、面積 } EDT \text{ は時間の各部分に対して等しく減少し、}
\]

\begin{center}
\includegraphics[width=1168pt,height=1676pt,offset=-24pt]{287.png}
\end{center}

\textit{例3.} \textbf{} 速度  AP \text{ は物体の下降における速度を表し、} APq + 2ABP \text{ は抵抗を表し、} DBq - ABq \text{ は角度 } DAB \text{ が直角であるときの重力力を表します。中心点 } D, \text{ 頂点 } B \text{ を持つ双曲線}
\[
\textbf{D} 
\]

\textbf{直角} \( BETV \) が直線 \( DA \), \( DP \), および \( DQ \) の積 \( DA \cdot DP \), \( DA \cdot DQ \) を分割し, \( E \), \( T \), および \( V \) に交わるとき, この放物線のセクター \( DET \) は下降時間に比例する。

なぜなら, 速度の増加量 \( PQ \) とその比となる面積 \( DPQ \) は, 重力の超過部分（抵抗を除いた）に比例し, 即ち
\[ DB^2 - AB^2 - 2ABP - AP^2 \]
または
\[ DB^2 - BP^2 \]
と等しい。また, 面積 \( DTV \) は面積 \( DPQ \) に対して
\[ DT^2 : DP^2 \]
であり, さらに
\[ GT^2 \] すなわち
\[ GD^2 - BD^2 \] と \( BP^2 \) の比も等しく,
\[ GD^2 \] と \( BD^2 \) の比も同様に
\[ BD^2 \] と \( BD^2 - BP^2 \) の比で割ったものと等しい。したがって, 面積 \( DPQ \) は \( BD^2 - BP^2 \) に比例するとき, 面積 \( DTV \) も与えられた \( BD^2 \) とともに増加する。

\[
\textbf{よって} \, EDT \, の各時間部分に対する面積 \, EDT \, \textbf{は、} \, DTV \, \textbf{の} \, n \, \textbf{個の等しい部分を加算することで得られ、したがって} \, EDT \, \textbf{に比例する。} \, \Box
\]

\textit{\textbf{命題}} \, I. \, よって、速度 \, AP \, \textbf{は、非抵抗媒質中で} \, EDT \, \textbf{の時間に上昇または下降するとき獲得または失うことができる速度} \, \textit{AP} \, \textbf{に対して、三角形} \, DAP \, \textbf{の面積に対するセクター} \, D \, \textbf{中心} \, DA \, \textbf{から} \, ADT \, \textbf{の角度で描かれた面積} \, \textit{sectoris} \, \textbf{に比例する。したがって、与えられた時間から速度が得られる。非抵抗媒質中での速度は、時間と同様にセクター} \, D \, \textbf{に比例する。}

\[
\textbf{抵抗は三角形に比例し、}
\]
\[
\textbf{中点では最小値で、}
\]
\[
\textbf{そのとき抵抗は質量の等しいものに近づく。}
\]
\[
\textbf{これは扇形と三角形の慣習に従います。}
\]

\textbf{(第14命題. 問題 IV.)}

\textit{\textbf{同様にして、私は言います：上昇または下降中描かれた距離は、}
\[
\textbf{時間によって指数化された面積の和または差に比例し、}
\]
\[
\textbf{ある他の面積が算術級数で増加または減少するとき。}
\]}

\textbf{もし重力と抵抗の合力が幾何級数で増加または減少する場合。}


\[
\textit{\textbf{AC (図3最後の3つ)は重力、AKは抵抗を表す}}
\]

\textbf{第15命題}

物体が上昇するとき、または下降するとき、その速度の二乗に比例する抵抗を受ける運動について考える。点 \( A \) から等しい部分に分割すると、物体が上昇するときには反対方向へ分割される。引く線 \( Ab \) を引き、\( Ab \) は \( DB \) に対して \( DBq. \) と \( 4BAC \) の比で比例し、領域 \( AbNK \) は進歩的な算術級数において増加または減少する。ここで力 \( CK \) が進歩的な幾何級数に従って取り扱われる。したがって、物体の位置からその最大高度までの距離は、領域 \( AbNK \) の余分部分 \( DET \) よりも大きい領域 \( AbNK \) の差に比例すると主張する。

なぜなら、線 \( AK \) は抵抗に対して比例し、すなわち \( APq. + 2BAP \) に対して比例する。任意の量 \( Z \) を仮定し、\( AK \) を \( \frac{APq. + 2BAP}{Z} \) と設定すると（この命題の補題IIによると）、線 \( AK \) の動量は次のようになる。

$KL = \frac{2APQ + 2BA \times PB}{Z} \text{ または } \frac{2BPQ}{Z},$
これらの量は $\textbf{AbNK}$ の運動量 \(KLON\) が
$\frac{2BPQ \times LO}{Z} \text{ または } \frac{BPQ \times BD^{\text{cub.}}}{2Z \times CK \times AB}.$

\textbf{1. 案：} 働体が上昇するとき、重力 \(ABq.\) + \(BDq.\) が存在し、円錐 \(BET\) の中心軸である場合（図案1, 第十三命題参照），線分 \(AC\) は重力に比例しているので，
$$\frac{ABq. + BDq.}{Z}$$
そして \(DPq.\) または \(APq. + 2BAP + ABq. + BDq.\) は \(AK \times Z + AC \times Z\) すなわち \(CK \times Z\) となる。従って面積 \(DTV\) は面積 \(DPQ\) に比して
$$\frac{DTq.}{1}.$$

\[
\textbf{案}\textbf{: 2.} \quad \text{物体が上昇し、重力は} \, AB^2 - BD^2 \, \text{に比例するとき、直線} \, AC \, (\text{図: 案 2. 定理 XIII}) \, \text{は} \, \frac{AB^2 - BD^2}{Z} \, \text{であり、} \, DT^2 \, \text{は} \, DP^2 \, \text{に対して} \, \frac{DF^2}{DB^2 + BP^2 - BD^2} \, \text{である。}
\]
ここで、
\[
DT^2 : DP^2 = DB^2 : (BP^2 - BD^2) \quad \text{または} \quad DT^2 : DP^2 = DB^2 : (AP^2 + 2BAP + AB^2 - BD^2)
\]
\[
= DB^2 : (AK \times Z + AC \times Z) \, \text{または} \, DB^2 : CK \times Z.
\]
したがって、面積 $DTV$ は面積 $DPQ$ に対して
\[
DB^2 : CK \times Z
\]
に比例します。

\[
\textbf{案}\textbf{: 3.} \quad \text{同様の議論により、物体が下落し、重力は} \, BD^2 - AB^2 \, \text{に比例するとき、直線} \, AC \, (\text{図: 案 3. 定理 XIII}) \, \text{は} \, DB^2 : CK \times Z \, \text{に比例します。}
\]

次に、前節で述べたとおり、

\[
\frac{BD^q - AB^q}{Z}
\]

は領域 \(DTV\) が領域 \(DPQ\) に比して \(DB^q : CK \times Z\) の比である。すなわち、上記の通り。

したがって、これらの領域は常にこの比で等しい。もし領域 \(DTV\) を時間の動量とみなして定数の矩形、例えば \(BD \times m\) と表すならば、

\[
\textbf{領域} DPQ, \textit{i.e.}, \frac{1}{2} BD \times PQ; \quad \text{は} \quad BD \times m \text{に対する比}
\]

が \(CK : Z \times DB^q\) である。従って、

\[
PQ \text{ in } BD = CK \times Z
\]

と等しいことがわかり、さらに領域 \(AbNK\) の動量は上部の \(KLON\) に対する

\[
2BD \times m \times CK \times Z
\]

$\frac{BP \times BD \times m}{AB}$ が得られます。この量から面積 $DET$ の動量 $DTV$，すなわち $BD \times m$ を引くと，
$\frac{AP \times BD \times m}{AB}$ が残ります。したがって，面積の差の動量の差は
$\frac{AP \times BD \times m}{AB}$ に等しくなります；従って（与えられた $\frac{BD \times m}{AB}$ の値から）
上昇または下降するときの物体が描く空間の速度 $AP$，すなわちその動量と比例します。したがって，面積の差は増加または減少する動量に比例し，同時に発生したり消えたりすることもあります。

\textbf{証明}

したがって、任意の長さ \( V \) を、線分 \( DA \) に対する弧 \( ET \) の比、すなわち線分 \( DA \) に対する線分 \( DE \) の比で取るとき、物体が完全に抵抗のある媒質中で上昇または下降する間に描く距離は、同じ時間内で完全に抵抗のない媒質中で描ける距離に対して、ある面積の差と比例します。この差は \(\frac{BD \times V^2}{4AB}\) であり、したがって与えられた時間から与えられる距離が得られます。

実際、完全に抵抗のない媒質中での距離は時間の二倍比に比例し、すなわち \( V^2 \) と比例します。そして、与えられた \( BD \) および \( AB \) の値から、\(\frac{BD \times V^2}{4AB}\) が得られます。また、時間は \( DET \) に比例します。


\[
\textbf{すなわち}\, \frac{1}{2}BD \times ET \,\textbf{の面積のモーメントは}\, \frac{BD \times V}{2AB} \,\textbf{と}\, \frac{1}{2}BD \,\textbf{のモーメントに比例し、}
\]

\[
\textbf{すなわち}\, \frac{BD \times V}{2AB} \,\textbf{と}\, \frac{DAq. \times 2m}{DEq.} \,\textbf{および}\, \frac{1}{2}BD \times 2m,
\]

\[
\textbf{すなわち}\, \frac{BD \times V \times DAq. \times m}{AB \times DEq.} \,\textbf{と}\, BD \times m.
\]

\[
\textbf{したがって}\, 面積\, V^2 \,\textbf{のモーメントは}\, AKNb \,\textbf{および}\, DET \,\textbf{の差分面積のモーメントに比例し、}
\]

\[
\frac{BD \times V \times DA \times m}{AB \times DE} \,\textbf{と}\, \frac{AP \times BD \times m}{AB}
\]

\[
\textbf{すなわち}\, \frac{BD \times V \times DA \times m}{AB \times DE} \,\textbf{と}\, \frac{AP \times BD \times m}{AB}.
\]

\[
\frac{\textbf{V} \times \textbf{DA}}{\textbf{DE}} \propto \textbf{AP}; \text{ したがって、}\, \textbf{V} \& \textbf{AP} \,\text{が最小のとき、}
\]
これらは等しい比を成す。

したがって、最小の面積
\[
\frac{\textbf{BD} \times \textbf{V}^2}{4\textbf{AB}}
\]
は、$\textbf{DET}$ および $\textbf{AKNb}$ の面積の差に等しい。

ここで、下り降りや上昇の始点または終点で描かれた空間が両方とも中間において等しくなるとき、これらは等しい比を成す。したがって、そのときこれらの面積は互いに
\[
\frac{\textbf{BD} \times \textbf{V}^2}{4\textbf{AB}} \quad \text{および}\quad \textbf{DET} \& \textbf{AKNb} の差
\]
と等しい。これらの比は類似の増加を伴うため、任意の等しい時間においてこれらは互いにその面積と等しい。

\[
\frac{BD \times V^2}{4AB} \textbf{ と } \textit{area} \, DET \textbf{ および } \textit{area} \, AKNb \textbf{ の差} \textbf{ は等しい。}
\]



この証明により、証明終了を表します。


\section{第2編 第4節：抵抗媒体中の物体の円運動および螺旋運動（命題XVII〜XVIII）}


% Section: 第2編 第4節：抵抗媒体中の物体の円運動および螺旋運動（命題XVII〜XVIII）



\sectpage{IV.}

\begin{center}\textbf{抵抗媒体中の物体の円運動および螺旋運動について。}\end{center}

\condpagelarge{命題 III.}

\textbf{設 $PQRr$ を螺旋とし、この螺旋がすべての半径 $SP$, $SQ$, $SR$, \ldots{} を等角的に交わすとする。直線 $PT$ が任意の点 $P$ でこの螺旋を接し、半径 $SQ$ との交点を $T$ とし、螺旋に垂直な直線 $PO$ および $QO$ の交点を $O$ とする。直線 $SO$ を引く。このとき、点 $P$ と $Q$ が互いに近づき、重なる場合、角 $\textbf{\textit{PSO}}$ は直角となり、最後の比 \(\frac{TQ \times PS}{PQ^2}\) は比 \(1\) となる。}

\begin{proof}
点 $P$ と $Q$ が互いに近づき、重なる場合を考える。螺旋 $PQRr$ の性質により、直線 $PT$ は螺旋を接しており、半径 $SQ$ との交点を $T$ とする。直線 $PO$ および $QO$ が垂直であり、これらは交点 $O$ を持つ。

まず、角 $\textbf{\textit{PSO}}$ の性質について考察する。螺旋の性質により、半径 $SP$, $SQ$, $SR$, \ldots{} は等角的に交わるため、直線 $PT$ が螺旋を接し、点 $P$ と $Q$ が近づくとき、直線 $SO$ は垂直に傾き、角 $\textbf{\textit{PSO}}$ は直角となる。

次に、比 \(\frac{TQ \times PS}{PQ^2}\) の性質について考察する。螺旋の性質により、点 $P$ と $Q$ が近づくとき、線分 $TQ$ および $PS$ の長さは一定であり、線分 $PQ$ の長さは減少する。したがって、比 \(\frac{TQ \times PS}{PQ^2}\) は比 \(1\) に収束する。

以上より、点 $P$ と $Q$ が互いに近づき、重なる場合、角 $\textbf{\textit{PSO}}$ は直角となり、最後の比 \(\frac{TQ \times PS}{PQ^2}\) は比 \(1\) となる。
\end{proof}

\textbf{等角性。}


\par\begin{center}
    \centering
    \includegraphics[width=1222px,height=1162px,trim=-24 0 0 0]{292.png}

\end{center}\par

よって、直角 $\angle OPQ$ や $\angle OQR$ から等しい直角 $\angle SPQ$ や $\angle SQR$ を引くと、$\angle OPS$ および $\angle OQS$ は依然として等しい直角となる。したがって、点 $O$, $S$, $P$ を通る円は点 $Q$ も通過する。点 $P$ と点 $Q$ の接近により、この円は接地点 $PQ$ 上に存在し、螺旋線を垂直に切る。従って直線 $OP$ は円の直径となり、$\angle OSP$ は半円内の直角となる。

\textbf{証明終了。}

点 $O$ から直線 $OP$ に垂直な線分 $QD$ や $SE$ を引き、これらの線分の比について述べる。

\[
\textbf{最後のものは次のようになります: } TQ \textbf{ と } PD \textbf{ の比は } TS \textbf{ や } PS \textbf{ と } PE \textbf{、または } PO \textbf{ と } PS \textbf{ の比に等しいです。また、} PD \textbf{ と } PQ \textbf{ の比は } PQ \textbf{ と } PO \textbf{ の比に等しいです。したがって、} TQ \textbf{ と } PQ \textbf{ の比は } PQ \textbf{ と } PS \textbf{ の比に等しくなります。これより、} PQq. \textbf{ は } TQ \times PS \textbf{ に等しいです。}
\]


\textbf{命題 XVII. 理論 XI.}

\textit{\textbf{もし抵抗媒体の密度が各点で互易比である中心からの距離と反比例し、かつ向心力が密度の二倍比であるとき、物体は円錐螺旋線を描きつつ中心から引かれたすべての半径を一定の角度で交わることができる。}}


\textbf{命題 XVII. 定理 17}

上記第14節の命題において、$SQ$ を $V$ まで延長し、$SV = SP$ とします。等時間に弧 $PQ$ および $QR$ を描き、円錐曲線 $PSQ$ および $QSr$ の面積を等しくします。そして、中心力は物体が $P$ に引き寄せられる力を表し、この力は $SPq$ と反比例します（第1巻の命題 X）。

したがって、力によって生成された線分 $TQ$ は、その力と弧 $PQ$ の描画時間の二倍の比から成る複合的な比にあります。すなわち、

\[
TQ \times SPq
\]

これは

\[
(\text{第1巻の命題 X})
\]

により等しいです。

ここで、抵抗力はこの状況において無視することができます（すなわち、その影響が中心力に対して無限に小さいと仮定します）。


\textbf{証明} $\quad$ \text{最後の補題} $PQq. \times SP$ は時間の二倍比に比例し、従って時間は

$PQ \times \sqrt{SP}$

と等しい。すなわち、物体が時間 $t$ で弧 $PQ$ を描く速度は
 
$\frac{PQ}{PQ \times \sqrt{SP}} = \frac{1}{\sqrt{SP}}$

すなわち、\(SP\) の半分の比に反比例する。同様の議論により、弧 \(QR\) を描く速度は \(SQ\) の半分の比に反比例する。これらの弧 \(PQ\) と \(QR\) は互いの速度に対する比であり、すなわち
\[
\frac{1}{\sqrt{SP}} \quad \text{と} \quad \frac{1}{\sqrt{SQ}}
\]
が互いの比である。


\par\begin{center}
    \centering
    \includegraphics[width=0.8\linewidth]{292.png}

\end{center}\par

\textbf{命題 XVII}

円運動中の物体が抵抗媒体中を螺旋運動するとき、その物体の速度と半径との関係は、直線$SQ$に対する直線$SP$の比、すなわち$\sqrt{SQ}$に対する$\sqrt{SP} \times \sqrt{SQ}$の比に等しい。また、角$SPQ$と$SQr$が等しく、面積$PSQ$と$QSr$も等しいとき、弧$PQ$は弧$Qr$に対して直線$SQ$に対する直線$SP$の比で比例する。

\textbf{証明}

\begin{align*}
\text{arcus } PQ \colon \text{arcum } Qr &= SQ \colon SP, \\
&\text{すなわち} \quad \sqrt{SQ} \colon \sqrt{SP} \times \sqrt{SQ}.
\end{align*}

角$SPQ$と$SQr$が等しく、面積$PSQ$と$QSr$も等しいとき、弧$PQ$は弧$Qr$に対して直線$SQ$に対する直線$SP - SP^{\frac{1}{2}} \times SQ^{\frac{1}{2}}$の比で比例する。これは、点$P$と$Q$が重なるとき、最終的な比$SP - SP^{\frac{1}{2}} \times SQ^{\frac{1}{2}}$は$\frac{1}{2}VQ$に等しいからである。

\textbf{命題 XVIII}

抵抗媒体中で螺旋運動する物体の速度と半径との関係は、直線$SQ$に対する直線$SP - SP^{\frac{1}{2}} \times SQ^{\frac{1}{2}}$の比に等しい。すなわち、面積$PSQ$と$QSr$が等しいとき、弧$PQ$は弧$Qr$に対して直線$SQ$に対する$\frac{1}{2}VQ$の比で比例する。

\textbf{証明}

抵抗媒体中で螺旋運動する物体の速度と半径との関係は、面積$PSQ$と$QSr$が等しいとき、弧$PQ$は弧$Qr$に対して直線$SQ$に対する$\frac{1}{2}VQ$の比で比例することを示す。これは、点$P$と$Q$が重なるとき、最終的な比$SP - SP^{\frac{1}{2}} \times SQ^{\frac{1}{2}}$は$\frac{1}{2}VQ$に等しいからである。


\par\begin{center}
  \centering
  

\end{center}\par

\textbf{注記}

抵抗媒体中で螺旋運動する物体の速度と半径との関係は、面積$PSQ$と$QSr$が等しいとき、弧$PQ$は弧$Qr$に対して直線$SQ$に対する$\frac{1}{2}VQ$の比で比例することを示す。これは、抵抗媒体中で螺旋運動する物体の速度と半径との関係が、面積$PSQ$と$QSr$が等しいとき、弧$PQ$は弧$Qr$に対して直線$SQ$に対する$\frac{1}{2}VQ$の比で比例することを示す。

\textbf{時間等しい間隔において描かれるべき円周上の弧の差分}は、抵抗によって生じる面積 \(RSr\) から導き出される。従って、抵抗は線素 \(Qr\) の減少量と生成に必要な時間の平方との比である。

実際、線素 \(Rr\)（命題I, 葉Xによる）は時間の二倍比に比例する。したがって、抵抗は \(\frac{Rr}{PQq. \times SP}\) となる。

また、\(PQ\) に対する \(Rr\) の比は \(SQ\) に対する \(\frac{1}{2}VQ\) の比であるから、\(\frac{Rr}{PQq. \times SP}\) は \(\frac{\frac{1}{2}VQ}{PQ \times SP \times SQ}\) と等しくなり、さらにこれは \(\frac{\frac{1}{2}OS}{OP \times SPq.}\) となる。ここで \(P\) および \(Q\) の点は考慮される。

coeuntibus, \textbf{$SP$} \& \textbf{$SQ$} coincidunt; \& ob similia triangula \textbf{$PVQ$}, \textbf{$PSO$}, fit \textbf{$PQ$} ad \(\frac{1}{2}VQ\) ut \textbf{$OP$} ad \(\frac{1}{2}OS\). Est igitur \(\frac{OS}{OP \times SP}\) ut
resistentia, id est in ratione densitatis Medii in \(P\) \& ratione
duplicata velocitatis conjunctim. Auferatur duplicata ratio velocitatis,
nempe ratio \(\frac{1}{SP}\), \& manebit Medii densitas in \(P\) ut
\(\frac{OS}{OP \times SP}\).

Detur Spiralis, \& ob datam rationem \textbf{$OS$} ad \textbf{$OP$}, densitas
Medii in \(P\) erit ut \(\frac{1}{SP}\). In Medio igitur cujus densitas
est reciproce ut distantia a centro \(SP\), corpus gyrari potest in hac



螺旋運動が可能となる媒体の密度は、中心からの距離 \(SP\) に反比例するとき、物体はこの媒体中を円運動を行うことができます。

Spirali.

\textbf{命題 XVII.}\wsp{}証明終了。

\textit{推論 1.}\wsp{}任意の点 \( P \)において物体が抵抗媒体に阻まれない場合、その速度は常に円周上での運動と同様である。すなわち、中心から点 \( P \)までの距離 \( SP \) に等しい。

\textit{推論 2.}\wsp{}媒体の密度は、距離 \( SP \) が与えられた場合、\( \frac{OS}{OP} \) と比例する。もしその距離が与えられなければ、\( \frac{OS}{OP \times SP} \) と比例する。したがって、任意の媒体密度に対応する螺旋曲線を描くことが可能である。

\textit{推論 3.}\wsp{}点 \( P \)における抵抗力は、中心指向力に対してその比として \( \frac{1}{2}OS \) に比例する。これらの力は直線と同様に比例するためである。

\[
Rr \& TQ \text{ または } \frac{\frac{1}{2}VQ \times PQ}{SQ} \& \frac{PQq.}{SP} \text{ が同時に生成される}
\]
すなわち，
\[
\frac{1}{2}VQ \& PQ, \text{ または } \frac{1}{2}OS \& OP.
\]
 Spiralisの与えられたとき，抵抗する力に対する向心力の比は与えられ，逆に，その比が与えられたとき， Spiralisも与えられる。

\textbf{命題 Corol.\,4.}
物体はこの Spiralis上を円運動を行うことができない場合，抵抗する力が向心力の半分未満である点においてのみ可能である。抵抗する力が向心力の半分と等しいとき， Spiralisは直線 $PS$ と一致し，その直線上で物体は中心に向かって半分の距離を下る。

\[
\text{ここで，物体は中心に近づく途中で半分の距離を下ります。}
\]


物体が抵抗媒体中でない場合、上記で証明したパラボラ軌道（第1巻第10章の理論X）における下降速度と同じ速度で円運動を行うことから、ここでの下降時間はそれらの時間の二倍以上であり、従って与えられるものとなる。

\textit{補則5.} そして、半径から等距離にある点での速度が螺旋 $PQR$ と直線 $SP$ で同じであることに注意し、螺旋 $PQR$ の長さと直線 $PS$ の長さの比は与えられた比であり、つまり $OP$ に対する $OS$ の比であるため、螺旋 $PQR$ 上での下降時間は直線 $SP$ 上での下降時間に対して与えられた比で与えられる。


\textbf{命題 6.} 与えます。

円周 \( S \) を中心として二つの円が描かれるとき、物体がこれらの円の円周間を回ることができる回転数は、\(\frac{PS}{OS}\) の比に相当し、すなわち、 Spiralis と半径 \( PS \) の角度の接線角の正切の比に相当します。また、これらの回転の時間は \(\frac{OP}{OS}\) の比であり、これは中間密度の逆数の比に相当します。

\pngright{295.png}{}{}{}

\textbf{命題 7.} 与えます。

物体が密度が位置から逆比例する媒質内を任意の曲線 \( AEB \) で回転するとき、その時間は物体が中心から離れる距離の逆数の比に相当します。

\textbf{その中心に近づき、}そして円周上の点$AS$が同じ角度で点$B$を通過し、かつその速度が点$A$での初速度と互いに反比例する半径の比率（すなわち$BS$は$AS$と$CS$の中間比）であるとき、その物体は無数の類似した回転$BFC$, $CGD$, \&c. を行い、そして交点で円周上の放射$AS$を連続的に比例する部分$AS$, $BS$, $CS$, $DS$ \&c. に区別する。

しかしながら、これらの回転の時間は円周$AEB$, $BFC$, $CGD$ \&c. の周長と直接比例し、そして速度は反比例である。

\textbf{命題 XVII. 定理 4.}

円運動または螺旋運動中の物体が中心に近づく時間は、その運動の最初の時間と反比例します；すなわち、$AS^{\frac{1}{2}}$, $BS^{\frac{1}{2}}$, $CS^{\frac{1}{2}}$ との関係で逆比します。

また、物体が中心に到達する全体の時間を求めるには、最初の時間 $AS^{\frac{1}{2}}$ と、無限大まで継続的に比例する各段階の和 $AS^{\frac{1}{2}}$, $BS^{\frac{1}{2}}$, $CS^{\frac{1}{2}}$ の比を取ることで得られます。すなわち、最初の項 $AS^{\frac{1}{2}}$ と、二つの最初の項との差 $AS^{\frac{1}{2}} - BS^{\frac{1}{2}}$ の比、そして近似して $\frac{2}{3} AS$ と $AB$ の比となります。

したがって、全体の時間を容易に求めることができます。

\textbf{命題 XVII. 補遺 8.}

これらの事実から、さらに次のように結論を引き出すことが可能です。密度が一定であるか、またはその他の任意の法則に従う抵抗媒体中の物体の運動についても同様です。中心 \( S \) から半径が連続的に比例する \( SA, SB, SC \&c. \) の円周上を描く場合、これらの円周間で回転数の比は、当該媒体の密度がこれらの円周間で平均的なものであるとき、その他の媒体における同様な回転数の比に近似されます。

\[
\textbf{訳文}
\]


\[
\textbf{証明}
\]

中心 \( S \) から半径 \( SA, SB, SC \&c. \) の円周上を描く場合、これらの円周間で物体が回転する回数の比は、当該媒体の密度がこれらの円周間で平均的なものであるとき、その他の媒体における同様な回転数の比に近似されます。

\[
\textbf{補足}
\]

この補遺は、抵抗媒体中の物体の運動についての更なる結論を導き出すためのものです。具体的には、密度が一定または特定の法則に従う媒体での円運動において、これらの円周間で回転数の比は、当該媒体の密度がこれらの円周間で平均的なものであるとき、その他の媒体における同様な回転数の比に近似されるという事実を示しています。




$\text{しかし、同一の比率において Spiralis 前定されたが、その中心点である中間所で切る半径 } AS \text{ と Spiralis 新たなるが、提案された中間所で半径を同じ角度で切る。}$

$\text{そしてまた、これらの角度切断は等しい角であるから、これらが等しい2つの円の間のすべての回転時間をも等しくする。}$

$\text{もしこれらが二つの円間で普遍的に成立すれば、運動はすべての円を継続する。}$

$\text{そしてこの方法で、物体が任意の規則的な中間所でどのようにかつどの時間に回転すべきかを想像することは難しくない。}$

\textbf{補則9.} そして、円錐曲線に近い螺旋運動においても、物体の運動は進行しますが、螺旋線の各回転を等しい間隔で配置し、それらが螺旋線を描く中心に向かって等しい速度で進むと理解できるので、物体がこのような螺旋線上をどのように運動するかを理解することができます。

\textbf{命題XVI. 理論XII.}

\textbf{もし}抵抗媒体の密度が各点から中心までの距離に反比例し、かつ向心力も同様に距離に反比例するとすれば、


\textbf{ dignitas illam ducta: } ここでは、物体が円錐螺旋線 \( \) 上で回転する場合について述べます。すなわち、中心から各放射線を引いた点 \( S \) において、すべての放射線が一定の角度で交差することを主張します。

\textbf{Demonstratur eadem methodo cum Propositione superiore. } 上記は前節の方法と同様に証明できます。もし中心に向かう力 \( F \) が点 \( P \) において各放射線 \( SP \) の距離のべき乗 \( SP^{n+1} \) に比例するものであるとするとき、上記の議論から、物体が任意の弧 \( PQ \) を描くのにかかる時間は \( PQ \times SP^{\frac{1}{2}n} \) と比例することが分かります。

\textbf{nam si vis centripeta in } \( P \) \textbf{ sit reciproce ut distantiis } \( SP \) \textbf{ dignitas quae libet } \( SP^{n+1} \) \textbf{ cujus index est } \( n \) \textbf{ + 1; colligetur ut supra, quod tempus quo corpus describit arcum quemvis } \( PQ \) \textbf{ erit ut } \( PQ \times SP^{\frac{1}{2}n} \).

\textbf{sive ut } $Rr / (PQq. \times SP^n)$ \textbf{ sive ut } $\frac{\frac{1}{2}nVQ}{PQ \times SP^n \times SQ}$, adeoque

物体が点 \( P \) において受ける抵抗は、$ Rr / (PQq. \times SP^n) $ と比例し、これはさらに $\frac{\frac{1}{2}nVQ}{PQ \times SP^n \times SQ}$ と比例します。

\[
ut \frac{\frac{1}{2}OS}{OP \times SP^{n + 1}}.
\]
よって、$P$ 点における密度は $SP^n$ に反比例する。

\textit{補遺。}

他にもこの命題と前節で述べた命題は、不均一な抵抗媒体中の物体の運動について述べており、これらは極めて小さな物体の運動を考慮することを意味します。すなわち、一方の媒体の一部よりも別の部分の方がより密度が高い場合、その差は無視されるべきです。また、他の条件が等しいとき、抵抗は密度に比例すると仮定します。したがって、抵抗力のない媒体では、密度はその程度まで増加または減少し、抵抗力の余剰または不足を除去する必要があります。

suppleatur.

\textbf{命題 XVII. 問題 V.}

\textit{与えられた法則で螺旋上回る物体の向心力と抵抗媒体の力を求める。} 見 \textit{図 Prop.\ XV.}

 Spiral $PQR$ とする。物体が弧 $PQ$ を最短時間で通過する速度、および点 $TQ$ の高さから、向心力と時間の平方に比例する力を得る。次に、面積 $\triangle PSQ$ と $\triangle QSR$ の差 $RSr$ から物体の減速を求めて、その減速から抵抗媒体の力を求めよ。

\textbf{証明}

まず、螺旋 $PQR$ 上で物体が弧 $PQ$ を最短時間で通過する速度を求めると、この速度は向心力と密接に関連している。また、点 $TQ$ の高さ $TQ$ から、向心力と時間の平方に比例する力を得ることができる。

次に、面積 $\triangle PSQ$ と $\triangle QSR$ の差 $RSr$ を考慮し、物体の減速を求める。この面積の差は、物体が螺旋上移動する間に体が受ける減速を表している。したがって、この差から物体の減速を求めることができる。

最後に、得られた減速から抵抗媒体の力を求め、その密度も同様に計算することができる。

\textbf{証明終了}

\textbf{命題 XVIII. 問題 VI.}

\textit{\textbf{中心力の法則が与えられたとき、各点での物体の密度を求め、その密度によって物体が螺旋軌道を描くための速度 retardation を求めよ。}}

中心力から各点での速度を求め、次にその速度の retardation から媒介体の密度を求めます。これは前節で述べた通りです。

しかしながら、このような議論を行う方法については、本命題の第十節と補題第二節において開示しました。また、このような複雑な探討を読者に長時間留まらせるつもりはありません。物体が進行するための力に関する追加事項はすでにあります。

\textbf{次に、中間媒質の密度と抵抗について}。これまで展開した運動及びその類似運動は、これらの媒質においても同様に行われます。


\section{第2編 第5節：流体の密度と圧力、および静水力学（命題XIX〜XXIII）}


% Section: 第2編 第5節：流体の密度と圧力、および静水力学（命題XIX〜XXIII）



\sectpage{V.}

\begin{center}\textbf{流体の密度と圧力、および静水力学}\end{center}

\condpagelarge{流体の定義。}

流体とは、いかなる部分も任意の力を受けて容易に移動することができ、かつ移動すると容易に他の部分との間に隙間が生じる物質である。

\propnopage{命題XIX. 理論XIII.}

\textbf{均一で静止している流体は、どの容器でも静止している状態で閉じ込められ、四面から圧迫されるとき、その部分すべて（凝縮、重力および全ての向心力を考慮に入れない限り）は等しく圧迫され、その圧迫によって生じる動揺は全くない。}

永久的にその位置に留まるものである。

\pngright{299.png}{}{}{}
%Illustration

\textbf{1.} 在球形容器 \(ABC\) 中，流体従四面均匀地被圧縮：我説ないどんな一部分会因これ种圧力而移動。なぜならもし某部分 \(D\) 移動了，那所有与中心等距离的类似部分必須同時に以相同的方式移動；これ是なぜなら所有的圧力是相似且相等的，除了由該圧力引起的运動外，どんな其他运動都被假定為不存在。しかし，これ些部分不可能都移動。

\textbf{証明：} 假設某部分 \(D\) 移動了。〜により流体従四面均匀圧縮，所有与中心等距离的部分必須以相同的方式移動。したがって，もし \(D\) 向某個方向移動，則与其等距离的其他部分也必須向相同的方向移動。しかしながら，これ种假設会導致矛盾，なぜならもし所有これ些部分都移動，那それら之間的相対位置将発生変化，これ違反了流体均匀圧縮的条件。

\textbf{2.} 由此可知，在静水中，深度相同的各点具有相等的圧力。なぜならもし某一点受到比其他点更大的圧力，則該点的一部分会向圧力より小さな方向移動，従而破坏均匀性。したがって，所有与中心等距离的部分必須以相同的方式移動，これ表明在静水中，深度相同的各点具有相等的圧力。

\begin{proof}
静水中において、ある一点 \( P \) の圧力が他の一点 \( Q \) の圧力より大きいと仮定すると、上記の論証により、中心から等距離にある点は同様な移動を行うこととなる。これにより、点 \( P \) および点 \( Q \) 間の相対的位置が変化し、結果として均一性が破綻する。したがって、同じ深さにあるすべての点には等しい圧力が作用しなければならない。
\end{proof}

\textbf{3.} 在静水中，各点的圧力与其深度成正比。なぜならもし某点的深度比較深，則該点受到的圧力より大き；反之亦然。これできる〜を通して考慮流体柱的重量来証明：流体柱越長（即深度越大），其重量越大，従而対底部的圧力也越大。

\textbf{証明：} 考慮一つの高度為 \(h\) 的流体柱，其底面積為 \(A\)。該流体柱的体積為 \(Ah\)，質量為 \(\rho Ah\)（其中 \(\rho\) 為流体密度）。したがって，該流体柱対底部的圧力為 \(\rho Ahg\)（\(g\) 為重力加速度），即圧力与深度成正比。

\textbf{4.} 在静水中，各点的圧强与其深度成正比。なぜなら圧强定義為単位面積上的圧力。〜によれば上述論証，各点的圧力与其深度成正比，したがって各点的圧强也与其深度成正比。

\textbf{証明：} 圧强 \(P\) 定義為 \(\frac{\text{圧力}}{\text{面積}} = \frac{\rho Ahg}{A} = \rho hg\)。したがって，圧强与深度成正比。

近接中心に移動するためには、流体が中心に濃縮される必要がある；これは仮説に反します。また、遠離中心に移動するためには、流体が円周に濃縮される必要があります；これも仮説に反します。さらに、中心から等しい距離を保ちながら同じ時間内に異なる方向へ移動することは不可能です。したがって、流体の部分はどこにも動かない。$\blacksquare$

\textit{例2.} これらの流体の各部分は等しく球対称に圧迫されると言える。


\textbf{そして：} EF は球体の一部であり、もしこの流体が均一に圧迫されないなら、圧力は小さくなる。そしてその流体が均一に圧迫されるまで、圧力を増加させると、最初の場合において各部分はその位置を維持する。しかし圧力が増加した後でも、最初の場合において各部分はその位置を維持し、新しい圧力の追加により移動する。これは定義によって流体の性質である。これら二つは矛盾している。従って、球体 EF が均一に圧迫されることを誤った主張として否定された。$\blacksquare$

\textit{場合3.} さらに、異なる部分の球体について述べる。


圧力はいかなるものか。実際、接する球形の部分は接触点において互いに等しく押し合うことになる（法則III）。しかし、第2の場合によれば、液体の部分も等しい力をで押し合います。したがって、互いに接していない任意の2つの球形の部分は、中間にある球面部分が双方を接することができるため、これらも等しい力をで押し合うことになる。\(\blacksquare\)

\textit{案4.} すなわち、液体のすべての部分はどこでも等しく押し合うことをすでに述べた。なぜなら、任意の2つの部分は球形の部分からどの点でも接されることができ、その点においてはこれらの球形の部分も接することができるためである。\(\blacksquare\)

{\ae}質的に圧迫される場合、それは第3運動法則により他の流体から{\ae}質的に圧迫されます。$\blacksquare$

\textit{案5.} したがって、流体の部分$GHI$が流体その他の部分に似たように密閉された容器内にあり、かつ四面から等しく圧迫されるとき、その部分は互いに等しく圧迫され、また静止します。これは明らかです。どの流体$GHI$も四面から等しく圧迫される場合、そのすべての部分は互いに等しく圧迫され、また静止します。$\blacksquare$

\textit{案6.} したがって、もし流体が柔軟な容器で密閉されているとき、$\blacksquare$

\textbf{そして、流体は等しく押し付けられても等しい圧力に屈し、強い圧力には強い圧力を受ける。これは流体の定義による。}

\textit{\textbf{案7.}} 固い容器内において、流体は一方より他方よりも強い圧力を受けた場合でも、等しい圧力に屈するが、瞬時に屈することである。なぜなら固い容器の側面は流体が移動している間に追随しないからである。一方が屈すると、対向側の側面が押し寄せ、その結果、全方向からの圧力は等しくなります。そして、最初に強く押された部分から流体が退くことを試みるとき、容器の抵抗により対向側の側面が制限されると、圧力は減少します。




\section*{第2編 第5節：流体の密度と圧力、および静水力学（命題XIX〜XXIII）}


\textbf{命題 XVII.} 流体中の各部分は時間の瞬間において等しい密度を有し、かつ局部的な運動なしに等しく圧迫される。そしてさらに、流体の部分は第五の場合において互いに均等な圧力を加え合い、また相互に静止する。

\textbf{証明} これらが真であることは明らかである。\QED

\textbf{推論} よって、流体中の各部分間の運動は、どこからか流体全体を外部表面に押し付けた圧力によって変更されることがない。これは、表面の形状がどこかで変化するか、またはすべての流体の部分が互いにより強く或いはより緩く押しつけ合うことで困難に或いは容易に移動するときを除き。

\section*{命題 XX. 定理 XIV}

\textbf{もし球形の流体が中心から異なる距離にある均質なものであるとすれば、}




流体が中心に重なる部分を覆う円錐面の各部分は全体に対して重力で支持される。円錐の底面の質量は、その底面と同形状かつ同等の深さを持つ流体の層によって支持される。

\pngright{301.png}{}{}{}

次に、$DHM$を底面とする平面と、$AEI$を上表面とする平面を考える。円錐形の流体を無数の同心円錐面 $BFK$, $CGL$ で区分し、各円錐面上部の表面での重力作用を考慮する。そして、各円錐面上部の表面での作用は等しいとする。





\pngright{301.png}{}{}{}




\textbf{すべてのもの。} よって、最高層の表面 $AE$ は単純な自己重力による力で押さえつけられます。そして、最高層全体と次層の表面 $BFK$（命題XIXにより）も同様に押さえつけられます。さらに、次層の表面 $BFK$ も自己重力による力で押さえつけられ、この力を前の力を加えると二倍の圧力が生じます。

この二倍の圧力 \& さらに自己重力による力、すなわち三倍の圧力により、第三層の表面 $CGL$ は押し付けられます。同様に、四倍の圧力と自己重力による力、つまり五倍の圧力により、第五層の表面も押し付けられます。このようにして、各々の表面にはそれぞれの圧力によって押し付けられます。

流体の密度と圧力、および静水力学に関する命題XIXからXXIIIまでを証明します。

\textbf{命題XIX. 定理19}

流体に重ねられた固体の量は、流体が及ぶ球の数に比例し、その頂点までの距離とは関係ありません。したがって、最低限の球の重さを増加させ、密度を無限大まで小さくすると、固体の最後の比は有限の円柱に対等な比になります。つまり、流体の最低限の球の重さと有限の円柱との比は、球の数が増加し、中心からの距離が無限に減少するとき、流体の重力作用の効果が表面から表面まで継続的に保たれるため、等しい比になります。したがって、有限の円柱の底面はその重量を支えます。

\textbf{命題XX. 定理20}

同様な議論により、流体の重力が中心からの距離に比例するとき、固体の最後の比は有限の円柱に対等な比になります。つまり、球の数が無限大になり、中心からの距離が無限に減少すると、流体の重力作用の効果が表面から表面まで継続的に保たれるため、等しい比になります。

\textbf{命題XXI. 定理21}

同様な議論により、流体の重力が中心からの距離に反比例するとき、固体の最後の比は有限の円柱に対等な比になります。つまり、球の数が無限大になり、中心からの距離が無限に減少すると、流体の重力作用の効果が表面から表面まで継続的に保たれるため、等しい比になります。

\textbf{命題XXII. 定理22}

同様な議論により、流体の重力が中心からの距離に任意の比例関数で減少するとき、固体の最後の比は有限の円柱に対等な比になります。つまり、球の数が無限大になり、中心からの距離が無限に減少すると、流体の重力作用の効果が表面から表面まで継続的に保たれるため、等しい比になります。

\textbf{命題XXIII. 定理23}


\textbf{証明}

以上の命題は、流体の密度と圧力、および静水力学に関する重要な原理を示しています。これらの命題は、流体の重力が中心からの距離にどのように依存するかによって、固体の最後の比が有限の円柱に対等な比になることを証明します。

\textbf{終}

\[
\textbf{ut \& ubi Fluidum sursum rarius est, deorsum densius.} \qquad \text{\QED}
\]

\textbf{命題1.}\wsp{}したがって、液体の上面は稀薄であり、下面は密度が高い。すなわち、液体の底板は全体の液体に覆われている質量によって直接押し圧されず、その一部のみを支持し、それが命題で説明されている通りである。残りの質量は液体の弓形形状により支えられている。

\textbf{命題2.}\wsp{}また、中心から等距離にある平面において圧力の量が常に同じであれば、その平面が水平面に平行であれ垂直であれ斜めであっても、または液体がその表面を押し圧して上面へ垂直に延びる場合でも、圧力はその平面全体を支えている。

\section*{第2編 第5節：流体の密度と圧力、および静水力学（命題XIX〜XXIII）}

\textbf{命題XIX. 定理.} 直線または斜めに曲折した穴や管を通って流れる液体，これらの穴や管は規則的または極めて不規則で，広くまたは非常に狭い．これらの状況において圧力は変化しないと結論付けることができ，この定理の証明を個々の流体の場合に適用すればよい．

\textbf{\textit{補則3.}} 同様な証明により（命題XIXによる），同種の密度の液体の部分は，重力圧力によって作用する力を除き，相互に動かないことが確認される．ただし，その動機が凝縮から生じる動機を排除した場合である．

\textbf{\textit{補則4.}} そして従って，同一密度の別の物体も同様である．

流体に密度や圧力の影響を受けない物体は、その物体がこの流体中に没せば、圧力によって得られる運動を全く得ない。したがって、それは下降しないし、上昇しないし、形状も変化させられない。もし球形であるなら、圧力にかかわらず球形のままであり続ける。方形であっても同様である。無論、その物体が柔らかいものでも、流体が自由に浸透するものでもなく、あるいは底部に接しているとしても同じである。なぜなら、流体内の任意の部分は没した物体とその形状、質量、比重の比を保持し、これらすべてが等しいからである。もし没した物体がその形状を保つならば、

\[
\textbf{式}\ref{wasp294}
\]

流体の各部分との比もまた同じであり、これは物体の質量と比重の特定された関係によるものである。

流体に液体が変化し、その形状を液体の形に与える。これについては、以前は上昇したり下降したり、または圧力により新しい形状を与えられたとしても、現在でも上昇したり下降したり、または新しい形状を与えることを強制される。これはなぜなら其重力及其他运動原因保持不変。しかし，〜によれば第5命題的推論5，它現在将静止并保持原有形状。

\textbf{推論5.} それゆえ、特定の密度が流体より大きい物体はその流体に接触して沈み、特定の密度が流体より小さい物体は上昇し、運動および形状変化を伴う。

\textbf{第2編 第5節：流体の密度と圧力、および静水力学（命題XIX〜XXIII）}

\textbf{【第16節】}

その重力の余剰または不足がいかに作用するかを考察する。実際、その重力の余剰または不足は、流体の部分と均衡している物体が受ける衝撃の比に他ならない。そしてこれは、ある棒の一方または他方に置かれた物体の重量の余剰または不足と比較されることが可能である。

\textit{命題6.}

したがって、流体中に構成された物体の重力は二つ存在する：一つは真実で絶対的な重力であり、もう一つは表面的で一般的な重力である。絶対的な重力は物体が下に傾きたいとする全力を表す。一方、表面的で一般的な重力は、物体がより強く流体の部分に対して下に向かう余剰の重力を表す。

流体や他の物体は、包囲する流体に対して下方に作用します。第一種の重力により、流体の各部分およびすべての物体はその位置において重力に影響を受けます。したがって、結合された質量は全体の質量を形成します。

全物体は重いものであり、これについては満ちた液体容器で実験することも可能です；また、全体の質量は各部分の質量と等しく、従ってこれらから構成されます。第二種の重力により、物体はその位置において重力を受けることはありません。つまり、互いに配置された物体は相互に重くならないが、下降する傾向を阻むために互いに妨げ合い、その位置で静止します。これは空気中の物体についても同様です。

\[
\textbf{証明：}
\]

流体の各部分は包囲する流体に対して下方に作用し、物体も同様です。第一種の重力により、これらの質量はその位置において重力に影響を受けます。したがって、結合された質量は全体の質量を形成します。


\[
\textbf{命題XIX}
\]


全物体は重いものであり、これについては満ちた液体容器で実験することも可能です；また、全体の質量は各部分の質量と等しく、従ってこれらから構成されます。第二種の重力により、物体はその位置において重力を受けることはありません。つまり、互いに配置された物体は相互に重くならないが、下降する傾向を阻むために互いに妨げ合い、その位置で静止します。これは空気中の物体についても同様です。

\[
\textbf{命題XX}
\]

流体の各部分は包囲する流体に対して下方に作用し、物体も同様です。第一種の重力により、これらの質量はその位置において重力に影響を受けます。したがって、結合された質量は全体の質量を形成します。


\[
\textbf{命題XXI}
\]


全物体は重いものであり、これについては満ちた液体容器で実験することも可能です；また、全体の質量は各部分の質量と等しく、従ってこれらから構成されます。第二種の重力により、物体はその位置において重力を受けることはありません。つまり、互いに配置された物体は相互に重くならないが、下降する傾向を阻むために互いに妨げ合い、その位置で静止します。これは空気中の物体についても同様です。

\[
\textbf{命題XXII}
\]

流体の各部分は包囲する流体に対して下方に作用し、物体も同様です。第一種の重力により、これらの質量はその位置において重力に影響を受けます。したがって、結合された質量は全体の質量を形成します。


\[
\textbf{命題XXIII}
\]


全物体は重いものであり、これについては満ちた液体容器で実験することも可能です；また、全体の質量は各部分の質量と等しく、従ってこれらから構成されます。第二種の重力により、物体はその位置において重力を受けることはありません。つまり、互いに配置された物体は相互に重くならないが、下降する傾向を阻むために互いに妨げ合い、その位置で静止します。これは空気中の物体についても同様です。





\textbf{命題 XIX.} 一般に、大衆は重い物を正しく判断しない。しかし、大衆が重いと判断する物は、空気の重量によって支えられないものである。

大衆の重量とは真実の重量より多いものを指すもので、それらは空気の重量を超える差異によるものである。従って、空気を圧迫しているときには軽いとみなされる物も存在し、上部に移動する。比較的軽いとは真の意味での軽さではなく、真空へと下降することを示すものである。

同様に、水において、より大きな重量を持つ物体は沈み、小さな重量を持つ物体は浮き上がるとするとき、これらの物体は比較的重く、または比較的軽く、その重さや軽さが比較的であり、真の重さとは異なる。

\textbf{命題 XX.} 任意の液体において、液体の各部分が他の部分よりも深くあるとき、それらの部分はより大きな圧力を受け、浅い部分は小さな圧力を受ける。

この命題を証明するために、まず液体中の任意の点を選び、その点から垂直に下った線分と水平面との交点を定める。この点が液体の表面よりも深いとき、その点より深いすべての点は、より大きな圧力を受けていることが示される。また、より浅い点は小さな圧力を受ける。

\textbf{命題 XXI.} 任意の液体において、液体の各部分が他の部分よりも深くあるとき、それらの部分はより大きな重量を持ち、浅い部分は小さな重量を持つ。

この命題を証明するために、まず液体中の任意の点を選び、その点から垂直に下った線分と水平面との交点を定める。この点が液体の表面よりも深いとき、その点より深いすべての点は、より大きな重量を持ち、より浅い点は小さな重量を持つことが示される。

\textbf{命題 XXII.} 任意の液体において、液体の各部分が他の部分よりも深くあるとき、それらの部分はより大きな面積を占め、浅い部分は小さな面積を占める。

この命題を証明するために、まず液体中の任意の点を選び、その点から垂直に下った線分と水平面との交点を定める。この点が液体の表面よりも深いとき、その点より深いすべての点は、より大きな面積を占め、より浅い点は小さな面積を占めることが示される。

\textbf{命題 XXIII.} 任意の液体において、液体の各部分が他の部分よりも深くあるとき、それらの部分はより長い時間にわたって圧力を受ける、浅い部分は短い時間にわたって圧力を受ける。

この命題を証明するために、まず液体中の任意の点を選び、その点から垂直に下った線分と水平面との交点を定める。この点が液体の表面よりも深いとき、その点より深いすべての点は、より長い時間にわたって圧力を受けることが示される。また、より浅い点は短い時間にわたって圧力を受ける。

これらの命題は、流体の密度と圧力、および静水力学に関する重要な原理を述べています。

\& その質量が水の質量に勝ち越すか、あるいは負けるかする差異が明らかである。しかしながら、物体は重くなることなく下降せず、また重くなることのある物体も、それによって上昇することもない。ただし、実際的な質量により全体の質量が増加しても、比較的見れば一般には水に浮かばないままである。これらの状況についての証明は似ている。

\textbf{補則7.} その重さが示された物体は、他のどの中心指向力においても同様の性質を有する。

\textbf{補則8.} よって、ある媒質において物体が動いている場合、

この部分では、ニュートンが流体の密度と圧力に関する証明について述べています。物体の重さが水の質量に勝ち越すか負けるかの差異、物体が重くなることなく下降しないことや重くなることのある物体も上昇しないこと、そして全体の質量が増えても比較的見れば水に浮かばないままであるという状況について説明しています。さらに、補則7と8では、これらの性質が他の中心指向力においても同様であることを述べています。

\textbf{命題XIXの補則9．} また、流体を圧迫して閉じ込められた物体がその外観形状を変えないとき、Corollaria Prop. XIXによれば、内部部分の相対的位置は変化しないことが明らかとなる。したがって、動物が浸かるとし、すべての感覚が運動から生じるという仮定をすれば、

\textbf{命題XX．} 動物が流体中に浸かっているとき、各部分の密度はその部分が流体に与える圧力に対して反比例する。

\textbf{証明．} ある物体が流体中に完全に浸かっており、その内部の各部分が等しい面積速度で移動していると仮定する。この場合、各部分は流体から受ける圧力によって均衡しており、各部分の密度はその部分が流体に与える圧力に対して反比例する。

\textbf{命題XXI．} 流体中の物体の各部分の密度は、その部分が流体に与える圧力と面積速度との積に正比する。

\textbf{証明．} ある物体が流体中に完全に浸かっており、その内部の各部分が等しい面積速度で移動していると仮定する。この場合、各部分は流体から受ける圧力によって均衡しており、各部分の密度はその部分が流体に与える圧力と面積速度との積に正比する。

\textbf{命題XXII．} 流体中の物体の各部分の密度は、その部分が流体に与える圧力と面積速度との積に反比例する。

\textbf{証明．} ある物体が流体中に完全に浸かっており、その内部の各部分が等しい面積速度で移動していると仮定する。この場合、各部分は流体から受ける圧力によって均衡しており、各部分の密度はその部分が流体に与える圧力と面積速度との積に反比例する。

\textbf{命題XXIII．} 流体中の物体の各部分の密度は、その部分が流体に与える圧力と面積速度との積に正比し、またその部分の体積と面積速度との積に反比例する。

\textbf{証明．} ある物体が流体中に完全に浸かっており、その内部の各部分が等しい面積速度で移動していると仮定する。この場合、各部分は流体から受ける圧力によって均衡しており、各部分の密度はその部分が流体に与える圧力と面積速度との積に正比し、またその部分の体積と面積速度との積に反比例する。

\textbf{命題XXI. 定理XV.} 流体中に浸かっている部分から生じる圧力は、物体が被圧縮されない限り、物体に感覚をもたらすことはありません。これは、流体によって圧迫されるシステムの各部分間の比率が等しいことを意味します。システムのすべての部分は、流体がない状態で構成されているかのように、同じ運動を受けるとともに、比較的自重だけを保持し続けます。ただし、流体がそれらの運動に対して多少抵抗する場合や、圧縮のために結合される必要がある場合を除き。

\begin{center}
\textbf{Prop.\ XXI. Theor.\ XV.}
\end{center}

流体の密度が圧縮に比例し、かつその部分が中心から距離によって逆比例的に下方に引き寄せられるとき、これらの距離が連続的な比を成すならば、流体の密度も同様に連続的な比を成ることを主張する。

球形底面 \(ATV\) において流体が接しているとし、中心を \(S\) とする。また、\(SA\), \(SB\), \(SC\), \(SD\), \(SE\) 等は連続的な比を成す距離とする。直線 \(AH\), \(BI\), \(CK\), \(DL\), \(EM\) をそれぞれ点 \(A\), \(B\), \(C\), \(D\), \(E\) における流体の密度と比例するように垂直に引き上げる。

\[
\textbf{証明}
\]

各点 \(A\), \(B\), \(C\), \(D\), \(E\) における流体の密度をそれぞれ \(\rho_A\), \(\rho_B\), \(\rho_C\), \(\rho_D\), \(\rho_E\) とし、これらは連続的な比を成すとする。つまり、

\[
\rho_A : \rho_B : \rho_C : \rho_D : \rho_E = k_1 : k_2 : k_3 : k_4 : k_5
\]

ここで \(k_1, k_2, k_3, k_4, k_5\) は連続的な比を成す定数である。

各点における密度が連続的な比を成ることから、流体の密度と距離との関係は連続的な比を成すことが示される。したがって、流体の密度も連続的な比を成すことになる。

\[
\blacksquare
\]

\textbf{第2編 第5節：流体の密度と圧力、および静水力学}

\medskip

命題 XIX. 問題.

同様に、各点において流体が重力を受ける量は、その点から流体の表面までの距離に対する比例である。すなわち、点 \( A \)、\( B \)、\( C \)、および \( D \) における重力は、\(\frac{AH}{AS}\)、\(\frac{BI}{BS}\)、\(\frac{CK}{CS}\) と同様に、または、これらが等しく \(\frac{AH}{AB}\)、\(\frac{BI}{BC}\)、\(\frac{CK}{CD}\) であるとする。


\par\begin{center}
    \centering
    \includegraphics[width=639pt,height=1084pt,keepaspectratio]{304.png}

\end{center}\par

次に、これらの重力を点 \( A \) から点 \( B \)，点 \( B \) から点 \( C \)，点 \( C \) から点 \( D \) と等級的に減少させた点 \( B \)，\( C \)，\( D \) を通じて連続的に引き延ばす。そして、これらの重力が高さ \( AB \)，\( BC \)，\( CD \) に達したときの圧力を \( AH \)，\( BI \)，\( CK \) とし、これらは平面 \( ATV \)（定理 XIV. 根拠による）に対する圧力である。

\medskip

証明.

各点における重力が等級的に減少するという仮定により、各段階の重力は直線 \( AH \)，\( BI \)，および \( CK \) に沿って連続的に引き延ばされる。これらの直線は平面 \( ATV \) 上で一定の圧力を生じるため、それぞれが高さ \( AB \)，\( BC \)，および \( CD \) まで達するときの圧力 \( AH \)，\( BI \)，および \( CK \) は、点 \( A \)，\( B \)，および \( C \) の重力に対する比例である。

流体の粒子は無限に及んで圧力を受けるため、粒子 \( A \) はすべての \( AH, BI, CK, DL \) の圧力を受け、粒子 \( B \) は最初の \( AH \) を除いてすべての圧力を受ける。粒子 \( C \) は最初の二つの \( AH, BI \) を除いてすべての圧力を受ける。これにより、粒子 \( A \) の密度 \( AH \) は粒子 \( B \) の密度 \( BI \) に比して、すべての \( AH + BI + CK + DL \) の総和（無限大まで）に対して、すべての \( BI + CK + DL, \&c. \) の総和に対して比が成り立つ。また、粒子 \( B \) の密度 \( BI \) は粒子 \( C \) の密度 \( CK \) に比して、すべての \( BI + CK + DL, \&c. \) の総和に対して、すべての \( CK + DL, \&c. \) の総和に対して比が成り立つ。

\textbf{よって、}これらの密度\textit{$AH$、$BI$、$CK$ など}はその差分に比例し、さらにこの Lem.\ I の定理により連続的に比例します。したがって、密度\textit{$AH$、$BI$、$CK$ など}の総和が連続的に比例することになります。よって、位置\textit{A、B、C}での密度が\textit{$AH$、$BI$、$CK$ など}であるなら、これらも連続的に比例します。

また、一気に移動する代わりに、距離\textit{$SA$、$SC$、$SE$}が連続的に比例するとき、密度\textit{$AH$、$CK$、$EM$}も連続的に比例します。同様の議論により、任意の距離での密度も連続的に比例します。

\[
SA, SD, SQ
\]
の比が連続比例であるとき、流体の密度 \( AH, DL, QT \) も連続比例となります。点 \( A, B, C, D, E, \ldots \) が進むと、底 \( A \) から流体の頂部まで連続的な深さを表すとき、任意の距離 \( SA, SD, SQ \) における密度 \( AH, DL, QT \) も常に連続比例で維持されます。証明終了。


\QED



\par\begin{center}
    \centering
    \includegraphics[width=1183px,height=1248px]{305.png}

\end{center}\par

\textbf{補遺}\wsp{}したがって、流体の密度が任意の二地点 \( A \) で与えられたとします。

\textbf{第27節}

点 \( E \) の密度を他の任意の点 \( Q \) における密度に結びつけることができるとします。中心 \( S \) を取って、直線 \( SQ \), \( SX \) で双曲線を描き、この双曲線は直線 \( AH \), \( EM \), \( QT \) と交わり、それぞれ点 \( a \), \( e \), \( q \) を定めます。また、直線 \( HX \), \( MY \), \( TZ \) は直線 \( SX \) の漸近線に垂直で、それぞれ点 \( h \), \( m \), \( t \) を定めます。領域 \( ZYmtZ \) が与えられた領域 \( YmhX \) と同一比であるとき、領域 \( EeqQ \) が与えられた領域 \( EeaA \) と同一比であるとする。すると直線 \( Zt \) の延長は直線 \( QT \) の密度に比例する点で切断します。

なぜなら、もし直線 \( SA \), \( SE \), \( SQ \) が連続的に比例している場合、領域 \( EeqQ \), \( EeaA \) は等しく、したがってこれらの領域は比例します。


流体の密度と圧力、および静水力学（命題XIX〜XXIII）【第28段落】

$YmtZ$, $XhmY$ も同様に等しい \& 直線 $SX$, $SY$, $SZ$ は $AH$, $EM$, $QT$ であり、これらも連続比に属するものでなければなりません。そして直線 $SA$, $SE$, $SQ$ が連続比に属する他の任意の順序を有する場合でも、直線 $AH$, $EM$, $QT$ は連続比の面積双曲線により同じ順序を得ます。

\textbf{命題XXII.} \textit{理論XVI.}

流体の密度が圧縮に比例し、その部分が重心からの距離の平方から逆比例するとき、


\section*{命題XXII. 理論XVI.}

ある流体の密度が圧縮に比例し、その部分が重心からの距離の平方から逆比例するとき、



\textbf{命題 XIX}：異なる点から流体の密度が等比数列で変化するとき、それら点までの距離は逆に等比数列で減少します。すなわち、もし異なる点の距離が音楽的比例に選ばれるなら、その点での流体の密度も同様に等比数列になります。


\par\begin{center}
    \centering
    \includegraphics[width=18cm,height=13.64cm,keepaspectratio]{306.png}

\end{center}\par

中心を \( S \) とし、\( SA \), \( SB \), \( SC \), \( SD \), \( SE \) を等比数列で選んだ距離とする。これらの距離に対して、直線 \( AH \), \( BI \), \( CK \), \ldots\ は垂直に引かれるものとする。これらは流体の密度と同一の比率を有し、それぞれの点 \( A \), \( B \), \( C \) における流体の密度は \( SA \), \( SB \), \( SC \) の逆数に比例します。また、各点での重力比は次のようになります：

\[
\frac{AH}{SA^q}, \quad \frac{BI}{SB^q}, \quad \frac{CK}{SC^q}, \quad \ldots
\]

ここで \( q \) は適切な指数です。

\textbf{これらの物質は、等しく均一に重力を持ち、}点$A$から$B$への初めの引力、点$B$から$C$への二番目の引力、点$C$から$D$への三番目の引力、そして類推します。これらの線分$AB$, $BC$, $CD$, $DE$などは、またはこれらと等価な高さ$AB$, $BC$, $CD$, $DE$などは、それぞれの高さに比例する距離$SA$, $SB$, $SC$, などに相当します。したがって、これらの圧力$\frac{AH}{SA}$, $\frac{BI}{SB}$, $\frac{CK}{SC}$, そして類推する指数により、これらの線分の和は等しい密度の差$AH - BI$, $BI - CK$, そして類推します。したがって、密度の差$AH - BI$, $BI - CK$, そして類推は、それぞれの指数$\frac{AH}{SA}$, $\frac{BI}{SB}$, $\frac{CK}{SC}$, そして類推する和に比例します。


\par\begin{center}
    \centering
    \includegraphics[width=0.8\textwidth]{307.png}

\end{center}\par

中心 \( S \) から漸近線 \( SA \), \( SX \) 上に双曲線を描くと、この双曲線は垂直線 \( AH \), \( BI \), \( CK \) をそれぞれ点 \( a \), \( b \), \( c \) で交わり、垂線 \( Ht \), \( Iu \), \( Kw \) が点 \( h \), \( i \), \( k \) に達する。また、密度の差分 \( tu \), \( uw \) 等はそれぞれ \(\frac{AH}{SA}\), \(\frac{BI}{SB}\) 等となる。

さらに、矩形 \( tu \times th \), \( uw \times ui \) すなわち \( tp \), \( uq \) も同様に \(\frac{AH \times th}{SA} : \frac{BI \times ui}{SB}\) 等となり、これは \( Aa \), \( Bb \) 等と等しい。

実際、双曲線の性質から \( SA \) に対する \( AH \) や \( St \) の比は \( th \) に対する \( Aa \) の比であり、従って \(\frac{AH \times th}{SA}\) は \( Aa \) に等しい。

\[
\text{同様の議論により、}\frac{BI \times ui}{SB} = Bb, Cc, \&c. です。
\]
\[
\text{ここで、}Aa, Bb, Cc, \&c.\ は連続的に比例していますので、差分
\]
\[
Aa - Bb, Bb - Cc, \&c.
\]
\[
\text{も比例します。したがって、これらの差分の積}
\]
\[
tp, uq, \&c.\ は比例し、また、これらの差分の総和
\]
\[
Aa - Cc \text{または} Aa - Dd
\]
\[
\text{の総和}
\]
\[
tp + uq,
\]
\[
\text{または} tp + uq + wr
\]
\[
\text{も比例します。このような項を最大限に取り、すべての差分の総和}
\]
\[
Aa - Ff
\]
\[
\text{は総和}
\]
\[
zthn
\]
\[
\text{と比例します。項の数を増やすことにより}
\]

\textbf{終点} \(A\), \(B\), \(C\) などと点間の距離を無限に小さくするとき、これらの直角三角形は双曲線弧 \(zthn\) となる等面積を有し、したがって差分 \(Aa - Ff\) はその面積に対応します。任意の距離 \(SA\), \(SD\), \(SF\) を調和級数に従って選んだとすると、差分 \(Aa - Dd\), \(Dd - Ff\) は等しいことになります。したがってこれらの差分に対する面積比 \(thlx\), \(xlnz\) も等しくなります。また、密度 \(St\), \(Sx\), \(Sz\) すなわち \(AH\), \(DL\), \(FN\) は連続する比例関係にあることが分かります。\(\blacksquare\)

\textbf{推論}：したがって、任意の二つの流体の密度 \(AH\) と \(CK\) が与えられた場合、その差分 \(tw\) の面積 \(thkw\) も得られます。そして、任意の高さ \(SF\) における密度 \(FN\) を求めるためには、面積 \(thnz\) を与えられた面積 \(thkw\) に比例させて求めます。すなわち、
\[
FN = \frac{Aa - Ff}{Aa - Cc} \cdot thnz.
\]

\textbf{補釈}：同様の議論により、流体の各部分の重力が中心からの距離の三倍比で減少する場合、または各点の重力が中心からの距離の平方の逆数に比例する（すなわち \(SA, SB, SC, \ldots\) の各点での重力はそれぞれ \(\frac{1}{SA^2}, \frac{1}{SB^2}, \frac{1}{SC^2}\) など）場合でも、その性質が証明できます。

\[
\frac{SA \textbf{ cub.}}{SAq.}, \frac{SA \textbf{ cub.}}{SBq.}, \frac{SA \textbf{ cub.}}{SCq.}
\]
を算術級数としてとるとき、密度 \( AH \), \( BI \), \( CK \), \&c.\ は幾何級数となる。また、質量が距離の四倍比で減少し、距離の立方の逆比例（すなわち
\[
\frac{SA \textbf{ qq.}}{SA \textbf{ cub.}}, \frac{SA \textbf{ qq.}}{SB \textbf{ cub.}}, \frac{SA \textbf{ qq.}}{SC \textbf{ cub.}}
\]
を算術級数としてとるとき、密度 \( AH \), \( BI \), \( CK \), \&c.\ は幾何級数となる。そして無限にまでこの論証を延長することができる。さらに、流体の各部分の質量が全ての距離において

\[
\frac{SA \textbf{ cub.}}{SAq.}, \frac{SA \textbf{ cub.}}{SBq.}, \frac{SA \textbf{ cub.}}{SCq.}
\]
を算術級数としてとるとき、密度 \( AH \), \( BI \), \( CK \), \&c.\ は幾何級数となる。そして無限にまでこの論証を延長することができる。

\[
\text{同様に、距離が等差数列で増加するとき、}
\]
\[
\text{密度は等比数列で減少し、}
\]
\[
\text{これは Vir Cl. \textbf{エドワード・ハレ}氏が発見しました。}
\]
\[
\text{もし重力が距離に比例し、距離の平方が等差数列で増加するとき、}
\]
\[
\text{密度は等比数列で減少します。}
\]
\[
\text{そして無限大までこの関係が続きます。}
\]
\[
\text{これらは流体が圧縮されるとき、}
\]
\[
\text{密度が圧縮力に比例するか、つまり、}
\]
\[
\text{あるいは等価的に、流体が占める空間がこの力の逆数に比例するとき、}
\]
\[
\text{このような関係が成立します。}
\]
\[
\text{他の圧縮法則を仮定することも可能です、たとえば、}
\]
\[
\text{圧縮力を流体密度の平方の立方に比例させるとき、}
\]

\textbf{第2編 第5節：流体の密度と圧力、および静水力学（命題XIX〜XXIII）}


三倍の比において等しい力を四倍の比で表す。この場合、もし重力が中心からの距離の平方に反比例するなら、密度はその距離の立方に反比例します。仮定すると、圧縮力が密度の立方乗に比例し、もし重力が中心からの距離の平方に反比例するなら、密度はその距離の1.5倍の比で反比例します。また、圧縮力が密度の二倍の比で比例し、重力が距離の二倍の比で反比例するなら、密度は距離に反比例します。

\textbf{証明}

\begin{enumerate}
    \item \textit{命題XIX}：三倍の比において等しい力を四倍の比で表す。
    
    ここで、力$F$が等しい力を三倍の比で表し、密度$\rho$がその距離の立方に反比例するとします。つまり，
    \[
    F \propto \frac{1}{\rho^3}
    \]
    ここで、$r$を中心からの距離とすると、
    \[
    F = k \cdot \frac{1}{\rho^3} = k \cdot \left(\frac{1}{r}\right)^3 = \frac{k}{r^3}
    \]
    ここで$k$は比例定数です。したがって、密度$\rho$は距離$r$の立方に反比例します。

    \item \textit{命題XX}：重力が中心からの距離の平方に反比例するとき、密度はその距離の1.5倍の比で反比例します。
    
    ここで、圧縮力$F$が密度$\rho$の立方乗に比例し、重力が中心からの距離$r$の平方に反比例するとします。つまり，
    \[
    F = k_1 \cdot \rho^3, \quad g = k_2 \cdot \frac{1}{r^2}
    \]
    ここで、$g$は重力であり、$k_1$と$k_2$は比例定数です。圧縮力$F$は重力$g$によって平衡されるため、
    \[
    F = g \cdot A, \quad A = k_3 \cdot r^2
    \]
    ここで、$A$は圧迫面積であり、$k_3$は比例定数です。したがって，
    \[
    k_1 \cdot \rho^3 = k_2 \cdot \frac{1}{r^2} \cdot k_3 \cdot r^2
    \]
    ここで、$k_2 \cdot k_3$を新たな比例定数とすると、
    \[
    \rho^3 = C \cdot \frac{1}{r^{1.5}}
    \]
    ここで$C$は新たな比例定数です。したがって、密度$\rho$は距離$r$の1.5倍の比で反比例します。

    \item \textit{命題XXI}：圧縮力が密度の二倍の比で比例し、重力が距離の二倍の比で反比例するとき、密度は距離に反比例します。
    
    ここで、圧縮力$F$が密度$\rho$の二倍の比で比例し、重力$g$が距離$r$の二倍の比で反比例するとします。つまり，
    \[
    F = k_1 \cdot 2\rho, \quad g = k_2 \cdot \frac{1}{2r}
    \]
    ここで、$k_1$と$k_2$は比例定数です。圧縮力$F$は重力$g$によって平衡されるため、
    \[
    F = g \cdot A, \quad A = k_3 \cdot r^2
    \]
    ここで、$A$は圧迫面積であり、$k_3$は比例定数です。したがって，
    \[
    k_1 \cdot 2\rho = k_2 \cdot \frac{1}{2r} \cdot k_3 \cdot r^2
    \]
    ここで、$k_2 \cdot k_3$を新たな比例定数とすると、
    \[
    2\rho = C \cdot \frac{r}{2}
    \]
    ここで$C$は新たな比例定数です。したがって、密度$\rho$は距離$r$に反比例します。
\end{enumerate}

\textbf{終}

すべての粒子が中心距離に反比例して互いに斥力を持つとき、これらは弾性流体を形成します。その密度は圧縮度に比例します。逆もまた然りで、流体が互いに斥力を持つ粒子から構成され、その密度が圧縮度に比例する場合、各粒子の中心斥力は中心距離に反比例します。


\par\begin{center}
    \centering
    \includegraphics[width=1169px,height=783px,trim=-24 0 0 0]{309.png}

\end{center}\par

次に、立方体空間 $ACE$ 内の流体を圧縮するとします。

流体の密度と圧力、および静水力学に関する論証を次のように再構成する。

空間の立方体積が $ace$ 未満に縮小し、両方の空間で占める位置は相似であるとき、距離比は立方体の辺の比 $AB$, $ab$ と等しいです。また、大きい立方体 $ABCD$ の体積 $AB^3$ と小さい立方体 $db$ の体積 $ab^3$ の比は、中間密度は互いに反比例します。すなわち、$AB^3 : ab^3$ です。さらに、大きい立方体の辺 $AB$ を基準として、小さい立方体の辺 $db$ に対応する正方形 $DP$ の面積を等しいものとします。仮定により、正方形 $DP$ が含む流体に作用する圧力は、同じく正方形 $DB$ が含む流体に作用する圧力に対して、中間密度の比 $ab^3 : AB^3$ です。

\[
\textbf{証明：} \quad
\text{大きい立方体の辺 } AB \text{ と小さい立方体の辺 } ab \text{ の比は }
AB : ab = \sqrt[3]{AB^3} : \sqrt[3]{ab^3}
\]

ここで、$DP$ が $db$ に等しい正方形とします。このとき、流体に対する圧力の比は中間密度の比に比例します。

\[
\textbf{従って：} \quad
\frac{\text{正方形 } DP \text{ が含む流体に作用する圧力}}{\text{正方形 } DB \text{ が含む流体に作用する圧力}}
= \frac{ab^3}{AB^3}
\]

以上より、流体の密度と圧力の関係を明確に示すことができます。

流体が包まれているとき、平面 \(DP\) によって流体に加わる圧力は、平面 \(DB\) によって流体に加わる圧力と同じである。これは、\(DP^2\) に対する \(DB^2\) の比と等しいからで、すなわち \(AB^2\) に対する \(ab^2\) の比と等しい。従って、平面 \(DB\) によって流体に加わる圧力は、平面 \(db\) によって流体に加わる圧力を \(AB\) に対する \(ab\) の比で表すことができる。

平面 \(FGH\) および \(fgh\) を立方体の断面として用いて、流体を二つに分けると、これらは互いに等しい力で押さえつけ合う。これは、平面 \(AC\) 及び \(ac\) によって加わる圧力において、それぞれが \(ab\) に対する \(AB\) の比であるからである。したがって、これらの向心力は、これらの圧力を支えるために働く力と等しい。% 

\textbf{同様にして、各粒子が平面} \( FGH \) と \( fgh \) 上で作用する力を考慮すると、これらの力はそれぞれの粒子が単独で作用する力を同一比率に持つ。これは、両方の立方体において粒子の配置が同じであるためです。したがって、大きい立方体内における各粒子が平面 \( FGH \) 上で作用する力を、小さい立方体内における各粒子が平面 \( fgh \) 上で作用する力に対して、\( ab \) に対する \( AB \) の比率と同じになります。すなわち、粒子間の距離の逆数の比率と一致します。\(\blacksquare\)

また、逆に、もし各粒子の力を粒子間の距離の逆数に比例すると仮定する場合、つまり、立方体の辺 \( AB \) に対する \( ab \) の比率に比例する場合、これらの力の総和は同じ比率になります。

\textbf{命題 XIX. 定理 10.}

中心からの力が半径と反比例するとき、各側面の圧力をその半径の力の総和に比例させると、正方形 \(DP\) の圧力は側面 \(DB\) の圧力に対して \(ab \text{quad.}\) に、また \(AB\) の平方に対して比例します。したがって、正方形 \(DP\) の圧力と側面 \(db\) の圧力との比は \(ab \text{cub.}\) と \(AB\) の立方との比に等しくなります。これはつまり、圧縮力の比は密度の比と同じであることを意味します。

\textbf{\textit{補遺.}}

同様の議論により、中心からの力が半径に対して反比例する場合、圧縮力の立方は密度の平方に比例します。もし力が

\begin{equation*}
\textbf{命題 XX. 定理 11.}
\end{equation*}

中心からの力が半径と反比例し、また各側面の圧力をその半径の力の総和に比例させるとき、正方形 \(DP\) の圧力は密度の平方に対して比例します。

\textbf{向心力は，半径と立方または四次比に反比例し，圧縮力は密度の立方または四次方根に比例する。そして一般に，$D$ を半径，$E$ を圧縮された流体の密度とし，向心力が半径との任意のべき $D^n$ と比例する場合，圧縮力は密度の $n+2$ 次方根の三次方根に比例する。逆もまた然りである。これらのすべては，特に近傍にある個々の向心力を表すものであり，またはそれから短時間に広がるものと理解されるべきである。}

\textbf{証明：} まず，半径 $D$ に対する向心力が $D^n$ と比例する場合を考える。円周上の任意の点において，向心力を表す直線はその点から中心に引かれた半径に対して垂直である。したがって，圧縮力はこの直線上で作用し，その量は半径に対する密度の $n+2$ 次方根の三次方根に比例する。

次に，円周上の点から中心への半径を $D$ とし，その点における密度を $E$ とする。向心力が $D^n$ と比例すると仮定した場合，圧縮力は $E^{n+2}$ の三次方根に比例する。

逆の関係も同様に証明できる。つまり，圧縮力が密度の $m$ 次方根の三次方根と比例する場合，向心力は半径の $(m-2)/n$ 乗と比例する。

これらの結果は，特に近傍にある個々の向心力を表すものであり，またはそれから短時間に広がるものと理解されるべきである。

\textbf{例を挙げると磁気体の場合です。}

これらの物質の引力は、その種類に近接する物質に対してほぼ限界されます。磁石の力は鉄板によって縮小され、鉄板の表面でほぼ限界を迎えます。つまり、より遠くにある物体は磁石よりも鉄板から強く引き寄せられるわけではありません。同様に、もし粒子が他の同じ種類の粒子を離れる場合、それほど遠い粒子に対しては力はほとんど働かないか、あるいは僅かに働くかもしれませんが、その効果は中間の粒子によって増幅されたものであるとします。このような粒子から構成される流体について述べたのがこの命題です。


\par\begin{center}
  \centering
  \includegraphics[width=0.8\textwidth]{311.png}

\end{center}\par

流体の密度と圧力、および静水力学に関する議論において、もし個々の粒子の作用が無限遠方まで拡散するとするとき、より大きな量の流体を等しく密にすることには、より大きい力を必要とする。すなわち、一つの粒子が自身からどの場所までの距離に反比例して他のすべての粒子を無限遠方に放つとする場合、流体を相似の容器で等しく圧縮し、密にする力は直径の平方に比例する。したがって、同じ容器内で流体を圧縮する力は、立方根の長さと密度の二乗との比に反比例する。

また、実際には流体が弾性を持つ場合、個々の粒子が互いに放たれるものかどうかという点については、物理学的な問題である。

\textbf{流体の性質について，その構成する粒子から数学的に証明し，哲学家たちにその問題を扱う手がかりを与えました．}

\textit{命題XIX. 定理19．} 流体の密度は，その部分の面積に対する質量比で定義される．

\textbf{設流体の部分が \( A \) の面積を占め，その質量を \( m \) とする．このとき，流体の密度 \( \rho \) は}

\[ \rho = \frac{m}{A} \]

\textit{命題XX. 定理20．} 流体の部分が等しい面積に対して等しい質量を有するとき，その部分は等しい密度を持つ．

\textbf{証明．} 仮定により，流体の部分 \( A \) と \( B \) の質量が等しく，それぞれの面積も等しい．したがって，

\[ \rho_A = \frac{m}{A}, \quad \rho_B = \frac{m}{B} \]

ここで \( A = B \)，よって

\[ \rho_A = \rho_B \]

\textit{命題XXI. 定理21．} 流体の部分が等しい面積に対して異なる質量を有するとき，その部分は異なる密度を持つ．

\textbf{証明．} 仮定により，流体の部分 \( A \) と \( B \) の面積が等しく，それぞれの質量が異なる．したがって，

\[ \rho_A = \frac{m_1}{A}, \quad \rho_B = \frac{m_2}{B} \]

ここで \( m_1 \neq m_2 \)，よって

\[ \rho_A \neq \rho_B \]

\textit{命題XXII. 定理22．} 流体の部分が等しい質量に対して異なる面積を有するとき，その部分は異なる密度を持つ．

\textbf{証明．} 仮定により，流体の部分 \( A \) と \( B \) の質量が等しく，それぞれの面積が異なる．したがって，

\[ \rho_A = \frac{m}{A}, \quad \rho_B = \frac{m}{B} \]

ここで \( A \neq B \)，よって

\[ \rho_A \neq \rho_B \]

\textit{命題XXIII. 定理23．} 流体の部分が等しい質量と面積に対して異なる密度を有するとき，その部分は流体の密度が変化していることを示す．

\textbf{証明．} 仮定により，流体の部分 \( A \) と \( B \) の質量と面積が等しく，それぞれの密度が異なる．したがって，

\[ \rho_A = \frac{m}{A}, \quad \rho_B = \frac{m}{B} \]

ここで \( \rho_A \neq \rho_B \)，よって流体の密度は変化している．


\chapter{第2編 抵抗媒質中の運動（第6節〜第9節：振子・流体運動・波動音速・渦動説批判）}



\section{第2編 第6節：振子の運動と抵抗（命題XXIV〜XXXI）}


% Section: 第2編 第6節：振子の運動と抵抗（命題XXIV〜XXXI）



【第6節：振子の運動と抵抗】

\textbf{命題XXIV. 理論XVIII.}

\textit{自由に吊るされた物体中の質量は、その中心が懸挂点から等距離にある場合、これらの質量は空気中での往復運動の時間の二乗比と重力の積に合成した比で比例する。}

なぜなら、与えられた力が与えられた時間中に与えられた質量を生成できる速度は、力と時間が直接的であり、質量は反比例します。
\begin{center}
\textbf{図1}\\
\includegraphics[width=0.8\linewidth]{312.png}
\end{center}

大きい力や長い時間、または小さい質量があるほど、その速度も大きくなります。

$\text{速度が生じます。これは運動の第2法則から明らかです。}$

$\text{もしこれらの振動体が同じ長さであれば、}$
\[
\textbf{推進力} \, F \propto 1/r^2
\]
$\text{が直角に等距離にある場所で等しくなるようにします。したがって、}$
\[
\text{二体が等しい弧を描くとき、これらの弧が等分割されると、}
\]
\[
\text{各部分を描く時間は全振動の時間を直接比例して持つことから、}
\]
\[
\textbf{速度} \, v
\]
$\text{は対応する振動部分に対して推進力と全振動の時間の量に直接比例します。}$


$\textbf{逆に、したがって}$
質量と力および振動時間は直接的であり、速度は逆に比例します。しかし速度は時間に対して逆に比例し、したがって時間は直接的であり、速度は逆に比例します。そして平方の時間は直接的であり、質量は運動する力をとく時間の平方に比例し、すなわち質量は力と時間の平方に比例します。つまり質量は力と時間の平方に比例します。\QED

\textbf{補則1.}したがって、もし時間が等しい場合、各物体の質量はその質量に比例します。

\textbf{補則2.}もし質量が等しい場合、質量は等しくなります。


\textbf{3.} すなわち時間の平方。

\textbf{3.補則}\wsp{}もし質量が等しいならば、その重さは逆に時間の平方に比例する。

\textbf{4.} すなわち、他の条件が一定であるとき、時間の平方は吊り物の長さに比例し、したがって質量と時間とが等しい場合、重さは吊り物の長さに比例する。

\textbf{5.} そして普遍的に、吊り物の質量はその重さと時間の平方との直接比および吊り物の長さとの逆比である。

\textbf{6.} しかし、中身が抵抗しない場合、


\textbf{第2編 第6節：振子の運動と抵抗（命題XXIV〜XXXI）}

\textbf{命題 XXVIII. 定理 15.} 振動体の運動は、比較的質量と時間の平方に直接比例し、そして懸垂長さには反比例する。

すなわち、比較的質量は、任意の重力場中の物体を移動させる力であり、これについては先に述べた通りです。従って、そのような重力場が抵抗しない場合と、真空における絶対質量は等しい効果をもたらします。

\textbf{命題 XXIX. 定理 16.} 比較的質量の比は、各物体間の質量の比に直接比例する。

\textbf{命題 XXX. 定理 17.} 同じ物体における比較的質量と絶対質量との比は、その位置によって変化する。

\textbf{命題 XXXI. 定理 18.} 物体の重さは、各物体間の質量の比に直接比例し、そして物体の長さにより反比例する。

実際、最も精密な実験により明らかになった結果によれば、

\textit{補則 7.} これらの関係を用いて、物体間の質量の比較と、同一物体における異なる位置での重力の変化を理解することができます。

\pngright{313.png}{}{}{}

以上が第2編 第6節：振子の運動と抵抗（命題XXIV〜XXXI）の翻訳です。

\textbf{命題XXV.} \textbf{理論XIX.}

\textit{\textbf{抵抗のある媒質中で運動する振子と、同一重力密度の媒質中で運動しない振子が、 Cycloid の弧上を等時間に振動し、その弧の部分も比例的に描く。}}


\par\begin{center}
    \centering
    \includegraphics[width=1429px,height=1082px,keepaspectratio]{313.png}

\end{center}\par

設 $AB$ を Cycloid の弧とする。その中で任意の時間に振動する物体 $D$ が描く弧を描写し、この弧は $C$ 点で二等分される。

\textbf{証明.} 物体 $D$ が点 $A$ から点 $B$ へと移動するとき、その速度は常に一定である。したがって、物体 $D$ の運動時間は弧の長さに比例し、また媒質が抵抗しない場合でも同様である。

物体 $D$ が Cycloid 上を振動するとき、点 $C$ を中心とする円周上の任意の点 $P$ から点 $A$ までの線分 $AP$ の長さは、点 $P$ から点 $B$ までの線分 $BP$ の長さと等しい。これは Cycloid の性質によるものである。

したがって、物体 $D$ が Cycloid 上を振動する時間も、媒質が抵抗しない場合でも同じである。また、物体 $D$ が Cycloid 上の任意の点から次の点へ移動するとき、その移動距離は弧の部分と比例する。

したがって、物体 $D$ の振動は Cycloid の弧上を等時間に進行し、その弧の部分も比例的に描く。

$\textbf{C}$がその点の下限とし、$\textbf{C}$における力は任意の位置$\textbf{D}$や$\textbf{d}$や$\textbf{E}$において$\textbf{D}$から$\textbf{C}$までの弧$\textbf{CD}$や$\textbf{Cd}$や$\textbf{CE}$の長さに比例する力である。この力を同じ弧で表すと、抵抗が時間の量によって与えられるので、それを表すために Cycloid の部分$\textbf{CO}$を同じ量で表し、$\textbf{Od}$の弧を弧$\textbf{CD}$に対する比として弧$\textbf{OB}$に対する比で定める。すると、抵抗がある中での$\textbf{d}$における力は、抵抗$\textbf{CO}$を超える力$\textbf{Cd}$の余分部分であり、これを弧$\textbf{Od}$で表すので、中が抵抗がないときの$\textbf{D}$での力は、

$$
\textbf{Od}
$$

と比例する。

\textbf{第2編 第6節：振子の運動と抵抗（命題XXIV〜XXXI）}

\medskip

弧 $Od$ は弧 $CD$ に等しい; よって、点 $B$ においても弧 $OB$ は弧 $CB$ に等しい。したがって、二体の物体 $D$, $d$ が点 $B$ から発生し、その力で推進されるとき、初めの力を $CB$ および $OB$ の弧とすると、最初の速度も同じ比率で与えられる。仮定するに弧 $BD$ および $Bd$ があり、残りの弧 $CD$, $Od$ も同じ比率で与えられる。したがって、力 $CD$ および $Od$ は始まりにおいて同じ比率を保ち、よって物体は同時に同じ比率で描く弧を進む。

\medskip
\begin{center}

\end{center}


$$\textbf{力、速度、および残りの弧} CD, Od \textbf{は常に} CD, OB \textbf{の全弧に比例します。} \therefore \textbf{したがって，物体} D, d \textbf{同時に到达点} C, O \textbf{。在介質不抵抗的情况下，物体} D \textbf{到达点} C; 在介質存在阻力的情况下，物体 d \textbf{到达点} O \textbf{。〜により在} C, O \textbf{的速度与弧} CB, OB \textbf{成比例；したがって，在継続运動時，それら所描描的弧也将以相同的比例進行。設} CE, Oe \textbf{為これ些弧。物体} D \textbf{在介質不抵抗的情况下被减速至点} E \textbf{的速度与} CE \textbf{成比例；而物体} d \textbf{在介質存在阻力的情况下被减速至点} e \textbf{的速度与} Oe \textbf{成比例。}$$


\textbf{合計した力} \(Ce\) と抵抗力 \(CO\), 即ち \(Oe\) に等しいので、物体を retardate（遅延させる）力は弧 \(CE\) および \(Oe\) の比に比例します。したがって、物体の retardation（遅延）は弧 \(CB\) および \(OB\) の比に従います。これにより、与えられた一定な比率において速度は遅延し続けます。したがって、与えられた一定な比率において、速度と描かれた弧は常に互いに弧 \(CB\) および \(OB\) の比で比例します。よって、もし全体の弧 \(AB\) および \(aB\) を与えられた同一の比率で加える場合、物体 \(D\) および \(d\) はこれらの弧を描き、位置 \(A\) および \(a\) での運動全体を同時に失います。したがって、振動 \(BA\) および \(BE\) は等時的であり、これらは与えられた全体の弧 \(BA\) および \(BE\) の比に比例します。

\[
\textbf{sunt arcuum partes qu{\ae}libet } BD, Bd \text{ vel } BE, Be \text{ qu{\ae} simul describuntur.}
\]

\textbf{Corol.} よって、抵抗のある媒質において最も速い運動は点 $C$（最下点）に到達しないが、点 $O$ で全弧 $aB$ を二等分する点を見つけられる。そして物体はさらに点 $a$ に向かう途中で、前回の下降から $B$ へ $O$ への移動と同じ程度の速度減速を経験する。

\condpagelarge{命題 XXVI. 定理 XX.}

\textit{速度に比例する抵抗を受ける、糸でつるされた物体のサイクロイド上の振動は、等時的である。}

証明。二つの物体が、つり点から等しい距離にあって、振動しながら不等な弧を描き、弧の対応する部分での速度が、互いに全弧の比をなすとする。すると、速度に比例する抵抗も、互いにそれらの弧の比をなす。そこで、重力から生じる起動力（これも同じ弧に比例する）からこれらの抵抗を引くか、あるいは足せば、その差または和は、互いに同じ弧の比をなす。そして速度の増加または減少はこの差または和に比例するから、速度はつねに全弧に比例する。したがって、速度がある場合に全弧に比例していれば、つねに同じ比を保つ。ところで運動の初め、物体が下降して弧を描き始めるときには、力が弧に比例しているから、弧に比例した速度が生じる。ゆえに速度は、描かれるべき全弧につねに比例し、したがってそれらの弧は同時に描かれる。\QED

\condpagelarge{命題 XXVII. 定理 XXI.}

\textit{糸でつるされた物体が速度の二乗の比で抵抗を受けるならば、抵抗のある媒質中での振動の時間と、同じ比重の抵抗のない媒質中での振動の時間との差は、振動で描かれる弧に、ほぼ比例する。}

証明。等しい長さの振り子が、抵抗のある媒質中で不等な弧 $A$, $B$ を描くとする。弧 $A$ における物体の抵抗は、弧 $B$ の対応する部分における抵抗に対して、速度の二乗の比、すなわちほぼ $A^2$ 対 $B^2$ の比をなす。もし弧 $B$ における抵抗が、弧 $A$ における抵抗に対して、長方形 $AB$ 対 $A^2$ の比をなすならば、前の命題により、弧 $A$ と $B$ における時間は等しくなるであろう。したがって、弧 $A$ における抵抗 $A^2$、あるいは弧 $B$ における抵抗 $AB$ は、弧 $A$ での時間が抵抗のない媒質中の時間を超える分をもたらし、また抵抗 $BB$ は、弧 $B$ での時間が抵抗のない媒質中の時間を超える分をもたらす。ところが、これらの超過分は、それをもたらす力 $AB$ と $BB$ に、したがって弧 $A$ と $B$ に、ほぼ比例する。\QED

\textit{系1.}\wsp{}これから、抵抗のある媒質中で不等な弧について行った振動の時間から、同じ比重の抵抗のない媒質中での振動の時間を知ることができる。たとえば、一方の弧が他方の2倍の大きさであるなら、時間の差は、小さいほうの弧での時間が抵抗のない媒質中の時間を超える分に対して、弧の差対小さいほうの弧の比をなす。

\textit{系2.}\wsp{}短い振動ほど等時性が高く、きわめて短い振動は、抵抗のない媒質中とほぼ同じ時間で行われる。大きな弧での振動の時間はやや長い。それは、時間を延ばす下降中の抵抗のほうが、時間を縮める次の上昇中の抵抗よりも、下降で描かれた長さに応じて大きいからである。また、短い振動も長い振動も、媒質の運動によって時間がいくらか延びるように見える。というのは、減速する物体は、速度の比から見積もるよりわずかに小さい抵抗を受け、加速する物体は、一様に進む物体よりわずかに大きい抵抗を受けるからである。これは、媒質が物体から受け取った運動のまま同じ向きに進むとき、前の場合（減速）にはよりかき乱され、後の場合（加速）にはかき乱されることが少ないので、運動している物体とより強く、あるいはより弱く歩調が合うためである。こうして振り子には、下降では速度の比から見積もるより大きく、上昇では小さく抵抗が働き、どちらの原因によっても時間は延びる。

\condpagelarge{命題 XXVIII. 定理 XXII.}

\textit{サイクロイド上で振動する、糸でつるされた物体が、時間の瞬間に比例する抵抗（一定の抵抗）を受けるならば、その抵抗の重力に対する比は、全下降で描かれた弧が次の上昇で描かれた弧を超える分の、振り子の長さの2倍に対する比に等しい。}

$BC$ を下降で描かれた弧、$Ca$ を上昇で描かれた弧、$Aa$ をこの二つの弧の差とする。命題XXVで作図し証明したことをそのまま用いると、振動する物体が任意の位置 $D$ で押される力は、抵抗に対して、弧 $CD$ 対弧 $CO$ の比をなす。ここで $CO$ は差 $Aa$ の半分である。したがって、振動する物体がサイクロイドの始め、すなわち最も高い点で押される力、つまり重力は、抵抗に対して、その最高点と最低点 $C$ との間のサイクロイドの弧が弧 $CO$ に対するのと同じ比をなす。すなわち（弧を2倍にして）サイクロイド全体の弧、つまり振り子の長さの2倍が、弧 $Aa$ に対する比をなす。\QED

\condpagelarge{命題XXIX． 問題 VII．}

\textbf{置かれている条件}：振動する物体が円錐曲線（ Cycloide）上を運動し、その運動には抵抗があるとします。このとき、各点での抵抗力を求めます。


\par\begin{center}
    \centering
    \includegraphics[width=1500px,height=1143px,trim=-24 0 0 0,clip]{317.png}

\end{center}\par

円錐曲線（ Cycloide）の弧 $Ba$ を描いたとき、点 $C$ は円錐曲線の最低点であり、半径 $CZ$ の長さが振り子の長さと等しいとします。任意の点 $D$ での抵抗力を求めます。

無限大の直線 $OQ$ が点 $O, C, P, Q$ を通るとき、垂線 $OK, CT, PI, QE$ は中心 $O$ から引かれます。

\textbf{近接する直線} \( OK \), \( OQ \)を描くと、双曲線 \( TIGE \)が垂線 \( CT \), \( PI \), \( QE \)と交わり、点 \( I \)を通る。また、点 \( I \)から \( KF \)を引き、この直線が近接する直線 \( OK \)と交わる点を \( K \)とする。さらに、垂線 \( CT \), \( QE \)はそれぞれ点 \( L \)と \( F \)で交差する。

双曲線 \( PIEQ \)の面積は、点 \( I \)から描かれた垂線 \( CT \)と \( QE \)が点 \( L \)と \( F \)で交差するとき、双曲線 \( PITE \)の面積に対する比として表される。すなわち、双曲線 \( PIEQ \)の面積は、体が下降した弧 \( BC \)に対する比として、双曲線 \( PITC \)の面積と弧 \( Ca \)を上昇した弧に対する比として、また双曲線 \( IEF \)の面積と点 \( OQ \)から点 \( OC \)までの距離に対する比として表される。

次に、垂線 \( MN \)で双曲線 \( PINM \)を切り取った面積が、双曲線 \( PIEQ \)に対する比として双曲線 \( PINM \)の弧 \( CZ \)に対する比と等しい。そして、垂線 \( RG \)で双曲線 \( PIRG \)を切り取った面積も同様に定義される。


\par\begin{center}
  \centering
  \includegraphics[width=0.8\textwidth]{318.png}

\end{center}\par

\textbf{双曲線} \( PI\Gamma \), が面積 \( PIEQ \) に比例し、任意の弧 \( CD \) が全落下中に描く弧 \( BC \) に比例するとき、点 \( D \) における抵抗は重力に対する比として、面積 \( \frac{OR}{OQ} IEF - IGH \) が面積 \( PIENM \) に対して比となる。

なぜなら、落下中に点 \( Z \), \( B \), \( D \), \( a \) に引き寄せられる力はそれぞれ弧 \( CZ \), \( CB \), \( CD \), \( Ca \) に対応し、これらの力はそれぞれ面積 \( PINM \), \( PIEQ \), \( PI\Gamma \), \( PITC \) に対応するとき、各弧とその対応する力を表す面積を定義すると、点 \( D \) から最小距離 \( Dd \) を描く間隔は平行四辺形 \( RGgr \) の面積として表される。

\& \( rg \)を\( h \)に引いて、\( GHhg \), \& \( RGgr \)が円弧面積\( IGH \), \( PIGR \)の contemporanea arearum decrementaとします。また、
\[ \frac{OR}{OQ} IEF - IGH \]
は\( GHhg - \frac{Rr}{OQ} IEF \),
すなわち
\[ Rr \times HG - \frac{Rr}{OQ} IEF, \]
が円弧面積\( PIGR \)の decrementum \( RGgr \)または\( Rr \times RG \)に等しく、
\[ HG - \frac{IEF}{OQ} \]
は\( RG \)に対して比例します。

したがって、
\[ OR \times HG - \frac{OR}{OQ} IEF \]
は
\[ OR \times GR \]
または
\[ OP \times PI \]
に等しくなります。これは以下の理由により成立します：
\[ OR \times HG = OR \times HR - OR \times GR, \]
\[ ORHK = OPIK, \]
\[ PIHR = PIGR + IGH. \]

したがって、
\[ PIGR + IGH - \frac{OR}{OQ} IEF \]
は\( OPIK \)に対して比例します。

したがって、面積 \(\frac{OR}{OQ} IEF - IGH\) を \(Y\) とし、面積 \(PIGR\) の減少量 \(RGgr\) を与えると、面積 \(Y\) の増加量は \(PIGR - Y\) となる。

もし \(V\) を重力から生じる弧 \(CD\) 上の運動を比例する力を表し、その力が点 \(D\) に作用するとし、\(R\) を抵抗力とするとき、全体的な作用力 \(V - R\) は点 \(D\) に作用する。従って、与えられた時間における粒子の速度増加量は次のようになる。

また、抵抗力 \(R\)（仮定により）は速度の平方と比例し、したがって（補題 II. による）、抵抗力の増加量は次のようになる：

\[
\textbf{補題 II.} \quad \text{抵抗の増加量は速度の増加量の平方に比例する。}
\]

この式を用いて、与えられた時間における粒子の速度増加量を求めることができる。

$$
\textbf{ut velocitas \& incrementum velocitatis conjunctim, id est ut spatium data temporis particula descriptum \& } V - R \textbf{ conjunctim; atque adeo, si momentum spatii detur, ut } V - R; \textbf{id est, si pro vi } V \textbf{ scribatur ejus exponens } PIGR, \& resistentia  R \textbf{ exponatur per aliam aliquam aream } Z, \textbf{ut } PIGR - Z.
$$

したがって、領域 $PIGR$ は与えられた動量の減少に伴い一様に減少し、領域 $Y$ は比率 $PIGR - Y$ で増加し、領域 $Z$ は比率 $PIGR - Z$ で増加します。したがって、領域 $Y$ と $Z$ が同時に始まり、初期において等しい場合、これらの領域は加算によって
$$
\textbf{per additionem}
$$
ます。

{\ae}qualium momentorum pergent esse {\ae}quales, \& {\ae}qualibus itidem momentis subinde decrescentes simul evanescent. Et vicissim, si simul incipiunt \& simul evanescunt, {\ae}qualia habebunt momenta \& semper erunt {\ae}quales: id adeo quia si resistentia $Z$ augeatur, velocitas una cum arcu illo $Ca$, qui in ascensu corporis describitur, diminuetur; \& puncto in quo motus omnis una cum resistentia cessat propius accedente ad punctum $C$, resistentia citius evanescet quam area $Y$. Et contrarium eveniet ubi resistentia diminuitur.

Jam vero area $Z$ incipit desinitque ubi resistentia nulla est, hoc \textbf{は}、抵抗が存在しない点で面積 $Z$ が始まり、その終了するとき、抵抗のない点においてです。

この状況では、物体の上昇中に描かれる弧 $Ca$ と共に速度が減少します。そして、運動と抵抗ともに終焉する点 $C$ に近づくにつれて、抵抗は面積 $Y$ より早く消失します。逆もまた真で、抵抗が減少すると異なります。

したがって、抵抗 $Z$ が増加すると、物体の速度と共に弧 $Ca$ の大きさも減少し、運動と抵抗ともに終焉する点 $C$ に近づくにつれて、抵抗は面積 $Y$ より早く消失します。逆の場合でも同様です。

抵抗が存在しない点で面積 $Z$ が始まり、その終了するとき、抵抗のない点において、速度と弧 $Ca$ の大きさは常に等しくなります。これは、抵抗 $Z$ が増加すると物体の速度と共に弧 $Ca$ の大きさも減少し、運動と抵抗ともに終焉する点 $C$ に近づくにつれて、抵抗は面積 $Y$ より早く消失することを意味します。

逆の場合でも同様です。抵抗が減少すると、物体の速度と共に弧 $Ca$ の大きさも増加し、運動と抵抗ともに終焉する点 $C$ に近づくにつれて、抵抗は面積 $Y$ より遅く消失します。

\textbf{そのとき、運動の開始と終了において、弧 \( CD \), \( CD \) は弧 \( CB \) および \( Ca \) に等しくなり、よって直線 \( RG \) が直線 \( QE \) および \( CT \) に交わる点で入射する。}

そして面積 \( Y \) すなわち \( \frac{OR}{OQ} IEF - IGH \) は始まりと終わりにおいて消失し、よって \( \frac{OR}{OQ} IEF \) および \( IGH \) が等しくなる点で始まりと終わりとなる：これは（図の作図により）直線 \( RG \) が直線 \( QE \) および \( CT \) に交わる点で入射することを意味する。従って
これらの面積は同時に始まり、同時に消失し、よって常に等しい。したがって面積 \( \frac{OR}{OQ} IEF - IGH \) は抵抗 \( Z \) と等しくなり、よってそれは

\begin{equation}
\textbf{\(\frac{OR}{OQ} IEF - IGH = Z\).}
\end{equation}




\textbf{命題 XXIV. Corol.\,1.} 重力に対して抵抗が存在する地点 \( C \)において、抵抗は重力に比例し、領域 \( PINM \)に対する領域 \( IEF \)と \( OP \)に対する \( OQ \)の比である。

\textbf{命題 XXIV. Corol.\,2.} この抵抗は最大となる地点で、領域 \( PIHR \)が領域 \( IEF \)に対して \( OR \)が \( OQ \)に対しての比である。なぜならその場合、その運動量（すなわち \( PIGR - Y \)) はゼロになるからである。

\textbf{命題 XXIV. Corol.\,3.} したがって、各地点における速度も明らかとなる：実際、それは抵抗に対する半分の比であり、そして運動そのものに対する半分の比でもある。




初めに，振動体の速度を Cycloide 上で等しくし，抵抗なしに置く。

しかしながら，この命題において抵抗と速度を見つけるための計算は困難であるから，次なる命題を追加する。それはより一般的であり，哲学的な用途に充分な正確さを持つ。

\condpagelarge{命題 XXX. 理論 XXIII.}

もし直線 $aB$ が Cycloide の弧と等しくし，振動体がその点 $D$ を描く垂直線 $DK$ が pendulum の長さに対して抵抗の比を表すとするならば，

\[
\textbf{直線} \, aB \, \textbf{は} \, D \, \textbf{点における} \, DK \, \textbf{の長さと} \, \textbf{pendulum の抵抗} \, \textbf{との比を示す。}
\]

この命題はより一般的であり，哲学的な用途に充分な正確さを持つ。

\textbf{命題 XXIV. 定理 1.}

重力による力を考慮すると、自由落下と上昇の全弧との差をその半和に引くと、得られる値は直線 $aB$ 上のすべての垂線によって占領される面積 $BKaB$ に近似して等しくなる。


\par\begin{center}
    \centering
    \includegraphics[width=1334px,height=869px,trim={0 -24 0 0},clip]{321.png}

\end{center}\par

次に、全振動の Cycloid の弧 $aB$ を直線 $aB$ で表し、真空中で描かれる弧 $AB$ を考えます。$AB$ の中点 $C$ を取ると、点 $C$ は Cycloid の最低点を表します。

\textbf{命題 XXV. 定理 2.}

全振動の Cycloid の弧 $aB$ と真空中で描かれる弧 $AB$ の差をその半和に引くと、得られる値は直線 $aB$ 上のすべての垂線によって占領される面積 $BKaB$ に近似して等しくなる。

\textbf{命題 XXVI. 定理 3.}


\textbf{命題 XXVII. 定理 4.}


\textbf{命題 XXVIII. 定理 5.}


\textbf{命題 XXIX. 定理 6.}

全振動の Cycloid の弧 $aB$ と真空中で描かれる弧 $AB$ の差をその半和に引くと、得られる値は直線 $aB$ 上のすべての垂線によって占領される面積 $BKaB$ に近似して等しくなる。

\textbf{命題 XXX. 定理 7.}


\textbf{命題 XXXI. 定理 8.}


\[
CD \text{ ut vis a gravitate oriunda, qua corpus in } D \text{ secundum Tangentem Cycloidis urgetur, eamque habebit rationem ad longitudinem Penduli quam habet vis in } D \text{ ad vim gravitatis.}
\]

したがって、
\[
CD \text{ を重力から生じる力を表し、} D \text{ における Cycloidis の接線方向に物体が作用される力を } D \text{ における Penduli の長さに対する比で持つとします。}
\]

したがって、
\[
\]

したがって、
\[
\]

したがって、
\[
CD \text{ を重力から生じる力を表し、} D \text{ における Cycloidis の接線方向に物体が作用される力を } D \text{ における Penduli の長さに対する比で持つとします。}
\]

したがって、
\[
\]

したがって、
\[
\]

したがって、
\[
CD \text{ を重力から生じる力を表し、} D \text{ における Cycloidis の接線方向に物体が作用される力を } D \text{ における Penduli の長さに対する比で持つとします。}
\]

したがって、
\[
\]

したがって、
\[
\]

したがって、
\[
CD \text{ を重力から生じる力を表し、} D \text{ における Cycloidis の接線方向に物体が作用される力を } D \text{ における Penduli の長さに対する比で持つとします。}
\]

したがって、
\[
\]

したがって、
\[
\]

したがって、
\[
CD \text{ を重力から生じる力を表し、} D \text{ における Cycloidis の接線方向に物体が作用される力を } D \text{ における Penduli の長さに対する比で持つとします。}
\]

したがって、
\[
\]

したがって、
\[
\]

したがって、
\[
CD \text{ を重力から生じる力を表し、} D \text{ における Cycloidis の接線方向に物体が作用される力を } D \text{ における Penduli の長さに対する比で持つとします。}
\]

したがって、
\[
\]

したがって、
\[
\]

したがって、
\[
CD \text{ を重力から生じる力を表し、} D \text{ における Cycloidis の接線方向に物体が作用される力を } D \text{ における Penduli の長さに対する比で持つとします。}
\]

したがって、
\[
\]

したがって、
\[
\]

したがって、
\[
CD \text{ を重力から生じる力を表し、} D \text{ における Cycloidis の接線方向に物体が作用される力を } D \text{ における Penduli の長さに対する比で持つとします。}
\]

したがって

\textbf{円周} \( Ee \) に当たる速度は、点 \( B \) から自由落下して \( D \) と \( d \) の位置で得られる速度と等しくなります。これは命題 LII（第1編）により明らかです。

これらの速度を次のように表します：
\[
\textbf{DE}, \, de
\]
点 \( B \) から抵抗のある媒質を通じて落下して \( D \) に到達するときの速度を \( DF \) とします。円心 \( C \) を中心とする半径 \( CF \) の円 \( FfM \) が直線 \( de \) および \( AB \) と交わる点 \( f \) および \( M \) を通ります。このとき、点 \( M \) は無抵抗条件下でさらに上昇する軌道の一部となります。また、\( df \) は点 \( d \) に到達するときに得られる速度です。

したがって、速度 \( Fg \) を表す場合も同様の論理が適用されます。

\textbf{動量} \( D \)を有する物体が、抵抗の作用により最短距離 \( Dd \)だけ進むとき、その物体は抵抗なしで到達する点 \( N \)まで上昇します。ここで、\( CN = Cg \)とします。すると、物体は点 \( N \)からさらに上昇することはありません。したがって、点 \( M \)から点 \( N \)までの距離 \( MN \)は、速度の減少によって生じる上昇量の減少を表します。

次に、直線 \( df \)に対して垂直な直線 \( Fm \)を引きます。物体の速度 \( DF \)が抵抗 \( DK \)により減少した量 \( fg \)は、同様に抵抗 \( CD \)によって生じた速度の増加量 \( fma \)と等しいです。これは、生成する力 \( DK \)と生成される力 \( CD \)との比から導かれます。

さらに、相似な三角形 \( Fmf \), \( Fhg \), \( FDC \)より、

\[
\frac{fm}{Fm} = \frac{Fh}{Fg} = \frac{FD}{FC}
\]

となります。

\[\begin{gathered}
\text{次に、} CD \text{に対する} DF, \text{そして等しく} Fg \text{に対する} Dd \text{は} DK \text{に対する} DF. \\
\text{また、} Fg \text{に対する} Fh \text{は} CF \text{に対する} DF; \text{そして等しく} Fh \text{すなわち} MN \text{に対する} Dd \text{は} DK \text{に対する} CF. \\
\text{さらに、} DR \text{に対する} \frac{1}{2}aB \text{は} DK \text{に対する} CF; \text{したがって} MN \text{に対する} Dd \text{は} DR \text{に対する} \frac{1}{2}aB. \\
\text{よって、} MN \times \frac{1}{2}aB, \text{すなわち} Aa \times \frac{1}{2}aB \text{は} Dd \times DR \text{の総和に等しい;} \\
\text{すなわち、} BRrSa \text{の面積と等しくなります。または} DRrd \text{が合成します。} Aa \text{と} aB \text{は} P \text{と} O \text{で二等分されます;} 
\end{gathered}\]



\[
\text{したがって、}
\]


$\frac{1}{2}aB$ \textbf{ または } $OB = CP$，従って $DR$ は $DK$ に比して $CP$ は $CF$ または $CM$ に比する。
したがって、$KR$ は $DR$ に比して $PM$ は $CP$ に比する。

したがって、点 M が振動の中心 O の位置で体が振幅 A 内で振動しているとき、近傍の点 P に近づき、前半の振動では A と P の間を、後半の振動では P と a の間を等しく移動し、それぞれの場合において P から反対方向へ等しい距離に移動する場合、点 K は振動の中心 O の近傍、すなわち点 O の周辺、例えば V に P から R に移動します。前半の振動では


\textbf{第2編 第6節：振子の運動と抵抗（命題XXIV〜XXXI）}


したがって、点 \( R \) から \( E \) へおよび後方では \( R \) から \( D \) へは等しく反対方向に移動します。よって、線分 \( KR \) によって描かれる面積は、前半の振動部分において点 \( R \) の外側にあり、後半の部分では点 \( R \) の内側にあります；そしてこれらはほぼ等しい面積を有します。したがって、前半の部分で面積 \( BRSa \) が加えられ、後半の部分で同様の面積が引かれると、残る面積 \( BKTa \) は面積 \( BRSa \) にほぼ等しいものとなります。したがって、直角三角形 \( Aa \times \frac{1}{2}aB \) や \( AaO \) の面積は、面積 \( BRSa \) にほぼ等しく、また面積 \( BKTa \) にもほぼ等しいものとなります。証明終了。\(\blacksquare\)


\textbf{補遺}．したがって、抵抗と弧 $Ca$, $CB$ の差 $Aa$ から、抵抗に対する重力の比を近似して求めることができます。

なぜなら、もし抵抗が一様であるとするとき、図形 $aBKkT$ は直角三角形であり、辺 $Ba$ と $DK$ 上に存在します。よって、半分の辺 $Ba$ と $Aa$ の直角三角形の面積は、辺 $Ba$ と $DK$ の直角三角形の面積と等しくなります。これにより、$DK$ が抵抗を表し、振り子の長さが重力を表すとき、抵抗に対する比は $\frac{1}{2}Aa$ に対して振り子の長さとなります。これは命題 XXVIII で証明されています。

\textbf{訳注}：補遺部分は、抵抗と重力との関係を近似して求めることについて述べています。具体的には、一様な抵抗を考えると、直角三角形の面積の比から、抵抗に対する重力の比を求めることができます。この結果は、命題 XXVIII で詳細に証明されています。



\text{もし抵抗が速度に比例する場合、図形} \ aBKkT \ \text{は近似して楕円となります。なぜなら、非抵抗媒質中で完全な振動を描く物体が} \ BA \ \text{の長さを描くとき、任意の点} \ D \ \text{での速度は} \ AB \ \text{に描かれた円の直径} \ DE \ \text{と比例します。}

\text{したがって、抵抗媒質中} \ Ba \ \text{、非抵抗媒質中} \ BA \ \text{の場合で、近似して同じ時間内に描かれるとき、}
Ba \ \text{上の各点での速度は} \ BA \ \text{上の対応する点の速度とほぼ比例し、それが} \ Ba \ \text{です。}





\textbf{BA}方向に運動するとき、抵抗が中心を通る円または楕円上に描かれた弦の逆数に比例すると仮定すると、速度は点Oで半径DKと中間抵抗OVとの関係で与えられる。よって、図形BKVTaは近似して楕円となる。抵抗が速度に対して比例するとして、点Oでの抵抗を表すOVを用いる。中心Oを持つ楕円の半長軸OBと半短軸OVから描かれた楕円 BKVT は、直角三角形AaBOの面積に近似される。したがって、
\[
Aa \times BO : OV \times BO = \text{この楕円の面積} : OV \times BO.
\]
すなわち
\[
Aa : OV = \text{半円の面積} : OV,
\]
と表される。


ニュートンの『自然哲学の数学的諸原理』第2編 第6節：振子の運動と抵抗（命題XXIV〜XXXI）の一部です。

\textbf{証明}：
半径の平方に比例し、すなわち約11対7の比率にあります。従って：
\[
\frac{7}{11}Aa \text{ は } 振子の抵抗在点  O \text{ における長さと、振動体の重力との比に相当します。}
\]

もし抵抗 $DK$ が速度の二乗比で与えられる場合、図形 $BKTVa$ は頂点 $V$ を持つ放物線であり、軸は $OV$ です。したがって、この放物線は直角三角形 $Ba$ と $OV$ の約三分の一に近い二つの部分に等しい面積を持ちます。よって、半径 $Aa$ の半分 $Ba$ と $Aa$ の矩形の面積は、$Ba$ と $OV$ の約三分の二の矩形の面積に等しくなります。したがって：
\[
OV = \frac{3}{4}Aa
\]
従って、振子の抵抗在点 $O$ における重力との比は：
\[
\textbf{OV} = \frac{3}{4}\textbf{Aa}
\]

以上が翻訳です。必要に応じて画像を配置することも可能です。

$\frac{3}{4}Aa$ は振子の長さに対する比である。

これらの結論は実践的な問題において十分に正確であると考える。なぜなら、楕円や放物線が点$V$で半分の位置に一致する場合、それらが一方へ超えると他方へ欠如し、その結果、近傍では同じ状態になる。

\condpagelarge{命題XXXI. 理論XXIV.}

\textbf{もし振動体の抵抗が描かれた各弧の部分に比例して増加または減少するなら、その差異}





\text{下降弧とその後の上升弧を描く角度差は、近似して同一の比率で増加または減少します。}

\text{これは、媒質による振子の retardation から来ますので、全 retardation はその retardation に比例する抵抗に比例します。}

\text{先の命題において、直線} $\frac{1}{2}aB$ \text{と弧} $CB, Ca$ \text{の差異} $Aa$ \text{は面積} $BKT$ \text{に等しいでした。}
\text{もし長さ} $aB$ \text{が維持されるなら、その面積は} $DK$ \text{を順次適用した比率で増加または減少します;} 


この部分は以下のようになります。

\[
\textbf{これにより、抵抗の量は} \frac{\mathrm{d}}{\mathrm{dt}}(aB) \& \textbf{抵抗と連続的に比例する。したがって} aB \& \textbf{抵抗との積} \textbf{は} Aa \& \textbf{抵抗と連続的に比例するので、} Aa \textbf{は抵抗に比例する。}
\]
\(\blacksquare\)

\textit{\textbf{補題1.}}\,従って、もし抵抗が速度と比例するなら、同じ媒質における弧の差は全弧と比例する：逆もまた然り。

\textit{\textbf{補題2.}}\,もし抵抗が速度の二倍比に比例するなら、その差は全弧の二倍比に比例する：逆もまた然り。

\textit{\textbf{補題3.}}\,そして一般に、もし抵抗が速度の三倍比または他の何かに比例するなら、

（訳文をここで中断します。続きを翻訳する場合はご通知ください。）


$\textbf{第2編 第6節：振子の運動と抵抗（命題XXIV〜XXXI）}$


任意の速度の比において、弧の全長に対する差異もそのような比にあります。逆もまた然り。

$\textbf{補則 4.}$
そしてもし抵抗が単純な速度の比と部分的に同一であり、部分的にその二倍の比であるとき、差異も弧の全長に対する比と部分的にその二倍の比にあります。逆もまた然り。抵抗に関する法則およびその速度に対する比率は、弧の長さに対する差異と同じ比率を持つことになります。

$\textbf{補則 5.}$
従って、不等長の弧を順次描く振子において、抵抗の増加量と減少量の比が求められることが可能です。

この弧を描く長さに対する、抵抗の増加量と減少量との比は、速度がより大きくなるかより小さくなるかによります。



この論証は、運動の法則と抵抗の性質から導き出されます。まず、物体が一定の時間内に移動する距離に対する抵抗の変化を考慮します。速度が大きくなるほど、抵抗も増加し、速度が小さくなるほど、抵抗も減少します。したがって、弧を描く長さに対する抵抗の増加量と減少量との比は、速度の大小によります。




\section{第2編 第7節：流体の運動と投射体の抵抗（命題XXXII〜XL）}


% Section: 第2編 第7節：流体の運動と投射体の抵抗（命題XXXII〜XL）



\condpagelarge{命題 XXXII. 定理 XXV.}

\textit{二つの粒子系があって、それぞれ等しい数の粒子から成り、一方の系の各粒子が他方の系の対応する各粒子と相似で、互いに与えられた密度の比をもち、また時間に比例して互いに（一方の系の中どうし、他方の系の中どうしで）相似に動き始めるとする。さらに、同じ系の中の粒子は、反射の瞬間以外は互いに接触せず、また、対応する粒子の直径に反比例し速さの二乗に正比例する加速力によるほかは、互いに引き合うことも退け合うこともないとする。このとき、二つの系の粒子は、時間に比例して互いに相似に動きつづける。逆もまた成り立つ。}

時間に比例して互いに相似に動く物体とは、それらの時間の終わりごとに、互いの位置がつねに相似であるような物体をいう。たとえば、一方の系の粒子を他方の系の対応する粒子と比べるときがそうである。したがって、相似な図形の相似で比例した部分が、対応する粒子によって描かれる時間は、比例する。

そこで、このような二つの系があれば、対応する粒子は、動き始めた運動が相似であるから、互いに出会うまで、相似に動きつづける。というのは、もし粒子がどんな力にも動かされなければ、運動の法則Iにより、直線上を一様に進むからである。もし何らかの力で互いに動かし合うとして、その力が対応する粒子の直径に反比例し速さの二乗に正比例するなら、粒子の位置は相似で、力は比例しているから、対応する粒子を動かす全体の力は、（法則の系2により）個々の動かす力から合成されて、相似な向きをもつ。それはあたかも、粒子の間の相似な位置にある中心を向いているようである。そしてこれらの全体の力は、互いに、それを合成する個々の力と同じ比をなす。すなわち、対応する粒子の直径に反比例し、速さの二乗に正比例する。したがって、対応する粒子は相似な図形を描きつづける。このことは、中心が静止しているならば、命題IVの系1, 2, 7により成り立つ。もし中心が動くなら、移動が相似であるから、系の粒子の間の位置は相似なままであり、粒子が描く図形には相似な変化が生じる。それゆえ、対応する相似な粒子の運動は、最初の出会いまで相似であり、したがって出会いも相似で、反射も相似であり、そののち（すでに示したことにより）、ふたたび互いにぶつかるまで、互いに相似な運動をする。こうして、限りなくつづく。\QED

\textit{系1.}\wsp{}したがって、もし相似で、系の対応する粒子に対して相似な位置にある二つの物体が、時間に比例して相似に動き始め、それらの密度が、対応する粒子の密度の比に等しいならば、これらの物体は、時間に比例して相似に動きつづける。なぜなら、どちらの系でも、大きな部分の比は粒子の比と同じだからである。

\textit{系2.}\wsp{}また、系の相似で相似な位置にある部分がすべて互いに静止していて、そのうちの二つ、他の部分より大きく、両方の系で互いに対応するものが、相似な位置の線に沿って、何らかの相似な運動で動き始めるならば、これらは系の残りの部分に相似な運動を引き起こし、互いに時間に比例して相似に動きつづけ、したがって、自分の直径に比例する距離を描く。

\condpagelarge{命題 XXXIII. 定理 XXVI.}

\textit{同じ仮定のもとで、系の大きな部分が受ける抵抗は、それらの速さの二乗の比、直径の二乗の比、および系の部分の密度の比から合成された比をなす。}

抵抗は、一部は、系の粒子が互いに動かし合う向心力や遠心力から、一部は、粒子と大きな部分との衝突と反射から生じる。前の種類の抵抗は、それを生じるすべての起動力に比例する。すなわち、すべての加速力と、対応する部分の物質の量との積に比例する。これは（仮定により）速さの二乗に正比例し、対応する粒子の間の距離に反比例し、対応する部分の物質の量に正比例する。そして（一方の系の粒子の距離が、他方の系の対応する距離に対して、前の系の粒子または部分の直径が、他方の系の対応する粒子または部分の直径に対するのと同じ比をなし、物質の量は部分の密度と直径の三乗に比例するから）抵抗は互いに、速さの二乗、直径の二乗、および系の部分の密度に比例する。\QED

後の種類の抵抗は、対応する反射の回数と力との積に比例する。ところが反射の回数は、互いに、対応する部分の速さに正比例し、その反射の間の空間に反比例する。そして反射の力は、対応する部分の速さ、大きさ、密度の積に比例する。すなわち、速さと、直径の三乗と、部分の密度に比例する。これらすべての比を合わせれば、対応する部分の抵抗は互いに、速さの二乗、直径の二乗、部分の密度の積に比例する。\QED

\textit{系1.}\wsp{}したがって、もしこれらの系が、空気のような弾性のある二つの流体で、その部分が互いに静止しており、また大きさと密度が流体の部分に比例する相似な二つの物体が、それらの部分の間に相似に置かれて、相似な位置の線に沿って任意に投げ出され、さらに流体の粒子が互いに動かし合う起動力が、投げ出された物体の直径に反比例し速さの二乗に正比例するならば、これらの物体は、時間に比例して流体に相似な運動を引き起こし、相似で、自分の直径に比例する距離を描く。

\textit{系2.}\wsp{}したがって、同じ流体の中では、速い投射体は、ほぼ速さの二乗の比で抵抗を受ける。なぜなら、もし離れた粒子が互いに動かし合う力が速さの二乗の比で増すなら、投射体は正確に同じ二乗の比で抵抗を受けるであろう。それゆえ、離れた部分が互いに力をまったく及ぼさない媒質では、抵抗は正確に速さの二乗の比になる。そこで、相似で等しい部分から成り、等しい距離にしたがって規則的に配置された三つの媒質 $A$, $B$, $C$ を考える。媒質 $A$ と $B$ の部分は、互いに $T$ 対 $V$ の比をなす力で退け合い、媒質 $C$ の部分は、そのような力をまったくもたないとする。そして等しい四つの物体 $D$, $E$, $F$, $G$ がこれらの媒質の中を動くとし、前の二つ $D$, $E$ は前の二つの媒質 $A$, $B$ の中を、あとの二つ $F$, $G$ は第三の媒質 $C$ の中を動くとする。物体 $D$ の速さの物体 $E$ の速さに対する比、および物体 $F$ の速さの物体 $G$ の速さに対する比が、力 $T$ の力 $V$ に対する比の平方根であるとすれば、物体 $D$ の抵抗の物体 $E$ の抵抗に対する比と、物体 $F$ の抵抗の物体 $G$ の抵抗に対する比は、速さの二乗の比になり、したがって物体 $D$ の抵抗の物体 $F$ の抵抗に対する比は、物体 $E$ の抵抗の物体 $G$ の抵抗に対する比に等しい。いま物体 $D$ と $F$ が等しい速さで、$E$ と $G$ も等しい速さであるとする。そして物体 $D$ と $F$ の速さを任意の比で増し、媒質 $B$ の粒子の力を同じ二乗の比で減らしていくと、媒質 $B$ は、思いのままに媒質 $C$ の形と状態に近づく。それゆえ、これらの媒質の中の、等しく等速の物体 $E$ と $G$ の抵抗は、たえず等しさに近づき、その差はついには任意に与えられた量より小さくなる。したがって、物体 $D$ と $F$ の抵抗の比は、物体 $E$ と $G$ の抵抗の比に等しいから、これもまた同様に等しさの比に近づく。ゆえに、物体 $D$ と $F$ がきわめて速く動くとき、それらの抵抗はほぼ等しい。そして物体 $F$ の抵抗は速さの二乗の比であるから、物体 $D$ の抵抗もほぼ同じ比になる。\QED

\textit{系3.}\wsp{}したがって、任意の弾性流体の中をきわめて速く動く物体の抵抗は、もし流体の弾性が粒子の遠心力から生じるならば、流体の部分がその遠心力をもたず、互いに退け合わないとしたときの抵抗とほとんど同じである。

\textit{系4.}\wsp{}したがって、部分が離れていても互いに退け合わない媒質の中で、相似で等速の物体の抵抗は直径の二乗に比例するから、弾性流体の中を、等しい速さで、きわめて速く動く物体の抵抗も、ほぼ直径の二乗に比例する。

\textit{系5.}\wsp{}また、粒子が互いに退け合わない、同じ密度の媒質の中で、相似で等しく等速な物体は、その粒子が、より多くて小さくても、より少なくて大きくても、等しい時間に等しい量の物質にぶつかり、それに等しい量の運動を与え、逆に（運動の法則IIIにより）それから等しい反作用を受ける。すなわち、等しく抵抗を受ける。したがって、同じ密度の弾性流体の中でも、きわめて速く動くとき、それらの抵抗はほぼ等しい。その流体が粗い粒子から成っていても、きわめて微細な粒子から成っていてもかまわない。媒質の微細さによって、きわめて速く動く投射体の抵抗は、あまり減らない。

\textit{系6.}\wsp{}ところで、流体の粒子は、互いに退け合う力のために、まわりの他の粒子を同時に動かさずには動けず、したがって、そのような力をもたない場合よりも、互いに動きにくい。そして遠心力が大きいほど、互いにいっそう動きにくい。そこで、そのような流体の中の投射体は、その力が強いほど動きにくいことは明らかである。したがって、前の系で扱った、きわめて速い物体の速さが小さくなるとき、もし粒子の力が同じ二乗の比で小さくなるなら、抵抗は速さの二乗の比で小さくなるはずである。ところが力はまったく小さくならないのだから、抵抗は、速さの二乗より小さい比で小さくなることは明らかである。

\textit{系7.}\wsp{}さらに、遠心力が抵抗を増すのは、粒子が、その力によって自分の運動を、流体を通してより遠くまで伝えるからである。そしてその距離は、大きな物体に対しては小さな比をもつ。したがって、その力から生じる抵抗の増加は、大きな物体ではより小さなことは明らかである。それゆえ、物体が大きいほど、減速する物体の抵抗は、より正確に速さの二乗の比で減少する。

\textit{系8.}\wsp{}それゆえ、この二乗の比は、密度と弾性が同じで、より小さな粒子から成る流体で、より正確に成り立つ。というのは、もし大きな物体が小さくなり、流体の粒子も、密度と弾性を保ったまま同じ比で小さくなるなら、抵抗の比は前と同じままである。これは、前のことから容易に導かれる。

\textit{系9.}\wsp{}これらのすべては、弾性が粒子の遠心力に由来する流体について成り立つ。もしその力が別のところから生じるなら、たとえば粒子が羊毛や木の枝のように広がることから、あるいはほかの何らかの原因で、粒子の運動が互いにあまり自由でなくなるなら、媒質の流動性が小さいために、抵抗は、前の系で述べたより大きくなる。

\condpagelarge{命題 XXXIV. 定理 XXVII.}

\textit{前の二つの命題で証明したことは、系の粒子が互いに接している場合にも、その粒子がきわめてなめらかであるならば、成り立つ。}

粒子が何らかの力で互いに退け合い、その力が粒子の表面に近づくにつれて限りなく増し、逆に、そこから遠ざかるにつれて急速に減って、たちまち消えるものと考える。また、系が圧縮されて、その部分が、この力が接触を妨げる分を除いて、互いに接していると考える。ところで、粒子の力が広がる空間をきわめて狭くして、粒子がほとんど互いに接するようにすれば、粒子どうしの運動は、互いに接しているときとほぼ同じになる。粒子は、きわめてなめらかな場合と同じ容易さで互いにすべり、互いにぶつかれば、前に述べた力によって、あたかも弾性があるかのように、はね返される。したがって、粒子が接していないことによるごくわずかな間隔が違いを生む分を除いて、どちらの場合も運動は同じである。この違いは、粒子の間隔を小さくすることで、限りなく小さくできる。ところで、前の二つの命題で証明したことは、粒子が接していない場合に成り立ち、それは、力の広がる空間を小さくして粒子の間隔を限りなく小さくしても成り立つ。したがって、同じことは、粒子が接している場合にも、ついには任意に与えられた差より小さくなる違いを除いて、成り立つ。そこで私は、それが正確に成り立つと言う。もしそうでないというなら、任意の場合にその違いを示してみよ。ところが、その違いが任意に与えられた量より小さいことはすでに示した。ゆえに、違いは誤って示されているのであって、違いはない。\QED

\textit{系1.}\wsp{}したがって、もし二つの系の部分がすべて互いに静止していて、ただし、他の部分より大きく、互いに対応し、他の部分に対して相似に置かれた二つの部分だけは動くとして、これらが相似な位置の線に沿って任意に投げ出されるならば、これらは系の中に相似な運動を引き起こし、時間に比例して、相似で自分の直径に比例する距離を描きつづける。そして速さの二乗の比、直径の二乗の比、および系の密度の比から合成された比で抵抗を受ける。

\textit{系2.}\wsp{}それゆえ、もしこれらの系が相似な二つの流体で、その二つの大きな部分がその中に投げ入れられた物体であるなら、そして流体の粒子がきわめてなめらかで、大きさと密度が物体に比例するなら、物体は時間に比例して、相似で自分の直径に比例する距離を描きつづけ、前の系で定めた比で抵抗を受ける。

\textit{系3.}\wsp{}したがって、同じ流体の中では、大きさの与えられた投射体は、速さの二乗の比で抵抗を受ける。

\textit{系4.}\wsp{}しかし、もし流体の粒子がきわめてなめらかではなく、あるいは運動の自由を減らす何らかの力で互いに動かし合っているなら、遅い投射体は抵抗に打ち勝ちにくく、したがって、速さの二乗の比より大きく抵抗を受ける。

\condpagelarge{命題 XXXV. 定理 XXVIII.}

\textit{直径の等しい球と円柱が、希薄で弾性のある媒質の中を、円柱の軸の向きに、等しい速さできわめて速く動くとき、球の抵抗は円柱の抵抗の半分である。}

\pngright{333.png}{}{}{}

実際、（命題XXXIIIの系3により）抵抗は、流体の部分が互いに力を及ぼして退け合わないときとほぼ同じであるから、流体の部分は、そのような力をもたず、空間全体に一様に散らばっていると仮定しよう。そして、媒質が物体に及ぼす作用は（法則の系5により）、物体が静止した媒質の中を動くときも、媒質の粒子が同じ速さで静止した物体にぶつかるときも、同じであるから、物体を静止したものとみなして、動く媒質からどれだけの衝撃を受けるかを見ることにしよう。

そこで $ABKI$ を、中心 $C$、半径 $CA$ で描いた球とし、媒質の粒子が、与えられた速さで、$AC$ に平行な直線に沿って、この球にぶつかるとする。$FB$ をそのような直線とする。その上に、半径 $CB$ に等しく $LB$ をとり、球に $B$ で接する $BD$ を引く。$AC$ と $BD$ に垂線 $BE$, $DL$ を下ろす。すると、媒質の粒子が直線 $FB$ に沿って斜めに入射して球を $B$ で打つ力は、同じ粒子が球のまわりに軸 $ACI$ で描いた円柱 $ONGQ$ を $b$ で垂直に打つ力に対して、$LD$ 対 $LB$、すなわち $BE$ 対 $BC$ の比をなす。さらに、この力が、球を入射の向き $FB$ すなわち $AC$ の向きに動かす効力は、その同じ力が球を、その作用の向き、すなわち球を直接押す直線 $BC$ の向きに動かす効力に対して、$BE$ 対 $BC$ の比をなす。これらの比を合わせれば、直線 $FB$ に沿って斜めに球に入射する粒子が、球を入射の向きに動かす効力は、同じ粒子が円柱に垂直に入射して、円柱を同じ向きに動かす効力に対して、$BE^2$ 対 $BC^2$ の比をなす。ゆえに、円柱の円形の底面 $NAO$ に垂線 $bHE$ を立て、$bE$ を半径 $AC$ に等しくし、$bH$ を $\frac{BE^2}{CB}$ に等しくすれば、$bH$ 対 $bE$ は、粒子の球に対する効果の、粒子の円柱に対する効果に対する比に等しい。したがって、すべての直線 $bH$ が占める立体は、すべての直線 $bE$ が占める立体に対して、すべての粒子の球に対する効果の、すべての粒子の円柱に対する効果に対する比をなす。ところが、前の立体は、頂点 $V$、軸 $CA$、通径 $CA$ で描かれた放物線体であり、後の立体は、この放物線体に外接する円柱である。そして放物線体が外接円柱の半分であることはよく知られている。ゆえに、媒質が球に及ぼす全体の力は、同じ媒質が円柱に及ぼす全体の力の半分である。したがって、もし媒質の粒子が静止していて、円柱と球が等しい速さで動くとすれば、球の抵抗は円柱の抵抗の半分になる。\QED

\condpagelarge{注.}

\pngright{334.png}{}{}{}

同じ方法で、ほかの図形も、抵抗について互いに比べることができ、また、抵抗のある媒質の中で自分の運動を続けるのにより適した図形も見いだせる。たとえば、中心 $O$、半径 $OC$ で描いた円形の底面 $CEBH$ と、高さ $OD$ とによって、同じ底面と高さで作られ、軸の向きに $D$ のほうへ進む円錐台のうち、最も抵抗を受けないものを作るなら、高さ $OD$ を $Q$ で二等分し、$OQ$ を $S$ まで延長して $QS$ を $QC$ に等しくすれば、$S$ は求める円錐台のもとの円錐の頂点である。

ここから、ついでに言えば、角 $CSB$ はつねに鋭角であるから、もし立体 $ADBE$ が、楕円形または卵形の図形 $ADBE$ を軸 $AB$ のまわりに回転して生じるとし、その生成する図形に三つの直線 $FG$, $GH$, $HI$ が点 $F$, $B$, $I$ で接し、$GH$ は接点 $B$ で軸に垂直で、$FG$, $HI$ が $GH$ とつくる角 $FGB$, $BHI$ が 135 度であるとすれば、図形 $ADFGHIE$ を同じ軸 $CB$ のまわりに回転して生じる立体は、前の立体より抵抗が小さい。ただし、どちらも自分の軸 $AB$ の向きに進み、どちらも端 $B$ が先頭にあるとする。この命題は、船を造るのに役立たなくはないだろうと私は考える。

\pngright{335.png}{}{}{}

もし図形 $DNFB$ が次のようなものであるとする。すなわち、その上の任意の点 $N$ から軸 $AB$ に垂線 $NM$ を下ろし、与えられた点 $G$ から、$N$ で図形に接する直線に平行な直線 $GR$ を引いて、延長した軸と $R$ で交わらせるとき、$MN$ 対 $GR$ が $GR^3$ 対 $4BR \times GB^2$ に等しい。この図形を軸 $AB$ のまわりに回転して描かれる立体は、希薄で弾性のある媒質の中を $A$ から $B$ のほうへきわめて速く動くとき、同じ長さと幅で描かれた他のどの円形の立体よりも、抵抗が小さい。

\condpagelarge{命題 XXXVI. 問題 VIII.}

\textit{希薄で弾性のある流体の中を、きわめて速く進む球形の物体の抵抗を求めること。}（図は、命題XXXVの図を見よ。）

$ABKI$ を、中心 $C$、半径 $CA$ で描いた球形の物体とする。$CA$ を、まず $S$ まで、ついで $R$ まで延長して、$AS$ を $CA$ の3分の1にし、$CR$ 対 $CS$ が球形の物体の密度の媒質の密度に対する比になるようにする。$CR$ に垂線 $PC$, $RX$ を立て、中心 $R$、漸近線 $CR$, $RX$ で任意の双曲線 $PVY$ を描く。$CR$ 上に任意の長さの $CT$ をとり、垂線 $TV$ を立てて双曲線の面積 $PCTV$ を切り取り、$CZ$ をこの面積を直線 $PC$ に当てはめた辺（面積を $PC$ で割った長さ）とする。私は次のように言う。球が距離 $CZ$ を描くあいだに、媒質の抵抗によって失う運動は、その最初の全運動に対して、ほぼ長さ $CT$ 対長さ $CR$ の比をなす。

実際、（運動の法則IIIにより）球のまわりに描いた円柱 $GNOQ$ が、媒質の粒子にぶつかって失う運動は、それが同じ粒子に与える運動に等しい。粒子がそれぞれ円柱によってはね返され、円柱が近づいたのと同じ速さで、円柱から跳ね返るものとしよう。そのような反射は、（法則の系3により）粒子がきわめて小さく、弾性によって最大にはね返されるなら、そうなるからである。したがって、粒子が円柱から跳ね返る速さは、円柱の速さに加わって、円柱の速さの2倍の全体の速さをつくる。それゆえ、円柱が各粒子の反射によって失う運動は、円柱の全運動に対して、粒子の2倍対円柱の比になる。そこで、媒質の密度の円柱の密度に対する比は $CS$ 対 $CR$ であるから、もし $Ct$ が、円柱がきわめて短い時間に描く長さであるなら、その時間に失われる運動は、円柱の全運動に対して $2Ct \times CS$ 対 $AI \times CR$ となる。実際、これが、円柱が押しのけてはね返す媒質の物質の、円柱の質量に対する比である。ところで、球は円柱の3分の2であり、球の抵抗は（前の命題により）円柱の抵抗の半分であるから、球が長さ $L$ を描くあいだに失う運動は、球の全運動に対して、$Ct \times CS$ 対 $\frac{2}{3}AI \times CR$、すなわち $Ct$ 対 $CR$ になる。垂線 $tv$ を立てて双曲線と $v$ で交わらせれば、（第2編 命題Vの系1により）物体が、面積 $CtvP$ に比例する長さを描くあいだに、その全運動 $CR$ のうちの任意の部分 $Ct$ を失うとき、面積 $CTVP$ に比例する長さを描くあいだには、その運動の部分 $CT$ を失う。ところが、長さ $Ct$ は $\frac{CPvt}{CP}$ に等しく、長さ $CZ$ は（仮定により）$\frac{CPTV}{CP}$ に等しい。したがって、長さ $Ct$ 対長さ $CZ$ は、面積 $CPvt$ 対面積 $CPVT$ に等しい。それゆえ、球がきわめて短い長さ $Ct$ を描くあいだに、全体に対して $Ct$ 対 $CR$ の部分の運動を失うなら、任意の別の長さ $CZ$ を描くあいだには、全体に対して $CT$ 対 $CR$ の部分の運動を失う。\QED

\textit{系1.}\wsp{}もし物体の最初の速さが与えられていれば、物体が距離 $Ct$ を描くあいだに、その運動の部分 $Ct$ を失う時間が与えられる。そしてそこから、抵抗の重力に対する比が、この失われた運動の部分の、球の重力が同じ時間に生じさせる運動に対する比に等しいとして、抵抗の球の重さに対する比が与えられる。

\textit{系2.}\wsp{}これらを定めるにあたって、私は、流体の粒子がその弾性によって球から最大にはね返されると仮定し、はね返された粒子が球に与える衝撃は、はね返されない場合の2倍であるとした。したがって、粒子が弾性その他のあらゆる反射する力をもたない流体の中では、球形の物体は抵抗が半分になり、したがって、同じ速さの部分を失うあいだに、この問題の作図で述べた距離の2倍遠くまで進むことは明らかである。

\textit{系3.}\wsp{}もし粒子の反射の力が、最大でもなく、まったくないのでもなく、中間のある比をもつなら、抵抗も同様に、この問題の作図と前の系で定めた限界のあいだの中間の比をもつ。

\textit{系4.}\wsp{}遅い物体は、速さの二乗の比から見積もるよりわずかに大きく抵抗を受けるから、これらは任意の長さ $CZ$ を描くあいだに、その運動のうち、全運動に対して $CT$ 対 $CR$ となる部分より大きな部分を失う。

\textit{系5.}\wsp{}きわめて速い物体の抵抗がわかれば、遅い物体の抵抗もわかる。ただし、速さに応じた抵抗の減少の法則が見いだせればである。

\condpagelarge{命題 XXXVII. 問題 IX.}

\textit{与えられた容器から孔を通って流れ出る水の運動を定めること。}

もし容器が水で満たされ、底に孔があけられて水が孔から流れ出るなら、容器は、孔の真上にある部分の重さを除いた、水全体の重さを支えることは明らかである。というのは、もし孔が何かの栓でふさがれていれば、栓は自分の真上にある水の重さを支え、容器の底は残りの水の重さを支えるからである。ところが栓を取り去れば、容器の底は、前と同じ水の圧力、同じ重さで押され、栓が支えていた重さは、もはや支えられないから、水を下降させ、孔から流れ出させる。

ここから、流れ出る水全体の運動は、孔の真上にある水の重さが生じさせうる運動であることがわかる。というのは、水の粒子はそれぞれ、妨げられない限り、自分の重さで下降し、それも一様に加速された運動で下降する。そして、妨げられる限りは、妨げるものを押すからである。この妨げるものは、容器の底か、下を流れ下る水である。したがって、容器の底が支えないその重さの部分は、流れ下る水を押し、それに比例する運動を生じさせる。

そこで、$F$ を孔の面積、$A$ を孔の真上にある水の高さ、$P$ をその重さ、$AF$ をその量とし、$S$ を、与えられた時間 $T$ のあいだに、空虚の中で自由に落ちて描く距離、$V$ を、その時間の終わりに落下によって得た速さとする。すると、その得られた運動 $AF \times V$ は、同じ時間に流れ出る水全体の運動に等しい。水が孔から流れ出るときの速さの速さ $V$ に対する比を $d$ 対 $e$ とする。水は速さ $V$ で距離 $2S$ を描きうるから、流れ出る水は、同じ時間に、自分の速さ $\frac{d}{e}V$ で距離 $\frac{2d}{e}S$ を描きうる。したがって、長さ $\frac{2d}{e}S$、幅が孔と同じ水の柱が、その時間に、流れ下って容器から出ることができる。すなわち柱 $\frac{2d}{e}SF$ である。ゆえに、流れ出る水の量にその速さを掛けて得られる運動 $\frac{2dd}{ee}SFV$、すなわち、その流出の時間に生じるすべての運動は、運動 $AF \times V$ に等しい。これらの等しい運動を $FV$ で割れば、$\frac{2dd}{ee}S$ が $A$ に等しくなる。したがって $dd$ 対 $ee$ は $A$ 対 $2S$ であり、$d$ 対 $e$ は $\frac{1}{2}A$ 対 $S$ の比の平方根である。ゆえに、水が孔から流れ出る速さは、水が落下して時間 $T$ のあいだに距離 $S$ を描いて得る速さに対して、孔の真上にある水の高さが、その高さの2倍と、物体が時間 $T$ で落下して描く距離 $S$ との比例中項に対する比をなす。

そこで、これらの運動を上向きにするならば、水は速さ $V$ で、落ちてきた高さ $S$ まで上るはずであり、高さは（よく知られているように）速さの二乗に比例するから、流れ出る水は高さ $\frac{1}{2}A$ まで上るであろう。したがって、物体が落下して高さ $\frac{1}{2}A$ を描くことができる時間に流れ出る水の量は、孔の真上にある水の柱全体 $AF$ に等しい。

ところで、流れ出る水が、その運動を上向きにすれば、孔の上にある水の高さの半分まで垂直に上るのだから、もし水が容器の側面の管から斜めに流れ出るならば、抵抗のない空間で、放物線を描くことになる。その通径は、管の口の上の容器の中の水の高さであり、その直径は、その口から水平に垂直に引かれ、縦線は管の軸に平行である。

これらはすべて、きわめて微細な流体について理解すべきである。というのは、もし水が粗い部分から成っているなら、水は上に述べた比から見積もるより、ゆっくり流れ出るからである。とくに、流れ出る孔が狭い場合にはそうである。

最後に、水が水平な管から流れ出るなら、容器の底は完全であり、水が流れ出ない場合と同じ、上にある水の圧力をいたるところで受ける。したがって、容器は、流出にかかわらず、水全体の重さを支える。しかし、水が流れ出る側の壁は、水が流れ出ないとしたら支えるはずの圧力のすべてを支えない。というのは、孔があけられたところの圧力が取り除かれるからである。その圧力は、底面が孔に等しく、高さが孔の上の水全体の高さに等しい水の柱の重さに等しい。したがって、もし容器が、振り子のように、長い糸で釘からつるされているなら、水が水平な線に沿って任意の向きに流れ出る場合、容器はつねに垂線から反対の向きに離れる。同じ理由は、湿った火薬を詰めた球の運動にもあてはまる。球は、中身が孔から少しずつ炎となって吹き出すにつれて、炎のある側から離れ、反対の向きに勢いよく運ばれる。

\condpagelarge{命題 XXXVIII. 定理 XXIX.}

\textit{任意のきわめて流動的な媒質の中での、球形の物体の前面での抵抗を定めること。}

\pngright{341.png}{}{}{}

水が円柱形の容器 $ABCD$ から、円柱形の管 $EFGH$ を通って、下の容器 $IKLM$ に流れ込み、そこから容器の縁 $IM$ を越えて流れ出るとする。その縁は上の容器の底 $CD$ と同じ高さにあり、水は管の全体を一様な運動で降りるとする。そして管の中央に球 $P$ を置き、$PR$ を球の上の水の高さ、$SR$ をその容器の底の上の高さとする。球は、管の両側の壁に取り付けた、ごく細い糸 $TV$ で支えられている。すると、上の比例によって、与えられた時間に流れ下る水の量は、それが通る孔の広さに比例することは明らかである。すなわち、球を取り去れば、管の口に比例し、球があれば、球と管との間のまわり全体の空間に比例する。というのは、流れ下る水の速さは（前の命題により）、物体が落下して水の高さ $SR$ の半分を描くことによって得る速さであるから、球を取り去っても、あっても同じだからである。したがって、流れ下る水は、それが通る空間の広さに比例する。たしかに、水が狭い空間を通るのは、広い空間を通るより容易であることはありえない。したがって球があるときの水の速さは、球を取り去ったときより大きくなりえない。それゆえ、球があるとき、流れ出る水の量は、それが通る空間の比から見積もるより大きくはならない。もし水がきわめて微細でかつ流動的な液体でないなら、これが狭い空間を通るのは、粒子の粗さのために、いくぶん遅くなるだろう。だが、ここでは液体がきわめて流動的であると仮定している。

そこで、球が与えられた時間に下降を妨げる水の量は、もし球を取り去れば同じ時間に下降するはずの水の量に対して、球のまわりに描いた円柱の底面が管の口に対するのと同じ比、すなわち球の直径の平方の管の内径の平方に対する比をなす。したがって、球が下降を妨げる水の量は、同じ時間に、その円柱の底面に等しい、容器の底の円形の孔を通って下降しうる水の量に等しく、その下降は、底のどの円形の部分でも、その底面に等しいところで妨げられる。

さて、容器と球が合わせて支える水の重さは、容器の中の水全体の重さから、流れ下る水を加速してその運動を生じさせるのに足りる部分を除いたものである。この部分は、前の命題により、底面が球と管の間の水が流れ下る空間に等しく、高さが、線 $SR$ で表される容器の底の上の水の高さに等しい水の柱の重さに等しい。したがって、容器の底と球は、合わせて、それらの真上にある容器の中の水全体の重さを支える。ところで、容器の底は自分の真上にある水の重さを支えるのだから、残りは、球もまた自分の真上にある水の重さを支える、ということになる。球は、その停滞して何の運動もなく自分の上にある水の重さを支えるのではなく、流れ下る水に抵抗して、それだけの重さの効果を妨げるのである。したがって、球は、その重さに等しい、流れ下る水の力を支える。実際、球は、もし球を取り去れば、その重さが正確に生じさせるであろう水の量の下降と流出を妨げるのである。水は、その下降が妨げられるかぎり、あらゆる妨げるものを自分の重さで押す。したがって、妨げるものは、水の下降を妨げるかぎり、その下降を生じさせるはずの重さに等しい力を支える。ところで球は、自分の真上にある水の柱の重さが生じさせうる量の水の下降を妨げる。したがって、流れ下る水の力を、その重さに等しいだけ支える。作用と反作用は、運動の法則IIIにより、互いに等しく、反対の向きに向けられる。球が下降する水に及ぼす作用は、その下降を妨げるために、上向きに向けられ、妨げられる下降の運動に比例し、それを取り去るのにちょうど足りる。したがって、水の球に対する反対の作用は、同じ運動を取り去る、あるいは生じさせうる力に等しい。すなわち、球の真上にあって高さが $RS$ である水の柱の重さに等しい。

いま、管の上の口をふさいで、水が降りられないようにするなら、球は、管と下の容器 $IKLM$ に停滞する水の重さによって、四方から押される。しかし、その圧力があっても、球が水と同じ比重であれば、静止する。その圧力は球をどちらの向きにも押さない。したがって、管が開かれて、上の容器から水が降りるとき、球を下向きに押す力のすべては、その水の下降から生じ、したがって、高さが $RS$ で直径が球と同じ水の柱の重さに等しい。ところで、この重さは、球のまわりに描いた円柱の底面に等しい孔から、球を取り去れば、与えられた任意の量の水が流れ出ることができる時間に、その水の運動のすべてを生じさせるのに足りる。したがって、水が円柱の中を一様に降りて、球の直径の3分の2を描く時間に、球に等しい水の運動のすべてを生じさせるのに足りる。というのは、球の幅をもち、高さが3分の2である水の円柱は、球に等しいからである。したがって、流れる水が静止した球に及ぼす衝撃は、水が流れて球の直径の3分の2を描く時間に、もしそれが一様につづくならば、球に等しい流体の部分の運動のすべてを生じさせるであろう。

ところで、管の中の水について証明したことは、任意の流れる水についても、その中に静止している任意の球が押される場合にあてはまる。そして水について証明したことは、あらゆるきわめて微細な流体でも成り立つ。これらすべてについて、同じ論証が成り立つ。

さて、法則の系5により、流体の球に対する力は、球が静止して流体が一様な速さで動く場合も、流体が静止して球が同じ速さで反対の向きに進む場合も、同じである。したがって、任意のきわめて流動的な媒質の中を一様に進む球の抵抗は、球が自分の直径の3分の2を描く時間に、球と同じ大きさで媒質と同じ密度の物体に一様に加えられて、球が進みながら自分の直径の3分の2を描く時間に、その物体の中に球の速さを生じさせうる力に等しい。これが、球の前に進む表面部分での抵抗である。\QED

\textit{系1.}\wsp{}もし球形の固体が、自分と同じ密度のきわめて微細な流体の中を自由に動き、動くあいだ、静止しているときと同じ力で後ろから押されるならば、その抵抗は、命題XXXVIの系2で述べたものになる。したがって、計算してみれば、固体は、自分の直径の長さを進んで描くより前に、運動の半分を失うことがわかる。もし動くあいだ、後ろからの押す力が小さければ、いっそう遅くなる。逆に、より大きく押されれば、遅くなり方は小さい。

\textit{系2.}\wsp{}したがって、投射体の抵抗が流体の部分を無限に分割することで無限に小さくなる、と信じる人々は誤っている。もし流体がたいへん粗ければ、その部分の分割によって、抵抗はいくらか小さくなる。しかし、適当な流動性の度合い（たとえば空気や水や水銀の流動性ほど）に達してしまえば、固体の前面での抵抗は、部分をさらに分割しても、あまり小さくならない。実際、抵抗は、前の系で定めた限界より小さくなることはけっしてない。

\textit{系3.}\wsp{}したがって、投射体が運動の目立った減少なしにきわめて遠くまで進む媒質は、きわめて流動的であるだけでなく、その中を動く物体よりもはるかに希薄である。ただし、きわめて流動的なあらゆる媒質は、投射体の後ろの部分に絶え間なく衝撃を与えて、前の部分で妨げ抵抗する分だけ、その運動を進めるのだと言う人があれば、別である。そして、投射体が媒質に与える運動のうち、ある部分が、媒質が円を描いて運ばれて物体の後ろに返されることは、本当らしい。というのは、私が行ったいくつかの実験でも、十分に圧縮された流体では、ある部分が返されることがわかったからである。しかし、どんな場合にもすべてが返されるというのは、道理にも合わないように見え、私がこれまでに試した実験ともよく合わない。というのは、どんなに微細でも、密度のある流体の、固体を動かし、またこれに抵抗する力はきわめて大きいということ、そしてその力の大きさが実験でどう定められるかは、つづく二つの命題で、いっそう十分に明らかになるだろう。

\condpagelarge{補助定理 IV.}

\textit{球形の容器が、静止した均質な流体で満たされていて、加えられた力で直線に動き、つねに加速される前進運動でつづき、そのあいだに円を描いて動かないならば、中に入っている流体の部分は、容器の運動に等しく加わって、互いに静止している。同じことは、任意の形の容器でも成り立つ。事柄は明らかで、証明を要しない。}

\condpagelarge{命題 XXXIX. 定理 XXX.}

\textit{強まる風のように加速された運動で進み、その部分が互いに静止している流体は、同じ密度の浮かんでいる物体をすべてさらい、それらを同じ速さで一緒に運ぶ。}

実際、上の補助定理により、剛体の球形の容器が、静止した均質な流体で満たされていて、少しずつ与えられた運動で進むならば、容器の運動に加わる流体の部分は、すべて互いに静止している。したがって、もし流体のいくつかの部分が凍ったとしても、これらは残りの部分に対して静止しつづける。というのは、部分がすべて互いに静止しているのだから、それらが流体であっても、そのいくつかが固まっていても、同じことだからである。ゆえに、もし容器が外から加えられた何らかの力で動き、その運動を流体に与えるなら、流体はまた、その運動を、自分自身の凍った部分に与え、それらを一緒にさらう。ところが、これらの凍った部分は、流体と同じ密度の固体である。そして、流体が動く容器に閉じ込められている場合も、自由な空間の中で風のように吹く場合も、理由は同じである。ゆえに、加速された前進運動で運ばれ、その部分が互いに静止しているあらゆる流体は、はじめ静止していた、同じ密度の、中に含まれたどんな固体でも、自分と一緒にさらって、一緒に動くように強いる。\QED

\condpagelarge{命題 XL. 問題 X.}

\textit{密度の与えられたきわめて流動的な媒質の中での、球形の固体の抵抗を求めること。}

与えられた任意の流体の中で、種類の与えられた固体の、大きさを限りなく増したときの最終的な抵抗を求める。ついで、次のように言う。その固体が失う運動が、進んで自分の半径の長さを描く時間に、その最初の全運動に対して占める比は、同じ流体がきわめて微細になったとき、任意の与えられた固体が、自分の直径の長さを描くあいだに失う運動の、その最初の全運動に対する比に、ほぼ等しい。というのは、もし微細にされた流体の最小の粒子が、その中を動く与えられた固体に対して、微細にされていない流体の同じ数の最小の粒子が、大きくなった固体に対してもつのと同じ比と位置をもち、どちらの流体の粒子もきわめてなめらかで、遠心力・向心力をまったくもたず、固体がそれぞれの流体の中で、任意の比例する時間に相似に動き始めるなら、これらは相似に動きつづけ、したがって、自分の半径に等しい距離を描く時間に、全体に比例する運動の部分を失う。それは、微細にされた媒質の部分が小さくなり、微細にされていない媒質の中を動く固体の大きさが限りなく大きくなっても、そうである。ゆえに、微細にされていない媒質の中での大きくなった固体の抵抗から、上の比例によって、微細にされた媒質の中での大きくなっていない固体の抵抗が与えられる。\QED

もし粒子がきわめてなめらかでないなら、どちらの流体でも同じようになめらかである、と仮定しなければならない。そうすれば、なめらかさの不足から生じる抵抗は、どちらでも同じように増し、命題はなおも成り立つ。

\textit{系1.}\wsp{}したがって、球形の固体の大きさを増すと、その抵抗が二乗の比で増すなら、与えられた球形の固体の抵抗は、流体の粒子の大きさを小さくしても、まったく小さくならない。

\textit{系2.}\wsp{}もし、球形の固体を大きくするとき、抵抗が直径の二乗より小さい比で増すなら、流体の粒子を小さくすれば、抵抗は、増加した抵抗が直径の二乗の比から不足する比で小さくなる。

\textit{系3.}\wsp{}ここから、与えられた固体の抵抗は、流体の部分を分割しても、あまり小さくできないことが明らかである。というのは、大きくなった固体の抵抗は、その固体が動きながら押しのけ、自分が占めた場所から押し出す、抵抗する流体の物質の量に、ほぼ比例するはずだからである。すなわち、固体が動く円柱形の空間に、したがって、固体の半径の二乗に、ほぼ比例する。

\textit{系4.}\wsp{}したがって、二つの流体が与えられ、そのうち一方が他方より抵抗する力でずっと勝るならば、抵抗の小さいほうが他方より希薄である。そして、すべての流体の抵抗する力は、それらの密度に、ほぼ比例する。とくに、固体が大きく、速く動き、流体の圧縮が等しい場合はそうである。
\condpagelarge{一般的な注}

以上に証明したことを、私は次のようにして試した。重さがローマのオンスで $57\frac{7}{22}$、直径がロンドンのインチで $6\frac{7}{8}$ の木の球を作り、細い糸で、じゅうぶん丈夫な鉤からつるした。鉤と球の振動の中心との距離が $10\frac{1}{2}$ フィートになるようにした。糸の上に、つり点から $10$ フィート $1$ インチ離れた点に印をつけ、その点と向かい合わせに、インチで目盛りをつけた定規を置いて、それによって振り子の描く弧の長さを読み取れるようにした。そして、球が自分の運動の4分の1を失うまでの振動の回数を数えた。

振り子を垂直から2インチの距離まで引いてから放し、下降のあいだに2インチの弧を、下降とそれに続く上昇から成る最初の振動全体で、ほぼ4インチの弧を描くようにしたとき、同じ振り子は164回の振動で運動の8分の1を失い、最後の上昇では $1$ インチと $\frac{3}{4}$ インチの弧を描いた。最初の下降で4インチの弧を描いたときは、121回の振動で運動の8分の1を失い、最後の上昇で描いた弧は $3\frac{1}{2}$ インチであった。最初の下降で、8, 16, 32 または 64 インチの弧を描いたときは、それぞれ 69, $35\frac{1}{2}$, $18\frac{1}{2}$, $9\frac{2}{3}$ 回の振動で運動の8分の1を失った。したがって、最初の下降と最後の上昇で描かれた弧の差は、第1, 2, 3, 4, 5, 6の場合に、それぞれ $\frac{1}{4}$, $\frac{1}{2}$, 1, 2, 4, 8 インチであった。これらの差を、それぞれの場合の振動の回数で割ると、弧 $3\frac{3}{4}$, $7\frac{1}{2}$, 15, 30, 60, 120 インチを描く平均の1回の振動で、下降とそれに続く上昇で描かれた弧の差は、それぞれ $\frac{1}{656}$, $\frac{1}{242}$, $\frac{1}{69}$, $\frac{4}{71}$, $\frac{8}{37}$, $\frac{24}{29}$ インチになる。これらは、大きな振動ではほぼ描かれた弧の二乗に比例し、小さな振動では、その比よりわずかに大きい。したがって（この編の命題XXXIの系2により）、球の抵抗は、速く動くときにはほぼ速さの二乗に比例し、遅く動くときには、その比よりわずかに大きい。これは、命題XXXIIの系で証明したとおりである。

いま、$V$ を任意の振動における最大の速さとし、$A$, $B$, $C$ を与えられた量として、弧の差が $AV + BV^{\frac{3}{2}} + CV^2$ であると仮定しよう。上の六つの場合の最大の速さは、半分の弧 $1\frac{7}{8}$, $3\frac{3}{4}$, $7\frac{1}{2}$, 15, 30, 60 の弦に比例し、したがってほぼ弧そのものに、すなわち数 $\frac{1}{2}$, 1, 2, 4, 8, 16 に比例する。そこで、第2、第4、第6の場合に、$V$ のかわりに数 1, 4, 16 を書くと、弧の差は、第2の場合に $\frac{1}{242}$ に等しく $A + B + C$、第4の場合に $\frac{2}{35\frac{1}{2}}$ に等しく $4A + 8B + 16C$、第6の場合に $\frac{8}{9\frac{2}{3}}$ に等しく $16A + 64B + 256C$ となる。これらの式から $A$, $B$, $C$ を決めれば、任意に与えられた速さに対する弧の差を求める規則が得られる。

ところで、最大の速さは、サイクロイドでは振動で描かれる弧に比例し、円では、その弧の半分の弦に比例する。したがって、弧が等しければ、サイクロイドでのほうが円でよりも、弧と弦との比だけ大きい。一方、時間は円のほうがサイクロイドより、速さの逆の比だけ長い。そこで、円での抵抗からトロコイド（サイクロイド）での抵抗を求めるには、速さが単純な比で増すために、抵抗をほぼ弧対弦の二乗の比で増し、時間（すなわち、前に述べた弧の差を生む抵抗が働く持続）が同じ比で縮むために、弦対弧の比で減らさなければならない。すなわち、（比を合わせれば）抵抗を、ほぼ弧対弦の比で増さなければならない。この比は、第2の場合に $6283$ 対 $6279$、第4の場合に $12566$ 対 $12533$、第6の場合に $25132$ 対 $24869$ である。そこから、抵抗 $\frac{1}{242}$, $\frac{2}{35\frac{1}{2}}$, $\frac{8}{9\frac{2}{3}}$ は、$\frac{6283}{6279 \times 242}$, $\frac{25132}{12533 \times 35\frac{1}{2}}$, $\frac{201056}{24869 \times 9\frac{2}{3}}$、すなわち小数で $0.004135$, $0.056486$, $0.8363$ になる。ここから、式 $A + B + C = 0.004135$, $4A + 8B + 16C = 0.05648$, $16A + 64B + 256C = 0.8363$ が得られ、これらから、項を適切に比べて解析的に整理すれば、$A = 0.0002097$, $B = 0.0008955$, $C = 0.0030298$ となる。したがって弧の差は $0.0002097V + 0.0008955V^{\frac{3}{2}} + 0.0030298V^2$ に比例する。そして（この編の命題XXXの系により）、弧が振動で描かれる場所での球の抵抗は、速さが $V$ のとき、球の重さに対して、$\frac{7}{11}AV + \frac{16}{23}BV^{\frac{3}{2}} + \frac{3}{4}CV^2$ 対振り子の長さ、の比をなすから、$A$, $B$, $C$ に見いだした数を代入すれば、球の抵抗の重さに対する比は、$0.0001334V + 0.000623V^{\frac{3}{2}} + 0.00227235V^2$ 対、つり点の中心と定規との間の長さ、すなわち 121 インチとなる。ゆえに、$V$ が第2の場合に 1、第4の場合に 4、第6の場合に 16 であるから、抵抗の球の重さに対する比は、第2の場合に $0.003029$ 対 121、第4の場合に $0.042875$ 対 121、第6の場合に $0.63013$ 対 121 となる。

第6の場合に、糸に印をつけた点が描いた弧は $120 - \frac{8}{9\frac{2}{3}}$、すなわち $119\frac{5}{29}$ インチであった。そして、半径が 121 インチ、つり点と球の中心との間の振り子の長さが 126 インチであったから、球の中心が描いた弧は $124\frac{3}{31}$ インチであった。振動する物体の最大の速さは、空気の抵抗のために、描かれた弧の最低点には来ず、弧全体のほぼ中央にくるから、これは、球が下降全体を抵抗のない媒質の中で、その弧の半分 $62\frac{3}{62}$ インチを描いたのと、ほぼ同じである。それも、上で振り子の運動を還元したサイクロイドの上でのことである。したがって、その速さは、球が垂直に落下し、その落下でその弧の正矢に等しい高さを描くことによって得られる速さに等しい。ところで、サイクロイドでは、その正矢は、弧 $62\frac{3}{62}$ に対して、同じ弧が振り子の長さの2倍 $252$ に対するのと同じ比をなし、したがって $15.278$ インチに等しい。ゆえに、その速さは、物体が落下して $15.278$ インチの距離を描くことによって得られる速さにほかならない。ところで、物体が1秒の時間に落下して描く距離は（ホイヘンスが振り子の実験から定めたように）パリのフィートで $15\frac{1}{12}$、すなわちイギリスのフィートで $16\frac{11}{24}$、すなわち $197\frac{1}{2}$ インチであり、時間は距離の平方根の比をなす。したがって球は、時間 $0\second 16\third 38\fourth$ のあいだ落下して $15.278$ インチを描き、前述の速さを得る。そしてその同じ速さを一様に続ければ、同じ時間に、2倍の長さ $30.556$ インチを描く。このような速さで、球は、その重さに対して $0.63013$ 対 $121$、あるいは（抵抗のうち速さの二乗に比例する部分だけを見れば）$0.58172$ 対 $121$ の抵抗を受ける。

ところで、静水力学の実験によって、この木の球の重さが、同じ大きさの水の球の重さに対して 55 対 97 であることがわかった。したがって、121 対 $213.4$ が同じ比であるから、前述の速さで進む水の球の抵抗は、その重さに対して $0.58172$ 対 $213.4$、すなわち 1 対 $366\frac{5}{6}$ になる。ゆえに、水の球の重さが、球が一様に続く速さで $30.556$ インチ〔原文は pedum（フィート）とあるが、前の計算 $2 \times 15.278$ インチからインチと解した〕の長さを描く時間のうちに、落下する球にその速さのすべてを生じさせうるとすれば、一様に続く抵抗の力は、1 対 $366\frac{5}{6}$ の比で小さい速さを、すなわち速さ全体の $\frac{1}{366\frac{5}{6}}$ を取り去ることができる、ということは明らかである。したがって、球がその速さを一様に続けて、自分の半径、すなわち $3\frac{7}{16}$ インチの長さを描く時間のうちに、運動の $\frac{1}{3262}$ の部分を失う。

私はまた、振り子が運動の4分の1を失うまでの振動の回数も数えた。次の表で、上段の数は最初の下降で描かれた弧の長さをインチとインチの端数で表し、中段の数は最後の上昇で描かれた弧の長さを表し、下段には振動の回数が立つ。この実験は、運動の8分の1だけを失わせた場合よりも正確なものとして記した。計算は、望む人が試みるとよい。

\begin{center}
\begin{tabular}{lcccccc}
最初の下降 & 2 & 4 & 8 & 16 & 32 & 64 \\
最後の上昇 & $1\frac{1}{2}$ & 3 & 6 & 12 & 24 & 48 \\
振動の回数 & 374 & 272 & $162\frac{1}{2}$ & $83\frac{1}{3}$ & $41\frac{2}{3}$ & $22\frac{2}{3}$ \\
\end{tabular}
\end{center}

そののち、直径2インチ、重さがローマのオンスで $26\frac{1}{4}$ の鉛の球を、同じ糸で、球の中心とつり点との間隔が $10\frac{1}{2}$ フィートになるようにつるし、与えられた運動の一部が失われるまでの振動の回数を数えた。続く二つの表のうち、前の表は運動全体の8分の1がなくなるまでの振動の回数を示し、後の表は、同じ運動の4分の1が失われるまでの振動の回数を示している。

\begin{center}
\begin{tabular}{lccccccc}
最初の下降 & 1 & 2 & 4 & 8 & 16 & 32 & 64 \\
最後の上昇 & $\frac{7}{8}$ & $\frac{7}{4}$ & $3\frac{1}{2}$ & 7 & 14 & 28 & 56 \\
振動の回数 & 226 & 228 & 193 & 140 & $90\frac{1}{2}$ & 53 & 30 \\
 & & & & & & & \\
最初の下降 & 1 & 2 & 4 & 8 & 16 & 32 & 64 \\
最後の上昇 & $\frac{3}{4}$ & $1\frac{1}{2}$ & 3 & 6 & 12 & 24 & 48 \\
振動の回数 & 510 & 518 & 420 & 318 & 204 & 121 & 70 \\
\end{tabular}
\end{center}

前の表から、第3、第5、第7の観測を選び、これらの観測での最大の速さを、それぞれ数 1, 4, 16、一般には上と同じく量 $V$ で表すと、最初の観測では $\frac{2}{193} = A + B + C$、二番目では $\frac{2}{90\frac{1}{2}} = 4A + 8B + 16C$、三番目では $\frac{8}{30} = 16A + 64B + 256C$ となる。これらの式は、上に示した整理によって、$A = 0.000145$, $B = 0.000247$, $C = 0.0009$ を与える。そこから、速さ $V$ で動く球の抵抗は、その重さ $26\frac{1}{4}$ オンスに対して、$0.000923V + 0.000172V^{\frac{3}{2}} + 0.000675V^2$ 対、振り子の長さ 121 インチ、の比をなすことが得られる。そして、抵抗のうち速さの二乗に比例する部分だけを見れば、これは球の重さに対して $0.000675V^2$ 対 121 インチとなる。ところで、この部分の抵抗は、第1の実験では、重さ $57\frac{7}{22}$ オンスの木の球に対して $0.00227235V^2$ 対 121 であった。ここから、木の球の抵抗の鉛の球の抵抗に対する比は（速さが等しいとして）$57\frac{7}{22} \times 0.00227235$ 対 $26\frac{1}{4} \times 0.000675$、すなわち $130309$ 対 $17719$、あるいは $7\frac{1}{3}$ 対 1 となる。二つの球の直径は $6\frac{7}{8}$ と 2 インチであり、これらの平方は、互いにほぼ $47\frac{1}{4}$ 対 4、すなわち $11\frac{13}{16}$ 対 1 である。ゆえに、等速の球の抵抗は、直径の二乗の比より小さい比をなしていた。しかし、私たちはまだ糸の抵抗を考えていなかった。これは確かにきわめて大きく、振り子について求めた抵抗から引かなければならない。私はこれを正確に定めることはできなかったが、小さいほうの振り子の全抵抗の3分の1より大きいことがわかった。そこから、糸の抵抗を除けば、球の抵抗は、直径の二乗の比にほぼ比例することを知った〔原文には dimidiata（平方根の比）とあるが、直後の数値 $10\frac{1}{2}$ 対 1 が $11\frac{13}{16}$ 対 1 に近いと述べているから、二乗の比の意と解した〕。というのは、比 $7\frac{1}{3} - \frac{1}{3}$ 対 $1 - \frac{1}{3}$、すなわち 7 対 $\frac{2}{3}$、あるいは $10\frac{1}{2}$ 対 1 は、直径の二乗の比 $11\frac{13}{16}$ 対 1 から遠くないからである。

糸の抵抗は、大きな球ではそれほど重要でないから、直径 $18\frac{1}{4}$ インチの球でも実験を試みた。つり点と振動の中心との間の振り子の長さは $122\frac{3}{4}$ インチ、つり点と糸の結び目との間は $109\frac{1}{2}$ インチであった。最初の振り子の下降で結び目が描いた弧は 32 インチ、5回の振動ののち、最後の上昇で同じ結び目が描いた弧は 28 インチであった。弧の和、すなわち平均の振動で描かれる全弧は 30 インチ〔原文のまま。$32 + 28 = 60$ であり、直後の文でも 60 インチを用いている〕、弧の差は 4 インチである。その10分の1、すなわち平均の振動における下降と上昇の差は $\frac{2}{5}$ インチである。半径 $109\frac{1}{2}$ 対半径 $122\frac{1}{2}$ が、平均の振動で結び目が描く全弧 60 インチ対、平均の振動で球の中心が描く全弧 $67\frac{1}{8}$ に等しいように、差 $\frac{2}{5}$ は新しい差 $0.4475$ に対応する。もし振り子の長さが、描かれる弧の長さは変わらずに、126 対 $122\frac{1}{2}$ の比で増すならば、その速さはこの比の平方根の比で減り、下降とそれに続く上昇で描かれる弧の差 $0.4475$ は、速さの比で減って $0.4412$ になる。ついで、描かれる弧が $67\frac{1}{8}$ 対 $124\frac{3}{31}$ の比で増すならば、この差 $0.4412$ はその二乗の比で増して $1.509$ になる。これは、振り子の抵抗が速さの二乗に比例するという仮定の上でのことである。したがって、もし振り子が全弧 $124\frac{3}{31}$ インチを描き、そのつり点と振動の中心との間の長さが 126 インチならば、下降とそれに続く上昇で描かれる弧の差は $1.509$ インチになるはずである。この差に、振り子の球の重さ（208オンス）を掛けると $313.9$ になる。

また、先の木の球で作った振り子が、つり点から 126 インチ離れた振動の中心で、全弧 $124\frac{3}{31}$ インチを描いたとき、下降と上昇で描かれた弧の差は、$\frac{126}{121}$ に $\frac{8}{9\frac{2}{3}}$ を掛けたもの、すなわち $\frac{25}{29}$ であった。これに、球の重さ $57\frac{7}{22}$ オンスを掛けると $48.55$ になる。私がこれらの差を球の重さに掛けたのは、球の抵抗を求めるためである。というのは、差は抵抗から生じ、抵抗に正比例し、重さに反比例するからである。したがって、抵抗は数 $313.9$ と $48.55$ に比例する。ところで、小さいほうの球の抵抗のうち、速さの二乗に比例する部分は、全抵抗に対して $0.58172$ 対 $0.63013$、すなわち $44.4$ 対 $48.55$ であり、大きいほうの球の抵抗のうちその部分は、全抵抗にほとんど等しい。したがって、これらの部分は、ほぼ $313.9$ 対 $44.4$、すなわち $7.07$ 対 1 の比をなす。ところで、球の直径は $10\frac{3}{4}$ と $6\frac{7}{8}$ であり〔前の文では $18\frac{1}{4}$ とあるが、原文のまま〕、これらの平方 $351\frac{1}{2}$ と $47\frac{17}{64}$ は $7.438$ 対 1 であって、球の抵抗の比 $7.07$ 対 1 にほぼ等しい。この比の差は、糸の抵抗から生じうる差より大きくない。ゆえに、抵抗のうち速さの平方に比例する部分は、（球が等しいとき）速さの平方に比例し、（速さが等しいとき）球の直径の平方にも比例する。したがって（この編の命題XLの系により）、大きくて速い球が空気の中を動くときに受ける抵抗は、空気を無限に分割し微細にしても、あまり減らすことはできない。ゆえに、物体がはるかに小さい抵抗しか受けない媒質は、すべて空気より希薄である。

ところで、これらの実験に使った球のうち、最大のものは完全な球ではなかった。そのため、ここに示した計算では、簡潔にするために細かい点をいくつか無視した。十分に正確でない実験に対して、正確な計算にはまったくこだわらない。そこで、（真空の証明はこれらに依存しているので）、もっと多くの、もっと大きな、もっと正確な球で実験が行われることを望む。もし球を等比数列にとって、たとえば直径が 4, 8, 16, 32 インチであるようにすれば、実験の進み方から、さらに大きな球でどうなるはずかが推定されるであろう。

さて、さまざまな流体の抵抗を互いに比べるために、私は次のことを試みた。長さ4フィート、幅と高さがそれぞれ1フィートの木の箱を用意した。このふたをとって泉の水を満たし、水に沈めた振り子が、水の中ほどで振動するようにした。重さ $166\frac{1}{6}$ オンス、直径 $3\frac{5}{8}$ インチの鉛の球は、次の表に記したように動いた。つり点から糸に印をつけた点までの振り子の長さは 126 インチ、振動の中心までは $134\frac{1}{8}$ インチであった。

\begin{center}
\begin{tabular}{p{5.5cm}ccccccccc}
最初の下降で糸の印の点が描いた弧（インチ） & 64 & 32 & 16 & 8 & 4 & 2 & 1 & $\frac{1}{2}$ & $\frac{1}{4}$ \\
最後の上昇で描いた弧（インチ） & 48 & 24 & 12 & 6 & 3 & $1\frac{1}{2}$ & $\frac{3}{4}$ & $\frac{3}{8}$ & $\frac{3}{16}$ \\
失われた運動に比例する弧の差（インチ） & 16 & 8 & 4 & 2 & 1 & $\frac{1}{2}$ & $\frac{1}{4}$ & $\frac{1}{8}$ & $\frac{1}{16}$ \\
水中の振動の回数 & & & $\frac{29}{60}$ & $1\frac{1}{5}$ & 3 & 7 & $11\frac{1}{4}$ & $12\frac{2}{3}$ & $13\frac{1}{3}$ \\
空気中の振動の回数 & $85\frac{1}{2}$ & & 287 & 535 & & & & & \\
\end{tabular}
\end{center}

第4列の実験では、等しい運動が、空気中では 535 回、水中では $1\frac{1}{5}$ 回の振動で失われた。ただし、空気中の振動は、水中のものよりわずかに速く、44 対 41 の比であった。というのは、水中で $14\frac{2}{3}$ 回、空気中で $13\frac{2}{3}$ 回の振動が同時に行われたからである。したがって、水中の振動を、二つの媒質での振り子の運動が等速になるように、この比で速めたとすれば、同じ運動が失われる水中の振動の回数 $1\frac{1}{5}$ は、（抵抗がこの比の二乗で増し、時間が同じ比の単純な比で縮むから）同じ比 44 対 41 で減って、$1\frac{1}{5} \times \frac{41}{44}$、すなわち $\frac{123}{110}$ になる。したがって、振り子の速さが等しいとき、等しい運動は、空気中では 535 回の振動で、水中では $\frac{123}{110}$ 回の振動で失われる。ゆえに、振り子の水中での抵抗は、空気中での抵抗に対して 535 対 $\frac{123}{110}$ である。これが、第4列の場合の全抵抗の比である。

いま、$AV + CV^2$ を、速さ $V$ で動く球の空気中での抵抗とする。そして、第4列の場合の最大の速さは、第1列の場合のそれに対して 1 対 8 であり、第4列の場合の抵抗は、第1列の場合の抵抗に対して、これらの場合の弧の差を振動の回数で割ったものの比、すなわち $\frac{2}{535}$ 対 $\frac{16}{85\frac{1}{2}}$、すなわち $85\frac{1}{2}$ 対 $4280$ であるから、これらの場合に、速さとして 1 と 8、抵抗として $85\frac{1}{2}$ と $4280$ を書けば、$A + C = 85\frac{1}{2}$ および $8A + 64C = 4280$、すなわち $A + 8C = 535$ となり、方程式の整理から $7C = 449\frac{1}{2}$、$C = 64\frac{3}{14}$、$A = 21\frac{2}{7}$ が得られる。したがって抵抗はほぼ $21\frac{2}{7}V + 64\frac{3}{14}V^2$ に比例する。ゆえに、速さが 1 であった第4列の場合には、全抵抗は、速さの平方に比例するその部分に対して、$21\frac{2}{7} + 64\frac{3}{14}$、すなわち $85\frac{1}{2}$ 対 $64\frac{3}{14}$ の比をなす。それゆえ、水中での振り子の抵抗は、空気中の抵抗のうち速さの平方に比例する部分（速い運動ではこれだけを考えればよい）に対して、$85\frac{1}{2}$ 対 $64\frac{3}{14}$ と 535 対 $\frac{123}{110}$ との積の比、すなわち 637 対 1 になる。もし、水中で振動する振り子の糸全体が沈められていたならば、その抵抗はいっそう大きかったであろう。そこで、空気中で振動する振り子の抵抗のうち、速さの平方に比例する部分（速い物体ではこれだけを考えればよい）は、同じ速さで水中で振動する同じ振り子全体の抵抗に対して、およそ 800 または 900 対 1、すなわち水の密度の空気の密度に対する比に、ほぼ等しい。

この計算では、水中の振り子の抵抗のうち、速さの平方に比例する部分だけをとるべきであるが、（おそらく不思議に思われるだろうが）水中では、抵抗が速さの二乗より大きい比で増していた。その原因を調べたところ、箱が振り子の球の大きさに対して狭すぎて、道を譲る水の運動を、その狭さのためにあまりに妨げていたことにつきあたった。実際、直径 1 インチの振り子の球を沈めたときには、抵抗は、ほぼ速さの二乗の比で増した。私はこれを、二つの球から成る振り子をつくって試みた。そのうち、下の小さいほうの球は水中で振動し、上の大きいほうの球は、水面のすぐ上で糸に取り付けられていて、空気中で振動しながら振り子の運動を助け、それを長持ちさせるようにした。この方法で行った実験は、次の表に記したとおりである。

\begin{center}
\begin{tabular}{p{5.5cm}ccccccc}
最初の下降で描かれた弧 & 16 & 8 & 4 & 2 & 1 & $\frac{1}{2}$ & $\frac{1}{4}$ \\
最後の上昇で描かれた弧 & 12 & 6 & 3 & $1\frac{1}{2}$ & $\frac{3}{4}$ & $\frac{3}{8}$ & $\frac{3}{16}$ \\
失われた運動に比例する弧の差 & 4 & 2 & 1 & $\frac{1}{2}$ & $\frac{1}{4}$ & $\frac{1}{8}$ & $\frac{1}{16}$ \\
振動の回数 & $3\frac{3}{8}$ & $6\frac{1}{2}$ & $12\frac{1}{12}$ & $21\frac{1}{5}$ & 34 & 53 & $62\frac{1}{5}$ \\
\end{tabular}
\end{center}

ここでは、抵抗が速さの二乗より大きい比で増すことは、まったくない。そして、もし箱が振り子に応じて大きくされるならば、大きな振り子でも同じことが起こると考えられる。しかし、速さが段階的に限りなく増すならば、抵抗は、空気中でも水中でも、最後にはわずかに二乗より大きい比で増さなければならない。というのは、ここに記した実験での抵抗は、この編の命題XXXVIとXXXVIIIで、きわめて速い物体について証明した比から見積もるより、小さいからである。実際、きわめて速い物体は、背後に空虚な空間を残すから、その前の部分で感じる抵抗は、後ろの部分での媒質の圧力によって、少しも減らされない。

媒質の抵抗を互いに比べるために、私は鉄の振り子を水銀の中で振動させてもみた。鉄の糸の長さはおよそ3フィート、振り子の球の直径はおよそ $\frac{1}{3}$ インチであった。糸の、水銀のすぐ上のところには、振り子の運動をいっそう長く続けさせるために、十分に大きな別の鉛の球をつけた。ついで、およそ3ポンドの水銀を入れられる小さな器に、水銀と普通の水とを交互に入れて、振り子が二つの流体の中で次々に振動するようにして、抵抗の比を求めた。その結果、水銀の抵抗の水の抵抗に対する比は、およそ 13 または 14 対 1、すなわち水銀の密度の水の密度に対する比にほぼ等しかった。やや大きめの振り子の球、たとえば直径がおよそ $\frac{1}{2}$ か $\frac{2}{3}$ インチの球を使ったときは、水銀の抵抗の水の抵抗に対する比は、およそ 12 または 10 対 1 となった。しかし、前の実験のほうが信頼できる。それは、後の実験では、器が、沈めた球の大きさに対して狭すぎたからである。球を大きくすれば、器も大きくしなければならない。私は、この種の実験を、もっと大きな器で、また溶かした金属の液体や、そのほかの熱い液体・冷たい液体で、繰り返すつもりでいたが、すべてを試みる暇はない。ただ、すでに記したことから、速く動く物体の抵抗は、それが動く流体の密度に、ほぼ比例することは十分に明らかである。正確に比例するとは言わない。というのは、同じ密度なら、粘りけのある流体は、さらさらした流体より、まちがいなく大きな抵抗を及ぼすからである。たとえば、冷たい油は熱い油より、熱い油は雨水より、水は酒精より大きい。しかし、感覚で見て十分に流動的な液体、たとえば、空気、真水や塩水、ぶどう酒の酒精、テレビン油や塩類の酒精、蒸留して澱を除き熱した油、硫酸油、水銀、溶けた金属、そのほか、容器の中でかき回されると与えられた運動を長く保ち、自由に注ぎ出されると流れ落ちるあいだに滴に分かれるほど流動的な液体では、私は、いま示した規則がかなり正確に成り立つことを疑わない。とくに、実験を、もっと大きくて、もっと速く動かした振り子の物体で行った場合にはそうである。

ゆえに、水の球が空気中を動くとき、その半径の長さを描くあいだに運動の $\frac{1}{3261}$ を取り去られるほどの抵抗を受けること（すでに先に示したとおり）、そして空気の密度の水の密度に対する比が、およそ 800 または 850 対 1 であることから、次の規則が一般に成り立つことになる。任意の球形の物体が、十分に流動的な任意の媒質の中を動くとき、抵抗のうち速さの二乗に比例する部分だけを見れば、この部分は、同じ物体が、その運動を一様に続けて自分の半径の2倍の長さを描くあいだに、物体の運動全体を取り去りうる、あるいは生じさせうる力に対して、媒質の密度の物体の密度に対する比に、ほぼ等しい。したがって、この抵抗は、この編の命題XXXVIIIの系1で与えた法則から見積もるより、およそ3倍大きい。そのため、球の前の部分が動きながら媒質に与える運動全体の、ほぼ3分の2は、媒質が円を描いてもどり、球が背後に空虚なまま残すはずの空間に流れ込むことによって、球の後ろの部分に返される。それゆえ、球の速さが、媒質がその空間にそれほど速く流れ込めなくなって、球の背後にいつも多少の空虚が残るまでに増すならば、抵抗は、いま述べた一般の規則から見積もるより、結局、およそ3倍大きくなる。

ここまで私たちは、振動する振り子の実験を用いてきた。それは、その運動のほうが、より容易に、より正確に観察し測定できるからである。しかし、円を描く運動をさせた振り子、すなわち元にもどって円を描く振り子の運動も、一様であって、与えられた速さに応じた抵抗を調べるのにはるかに適していると思われるので、これも試みた。すなわち、円運動をする振り子を12回転させて、最初の回転と最後の回転で描いた二つの円の大きさを記した。そしてそこから、物体の初めと終わりの速さを求めた。ついで、物体が平均の速さで12個の平均の円を描くあいだに、これらの速さの差を失うと考えて、その差が、物体が12個の円を通る全行程で失われうるような抵抗を求めた。そして、この種の実験はあまり正確には行えなかったけれども、求めた抵抗は、前の実験とよく一致した。

最後に、この時代の哲学者たちの最も受け入れられた意見は、あるきわめて微細なエーテル状の媒質が存在して、すべての物体のすべての孔と通路を自由に貫いているというものである。そしてまた、物体の孔を流れるそのような媒質から、抵抗が生じるはずだというものである。そこで、動く物体で私たちが経験する抵抗が、すべてその外側の表面にあるのか、それとも内部の部分もそれぞれの表面で目立った抵抗を感じるのかを確かめるために、私は次のような実験を考えた。長さ11フィートの糸で、十分に丈夫な鋼の鉤から、鋼の環を介して、モミ材の丸い箱をつるし、前述の長さの振り子をつくった。鉤は上向きに尖っていて、その刃は凹んでおり、環が上の弧でその刃にのって、自由に動くようにした。環の下の弧には糸をつないだ。こうして作った振り子を、垂直からおよそ6フィートの距離まで、鉤の刃に垂直な面に沿って引いた。それは、振り子の振動のとき、環が鉤の刃の上を行ったり来たりすべらないようにするためである。というのは、環が鉤に触れるつり点は、動かないままでなければならないからである。そこで、振り子を引いた位置を正確に記録し、振り子を放してから、最初、2回目、3回目の振動の終わりにもどってくる三つの別の位置も記録した。これを何度も繰り返して、それらの位置をできるだけ正確に求めた。ついで、箱を鉛と、手元にあったより重い金属で満たした。だがその前に、空の箱を、その箱のまわりに巻かれた糸の部分と、鉤と箱の間に張られた残りの部分の半分とともに、秤にかけた。（張られた糸は、垂直から離れた振り子を、いつも自分の重さの半分で押すからである。）この重さに、箱が収めている空気の重さを加えた。全部で、金属で満たされた箱の、およそ78分の1の重さであった。つぎに、金属で満たされた箱は、その重さで糸を引き延ばし、振り子の長さを増したので、すでに振動している振り子の長さが前と同じになるように、糸を縮めた。それから、最初に記した位置まで振り子を引いて放し、箱が2番目に記した位置にもどるまで、およそ77回の振動を数え、ついで同じだけの回数で3番目に記した位置にもどるまで、またさらに同じだけの回数で、もどって4番目の位置に達するまで数えた。ここから私は、満ちた箱の全抵抗は、空の箱の全抵抗に対して、78 対 77 より大きい比をもたないと結論する。というのは、もし二つの抵抗が等しいならば、満ちた箱は、その固有の力が空の箱の固有の力の78倍であるために、自分の振動運動を、それだけ長く保たなければならず、したがって、78回の振動が完了するたびに、記した位置にもどるはずだからである。ところが実際には、77回の振動が完了したときに、同じ位置にもどった。

そこで、$A$ を箱の外側の表面での抵抗、$B$ を空の箱の内部での抵抗とする。そして、もし等速の物体の内部での抵抗が、物質、すなわち抵抗を受ける粒子の数に比例するならば、満ちた箱の内部での抵抗は $78B$ であるから、空の箱の全抵抗 $A + B$ は、満ちた箱の全抵抗 $A + 78B$ に対して 77 対 78 であり、分けて取れば $A + B$ 対 $77B$ は 77 対 1、したがって $A + B$ 対 $B$ は $77 \times 77$ 対 1、分けて取れば $A$ 対 $B$ は 5928 対 1 である。ゆえに、空の箱の内部での抵抗は、その外側の表面での抵抗の 5000 分の 1 より、さらに小さい。これは、満ちた箱のより大きい抵抗が、ある微細な流体が封じ込められた金属に及ぼす作用から生じる、という仮定から論じたものである。しかし私は、原因はまったく別だと考える。というのは、満ちた箱の振動の時間は、空の箱の振動の時間より短く、したがって、満ちた箱の外側の表面での抵抗は、その速さと、振動で描く空間の長さに比べて、空の箱のそれより大きいからである。そうであるならば、箱の内部での抵抗は、ないか、まったく感じられないかのどちらかである。

この実験は、記憶によって述べた。というのは、かつてそれを書き記した紙が失われたからである。そのため、記憶から失われた数の端数のいくつかを、省かざるをえなかった。すべてをもう一度試みる暇はない。最初のとき、私は弱い鉤を使ったので、満ちた箱のほうが早く減速した。その原因を探すと、弱い鉤が箱の重さに屈し、その振動に従って、あらゆる方向に曲がっていたことがわかった。そこで、丈夫な鉤を用意して、つり点が動かないままにしたところ、すべてが上に記したとおりになった。

球形の物体の水や水銀の中での抵抗を求めたのと同じ方法で、ほかの形の物体の抵抗も求めることができる。こうして、さまざまな船の形を、小さな模型でつくって互いに比べれば、どれが航海にもっともふさわしいかを、少ない費用で試すことができる。

\section{第2編 第8節：流体内を伝播する運動（波動・音速の理論、命題XLI〜L）}


% Section: 第2編 第8節：流体内を伝播する運動（波動・音速の理論、命題XLI〜L）



\begin{center}\textbf{流体を通じて伝わる運動について}\end{center}

\condpagelarge{命題 XLI. 定理 XXXI.}

\textit{圧力は、流体の粒子が一直線上に並んでいる場合を除いて、流体を通じて直線に沿っては伝わらない。}

\pngright{362.png}{}{}{}

もし粒子 $a$, $b$, $c$, $d$, $e$ が一直線上に並んでいるなら、圧力は $a$ から $e$ へまっすぐに伝わりうる。しかし粒子 $e$ は、斜めに置かれた粒子 $f$ と $g$ を斜めに押し、その粒子 $f$ と $g$ は、さらに先の粒子 $h$ と $k$ に支えられない限り、与えられた圧力を支えられない。そして支えられる限り、支えている粒子を押し、これらは、さらに先の $l$ と $m$ に支えられ、それらを押さない限り圧力を支えられず、以下、限りなくつづく。したがって圧力は、一直線上にない粒子に伝わると同時に、広がり始め、斜めに限りなく伝わっていく。そして斜めに伝わり始めたのち、一直線上にないさらに先の粒子にぶつかれば、ふたたび広がり、それは、正確に一直線上にない粒子にぶつかるたびに起こる。\QED

\pngright{365.png}{}{}{}

\textit{系.}\wsp{}もし、与えられた点から流体を通じて伝わる圧力の一部が、障害物にさえぎられるなら、さえぎられなかった残りの部分は、障害物の背後の空間に広がる。このことは、次のようにも証明できる。点 $A$ から圧力があらゆる方向に、もし可能なら直線に沿って伝わり、孔のあいた障害物 $NBCK$ によって、孔 $BC$ を通る円錐形の部分 $APQ$ を除いたすべてがさえぎられるとする。円錐 $APQ$ を、横断する平面 $de$, $fg$, $hi$ で円錐台に分ける。そして、円錐 $ABC$ が圧力を伝えながら、さらに先の円錐台 $degf$ を面 $de$ で押し、この円錐台が次の円錐台 $fgih$ を面 $fg$ で押し、その円錐台が第三の円錐台を押し、以下限りなくつづく間に、（運動の法則IIIにより）第一の円錐台 $defg$ は、第二の円錐台 $fghi$ の反作用によって、面 $fg$ で、それを押し圧するのと同じだけ押され圧されることは明らかである。ゆえに、円錐 $Ade$ と円錐台 $fhig$ の間の円錐台 $degf$ は、両側から圧縮され、したがって（命題XIXの系6により）、四方から同じ力で圧縮されない限り、自分の形を保てない。それゆえ、面 $de$, $fg$ で圧されるのと同じ衝撃で、側面 $df$, $eg$ のほうへ逃れようとし、そこでは（剛体ではなく、あらゆる点で流体だから）、その努力を抑えるまわりの流体がなければ、外へ出て広がる。したがって、外へ出ようとする努力によって、まわりの流体を側面 $df$, $eg$ で、円錐台 $fghi$ を押すのと同じ衝撃で押す。そのため、圧力は、側面 $df$, $eg$ から両側の空間 $NO$, $KL$ へも、面 $fg$ から $PQ$ のほうへ伝わるのと同じだけ伝わる。\QED

\condpagelarge{命題 XLII. 定理 XXXII.}

\textit{流体を通じて伝わる運動は、すべて、まっすぐな道からそれて、動いていない空間へ広がる。}

\textit{場合1.}\wsp{}運動が点 $A$ から孔 $BC$ を通って伝わり、（もし可能なら）円錐形の空間 $BCQP$ の中を、点 $C$ から広がる直線に沿って進むとしよう。そしてまず、この運動が、静かな水面の波であると仮定する。$de$, $fg$, $hi$, $kl$ などを個々の波の最も高い部分とし、これらは同数の間の谷で互いに隔てられている。すると、水は、波の峰では流体の動いていない部分 $LK$, $NO$ より高いから、その水は、峰の端 $e$, $g$, $i$, $l$ など、$d$, $f$, $h$, $k$ などから、両側へ $KL$ と $NO$ のほうへ流れ下る。また、波の谷では、流体の動いていない部分 $KL$, $NO$ より低いから、その動いていない部分から、波の谷へ水が流れ込む。前の流れ下りによって波の峰が、あとの流れ込みによって谷が、両側へ広がり、$KL$ と $NO$ のほうへ伝わる。そして、$A$ から $PQ$ のほうへの波の運動は、峰が次の谷へ絶えず流れ下ることによって起こり、したがって、降下の速さより速くはなく、水が両側へ $KL$ と $NO$ のほうへ流れ下ることも、同じ速さで行われなければならないから、波の広がりは、両側へ $KL$ と $NO$ のほうへ、波そのものが $A$ から $PQ$ のほうへまっすぐ進むのと同じ速さで伝わる。したがって、両側の $KL$ と $NO$ までの空間全体が、広がった波 $rfgr$, $shis$, $tklt$, $vmnv$ などによって占められる。\QED これがそのとおりであることは、静かな水の中で誰でも試すことができる。

\textit{場合2.}\wsp{}いま、$de$, $fg$, $hi$, $kl$, $mn$ が、点 $A$ から弾性のある媒質を通じて次々に伝わる脈動を表すとする。脈動は、媒質の次々の凝縮と希薄化によって伝わると考える。すなわち、各脈動の最も密な部分は、中心 $A$ のまわりに描いた球面を占め、続く脈動の間には等しい間隔がはさまる。線 $de$, $fg$, $hi$, $kl$ などが、孔 $BC$ を通って伝わった脈動の最も密な部分を表すとする。そして、媒質は、そこでは、両側の $KL$ と $NO$ のほうの空間より密であるから、両側にあるその空間 $KL$, $NO$ のほうへも、脈動のより希薄な間隔のほうへも広がる。そのようにして、間隔のところではつねにより希薄になり、脈動のところではより密になって、それらの運動に加わる。そして、脈動の前進運動は、より密な部分が、先のより希薄な間隔のほうへ絶えずゆるむことから生じ、脈動は同じ速さで、媒質の静止した部分 $KL$, $NO$ のほうへ両側にゆるまなければならないから、これらの脈動は、動いていない空間 $KL$, $NO$ へ、中心 $A$ からまっすぐ伝わるのと同じ速さで四方に広がり、したがって空間全体 $KLON$ を占める。\QED

私たちはこれを音で経験する。音は、間に家があっても聞こえ、窓から部屋に入れば部屋のあらゆる部分に広がって、すべての隅で聞こえる。それは向かい側の壁で反射されたのではなく、窓からまっすぐ伝わったのである。

\textit{場合3.}\wsp{}最後に、どのような種類の運動であっても、$A$ から孔 $BC$ を通って伝わるとする。この伝わり方は、中心 $A$ により近い媒質の部分が、さらに先の部分を押し動かす限りでしか起こらない。そして押される部分は流体であり、したがって、圧されるのがより少ない地域へ四方に退く。それゆえ、それらは、媒質の静止した部分のすべてへ、側方の $KL$ と $NO$ へも、前方の $PQ$ へも退き、そのようにして、運動は、孔 $BC$ を通るやいなや、広がり始め、そこから、あたかも始まりであり中心であるかのように、すべての向きへまっすぐ伝わる。\QED

\condpagelarge{命題 XLIII. 定理 XXXIII.}

\textit{振動するすべての物体は、弾性のある媒質の中では、脈動の運動を四方へまっすぐ伝え、弾性のない媒質の中では、円を描く運動を引き起こす。}

\textit{場合1.}\wsp{}実際、振動する物体の部分は、かわるがわる行ったり来たりしながら、行くときに、自分に隣り合う媒質の部分を押し進め、押しながらそれらを圧縮して凝縮し、ついで、もどるときに、圧縮された部分が退き、広がるのを許す。したがって、振動する物体に近い媒質の部分は、その物体の部分に似て、かわるがわる行き、もどる。そして、この物体の部分がこれらの媒質の部分を揺り動かしたのと同じ仕方で、これらは、似た振動で揺り動かされて、自分に隣り合う部分を揺り動かし、これらも同様に揺り動かされて、さらに先の部分を揺り動かし、以下、限りなくつづく。そして、媒質の最初の部分が行くときに凝縮し、もどるときにゆるむように、ほかの部分も、行くたびに凝縮し、もどるたびに広がる。したがって、すべてが同時に行き、同時にもどるのではなく（そうすれば、互いの距離を一定に保つことになり、かわるがわる希薄になり凝縮することがない）、凝縮するときには互いに近づき、希薄になるときには互いに遠ざかって、そのうちのいくつかが行くあいだに、ほかのものがもどり、それがかわるがわる限りなくつづく。そして、行く部分、行くことによって凝縮された部分は、障害物を打つ前進運動によって、脈動である。したがって、すべての振動する物体から、次々の脈動がまっすぐ伝わり、それは、物体が一つ一つの振動で一つ一つの脈動を引き起こす時間の間隔が等しいから、ほぼ互いに等しい距離をおいて伝わる。\QED そして、振動する物体の部分が、ある確かな決まった向きに行き、もどるにしても、そこから媒質を通じて伝わる脈動は、前の命題により、側方に広がり、その振動する物体を共通の中心として、ほぼ球面で同心の面に沿って、四方に伝わる。このことの例は、波に見られる。波は、もし振動する指によって起こされるなら、指の運動の向きに沿って両側に進むだけでなく、同心円のようにたちまち指を取り囲み、四方に伝わる。というのは、波の重さが弾性の力の代わりをするからである。

もし媒質が弾性のないものなら、振動する物体の揺れる部分に圧された媒質の部分は凝縮できないから、運動は、媒質が最も容易に道を譲るところへ、すなわち、振動する物体が、そうしなければ背後に空虚なまま残すはずの部分へ、瞬時に伝わる。これは、任意の媒質の中に投げ出された物体の場合と同じである。媒質は、投射体に道を譲るとき、限りなくは退かず、円を描くように進んで、物体が背後に残す空間へ向かう。したがって、振動する物体が任意の向きに進むたびに、媒質は道を譲りながら、円を描いて、物体が残す部分へ進み、物体がもとの位置へもどるたびに、媒質はそこから押しもどされて、もとの位置にもどる。そして、振動する物体が固くなく、あらゆる点でたわむとしても、大きさが与えられたままなら、その振動によって、媒質をどこかで押せば、同時に別のところでそれが道を譲らざるをえないから、媒質は、押されるところから退きながら、つねに、道を譲る部分へ円を描いて進む。

\textit{系.}\wsp{}したがって、炎の部分の動揺が、まわりの媒質を通じて直線に沿って圧力を伝えるのに役立つと信じる人々は、誤っている。そのような圧力は、炎の部分の動揺だけから生じるのではなく、全体の膨張から導かれなければならない。

\condpagelarge{命題 XLIV. 定理 XXXIV.}

\textit{水が、まっすぐ立った管の脚 $KL$, $MN$ の中を、かわるがわる上り下りするとする。そして、つり点と振動の中心との間の長さが管の中の水の長さの半分に等しい振り子をつくる。このとき、水は、振り子が振動するのと同じ時間で上り下りする、と私は言う。}

\pngright{369.png}{}{}{}

水の長さを、管と脚の軸に沿って測り、それをこれらの軸の和に等しくする。そこで $AB$, $CD$ を、それぞれの脚の中の水の平均の高さとする。そして、脚 $KL$ の中の水が高さ $EF$ まで上るとき、脚 $MN$ の中の水は高さ $GH$ まで下りるとする。ところで $P$ を振り子の物体、$VP$ を糸、$V$ をつり点、$SPQR$ を振り子が描くサイクロイド、$P$ をその最低点、$PQ$ を高さ $AE$ に等しい弧とする。水の運動をかわるがわる加速し減速する力は、どちらかの脚の水の重さが他方の脚の水の重さを超える分である。したがって、水が脚 $KL$ で $EF$ まで上り、他方の脚で $GH$ まで下りるとき、その力は水 $EABF$ の重さの2倍であり、したがって水全体の重さに対して $AE$、すなわち $PQ$ 対 $VP$、すなわち $PR$ の比をなす。また、重りの物体 $P$ が、サイクロイドの任意の位置 $Q$ で加速され減速される力は、その全体の重さに対して、最低の位置 $P$ からの距離 $PQ$ 対サイクロイドの長さ $PR$ の比をなす。ゆえに、等しい距離 $AE$, $PQ$ を描く水と振り子の起動力は、動かされるべき重さに比例する。したがって、これらの力は、水と振り子がはじめに等しい速さで動くなら、等しい時間に等しく動かしつづけ、往復運動で同時に行きかつもどるようにする。\QED

\textit{系1.}\wsp{}したがって、上り下りする水の、すべての往復は、運動が強くても弱くても、等時的である。

\textit{系2.}\wsp{}もし管の中の水全体の長さがパリのフィートで $6\frac{1}{9}$ なら、水は1秒の時間で下り、次の1秒の時間で上る。以下、かわるがわる限りなくつづく。実際、長さ $3\frac{1}{18}$ フィートの振り子は、1秒の時間で振動するからである。

\textit{系3.}\wsp{}水の長さを増すか減らすと、往復の時間は、長さの比の平方根の比で増減する。

\condpagelarge{命題 XLV. 定理 XXXV.}

\textit{波の速さは、幅の比の平方根の比をなす。}

これは次の命題の作図から導かれる。

\condpagelarge{命題 XLVI. 問題 XI.}

\textit{波の速さを求めること。}

つり点と振動の中心との間の長さが波の幅に等しい振り子をつくる。すると、その振り子が一つ一つの振動を行う時間に、波は進んで、ほぼ自分の幅を進み終える。

波の幅とは、最も低い谷か最も高い峰かの間にはさまる、横の長さをいう。$ABCDEF$ を、次々の波で上り下りする静かな水の表面とし、$A$, $C$, $E$ などを波の峰、$B$, $D$, $F$ などを間の谷とする。波の運動は、水が次々に上り下りすることで起こり、いま最も低い部分 $A$, $C$, $E$ などが、まもなく最も高くなる。そして、最も高い部分が下り、最も低い部分が上る起動力は、持ち上げられた水の重さである。したがって、この上り下りのかわるがわるは、管の中の水の往復運動に似ており、同じ時間の法則にしたがう。それゆえ（命題XLIVにより）、波の最も高い位置 $A$, $C$, $E$ と最も低い位置 $B$, $D$, $F$ の間の距離が、振り子の長さの2倍に等しいならば、最も高い部分 $A$, $C$, $E$ は、一つの振動の時間で最も低くなり、もう一つの振動の時間でふたたび上る。したがって、個々の波が通過する時間は、二つの振動の時間である。すなわち、波は、その振り子が2回振動する時間に、自分の幅を進む。ところが同じ時間に、長さがその4倍で、波の幅に等しい振り子は、1回振動する。\QED

\textit{系1.}\wsp{}したがって、幅がパリのフィートで $3\frac{1}{18}$ の波は、1秒の時間に進んで自分の幅を進み終える。したがって1分の時間には $183\frac{1}{3}$ フィートを進み、1時間の間にはほぼ 11000 フィートを進む。

\textit{系2.}\wsp{}そして、より大きな波やより小さな波の速さは、幅の比の平方根の比で増減する。

これらは、水の部分がまっすぐ上りまっすぐ下るという仮定の上でのことである。しかし、その上り下りは、実際には円を描いて行われるから、私は、この命題で時間は、ほぼ定められるにすぎないと断言する。

\condpagelarge{命題 XLVII. 定理 XXXVI.}

\textit{弾性のある流体の中を伝わる脈動の速さは、弾性の力の平方根に正比例し、密度の平方根に反比例する、という合成された比をなす。ただし、流体の弾性の力は、その凝縮に比例すると仮定する。}

\textit{場合1.}\wsp{}媒質が均質で、これらの媒質中の脈動の距離が互いに等しく、一方の媒質で運動がより強いとする。すると、対応する部分の収縮と膨張は、その運動に比例する。この比例は正確ではない。しかし、収縮と膨張がきわめて強くなければ、感じられるほどの誤りはないから、物理的に正確なものとみなせる。ところで、起動する弾性の力は、収縮と膨張に比例し、等しい部分に同時に生じる速さは力に比例する。したがって、対応する脈動の等しく対応する部分は、収縮と膨張に比例する距離を、距離に比例する速さで、同時に行き、もどる。そのため、行きともどりの一回の時間に、進んで自分の幅を進み終え、すぐ前の脈動の位置につねに続く脈動は、距離が等しいから、どちらの媒質でも等しい速さで進む。

\textit{場合2.}\wsp{}しかし、脈動の距離、すなわち長さが、一方の媒質で他方より大きいなら、対応する部分が、行き、もどるごとに、脈動の幅に比例する距離を描くと仮定する。すると、それらの収縮と膨張は等しい。したがって媒質が均質なら、それらを往復運動で動かす弾性の起動力も等しい。ところが、この力で動かされる物質は、脈動の幅に比例し、行き、もどるごとに動かなければならない距離も、同じ比である。そして、一回の行きともどりの時間は、物質の平方根の比と距離の平方根の比との合成された比、したがって距離に比例する。ところで、脈動は、一回の行きともどりの時間に、自分の幅を進み終える。すなわち、時間に比例する距離を進む。したがって、等速である。

\textit{場合3.}\wsp{}ゆえに、密度と弾性が等しい媒質では、すべての脈動は等速である。そして、もし媒質の密度か弾性が強められるなら、起動力は弾性の比で、動かされる物質は密度の比で増すから、前と同じ運動が行われる時間は、密度の比の平方根の比で増え、弾性の比の平方根の比で減る。したがって、脈動の速さは、媒質の密度の平方根の逆の比と弾性の平方根の比との合成された比をなす。\QED

\condpagelarge{命題 XLVIII. 定理 XXXVII.}

\pngright{374.png}{}{}{}

\textit{脈動が流体を通じて伝わるとき、流体の個々の粒子は、ごく短い往復運動で行きもどりしながら、つねに、振動する振り子の法則にしたがって、加速され減速される。}

$AB$, $BC$, $CD$ などを、次々の脈動の等しい距離とし、$ABC$ を、$A$ から $B$ のほうへ伝わる脈動の運動の向きとする。$E$, $F$, $G$ を、静止した媒質の三つの物理的な点で、直線 $AC$ 上に互いに等しい距離をおいて置かれているとする。$Ee$, $Ff$, $Gg$ を、これらの点が一回ごとの振動で往復運動によって行きもどりする、等しいごく短い距離とし、$\epsilon$, $\varphi$, $\gamma$ をそれらの点の中間の任意の位置、$EF$, $FG$ をそれらの点の間にはさまる物理的な小線、すなわち媒質の線状の部分で、位置 $\epsilon\varphi$, $\varphi\gamma$ および $ef$, $fg$ へ次々に移されるものとする。直線 $Ee$ に等しく直線 $PS$ を引く。これを $O$ で二等分し、中心 $O$、半径 $OP$ で円 $SIPi$ を描く。この円周の全体とその部分によって、一回の振動の全時間とその比例する部分を表す。すると、任意の時間 $PH$ または $PHSh$ が経過したとき、$PS$ に垂線 $HL$ または $hl$ を下ろし、$Ee$ を $PL$ または $Pl$ に等しくとれば、物理的な点 $E$ は $\epsilon$ にある。この法則で、任意の点 $E$ は、$E$ から $\epsilon$ を通って $e$ へ行き、そこから $\epsilon$ を通って $E$ へもどりながら、同じ加速と減速の度合いで、振動する振り子と一回ごとの振動を行う。媒質の物理的な点がそれぞれ、このような運動で動かされなければならないことを証明しなければならない。そこで、媒質が何らかの原因でこのような運動をさせられると仮定して、そこから何が導かれるかを見よう。

円周 $PHSh$ 上に、等しい弧 $HI$, $IK$ または $hi$, $ik$ をとる。これらは、等しい直線 $EF$, $FG$ が脈動の全間隔 $BC$ に対するのと同じ比を、円周全体に対してもつ。そして垂線 $IM$, $KN$ または $im$, $kn$ を下ろす。点 $E$, $F$, $G$ は相似な運動で次々に動かされるから、もし $PH$ または $PHSk$ が点 $E$ の運動のはじめからの時間であるなら、$PI$ または $PHSi$ は点 $F$ の運動のはじめからの時間、$PK$ または $PHSh$ は点 $G$ の運動のはじめからの時間である。したがって、$E\epsilon$, $F\varphi$, $G\gamma$ は、点が行くときには $PL$, $PM$, $PN$ に、点がもどるときには $Pn$, $Pm$, $Pl$ に、それぞれ等しい。ゆえに、$\epsilon\gamma$ は、点が行くときには $EG - LN$ に等しく、もどるときには $EG + ln$ に等しい。ところが $\epsilon\gamma$ は、媒質の部分 $EG$ の位置 $\epsilon\gamma$ での幅、すなわち膨張であり、したがってその部分の膨張は、行くときには、その平均の膨張に対して $EG - LN$ 対 $EG$、もどるときには $EG + ln$、すなわち $EG + LN$ 対 $EG$ である。ここで、$LN$ 対 $KH$ は $IM$ 対半径 $OP$ であり、$EG$ 対 $BC$ は $HK$ 対円周 $PHShP$ であり、入れかえれば、$EG$ 対 $HK$ は $BC$ 対円周 $PHShP$ である。すなわち（円周を $Z$ とすれば）$\frac{OP \times BC}{Z}$ 対 $OP$ であり、等比により、$LN$ 対 $EG$ は $IM$ 対 $\frac{OP \times BC}{Z}$ である。したがって、位置 $\epsilon\gamma$ での部分 $EG$ の膨張の、はじめの位置 $EG$ での平均の膨張に対する比は、行くときには $\frac{OP \times BC}{Z} - IM$ 対 $\frac{OP \times BC}{Z}$、もどるときには $\frac{OP \times BC}{Z} + im$ 対 $\frac{OP \times BC}{Z}$ である。そこで $\frac{OP \times BC}{Z}$ を $V$ とおけば、部分 $EG$、すなわち物理的な点 $F$ の膨張の、その平均の膨張に対する比は、行くときには $V - IM$ 対 $V$、もどるときには $V + im$ 対 $V$ であり、同じ部分の弾性の力の、その平均の弾性の力に対する比は、行くときには $\frac{1}{V - IM}$ 対 $\frac{1}{V}$、もどるときには $\frac{1}{V + im}$ 対 $\frac{1}{V}$ である。

\pngright{373.png}{}{}{}

同じ論法で、物理的な点 $E$ と $G$ の弾性の力は、行くときには、$\frac{1}{V - HL}$ と $\frac{1}{V - KN}$ の $\frac{1}{V}$ に対する比であり、力の差の媒質の平均の弾性の力に対する比は、$\frac{HL - KN}{VV - V \times HL - V \times KN + HL \times KN}$ 対 $\frac{1}{V}$ である。すなわち（脈動が短いので $HK$ と $KN$ が量 $V$ よりかぎりなく小さいと仮定すれば）$\frac{HL - KN}{VV}$ 対 $\frac{1}{V}$、すなわち $HL - KN$ 対 $V$ である。そこで量 $V$ は与えられているから、力の差は $HL - KN$ に比例する。すなわち（$HL - KN$ 対 $HK$ が $OM$ 対 $OI$ または $OP$ に比例し、$HK$ と $OP$ が与えられているから）$OM$ に比例する。すなわち、$Ff$ を $\Omega$ で二等分すれば、$\Omega\varphi$ に比例する。同じ論法で、物理的な点 $\epsilon$ と $\gamma$ の弾性の力の差は、物理的な小線 $\epsilon\gamma$ のもどるときにも、$\Omega\varphi$ に比例する。ところが、その差（すなわち点 $\epsilon$ の弾性の力が点 $\gamma$ の弾性の力を超える分）は、はさまれた媒質の物理的な小線 $\epsilon\gamma$ を加速する力である。したがって、物理的な小線 $\epsilon\gamma$ の加速力は、その振動の中央の位置 $\Omega$ からの距離に比例する。それゆえ、時間は（第1編 命題XXXVIIIにより）弧 $PI$ で正しく表され、媒質の線状の部分 $\epsilon\gamma$ は、前に述べた法則で、すなわち振動する振り子の法則で動く。そして、媒質全体を構成するすべての線状の部分についても、理由は同じである。\QED

\textit{系.}\wsp{}ここから、伝わる脈動の数は、振動する物体の振動の数と同じであり、それらが進むうちに増えることはないことが明らかである。というのは、物理的な小線 $\epsilon\gamma$ は、はじめの位置にもどるやいなや静止し、そののちは、振動する物体の衝撃か、振動する物体から伝わる脈動の衝撃によって、新しい運動を起こされない限り、動かないからである。したがって、脈動が振動する物体から伝わらなくなるやいなや、静止する。

\condpagelarge{命題 XLIX. 問題 XII.}

\textit{媒質の密度と弾性が与えられたとき、脈動の速さを求めること。}

媒質が、私たちの空気のように、上にのしかかる重さによって圧縮されているとし、$A$ を、その重さが上にのしかかる重さに等しく、密度が、脈動が伝わる圧縮された媒質の密度と同じである均質な媒質の高さとする。そして、つり点と振動の中心との間の長さが $A$ である振り子をつくるとする。すると、その振り子が、行きともどりから成る一回の完全な振動を行う時間に、脈動は進んで、半径 $A$ の円の円周に等しい距離を進み終える。

実際、前の命題で作図したことをそのまま用いて、もし任意の物理的な線 $EF$ が、一回ごとの振動で距離 $PS$ を描くとき、行きともどりの両端の位置 $P$, $S$ で、自分の重さに等しい弾性の力に押されるとするなら、この線は、同じ力がサイクロイドの上で振動を行いうる時間に、一回ごとの振動を行う。ここでサイクロイドの全周は長さ $PS$ に等しい。それは、等しい力が、等しい小物体を等しい距離だけ同時に押すからである。したがって、振動の時間は振り子の長さの平方根の比をなし、振り子の長さは全サイクロイドの弧の半分に等しいから、一回の振動の時間は、長さ $A$ の振り子の振動の時間に対して、長さ $\frac{1}{2}PS$、すなわち $PO$ の長さ $A$ に対する比の平方根の比になる。ところが、物理的な小線 $EG$ が両端の位置 $P$, $S$ にあるときに押される弾性の力は、（前の命題の証明で）その全体の弾性の力に対して $HL - KN$ 対 $V$、すなわち（点 $K$ がいま $P$ に落ちるから）$HK$ 対 $V$ であった。そしてその全体の力、すなわち小線 $EG$ を圧縮する上の重さは、小線の重さに対して、上の重さの高さ $A$ 対 小線の長さ $EG$ の比をなす。したがって、等比により、小線 $EG$ がその位置 $P$ と $S$ で押される力は、その小線の重さに対して $HK \times A$ 対 $V \times EG$ の比をなす。そして、等しい小物体が等しい距離を押される時間は、力の平方根の逆の比をなすから、その弾性の力が押すときの一回の振動の時間は、重さの力が押すときの一回の振動の時間に対して、$V \times EG$ 対 $HK \times A$ の比の平方根の比になり、したがって、長さ $A$ の振り子の振動の時間に対しては、$V \times EG$ 対 $HK \times A$ と $PO$ 対 $A$ との合成された比の平方根の比になる。すなわち（前の命題で $V$ は $\frac{PO \times BC}{Z}$ に等しく、$HK$ は $\frac{EG \times Z}{BC}$ に等しかったから）$\frac{PO^2 \times BC \times EG}{Z}$ 対 $\frac{EG \times Z \times A^2}{BC}$ の比の平方根の比、すなわち $PO^2 \times BC^2$ 対 $Z^2 \times A^2$ の平方根の比、すなわち $PO \times BC$ 対 $Z \times A$、いいかえれば $BC$ 対 $\frac{Z \times A}{PO}$ の比である。ところで、行きともどりから成る一回の振動の時間に、脈動は進んで自分の幅 $BC$ を進み終える。ゆえに、脈動が距離 $BC$ を進む時間は、行きともどりから成る一回の振動の時間に対して、$BC$ 対 $\frac{Z \times A}{PO}$、すなわち $BC$ 対 半径が $A$ である円の円周、の比をなす。そして、脈動が距離 $BC$ を進む時間は、この円周に等しい長さを進む時間に対して、同じ比をなす。したがって、そのような振動の時間に、脈動はこの円周に等しい長さを進む。\QED

\condpagelarge{命題 L. 問題 XIII.}

\textit{脈動の間隔を求めること。}

脈動を、その振動で起こす物体の、与えられた時間での振動の数を求める。その数で、脈動が同じ時間に進みうる距離を割れば、得られた部分が一つの脈動の幅である。\QED

\condpagelarge{注.}

最後の命題は、光と音の運動に関係する。光は直線に沿って伝わるから、（命題XLIとXLIIにより）作用だけではありえない。音は、振動する物体から生じるのだから、（命題XLIIIにより）空気の伝わる脈動にほかならない。これは、音が、そばにある物体に引き起こす振動から確かめられる。ただし、それは激しく低い音、たとえば太鼓の音の場合である。というのは、より速く短い振動は、より引き起こしにくいからである。また、どんな音でも、鳴る物体と同じ音（同音）の弦に当たれば、振動を引き起こすことはよく知られている。さらに音の速さからも確かめられる。実際、雨水と水銀の比重は、互いにほぼ 1 対 $13\frac{2}{3}$ であり、気圧計で水銀の高さがイギリスのインチで 30 に達するとき、空気と雨水の比重は互いにほぼ 1 対 850 である。したがって、空気と水銀の比重は 1 対 11617 になる。そこで水銀の高さが 30 インチであるから、その重さが下にある私たちの空気を圧縮しうる、一様な空気の高さは 348500 インチ、すなわちイギリスのフィートで 29042 になる。これが、前の問題の作図で $A$ と呼んだ高さである。半径 29042 フィートの円の円周は 182476 フィートである。そして、長さ $39\frac{1}{5}$ インチの振り子は、行きともどりから成る一回の振動を、よく知られているように、2秒の時間で行うから、長さ 29042 フィート、すなわち 348500 インチの振り子は、同様の振動を $188\frac{4}{7}$ 秒の時間で行うはずである。したがって、その時間に、音は進んで 182476 フィートを進み、それゆえ1秒の時間には 968 フィートを進む。メルセンヌは、その『弾道学』の命題XXXVに、実験によって、音が5秒の時間に、フランスのトワーズで 1150（すなわちフランスのフィートで 6900）進むことを見いだしたと書いている。そこから、フランスのフィートとイギリスのフィートの比が 1068 対 1000 であるから、音は1秒の時間にイギリスのフィートで 1474 進むことになる。同じメルセンヌは、幾何学者のロベルヴァルが、テオドンの包囲戦のとき、大砲の轟音が、火を見てから 13 または 14 秒のちに聞こえたことを観察したが、そのとき彼は、その大砲から半リューもほとんど離れていなかったと書いている。フランスのリューは 2500 トワーズを含むから、ロベルヴァルの観察によれば、音は 13 または 14 秒の時間にパリのフィートで 7500 を進み、したがって1秒の時間にパリのフィートで 560、イギリスのフィートで約 600 を進んだ。これらの観察は互いに大きく違い、私たちの計算は中間にある。私たちの学寮の回廊は長さ 208 フィートで、その一方の端で出した音は、4回の反響でこだまを4重にする。私は実験で、音が1回反響するごとに、長さがおよそ 6 または 7 インチの振り子が、先の音の反響へ行き、あとの反響へもどりながら振動することを見いだした。振り子の長さをあまり正確には決められなかったが、長さ 4 インチでは振動が速すぎ、9 インチでは遅すぎると判断した。したがって、音が行ってもどる 416 フィートの距離を進む時間は、9 インチの振り子が振動する時間より短く、4 インチの振り子が振動する時間より長い。すなわち、$28\frac{3}{4}$ 第3分（秒の60分の1）より短く、$19\frac{1}{6}$ 第3分より長い。そのため、1秒の時間に音が進むイギリスのフィートは、866 より多く 1272 より少ない。したがって、ロベルヴァルの観察から見積もるより速く、メルセンヌの観察から見積もるより遅い。さらにその後のより正確な観察で、私は、振り子の長さは 5 インチ半より大きく 8 インチより小さくなければならないことを定めた。そのため、音が1秒の時間に進むイギリスのフィートは 920 より多く 1085 より少ない。ゆえに、上に述べた幾何学的な計算による音の運動は、これらの限界の間にあって、これまで試すことのできた範囲では、現象とよく合う。そして、この運動は空気全体の密度に依存するから、音は、エーテルやある微細な空気の運動にあるのではなく、空気全体の動揺にあることになる。

空気を抜いた容器の中を音が伝わることについての、いくつかの実験は、これに反するように見える。しかし、容器から空気をすっかり抜くことはほとんどできず、十分に抜けば、音は目立って弱まるのがふつうである。たとえば、もし空気全体のわずか100分の1が容器に残るなら、音は100分の1の力で弱くなるはずで、したがって、同じ音を自由な空気の中で出して聞いている人が、そのうち鳴る物体から10倍の距離まで遠ざかった場合に劣らず、聞こえるはずである。そこで、等しく鳴る二つの物体を比べなければならない。その一方は空気を抜いた容器の中に、他方は自由な空気の中にあって、聞く人からの距離が、空気の密度の平方根の比をなすようにする。そして、もし前の物体の音が、あとの物体の音を超えなければ、異論はなくなる。

音の速さがわかれば、脈動の間隔もわかる。メルセンヌは（『和声』第1巻 命題IVに）、そこで述べるいくつかの実験によって、張った弦が、4フィートの開いたオルガンの管、あるいは2フィートの閉じた管（オルガン製作者がC fa ut と呼ぶもの）と同音になるとき、1秒の間に104回往復することを見いだしたと書いている。したがって、音が1秒の時間に描く 968 フィートの距離の中に、104 の脈動があり、それゆえ一つの脈動は、およそ $9\frac{1}{4}$ フィート、すなわち管の長さのおよそ2倍の空間を占める。そこから、すべての開いた管の音で、脈動の幅は、管の長さの2倍に等しいことが本当らしい。

さらに、鳴る物体の運動がやむと、音がただちにやみ、鳴る物体からいちばん遠く離れていても、最も近いときより長く聞こえつづけない理由は、この編の命題XLVIIIの系から明らかである。また、拡声管で音が大きく強まる理由も、述べた原理から明らかである。実際、すべての往復運動は、一回の往復ごとに、それを生じる原因によって強められるのがふつうである。ところが、管の中の運動は、音が広がるのを妨げられるので、よりゆっくりと失われ、より強くもどってくる。そのため、一回の往復ごとに加えられる新しい運動によって、いっそう強められる。以上が、音のおもな現象である。
\section{第2編 第9節：流体の渦運動（デカルト渦動説の批判的論駁、命題LI〜LIII）}


% Section: 第2編 第9節：流体の渦運動（デカルト渦動説の批判的論駁、命題LI〜LIII）



\section*{IX. 流体の円運動}

\centering
\textbf{流体の渦運動に関する考察。}

\bigskip

\textit{流体の部分間の相対的な速度が一定であるとき、その摩擦抵抗は他の条件が等しい場合に、流体部分が互いに分離される速度と比例する。}

\bigskip

\textbf{命題 LI. 理論 XXXVIII.}

\textit{無限長い固形円柱が無限大の均一な流体中に、軸を中心に等速で回転させられ、その衝撃のみによって流体を球運動に導くと、流体の部分はすべて一定の速度で球運動を続ける。}

\textbf{第2編 第9節：流体の渦運動}

\textit{命題 LI. 定理 1.} 流体が均一な速度で回転している場合、その流体の各部分の周期的時間は、それらの部分から円筒軸までの距離に比例します。


\par\begin{center}
    \centering
    \includegraphics[width=0.85\textwidth]{381.png}

\end{center}\par

円筒 $AFL$ が軸 $S$ 回りに均一な速度で回転しているとします。この円筒を同心円周 $BGM$, $CHN$, $DIO$, $EKP$ 等に分けて、流体を無数の同心円筒形の固体層に区分けします。そして、流体が均質であるとする仮定により、隣接する円筒面同士での相互作用は、それらの円筒の移動と接触面積に比例すると考えます。

\textbf{命題 LII. 定理 2.} 各部分の周期的時間は、その部分から軸までの距離に比例します。

各部分が均一な速度で回転していると仮定し、流体の各部分の周期的時間を $t$ とします。このとき、各部分から円筒軸までの距離を $r$ とすると、命題 LI の仮定により、各部分の周期的時間はその距離 $r$ に比例します。

\textbf{命題 LIII. 定理 3.} 流体が均一な速度で回転している場合、流体の各部分から円筒軸までの距離とその部分の周期的時間との間には、平方反比関係があります。

各部分から円筒軸までの距離 $r$ とその部分の周期的時間 $t$ の関係を考察します。命題 LII の結果により、$t \propto r$ が成り立ちます。また、流体の各部分は均一な速度で回転しているため、その角加速度 $\alpha$ は一定です。したがって、周期的時間 $t$ と距離 $r$ の関係は平方反比関係であることが示されます。

\textbf{証明}：各部分の周期的時間 $t$ と距離 $r$ の間には、
\[ t \propto \frac{1}{r^2}. \]

以上が第9節「流体の渦運動」における主な内容です。デカルトの渦動説に対する批判的な論駁は、これらの命題を通じて行われています。

これらの印象が生じる。もし円周に対する印象が凹面より凸面側で大きいか小さいかによりますと、強力な印象が優位となることになり、その結果、円周の運動は加速されるか遅延される。これは円周の運動方向と同一または反対の領域に向かうように影響します。

したがって、各円周が自らの運動において均一に perseveret（継続する），両側からの印象は互いに等しくなり、反対の領域へと生じるべきです。ここでは、印象は接面近傍であり、これらの移動は互いに対向するため、表面に対する移動は表面に対して逆方向となります。

すなわち、
\[
\textbf{translationes inverse ut superficies, hoc est inverse ut superficierum}
\]

円周が均一に運動を継続するためには、両側からの印象は等しくなり、反対の領域へと作用すべきである。そして、これらの印象は近傍の表面から生じ、互いに対向する移動により、表面に対する移動は逆方向となる。

\textbf{離心から軸までの距離。} しかし、角運動量の差は軸からの距離と直線的な移動が適用される距離との関係で異なります。すなわち、これらは（理由を結びつけると）逆に平方根比例します。したがって、無限遠方まで延びる直線 \( SABCDEQ \) の各部分に対して、垂直線 \( Aa \), \( Bb \), \( Cc \), \( Dd \), \( Ee \) などからそれぞれの \( SA \), \( SB \), \( SC \), \( SD \), \( SE \) の平方が逆比例するものとし、これらの垂直線の終点を通る曲線を線分として理解すると、離心距離の総和は、

 
\textbf{双曲線} となるでしょう。

\textbf{これが全体の角運動量は、直線的動径 \(Aa\)，\(Bb\)，\(Cc\)，\(Dd\)，\(Ee\) の合計に比例することである。すなわち、流体を均質とし、球体の数を無限大に増加させ、半径を無限小に縮小させる場合、円錐曲線 \(AaQ\)，\(BbQ\)，\(CcQ\)，\(DdQ\)，\(EeQ\) などの面積と角運動時間は逆比例する。したがって、任意の粒子 \(D\) の周期的な時間は面積 \(DdQ\) に直接比例し、これは（曲線の quadratura による符号付け）半径 \(SD\) から導かれる。}

\textbf{訳注：}
\begin{itemize}
    \item \textit{summa linearum}：直線的動径の合計
    \item \textit{Medium uniformiter fluidum}：均質な流体
    \item \textit{orbium numerus augeatur \& latitudo minuatur in infinitum}：球体の数を無限大に増加させ、半径を無限小に縮小させる
    \item \textit{areae Hyperbolicae his summis Analogae}：円錐曲線 \(AaQ\)，\(BbQ\)，\(CcQ\)，\(DdQ\)，\(EeQ\) などの面積
    \item \textit{tempora motibus angularibus reciproce proportionalia erunt etiam his areis reciproce proportionalia}：角運動時間は逆比例する
    \item \textit{(per notas Curvarum quadraturas)}：（曲線の quadratura による符号付け）
\end{itemize}

\textbf{証明}

\textit{補題 1.}\quad したがって、流体の部分的な角運動量はその各部から軸までの距離に反比例し、絶対的速度は等しいです。

\textit{補題 2.}\quad もしこの流体が無限長な円筒形容器内に存在し、さらに内部に別の円筒形容器を含んでおり、両方の円筒が共通の軸を中心に回転し、各円筒の回転時間はその半径に比例し、また流体の各部分が自運動を維持する場合、個々の部分の周期的な時間はそれぞれの部から円筒軸までの距離に比例します。


\textbf{補則3.} もし円柱と流体がこのような形で回転運動を加えられたり除去されたりすると、新たな回転運動によって互いの摩擦が変化しないので、流体内部の各部分の運動も相互に変化しない。なぜなら、各部分の相対移動は互いの摩擦によるものであるから、一方の部分が他方の部分に対して摩擦を受けた場合、その部分はより速く回転することも遅くなることもなく、常に同じ速度で回転する。

\textbf{補則4.} よって、円柱全体と流体全体のシステムから外周円柱のすべての角運動を除去すると、流体の運動が得られる。

在静止的円柱体内。

\textbf{補則5.} よって、外円柱体和流体均処于静止状態時，若内円柱体均匀旋転，則流体会獲得円周运動，并逐漸伝播至整個流体中；直到流体各部分〜によれば第四補則所定義的条件獲得動量之前，其動量不会停止增長。

\textbf{補則6.} 〜により流体試図進一步拡展其動量，除非外円柱体被强力阻止，否則它也会在外円柱体周囲旋転，かつ其运動速度会持続增加直至某一時刻。

\textbf{第2編 第9節：流体の渦運動}

\textit{命題 LI. 仮説 XXXIX.}

円筒の両方の周期は互いに等しい。もし外側の円筒が強制的に留め置かれるならば、その円筒は流体の運動を遅くするよう努め、もし内側の円筒が外部からの何らかの力によって運動を受け持つことなく留め置かれれば、同じものは徐々に停止する。

これらすべては静止している深い流体中に試すことができる。

\textbf{命題 LII.}

もし固体球が均一で無限大の流体中に、軸の位置が与えられた一定の速度で回転し、その運動は流体に単に衝突によって与えられるならば、流体の一部は常に均一に保持される。

\textbf{証明}

球体 \(S\) が軸 \(A\) の周りに一定の速度で回転しているとする。この球体は流体 \(F\) 中で運動しており、流体は無限大であり均一であると仮定する。球体 \(S\) の表面に接する流体部分 \(PQ\) が球体 \(S\) の運動によって動いている。

球体 \(S\) の中心から一定の距離にある点 \(P\) における流体の速度は、球体 \(S\) の回転速度と等しい。したがって、流体 \(PQ\) の各部分は常に均一な速度で移動する。

\textbf{命題 LIII.}

もし固体球が均一で無限大の流体中に、軸の位置が与えられた一定の速度で回転し、その運動は流体に単に衝突によって与えられるならば、流体の各部分は常に等しい速度で移動する。

\textbf{証明}



\textbf{補足}

これらの命題は、デカルトの渦運動説を批判的に論駁するために用いられる。デカルトは流体中の渦運動を説明するために、流体の各部分が一定の速度で移動することを主張していたが、ニュートンの証明によれば、流体の各部分は常に均一な速度で移動することが示される。

\textbf{第2編 第9節：流体の渦運動}

\textit{命題LI.} 任意の部分がその周期的な時間は、中心から距離に比例する平方に等しい。

図\ Prop.\ LI.

\textit{1.\ 案.} 球体 \( A F L \)が中心 \( S \)を中心に均質な軸 \( S \)周りに回転し、同心円の環状領域 \( B G M \), \( C H N \), \( D I O \), \( E K P \)などに流体を分けるとする。これらの環状領域は仮定によって固体とみなす。均質な流体であるため、隣接する環状領域間の相互作用により生じる衝撃は（仮説によれば）、その相互移動に比例するものとなる。

\[
\textbf{訳注：}
\]

この部分では、ニュートンはデカルトの渦運動理論を批判的に論駁しています。デカルトは流体が円錐形の渦を形成すると主張していましたが、ニュートンはその仮説に対して反論し、流体の各部分が中心からの距離に比例する周期を持つことを示唆しています。

\& 接近する面において作用されるもの。もし凹面より凸面よりもいずれか一方の部分で作用が大きいならば、強力な作用は惑星の速度を加速または減速させる。作用方向によっては運動の同一領域へ向けるか反対方向へ向けるかにより異なる。したがって、各惑星が自らの運動において均一に perseveret するためには、双方の部分で作用が等しくなり、互いに対極的な領域で発生しなければならない。ゆえに、これらの作用は接続する面であり、それらの相互移動は

\[
\textbf{ inverse ut superficies, hoc est}
\]

すなわち、面に対して反対方向となる。

\textbf{逆に、面積から中心までの距離の平方に反比例する。しかし、角速度の変位は、これらの距離に対する直線的な応用と等しいか、あるいはこれらが直接的にそして距離反比例的に；すなわち（理由を結びつけると）、距離の立方反比例である。}

したがって、無限長な直線 \( SABCDEQ \) 上の各部分 \( SA, SB, SC, SD, SE, \&c. \) に対して、それぞれの垂線 \( Aa, Bb, Cc, Dd, Ee, \&c. \) を引き、これら各々が距離の逆立方比に比例するとき、その総距離は、すなわち全体の角速度は、対応する直線の総和と等しい。

\textbf{第9節：流体の渦運動（デカルト渦動説の批判的論駁、命題LI〜LIII）}

\textit{命題 LI.}\wsp{}もし流体を均質とし、円錐曲線に似た面積比で円軌道の数を増加させ、半径を無限大まで小さくする場合、各円軌道に対する周期時間は逆比例し、これらの面積もまた逆比例します。

すなわち、任意の円軌道 \( DIO \) の周期時間は面積 \( DdQ \) に反比例し、これは（曲線の quadratura により）距離 \( SD \) の平方に直接比例することを示すことができます。これが最初に証明したいことでした。

\textit{2. 案例}\wsp{}中心から無数の直線が球面へと引かれます。

流体を含む軸が与えられており、その角度は互いに微分して異なるものである。そして、これらの直線が軸を中心に回転すると無数の円周が形成される。各円周はさらに4つの円周で囲まれており、一つは内側、もう一つは外側であり、残りの二つは側面である。内側と外側の円周との摩擦により、各円周は互いに反対方向へ等しく推進される。これは初等運動の証明から明らかである。したがって、任意の球体から無限遠方に至る円周列は、

\[
\textbf{等しく推進される}
\]

となります。



この論証を更に明確にするため、図を用います。球体から無限遠方に至る円周列は、

\[
\textbf{等しく推進される}
\]

と示されます。



\textbf{第2編 第9節：流体の渦運動}

直線方向に進むものは、第一法則の法則により動くが、円筒面の側壁による摩擦によって妨げられる限りは例外である。しかし、この法則に基づいて運動を考慮した場合、円筒面の側壁による摩擦は存在せず、従って、その摩擦によって運動が妨げられるものも存在しない。もし中心から等距離にある環が、極周囲より近接する部分で速く回転したり遅く回転したりすると、遅いものは加速され、早いものは減速され、互いの摩擦により常に周期的に中心に向かって進むことになる。したがって、この摩擦は第一法則に基づく運動を妨げるものではない、そして

\begin{equation}
r
\end{equation}
\textbf{注：} ここで $r$ は半径を表します。

\textit{注記：} この部分では、ニュートンはデカルトの渦動説に対する批判的な論駁を行っています。円筒形容器内の流体が回転する場合に、側壁による摩擦や互いの摩擦がどのように影響するかを考察しています。ニュートンは、これらの摩擦が第一法則に基づく運動を妨げないことを主張しています。

\textbf{よって、その法則は次のようになります：}各環状体の周期的な時間は、その距離から中心までの平方に比例します。これについては二つ目の証明を行います。

\textit{\textbf{案 3.}}\wsp{}各環状体を無数の粒子に分割し、これらは絶対的に均一な流体であるとします；そしてこれらの横断面は円運動の法則に関与せず、流体の構成のみに関与するため、円運動は以前と同じように持続します。すべての環状体のこれらの横断面が微細な不規則性または相互作用の力を失うまで、

\[
\textbf{この証明は続く}
\]

非均質性は変化することなく、あるいは均質性を保つ。そして原因の比が一定である限り、効果の比もまた一定となる。すなわち運動と周期的な時間の比も一定となる。\(\blacksquare\)

他に、円周運動とその起源する心向力が日食方向より極方向に向かう場合、ある原因が必要であり、個々の粒子が自転円周上に留まることを保証し、太陽系中の物質が常に中心から遠ざかり、外側の渦巻きを通じて極へ移動し、その後極軸を通じて日食方向に戻るという現象が起こらないようにする。

\textit{補則 1.} したがって、流体の部分が球軸周りに運動する角運動は、


\textbf{訳文：}


\textit{補則 1.} したがって、流体の部分が球軸周りに運動する角運動は、


\textbf{訳文：}
流体の各部分が球軸周りに運動する角運動は、その均質性が保たれることを意味します。円周運動と心向力が日食方向より極方向に向かう場合、個々の粒子が自転円周上に留まることを確保するための原因が必要です。太陽系中の物質が常に中心から遠ざかり、外側の渦巻きを通じて極へ移動し、その後極軸を通じて日食方向に戻るという現象が起こらないようにします。\(\blacksquare\)

\textbf{第2編 第9節：流体の渦運動}

\textit{命題 LI.} 二つの球体が中心から等しい距離にあるとき、その間の距離の平方に反比例する力が働きます。また、それぞれの球体の絶対速度は、その距離から派生した同じ距離の平方に反比例します。

\textit{推論 2.} もし静止流体中に中心軸を中心に無限大で等速回転する球体が与えられた場合、流体には渦動運動が生じ、その運動は無限遠方にまで徐々に拡散します。各流体部分の速度が加速する前に、各部分の周期時間は中心から球までの距離の平方と等しくなることがありません。

\textbf{訳文}

\textit{命題 LI.} 二つの球体が中心から等しい距離にあるとき、その間の距離の平方に反比例する力が働きます。また、それぞれの球体の絶対速度は、その距離から派生した同じ距離の平方に反比例します。

\textit{推論 2.} もし静止流体中に中心軸を中心に無限大で等速回転する球体が与えられた場合、流体には渦動運動が生じ、その運動は無限遠方にまで徐々に拡散します。各流体部分の速度が加速する前に、各部分の周期時間は中心から球までの距離の平方と等しくなることがありません。


\textbf{sunt reciproce ut quadrata distantiarum a centro globi, \& velocitates
absolut{\ae} reciproce ut eadem quadrata applicata ad distantias
ab axe.}


 
\textbf{推論 2.}$\wsp{}$もし静止流体中に中心軸を中心に無限大で等速回転する球体が与えられた場合、流体には渦動運動が生じ、その運動は無限遠方にまで徐々に拡散します。各流体部分の速度が加速する前に、各部分の周期時間は中心から球までの距離の平方と等しくなることがありません。


\[
\textbf{Corol.\fsp{}2.}\wsp{}Si globus in fluido quiescente similari \& infinito circa
axem positione datum uniformi cum motu revolvatur, communicabitur
motus fluido in morem Vorticis, \& motus iste paulatim propagabitur
in infinitum; neque prius cessabit in singulis fluidi partibus
accelerari, quam tempora periodica singularum partium sint
ut quadrata distantiarum a centro globi.
\]

\textbf{補則3.} なぜなら、渦巻の内部部分はより大きい速度で外側部分を引き寄せ、その結果、内部部分はその作用により永遠に動いていることになり、さらに外部部分は他の更なる外部部分へ同じ運動量を同時に転送します。この作用によって、彼ら自身の運動量は明らかに常に一定であることを保証します。したがって、運動は渦巻の中心から円周へと永遠に移され、無限の円周上に吸収されます。二つの任意の同心球面間の素材は、内部素材から受け取ったすべての運動を転送するため、二つの球面間の任意の表面は加速されることはありません。

\begin{center}
\includegraphics[width=0.8\textwidth]{fig_19-1.png}
\end{center}

【第2編 第9節：流体の渦運動】

4. 由此可知，〜のために渦旋保持恒定的运動状態，必須存在某种活躍原理，使得球体総是接受与施加于渦旋物質相同的動量。若無此原理，則球体及其内部渦旋部分将逐漸减速，并最終停止旋転。

5. 假設另一球体相対于該渦旋中心以固定距离存在，并回着与定軸以某种傾角运動，此時渦旋的内部分子会不断向外部伝逓動量。若これ些内部分子不接受新的動量，則それら也会逐漸减速并最終停止旋転。

\textbf{常に回転し、その運動で流体を渦に引き込む：最初は新しい小さな渦が中心点周りの球体と共に回転し、その運動は徐々に広がり、無限大まで拡散していく。}これが第1の渦の様子と同じです。同様にして、この球体が第2の渦の運動によって引き込まれるとき、逆に第2の球体もこの運動によって引き込まれます。つまり、二つの球体は互いの中心点を中心に回転し、互いに円運動を避けるかのように動き続けます。ただし、その動きが何かの力によって制限されない限りです。

その後、もし球体が自己の運動で一定の力を受けると、

\begin{equation}
\textbf{これらの力を常に印加すると、}
\end{equation}

球体はその運動を維持し続けます。

\textbf{もし}機械的法則に従ってすべてが許されるとすれば、球体の運動は徐々に衰え、（コロラリウム3と4で説明した理由により）ついには渦は静止するでしょう。

\textit{コロラリウム6.} 仮定すると、いくつかの球が与えられた位置にある軸周りを与えられた速度で恒久的に回転させると、無限に続く無数の渦が生じます。なぜなら、個々の球は、一つの球がその運動を無限大まで拡散させる理由と同じく、自らの運動も無限大まで拡散させることになるからです。したがって、流体の一部は、すべての球の作用によって同じように動かされます。

\begin{center}
\includegraphics[width=0.8\textwidth]{fig22}{幅}{高さ}{オフセット}
\end{center}

結果は、渦巻く範囲が一定ではなく、互いにゆっくりと移動する。球体は、渦巻きの作用により、自身の中でゆっくりと永遠にその位置から動く。これは上位の命題で述べた通りである。また、これらの球体は互いの位置関係を一定に保つことはないが、何かしらの力によって保持される限りは例外となる。

力を保証されないとき、常に球体に作用している渦巻きの力が消失すると、物質は第三コロラリオと第四コロラリオで示された理由によりゆっくりと静止し、渦巻きを生じなくなります。

\textbf{コロラリオ} \textit{7.} もしこの流体が球形の容器内で閉じ込められた場合、球体は互いに位置関係を保ちつつ動くことになります。

在中心処進行均匀旋転的渦旋，若球体与容器囲回同一軸線以相同部分進行旋転，則それら的周期時間応当是半径的平方。流体各部分不会在其运動中保持不変，除非其周期時間是相対于漩渦中心距离的平方。除此之外，どんな其他构成漩渦的方式都是不可能持久存在的。

\textbf{補遺 8.} 若容器包含流体且球体維持此运動，かつ相対于任意与定軸線進行額外的普遍角运動，

$$
\textit{si vas, fluidum inclusum \& globus servent hunc motum, et motu praeterea communi angulari circa axem quemvis datum}
$$

則それら将保持これ种状態。

\textbf{訳文}

回転すると、この新しい運動によって流体の部分間の摩擦は変わらないため、流体の部分同士の運動も変化しない。すなわち、部分同士の移動は相互の摩擦に依存する。

任意の部分はその運動を維持し、これにより一方から他方への摩擦はもう片方からの摩擦と同じ程度に遅れたり速くならない。

\textit{補則9.} したがって、もし容器が静止しており球体が動いている場合、流体も動く。すなわち、球体を通過する平面を考え、その軸周りに反対方向で回転すると、この回転の時間と球体の回転時間との合計時間を平方根として取ることにする。

\begin{center}
\includegraphics[width=0.8\textwidth]{388.png}
\end{center}

流体の渦運動について、デカルトの渦動説を批判的に論駁する【第2編 第9節】において、以下はその一部です。

半径の比：球体の半径の平方に比例し、この平面に対する流体部の周期時間は、各距離の平方から球体中心までの差に比例します。

\textit{系10.}\wsp{}したがって、容器が球と同じ軸のまわりに、あるいは別の軸のまわりに、任意の速度で動かされるならば、流体の運動は決まる。なぜなら、系全体から容器の角運動を取り去っても、すべての運動は、系の内部では以前と同じ互いの関係を保つからである（系8による）。そして、それらの運動は系9によって与えられる。

\textit{系11.}\wsp{}容器と流体が静止していて、球が一様な運動で回転するならば、運動はしだいに流体全体に、さらに容器にまで伝わり、容器は、強く押さえつけられていない限り回り出す。そして流体と容器は、それらの周期が球の周期に等しくなるまでは、加速をやめない。もし容器がある力で押さえつけられるか、あるいは一定で一様な何らかの運動で回されるならば、媒質はしだいに系8, 9, 10で定められた運動の状態に至り、他のどんな状態にもとどまることはない。さらに、容器と球をそれぞれ決まった運動で回していた力がなくなり、系全体が力学の法則にゆだねられるならば、容器と球は流体を介して互いに作用し合い、互いの周期が等しくなって系全体がひとつの固体のように一緒に回るようになるまでは、流体を介して互いに運動を伝え合うことをやめない。

\condpagelarge{注.}

以上のすべてにおいて、私は流体が、密度と流動性について一様な物質でできていると仮定している。そのような流体とは、同じ球が同じ運動で、同じ時間のあいだに、流体のどこに置かれても、自分から等しい距離のところへ、似た等しい運動を伝えることができるものである。物質は、その円運動によって渦の軸から遠ざかろうとし、そのために、それより外側の物質をすべて押す。この圧力によって、部分どうしの摩擦はより強くなり、互いに分離しにくくなる。その結果、物質の流動性は減る。また、流体の部分がどこかでより粗い、あるいは大きいなら、そこでは部分が互いに分離する面が少ないので、流動性は小さくなる。このような場合、不足した流動性は、部分の滑らかさ、あるいは粘りけ、あるいは何か別の条件によって埋め合わされると、私は仮定する。そうならなければ、物質は流動性の小さいところでは、よりくっつき合って鈍くなり、したがって運動をよりゆっくり受け取り、前に述べた比より遠くまで伝えることになるだろう。容器の形が球形でなければ、粒子は円ではなく、容器の形に沿った線の上を動き、周期は、中心からの平均の距離の平方にほぼ比例するだろう。中心と周の間で、空間が広いところでは運動はより遅く、狭いところではより速いだろう。それでも、速い粒子が周に向かうわけではない。それらはより曲がり方の小さい弧を描くので、中心から遠ざかろうとする傾向は、速さが増す分と同じだけ、この曲がりの減少によっても減るからである。狭い空間から広い空間に進むときには、中心からやや遠くへ離れるが、この後退によって遅くなる。そして広い空間から狭い空間へ近づくときには加速される。こうして個々の粒子は、かわるがわる遅くなり、速くなることを永遠にくりかえす。これらのことは、固い容器の中ではこのとおりになる。無限の流体の中では、渦の構造は、この命題の系6によって知られる。

私がこの命題によって渦の性質を調べようと試みたのは、天体の現象が渦によって説明できるかどうかを確かめるためである。というのは、木星のまわりを回る惑星の周期は、木星の中心からの距離の3/2乗の比をなすという現象があり、同じ規則が太陽のまわりを回る惑星にも成り立つからである。これらの規則は、これまでの天文観測が明らかにしたかぎり、どちらの惑星にもきわめて正確に成り立つ。それゆえ、もしそれらの惑星が木星や太陽のまわりを回る渦に運ばれているのなら、これらの渦もまた同じ法則で回らなければならない。ところが、渦の部分の周期は、運動の中心からの距離の二乗の比になることが示された。この比は、渦の物質が中心から遠いほど流動的でない限り、あるいは、流体の部分の滑らかさの不足から生じる抵抗が、部分が互いに分離する速さの増加に応じて、速さの増える比より大きな比で増さない限り、減らして3/2乗の比にすることはできない。だが、そのどちらも理にかなっているようには見えない。より粗くて流動性の小さい部分は（中心に向かって重くなければ）周に向かうだろう。そして、この節のはじめに、証明のために抵抗が速さに比例するという仮説を立てたとしても、抵抗は速さの比より小さい比であるほうが、本当らしい。これを認めるならば、渦の部分の周期は、その中心からの距離の二乗より大きい比になる。もしも（ある人たちの意見のように）渦が中心の近くで速く動き、ある限界まで遅くなり、それから周の近くでふたたび速くなるのだとすれば、3/2乗の比も、そのほかのどんなきまった比も成り立ちえない。それゆえ、この3/2乗の比という現象が渦によってどう説明できるのかは、哲学者たちに考えてもらうことにしよう。

\condpagelarge{命題 LIII. 定理 XL.}

\textit{渦の中に運ばれて、円を描いてもとにもどる物体は、渦と同じ密度をもち、速さと進路の定まり方について、渦の部分と同じ法則で動く。}

渦の小さな一部分を考え、その粒子すなわち物理的な点が、互いに与えられた位置を保っているものとして、これが凍っていると仮定する。この部分は、密度も、固有の力も、形も変わらないのだから、以前と同じ法則で動くだろう。逆に、渦の凍って固くなった部分が、渦の残りと同じ密度をもち、流体にとかされたとすれば、その部分は、流体となった粒子が互いに動き合うことを除いて、以前と同じ法則で動くだろう。そこで粒子どうしの運動は、全体の進行運動には関係しないものとして無視すれば、全体の運動は以前と同じになる。そしてその運動は、中心から等しい距離にある渦の他の部分の運動と同じになる。これは、流体にとかされた固体が、渦の他の部分とまったく同様な渦の部分になるからである。ゆえに固体は、渦の物質と同じ密度ならば、すぐまわりにある物質に対して相対的に静止しつつ、その部分と同じ運動で動くだろう。もし固体のほうが密ならば、以前よりも渦の中心から遠ざかろうとする傾向が強い。したがって、それまでその軌道にいわばつり合いの状態でとどめていた渦の力にうち勝って、中心から遠ざかり、回りながらうず巻線を描き、もはや同じ軌道にはもどらない。同じ論法で、もしそれがまばらならば、中心に近づく。ゆえに、流体と同じ密度でない限り、同じ軌道にはもどらない。ところが、その場合には、渦の中心から等しい距離にある流体の部分と同じ法則で回ることが示されている。\QED

\textit{系1.}\wsp{}したがって、渦の中を回り、いつも同じ軌道にもどる固体は、それが浮かんでいる流体に対して、相対的に静止している。

\textit{系2.}\wsp{}また、渦が密度について一様ならば、同じ物体は、渦の中心からどんな距離でも回ることができる。

\condpagelarge{注.}

\pngright{392.png}{}{}{}

ここから、惑星が物質的な渦に運ばれているのではないことが明らかになる。というのは、コペルニクスの仮説によれば、太陽のまわりを運ばれる惑星は、太陽に焦点をもつ楕円を描いて回り、太陽に引いた半径は、時間に比例する面積を描くからである。ところが、渦の部分はこのような運動で回ることができない。$AD$, $BE$, $CF$ を太陽 $S$ のまわりに描いた三つの軌道とし、その最も外側の $CF$ は太陽と同心の円で、内側の二つの遠日点を $A$, $B$、近日点を $D$, $E$ とする。すると、軌道 $CF$ を回る物体は、太陽に引いた半径が時間に比例する面積を描くのだから、一様な運動で動く。一方、軌道 $BE$ を回る物体は、天文学の法則に従えば、遠日点 $B$ ではゆっくり、近日点 $C$ では速く動く。ところが力学の法則に従えば、渦の物質は、$A$ と $C$ の間の狭い空間では、$D$ と $F$ の間の広い空間でよりも速く動かなければならない。すなわち、遠日点のほうが近日点よりも速くなければならない。この二つは互いに矛盾する。

同じように、いま火星の遠日点がある乙女座の初めでは、火星と金星の軌道の間の距離は、魚座の初めでのこれらの軌道の間の距離のおよそ3対2である。したがって、これらの軌道の間の渦の物質は、魚座の初めには乙女座の初めよりも、3対2の比で速くなければならない。というのは、同じ量の物質が、同じ一回転の時間で通り抜ける空間が狭いほど、より大きな速さで通り抜けなければならないからである。そこで、もし地球がこの天空の物質に対して相対的に静止していて、それに運ばれながら太陽のまわりを回るとすれば、その速さは、魚座の初めでは乙女座の初めに対して、3対2の比になるはずである。すると、太陽の見かけの日周運動は、乙女座の初めでは70分より大きく、魚座の初めでは48分より小さくなってしまう。ところが（経験が示すところでは）この太陽の見かけの運動は、乙女座の初めよりも魚座の初めのほうが大きく、したがって地球は、魚座の初めよりも乙女座の初めのほうが速いのである。

こうして、渦の仮説は天文学の現象とまったく相いれず、天体の運動を説明するというよりは、かき乱すものである。しかし、これらの運動が渦のない自由な空間でどのように行われるかは、第1編から理解でき、『世界の体系』でさらに十分に教えよう。

\section*{\textbf{DE}}
\section*{\textbf{Mundi Systemate}}






\chapter{第3編 世界の体系（推理の規則・太陽系の諸現象・万有引力の確立）}



\section{第3編 哲学における推理の規則および太陽系現象（仮説I〜IX）}


% Section: 第3編 哲学における推理の規則および太陽系現象（仮説I〜IX）
\condpagelarge{第3編\quad 世界の体系について}

前の二つの編で、私は哲学の原理を述べたが、それは哲学的なものではなく、数学的なものにとどまる。すなわち、それによって哲学上の事柄が論じられうるような原理である。これらは、運動と力の法則と条件であって、哲学に最も関係するものである。それでも、これらが実りのないものに見えないように、私はいくつかの哲学的な注で、それを説明した。そこでは、一般的で、哲学が最も基礎をおくように見える事柄、たとえば物体の密度と抵抗、物体のない空間、光と音の運動を扱った。残るのは、同じ原理から、世界の体系の構成を教えることである。この主題について、私は第3編を、より多くの人に読まれるように、通俗的な方法で書いていた。ところが、先に置いた原理が十分に理解されていない人は、帰結の力をまったく感じとれず、長い年月にわたって慣れ親しんできた先入観を捨てないであろう。そこで、事柄が論争に引き込まれないように、その編の要点を、数学の流儀で命題に移した。それは、先に原理を調べ終えた人だけに読まれるようにするためである。しかし、そこには、数学に通じた読者にも、あまりに長い足止めとなりうる命題が、非常に多く出てくるから、私は、誰もがそれらをすべて調べることを望まない。もし誰かが、定義、運動の法則、そして第1編の最初の三つの節を注意深く読むなら、それで足りるであろう。そのうえで、世界の体系についてのこの編に進み、前の編のほかの命題で、ここに引用されたものは、好みに応じて参照すればよい。

\condpagelarge{仮説}

\textbf{仮説 I.}\wsp{}\textit{自然の事物の原因は、真であって、かつその現象を説明するのに十分なもの以上のものを、認めてはならない。}

実際、自然は単純であって、余計な原因をほしいままにしない。

\textbf{仮説 II.}\wsp{}\textit{したがって、同じ種類の自然の結果には、同じ原因がある。}

たとえば、人間と獣の呼吸、ヨーロッパとアメリカでの石の落下、台所の火と太陽の光、地球と惑星での光の反射。

\textbf{仮説 III.}\wsp{}\textit{すべての物体は、別の任意の種類の物体に変わることができ、性質のあらゆる中間の度合いを次々に帯びる。}

\textbf{仮説 IV.}\wsp{}\textit{世界の体系の中心は静止している。}

これは、ある人は地球が、ある人は太陽が中心で静止していると主張するが、すべての人に認められている。

\textbf{仮説 V.}\wsp{}\textit{木星のまわりを回る惑星は、木星の中心へ引いた半径で、時間に比例する面積を描き、その周期は、木星の中心からの距離の3/2乗の比をなす。}

これは天文観測から明らかである。これらの惑星の軌道は、木星に同心の円とは感じられるほど違わず、それらの円上のそれらの運動は一様であることがわかる。そして、周期が軌道の半径の3/2乗の比をなすことは、天文学者たちが一致して認めている。そして、これらすべてを、マイクロメーターと衛星の食によって、より正確に定めたフラムスティードは、私への手紙で、さらに彼の数値を私に知らせて、この3/2乗の比が、感覚でとらえられる限りでは、できるかぎり正確に成り立つことを示した。このことは、次の表から明らかである。

\condpagelarge{衛星の周期}

\begin{center}
\begin{tabular}{llll}
1日 18時 $28\minute\frac{3}{5}$ & 3日 13時 $17\minute\frac{9}{10}$ & 7日 3時 $59\minute\frac{3}{5}$ & 16日 18時 $5\minute\frac{1}{5}$ \\
\end{tabular}
\end{center}

\condpagelarge{衛星の木星の中心からの距離}

\begin{center}
\begin{tabular}{l|c|c|c|c|l}
観測による & 1 & 2 & 3 & 4 & \\
\hline
カッシーニ & 5 & 8 & 13 & 23 & \\
ボレッリ & $5\frac{2}{3}$ & $8\frac{2}{3}$ & 14 & $24\frac{2}{3}$ & \\
タウンリー（マイクロメーターによる） & $5.51$ & $8.78$ & $13.47$ & $24.72$ & 木星の半径 \\
フラムスティード（マイクロメーターによる） & $5.31$ & $8.85$ & $13.98$ & $24.23$ & を単位とする \\
フラムスティード（衛星の食による） & $5.578$ & $8.876$ & $14.159$ & $24.903$ & \\
\hline
周期から求めたもの & $5.578$ & $8.878$ & $14.168$ & $24.968$ & \\
\end{tabular}
\end{center}

\textbf{仮説 VI.}\wsp{}\textit{五つの主要な惑星、水星、金星、火星、木星、土星は、それぞれの軌道で太陽をめぐる。}

水星と金星が太陽のまわりを回ることは、それらの月のような位相から示される。満ちた顔で輝くときは太陽の向こう側にあり、半分のときは太陽と同じ方向にあり、三日月形のときは太陽のこちら側にあって、ときには、その円盤の上を斑点のように通り過ぎる。火星についても、太陽との合の近くでは満ちた顔、矩では膨らんだ形であることから、それが太陽をめぐることは確かである。木星と土星についても、それらがつねに満ちた位相であることから、同じことが示される。

\textbf{仮説 VII.}\wsp{}\textit{五つの主要な惑星と、（太陽の地球のまわりの、あるいは）地球の太陽のまわりの周期は、太陽からの平均の距離の3/2乗の比をなす。}

ケプラーが見いだしたこの比は、すべての人に認められている。惑星が地球のまわりを回るにせよ、太陽のまわりを回るにせよ、周期も軌道の大きさも同じである。そして周期の測定については、すべての天文学者が一致している。ところが軌道の大きさは、ケプラーとブリアルドゥスが、すべてのうちで最も注意深く観測から決めた。そして、周期に対応する平均の距離は、彼らが見いだした距離と、感じられるほどは違わず、それらの間にあるのがふつうである。次の表で見るとおりである。

\condpagelarge{惑星と地球の太陽からの平均の距離}

\begin{center}
\begin{tabular}{lcccccc}
 & 土星 & 木星 & 火星 & 地球 & 金星 & 水星 \\
ケプラーによる & 951000 & 519650 & 152350 & 100000 & 72400 & 38806 \\
ブリアルドゥスによる & 954198 & 522520 & 152350 & 100000 & 72398 & 38585 \\
周期による & 953806 & 520116 & 152399 & 100000 & 72333 & 38710 \\
\end{tabular}
\end{center}

水星と金星の太陽からの距離については、論ずる余地がない。それらは太陽からの離角によって決められるからである。また上の惑星の太陽からの距離についても、木星の衛星の食によって、すべての論争が取り除かれる。実際、その食によって、木星が投げる影の位置が定められ、そのようにして木星の太陽中心の経度が得られる。そして太陽中心の経度と地球中心の経度を照らし合わせることによって、木星の距離が決められる。

\textbf{仮説 VIII.}\wsp{}\textit{主要な惑星は、地球へ引いた半径では、時間にまったく比例しない面積を描くが、太陽へ引いた半径では、時間に比例する面積を描く。}

実際、地球に対しては、惑星はあるときは進み、あるときは留まり、またあるときは退く。ところが太陽に対しては、つねに進み、それもほぼ一様な運動であるが、近日点ではやや速く、遠日点ではやや遅く、そのため面積の描かれ方は一様である。この命題は天文学者にとってきわめてよく知られていて、とくに木星では、衛星の食によって最もよく示される。その食によって、この惑星の太陽中心の経度と太陽からの距離が決められると、私たちは述べた。

\textbf{仮説 IX.}\wsp{}\textit{月は、地球の中心へ引いた半径で、時間に比例する面積を描く。}

これは、月の見かけの運動を、その見かけの直径と照らし合わせることから明らかである。ただし、月の運動は、太陽の力によって、いくらか乱される。しかし、私は、この仮説では、誤差の感じられないごくわずかな点は、物理学者として無視する。
\section{第3編 万有引力の確立、惑星・月・潮汐の理論（命題I〜XXIV）}


% Section: 第3編 万有引力の確立、惑星・月・潮汐の理論（命題I〜XXIV）



\condpagelarge{命題 I. 定理 I.}

\textit{木星のまわりを回る惑星を、たえず直線運動から引きもどして、それぞれの軌道にとどめている力は、木星の中心に向かい、その中心からの場所の距離の平方に反比例する。}

命題の前半は、仮説Vと第1編 命題IIまたはIIIから明らかであり、後半は、仮説Vと同じ編 命題IVの系6から明らかである。

\condpagelarge{命題 II. 定理 II.}

\textit{主要な惑星を、たえず直線運動から引きもどして、それぞれの軌道にとどめている力は、太陽に向かい、その中心からの距離の平方に反比例する。}

命題の前半は、仮説VIIIと第1編 命題IIから明らかであり、後半は、仮説VIIと同じ編 命題IVから明らかである。この命題の後半は、遠日点が静止していることによって、最も正確に示される。というのは、二乗の比からのごくわずかなずれも、（第1編 命題XLVの系1により）一回転ごとにアプシスの運動を著しいものにし、多くの回転では途方もないものにするはずだからである。

\condpagelarge{命題 III. 定理 III.}

\textit{月をその軌道にとどめている力は、地球に向かい、その中心からの場所の距離の平方に反比例する。}

主張の前半は、仮説IXと第1編 命題IIまたはIIIから明らかであり、後半は、月の遠地点のきわめて遅い運動から明らかである。というのは、その運動は、一回転ごとにわずか3度の順行であって、無視できるからである。実際、第1編 命題XLVの系1により、月の地球の中心からの距離を $D$ とすれば、そのような運動を生じる力は $D^{2\frac{4}{243}}$ に反比例する、すなわち $D$ のその指数が $2\frac{4}{243}$ であるべきに反比例する。すなわち、距離の逆の二乗よりわずかに大きい比であるが、それは三乗よりも二乗に $60\frac{3}{4}$ 倍近い。このようなわずかな近づきは、当然無視してよい。それは（あとで述べるように）太陽の作用から生じるから、ここでは無視すべきである。そこで、地球に向かうその力は、$D^2$ に反比例することになる。このことは、この力を重力と比べるとき、つぎの命題のように、いっそう十分に明らかになる。

\condpagelarge{命題 IV. 定理 IV.}

\textit{月は地球に向かって重く、重力によってたえず直線運動から引きもどされ、その軌道にとどめられている。}

月の地球の中心からの平均の距離は、地球の半径で、大部分の天文学者によれば59、ウェンデリヌスによれば60、コペルニクスによれば $60\frac{1}{3}$、キルヒャーによれば $62\frac{1}{2}$、ティコによれば $56\frac{1}{2}$ である。しかし、ティコと、その屈折表に従う人々は、太陽と月の屈折を（光の本性にまったく反して）恒星のそれより、およそ4分か5分大きく定めたので、月の視差を同じ分だけ増した。すなわち視差全体のおよそ12分の1か15分の1である。この誤りを正せば、距離はおよそ地球の半径の61倍となり、他の人々が定めたのとほぼ同じになる。そこで、平均の距離を半径の60倍とし、月の周期を、恒星に対して27日7時間43分で完了するとし（天文学者が定めたとおり）、地球の周囲をパリのフィートで 123249600 とする（近ごろフランス人が測って決めたとおり）。そして、もし月からすべての運動を取り去り、それが軌道にとどめられているその力のすべてに押されて、地球に落ちるように放たれるなら、月は1分の時間に落ちて、パリのフィートで $15\frac{1}{12}$ を描く。これは、計算によって、第1編 命題XXXVIによって、あるいは（同じことになるが）同じ編 命題IVの注によって求められる。ところで、その力は地球に近づくにつれて、距離の逆二乗の比で増すから、地球の表面では月におけるより $60 \times 60$ 倍大きく、したがって、その力で私たちの地方で落ちる物体は、1分の時間にパリのフィートで $60 \times 60 \times 15\frac{1}{12}$、1秒の時間にパリのフィートで $15\frac{1}{12}$ を描くはずである。ところが、私たちの地方で重力によって落ちる物体は、ホイヘンスが振り子の実験とそれにもとづく計算で示したように、1秒の時間にパリのフィートで $15\frac{1}{12}$ を描く。したがって、月を軌道にとどめている力は、私たちが重力と呼びならわしている、まさにその力である。というのは、もし重力がそれと違うなら、物体は両方の力を合わせて地球へ向かい、2倍の速さで落ち、1秒の時間に落ちて、パリのフィートで $30\frac{1}{6}$ を描くはずで、それは経験にまったく反するからである。

この計算は、地球が静止しているという仮定にもとづく。実際、もし地球と月が太陽のまわりを動き、その間に共通の重心のまわりも回るなら、月と地球の中心の互いの距離は、地球の半径の $60\frac{1}{2}$ 倍になる。これは、（第1編 命題LXにより）計算をしてみれば明らかである。

\condpagelarge{命題 V. 定理 V.}

\textit{木星のまわりを回る惑星は木星に向かって重く、太陽のまわりを回る惑星は太陽に向かって重く、その重力によってたえず直線運動から引きもどされ、曲線の軌道にとどめられている。}

というのは、木星のまわりを回る惑星の木星のまわりの回転、水星や金星やそのほかの太陽のまわりを回る惑星の太陽のまわりの回転は、月の地球のまわりの回転と同じ種類の現象である。したがって仮説IIにより、同じ種類の原因に依存する。とくに、それらの回転が依存する力が、木星と太陽の中心に向かい、木星や太陽から遠ざかるにつれて、重力が地球から遠ざかるにつれて減るのと同じ比、同じ法則で減ることが証明されているのだから、そうである。

\textit{系1.}\wsp{}したがって、重力はすべての惑星に及ぶ。というのは、金星、水星やそのほかが、木星と同じ種類の物体であることを、誰も疑わないからである。たしかに、ホイヘンスの惑星（土星の衛星）は、木星の衛星が木星に向かって重いのと同じ論拠で、土星に向かって重い。そして、すべての引力は（運動の法則IIIにより）相互的だから、土星も逆にホイヘンスの惑星に向かって重い。同じ論拠で、木星はそのすべての衛星に向かって、地球は月に向かって、太陽はすべての主要な惑星に向かって重い。

\textit{系2.}\wsp{}一つ一つの惑星に向かう重力は、その中心からの場所の距離の平方に反比例する。

\condpagelarge{命題 VI. 定理 VI.}

\textit{すべての物体は、個々の惑星に向かって重く、同じ任意の惑星に向かうそれらの重さは、惑星の中心から等しい距離では、それぞれの物体にある物質の量に比例する。}

すべての重い物体の地球への落下は（空気のごくわずかな抵抗から生じる不等な遅れを除けば）等しい時間に行われることは、すでに以前から他の人々が観察している。そして時間の等しさは、振り子で最も正確に見てとれる。私は金、銀、鉛、ガラス、砂、食塩、木、水、小麦で、このことを試みた。私は、丸くて等しい二つの木の箱を用意した。一方を木で満たし、同じ重さの金を、（できるだけ正確に）もう一方の振動の中心につるした。箱は、等しい長さ11フィートの糸につるされて、重さ、形、空気の抵抗についてまったく等しい振り子をつくった。そして、等しい振動で並べて置かれると、きわめて長い間、一緒に行き、一緒にもどった。したがって、金の中の物質の量は（第2編 命題XXIVの系1と6により）、木の中の物質の量に対して、全体の金への起動力の作用の、全体の木へのその作用に対する比、すなわち重さの重さに対する比であった。そしてほかのものについても同じであった。同じ重さの物体で、物質の違いが、全体の物質の1000分の1より小さければ、これらの実験によって明らかに見いだせたであろう。

ところで、惑星に向かう重力の本性が、地球に向かう重力のそれと同じであることは、疑いがない。実際、これらの地上の物体を月の軌道にまで持ち上げ、月とともにすべての運動を取り去って、同時に地球に落ちるように放つと仮定する。すでに示したことから、これらは等しい時間に月と等しい距離を描くことは確実であり、したがって、それらは、月にある物質の量に対して、それぞれの重さの月の重さに対する比をなす。さらに、木星の衛星は、木星の中心からの距離の3/2乗の比をなす周期で回るから、それらの木星に向かう加速的な重力は、木星の中心からの距離の平方に反比例し、したがって、木星から等しい距離では、それらの加速的な重力は等しくなる。そのため、等しい時間に等しい高さから落ちて、等しい距離を描くであろう。それは、この地球で重い物体について起こるのと同じである。そして同じ論拠で、太陽のまわりを回る惑星も、太陽から等しい距離から放たれれば、太陽への落下で、等しい時間に等しい距離を描くであろう。ところで、不等な物体を等しく加速する力は、物体に比例する。すなわち、重さは惑星の中の物質の量に比例する。さらに、木星とその衛星の太陽に向かう重さが、それらの物質の量に比例することは、衛星のきわめて規則的な運動から明らかである（第1編 命題LXVの系3による）。というのは、もしこれらのうちのあるものが、他のものより、物質の量に応じて、太陽のほうへより強く引かれるなら、衛星の運動は（第1編 命題LXVの系2により）引力の不等によって乱されるはずだからである。もし（太陽から等しい距離で）ある衛星が、木星がその物質の量に応じて重いより、太陽に向かって、その物質の量に応じて、ある与えられた比、たとえば $d$ 対 $e$ で重いならば、太陽の中心と衛星の軌道の中心との距離は、太陽の中心と木星の中心との距離より、つねにほぼその比の平方根の比だけ大きいはずである。これは、私が少し計算して見いだしたことである。そしてもし衛星が、太陽に向かってその $d$ 対 $e$ の比で軽いなら、衛星の軌道の中心の太陽からの距離は、木星の中心の太陽からの距離より、その平方根の比だけ小さいはずである。したがって、もし太陽から等しい距離で、ある衛星の太陽に向かう加速的な重力が、木星の太陽に向かう加速的な重力より、重力全体のわずか1000分の1だけ大きいか小さいなら、衛星の軌道の中心の太陽からの距離は、木星の太陽からの距離より、全距離の $\frac{1}{2600}$、すなわち最も外側の衛星の木星の中心からの距離の5分の1だけ、大きいか小さいはずである。この軌道の離心率は、きわめて感じられるものになる。ところが衛星の軌道は木星に同心であり、したがって、木星と衛星の太陽に向かう加速的な重力は互いに等しい。そして同じ論拠で、土星とその伴星の太陽に向かう重さは、太陽から等しい距離では、それらの中の物質の量に比例する。また、月と地球の太陽に向かう重さは、ないか、あるいはそれらの質量に正確に比例する。

さらに、各惑星の個々の部分の、他の任意の惑星に向かう重さは、互いに、個々の部分の物質に比例する。実際、もしある部分が、物質の量から見積もるより多く重く、ほかの部分が少なく重いなら、惑星全体は、最も多く含む部分の種類に応じて、全体の物質の量から見積もるより多くか少なく重いであろう。しかし、それらの部分が外側にあるか内側にあるかは関係ない。実際、たとえば私たちのところにある地上の物体を月の軌道に持ち上げて、月の物体と比べると考えてみよ。もしこれらの重さが、月の外側の部分の重さに対して、それらの中の物質の量に比例し、内側の部分の重さに対しては、より大きいか小さい比をなすなら、それらは月全体の重さに対して、より大きいか小さい比をなすであろう。上に示したことに反する。

\textit{系1.}\wsp{}ここから、物体の重さは、その形や織り方に依存しない。実際、もし形によって変わりうるなら、等しい物質でも、形の違いに応じて、より大きいか小さいはずである。経験にまったく反する。

\textit{系2.}\wsp{}したがって、地球のまわりにあるすべての物体は、地球に向かって重く、地球の中心から等しく離れたすべての物体の重さは、それらの中の物質の量に比例する。実際、もしエーテルやほかの何らかの物体が、重力をまったくもたないか、物質の量から見積もるより軽いなら、それは他の物体と、物質の形によってしか違わないのだから、（仮説IIIにより）その形を変えることで、物質の量から見積もって最も重い物体と同じ状態のものに、しだいに変えることができ、逆に、最も重い物体は、その形を帯びることによって、重力をしだいに失うであろう。したがって、重さは物体の形に依存し、形とともに変わりうることになり、前の系で証明したことに反する。

\textit{系3.}\wsp{}したがって、真空は必然的に存在する。実際、もしすべての空間が満たされているなら、空気の領域を満たす流体の比重は、物質の密度が最大であるため、水銀や金やほかの最も密な物体の比重に少しも劣らないであろう。したがって、金もほかのどんな物体も、空気の中を落ちることはできない。というのは、流体の中の物体は、比重がより大きくない限り、決して沈まないからである。

\textit{系4.}\wsp{}重力は、磁力とは種類が違う。実際、磁気の引力は、引かれる物質に比例しない。ある物体はより強く、ある物体は弱く引かれ、大部分は引かれない。そして磁力は、物質の量に比して、重力よりはるかに大きい。しかも同じ物体で強めたり弱めたりすることができ、また磁石から遠ざかるにつれて、距離の二乗より大きい比で減る（第1編 命題LXXXVによる）。それは、力が、接触しているときのほうが、引くものがわずかでも離れたときよりはるかに強いからである。

\condpagelarge{命題 VII. 定理 VII.}

\textit{重力はすべての物体に及び、それぞれの物体の中の物質の量に比例する。}

すべての惑星が互いに重いことは、すでに前に証明した。また、一つ一つを別々に見て、それに向かう重力が、惑星の中心からの場所の距離の平方に反比例することも。そこから（第1編 命題LXIXとその系により）、すべてに向かう重力が、それらの中の物質に比例することが導かれる。

さらに、惑星 $A$ の部分がすべて任意の惑星 $B$ に向かって重く、各部分の重力が、全体の重力に対して、部分の物質の全体の物質に対する比をなし、すべての作用に（運動の法則IIIにより）等しい反作用があるなら、惑星 $B$ は逆に惑星 $A$ のすべての部分に向かって重く、その一つ一つの部分に向かう重力は、全体に向かう重力に対して、部分の物質の全体の物質に対する比をなす。\QED

\textit{系1.}\wsp{}ゆえに、惑星全体への重力は、個々の部分への重力から生じ、合成される。このことの例は、磁気や電気の引力に見られる。実際、全体への引力はすべて、個々の部分への引力から生じる。この事柄は、重力について、多くの小さな惑星が一つの球に合して、より大きな惑星をつくると考えれば、理解される。というのは、全体の力は、それを構成する部分の力から生じるはずだからである。もし誰かが、私たちのところにあるすべての物体が、この法則にしたがって互いに重いはずなのに、そのような重力がまったく感じられない、と異議を唱えるなら、私は次のように答える。これらの物体への重力は、地球全体への重力に対して、これらの物体の地球全体に対する比をなすから、感じられうるよりはるかに小さい。

\textit{系2.}\wsp{}物体の等しい粒子への重力は、粒子からの場所の距離の平方に反比例する。これは、第1編 命題LXXIVの系3から明らかである。

\condpagelarge{命題 VIII. 定理 VIII.}

\textit{互いに重い二つの球の物質が、中心から等しい距離の領域でどこも均質であるなら、一方の球の他方の球に向かう重さは、中心間の距離の平方に反比例する。}

私は、惑星全体への重力が、部分への重力から生じ、合成されること、また個々の部分への重力が、部分からの距離の平方に反比例することを見いだしたのち、この逆二乗の比が、多くの力から合成された全体の力でも正確に成り立つのか、それともほぼ成り立つにすぎないのかを、疑った。実際、その比は、大きな距離では十分に成り立っても、惑星の表面の近くでは、粒子の距離が不等で位置が似ていないために、著しく外れるということが、ありえたからである。しかしついに、第1編 命題LXXVとその系によって、ここで扱う命題の真実を理解した。

\textit{系1.}\wsp{}ここから、さまざまな惑星に向かう物体の重さを求め、互いに比べることができる。実際、惑星のまわりを円を描いて回る等しい物体の重さは、（第1編 命題IVにより）円の直径に正比例し、周期の二乗に反比例する。そして惑星の表面やそのほか中心から任意の距離での重さは、（この命題により）距離の二乗の逆の比で、大きくも小さくもなる。こうして、金星の太陽のまわりの周期 $224\frac{2}{3}$ 日、木星の最も外側の衛星の木星のまわりの周期 $16\frac{3}{4}$ 日、ホイヘンスの衛星の土星のまわりの周期 15日と $22\frac{2}{3}$ 時間、月の地球のまわりの周期 27日7時間43分を、金星の太陽からの平均の距離と、最も外側の木星の衛星の太陽中心の最大離角（フラムスティードの観測による、木星の太陽からの平均の距離で $8\minute 13\second$）と、土星の衛星の太陽中心の最大離角 $3\minute 20\second$、および、太陽の水平視差、すなわち太陽から見た地球の半径がおよそ $20\second$ であるという仮定のもとでの月の地球からの距離と、比べて計算すると、太陽、木星、土星、地球から等しい距離にある等しい物体の、それぞれ太陽、木星、土星、地球に向かう重さは、互いに $1$, $\frac{1}{1100}$, $\frac{1}{2360}$, $\frac{1}{28700}$ になることがわかった。ところで、太陽の平均の見かけの半径はおよそ $16\minute 6\second$ である。木星の太陽から見た半径について、フラムスティードは、衛星の食による木星の影の直径から、最も外側の衛星の離角に対して 1 対 $24.9$ と定めたから、その離角が $8\minute 13\second$ であるとき、太陽から見た木星の半径は $19\second\frac{3}{4}$ になる。土星の直径は、その環の直径に対して 4 対 9 であり、（フラムスティードの測定による）太陽から見た環の直径は $50\second$、したがって太陽から見た土星の半径は $11\second$ である。私はむしろ $10\second$ か $9\second$ と言いたい。光の屈折率の不等によって、土星の球はいくらか広がるからである。これらから計算すれば、太陽、木星、土星、地球の真の半径は、互いに 10000, 1063, 889, 208 になる。そこで、太陽、木星、土星、地球の中心から等しい距離にある等しい物体の、太陽、木星、土星、地球に向かう重さは、それぞれ 1, $\frac{1}{1100}$, $\frac{1}{2360}$, $\frac{1}{28700}$ であり、距離を増減すれば、重さは二乗の比で減増するから、それらの中心から 10000, 1063, 889, 208 の距離にある、したがってそれらの表面にある同じ等しい物体の、太陽、木星、土星、地球に向かう重さは、それぞれ 10000, $804\frac{1}{2}$, 536, $805\frac{1}{2}$ になる。月の表面での物体の重さが、地球の表面での物体の重さのほぼ半分であることは、あとで述べる。

\textit{系2.}\wsp{}ゆえに、地球と惑星の表面での、等しい物体の重さは、ほぼ、太陽から見た見かけの直径の平方根の比をなす。地球の太陽から見た直径は、まだ確定していない。私はこれを $40\second$ と仮定した。ケプラー、リッチョーリ、ウェンデリヌスの観測は、それよりはるかに大きいことを許さず、ホロックスとフラムスティードの観測は、それよりやや小さいことを示すように見えるからである。そして私は、過大に誤るほうをとった。もし、その直径と表面での平均の重力が、惑星の直径と、その表面での重力の中間にあるなら、土星、木星、火星、金星、水星の太陽から見た直径はおよそ $18\second$, $39\second\frac{1}{2}$, $8\second$, $28\second$, $20\second$ だから、地球の直径はおよそ $24\second$、したがって太陽の視差はおよそ $12\second$ となり、ホロックスとフラムスティードがほぼ定めたとおりになる。しかし、やや大きな直径のほうが、この系の規則とよく合う。

\textit{系3.}\wsp{}個々の惑星の中の物質の量もわかる。実際、その量は、惑星の、自分から等しい距離での力に比例する。すなわち、太陽、木星、土星、地球で、それぞれ 1, $\frac{1}{1100}$, $\frac{1}{2360}$, $\frac{1}{28700}$ である。もし太陽の視差が $20\second$ より小さいと仮定するなら、地球の中の物質の量は、3乗の比で減らさなければならない。

\textit{系4.}\wsp{}惑星の密度もわかる。実際、等しく均質な物体の、均質な球の表面での球への重さは、第1編 命題LXXIIにより、球の直径に比例する。したがって、不均質な球の密度は、重さを直径で割ったものに比例する。ところで、太陽、土星、木星、地球の真の直径は、互いに 10000, 889, 1063, 208 であり、それらへの重さは 10000, 536, $804\frac{1}{2}$, $805\frac{1}{2}$ であったから、密度は 100, 60, 76, 387 に比例する。ここで求めた地球の密度は、太陽の視差に依存せず、月の視差によって定められ、したがってここでは正しく決まる。ゆえに、太陽は木星よりやや密で、地球は太陽よりはるかに密である。

\textit{系5.}\wsp{}ところで、惑星の密度は、互いに、ほぼ、太陽からの距離の比と、太陽から見た見かけの直径の平方根の比との、合成された比をなす。すなわち、土星、木星、地球、月の密度 60, 76, 387, 700 は、ほぼ、距離の逆数 $\frac{1}{9538}$, $\frac{1}{5201}$, $\frac{1}{1000}$, $\frac{1}{1000}$ に、見かけの直径 $18\second$, $39\second\frac{1}{2}$, $40\second$, $11\second$ の平方根を掛けたものに比例する。実際、私たちは系2で、惑星の表面での重力は、ほぼ、太陽から見た見かけの直径の平方根の比をなすと述べ、補助定理IVで、密度は、その重力を真の直径で割ったものに比例すると述べた。したがって、密度は、ほぼ、見かけの直径の平方根を真の直径で割ったもの、すなわち、惑星の太陽からの距離に、見かけの直径の平方根を掛けたものの逆数に比例する。ゆえに神は、各惑星が密度の度合いに応じて、太陽の熱をより多くまたはより少なく受けるように、惑星を太陽からのさまざまな距離に置いた。私たちの水は、もし地球が土星の軌道に置かれたなら凍り、水星の軌道に置かれたなら、たちまち蒸気となって去るであろう。実際、熱が比例する太陽の光は、水星の軌道では、私たちのところの7倍の密度であり、私は温度計で、夏の太陽の熱の7倍で水が沸騰することを、試して知った。そして、水星の物質が熱に適合していて、したがって私たちのこの物質より密であることは疑いない。それは、すべての密な物質が、自然の作用をなしとげるのに、より大きな熱を必要とするからである。

\condpagelarge{命題 IX. 定理 IX.}

\textit{重力は、惑星の表面から下へ進むにつれて、ほぼ中心からの距離の比で減る。}

もし惑星の物質が、密度について一様ならば、この命題は、第1編 命題LXXIIIにより、正確に成り立つ。したがって、誤差は、密度の不均等から生じうる分だけである。

\condpagelarge{命題 X. 定理 X.}

\textit{惑星の天空での運動は、きわめて長く保たれうる。}

第2編 命題XLの注で、私たちの空気中を自由に動いて自分の半径の長さを描く、凍った水の球は、空気の抵抗によって、運動の $\frac{1}{3200}$ を失うことが示された。そして同じ比は、（第2編 命題XLにより）どんなに大きく速い球でも、ほぼ成り立つ。ところで、私たちの地球の球が、全体が水から成る場合より密であることは、次のように推論する。もしこの球が全体水なら、水より希薄なものは、比重が小さいために、浮き出て水に浮かぶであろう。そのため、水で四方から覆われた土の球は、もし水より希薄なら、どこかに現れ出て、そこから流れ下る水はすべて、反対側の領域に集まるであろう。そして、大部分が海に囲まれた私たちの地球についても、理由は同じである。もし地球がより密でないなら、海から現れ出て、軽さの度合いに応じてその一部が水の上に出て、すべての海は反対側の領域に流れ込むであろう。同じ論拠で、太陽の黒点は、それが浮かぶ太陽の輝く物質より軽い。そして惑星がどのようにでもできたとき、全体が流体であった間、すべての重い物質は中心へ向かった。そこから、地球の地表の土は水のほぼ2倍重く、少し下の鉱山では水のほぼ3倍、4倍、あるいは5倍重いのが見いだされるから、地球の物質全体の量は、全体が水から成る場合の、おそらく5倍か6倍であろう。とくに、地球が木星よりほぼ5倍密であることは、すでに前に示された。したがって、もし木星が水よりやや密であるなら、木星は、320の半径の距離を描く21日の間に、私たちの空気と同じ密度の媒質の中で、運動のおよそ10分の1を失うであろう。ところが、媒質の抵抗は、重さと密度の比で減るから、水は水銀より $13\frac{2}{3}$ 倍軽く、同じ比でより小さく抵抗し、空気は水より800倍軽く、同じ比でより小さく抵抗する。そして、惑星が動く媒質の重さが限りなく小さくなる天空にのぼれば、抵抗はほとんどやむ。

\condpagelarge{命題 XI. 定理 XI.}

\textit{地球、太陽、およびすべての惑星の、共通の重心は、静止している。}

実際、その重心は（法則の系4により）静止するか、直線上を一様に進む。ところが、その重心がつねに進むなら、世界の中心も動くことになり、仮説IVに反する。

\condpagelarge{命題 XII. 定理 XII.}

\textit{太陽は絶え間なく動かされているが、すべての惑星の共通の重心から遠く離れることはない。}

実際、命題VIIIの系3により、太陽の中の物質は木星の中の物質に対して 1100 対 1 であり、木星の太陽からの距離は、太陽の半径に対して、ほぼ同じ比であるから、木星と太陽の共通の重心は、ほぼ太陽の表面に落ちる。同じ論拠で、太陽の中の物質は土星の中の物質に対して 2360 対 1 であり、土星の太陽からの距離は太陽の半径に対してそれよりやや小さい比であるから、土星と太陽の共通の重心は、太陽の表面のやや下の点に落ちる。そして同じ計算を続ければ、もし地球とすべての惑星が太陽の一方の側に集まるなら、すべての共通の重心は、太陽の中心から、太陽の直径ひとつ分ほども離れない。ほかの場合には、重心の距離はつねに小さい。したがって、その重心は絶えず静止しているから、太陽は、惑星のさまざまな位置に応じて、あらゆる向きに動くが、その重心から遠く離れることはない。

\textit{系.}\wsp{}ここから、地球、太陽、すべての惑星の共通の重心を、世界の中心とみなさなければならない。実際、地球、太陽、すべての惑星は互いに重く、したがって、重力の力に応じて、運動の法則にしたがって絶えず動かされているから、これらの動く中心を、静止する世界の中心とみなすことはできないことは明らかである。もし、すべての物体が最も重く向かう物体を中心に置くべきなら（俗説のように）、その特権は太陽に与えられるはずである。ところが、太陽は動くから、静止する点を選ばなければならず、太陽の中心はそこからできるだけ離れず、もし太陽がより密で大きければ動きが小さいから、そこからさらに離れないであろう。

\condpagelarge{命題 XIII. 定理 XIII.}

\textit{惑星は、太陽の中心に焦点をもつ楕円の上を動き、その中心へ引いた半径で、時間に比例する面積を描く。}

私たちは上で、これらの運動を現象から論じた。いま運動の原理が知られたから、これらから、天体の運動を先験的に導く。惑星の太陽に向かう重さは、太陽の中心からの距離の平方に反比例するから、もし太陽が静止していて、ほかの惑星が互いに作用しないなら、それらの軌道は、太陽を共通の焦点にもつ楕円であり、面積は時間に比例して描かれるであろう（第1編 命題IとXI、および命題XIIIの系1による）。ところで、惑星の互いの作用はごく小さく（無視できるほどで）、動く太陽のまわりの楕円上の惑星の運動を、（第1編 命題LXVIにより）静止した太陽のまわりでこれらの運動が行われる場合よりも、いっそう乱さない。

たしかに、木星の土星への作用は、まったく無視してよいわけではない。実際、木星への重力は、太陽への重力に対して（等しい距離で）1 対 1100 であり、したがって、木星と土星の合では、土星の木星からの距離が土星の太陽からの距離に対してほぼ 4 対 9 だから、土星の木星への重力は、土星の太陽への重力に対して 81 対 $16 \times 1100$、すなわちおよそ 1 対 217 になる。しかし、木星への重力から生じる、土星の太陽のまわりの運動のすべての誤差は、土星の軌道の焦点を、木星と太陽の共通の重心に置くことによって（第1編 命題LXVIIにより）ほぼ避けることができ、したがって、最大のときでも2分をほとんど超えない。そして、木星と土星の合では、太陽の土星への、木星の土星への、木星の太陽への加速的な重力は、ほぼ 16, 81, $\frac{16 \times 81 \times 2360}{25}$、すなわち 122342 であり、したがって、太陽の土星への重力と、木星の土星への重力との差は、木星の太陽への重力に対して 65 対 122342、すなわち 1 対 1867 になる。この差に、土星が木星の運動を乱す最大の効力は比例するから、木星の軌道の乱れは、土星のそれよりはるかに小さい。残りの軌道の乱れは、さらにはるかに小さい。

\condpagelarge{命題 XIV. 定理 XIV.}

\textit{軌道の遠日点と交点は静止している。}

遠日点は、第1編 命題XIによって静止し、軌道の面も、同じ編 命題Iによって静止し、面が静止すれば交点も静止する。しかし、回る惑星と彗星の互いの作用から、いくつかの不等が生じるが、それは小さいために無視できる。

\textit{系1.}\wsp{}恒星も静止している。それは、遠日点と交点に対する与えられた位置を保つからである。

\textit{系2.}\wsp{}したがって、地球の年周運動から生じる感じられる視差が恒星にはないから、それらの力は、物体の途方もない距離のために、私たちの体系の領域では感じられる効果をまったく生じない。

\condpagelarge{命題 XV. 問題 I.}

\textit{軌道の主軸を求めること。}

これらは、（第1編 命題XVにより）周期の3/2乗の比にとり、ついでそれぞれを、（第1編 命題LXにより）太陽と各周回惑星の質量の和の、その和と太陽との間にある二つの比例中項の第一のものに対する比で、増さなければならない。

\condpagelarge{命題 XVI. 問題 II.}

\textit{軌道の離心率と遠日点を求めること。}

問題は、第1編 命題XVIIIによって解かれる。

\condpagelarge{命題 XVII. 定理 XV.}

\textit{惑星の日周運動は一様であり、月の秤動は、その日周運動から生じる。}

これは、運動の法則Iと、第1編 命題LXVIの系22から明らかである。ところで、月は自分の軸のまわりを一様に回転し、その日は一か月であるから、その面は、つねにその軌道の遠いほうの焦点を向き、したがって、その焦点の位置に応じて、地球から両側にそれる。これが経度の秤動である。緯度の秤動は、月の軸の軌道面に対する傾きから生じる。そして、これらがそうであることは、現象から明らかである。

\condpagelarge{命題 XVIII. 定理 XVI.}

\textit{惑星の軸は、その軸に垂直に引いた直径より小さい。}

惑星は、すべての円運動の日周運動が取り去られれば、部分の重力がどこでも等しいために、球形をとるはずである。その円運動によって、軸から離れる部分は、赤道に沿って上ろうとする。したがって、もし物質が流体なら、その赤道への上昇によって直径を大きくし、極への下降によって軸を小さくするであろう。こうして木星の直径は、（カッシーニとフラムスティードの観測が一致して）極の間のほうが、東から西よりも短いことがわかる。同じ論拠で、もし私たちの地球が赤道の下で極よりやや高くなければ、海は極のほうへ沈んで、赤道に沿って上り、そこをすべて水没させるであろう。

\condpagelarge{命題 XIX. 問題 III.}

\textit{惑星の軸の、それに垂直な直径に対する比を求めること。}

\pngright{410.png}{}{}{}

この問題を解くには、多くの計算が必要であり、これは規則よりも例によって、より容易に学ばれる。そこで計算を行って、私は、第1編 命題IVにより、地球の赤道の下の部分の、日周運動から生じる遠心力は、重力に対して 1 対 $290\frac{4}{5}$ であることを見いだした。そこで、$APBQ$ を、楕円が短軸 $PQ$ のまわりに回転して生じる地球の形とし、$ACQqca$ を、極 $Qq$ から中心 $Cc$ へ、そこから赤道 $Aa$ へ進む、水で満たした管とする。すると、管の脚 $ACca$ の中の水の重さは、もう一方の脚 $QCcq$ の中の水の重さに対して 291 対 290 でなければならない。それは、円運動から生じる遠心力が、重さの291の部分のうち一つを支えて取り去り、もう一方の脚では、重さ290が残りの部分を支えるからである。さらに、（第1編 命題XCIの系2から）計算して、私は次のことを見いだした。もし地球が一様な物質から成り、すべての運動を取り去られ、その軸 $PQ$ の直径 $AB$ に対する比が 100 対 101 であるなら、地球に向かう場所 $Q$ での重力は、中心 $C$、半径 $PC$ または $QC$ で描いた球に向かう同じ場所 $Q$ での重力に対して、$126\frac{2}{15}$ 対 $125\frac{2}{15}$ である。そして同じ論拠で、楕円 $APBQ$ が軸 $AB$ のまわりに回転して描かれる回転楕円体に向かう場所 $A$ での重力は、中心 $C$、半径 $AC$ で描いた球に向かう同じ場所 $A$ での重力に対して、$125\frac{2}{15}$ 対 $126\frac{2}{15}$ である。ところで、地球に向かう場所 $A$ での重力は、前述の回転楕円体と球に向かう重力の比例中項である。というのは、球は、直径 $PQ$ を 101 対 100 の比で小さくすることによって、地球の形に変わり、この形は、二つの直径 $AP$, $PQ$ に垂直な第三の直径を同じ比で小さくすることによって、前述の回転楕円体に変わり、$A$ での重力は、どちらの場合も、ほぼ同じ比で減るからである。したがって、中心 $C$、半径 $AC$ で描いた球に向かう $A$ での重力は、地球に向かう $A$ での重力に対して 126 対 $125\frac{1}{2}$ であり、中心 $C$、半径 $QC$ で描いた球に向かう場所 $Q$ での重力は、中心 $C$、半径 $AC$ で描いた球に向かう場所 $A$ での重力に対して、（第1編 命題LXXIIにより）直径の比、すなわち 100 対 101 である。これら三つの比、$126\frac{2}{15}$ 対 $125\frac{2}{15}$、$125\frac{1}{2}$ 対 126、100 対 101 を合わせれば、地球に向かう場所 $Q$ での重力の、地球に向かう場所 $A$ での重力に対する比は、$126 \times 126 \times 100$ 対 $125 \times 125\frac{1}{2} \times 101$、すなわち 501 対 500 になる。

さて、第1編 命題XCIの系3により、管のどちらの脚 $ACca$, $QCcq$ の中でも重力は、場所の地球の中心からの距離に比例するから、これらの脚を横断する平行な面で全体に比例する部分に分けるならば、脚 $ACca$ の個々の部分の重さは、もう一方の脚の同数の部分の重さに対して、大きさと加速的な重力との積の比をなす。すなわち、101 対 100 と 500 対 501、すなわち 505 対 501 である。したがって、もし脚 $ACca$ の各部分の、日周運動から生じる遠心力が、同じ部分の重さに対して 4 対 505 で、各部分の重さを505の部分に分けたうちの4つを取り去るなら、重さは両方の脚で等しいままで、流体は平衡に立つであろう。ところが、各部分の遠心力は、その重さに対して 1 対 290 である。すなわち、向心力は、重さの $\frac{4}{505}$ の部分であるべきなのに、わずか $\frac{1}{290}$ の部分である。そこで、三数法によって、もし遠心力 $\frac{4}{505}$ が、脚 $ACca$ の水の高さを、脚 $QCcq$ の水の高さより、全体の高さの100分の1だけ高くするなら、遠心力 $\frac{1}{290}$ は、脚 $ACca$ の高さの超過を、もう一方の脚 $QCcq$ の高さの $\frac{3}{689}$ の部分だけにする、と私は言う。ゆえに、地球の赤道に沿う直径は、極を通る直径に対して 692 対 689 である。したがって、地球の平均の半径は、近ごろのフランス人の測定によれば、パリのフィートで 19615800、すなわち 3923 マイル（1マイルを 5000 フィートの測定とする）であるから、地球は、極より赤道のほうが、パリのフィートで 85200、すなわち17マイル高い。

もし惑星が、地球より大きいか密であるか、小さいか希薄であるなら、日周回転の周期が同じ限り、遠心力の重力に対する比は変わらず、したがって、極の間の直径の赤道に沿う直径に対する比も変わらない。しかし、もし日周運動が任意の比で速くなるか遅くなるなら、遠心力はその二乗の比で増減し、したがって直径の差も、同じ二乗の比で増える。そこで、地球は恒星に対して23時間56分で回転し、木星は9時間56分で回転して、時間の平方が 29 対 5 であるから、木星の直径の差は、その短い直径に対して $\frac{29 \times 3}{5 \times 689}$ 対 1、すなわち 1 対 $39\frac{3}{5}$ になる。したがって、木星の東から西へ引いた直径は、その極の間の直径に対して、ほぼ $40\frac{3}{5}$ 対 $39\frac{3}{5}$ である。これらは、惑星の物質が一様であるという仮定のもとでのことである。実際、もし物質が周よりも中心に密なら、東から西へ引いた直径は、さらに大きくなるであろう。

\condpagelarge{命題 XX. 問題 IV.}

\textit{さまざまな地方での物体の重さを求め、互いに比べること。}

水の管 $ACQqca$ の不等な脚の重さは等しく、脚全体に比例し、全体の中で似た位置にある部分の重さは、互いに全体の重さに比例し、したがって互いに等しい。そこで、等しく、脚の中で似た位置にある部分の重さは、脚に反比例する。すなわち 692 対 689 に反比例する。管の脚の中で似た位置にある、均質で等しいどんな物体についても、理由は同じである。これらの重さは、脚に反比例する。すなわち、物体の地球の中心からの距離に反比例する。したがって、もし物体が管の最上部、すなわち地球の表面にあるなら、それらの重さは、互いに、中心からの距離に反比例する。そして同じ論拠で、地球の表面全体のほかのどの地方でも、重さは、場所の中心からの距離に反比例し、したがって、地球が回転楕円体であるという仮定のもとで、比例で与えられる。そこから、次の定理がつくられる。赤道から極へ進むにつれての重さの増加は、ほぼ、緯度の2倍の正矢に、あるいは同じことだが、緯度の正弦の平方に比例する。たとえば、パリの緯度は $48$ 度 $45\minute$、カーボヴェルデ近くのゴレ島のそれは $14$ 度 $15\minute$、ギアナの海岸のカイエンヌのそれはおよそ $5$ 度、極の下の地方では $90$ 度である。これらの2倍の $97\frac{1}{2}$ 度、$28\frac{1}{2}$ 度、$10$ 度、$180$ 度の正矢は、11305, 1211, 152, 20000 である。そこで、極での重力の赤道の下での重力に対する比が 692 対 689、極の下での重力の超過の赤道の下での重力に対する比が 3 対 689 であるから、パリ、ゴレ島、カイエンヌの重力の超過は、赤道の下の重力に対して、$\frac{3 \times 11305}{20000}$, $\frac{3 \times 1211}{20000}$, $\frac{3 \times 152}{20000}$ 対 689、すなわち 33915, 3633, 456 対 13780000 になる。したがって、これらの場所での全体の重力は、互いに 13813915, 13783633, 13780456, 13780000 になる。

そこで、等しい時間に振動する振り子の長さは重力に比例し、パリで1秒ごとに振動する振り子の長さは、パリのフィートで3と $\frac{17}{24}$ インチであるから、ゴレ島、カイエンヌ、赤道の下で、1秒ごとに振動する振り子の長さは、パリの振り子の長さより、それぞれ $\frac{81}{1000}$, $\frac{89}{1000}$, $\frac{90}{1000}$ インチだけ短くなる。これらはすべて、地球が一様な物質から成るという仮定のもとでのことである。実際、もし物質が、表面より中心でやや密なら、これらの超過はやや大きくなる。それは、もし中心で余分な物質、すなわちそこの密度を大きくしている物質を取り去って、別に考えるなら、残りの一様な密度の地球に向かう重力は、重さのある物体の中心からの距離に反比例し、余分な物質に向かう重力は、その物質からの距離の平方に、ほぼ反比例するからである。したがって、赤道の下での重力は、その余分な物質に向かうものが、上の計算から見積もるより小さくなり、したがって地球は、そこで重力が不足するために、前に定めたより少し高く上るであろう。ところが、フランス人は実験で、1秒ごとに振動する振り子の長さが、パリのほうがゴレ島より10分の1インチ、カイエンヌより8分の1インチ長いことを見いだした。これらの差は、上の計算から出た差 $\frac{81}{1000}$ と $\frac{89}{1000}$ よりやや大きい。したがって（これらの粗い観測が十分信頼できるなら）、地球は、上の計算より赤道の下でいくらか高く、鉱山の表面近くよりも、中心でより密であることになる。もし、赤道の重力を超える、これらの北の地方での重力の超過が、いっそう注意深く行った実験で、ついに正確に決められ、ついでその超過がどこでも、緯度の2倍の正矢の比でとられるなら、そのとき、普遍的な尺度も、さまざまな地方で示される等しい振り子による時間の均差も、さらに地球の直径の比と、その中心での密度も、その密度が周へ向かって一様に減るという仮定のもとで、決まる。この仮定は、正確ではないが、計算を行うためにとることができる。

\condpagelarge{命題 XXI. 定理 XVIII.}

\textit{春分点は後退し、地球の軸は、一回転ごとに、揺れながら二度黄道に傾き、二度もとの位置にもどる。}

これは、第1編 命題LXVIの系20から明らかである。しかし、この揺れの運動はきわめて小さいはずで、ほとんど感じられないか、あるいはまったく感じられない。
\condpagelarge{命題 XXII. 定理 XIX.}

\textit{月のすべての運動と、運動のすべての不等は、述べた原理から導かれる。}

大きな惑星は、太陽のまわりを運ばれるあいだに、そのまわりを回る別の小さな惑星を連れていくことができ、それらの小さな惑星は、大きな惑星の中心に焦点をもつ楕円の上を回るはずであることは、第1編 命題LXVにより明らかである。ところが、太陽の作用によって、それらの運動はさまざまに乱され、私たちの月に認められる不等におかされる。月は（命題LXVIの系2, 3, 4, 5により）、離心率の運動が妨げる分を除いて、朔望では矩におけるより速く動き、地球に引いた半径で、時間に対してより大きな面積を描き、軌道の曲がりが小さく、したがって地球により近づく。離心率は（命題LXVIの系9により）、月の遠地点が朔望にあるとき最大で、矩にあるとき最小である。そこから、月は近地点では速くて私たちに近く、遠地点では遅くて、朔望では矩におけるより遠くなる。さらに遠地点は進み、交点は後退するが、その運動は不均等である。遠地点は（命題LXVIの系7, 8により）朔望では速く進み、矩ではゆっくり後退し、進みが後退を超える分によって、毎年順行の向きに運ばれる。交点は（命題LXVIの系11により）朔望では静止し、矩で最も速く後退する。また月の最大の緯度は、矩のほうが（命題LXVIの系10により）朔望より大きく、平均運動は、地球が近日点にあるときのほうが（系6により）遠日点にあるときより速い。これらが、天文学者が認めた、より著しい不等である。

そのほかに、まだ観測されていない不等もあって、それによって月の運動はひどく乱され、これまでどんな法則でも、確かな規則にまとめられなかった。実際、月の遠地点と交点の一時間あたりの運動（速さ）、およびそれらの均差、さらに、朔望での最大の離心率と矩での最小の離心率との差、そして出差と呼ばれる不等は、（命題LXVIの系14により）毎年、太陽の見かけの直径の3乗の比で増減する。そして出差はさらに、ほぼ矩の間の時間の二乗の比で（第1編 補助定理Xの系1, 2および命題LXVIの系16により）増減する。だが、この不等は、天文学の計算では、ふつう月の補正量（プロスタファイレシス）に帰せられ、これと混同される。

\condpagelarge{命題 XXIII. 問題 V.}

\textit{木星と土星の衛星の不等な運動を、月の運動から導くこと。}

私たちの月の運動から、木星の月、すなわち衛星の運動に類似した運動は、次のようにして導かれる。木星の最も外側の衛星の交点の平均運動は、私たちの月の交点の平均運動に対して、地球の太陽のまわりの周期の木星の太陽のまわりの周期に対する比の二乗と、衛星の木星のまわりの周期の月の地球のまわりの周期に対する比（単純な比）との、合成された比をなす（第1編 命題LXVIの系16により）。したがって、この交点は100年で、先行する向きに $9$ 度 $34\minute$ 進む。内側の衛星の交点の平均運動は、同じ系により、この運動に対して、それらの周期のこの衛星の周期に対する比をなし、そこから与えられる。また、各衛星の遠地点の後続の向きの運動は、その交点の先行する向きの運動に対して、私たちの月の遠地点の運動の、その交点の運動に対する比（同じ系による）をなし、そこから与えられる。ただし、こうして求めた遠地点の運動は、ここで述べる暇のない理由によって、5 対 9、すなわちおよそ 1 対 2 の比で減らさなければならない。

各衛星の交点と遠地点の最大の均差は、ほぼ、それぞれ月の交点と遠地点の最大の均差に対して、前の均差の一回転の時間における衛星の交点と遠地点の運動が、あとの均差の一回転の時間における月の交点と遠地点の運動に対するのと同じ比をなす。木星から見た衛星の出差は、月の出差に対して、同じ系により、衛星と月の周期における交点の全運動の互いの比をなし、したがって最も外側の衛星では $6\second 22\third$ を超えない。これらの不等が小さく、運動が遅いために、衛星の運動はきわめて規則的であり、近ごろの天文学者は、交点にはまったく運動がないとするか、きわめて遅く逆行すると主張する。実際、フラムスティードは、自分の観測をカッシーニの観測と照らし合わせて、交点がゆっくり後退することを見いだした。

\condpagelarge{命題 XXIV. 定理 XX.}

\textit{海の満ち引きは、太陽と月の作用から生じるはずである。}

海が、月日でも太陽日でも、毎日2回ふくらみ、2回引くはずであることは、第1編 命題LXVIの系19から明らかである。また、深く自由な海では、水の最大の高さは、光体が場所の子午線に達してから、6時間より短い間をおいて起こる。それは、大西洋とエチオピア海の東の全域（フランスと喜望峰の間）でも、太平洋のチリとペルーの海岸でも見られる。これらすべての海岸では、満潮はおよそ3時ごろに起こる。ただし、運動が浅い場所を通って伝わるときは、いくらか遅れる。私は、時刻を、どちらの光体も場所の子午線に達したときから、地平線の下でも上でも数え、月日の時刻とは、月が見かけの日周運動で場所の子午線にもどるまでの時間の24分の1をいう。

ところで、二つの光体が引き起こす二つの運動は、はっきりとは見分けられず、ある混合した運動をつくる。光体の合や衝では、それらの効果が合わさって、最大の満ち引きが合成される。矩では、太陽は、月が水を押し下げるところで水を持ち上げ、太陽が持ち上げるところで押し下げる。そして、効果の差から、すべての中で最小の潮が生じる。そして、経験の示すところでは、月の効果は太陽の効果より大きいから、水の最大の高さは月の3時ごろに起こる。朔望と矩以外では、最大の潮は、月の力だけなら、つねに月の3時に、太陽の力だけなら、太陽の3時に起こるはずだが、合わさった力では、月の3時により近い、中間のある時刻に起こる。したがって、月が朔望から矩へ移るとき、太陽の3時が月の3時に先立つなら、水の最大の高さも月の3時に先立ち、その間隔は、月の八分点のやや後で最大になる。そして、等しい間隔で、最大の潮は、月が矩から朔望へ移るときには、月の3時に続く。これらは、開けた海でのことである。河口では、他が同じなら、より大きな満ち潮は、より遅くそこへ達する。

ところで、光体の効果は、それらの地球からの距離に依存する。実際、距離が小さいほど効果は大きく、大きいほど小さく、それは見かけの直径の3乗の比で変わる。したがって、冬には太陽は近地点にあって、より大きな効果を生じ、朔望での潮を、夏よりやや大きく、矩での潮を、他が同じなら、やや小さくする。そして、月は、近地点にあるとき、遠地点にある15日前後よりも、毎月より大きな潮を起こす。そのため、続く朔望で、二つの最大の潮が互いに続くことはない。

また、どちらの光体の効果も、その赤緯、すなわち赤道からの距離に依存する。実際、もし光体が極に置かれるなら、それは水の各部分を絶えず、作用の強まりや弱まりなしに引くから、運動のかわるがわるをまったく起こさない。したがって、光体は、赤道から極のほうへ遠ざかるにつれて、その効果をしだいに失い、そのため、至の朔望では、分の朔望よりも小さな潮を起こす。ところが、至の矩では、分の矩よりも大きな潮を起こす。それは、いま赤道にある月の効果が、太陽の効果を最も大きく超えるからである。したがって、最大の潮は、光体の朔望で、最小の潮は矩で、どちらの分点の時期の前後にも起こる。そして、朔望の最大の潮には、つねに矩の最小の潮が続くことが、経験で確かめられている。また、太陽の地球からの距離が、冬のほうが夏より小さいために、最大の潮と最小の潮は、春分の前に起こることのほうが、あとに起こることより多く、秋分のあとに起こることのほうが、前に起こることより多い。

\pngright{419.png}{}{}{}

光体の効果は、場所の緯度にも依存する。$ApEP$ を、深い水で四方を覆われた地球とし、$C$ をその中心、$Pp$ を極、$AE$ を赤道、$F$ を赤道以外の任意の場所、$Ff$ をその場所の緯線、$Dd$ を、赤道の反対側のそれに対応する緯線、$L$ を月が3時間前に占めていた場所、$H$ をその真下にある地球の場所、$h$ をそれと反対の場所、$K$, $k$ をそこから 90 度離れた場所、$CH$, $Ch$ を地球の中心から測った海の最大の高さ、$CK$, $Ck$ を最小の高さとする。そして軸 $Hh$, $Kk$ で楕円を描き、ついでこの楕円を長軸 $Hh$ のまわりに回転して回転楕円体 $HPKhpk$ を描くなら、これはほぼ海の形を表し、$CF$, $Cf$, $CD$, $Cd$ が、場所 $F$, $f$, $D$, $d$ での海の高さになる。さらに、前述の楕円の回転で、任意の点 $N$ が円 $NM$ を描き、緯線 $Ff$, $Dd$ と任意の場所 $R$, $T$ で、赤道 $AE$ と $S$ で交わるなら、$CN$ は、この円上のすべての場所 $R$, $S$, $T$ での海の高さである。ここから、任意の場所 $F$ の日周運動では、満ち潮は、月が地平線の上の子午線に達して3時間後に、$F$ で最大になる。そののち、月が没して3時間後に、$Q$ で干潮が最大になる。ついで、月が地平線の下の子午線に達して3時間後に、$f$ で満ち潮が最大になり、最後に、月が昇って3時間後に、$Q$ で干潮が最大になる。そして、あとの $f$ での満ち潮は、先の $F$ での満ち潮より小さい。実際、海全体は、二つの半球状の波に分かれる。一つは半球 $KHkC$ にあって北に向かい、もう一つは反対の半球 $KhkC$ にある。したがって、これらを、北の波、南の波と呼んでよい。これらの波は、つねに互いに反対側にあって、月の12時間の間隔をおいて、かわるがわる、それぞれの場所の子午線にやってくる。そして北の地方は北の波により多く関与し、南の地方は南の波により多く関与するから、そこから、赤道以外の各場所で、かわるがわる大きく、小さい潮が生じる。そして、月が場所の頂点のほうへ傾いているとき、大きいほうの潮は、月が地平線の上の子午線に達して、およそ3時間後に起こり、月が赤緯を変えると、より小さい潮に変わる。そして、潮の最大の差は、至のころに起こる。とくに、月の昇交点が白羊宮のはじめにあるときはそうである。こうして、経験から次のことが確かめられている。冬には朝の潮が夕方の潮を、夏には夕方の潮が朝の潮を超える。プリマスでは高さがおよそ1フィート、ブリストルでは15インチである。これはコールプレスとスタームの観測による。

ところで、ここまでに述べた運動は、水の往復の力によって、いくらか変わる。この力によって、海の潮は、光体の作用がやんでも、しばらく続くことができる。この与えられた運動の保存は、かわるがわる起こる潮の差を小さくし、朔望の直後の潮をより大きくし、矩の直後の潮をより小さくする。そのため、プリマスとブリストルで、かわるがわる起こる潮は、1フィートや15インチの高さ以上には、互いにあまり違わず、また、同じ港で最大の潮は、朔望からの最初のものではなく、3番目のものになる。さらに、すべての運動は、浅瀬を通るときに遅れ、そのため、ある海峡や河口では、最大の潮は、朔望から4番目、あるいは5番目になる。

さらに、潮が大洋から異なる海峡を通って同じ港へ伝わり、ある海峡を他より早く通ることも起こりうる。この場合、同じ潮が、次々に着く二つか多くのものに分かれて、さまざまな種類の新しい運動を合成しうる。いま、等しい二つの潮が、異なる場所から同じ港に来るとし、その前のものが、もう一つに6時間先立ち、港の子午線に月が達してから3時ごろに起こるとする。もし月が、この子午線に達したとき赤道にあるなら、6時間ごとに等しい満ち潮が来て、それらは互いの干潮に落ち込んで、満ち潮と等しくなり、こうして、その日の間、水は静かに停滞する。もし月がそのとき赤道から傾いていたなら、潮は、すでに述べたように、大洋でかわるがわる大きく、小さくなり、そこから、この港には、かわるがわる、二つの大きな満ち潮と二つの小さな満ち潮が伝わる。そして二つの大きな満ち潮は、その中間で最も高い水を合成し、一つの大きいものと一つの小さいものは、その中間で水を平均の高さまで上げ、二つの小さな満ち潮の間では、水は最小の高さまで上る。こうして、24時間の間に、水は、ふつうのように2回ではなく、1回だけ最大の高さに達し、1回だけ最小の高さに達する。そして、月が場所の地平線の上の極のほうへ傾いているときは、最大の高さは、月が子午線に達してから6時間後か30時間後に起こり、月が赤緯を変えると、引き潮に変わる。

これらすべての例を、ハレーは、船乗りの観測から、北緯 20 度 50 分のトンキン王国のバトシャムの港で明らかにした。そこでは、月が赤道を通る日の次の日に水が停滞し、ついで月が北へ傾くと、他の港のように2回ではなく、毎日1回だけ満ち引きし始め、満ち潮は月の入りに、最大の引き潮は月の出に起こる。この潮は、月の赤緯とともに増し、7日目か8日目まで増し、ついでさらに7日間は、先に増したのと同じ度合いで減る。そして月が赤緯を変えるとやみ、まもなく引き潮に変わる。実際、そののち引き潮は月の入りに、満ち潮は月の出に起こって、月がふたたび赤緯を変えるまでつづく。この港とその近くの海峡への入り口は二つあり、一つは中国の海から大陸とルソン島の間を通り、もう一つはインド洋から大陸とボルネオ島の間を通る。インド洋から12時間をかけて、また中国の海から6時間をかけて、これらの海峡を通ってくる潮が、このようにして月の3時と9時に当たり、このような運動を合成するのか、あるいはこれらの海の条件が別なのかは、近くの海岸の観測で決めるべきこととして残しておく。

ここまで、私は月と海の運動の原因を述べた。運動の量については、いま少し付け加えるのがふさわしい。

\condpagelarge{命題 XXV. 問題 V.}

\textit{月の運動を乱す太陽の力を求めること。}

\pngright{422.png}{}{}{}

$Q$ を太陽、$S$ を地球、$P$ を月、$PADB$ を月の軌道とする。$QP$ 上に $QK$ を $QS$ に等しくとり、$QL$ 対 $QK$ が $QK$ 対 $QP$ の二乗の比になるようにし、$PS$ に平行に $LM$ を引く。そして、地球の太陽に向かう加速的な重力を、距離 $QS$ または $QK$ で表せば、$QL$ が月の太陽に向かう加速的な重力である。これは部分 $QM$, $LM$ から合成され、そのうち $LM$ と、$QM$ の部分 $SM$ が月の運動を乱す。これは第1編 命題LXVIとその系で述べたとおりである。地球と月が共通の重心のまわりを回る限り、地球のその中心のまわりの運動も、似た力で乱されるが、力と運動の和を月に帰して、力の和を、それらに類似する線 $SM$ と $ML$ で表してよい。力 $ML$ は（その平均の大きさで）、月が静止した地球のまわりを距離 $PS$ の軌道で回りうる重力に対して、（第1編 命題LXVIの系17により）、月の地球のまわりの周期と地球の太陽のまわりの周期との比の二乗、すなわち 27日7時間43分の 365日6時間9分に対する比の二乗、すなわち 1000 対 178725、あるいは 1 対 $178\frac{8}{11}$ である。月が静止した地球のまわりを、距離 $PS$、すなわち地球の半径の $60\frac{1}{2}$ 倍で、その軌道で回りうる力は、同じ時間に半径の60倍の距離で回りうる力に対して、$60\frac{1}{2}$ 対 60 であり、この力は私たちのところの重力に対して 1 対 $60 \times 60$ である。したがって、平均の力 $ML$ は、地球の表面での重力に対して、$1 \times 60\frac{1}{2}$ 対 $60 \times 60 \times 60 \times 178\frac{8}{11}$、すなわち 1 対 $638092.6$ である。そこから、線 $SM$ と $ML$ の比例によって、力 $SM$ も与えられる。そしてこれらが、月の運動を乱す太陽の力である。\QED

\condpagelarge{命題 XXVI. 問題 VI.}

\textit{月が地球に引いた半径で描く面積の増加を求めること。}

\pngright{424.png}{}{}{}

月が地球に引いた半径で描く面積は、月の運動が太陽の作用によって乱される限りを除いて、時間に比例すると、私たちは述べた。ここでは、その瞬間の増加（あるいは一時間あたりの増加）の不等を調べることにしよう。計算を容易にするために、月の軌道を円と仮定し、ここで問題にしている不等だけを除いて、すべての不等を無視する。また、太陽がきわめて遠いので、線 $QP$, $QS$ は互いに平行だと仮定する。こうすれば、力 $LM$ はつねにその平均の大きさ $SP$ に、力 $SM$ はその平均の大きさ $3PK$ に還元される。これらの力は、運動法則の系2により、力 $SL$ を合成する。この力は、半径 $SP$ に垂線 $LE$ を下ろせば、力 $SE$, $EL$ に分解される。そのうち $SE$ は、つねに半径 $SP$ に沿って働くから、その半径 $SP$ が描く面積 $QSP$ の描かれ方を、加速も減速もしない。一方 $EL$ は、垂線の方向に働くので、月を加速または減速するのと同じだけ、その面積の描かれ方を加速または減速する。月が矩 $C$ から合 $A$ へ進むあいだに、各瞬間に生じる月のこの加速は、加速する力 $EL$ に比例し、すなわち $\frac{3PK \times SK}{SP}$ に比例する。

時間を月の平均運動で表し、あるいは（ほとんど同じことになるが）角 $CSP$ で、あるいは弧 $CP$ でも表す。$CS$ に垂線 $CG$ を、$CS$ に等しくなるように立てる。四分円の弧 $AC$ を、無数の等しい部分 $Pp$ などに分け、それによって、やはり無数の等しい微小時間を表すことにする。$pk$ を $CS$ に垂線として引き、$SG$ を結んで、延長した $KP$, $kp$ と $F$, $f$ で交わらせると、$Kk$ 対 $PK$ は $Pp$ 対 $Sp$ に等しく、すなわち一定の比である。したがって $FK \times Kk$、すなわち面積 $FKkf$ は $\frac{3PK \times SK}{SP}$ に、すなわち $EL$ に比例する。合成すれば、全面積 $GCKF$ は、時間全体 $CP$ に月に加わったすべての力 $EL$ の和に比例し、したがってまた、この和によって生じた速さに、すなわち面積 $CSP$ の描かれ方の加速、あるいは瞬間の増加に比例する。

月が静止した地球のまわりを距離 $SP$ で、その周期 $CADBC$、$27$ 日 $7$ 時間 $43$ 分で回りうる力は、物体が時間 $CS$ のあいだ落下して長さ $\frac{1}{2}CS$ を描き、同時に月が軌道上を動く速さに等しい速さを得るようにする。これは第1編 命題IVの注から明らかである。ところで、$SP$ に下ろした垂線 $Kd$ は $EL$ の3分の1であり、八分点では $SP$、すなわち $ML$ の半分であるから、八分点では力 $EL$ が最大になって、力 $ML$ を $3$ 対 $2$ の比で超える。したがって、月がその周期で静止した地球のまわりを回りうる力に対しては、$100$ 対 $\frac{2}{3} \times 17872\frac{1}{2}$、すなわち $11915$ の比をなし、時間 $CS$ のあいだに、月の速さの $\frac{100}{11915}$ にあたる速さを生じるはずである。時間 $CPA$ のあいだには、$CA$ 対 $CS$、すなわち $SP$ の比だけ、より大きな速さを生じる。八分点での最大の力 $EL$ を、長方形 $\frac{1}{2}SP \times Pp$ に等しい面積 $FK \times Kk$ で表す。すると、最大の力が任意の時間 $CP$ に生じうる速さは、それより小さいすべての力 $EL$ が同じ時間に生じる速さに対して、長方形 $\frac{1}{2}SP \times CP$ 対面積 $KCGF$ の比をなす。全時間 $CPA$ では、生じる速さは互いに、長方形 $\frac{1}{2}SP \times CA$ と三角形 $SCG$ との比、すなわち四分円の弧 $CA$ 対半径 $SP$ の比をなす。したがって（ユークリッド原論 第5巻 命題9により）、後の速さは、全時間に生じて、月の速さの $\frac{100}{11915}$ になる。面積の平均の瞬間の増加に対応する、この月の速さに、もう一方の速さの半分を加え、また引くことにする。そして平均の瞬間の増加を数 $11915$ で表せば、和 $11915 + 50$、すなわち $11965$ は、朔望 $A$ における面積の最大の瞬間の増加を示し、差 $11915 - 50$、すなわち $11865$ は、矩におけるその最小の瞬間の増加を示す。ゆえに、朔望と矩で等しい時間に描かれる面積は、互いに $11965$ 対 $11865$ の比をなす。最小の増加 $11865$ に、ある増加を加える。この増加は、増加の差 $100$ に対して、台形 $FKCG$ 対三角形 $SCG$（これは、正弦 $PK$ の平方対半径 $SP$ の平方、すなわち $Pd$ 対 $SP$ に等しい）となるものとする。すると、この和が、月が中間の任意の位置 $P$ にあるときの面積の増加を示す。

これらのすべては、太陽と地球が静止していて、月が朔望周期ではなく $27$ 日 $7$ 時間 $43$ 分で回るという仮定の上でのことである。ところが、月の朔望周期は実際には $29$ 日 $12$ 時間 $44$ 分であるから、増加の増加分を、時間の比で増やさなければならない。こうすれば、平均の増加の $\frac{100}{11915}$ であった全増加は、その $\frac{100}{11023}$ になる。それゆえ、月の矩における面積の増加は、朔望におけるその増加に対して $11023 - 50$ 対 $11023 + 50$、すなわち $10973$ 対 $11073$ の比をなし、月が中間の別の位置 $P$ にあるときのその増加に対しては、$10973$ 対 $10973 + Pd$ の比をなす。ただし $SP$ は $100$ に等しいとする。

したがって、月が地球に引いた半径で、等しい微小時間ごとに描く面積は、数 $219.46$ と、月の最も近い矩からの距離の2倍の正矢（半径を1とする円における）との和に、ほぼ比例する。これは、八分点での変差が平均の大きさであるときのことである。もしそこでの変差がより大きいか小さいならば、その正矢を同じ比で増減しなければならない。

\condpagelarge{命題 XXVII. 問題 VII.}

\textit{月の一時間あたりの運動から、その地球からの距離を求めること。}

月が地球へ引いた半径で各瞬間に描く面積は、月の一時間あたりの運動と、月の地球からの距離の二乗との積に比例する。したがって、月の地球からの距離は、面積の平方根の比と、一時間あたりの運動の平方根の逆の比とから合成された比をなす。\QED

\textit{系1.}\wsp{}ここから月の見かけの直径が与えられる。それは地球からの距離に反比例するからである。天文学者は、この規則が現象とどれほどよく合うかを試みよ。

\textit{系2.}\wsp{}ここからまた、月の軌道を、これまでより正確に、現象から定めることができる。

\condpagelarge{命題 XXVIII. 問題 VIII.}

\textit{月が離心率なしに動くはずの軌道の直径を求めること。}

\pngright{427.png}{}{}{}

動く物体が、軌跡に垂直に引かれるときに描く軌跡の曲率は、引力に正比例し、速さの二乗に反比例する。私は、線の曲率は、等しい半径に対する接角の正弦または正接の、半径を限りなく小さくしたときの最後の比に、互いに比例するとする。ところで、朔望における月の地球への引力は、月の地球への重力が、太陽の力 $2PK$（命題XXVの図を見よ）を超える分であり、この力によって月の太陽に向かう加速的な重力が、地球の太陽に向かう加速的な重力を超える。矩では、その引力は、月の地球への重力と、月を地球に引く太陽の力 $KS$ との和である。そしてこれらの引力は、$\frac{AS + CS}{2}$ を $N$ と呼べば、ほぼ
\[ \frac{178725}{AS^2} - \frac{2000}{CS \times N} \quad\text{と}\quad \frac{178725}{CS^2} + \frac{1000}{AS \times N} \]
に比例する。すなわち $178725N \times CS^2 - 2000AS^2 \times CS$ と $178725N \times AS^2 + 1000CS^2 \times AS$ に比例する。実際、地球の太陽に向かう加速的な重力を数 178725 で表せば、矩では $PS$ または $SK$ に等しく月を地球に引く平均の力 $ML$ は 1000、朔望での平均の力 $SM$ は 3000 であり、これから平均の力 $ML$ を引けば、月が朔望で地球から引き離される力 2000 が残る。これを、私は前に $2PK$ と呼んだ。ところで、朔望 $A$, $B$ での月の速さは、矩 $C$, $D$ での月の速さに対して、$CS$ 対 $AS$ と、朔望での月が地球へ引いた半径で描く面積の瞬間的な増加の、矩での同じ面積の増加に対する比とから、合成された比をなす。すなわち $11073CS$ 対 $10973AS$ である。この比を二重に逆にとり、前の比を一度正にとれば、朔望における月の軌道の曲率の、矩におけるその曲率に対する比は、$120407 \times 178725AS^2 \times CS^2 \times N - 120407 \times 2000AS^4 \times CS$ 対 $122611 \times 178725AS^2 \times CS^2 \times N + 122611 \times 1000CS^4 \times AS$、すなわち $2151969AS \times CS \times N - 24081AS^3$ 対 $2191371AS \times CS \times N + 12261CS^3$ になる。

月の軌道の形が知られていないから、その代わりに、楕円 $DBCA$ をとる。その中心 $S$ に地球を置き、長軸 $DC$ は矩の間に、短軸 $AB$ は朔望の間にある。ところで、この楕円の面は、地球のまわりを角運動で回転し、曲率を考える軌跡は、すべての角運動をまったく欠く面の中で描かれなければならないから、月がその楕円の中を回りながらこの面の中で描く図形、すなわち図形 $Cpa$ を考えなければならない。その個々の点 $p$ は、楕円上に月の位置を表す任意の点 $P$ をとり、$Sp$ を $SP$ に等しく引くことで求められる。ただしそれは、角 $PSp$ が、矩 $C$ の時刻から経過した太陽の見かけの運動に等しくなるように、あるいは（ほぼ同じことだが）角 $CSp$ の角 $CSP$ に対する比が、月の朔望周期の恒星周期に対する比、すなわち 29日12時間44分の 27日7時間43分に対する比になるように行う。そこで、角 $CSa$ を、直角 $CSA$ に対して同じ比にとり、長さ $Sa$ を長さ $SA$ に等しくすれば、$a$ はこの軌道 $Cpa$ の近いアプシス、$C$ は遠いアプシスになる。そして、計算によって私は次のことを見いだした。頂点 $a$ での軌道 $Cpa$ の曲率と、中心 $S$、半径 $SA$ で描いた円の曲率との差は、頂点 $A$ での楕円の曲率とその同じ円の曲率との差に対して、角 $CSP$ の角 $CSp$ に対する比の二乗である。また $A$ での楕円の曲率は、その円の曲率に対して、$SA$ 対 $SC$ の比の二乗である。その円の曲率は、中心 $S$、半径 $SC$ で描いた円の曲率に対して $SC$ 対 $SA$ であり、この円の曲率は、$C$ での楕円の曲率に対して、$SA$ 対 $SC$ の比の二乗である。そして、頂点 $C$ での楕円の曲率と最後の円の曲率との差は、頂点 $C$ での図形 $Spa$ の曲率とその同じ円の曲率との差に対して、角 $CSP$ の角 $CSp$ に対する比の二乗である。これらの比は、接角の正弦と角の差から容易に求められる。これらの比を互いに突き合わせれば、図形 $Cpa$ の $a$ での曲率の、$C$ での曲率に対する比は、$AS^3 + \frac{16824}{100000}CS^2 \times AS$ 対 $CS^3 + \frac{16824}{100000}AS^2 \times CS$ となる。ここで数 $\frac{16824}{100000}$ は、角 $CSP$ と $CSp$ の平方の差を、小さいほうの角 $CSP$ の平方で割ったもの、すなわち（同じことだが）時間 27日7時間43分と 29日12時間44分の平方の差を、時間 27日7時間43分の平方で割ったものを表す。したがって、$a$ が月の朔望を、$C$ が月の矩を表すから、いま求めた比は、先に求めた、月の軌道の朔望での曲率の矩でのその曲率に対する比と、等しくなければならない。そこで、$CS$ 対 $AS$ の比を求めるために、外項と中項を互いに掛ける。出てきた項を $AS \times CS$ で割ると、
\[ 2062.79\,CS^4 - 2151969N \times CS^3 + 368682N \times AS \times CS^2 + 36342AS^2 \times CS^2 - 362046N \times AS^2 \times CS + 2191371N \times AS^3 + 4051.4\,AS^4 = 0 \]
となる。ここで、項 $AS$ と $CS$ の半和 $N$ のかわりに 1 を書き、それらの半差を $x$ とおくと、$CS = 1 + x$, $AS = 1 - x$ となる。これらを方程式に代入して、出てきた方程式を解くと、$x$ は $0.0072036$ に等しく、そこから半径 $CS$ は $1.0072$、半径 $AS$ は $0.9928$ となり、これらの数はほぼ $69\frac{11}{12}$ と $68\frac{11}{12}$ の比をなす。したがって、月の地球からの距離は、朔望では矩におけるその距離に対して（離心率の考察を除いて）$68\frac{11}{12}$ 対 $69\frac{11}{12}$、あるいは丸めた数で 69 対 70 である。
\chapter{第3編 世界の体系（月の摂動・潮汐論・分点歳差・彗星の幾何学軌道論）}



\section{第3編 月の摂動論、分点歳差、および彗星の幾何学軌道論（命題XXV〜XLII）}


% Section: 第3編 月の摂動論、分点歳差、および彗星の幾何学軌道論（命題XXV〜XLII）



月がその楕円軌道  $Cpa$  上で回転するとき、この軌道を描く図形は、この平面において  $Cpa$  のような楕円であり、各点  $p$  を特定するために、月の位置を表す任意の点  $P$  を楕円上に選んだ後、線分  $Sp$  を線分  $SP$  と等しく引く。ここで、角度  $SSp$  は太陽からの相対的な運動の角度であり、これは四分点から完成した時間  $C$  のときの角度と等しいものとする。または（ほぼ同じ結果を得るためには）、角度  $CSs$  は角度  $CSP$  に比例し、同様に太陽の同期回転周期の時間比  $29 \textbf{d.} 12 \textbf{h.} 44$\textbf{分}$\ )$と新月から$27$日$7$時間$43$分の周期回転期間の時間比  $27 \textbf{d.} 7 \textbf{h.} 43$\textbf{分}に比例するものとする。したがって、

\[ C \]

\begin{align*}
\textbf{角}CSaは直角CSAに同一の比で、かつ長さSaは長さSAに等しいとする；すると点aは軌道Cpaの極小焦点であり、点Cは同軌道の極大焦点となる。また、導入して求めると、軌道Cpaの頂点aにおける曲率と中心Sを通り半径SAで描かれた円の曲率との差は、角CSPに対する角CSpの二倍比に等しい差である；かつ、楕円Aにおける曲率は同じ円の曲率に対しても同様の二倍比である。
\end{align*}


\par\begin{center}
  \centering
  \includegraphics[width=0.8\textwidth]{429.png}

\end{center}\par

\textbf{SA}を\textbf{SC}に、かつその円の曲率を中心\textbf{S}と半径\textbf{SC}で描く円の曲率との比が\textbf{SC}に対して\textbf{SA}として与えられる。しかしながら、この曲率は楕円において点\textbf{C}における曲率に対して、\textbf{SA}と\textbf{SC}との比の二倍乗比に等しい。また、楕円において点\textbf{C}における曲率と最後の円の曲率との差は、図形\textbf{Spa}において点\textbf{C}における曲率と同円の曲率との差に対して、角\textbf{CSP}と角\textbf{CSp}の比の二倍乗比に等しい。これらの比率は接線の正弦と接線の差の正弦から容易に導かれる。これらを比較した結果、点\textbf{C}における曲率が明らかとなる。



図形 \( Cpa \) 在 \( a \) に対する其在 \( C \) 処的曲率，如 \( AS^3 + \frac{16824}{100000}CS^2 \times AS \) 対比 \( CS^3 + \frac{16824}{100000}AS^2 \times CS \)。其中数 \( \frac{16824}{100000} \) 表示角 \( CSP \) 与 \( CSp \) 的平方差，応用于より小さな角 \( CSP \) 的平方，即（等价地）時間 27 日 7 小時 43 分鐘与 29 日 12 小時 44 分鐘之間的時間差的平方差。

したがって，当 \( a \) 表示月之冲点，而 \( C \) 表示其四分之一円周位置，



\textbf{既に得られた比例は、月の軌道曲率在合点における比例と等しくなければならない。}

即ち、月の軌道曲率在合点における比例と、前述で求めた四分点における比例が一致しなければならない。従って、$CS$ に対する $AS$ の比例を求めるために、端点と中間比を互いに計算する。そして、$AS \times CS$ への延長線は以下の式を得る：

\[
2062.79 \, CS^2 - 2151969N \, CS^3 + 368682N \, AS \, CS^2 + 36342 \, AS^2 \, CS^2 - 362046N \, AS^2 \, CS + 2191371N \, AS^3 + 4051.4 \, AS^2 = 0
\]

ここで、$AS$ と $CS$ の半和を $N$ として書き換える。そして、

\[
N = \frac{AS + CS}{2}
\]

月の位置が対合点（Syzygiis）と四分点（Quadraturis）における太陽からの距離との比を求めるため、まず \( x \) を仮定します。すると、\( CS = 1 + x \) で、\( AS = 1 - x \) となります。これらの値を方程式に代入し解くと、得られる \( x \) の値は \( 0.\textbf{0072036} \) となり、従って半径 \( CS \) は \( 1.\textbf{0072} \) で、半径 \( AS \) は \( 0.\textbf{9928} \) となります。これらの数値は約 \( 69\frac{11}{12} \) と \( 68\frac{11}{12} \) の比に等しいです。したがって、対合点における月の太陽からの距離は四分点における距離に対して \( 68\frac{11}{12} \) と \( 69\frac{11}{12} \) の比、または丸めた数値では 69 の比となります。

\condpagelarge{命題 XXIX. 問題 IX.}

\textit{月の出差（変差）を求めること。}

\pngright{432.png}{}{}{}

この不等は、一部は月の軌道の楕円の形から、一部は、月が地球に引いた半径で描く面積の瞬間的な増加の不等から生じる。もし月 $P$ が、地球を中心として静止させた楕円 $DBCA$ の上を動き、地球に引いた半径 $SP$ で、時間に比例する面積 $CSP$ を描くとしよう。そして楕円の最大半径 $CS$ の最小半径 $SA$ に対する比が $69\frac{11}{12}$ 対 $68\frac{11}{12}$ であるとすれば、角 $CSP$ の正接は、矩 $C$ から数えた平均運動の角の正接に対して、楕円の半径 $SA$ 対その半径 $SC$、すなわち $68\frac{11}{12}$ 対 $69\frac{11}{12}$ の比をなす。ところが、面積 $CSP$ の描かれ方は、月が矩から朔望へ進むにつれて、朔望での面積の瞬間的な増加が矩でのそれに対して $11073$ 対 $10973$ になり、かつ中間の任意の位置 $P$ での瞬間的な増加が矩でのそれを超える分は、角 $CSP$ の正弦の平方に比例するように、加速されなければならない。これは、角 $CSP$ の正接を、数 $10973$ 対数 $11073$ の比の平方根の比で小さくすれば、十分に正確にそうなる。すなわち、数 $68\frac{5958}{10000}$ 対数 $68\frac{11}{12}$ の比で小さくするのである。こうすれば、角 $CSP$ の正接は、平均運動の正接に対して $68\frac{5958}{10000}$ 対 $69\frac{11}{12}$ となり、平均運動が45度である八分点では、角 $CSP$ は $44$ 度 $27\minute 29\second$ と求められる。これを平均運動の角45度から引くと、変差は $32\minute 31\second$ になる。

これは、月が矩から朔望へ進むあいだに、わずか90度の角 $CSA$ を描くのなら、そうなるはずの値である。しかし、地球の運動のために太陽が見かけの運動で前の方へ運ばれるので、月は太陽に追いつく前に、直角より大きい角 $CSa$ を描く。それは月の朔望周期の恒星周期に対する比、すなわち $29$ 日 $12$ 時 $44$ 分対 $27$ 日 $7$ 時 $43$ 分の比だけ大きい。こうして中心 $S$ のまわりの角はすべてこの比で広がり、そうでなければ $32\minute 31\second$ であるはずの変差は、同じ比で増して $35\minute 9\second$ になる。これを天文学者は $40$ 分と定め、より新しい観測からは $38$ 分である。ハレーは最近、これが、太陽の衝に向かう八分点では $38$ 分、太陽に向かう八分点では $32$ 分であることを見いだした。したがって、その平均の大きさは $35$ 分となり、私たちが求めた大きさ $35\minute 9\second$ とよく一致する。というのは、私たちは平均の大きさを、大きな軌道の曲がりや、三日月・新月に対する太陽の作用が、凸月・満月に対するそれより大きいことから生じうる差を、無視して計算したからである。

\condpagelarge{命題 XXX. 問題 X.}

\textit{円軌道上の月の、交点の一時間あたりの運動を求めること。}

$S$ を太陽、$T$ を地球、$P$ を月、$NPn$ を月の軌道、$Npn$ をその軌道の黄道面への射影、$N$, $n$ を交点、$nTNm$ を限りなく延長した交点線とする。$PI$, $PK$ を線 $ST$, $Qq$ に下ろした垂線、$Pp$ を黄道面に下ろした垂線、$Q$, $q$ を黄道面での月の矩、$pK$ を矩を結ぶ線 $Qq$ に下ろした垂線とする。すると、月の運動を乱す太陽の力は（命題XXVにより）二つあり、一つは線 $2IT$ または $2Kp$ に比例し、もう一つは線 $PI$ に比例する。月は前の力で太陽の方へ、後の力で線 $ST$ の方へ引かれる。

\pngright{432.png}{}{}{}

後の力 $PI$ は、力 $IT$ と $PT$ から合成され、そのうち $PT$ は月の軌道面に沿って働くから、面の位置をまったく変えない。したがってこれは無視してよい。一方、力 $IT$ は力 $2IT$ と合成されて全体の力 $3IT$ となり、これが月の軌道面を乱す。そしてこの力は（命題XXVにより）、月が静止した地球のまわりを円に沿って自分の周期で回りうる力に対して、$3IT$ 対、円の半径に数 $178.725$ を掛けたもの、すなわち $IT$ 対、半径に $59.575$ を掛けたもの、という比をなす。なお、この計算とこのあとのすべての計算で、私は月から太陽へ引いた線をすべて、地球から太陽へ引いた線に平行とみなす。それは、傾きが効果を減らす分は、別の場合にはほぼ同じだけ増やすからであり、また私たちが求めているのは交点の平均運動であって、計算をあまりにも複雑にするこのような細かい点は無視するからである。

さて、$PM$ を月が与えられた極めて短い時間に描く弧とし、$ML$ を、前述の力 $3IT$ に押されて、月が同じ時間に描きうる微小な線とする。$PL$, $MP$ を結び、これらを延長して $l$, $m$ で黄道面と交わらせ、$Tm$ に垂線 $PH$ を下ろす。$ML$ は $ST$ に平行であるから、$ml$ が $ML$ に平行なら $ml$ は黄道面内にあり、逆も成り立つ。ゆえに $ml$ は黄道面内にあるから $ML$ に平行であり、三角形 $LMP$, $Lmp$ は相似である。さて $MPm$ は、月が位置 $P$ で動いていた軌道面内にあるから、点 $m$ は、その軌道の交点 $N$, $n$ を通る線 $Nn$ の上にある。そして、微小な線 $LM$ を生じる力が、全部いちどきに位置 $P$ で加えられたとすれば、月は弦が $LP$ である弧に沿って動かされ、月を面 $MPmT$ から面 $LPlT$ へ移すことになる。したがって、この力から生じる交点の運動は角 $mTl$ に等しい。ところが $ml$ 対 $mP$ は $ML$ 対 $MP$ に等しく、$MP$ は時間が与えられているから与えられている。ゆえに $ml$ は長方形 $ML \times mP$、すなわち長方形 $IT \times mP$ に比例する。そして角 $mTl$ は、角 $Tml$ が直角であれば $\frac{ml}{Tm}$ に比例し、したがって $\frac{IT \times Pm}{Tm}$ に比例し、すなわち（$Tm$ と $mP$、$TP$ と $PH$ が比例することから）$\frac{IT \times PH}{TP}$ に比例し、$TP$ が与えられているので $IT \times PH$ に比例する。もし角 $Tml$、すなわち $STN$ が斜めなら、角 $mTl$ は、角 $STN$ の正弦対半径の比で、さらに小さくなる。ゆえに交点の速さは、$IT \times PH$ と角 $STN$ の正弦との積に比例する。すなわち、三つの角 $TPI$, $PTN$, $STN$ の正弦の積に比例する。

交点が矩にあり、月が朔望にあるとき、これらの角がすべて直角ならば、微小な線 $ml$ は無限に大きくなり、角 $mTl$ は角 $mPl$ に等しくなる。この場合、角 $mPl$ は、月が同じ時間に見かけの運動で地球のまわりに描く角 $PTM$ に対して、$1$ 対 $59.575$ の比をなす。というのは、角 $mPl$ は角 $LPM$、すなわち前述の太陽の力 $3IT$ が、その与えられた時間に生じさせうる、月の直線の道からのそれの角に等しく、角 $PTM$ は、月を軌道上にとどめている力が、同じ時間に生じさせる、月の直線の道からのそれの角に等しいからである。そしてこれらの力は、上で述べたように、互いに $1$ 対 $59.575$ の比をなす。ゆえに、月の（恒星に対する）一時間あたりの平均運動は $32\minute 56\second 27\third 12\fourth\frac{1}{2}$ であるから、この場合の交点の一時間あたりの運動は $33\second 10\third 33\fourth 12\fifth$ になる。ほかの場合には、この一時間あたりの運動は、$33\second 10\third 33\fourth 12\fifth$ に対して、三つの角 $TPI$, $PTN$, $STN$（すなわち月の矩からの距離、月の交点からの距離、交点の太陽からの距離）の正弦の積対半径の3乗、の比をなす。そして、どれかの角の正弦の符号が正から負へ、あるいは負から正へ変わるたびに、後退運動は前進運動に、前進運動は後退運動に変わらなければならない。したがって、月がどちらかの矩とそれに最も近い交点との間にあるたびに、交点は前進する。ほかの場合には後退し、後退が前進を超える分によって、毎月、先行する方向へ運ばれる。

\pngright{434.png}{}{}{}

\textit{系1.}\wsp{}ここから、与えられた極めて短い弧 $PM$ の端点 $P$, $M$ から、矩を結ぶ線 $Qq$ に垂線 $PK$, $Mk$ を下ろし、それらを延長して交点線 $Nn$ と $D$, $d$ で交わらせると、交点の一時間あたりの運動は、面積 $MPDd$ と線 $AZ$ の平方との積に比例する。実際、$PK$, $PH$, $AZ$ を前述の三つの正弦とする。すなわち $PK$ は月の矩からの距離の正弦、$PH$ は月の交点からの距離の正弦、$AZ$ は交点の太陽からの距離の正弦である。すると、交点の速さは積 $PK \times PH \times AZ$ に比例する。ところが $PT$ 対 $PK$ は $PM$ 対 $Kk$ に等しく、$PT$ と $PM$ は与えられているから、$Kk$ は $PK$ に比例する。また $AT$ 対 $PD$ は $AZ$ 対 $PH$ に等しく、したがって $PH$ は長方形 $PD \times AZ$ に比例する。これらの比を合わせれば、$PK \times PH$ は積 $Kk \times PD \times AZ$ に比例し、$PK \times PH \times AZ$ は $Kk \times PD \times AZ^2$ に比例する。すなわち、面積 $PDdM$ と $AZ^2$ との積に比例する。\QED

\textit{系2.}\wsp{}交点の位置が任意に与えられたとき、一時間あたりの平均運動は、月が朔望にあるときの一時間あたりの運動の半分であり、したがって $16\second 35\third 16\fourth 36\fifth$ に対して、交点の朔望からの距離の正弦の平方対半径の平方、すなわち $AZ^2$ 対 $AT^2$ の比をなす。なぜなら、月が一様な運動で半円 $QAq$ を通るとき、月が $Q$ から $M$ へ進む時間にすべての面積 $PDdM$ を足した和は、円の接線 $QE$ で限られる面積 $QMdE$ になり、月が点 $n$ に達するときには、この和は線 $PD$ が描く面積全体 $EQAn$ になる。ついで月が $n$ から $q$ へ進むと、線 $PD$ は円の外に出て、円の接線 $qe$ で限られる面積 $nqe$ を描く。交点は前には後退していたのに、今度は前進するから、これは先の面積から引かなければならず、面積 $QEN$ に等しいので、半円 $NQAn$ が残る。したがって、月が半円を描く時間にすべての面積 $PDdM$ を足した和は半円の面積であり、月が円を描く時間にすべての面積を足した和は円全体の面積である。ところで、月が朔望にあるとき、面積 $PDdM$ は、弧 $PM$ と半径 $MT$ とがつくる長方形であり、月が円を描く時間にこれに等しい面積をすべて足した和は、円周全体と円の半径とがつくる長方形である。この長方形は二つの円の面積に等しいから、前の長方形の2倍である。それゆえ交点は、月の朔望で持つ速さを一様に続けるとすれば、実際に描く空間の2倍の空間を描くことになる。そこで、一様に続けたとき、不均等な運動で実際に進んだ空間を描きうる平均の運動は、朔望での運動の半分である。したがって、交点が矩にあるとき、最大の一時間あたりの運動が $33\second 10\third 33\fourth 12\fifth$ であるから、この場合の平均の一時間あたりの運動は $16\second 35\third 16\fourth 36\fifth$ である。そして交点の一時間あたりの運動はつねに $AZ^2$ と面積 $PDdM$ との積に比例し、したがって月が朔望にあるときの交点の運動は、$AZ^2$ と面積 $PDdM$ との積に、すなわち（朔望で描かれる面積 $PDdM$ が一定であるから）$AZ^2$ に比例するから、平均の運動もまた $AZ^2$ に比例する。ゆえにこの運動は、交点が矩以外にあるとき、$16\second 35\third 16\fourth 36\fifth$ に対して $AZ^2$ 対 $AT^2$ の比をなす。\QED

\condpagelarge{命題 XXXI. 問題 XI.}

\textit{楕円軌道上の月の、交点の一時間あたりの運動を求めること。}

\pngright{436.png}{}{}{}

$Qpmaq$ を、長軸 $Qq$、短軸 $ab$ で描いた楕円とし、$QAq$ をそれに外接する円、$T$ を両者に共通な中心にある地球、$S$ を太陽、$p$ を楕円上を動く月、$pm$ を与えられたきわめて短い時間に月が描く弧、$N$ と $n$ を線 $Nn$ で結んだ交点とする。$pK$ と $mk$ を軸 $Qq$ に下ろした垂線とし、両側に延長して、円と $P$, $M$ で、交点線と $D$, $d$ で交わらせる。もし月が地球に引いた半径で時間に比例する面積を描くならば、楕円における交点の運動は面積 $pDdm$ に比例する。

さて、$PF$ が円に $P$ で接し、延長して $TN$ と $F$ で交わるとし、$pf$ が楕円に $p$ で接し、延長して同じ $TN$ と $f$ で交わるとし、これらの接線は軸 $TQ$ 上の $Y$ で出会うとする。また $ML$ を、円を回る月が弧 $PM$ を描くあいだに、前述の力 $3IT$ に押されて、横の運動によって描きうる距離とし、$ml$ を、楕円を回る月が同じ時間に、同じく力 $3IT$ に押されて描きうる距離とする。$LP$ と $lp$ を延長して黄道面と $G$, $g$ で交わらせ、$FG$ と $fg$ を結ぶ。延長した $FG$ が $pf$, $pg$, $TQ$ とそれぞれ $c$, $e$, $R$ で交わり、延長した $fg$ が $TQ$ と $r$ で交わるとする。円における力 $3IT$、すなわち $3PK$ は、楕円における力 $3IT$、すなわち $3pK$ に対して、$PK$ 対 $pK$、すなわち $AT$ 対 $aT$ の比をなすから、前の力で生じる距離 $ML$ は、後の力で生じる距離 $ml$ に対して、$PK$ 対 $pK$、すなわち（相似な図形 $PYKp$ と $FYRc$ により）$FR$ 対 $cR$ の比をなす。ところが $ML$ 対 $FG$ は、（相似な三角形 $PLM$, $PGF$ により）$PL$ 対 $PG$ に等しく、すなわち（平行線 $Lk$, $PK$, $GR$ により）$pl$ 対 $pe$ に等しく、すなわち（相似な三角形 $plm$, $cpe$ により）$lm$ 対 $ce$ に等しい。逆にいえば、$LM$ 対 $lm$、すなわち $FR$ 対 $cR$ は、$FG$ 対 $ce$ に等しい。したがって、もし $fg$ 対 $ce$ が $fY$ 対 $cY$、すなわち $fr$ 対 $cR$ に等しいなら（これは、$fr$ 対 $FR$ と $FR$ 対 $cR$ の積、すなわち $fT$ 対 $FT$ と $FG$ 対 $ce$ の積に等しい）、両辺から比 $FG$ 対 $ce$ を除くと、比 $fg$ 対 $FG$ と $fT$ 対 $FT$ が残るから、$fg$ 対 $FG$ は $fT$ 対 $FT$ に等しくなるはずである。なぜなら、$FG$ と $fg$ が地球 $T$ に張る角は、互いに等しくなるからである。ところが、それらの角は（前の命題で述べたことにより）、月が円では弧 $PM$ を、楕円では弧 $pm$ をたどる時間における交点の運動である。したがって、円と楕円での交点の運動は互いに等しくなる。これは、$fg$ 対 $ce$ が $fY$ 対 $cY$ に等しい場合に、すなわち $fg$ が $\frac{ce \times fY}{cY}$ に等しい場合に限って成り立つ。しかし、相似な三角形 $fgp$, $cep$ により、$fg$ 対 $ce$ は $fp$ 対 $cp$ に等しい。ゆえに $fg$ は $\frac{ce \times fp}{cp}$ に等しい。したがって、$fg$ が実際に張る角の、前の角（$FG$ が張る角）に対する比、すなわち楕円での交点の運動の円での交点の運動に対する比は、この $fg$、すなわち $\frac{ce \times fp}{cp}$ の、前の $fg$、すなわち $\frac{ce \times fY}{cY}$ に対する比、すなわち $fp \times cY$ 対 $cp \times fY$ に等しく、$fp$ 対 $fY$ と $cY$ 対 $cp$ の積に等しい。これは、$pb$ を $TN$ に平行に引いて $FP$ と $b$ で交わらせるとき、$Fb$ 対 $FY$ と $FY$ 対 $FP$ の積、すなわち $Fb$ 対 $FP$、すなわち $Dp$ 対 $DP$ に等しく、したがって面積 $Dpmd$ 対面積 $DPMd$ に等しい。そして後の面積は円における交点の運動に比例するから、前の面積は楕円における交点の運動に比例する。\QED

\textit{系.}\wsp{}したがって、交点の位置が与えられているとき、月が矩から任意の位置 $m$ に進む時間に、すべての面積 $pDdm$ を足した和は、楕円の接線 $QE$ で限られる面積 $mpQEd$ であり、一回転全体でこれらの面積を足した和は楕円全体の面積であるから、楕円における交点の平均運動は、円における交点の平均運動に対して、楕円対円の比、すなわち $Ta$ 対 $TA$、すなわち $68\frac{11}{12}$ 対 $69\frac{11}{12}$ の比をなす。ゆえに、円における交点の一時間あたりの平均運動は $16\second 35\third 16\fourth 36\fifth$ に対して $AZ^2$ 対 $AT^2$ の比をなすから、角 $16\second 21\third 2\fourth 36\fifth$ を、角 $16\second 35\third 16\fourth 36\fifth$ に対して $68\frac{11}{12}$ 対 $69\frac{11}{12}$ の比になるようにとれば、楕円における交点の一時間あたりの平均運動は、$16\second 21\third 2\fourth 36\fifth$ に対して $AZ^2$ 対 $AT^2$ の比、すなわち交点の太陽からの距離の正弦の平方対半径の平方の比をなす。

ところで、月は地球に引いた半径で、朔望では矩におけるより速く面積を描き、そのため朔望では時間が縮み、矩では延びる。そして時間とともに、交点の運動も増減する。月の矩における面積の瞬間的な増加は、朔望におけるそれに対して $10973$ 対 $11073$ であった。したがって、八分点における平均の増加は、朔望での超過と矩での不足に対して、これらの数の和の半分 $11023$ 対、その差の半分 $50$ の比をなす。ゆえに、軌道の等しい各部分に対する月の時間は、その速さに反比例するから、八分点での平均の時間は、この原因から生じる矩での時間の超過と朔望での時間の不足に対して、ほぼ $11023$ 対 $50$ の比になる。さて矩から朔望へ進むとき、各位置での面積の瞬間的な増加が矩での最小の増加を超える分は、月の矩からの距離の正弦の平方に、ほぼ比例することがわかる。したがって、任意の位置での増加と、八分点での平均の増加との差は、月の矩からの距離の正弦の平方と、45度の正弦の平方、すなわち半径の平方の半分との差に比例する。そして、八分点と矩の間の各位置での時間の増加、八分点と朔望の間での時間の減少も、同じ比である。ところで、月が軌道の等しい各部分を通る時間における交点の運動は、時間の二乗の比で加速または減速される。というのは、月が $PM$ を通るあいだのこの運動は（他が同じならば）$ML$ に比例し、$ML$ は時間の二乗に比例するからである。ゆえに、朔望での交点の運動は、月が軌道の与えられた部分を通る時間のうちに行われるものとして、数 $11073$ 対数 $11023$ の二乗の比で減る。この減少は、残りの運動に対して $100$ 対 $10973$、運動全体に対してはほぼ $100$ 対 $11073$ である。そして、八分点と朔望の間の位置での減少、および八分点と矩の間の位置での増加は、この減少に対して、ほぼ、それらの位置での全運動の朔望での全運動に対する比と、月の矩からの距離の正弦の平方と半径の平方の半分との差の、半径の平方の半分に対する比との積をなす。ゆえに、もし交点が矩にあるとして、八分点から両側に等しい距離にある二つの位置、および朔望と矩から同じ間隔だけ離れた別の二つの位置をとり、朔望と八分点の間の二つの位置での運動の減少から、八分点と矩の間の残りの二つの位置での運動の増加を引くならば、残りの減少は朔望での減少に等しくなる。これは、計算してみれば容易にわかる。したがって、交点の平均運動から引くべき平均の減少は、朔望での減少の4分の1である。朔望での交点の一時間あたりの全運動（月が地球に引いた半径で時間に比例する面積を描くと仮定した場合）は $32\second 42\third 5\fourth 12\fifth$ であった。そして月がより速くなって同じ空間を描く時間における交点の運動の減少は、この運動に対して $100$ 対 $11073$ であると述べたから、その減少は $17\third 43\fourth 10\fifth$ であり、その4分の1の $4\third 25\fourth 48\fifth$ を、上で求めた一時間あたりの平均運動 $16\second 21\third 2\fourth 36\fifth$ から引くと、補正した一時間あたりの平均運動 $16\second 16\third 36\fourth 48\fifth$ が残る。

もし交点が矩以外にあって、朔望から両側に等しい距離にある二つの位置に注目するならば、月がこれらの位置にあるときの交点の運動の和は、月が同じ位置にあって交点が矩にあるときの運動の和に対して、$AZ^2$ 対 $AT^2$ の比をなす。そして、いま述べた原因から生じる運動の減少は、運動そのものに比例するから、残りの運動も互いに $AZ^2$ 対 $AT^2$ の比をなし、平均運動も残りの運動に比例する。ゆえに、交点の位置が任意に与えられたとき、補正した一時間あたりの平均運動は、$16\second 16\third 36\fourth 48\fifth$ に対して $AZ^2$ 対 $AT^2$ の比、すなわち交点の朔望からの距離の正弦の平方対半径の平方の比をなす。
\condpagelarge{命題 XXXII. 問題 XII.}

\textit{月の交点の平均運動を求めること。}

年間の平均運動とは、一年間の一時間あたりの平均運動をすべて足し合わせたものである。交点が $N$ にあり、一時間ごとに、その時間が終わるたびに元の位置へ引きもどされて、自分自身の運動にかかわらず、恒星に対してつねに一定の位置を保つとしよう。その間に太陽 $S$ は、地球の運動によって交点から離れていき、見かけの年周の道を一様に完了するとする。そこで $Aa$ を、太陽へつねに引いた直線 $TS$ と円 $NAn$ との交点が、与えられた極めて短い時間に描く、与えられた微小な弧とする。すると、一時間あたりの平均運動は（すでに示したことにより）$AZ^2$ に比例し、すなわち（$AZ$, $ZY$ が比例するから）$AZ$ と $ZY$ の長方形に、すなわち面積 $AZYa$ に比例する。そして、最初からのすべての一時間あたりの平均運動の和は、すべての面積 $aYZA$ の和、すなわち面積 $NAZ$ に比例する。ところで最大の $AZYa$ は、弧 $Aa$ と円の半径との長方形に等しい。したがって、円全体にわたるすべての長方形の和は、同じ数の最大のものの和に対して、円全体の面積対、円周全体と半径との長方形、すなわち $1$ 対 $2$ の比をなす。最大の長方形に対応する一時間あたりの運動は $16\second 16\third 36\fourth 48\fifth$ であった。この運動は、恒星年の全体、365日6時間9分で、$39$ 度 $38\minute 5\second 39\third$ になる。したがって、その半分 $19$ 度 $49\minute 2\second 49\third\frac{1}{2}$ が、円全体に対応する交点の平均運動である。そして、太陽が $N$ から $A$ へ進む時間における交点の運動は、$19$ 度 $49\minute 2\second 49\third\frac{1}{2}$ に対して、面積 $NAZ$ 対円全体の比をなす。

\pngright{441.png}{}{}{}

これは、交点が一時間ごとに元の位置へ引きもどされ、太陽が一年を終えたときに、最初に離れたのと同じ交点にもどる、という仮定の上での話である。しかし、交点の運動のために、太陽はより早く交点にもどる。そこで、時間の短縮を計算しなければならない。太陽は一年で $360$ 度を進み、交点は最大の運動で同じ時間に $39$ 度 $38\minute 5\second 39\third$、すなわち $39.6349$ 度を進むはずである。そして、任意の位置 $N$ における交点の平均運動は、矩におけるその平均運動に対して $AZ^2$ 対 $AT^2$ の比をなすから、$N$ での太陽の運動の交点の運動に対する比は、$360\,AT^2$ 対 $39.6349\,AZ^2$、すなわち $9.0829032\,AT^2$ 対 $AZ^2$ になる。そこで、円全体の円周 $NAn$ を等しい微小部分 $Aa$ に分けるとき、円が静止している場合に太陽が微小部分 $Aa$ を通る時間は、円が交点とともに中心 $T$ のまわりを回っている場合に太陽が同じ微小部分を通る時間に対して、$9.0829032\,AT^2$ と $9.0829032\,AT^2 + AZ^2$ の逆の比になる。というのは、時間は微小部分を通る速さに反比例し、この速さは太陽と交点の速さの和だからである。ゆえに、交点が動かない場合に太陽が弧 $NA$ を通る時間を扇形 $NTA$ で表し、極めて小さな弧 $Aa$ を通る時間の微小部分を扇形の微小部分 $ATa$ で表すことにする。そして、$aY$ を $Nn$ に垂線として下ろしたとき、$AZ$ 上に $dZ$ を、長方形 $dZ \times ZY$ が扇形の微小部分 $ATa$ に対して $AZ^2$ 対 $9.0829032\,AT^2 + AZ^2$ の比になるように、すなわち $dZ$ 対 $\frac{1}{2}AZ$ が $AT^2$ 対 $9.0829032\,AT^2 + AZ^2$ になるようにとると、長方形 $dZ \times ZY$ は、弧 $Aa$ を通る全時間における、交点の運動から生じる時間の減少を表す。そして点 $d$ が曲線 $NdGn$ をたどるなら、曲線で囲まれた面積 $NdZ$ は、弧 $NA$ 全体を通る時間の全減少となり、したがって扇形 $NAT$ が面積 $NdZ$ を超える分が、その時間全体になる。また、交点の運動は時間が短いほど、時間の比で小さいのだから、面積 $AaYZ$ も同じ比で減らさなければならない。これは、$AZ$ 上に長さ $eZ$ を、$AZ$ に対して $AZ^2$ 対 $9.0829032\,AT^2 + AZ^2$ となるようにとれば、そうなる。実際、そのとき長方形 $eZ \times ZY$ は、面積 $AZYa$ に対して、弧 $Aa$ を通る時間の減少が、交点が静止している場合に通る全時間に対するのと同じ比になる。したがって、この長方形は交点の運動の減少に対応する。そして点 $e$ が曲線 $NeFn$ をたどるなら、減少の総和である面積全体 $NeZ$ は、弧 $AN$ を通る時間における全減少に対応し、残りの面積 $NAe$ は、残りの運動に対応する。それは、弧 $NA$ 全体を、太陽と交点の合わさった運動で通る時間における、交点の真の運動である。

さて、もし円の半径 $AT$ を1とすれば、半円の面積は $1.570796$ であり、図形 $NeFnT$ の面積は、無限級数の方法で求めると $0.1188478$ となる。ところで、円全体に対応する運動は $19$ 度 $49\minute 2\second 49\third\frac{1}{2}$ であったから、図形 $NeFnT$ の2倍に対応する運動は $1$ 度 $29\minute 57\second 51\third\frac{1}{2}$ である。これを先の運動から引くと、交点が太陽と会合してから次に会合するまでの交点の全運動 $18$ 度 $19\minute 4\second 58\third$ が残る。そしてこの運動を太陽の年周運動 $360$ 度から引けば、同じ会合のあいだの太陽の運動 $341$ 度 $40\minute 55\second 2\third$ が残る。この運動は年周運動 $360$ 度に対して、いま求めた交点の運動 $18$ 度 $19\minute 4\second 58\third$ が、その年周運動に対するのと同じ比をなす。したがって、後者は $19$ 度 $18\minute 0\second 22\third$ となる。これが恒星年における交点の平均運動である。同じ値は、天文表では $19$ 度 $20\minute 31\second 1\third$ である。差は全運動の400分の1より小さく、月の軌道の離心率と、黄道面に対する傾きから生じるように見える。軌道の離心率によって交点の運動は速くなりすぎ、その傾きによってかえってわずかに遅くなり、ちょうどよい速さにもどる。

\condpagelarge{命題 XXXIII. 問題 XIII.}

\textit{月の交点の真の運動を求めること。}

時間が面積 $NTA - NdZ$ に比例する時間で（前の図において）、この運動は面積 $NAeN$ に比例し、そこから与えられる。しかし計算が非常に難しいので、この問題の次の作図を用いるのがよい。

\pngright{444.png}{}{}{}

中心 $C$、任意の半径 $CD$ で円 $BEFD$ を描く。$DC$ を $A$ まで延長して、$AB$ 対 $AC$ が、平均運動対、交点が矩にあるときの真の平均運動の半分、になるようにする（すなわち $19$ 度 $18\minute 0\second 22\third$ 対 $19$ 度 $49\minute 2\second 49\third\frac{1}{2}$ で、したがって $BC$ 対 $AC$ は、運動の差 $0$ 度 $31\minute 2\second 27\third\frac{1}{2}$ 対後の運動 $19$ 度 $49\minute 2\second 49\third\frac{1}{2}$、すなわち $1$ 対 $38\frac{1}{3}$ である）。ついで、点 $D$ を通って円に $D$ で接する無限の直線 $Gg$ を引き、平均運動で求めた太陽と交点の位置との距離の半分に等しい角 $BCE$ または $BCF$ をとり、$AE$ または $AF$ を引いて垂線 $DG$ と $G$ で交わらせる。そして、交点の朔望のあいだの運動（すなわち $9$ 度 $10\minute 40\second$）に対して、接線 $DG$ が円 $BED$ の円周全体に対するのと同じ比をなす角をとり、その角を交点の平均運動に加えれば、その真の運動が得られる。実際、こうして求めた真の運動は、時間を面積 $NTA - NdZ$ で、交点の運動を面積 $NAeN$ で表して得られる真の運動に、ほぼ一致する。よく考えてみればわかる。これが交点の運動の年均差である。月均差もあるが、それは月の緯度を求めるうえでは、まったく必要でない。というのは、月の軌道の黄道面に対する傾きの変化には、二つの不等、一つは年周の、もう一つは月周のものがあるが、その月周の不等と、交点の月均差とは、互いに打ち消し合い補い合って、月の緯度を決めるときには、どちらも無視できるからである。

\textit{系.}\wsp{}この命題と前の命題から、交点は朔望では静止し、矩では一時間あたり $16\second 18\third 41\fourth\frac{1}{2}$ で後退することが明らかである。また、八分点での交点の運動の均差は $1$ 度 $30\minute$ である。これらはすべて天体の現象とよく合う。
\condpagelarge{命題 XXXIV. 問題 XIV.}

\textit{月の軌道の黄道面に対する傾きの、一時間あたりの変化を求めること。}

\pngright{445.png}{}{}{}

$A$ と $a$ を朔望、$Q$ と $q$ を矩、$N$ と $n$ を交点とし、$P$ を月の軌道上の位置、$p$ をその位置の黄道面への射影、$mTl$ を前のとおり交点の瞬間の運動とする。そして線 $Tm$ に垂線 $PG$ を下ろし、$pG$ を結んで、それが $Tl$ と $g$ で交わるまで延長し、$Pg$ も結ぶとする。すると角 $PGp$ は、月が $P$ にあるときの月の軌道の黄道面に対する傾きであり、角 $Pgp$ は時間の一瞬が過ぎたあとの同じ傾きであって、したがって角 $GPg$ は傾きの瞬間的な変化である。ところで、この角 $GPg$ は角 $GTg$ に対して、$TG$ 対 $PG$ と $Pp$ 対 $PG$ との積の比をなす。したがって、時間の一瞬のかわりに一時間をとれば、角 $GTg$ は（命題XXXにより）$33\second 10\third 33\fourth$ に対して $IT \times PG \times AZ$ 対 $AT^3$ の比をなすから、角 $GPg$（すなわち傾きの一時間あたりの変化）は $33\second 10\third 33\fourth$ に対して、$IT \times AZ \times TG \times \frac{Pp}{PG}$ 対 $AT^3$ の比をなす。\QED

これは、月が円軌道を一様に回るという仮定の上でのことである。もしその軌道が楕円ならば、交点の平均運動は、上で述べたように、短軸対長軸の比で小さくなる。また同じ比で正弦 $IT$ も小さくなる。ところが傾きの変化は、正弦 $IT$ の減少によって増す分だけ、交点の運動の減少によって減るから、前と変わらない。

\textit{系1.}\wsp{}$Nn$ に垂線 $TF$ を立て、$pM$ を月の黄道面での一時間あたりの運動とし、$QT$ に下ろした垂線 $pK$, $Mk$ を両側に延長して、$TF$ と $H$, $h$ で交わらせるとする。すると $Kk$ 対 $Mp$ は $pK$、すなわち $IT$ 対 $AT$ であり、$TZ$ 対 $AT$ は $TG$ 対 $Hp$ であるから、$IT \times TG$ は $\frac{Kk \times Hp \times TZ}{Mp}$ に等しい。すなわち、面積 $HpMh$ に比 $\frac{TZ}{Mp}$ を掛けたものに等しい。したがって、傾きの一時間あたりの変化は $33\second 10\third 33\fourth$ に対して、$HpMh$ に $AZ \times \frac{TZ}{Mp} \times \frac{Pp}{PG}$ を掛けたもの対 $AT^3$ の比をなす。

\textit{系2.}\wsp{}したがって、地球と交点が一時間ごとに、その時間が終わると新しい位置から引きもどされ、ただちに元の位置にもどされて、一か月の周期全体にわたってその位置が一定に保たれるとすれば、その一か月の傾きの変化の全体は、$33\second 10\third 33\fourth$ に対して、点 $p$ が回るあいだに生じ、それぞれの符号 $+$ と $-$ をつけて合わせた面積 $HpMh$ の総和に $AZ \times TZ \times \frac{Pp}{PG}$ を掛けたもの対 $Mp \times AT^3$ の比をなす。すなわち、円全体 $QAqa$ に $AZ \times TZ \times \frac{Pp}{PG}$ を掛けたもの対 $Mp \times AT^3$、あるいは円周 $QAqa$ に $AZ \times TZ \times \frac{Pp}{PG}$ を掛けたもの対 $2MP \times AT^2$ の比をなす。

\textit{系3.}\wsp{}したがって、交点の位置が与えられたとき、一か月にわたって一様に続ければその月ごとの変化が生じうる、一時間あたりの平均の変化は、$33\second 10\third 33\fourth$ に対して $AZ \times TZ \times \frac{Pp}{PG}$ 対 $2AT^2$ の比をなす。すなわち（$Pp$ 対 $PG$ が前述の傾きの正弦対半径であり、$\frac{AZ \times TZ}{AT}$ 対 $\frac{1}{2}AT$ が角 $ATn$ の2倍の正弦対半径であるから）その傾きの正弦に、交点の太陽からの距離の2倍の正弦を掛けたもの対、半径の平方の4倍、の比をなす。

\textit{系4.}\wsp{}交点が矩にあるときの傾きの一時間あたりの変化は、（上の命題により）角 $33\second 10\third 33\fourth$ に対して $IT \times AZ \times TG \times \frac{Pp}{PG}$ 対 $AT^3$、すなわち $\frac{IT \times TG}{AT} \times \frac{Pp}{PG}$ 対 $AT$、すなわち月の矩からの距離の2倍の正弦に $\frac{Pp}{PG}$ を掛けたもの対、半径の2倍、の比をなす。したがって、交点がこの位置にあるときに、月が矩から朔望へ進むあいだ（すなわち $177\frac{1}{6}$ 時間のあいだ）の一時間あたりの変化の総和は、同じ数の角 $33\second 10\third 33\fourth$ の総和、すなわち $5878\second\frac{1}{2}$ に対して、月の矩からの距離の2倍の正弦の総和に $\frac{Pp}{PG}$ を掛けたもの対、同じ数の直径の総和、すなわち、直径に $\frac{Pp}{PG}$ を掛けたもの対円周、の比をなす。すなわち、傾きが $5$ 度 $2\minute$ であれば、$7 \times \frac{876}{10000}$ 対 $22$、すなわち $279$ 対 $10000$ となる。したがって、先に述べた時間における一時間あたりの変化の総和から成る全変化は $164\second$、すなわち $2\minute 44\second$ である。

\condpagelarge{命題 XXXV. 問題 XV.}

\textit{与えられた時刻における、月の軌道の黄道面に対する傾きを求めること。}

\pngright{450.png}{}{}{}

$AD$ を最大の傾きの正弦、$AB$ を最小の傾きの正弦とする。$BD$ を $C$ で二等分し、中心 $C$、半径 $BC$ で円 $BGD$ を描く。$AC$ 上に $CE$ を、$CE$ 対 $EB$ が $EB$ 対 $2BA$ となるようにとる。そして、与えられた時刻に、交点の矩からの距離の2倍に等しい角 $AEG$ をつくり、$AD$ に垂線 $GH$ を下ろす。すると $AH$ が求める傾きの正弦である。

なぜなら
\[ GE^2 = GH^2 + HE^2 = BH \times HD + HE^2 = BH \times BD + HE^2 - BH^2 \]
\[ = BH \times BD + BE^2 - 2BH \times BE = BE^2 + 2EC \times BH = 2EC \times AB + 2EC \times BH = 2EC \times AH \]
だからである。ゆえに $2EC$ は与えられているので、$GE^2$ は $AH$ に比例する。いま、与えられた時間の一瞬が過ぎたあとの、交点の矩からの距離を $AEg$ とすれば、角 $GEg$ が与えられているから、弧 $Gg$ は距離 $GE$ に比例する。ところが $Hh$ 対 $Gg$ は $GH$ 対 $GC$ であるから、$Hh$ は積 $GH \times Gg$、すなわち $GH \times GE$ に比例する。すなわち $\frac{GH}{GE} \times GE^2$、すなわち $\frac{GH}{GE} \times AH$ に比例し、これは $AH$ と角 $AEG$ の正弦との積に比例する。それゆえ、もしある場合に $AH$ が傾きの正弦であるなら、それは（上の命題の系3により）傾きの正弦と同じだけ増すから、つねにその正弦に等しいままである。ところが、点 $G$ が点 $B$ か $D$ に落ちるときには、$AH$ はこの正弦に等しいから、つねにそれに等しいままである。\QED

この証明では、交点の矩からの距離である角 $BEG$ が、一様に増すと仮定した。不等のこまかい点をすべて検討する暇はないからである。いま、角 $BEG$ が直角であり、$Gg$ が交点と太陽との距離の一時間あたりの増加であるとすれば、傾きの一時間あたりの変化は（直前の命題の系3により）$33\second 10\third 33\fourth$ に対して、傾きの正弦 $AH$ と、直角 $BEG$（これは交点の太陽からの距離の2倍）の正弦との積対、半径の平方の4倍、の比をなす。すなわち、平均の傾きの正弦 $AH$ 対、半径の4倍、の比をなす。すなわち（その平均の傾きがほぼ $5$ 度 $8\minute\frac{1}{2}$ だから）その正弦 $896$ 対、半径の4倍 $40000$、あるいは $224$ 対 $10000$ である。ところで、正弦の差 $BD$ に対応する全変化は、その一時間あたりの変化に対して、直径 $BD$ 対弧 $Gg$ の比をなす。すなわち、直径 $BD$ 対半円周 $BGD$ と、交点が矩から朔望へ進む時間 $2080$ 時間対 $1$ 時間との積の比、すなわち $7$ 対 $11$ と $2080$ 対 $1$ との積の比である。したがってすべての比を合わせると、全変化 $BD$ は $33\second 10\third 33\fourth$ に対して $224 \times 7 \times 2080$ 対 $110000$、すなわち $2965$ 対 $100$ となり、そこから変化 $BD$ は $16\minute 24\second$ となる。

これは、月の軌道上の位置を考えないときの、傾きの最大の変化である。というのは、交点が朔望にあるとき、傾きは月の位置が変わっても変化しない。だが交点が矩にあるときは、傾きは、月が朔望にあるほうが、矩にあるときより、$2\minute 44\second$ だけ大きい。これは前の命題の系4で述べたとおりである。そしてこの超過の半分 $1\minute 22\second$ によって、平均の全変化 $BD$ は、月の矩では減って $15\minute 2\second$ に、月の朔望では増えて $17\minute 46\second$ になる。したがって、月が朔望にあるとすれば、交点が矩から朔望へ移るときの全変化は $17\minute 46\second$ であり、それゆえ、交点が朔望にあるときの傾きを $5$ 度 $17\minute 46\second$ とすれば、交点が矩にあり月が朔望にあるときの傾きは $5$ 度となる。そして、このことは観測によって確かめられる。実際、天文学者は、月の軌道の黄道面に対する傾きを、交点が矩にあって月が太陽の衝にあるときは、ほぼ $5$ 度としている。そして交点が朔望にあるときは、それが $5$ 度 $17\minute\frac{1}{2}$ か $5$ 度 $18\minute$ であると教えている。

いま、月が朔望にあり、交点がどこにあってもよいときの軌道の傾きを求めたいなら、$AB$ 対 $AD$ を、$5$ 度の正弦対 $5$ 度 $17\minute 46\second$ の正弦となるようにし、角 $AEG$ を交点の矩からの距離の2倍にとれば、$AH$ が求める傾きの正弦である。この軌道の傾きに等しいのが、月が交点から $90$ 度離れているときの同じ軌道の傾きである。月のほかの位置では、傾きの変化が許す月ごとの不等は、月の緯度の計算において、（上で述べたように）交点の運動の月ごとの不等によって補われ、ある意味で取り除かれるので、その緯度の計算では無視してよい。

\condpagelarge{注.}

ここまでは、軌道の離心率を考えない場合の月の運動について述べた。同様の計算によって、私は次のことを見いだした。遠地点は、太陽との合または衝にあるときは、恒星に対して毎日 $23\minute$ 進み、矩にあるときは毎日ほぼ $16\frac{1}{3}$ 分後退する。またその平均運動は、一年にほぼ $40$ 度である。フラムスティード氏がホロックスの仮説に合わせて作った天文表によれば、遠地点は朔望では毎日 $24\minute 28\second$ 進み、矩では毎日 $20\minute 12\second$ 後退し、年平均運動は $40$ 度 $41\minute$ で先行方向に運ばれる。朔望での遠地点の毎日の前進と、矩での毎日の後退との差は、表では $4\minute 16\second$、私たちの計算では $6\minute\frac{2}{3}$ であるが、これは表の欠陥によるものと私たちは疑う。しかし、私たちの計算も十分に正確だとは考えない。というのは、ある仕方で計算したところ、朔望での遠地点の毎日の前進と、矩での毎日の後退は、いくぶん大きく出たからである。だが、あまりにもこみいっていて近似に妨げられ、しかも十分に正確でない計算は、載せたくない。
\textbf{第3編 月の摂動論、分点歳差、および彗星の幾何学軌道論（命題XXXVI、証明XVI）}

月の apoapsis （遠地点）における日進運動は、表により \(4\) 分単位とされるが、私たびの計算によれば \(6\) 分 \(\frac{2}{3}\) である。この差異は表の誤りによるものと考えられる。しかし、私たちの計算も十分正確とは言えないと考えている。なぜなら、遠地点における日進運動と近地点における日退運動の比率が若干大きいことを発見したからである。しかしながら、計算結果があまりに複雑で、近似値を導くために妨げられるため、これらの計算を提供するのも好ましくない。

\textbf{Prop.\ XXXVI. Prob.\ XVI.}

\textbf{太陽の力による海面の動揺を求める。}

太陽の力 \( ML \) または \( PS \)，月食円において月の運動を攪乱する力は（この書籍の命題XXVにより），我々の地表における重力に対する比として \( 1 : 638092.6 \) である。そして，新月と満月において \( SM - LM \) または \( 2PK \) の力はその倍となる。

しかしながら，これらの力を地表まで下げるとき，それは地球の中心から距離に比例して減少し，すなわち \( 60.5 : 1 \) の比で減少する。したがって，前記の力は地表において重力に対する比として \( 1 : 38604600 \) となる。この力により海面は引き下がる。


\par\begin{center}
    \centering
    \includegraphics[width=1486px,height=775px,trim=-12 0 0 0,clip]{451.png}

\end{center}\par

\[
\text{在太阳正東或正西約90度的位置上，}
\]
太陽の引力以外に二倍大的力が潮汐に対して作用する。この力は、約90度の位置において水位を低下させるとき、対面の位置では水位を高める力と等しいものがある。したがって、これら力を総和としたものは、太陽に対する潮汐の引力に相当する；もし単に対面の位置における水位の上昇のみを考えるとしても、その効果は同様であり、約90度の位置においては作用しない。

\[
\textbf{訳文：}
\]

在太陽から約90度の位置にあるとき、太陽の引力以外に二倍大きな力が潮汐を引き上げます。この力を用いて水位が約90度の位置で低下する場合、同じ効果は太陽の反対側で水位を引き上げることによっても得られます。したがって、これらの総和の力は太陽が潮汐に作用する力となります；もしすべての力を太陽の反対側で水位を引き上げるためだけに使用すれば、その効果は同じであり、約90度の位置では全く働きません。

\[
\textbf{数学的注釈：}
\]

ここで、太陽の引力と二倍大きな力の比は1対12868200です。つまり、太陽の引力が$F_{\text{重力}}$であるとき、二倍大きな力は$2F_{\text{重力}}$となります。

\[
F_{\text{重力}} : 2F_{\text{重力}} = 1 : 12868200
\]

この比から、太陽の引力が約90度の位置で水位を引き下げるときと、二倍大きな力が太陽の反対側で水位を引き上げるときの効果は等しいことがわかります。


\textbf{補遺}$\wsp{}$従って、地球の自転によって生じた離心力は、重力に対して $1 : 291$ の比で作用し、赤道付近の水の高さが極地の水の高さを $85,200$ パリスの足長単位で測った場合、太陽の引力について述べたものでは、その引力に対して離心力は $1 : 12{,}868{,}200$ の比で作用し、従って水の高さが太陽の下または太陽に反対側にある地域でその高さを $90^\circ$ 遠く離れた場所よりも $44,221$ 倍超えることになる。

\textbf{訳注}\wsp{}パリスの足長単位は、当時のフランスの標準的な測定単位で、約 \(0.3248\) メートルに相当します。したがって、水の高さの差は非常に大きいことを示しています。

\textbf{図注}\wsp{} ここで示す図では、赤道付近と極地での水の高さの比較が示されています。

\textbf{第37命題. 問題17.}

\textbf{パリの足一本と11指で月の力を見つけ出す。}

月が海に及ぼす力を求めるためには、その力が太陽の力に対する比から得られることを確認する必要があります。この比は、これらの力によって引き起こされる海の運動の比から求めることができます。春分点と秋分点の交換時期および対立時期において、アヴォン川の入口付近、ブリストルより南第三の石碑まで、夏と秋の間に水位が上昇する全体的な量を考慮に入れます。

\[
\textbf{パリの足一本} \& 11\textbf{指} = 85200 \textbf{尺} : 44221
\]

サミュエル・シュトゥルミウスは、その大きさが約$45$フィートで、四平方では$25$フィートです。前後の高さは、最初の総力と最後の総力から生じます。したがって、太陽と月が赤道を通過し、地球から中程度に離れている場合、力を$S$と$L$とします。そして、月は四平方において、春分と秋分の時期には赤道から約$23$度半の傾斜で回転し、赤道から傾斜して光る力は海への移動力が小さくなります。これは（私の感覚では）、余弦補角の二乗比例に近いです。

\[
S \quad \text{および} \quad L
\]

そして、月は春分と秋分の時期には赤道から約23度半の傾斜で回転し、

\[
L \propto \sin^2(90^\circ - 23.5^\circ) = \cos^2(23.5^\circ)
\]

この式は、月の引力が海への移動力をどのように与えるかを示しています。

在四次方程式中，（当正弦値与半径之比為 91706 : 100000 時）
$\frac{841}{1000} L$ 為解，且在冲日前后両行星之力的総和為 \(L + S\)，其四次方差為 \(\frac{841}{1000} L - S\)。したがって，
$L + S$ 与 \(\frac{841}{1000} L - S\) 的比値為 45 : 25 或 9 : 5，従而
$$5L + 5S = \frac{7569}{1000} L - 9S$$，
且
$14S = \frac{2569}{1000} L$。したがって，\(L\) 与 \(S\) 的比値為 14000 : 2569 あるいは説是 \(5\frac{7}{15}\) : 1。
在普利茅斯港（〜によれば薩謬尔·科尔普雷斯的観測）的潮汐海面高度，在中等深度時，比較平均水位高出約十六英尺六英寸。

月の潮汐高さについて、春と秋には満弦や新弦の位置で太陽からの潮汐高さが約7または8フィートを上回ることがあります。もし中程度の余分な高さが春と秋に7.5フィートである場合、満弦や新弦の位置では潮汐高さは19.75フィートとなり、四分点では12.25フィートとなります。したがって、$L + S$ は $841/1000 L - S$ に比例し、$19\frac{3}{4}$ は $12\frac{1}{4}$ に比例します。さらに、$L$ は $S$ に対して734対100、すなわち7.33対1の比率となります。

したがって、月の引力と太陽の引力を前計算で比すると5.4667対1となり、後期の計算では5.3333対1となります。

\[
\textbf{ut } 7\frac{1}{3} \textbf{ ad } 1.
\]
Donec aliquid certius ex Observationibus accuratius institutis constiterit, usurpabimus proportionem mediocrem \(6\frac{1}{3}\) ad 1.

unde cum vis Solis sit ad vim gravitatis ut 1 ad 12868200, vis Lunae erit ad vim gravitatis ut 1 ad 2031821.

\textbf{Corol.\fsp{}1.}\wsp{}igitur cum aqua vi Solis agitata ad altitudinem pedis unius \& undecim digitorum ascendat, eadem vi Lunae ascendet ad altitudinem pedum duodecim.
Tanta autem vis ad omnes maris motus excitandos abunde sufficit, \& quantitati motuum probe respondet.

月の東から西へ広がる海域において、太平洋や大西洋およびエチオピア海の一部では、赤道を越えて太陽潮汐の影響を受け、水位は6フィート、9フィート、12フィート乃至15フィートほど上昇することがよくあります。しかし、太平洋はより深い海域であり広大であるため、その潮の高さは大西洋やエチオピア海よりも大きいとされています。実際、満潮が発生するときには、東から西にかけての海洋の幅が90度以上あることが必要です。一方で、エチオピア海において赤道内では水位の上昇はそれほど大きくないことが確認されています。

月の摂動論、分点歳差、および彗星の幾何学軌道論（命題XXV〜XLII）における第84部

このように、アフリカとアメリカ半島の南部との間の海が狭いことから、水は海沿岸両側、東方と西方同時に上昇し難く、しかし狭い海域においては、時々交替して沿岸に下流することを強いられる。その結果、遠く離れた島嶼における満潮と干潮は通常非常に微弱なものとなる。また、港湾において水が大きな勢いで流れ込む場所では、湾の内部へと定期的に充填され、排出されることにより、満潮と干潮の影響が現れる。



このようにして、海沿岸に近い地域では潮汐の効果が顕著となる一方で、遠く離れた島嶼ではその効果は微弱なものとなる。

\& 力が放散され、流入と流出は通常より大きいとされています。例えば、イギリスのプルマス港やチェプストウ橋周辺；ノルマンディーの聖ミヒール山とアーブランコワ（俗にオランジェ）市周辺；インド東部のカンバヤムとペグ周辺です。これらの場所では、大海が大きな速度で流入し、流出すると、沿岸は数万メートルの範囲でしばしば浸水したり乾燥したりします。流入と流出の力は、水位が30、40乃至50メートル以上上昇または下降するまでにその勢力を失うことはありません。そして、長方形湾の比例如しくなります。

\& このような港や海峡では、潮流が強力な流れと反復により増幅される。しかし、降下した潮は深いそして開けた海に注ぎ込むことで、勢いなく上昇し、再び下降することができる。これらの潮の高さは太陽と月の引力によって比例する。

\textbf{補則2.} 潮流を引き起こす月の引力は、 Pendulum 実験や静力学実験における力に比べて遥かに小さいことを理解している。

\textit{Corol.\,2.}\quad 月による海への引潮力は、重力に対する比が \(1 : 2031821\) であるため、その力は Pendulum 実験や静力学実験における力に比べて遥かに小さいことが明白である。

【第3編 月の摂動論、分点歳差、および彗星の幾何学軌道論（命題XXV〜XLII）】

\textbf{命題XXXIII}

水静力学においても同様に理解できる。海面上でのこの力は感覚的に観察される。

\textit{補遺3.} 月が海を動かす力は太陽の力を $6\dfrac{1}{3}$ に対して $1$ と似たものであり、これらの力は月と太陽の密度とその視直径の立方の積に比例する。したがって、月の密度に対する太陽の密度は $6\dfrac{1}{3}$ に対して $1$ で、太陽の直径の立方に対する月の直径の立方は逆比例し、すなわち（太陽と月の視直径はそれぞれ $31\minute.27\second.$ および $32\minute.12\second.$ である）：

\[
\textbf{月の密度} : \textbf{太陽の密度} = 6\dfrac{1}{3} : 1
\]
\[
(\textbf{太陽の直径})^3 : (\textbf{月の直径})^3 = 1 : 6\dfrac{1}{3}
\]


$\textbf{34}$ に至るまで。太陽の密度は地球の密度に対して $100 : 387$ であり、従って月の密度は地球の密度に対して $600 : 387$ すなわち $9 : 5$ 左右である。したがって、月の体はより密で、またより地球的なものとみなされるべきである。

\textit{補則 \textbf{4}．}従って、実際の月の直径は実際の地球の直径に対して $1 : 3.65$ であり、したがって月の質量は地球の質量に対して $1 : 26$ 左右である。

\textit{補則 \textbf{5}．}また、月面上での加速度引力は、


\par\begin{center}
    \centering
    \includegraphics[width=0.8\textwidth]{455.png}

\end{center}\par

このように推定される。

\textbf{第38命題. 問題 XVIII.}

\textit{月の形状を求める。}

もし月の体が海のように流体であると仮定すると、地球の引力はその流体部分を表面近くと遠くに急速に引き上げる力に対して、月の引力によって海が月の正面と反対側で引き上げられる力と等しいものであり、すなわち、月の向心加速度重力と地球の向心加速度重力との比は月の直径と地球の直径との積に比例し、具体的には26対1および5対18の積に相当する。

\textbf{訳文}

次に、月の海が12フィートまで上昇するとき、 lunar引力により液体月球は約90フィートまで上昇すべきである。この理由から、月の形は、その最大直径が地球の中心を通ってさらに180フィート余りを越える円錐体となるべきである。したがって、月はこのような形状を有するのであり、それは最初からそのような形状を有していたと考えるべきである。\QEIit

\textit{補遺} そこで、常に同じ月の面が地球に対して向かっていることが明らかになる。なぜなら、他の位置においては月体は静止できないし、常にこの位置に戻るための振動を行うからである。しかし、これらの振動は

\textbf{訳注}

\textit{補遺} ここでは、月の面が地球に対して常に同じ方向を向いていることが示される。なぜなら、月体は他の位置において静止することはなく、常にこの特定の位置に戻るために振動を行うからである。しかし、これらの振動は

（訳注：ニュートンはここで、月の自転と公転が一致しているため、同じ面が地球に対して常に向いていることを説明しています。）

【第3編 月の摂動論、分点歳差、および彗星の幾何学軌道論（命題XXV〜XLII）】

\textbf{命題XXXIII. §1.}

地球に対する引力が非常に小さいため、その影響は遅く及ばない。したがって、常に地球を向いている面は、月の軌道の中心から見れば、それほど速く変化せず、すぐに引き戻されず、また地球に戻ることもない。

\textbf{証明}


\par\begin{center}
  \centering
  \includegraphics[width=964pt,height=1252pt,trim={0 -24 0 0},clip]{458.png}

\end{center}\par

もし $APEp$ を均質な地球とし、その中心を $C$、極点を $P$, $p$ および赤道 $AE$ で表し、かつ中心 $C$ から半径 $CP$ で球面 $Pape$ を描くとする。また平面 $QR$ を考え、その中に中心から引かれた直線があると仮定する。

この証明は、地球に対する引力が非常に小さいため、常に地球を向いている面が月の軌道の中心から見ればすぐに引き戻されることなく、また地球に戻ることもないことを示す。

\textbf{太陽は地球の中心に垂直方向から作用し、} \& 地球全体の外側にある点 \( P \), \( A \), \( P' \), \( E \) が、近似して球面を描いた点群より高い位置にあるため、各個体は平面 \( QR \) から遠ざかる方向に作用しようとします。これらの個体のそれぞれが同じ距離から平面 \( QR \) に存在する場合、その全効果は地球を円周運動させる力として考えられます。

\[
\text{全個体の合力} = \text{各個体の効果}
\]

ここで、地球を円軌道上において均一に配置した環状物質が地球を同じ円運動させることを考えると、これらの合力は以下の関係になります：

\[
\text{合力} = 4 \times \text{環状物質の全効果}
\]

これは、地球を赤道円周 \( AE \) 上において均一に配置した場合の全効果の四倍となります。


月の運動が円周上にある軸 \( QR \) に属する平面 \( QR \) 上および直線 \( Pp \) に対して垂直である場合について述べます。

まず、赤道 \( AE \) と平行な円 \( IK \) を考えます。さらに、地球の粒子 \( L \) が円 \( IK \) の外側に球体 \( Pape \) 外部にあるとします。そして、平面 \( QR \) に対して垂直な直線 \( LM \) を引くと、その粒子 \( L \) に対する全引力は \( LM \) に比例します。この引力 \( LM \) を \( LN \), \( NM \) の二つに分解すると、地球の周囲を球体 \( Pape \) 外部で回るすべての粒子 \( L \) の効果的な力 \( MN \) は、これらの力を含む円周運動全体に関連します。

月をその中心を中心にした位置 $A p E P$ に変換するために必要な力の効果は、各個々の力 $L$ をその中心を中心にした位置 $L N$ に変換する効果に対する比が $3:2$ であるため、すべての力 $M N$ の効果は、この上記の効果 $L N$ よりも大きい効果に対する比が $3:1$ となるように変換されるべきです。したがって、すべての力 $M N$ の総効果は、すべての力 $L N$ の総効果より大きい効果の余剰効果に対する比が $3:4$ となります。

もしこれらの各個々の粒子が赤道に位置していた場合、すべての力 $L N$ の効果は消滅し、すべての力 $M N$ の効果は $4:3$ の比率で増加します。したがって、その余剰効果は、上記の効果よりも大きい効果に対する比が $3:4$ となります。

\textbf{絶対的な特定物体の位置における効力の四分之一は赤道における同じ特定物体の効力の一部を形成します。}しかし、赤道での昼夜運動はこの効力と等しいです。誰が欲しければ個々を検討してください。簡略化のため。

\textbf{命題II．補題}

地球全体の軸周りの運動は、すべての特定物体の運動から成る合算であり、その運動は環の軸周りの運動に対して、物質の比（地球における物質に対する環における物質）と、任意の円周の四分之一弧の平方比（二つの平方）との合成比に従います。


\par\begin{center}
  \centering
  \includegraphics[width=0.8\textwidth]{457.png}

\end{center}\par

\textbf{直径において；すなわち、材料の比に於いては 925275 と 1000000 の比です。}

蓋し、円筒が静止した軸を中心に回転する運動は、内接球体と共に回転する運動に対して、任意の四つの等しい正方形に対する三つの内接円の正方形の比に等しくなります；また、円筒の運動は薄い環の運動と、円筒と環との共通接触点を巡る運動に対し、二倍の円筒の材料に対して三倍の環の材料となります；そしてその環の運動が円筒の軸を中心に均一に継続される運動に対しては、同一時間における円筒の直径周りの均一な運動と等しくなります。




\section*{第3編 月の摂動論、分点歳差、および彗星の幾何学軌道論（命題XXV〜XLII）}


周期を考慮して、円周に対する直径の二倍と等しいものとする。

\bigskip

\textbf{引理III.}もし地球が除去され、その他の部分がすべて取り除かれた場合に、月のみが太陽を中心に地球の公転運動をし、また同時に自軸を中心に $ecliptic$面に対して $23^\circ 30'$ の傾斜角で日回り運動を行うとする。このとき、春分点と秋分点の運動は同様であり、もし該円環が流体であるか硬質な堅固な素材から構成されているかに関わらず同じである。

\bigskip

\textbf{命題XXXIX. 證明XIX.}求める春分点歳差。



月のノードの運動（中庸時間）は、ノードが四分点にあり、円軌道上にあるとき、$16^\circ.35'36''$であり、その半分は$\textbf{8}^\circ.17'38''18''$である（理由と先に説明した通り）；これは円軌道上のノードの運動の中庸時間運動であり、一回の太陽年では$20^\circ\ 11'\ 46''$となる。従って、月のノードがこのような軌道上にあるとき、一回の太陽年に$20^\circ\ 11'\ 46''$で進むことになる；もし各月のノードの運動が複数あれば、命題$\textbf{LXVI}$の補則$\textbf{16}$によれば。

第3編 第25命題

月の運動軸の傾きの年間変動は、太陽面に沿って1日を約1秒回る場合と比較して、20度11分46秒となる。これは、太陽1日の時間が23時56分であり、月の周期日数が27日7時間43分であるときと等しい。すなわち、1436対39343の比率である。また、環状軌道上にある月の運動軸の傾きの年間変動の比も同じである；すなわち、月同士が互いに接触しない場合、または液状となり連続した環となる場合、または最後にはその環が硬化し剛体となる場合においても等しい。

\[
\textbf{証明}
\]

まず、太陽面に沿って1日を約1秒回る月の運動軸の傾きの年間変動は20度11分46秒であると仮定する。これは、太陽1日の時間が23時56分であり、月の周期日数が27日7時間43分であるときの比を表している。

次に、環状軌道上にある月の運動軸の傾きの年間変動について考えると、以下の比率を考える：
\[
1436 : 39343
\]

これは、太陽面に沿って1日を約1秒回る場合と同様である。したがって、環状軌道上にある月の運動軸の傾きの年間変動も20度11分46秒となる。

最後に、月同士が互いに接触しない場合、または液状となり連続した環となる場合、または最後にはその環が硬化し剛体となる場合においても、この比率は変わらない。したがって、これらの条件下でも月の運動軸の傾きの年間変動は20度11分46秒である。

\[
\blacksquare
\]

\begin{center}
\textbf{第3編 月の摂動論、分点歳差、および彗星の幾何学軌道論}

\textit{命題XXV. 問題}

我々は次のように定める：円錐体 $PapAPepE$ の表面に等質な太陽系を配置し、その上部球体 $PapE$ が地球に対応するものとする。ここで、球体 $PapE$ は地球に対して直径 $aC$ に対する直径 $AC$ の $689:692$ の比で存在する。すなわち、小径 $PC$ や $aC$ が大径 $AC$ の $689:692$ の比であるとすると、球体の直径は $4143:474721$ の比で地球の直径に等しい。したがって、この比は $1000:114585$ となる。もしこの円錐体の表面が赤道線に沿って地球を包むとし、両者は円錐体の直径について同時に回転する場合、

\end{center}


\par\begin{center}
\centering
\includegraphics[width=964pt,height=1252pt,trim=-24 0 0 0,clip]{458.png}

\end{center}\par

月が回転すると、環の運動は内球体の運動（第2補題による）に比例し、すなわち \(4143 : 474721 \& 1000000 : 925275\) の比で、つまり \(4143 : 439248\) である。したがって、環の運動は環と内球体の総運動に比例し、すなわち \(4143 : 443991\) である。

もし環が内球体に追随し、その等分点（或いは赤道点）の回帰運動と共に運動する場合、残存する環の運動は以前の運動に対して \(4143 : 443391\) の比となる。したがって、赤道点の運動はこの比率で減少する。

よって、環の年間運動は

\[
\textbf{443391} \quad (\text{ラテン語原文: motus annuus punctorum})
\]

月の摂動論、分点歳差、および彗星の幾何学軌道論（命題XXV〜XLII）

\textbf{命題XXXIII}

地球と環から成る天体について、その運動は$20^\circ$11′46″であり、1436年39343日と4143年443391日に一致し、すなわち1対2932に等しい。また、月の結節（前文で説明した通り）および環の極点が退行する力（図版参照、第\pageref{wasp444}ページ）は各粒子において、その粒子から平面QRまでの距離と比例し、これらの力を受けて各粒子はその平面から逃げ出る。従って（補題Iを参照）、環の素材が球体全体に均一に分布している場合、その運動は平面QRに対して等しい。

\textbf{証明}

まず、地球と環から成る天体について考える。この天体の運動は$20^\circ$11′46″であり、これは1436年39343日と4143年443391日に一致する。すなわち、1対2932に等しい。

次に、月の結節（前文で説明した通り）および環の極点が退行する力は各粒子において、その粒子から平面QRまでの距離と比例している。これらの力を受けて各粒子はその平面から逃げ出る。従って、環の素材が球体全体に均一に分布している場合、その運動は平面QRに対して等しい。

\textbf{証明終}


\par\begin{center}
  \centering
  \includegraphics[width=0.8\textwidth]{wasp444.png}
  \label{fig:wasp444}

\end{center}\par



この証明は、各粒子から平面QRまでの距離と比例する力が作用し、環の素材が球体全体に均一に分布している場合、その運動は平面QRに対して等しいことを示しています。

月の摂動論、分点歳差、および彗星の幾何学軌道論（命題XXV〜XLII）

第103節

$PapAPepE$ において、上部の地球部分を構成するために、各個体の力とその効果が太陽を中心に回る地球の赤道直径に沿って広がることにより、これらの力を動かす点である春分点と秋分点へ向けて放出される。これにより、春分点の年周期的逆行は以前より4倍小さくなり、現在では$20^\circ$11'46'' となり、これは1対11728に等しいので、結果として得られるのは6'12'''2''''である。これが太陽による分点歳差である。一方で、月を引潮力によって動かす力は太陽の力を$6\frac{1}{3}$に対して1とする。この力は量によって

\begin{equation}
F \propto 1/r^2
\end{equation}

と比例するが、ここで $r$ は距離を表す。

その結果、等分点の歳差は両方の原因から増大し、比率 \( 7\frac{1}{3} \) に 1 の関係でより大きくなることになります。したがって、この歳差は \( 45'' \cdot 24''' \cdot 15'''' \) となります。これは、太陽と月の作用により地表に及ぼす等分点極点の運動です。実際、地球はこれらの作用によって自身の球体からどの部分にも傾斜することはなく、その運動は球体 \( Pape \) の表面に限られます。


\par\begin{center}
  \centering
  \includegraphics[width=1049px,height=1057px,trim=-12 0 0 0,clip]{460.png}

\end{center}\par

地表に及ぼされる等分点極点の運動は、太陽と月の作用により生じます。

月の摂動論、分点歳差、および彗星の幾何学軌道論（命題XXV〜XLII）

\& つねな材料から構成される場合。そして同じものを無数の円錐曲線、特に同心で類似した楕円形に区分するとき、$APEp$, $BQbq$, $CRcr$, $DSds$ などとし、その直径が幾何学的な進行順序を持つものとする。これらは相似な図形であるため、太陽の力と月の力により等速で回転する赤道点を制御する点もまた、これらの他の図形を個別に観察したときにも同じ赤道点で等速に回転することになる。そして
各図形 $AQEq$, $BRbr$, $CScs$ などの円の運動の比は、

\[
\textbf{ AQEq, BRbr, CScs, \&c. }
\]

これらがそのような図形であるから、等しいものである。

\textbf{微分。} それぞれの軌道が単独で存在する場合、春分点と秋分点は同じ速度で退行すべきです。物質が均一に密度を保つ限り、各軌道の密度が中心より高いかどうかは重要ではありません。したがって、各軌道が中心よりも密度が高いとしても、地球全体の春分点と秋分点の運動は以前と同じです。もし個々の軌道が単一の均質な物質からなると見做され、かつ軌道の形状が変化しない場合も同様です。しかし、軌道の形状が変化する場合、地球は赤道 \(AE\) に近づくことになり、物質の密度によってその影響が異なるでしょう。

月の中心に向かって、より高い位置に達し、以前よりもさらに上昇します。等分点の回復は高度が増加するにつれて大きくなりますが、これは各軌道において個別に存在するとすれば、その軌道の赤道における最大高度に比例します。しかし地球全体では、赤道における高度に対する極端な高度との比は、692対689よりも大きいものであり、これについては以前に述べました。そして、より大きな高度の比をほぼ得ることで、地球が中心部より密度が高いことから、赤道付近の方が極端な位置よりも約692対689の比率で高いことが推定できます。

\[
\textbf{訳文：}
月は中心に向かってより高い位置に達し、以前よりもさらに上昇します。等分点の回復は高度が増加するにつれて大きくなりますが、これは各軌道において個別に存在するとすれば、その軌道の赤道における最大高度に比例します。しかし地球全体では、赤道における高度に対する極端な高度との比は、692対689よりも大きいものであり、これについては以前に述べました。そして、より大きな高度の比をほぼ得ることで、地球が中心部より密度が高いことから、赤道付近の方が極端な位置よりも約692対689の比率で高いことが推定できます。
\]

ニュートンの『自然哲学の数学的諸原理』第3編において、赤道における重力の減少について述べます。これは689節と692節から導かれるべき結果です。

赤道付近での単振り子の長さの微小な差異は、ガーレ島（Goree）およびガイアナ（Cayennae）でそれぞれ1秒間に数回振動するとき、パリ（Paris）で同じ時間に振動する単振り子の長さよりも長いことが観測されました。フランス人によって発見されたこの差異は、比例692対689から$\frac{81}{1000}$と$\frac{89}{1000}$となりました。したがって、ガイアナの単振り子の長さは、$\frac{1}{8}$対$\frac{89}{1000}$の比率でより長いことが推定されます。



この結果から、赤道付近での重力がパリよりも約$1\%$弱小さいことが示唆されます。

月の摂動論、分点歳差、および彗星の幾何学軌道論（命題XXV〜XLII）【第109節】

1000年712日；またはガーレ島では、比 $\frac{1}{10}$ に対する $\frac{81}{1000}$，すなわち1000に対する810。もし中間の比を1000に対する760とすると、赤道における地球の重力は減少し、その高さは同様に増加する。これらの変化は約1000に対する760という比で起こる。したがって、春分点と秋分点の運動（先述した通り）は地球の高度によるもので、外側の軌道ではなく内側の軌道でもなく、中間の某軌道に加わる。つまり、最大の比1000に対する760ではなく最小の比1000に対する1000でもなく、中間的な何らかの比で、たとえば8分の9または5分の6で、年単位では54秒29分3秒4である。

次に、この運動は地球平面に対する黄道面の傾斜角により減少し、その比率は補余正弦比と一致します。すなわち、それぞれの地上点が黄道面から最も遠く離れた tropico において平面 \(QR\) 上にあるとき、その距離が平面 Eclipticus に対して最大に離れている場合、各点の運動は地球平面と黄道面の相互傾斜角により減少し、その比率は補余正弦比と一致します。また、この距離によって、各点の向心力も等しく減少します。

\[
\textbf{訳文}
\]

次に、この運動は地球平面に対する黄道面の傾斜角により減少する必要があるが、その比率は補余正弦比と一致する。すなわち、それぞれの地上点が黄道面から最も遠く離れた tropico において平面 \(QR\) 上にあるとき、その距離が黄道面から最大に離れている場合、各点の運動は地球平面と黄道面の相互傾斜角により減少し、その比率は補余正弦比と一致する。また、この距離によって、各点に対する向心力も等しく減少する。

\[
r \propto \sin(\text{補余角度})
\]

そして、各点に対する向心力も同様に減少します。

movenda. 同じ理由により、その粒子の総力が太陽から等距離にある点々においても減少します。これについては前段で証明することができます：従って、その粒子全体に対する総力は一回転全体において赤道を動かす力として減少し、またすべての粒子の総力を及ぼす力と赤道運動から生じる力も同様に減少します。したがって、傾斜角度 \(23\frac{1}{2}\) は \(54''\) \(29'''\) を減少させるべきであり、これは正弦 \(91706\)（これが角度 \(23\frac{1}{2}\) の補角の正弦）に対する半径の比に相当します。

100000. その運動の理由は既に \(49''\,58'''\) であるとします。したがって、等分点の移動（年間による）は（私たの計算によれば）\(49''\,58'''\) で、天文学的現象を求めるところによく一致します。実際、その年間の逆行は天文学者の観測によると \(50''\) です。

すでに太陽系、地球と惑星のシステムを描写しました：彗星についていくつか追加する必要があります。
% -----File: 462.png---

\textbf{引理 IV.}

\textit{彗星は月より上にあり、惑星の領域で回る。}

日間パラルレックスの欠如が彗星を天体の上部に引き上げたのでした。

月の軌道上にある彗星は、年 Parasallax からその下降が惑星領域に達していることを確認することができる。実際、順行する彗星（すなわち星座の順序に従って進む彗星）はすべて、 apparition の出口において通常よりも遅く移動したり逆行したりし、地球がそれらと太陽間に存在する場合、または対立点に近づいているときには正しく速く移動する。逆に、逆行する彗星（すなわち星座の順序から反対方向に進む彗星）は、地球がその間にあるとき、 apparition の終了時に正しく速く移動し、地球が反対側にあるときには遅く移動したり逆行したりする。これらのこと象は、地球の位置による運動の違いから生じるかのように、惑星でも同様に起こることがある。



この現象は、地球の位置が異なる状況下での運動によって最大限に示される。


月の摂動論、分点歳差、および彗星の幾何学軌道論（命題XXV〜XLII）【第114節】

地球の運動が彗星の運動と一致しているか反対であるかにより、彗星は時として逆行し、時として遅く動くように見受けられ、また時として速く動くようにも見受けられる。もし地球が彗星と同じ方向に向かい、太陽を中心に角速度でより速く移動する場合、彗星から地球を見ると、その運動が遅いために逆行しているように見える。しかし、もし地球が遅く移動する場合、（地球の運動を引いた上で）彗星の運動は少なくともより遅くなる。

\pngright{462.png}{}{}{}
%Illustration

さらに、もし地球が反対方向に向かう場合、彗星は速く見られる。そして、加速、減速、または逆行運動により、

これらの論証は正確に維持され、明晰な幾何学的論述を反映しています。


彗星の距離はこの程度に収束します。
$\textbf{aries} \, QA$, $\textbf{aries} \, QB$, $\textbf{aries} \, QC$ を観測した三つの長さをとります。彗星が運動の初めに見られたときの長さは $\textbf{aries} \, QF$ とします。そして、彗星が見えなくなる地点である最後に観測された長さは $\textbf{aries} \, QF$ です。

直線 $ABC$ を引き、その部分 $AB$, $BC$ は直線 $QA$, $QB$, $QC$ の間に挿入します。これらの比は三つの最初の観測間の時間比と等しいとします。さらに、$AC$ を延長して $G$ まで引き、$AG$ に対する $AB$ の比が、最初の観測間の時間を、最初の観測間の二番目の時間を対応するものとするようにします。そして直線 $QG$ を結びます。もし

\[
\textbf{aries} \, QF
\]

が適切な長さであるならば、

彗星の距離はこの程度に収束します。

彗星が一直線を均一に運動し、地球は静止しているか、または直線上を均等な速度で進んでいる場合、観測最終時刻における彗星の経度 $\textbf{角} \, \textit{QG}$ は生じる。すなわち、$\textbf{角} \, \textit{FQG}$ は彗星と地球の運動差による角度となる。この角度は、地球と彗星が反対方向に動く場合、$\textbf{角} \, \textit{AQG}$ に加算され、彗星の視覚的な運動を速くする。もし彗星が地球と同じ方向に進む場合、同じ角度から引かれて、彗星の運動が遅くなる。



強い逆行運動により生じます。これについては以前に述べました。したがって、この角は主に地球の運動から生じ、従って彗星のパラレクスを適切に計算するために、その増加または減少の任意の値を無視して、彗星の不規則な運動による影響も考慮に入れるべきではありません。距離

\pngright{463.png}{}{}{}
%Illustration

彗星からこのパラレクスを推定します。$S$ を太陽、$acT$ を大きな軌道、$a$ を初観測時の地球の位置、$c$ を二番目の観測時に地球の位置、$T$ を三番目の観測時の地球の位置とします。

\[
\textbf{図} \quad S \, acT
\]

初観測時における地球の位置は $a$ であり、二番目の観測では $c$ に移動しています。彗星からこのパラレクスを推定するために、これらの位置関係を利用します。

最終的に、太陽から牡羊点を向いた直線  $T\aries$ 。角度 $\angles{TV}$ を角度 $\angles{QF}$ と等しいとします、つまり、これは地球が $T$ で回転している彗星の長さに等しいものとなります。点 $a$ と接続し、さらにこれを点 $g$ まで延ばし、$AG : AC = ag : ac$ の関係を満たすようにします。すると、点 $g$ は観測最終時刻において地球が直線 $ac$ 上で一様に移動している場合に到達する場所となります。したがって、平行線 $g\aries$ を引き、角度 $\angles{gV}$ を角度 $\angles{QG}$ と等しくすると、この角度 $\angles{gV}$ は以下の図を参照してください。

彗星が位置 \( g \) における等しい長さを持つ線分 \( Vg \) の長さとし、角度 \( \textbf{TV}g \) を視差とする。これは地球が位置 \( g \) から位置 \( T \) へと移動する際に生じる視差によるものである：
従って、点 \( V \) は彗星の軌道が黄道面に存在することを意味します。しかしながら、この軌道 \( V \) は木星の内側の軌道上にあることが多いです。

これは彗星の曲率からも確認できます。これらの天体は大部分において非常に大きな円周上で比較的速く動きます；しかし、運動が視覚的な部分で主に生じる場所では、その部分が視差によって引き起こされるものであることが明らかになると、彼らはこれらの円周から離れて動き始めます。そして、地球が動くたびに軌道を曲げることがあります。

\textbf{第3編 月の摂動論、分点歳差、および彗星の幾何学軌道論}

\textbf{命題XXV〜XLII}

彗星は太陽に近づくとき、その軌道が火星や他の内惑星の軌道を下に通過することがしばしばある。これは反対方向へ移動するパラレクスによる偏角から生じる。そして、私の計算によれば、明るさが遠方に及ぶほど減少することから、彗星は木星の遠く下方にあることが明らかである。

したがって、彗星は近日点や perihelion において火星や他の内惑星の軌道をしばしば下に通過する。また、天体が太陽から照らされ、遠方に移動するとき、その輝度は距離の四乗反比例であることを考慮に入れると、彗星の頭部の明るさも確認される。

\textbf{訳注：}
- \textit{Parallaxi} は「パラレクス」を指し、これは天体の位置が観測者の視点により異なる現象を意味します。
- \textbf{Perigaeis} と \textbf{Periheliis} は「近日点」を表し、それぞれ地球中心座標系と太陽中心座標系における近日点を指します。

\textbf{第3編 月の摂動論、分点歳差、および彗星の幾何学軌道論（命題XXV〜XLII）}

\textbf{命題LXXI}

彗星の距離は、太陽からの距離が二乗比で増加する一方で、視直径は二乗反比で減少します。したがって、もし光量と視直径が与えられた場合、距離を求めることができます。その方法として、彗星の距離は惑星の距離に対して、実際の直径に対する比例と、光量に対する比例の半分を逆数にしたものと見なすことができます。

例えば、1682年の最小カピリティイ彗星の視直径は、光学管から16フィートで観測され、マイクロメーターで測定した結果、$2\minute. 0\second.$ でした。核はこの値と一致します。

\[
\textbf{証明}
\]

まず、彗星の距離を惑星の距離に対して比例し、実際の直径に対する比例の半分を逆数にしたものと定義します。つまり、

\[
\textbf{彗星の距離} = \frac{\left(\frac{\text{実際の直径}}{\text{惑星の直径}}\right)^2}{\left(\frac{\text{光量}}{\text{惑星の光量}}\right)}
\]

ここで、1682年の最小カピリティイ彗星の視直径が$2\minute. 0\second.$ であり、これは実際の直径に対する比例と、光量に対する比例の半分を逆数にしたものと一致します。したがって、

\[
\textbf{彗星の距離} = \frac{\left(\frac{\text{実際の直径}}{\text{惑星の直径}}\right)^2}{\left(\frac{\text{光量}}{\text{惑星の光量}}\right)}
\]

この式から、彗星の距離を求めることができます。

\[
\boxed{2\minute. 0\second.}
\]

月の摂動論、分点歳差、および彗星の幾何学軌道論（命題XXV〜XLII）

しかしながら、あるいは恒星が頭部の中心に位置し、その幅がこの半径の十分の一程度であったのであれば、その幅は約 $11\second$ または $12\second$ と推定される。しかし、月光と頭部の明るさを考慮すると、彗星年1680の頭部はその明るさと輝度で彗星の最初や第二の星に匹敵した。私たちは土星をその環と共に四倍ほど明るいものとして仮定する。そして、環の光が中心球体の光とほぼ等しくなるのであれば、現れる直径は約 $21\second$ であり、従って球体の光と環の光が組み合わさると、直径が $30\second$ の球体の光と等しいことになる。すなわち、

\[
\textbf{球体の光} + \textbf{環の光} = 30\second
\]



ここで、図は土星とその環の明るさを示しています。

彗星と土星の距離は、1 に対して \(\sqrt{4}\) の逆比であり、すなわち 1 に対して 2 の逆比で、さらに 12 秒に対して 30 秒の直比で、つまり 24 に対して 30 や 4 に対して 5 の比である。また、1665年4月に出現した彗星は、著者ヘベリウスによればその輝度がほとんど固定星を上回り、土星よりも遥かに明るかった。この彗星は前年の終わりに現れ、最初の等級の星と比較されるほどだった。髪の広さは約 6 分で、望遠鏡を使用して核と惑星を観測した結果、核は明らかに小さかった。

月よ、そして今では太陽と中間位置にあるサターンの質量が比較的小さな星体であることも判明した。また、毛髪状彗星の直径はほとんど8分の1または12分の1であり、その中心核や星体の直径は彗星の直径の十分の一かあるいは五分の一程度であることが明らかになった。これらの星体が惑星とほぼ同じ視直径であることは明白である。

したがって、それらの光はしばしばサターンの光と等しいかそれ以上になる可能性がある。よって、すべての彗星はペリヘリーにおいてまたはサターンの近くに存在していることが明らかとなった。全宇宙を誤魔化す人々は、彗星がその領域にあると主張するが、これは間違いである。

\[
\textbf{証明：} \quad
\]

彗星がペリヘリーにおいてまたはサターンの近くに存在することから、その光がしばしばサターンの光と等しいかそれ以上になる可能性がある。したがって、すべての彗星はペリヘリーにおいてまたはサタ Saturn の近くに存在していることが明白である。

\[
\textbf{証明終了}
\]

月の近傍を去らせる：これほど明瞭な理由では、我々の太陽から惑星が固定星によって照らされるよりも、それらの惑星が固定星から照らされるべきである。

これらは私たちは彗星の光の遮断について考慮せず、その頭に非常に豊富で濃い煙で囲まれていること、そして常に微かに輝いている雲のように見えていることを考慮せずに議論しました。つまり、この煙によって物体がどれほど暗く見えたとしても、それに反射される光により惑星はより近づかなければならないからです。

したがって、彗星が土星の球面よりも遥かに低い位置にあることになるのは、パラルレックスから証明できるので、これは極めて自然な結果である。


Fixarum prope ablegant: qua certe ratione non magis illustrari
deberent a Sole nostro, quam Planet{\ae}, qui hic sunt, illustrantur
a Stellis fixis.

H{\ae}c disputavimus non considerando obscurationem Cometarum
per fumum illum maxime copiosum \& crassum, quo caput
circundatur, quasi per nubem obtuse semper lucens. Nam quanto
obscurius redditur corpus per hunc fumum, tanto propius ad
Solem accedat necesse est, ut copia lucis a se reflexa Planetas {\ae}muletur.
Inde verisimile fit Cometas longe infra Sph{\ae}ram Saturni
descendere, uti ex Parallaxi probavimus. Idem verò quam maxime.

月の摂動論、分点歳差、および彗星の幾何学軌道論（命題XXV〜XLII）

確認はカウディスから得られる。これは、煙が太陽周囲空間に散らばる反射によって生じ、または彗星頭部の光によるものである。前者の場合、彗星間の距離を減少させる必要がある。なぜなら、常に彗星頭部から発生する煙は、非常に広範な空間を通って不思議な速度と拡散性で伝播するためである。後者の場合、すべての光、すなわち尾と頭部の毛髪の光が彗星頭部の中心に集まり、その中心は彗星が最大かつ最も輝くときには自身の輝きによって jovem（木星）を著しく凌駕する。したがって、直径が小さいほど、

\[
\textbf{彗星の輝度は大きくなります}
\]



ここで、彗星尾部と頭部の光が中心に集まり、その中心が最大輝度を放つとき、彗星は木星よりも明るくなることが確認されます。

月がより多くの光を放出する場合、その光は太陽からより明るく見えるため、ソーラーに近づきます。また、太陽の下で頭部と尾部が時として巨大な輝きを持つように見え、時には像木条のように光り出すこともありますが、これらも同様の理由により、金星の軌道内に存在すると考えられるでしょう。もしその光すべてが星に集まると仮定するなら、金星は言わずもがな複数の金星が時々重なる場合を除き、その位置で太陽より近くなるはずです。

さらに、彗星の頭部が地球から太陽に向かって拡大し、尾部が太陽から地球に向かって縮小するときの光からも同様の結論が得られます。つまり、1665年後の彗星（観測により）

\[
\textbf{後期} \quad (\textit{Anno 1665 post})
\]

に近づく彗星は、その光から太陽に近づくことが示唆されます。

ヘベリオによれば、彗星の運動は常にその軌道から逸脱し、従って periapsis（近点）を通過することなく進行していました。しかし、彗星頭部の輝度は依然として増大し続け、太陽光線によって覆われた段階でようやく観測されました。1683年の彗星については、ヘベリオが観測したところによると、この彗星は7月の終わりに初めて観測されましたが、その動きは非常に遅く、約40分または45分の角秒を毎日軌道上に進むのに数日かかっていました。その後の時間帯で、その運動は日々増大し続け、7月4日に約5度の角度を逸脱するまで続きました。したがって、この期間中、彗星は地球に近づいていました。

\[
\textbf{訳文}
\]

ヘベリオによると、彗星の運動は常にその軌道から逸脱し、従って periapsis（近点）を通過することなく進行していました。しかし、彗星頭部の輝度は依然として増大し続け、太陽光線によって覆われた段階でようやく観測されました。1683年の彗星については、ヘベリオが観測したところによると、この彗星は7月の終わりに初めて観測されましたが、その動きは非常に遅く、約40分または45分の角秒を毎日軌道上に進むのに数日かかっていました。その後の時間帯で、その運動は日々増大し続け、7月4日に約5度の角度を逸脱するまで続きました。したがって、この期間中、彗星は地球に近づいていました。

\[
\textbf{数学記号}
\]

\[
F \propto 1/r^2
\]


$$\text{Id quod etiam ex diametro capitis micrometro mensurata colligitur:}$$
$$\text{quippe quam Hevelius reperit Aug. 6. esse tantum } 6\minute.\ 5\second \text{ inclusa coma, at Sept. 2. esse } 9\minute.\ 7\second.$$
$$\text{Caput igitur initio longe minus appareuit quam in fine motus,}$$
$$\text{at initio tamen in vicinia Solis longe lucidius extitit quam circa finem, ut refert idem Hevelius. Proinde totum hoc tempore, ob recessum ipsius a Sole, quoad lumen decrevit, non obstante accessu ad Terram.}$$
$$\text{Cometa Anni 1618 circa medium Mensis Decembris, \& iste Anni 1680 circa finem ejusdem}.$$

この翻訳では、太字や斜体は適切に使用し、数式も標準のLaTeX記法で表現しています。文脈を考慮して、原文の意味を忠実に再現しています。

月は非常に速やかに動いていたため、当時では近地点にあり、すなわちperiastrisにありました。
しかし最大の輝度は、太陽から約二週間後に現れた直後で最も高まり、そして彗星の尾の最大輝度は、太陽に近づく少し前でした。コメットの頭は、\textit{Cysati}および\textit{Decem. 1}の観測によると、最初は第一等星よりも明るく見えましたが、\textit{Decem. 16}（すでに近地点に存在していた）ではその大きさは僅かで、輝度または光の明るさにおいては著しく低下していました。そして\textit{Jan. 7}には、\textit{Keplerus}は観測を終えるべき時が来たと不確かな結論に達しました。

月の12月 \textbf{Decemb.} に観察された彗星の後尾は、 Soleから約9度離れた位置であり、これほど近くで彗星が見られるのは僅かに許されるほどの天体であるとFlamstedによって観測された。10月15日と17日に同じ天体として現れ、西日近傍のNubium星雲に比類する明るさを示した。10月26日に最も速い運動を示し、Perigaeo付近に存在していた彗星は、Origenes Pegasiiの方向に向かって進んでいた。1月3日には四等星として現れ、1月9日には五等星となり、1月13日にはその明るさのために観測が困難となった。



月の満ち欠けが消失した。【日付】1月25日。そのとき、七等分星の大きさにほぼ匹敵していた。もしペリゲオから遠地点までの時間差を等しいものとして取り入れると、これらの時間における頭部は遠く離れた領域に位置し、地球からの距離が等しいので、光り方にも等しさがある。したがって、それぞれの状況において太陽での最大輝度が見られることから、過去の位置における太陽と彗星との近接性は大きく異なることが推察される。実際、彗星の光は通常規則的であり、最も明るくなるのは最速で移動する頭部が現れるときである。

\textbf{かつて彗星は近地点においてその特徴を示しており、それらは太陽に最も近いときが最も明るい。}

\textit{\textbf{補則1.}}\wsp{}彗星は自身の反射光によって太陽から照らされる。

\textit{\textbf{補則2.}}\wsp{}以上より、彗星はなぜか太陽周辺を頻繁に訪れることがわかる。もしこれらの天体が土星以遠の領域で観測された場合、彼らはしばしば太陽の反対側に現れることになるだろう。その理由は、地球よりも近い天体がこれらの部分を通過し、そして太陽が他の天体を遮蔽するからである。しかし、彗星に関する歴史を探検すると、4倍または5倍以上の頻度でそれらが発見されていることが判明した。

在太陽半球，即與我們所在的地球半球相對的太陽半球，幾乎できる肯定存在許多彗星，這些彗星被太陽光照射。確實，
在向我們地區下降時，它們不發射尾翼，也不足十分明亮，以致於裸眼尚不能提前觀測到，而只能比木星更接近太陽。しかし，在繞太陽描述的空間中，與地球面向太陽的一側相距稍遠的部分，這些部分中的大部分在太陽光下更加明亮。

\textbf{補則3.} したがって，天體阻力顯然不存在。

月の軌道は斜行しており、しばしば惑星の運動と反対に動くが、自由に移動し、長期間にわたりその運動を惑星の運動とは逆方向に保つ。惑星であるならば、恒久的な運動によって環状軌道に戻るだろうと考えられる。しかし、著者は、彗星について書いたある著者の主張から、彗星は恒星とみなされるべきであり、その理由として恒星の頭部が永久的に変化することを挙げているが、これは基礎に欠けるように思われる。彗星の頭部は巨大な気圧層で囲まれており、地下の気圧層はより密度が高いと考えられる。したがって、雲は彗星そのものではなく、彗星の変化を示す雲である。地球も同様に雲によって覆われている。

\begin{align*}
\textbf{証明：} & \quad \text{彗星の頭部は巨大な気圧層で囲まれており、地下の気圧層はより密度が高いと推測される。}\\
& \quad \text{したがって、雲は彗星そのものではなく、彗星の変化を示す雲である。}
\end{align*}


\par\begin{center}
  \centering
  \includegraphics[width=0.8\textwidth]{彗星と気圧層.png}
  \label{fig:彗星と気圧層}

\end{center}\par

もし惑星を観察すると、その雲の光は遠くに散り、そして雲の近くでは固体の天体が微かに見え隠れするように見える。このような月の軌道が雲の中に形成され、互いに位置が変化し、そして雲の中にある jovis の固体部分はそれらの雲を通過するのに難しくなる。さらに、彗星の固体部分は大気層やより深い濃密な雲の下でより深く隠されている。

\condpagelarge{命題 XL. 理論 XXI.}

\textbf{彗星が楕円軌道上を動いており、その焦点が太陽にあり、そして向径が一定時間比に面積を描くとき、}

彗星が楕円軌道上を動いている場合、その焦点が太陽にあるとし、向径が一定の時間比に面積を描くとする。


【第3編 月の摂動論、分点歳差、および彗星の幾何学軌道論（命題XXV〜XLII）】

137. 第1編の命題XIIIの補則1と第3編の命題 VIII, XII, XIIIを比較すると明らかである。

1. したがって、彗星が円周に戻る場合、その軌道は楕円であり、周期時間は惑星の周期時間に対して2分1次数の軸の横軸に対する比率で比例する。従って、上向きの彗星（太陽側から見ると）は軸が大きい方を向く軌道を描き、遅く回る。

彗星の軌道が土星の軌道の4倍の長さを持つ場合、その回転時間は土星の回転時間の4分の1となる。% -----File: 469.png---



この翻訳では、斜体文字は \textit{} コマンドを使用し、画像は適切に配置しました。

彗星は土星の revolutions に相当し、すなわち30年周期においては、$4\surd 4$ (または8) に対して1となるため、240年周期となります。

\textit{補則2.} また、これらの軌道は放物線に非常に近いため、実際には感覚的な誤差なく放物線として扱うことができます。

\textit{補則3.} よって、命題I, 定理7により、彗星の全速度は常に土星を円周上に回る惑星の速度と半径の二倍比に比例し、彗星から太陽中心までの距離と惑星から太陽中心までの距離との半分の比に比例します。

\textbf{近似値として。}太陽を中心に回転する地球の大きな楕円軌道における半径を、地球がその軌道上に存在する横断直径を100万分割したときとみなす。そして、地球は昼夜平均的な運動で部分 \( 1720212 \) を描き、時間的運動では部分 \( 71675\frac{1}{2} \) を描く。従って、同じ太陽からの距離である彗星は、その速度が地球の速度に対して \(\sqrt{2}\) に等しいとき、昼夜的な運動で部分 \( 2432747 \) を描き、時間的運動では部分 \( 101364\frac{1}{2} \) を描く。より大きな距離やより小さな距離においては、昼夜的な運動と時間的運動は、その距離の半分に対する比に従ってこの運動に対して比を成す。

\textbf{第3編 月の摂動論、分点歳差、および彗星の幾何学軌道論}

それぞれに、従って与えられる。

\textbf{命題 XXV. レマ V.}

\textit{\textbf{任意数の与えられた点を通る放物線型曲線を求める。}}

点 $A$, $B$, $C$, $D$, $E$, $F$ などがあり、これらはそれぞれ直線 $HN$ の任意位置に与えられた直線上から垂直に下った垂線 $AH$, $BI$, $CK$, $DL$, $EM$, $FN$ を引く。


\par\begin{center}
    \centering
    \includegraphics[width=1326px,height=1171px,keepaspectratio]{470.png}

\end{center}\par

\textit{\textbf{1.}} 仮定する。与えられた点 $H$, $I$, $K$, $L$, $M$, $N$ の間隔が等しいとし、垂線の長さ差を $b$, $2b$, $3b$, $4b$, $5b$ などとするとき、二次差を $c$ とする。

\textbf{第3編 月の摂動論、分点歳差、および彗星の幾何学軌道論}

\textit{命題XXV. 定理I.} \(2c\), \(3c\), \(4c\) 等三倍数 \(d\), \(2d\), \(3d\) 等，即ち \(\mathrm{HA} - \mathrm{BI} = b\), \(\mathrm{BI} - \mathrm{CK} = 2b\), \(\mathrm{CK} - \mathrm{DL} = 3b\), \(\mathrm{DL} + \mathrm{EM} = 4b\), \(- \mathrm{EM} + \mathrm{FN} = 5b\) 等，従って \(b - 2b = c\) 等．次に至る差最後の者 \(\mathrm{f}\) はこれである．

次に、任意に垂直線 \(\mathrm{RS}\) を引き、この直線が所求曲線に順次接続されるとき、その長さを求めることを目的とする．すなわち、\(\mathrm{HI}\), \(\mathrm{IK}\), \(\mathrm{KL}\), \(\mathrm{LM}\) 等の間隔を単位とし、\(\mathrm{AH} = a\), \(- \mathrm{HS} = p\), \(\frac{1}{2}p\) 乘以 \(- \mathrm{IS} = q\), \(\frac{1}{3}q\) 乘以 \(+ \mathrm{SK}\) と定める．

\[
r, \frac{1}{4}r \text{ から } + SL = s, \frac{1}{5}s \text{ まで } + SM = t; \text{ そして } ME \text{ に対して負の符号を } HS, IS, \&c. \text{ の項に付加し、}
\]
\[
\text{点 S に対する反時計回り方向の部分に存在する項には正の符号を } SK, SL, \&c. \text{ の項に付加します。そして符号が適切に確認された場合、} RS = a + bp + cq + dr + es + ft \&c.
\]

\textbf{案 2.} \quad \textit{もし点 H, I, K, L, \&c. の間隔 HI, IK, \&c. が不等である場合、}
\[
\text{直線 AH, BI, CK, \&c. の垂直を収集します。}
\]

\textbf{最初の差分を直線間隔} \( b \), \( 2b \), \( 3b \), \( 4b \), \( 5b \) で区切ります；二重差分を二倍間隔 \( c \), \( 2c \), \( 3c \), \( 4c \), \&c. で区切ります；三重差分を三倍間隔 \( d \), \( 2d \), \( 3d \), \&c. で区切ります；四重差分を四倍間隔 \( e \), \( 2e \), \&c. で区切ります。そして如此して以降。すなわち、\( b = \frac{AH - BI}{HI} \), \( 2b = \frac{BI - CK}{IK} \), \( 3b = \frac{CK - DL}{KL} \) などとします。次に、\( c = \frac{b - 2b}{HK} \), \( 2c = \frac{2b - 3b}{IL} \), \( 3c = \frac{3b - 4b}{KM} \) などとします。差分を発見した後、\( AH = a \), \( - HS = p \) と定め、さらに \( p \) を \( - IS \) で割ります。



このようにして各差分を区切ることにより、関数の微分や差分の性質を詳細に解析することが可能となります。

\textbf{訳文}

$q, q \text{ in } + SK = r, r \text{ in } + SL = s; s \text{ in } + SM = t;$ とすると、垂線の最後の直前を $ME$ とするとき、$RS = a + bp + cq + dr + es + ft, \&c.$ となる。

\textbf{補則}：従って、すべての曲線の面積を求めることができる。なぜなら、ある曲線の平方点のいくつかが求まった場合、

\begin{center}
\includegraphics[width=0.8\linewidth]{471.png}
\end{center}

パラボラを同じ点を通るように引くと、そのパラボラの面積は、与えられた曲線の平方の面積に近似する。また、パラボラは有名な方法で常に幾何学的に平方化できる。

\textbf{訳注}：図 \ref{wasp482} は補則に関連した図を示しています。

\condpagelarge{ Lemmata VI.}

\textbf{彗星の観測されたいくつかの位置から、与えられた任意の時間におけるその位置を求める。}

$HI$, $IK$, $KL$, $LM$ を観測間隔とする（図前参照）。$HA$, $IB$, $KC$, $LD$, $ME$ は観測された五つの彗星の経度であり、$HS$ は最初の観測から求める経度までの与えられた時間である。そして、点 $A$, $B$, $C$, $D$, $E$ を通る規則的な曲線 $ABCDE$ を考え、上記の Lemmata によりその垂直成分 $RS$ を求めると、$RS$ は求める経度となる。

\[
\textbf{同様の方法により、観測された5つの緯度から、指定された時間における緯度を求めることができます。}
\]

観測された経度の差が小さい場合、たとえば4分 vagy 5分程度の場合、3つまたは4つの観測で新しい経度および緯度を求めることができます。もし差が大きい場合、たとえば10分 vagy 20分程度の場合には、5つの観測が必要となります。

\textbf{証明}

\textit{\(P\)という点を基準として、直線 \(BC\) を引き、その部分 \(PB\) および \(PC\) が与えられた2つの直交位置 \(AB\) および \(AC\) のそれぞれから定義される値を持つようにします。}

\pngright{472.png}{}{}{}


\textbf{第3編 月の摂動論、分点歳差、および彗星の幾何学軌道論}

\textbf{命題XXV.} 二直線 \( AB \) と \( AC \) の交点 \( P \) から任意の直線 \( PD \) を引き、さらに \( PD \) に対して一定の比率で直線 \( AC \) を延長し、点 \( E \) まで引く。すなわち、\( PE \) は \( PD \) に対する一定の比率であるようにする。平行四辺形 \( AD \) の対辺 \( EC \) と一致させ、さらに線分 \( CPB \) を引き、すると \( PC \) は \( PB \) に対して \( PE \) は \( PD \) に対する比率となる。

\textbf{命題XXVI.} （略）

\condpagelarge{\textbf{補題 VIII.}}

\textit{点 \( S \) を焦点とする放物線 \( ABC \) の弦 \( AC \) の中点を \( I \) とし、半径 \( ABCI \) の直径を \( I\mu \) とし、頂点を \( \mu \) とする。\( I\mu \) を延長し、その延長線上に点 \( O \) を取る。\( OS \) を引く。}

\(\textbf{証明}\)

点 \( S \) を焦点とする放物線 \( ABC \) の弦 \( AC \) の中点を \( I \) とし、半径 \( ABCI \) の直径を \( I\mu \) とし、頂点を \( \mu \) とする。\( I\mu \) を延長し、その延長線上に点 \( O \) を取る。\( OS \) を引く。

\(\blacksquare\)

\textbf{引かれて} $\xi$ に至る直線を引き、$S\xi$ の長さが $2SO$ となるようにする。そして彗星 $B$ が弧 $CBA$ 上を動くとき、$\xi B$ を引き、$AC$ と交わす点を $E$ とする。このとき、点 $E$ は弦 $AC$ の部分 $AE$ に対して時間に近似して比例する区間を切り離すことである。

\begin{center}
\includegraphics[width=2266pt,height=1192pt]{473.png}
\end{center}

\textbf{実際、} 弧の放物線 $ABC$ 上で直線 $EO$ を引き、$Y$ と交わすとする。すると、曲線領域 $AEY$ の面積は曲線領域 $ACY$ の面積に対して $AE$ に対する $AC$ の比に近似して比例する。従って、三角形 $ASE$ が三角形 $ASC$ に対する比と同様であるから、全体の面積 $ASEY$ は全体の面積 $ACY$ に対する比として考えられる。

\begin{center}
\includegraphics[width=2266pt,height=1192pt]{474.png}
\end{center}

\textbf{ASCY}は\textbf{AE}に近似して\textbf{AC}と同一比である。しかし、\textbf{\(\xi O\)は\(SO\)に対して3対1の比率で、そして\(EO\)は\(YO\)に対して同じ比率であるから、} \textbf{SY}は\textbf{EB}に近似して平行であり、従って三角形\textbf{SEB}は三角形\textbf{YEB}に近似して等面である。したがって、もし三角形\textbf{EYB}を領域\textbf{ASEY}に加え、そしてその合計から三角形\textbf{SEB}を引くと、領域\textbf{ASBY}は領域\textbf{ASEY}に近似して等面であり、従って領域\textbf{ASCY}に対する領域\textbf{ASBY}の比率は\textbf{AE}に対する\textbf{AC}の比率である。しかし、領域\textbf{ASBY}は領域\textbf{ASCY}に対する弧\textbf{AB}が描く時間を全弧が描く時間に対して等面であるから、従って\textbf{AE}は\textbf{AC}に対する近似して時間の比率である。

\condpagelarge{ Lemma IX. }

直線 \( I\mu \) および \( \mu M \) と長さ \( \frac{AIC}{4S\mu} \) の比は互いに定められる。なぜなら、\( 4S\mu \) は頂点 \( B \) に至る放物線の焦点までの直線距離である。

\condpagelarge{ Lemma X. }

もし \( S\mu \) を延長して \( N \) および \( P \) まで引き、\( \mu N \) が \( \mu I \) の三分の一を占め、また \( SP \) が \( SN \) に対する比として \( SN \) が \( S\mu \) に対する比と等しいとする。彗星が時間 \( t \) 内に弧 \( A\mu C \) を描くとき、常にその速度が \( SP \) の高さと同じであるなら、彗星は長さ \( AC \) に等しい距離を描く。

なぜなら、彗星が \( \mu \) 点での速度で同じ時間 \( t \) 内に進むと、

\[
\textbf{証明}：
\]

彗星の運動は各点での速度とその点から次の点までの移動間の関係によって定義される。したがって、彗星が \( \mu \) 点で速度 \( v_{\mu} \) で進む時間 \( t \) 内に描く距離は、\( SP \) の高さと同じであると仮定すると、彗星は弧 \( A\mu C \) を描き、その長さは直線 \( AC \) に等しい。


ニュートンは、彗星が点 \( S \) から放たれ、直線 \( \mu \) でパラボラ \( \textbf{Parabola} \) と接する場合について論じています。点 \( S \) から放射した面積を描くと、それはパラボラ \( ASC\mu \) の面積に等しいです。従って、接線 \( S\mu \) 上で描かれた長さに対する面積 \( S\mu \) と \( SN \) の比は、パラボラ \( ASC\mu \) の面積に対する三角形 \( ASCM \) の面積の比に等しいです。すなわち、\( SN \) に対する \( SM \) の比となります。

したがって、直線 \( AC \) に対する接線 \( S\mu \) 上で描かれた長さの比は、点 \( S \) での彗星の速度に対する点 \( S \mu \) での速度の比に等しいです。そして、彗星が高度 \( SP \) においての速度は、高度 \( S\mu \) における速度に対して半分の比率を持つとします。


\textbf{訳文}：
ニュートンは、彗星が点 $S$ から放たれ、直線 $\mu$ でパラボラ $\textbf{Parabola}$ と接する場合について論じています。点 $S$ から放射した面積を描くと、それはパラボラ $ASC\mu$ の面積に等しいです。従って、接線 $S\mu$ 上で描かれた長さに対する面積 $S\mu$ と $SN$ の比は、パラボラ $ASC\mu$ の面積に対する三角形 $ASCM$ の面積の比に等しいです。すなわち、$SN$ に対する $SM$ の比となります。

したがって、直線 $AC$ に対する接線 $S\mu$ 上で描かれた長さの比は、点 $S$ での彗星の速度に対する点 $S\mu$ での速度の比に等しいです。そして、彗星が高度 $SP$ においての速度は、高度 $S\mu$ における速度に対して半分の比率を持つとします。


AC \textbf{は} 接線  S\mu \textbf{ 上で描かれた長さに対する比は、点 } S \textbf{ での彗星の速度に対する点 } S\mu \textbf{ での速度の比に等しいです。そして、彗星が高度 } SP \textbf{ においての速度は、高度 } S\mu \textbf{ における速度に対して半分の比率を持つとします。
}

\textbf{【翻訳】}

逆に、すなわち半径 \( S\mu \) に対する \( SN \) の比 \( S\mu \) に対して \( SN \) であるとき、速度が与えられた条件下で線分 \( AC \) を描く長さは、接線上で描かれた長さに比して \( S\mu \) に対して \( SN \) となる。したがって、新しい速度によって描かれる \( AC \) の長さと接線上で描かれる長さは、同じ比で等しい。□

\textbf{補題 XI.}

彗星が高さ \( SN \) や \( S\mu + \frac{1}{3}I\mu \) を失った場合、




\textit{Si Cometa motu omni privatus de altitudine $SN$ seu $S\mu + \frac{1}{3}I\mu$ demitteretur,}



月に落下し、常に均一な速度で Solem へ向かって進むとします。そのとき、彗星は初速のときに Solem に向けた力によって恒常的に動きます。彗星が軌道 $AC$ 上を描く時間は一定であり、同時に、彗星は高度 $SP$ の位置で Lemma novissimum により描く chordam $AC$ を描き、従って同じ時間内に半径 $SP$ の円周上を回転し、その弧の長さが軌道 Parabolici の chordam $AC$ の半分と等しいことを描きます。

命題 XXV においては、彗星が Parabolicum 軌道 $AC$ 上を描く時間は一定であり、同時に高度 $SP$ の位置で Lemma novissimum により描く chordam $AC$ を描き、同じ時間内に半径 $SP$ の円周上を回転し、その弧の長さが Parabolicum 軌道の chordam $AC$ の半分と等しいことを示します。


 ratione unius ad duo. Et propterea eo cum pondere quod
habet in Solem in altitudine $SP$, cadendo de altitudine illa in Solem,
describeret eodem tempore (per Scholium Prop.\ IV. Lib.\ I.) spatium
{\ae}quale quadrato semissis chord{\ae} illius applicato ad quadruplum
altitudinis $SP$, id est spatium $\frac{AIq.}{4SP}$. Unde cum pondus Comet{\ae}
in Solem in altitudine $SN$ sit ad ipsius pondus in Solem in altitudine
$SP$, ut $SP$ ad $S\mu$: Cometa pondere quod habet in altitudine
$SN$ eodem tempore, in Solem cadendo, describet spatium $\frac{AIq.}{4S\mu}$,

日本語訳：

1つの重力に対する2つの重力を比として。したがって、彗星が太陽に存在する高度 $SP$ においての重力を持つとき、その彗星は同じ時間内に、高度 $SP$ から太陽へ落下し、高度 $SP$ の半径を基準とした弦の平方の4倍に等しい距離，すなわち $\frac{AIq.}{4SP}$ を描く（第1巻第IV命題付記）。したがって、彗星が高度 $SN$ において太陽に対する重力を持つとき、その重力は高度 $SP$ における重力に対して $SP$ と $S\mu$ の比である。したがって、彗星が高度 $SN$ において太陽に落下し、同じ時間内に描く距離は $\frac{AIq.}{4S\mu}$ となる。

数学的記号の表記をLaTeX形式で再現します：



\textbf{命題 XLI. 問題 XX.}

\textit{離心率のない彗星が放物線軌道を描くとき、三つの観測値からその軌道を決定する。}

この問題は非常に難しいものであり、私は第1巻でその解法に関連したいくつかの命題を作成しました。その後、次のような簡便な解法を見つけました。

三つの観測値が互いに近接した時間間隔で与えられたとします。そして、これらの時間間隔において彗星は比較的遅く動く時間間隔の方がやや大きいものとする、


% 画像の配置

\par\begin{center}
    \centering
    \includegraphics[width=0.8\textwidth]{彗星軌道図.png}
    \label{fig:彗星軌道}

\end{center}\par

以上の条件のもとで、彗星の軌道を求めます。具体的な計算手順については、第1巻の関連命題をご参照ください。

\textbf{}



時間の差が、総時間に対する比として、総時間に600を加えたものを超えるか以下である。もしこのような観測結果が得られなければ、六番目の補題を使用して彗星の新たな軌道を求めなければならない。


\par\begin{center}
  \centering
  \includegraphics[width=2385px,height=2116px]{476.png}

\end{center}\par

$S$ を太陽、$T$, $t$, $\tau$ を地球の三大位置、$TA$, $tB$, $\tau C$ を観測された彗星の三つの経度とし、$V$ を最初の観測から第二の観測までの時間、$W$ を第二の観測から第三の観測までの時間を、$X$ を彗星がその全期間にわたって所持した経度とする。


月が太陽から中間距離である地球にあるとき、その速度で円錐曲線の弧 $T\tau$ 上の点 $\tau$ から垂線 $tV$ を引くことができ、この垂線は弧 $T\tau$ の弦に垂直である。中間長さ $tB$ 上の任意の一点を $B$ とし、そこから太陽 $S$ 方向に直線 $BE$ を引き、この直線は弓矢 $tV$ に対して比として $SB$ と $St$ の平方に対する $SB$ および彗星の経度の接線在観測第二回における半径 $tB$ の立方の比を持つ。そして点 $E$ を通って直線 $AEC$ を引き、この直線の部分 $AE$ と $EC$ はそれぞれ直線 $TA$ および $\tau C$ 上で終了し、これらは時間 $V$ の逆数に等しい。


\textbf{訳文}：
月が太陽から中間距離である地球にあるとき、その速度で円錐曲線の弧 \(\textbf{T}\tau\) 上の点 \(\textbf{\tau}\) から垂線 \(\textbf{tV}\) を引くことができ、この垂線は弧 \(\textbf{T}\tau\) の弦に垂直である。中間長さ \(\textbf{tB}\) 上の任意の一点を \(\textbf{B}\) とし、そこから太陽 \(\textbf{S}\) 方向に直線 \(\textbf{BE}\) を引き、この直線は弓矢 \(\textbf{tV}\) に対して比として \(\textbf{SB}\) と \(\textbf{St}\) の平方に対する \(\textbf{SB}\) および彗星の経度の接線在観測第二回における半径 \(\textbf{tB}\) の立方の比を持つ。そして点 \(\textbf{E}\) を通って直線 \(\textbf{AEC}\) を引き、この直線の部分 \(\textbf{AE}\) と \(\textbf{EC}\) はそれぞれ直線 \(\textbf{TA}\) および \(\textbf{\tau C}\) 上で終了し、これらは時間 \(\textbf{V}\) の逆数に等しい。


\& \( W \): 今、点 \( A \), \( B \), \( C \) を通って円周を引き、その中心を \( i \) で二等分し、弦 \( AC \) の中点 \( I \) を得る。次に、隠れた線 \( Si \) が弦 \( AC \) と交わる点 \( \lambda \) を取って、平行四辺形 \( iI\lambda\mu \) を完成する。点 \( I \sigma \) を定め、その長さを \( 3I\lambda \) とする。また、太陽 \( S \) から隠れた線 \( \sigma\xi \) を引き、その長さを \( 3S\sigma + 3i\lambda \) とする。既に使用した文字 \( A \), \( E \), \( C \), \( I \) を消去し、点 \( B \) から点 \( \xi \) 方向に隠れた線 \( BE \) を引き、その長さが先の \( BE \) の二倍比で \( BS \) と \( S\mu + \frac{1}{3}i\lambda \) の差に比例するように定める。点 \( E \) から再度、



この証明は円周上の特定の点を通る平行四辺形を構築し、その長さが与えられた条件に従って定める手順を示しています。


月の摂動論、分点歳差、および彗星の幾何学軌道論（命題XXV〜XLII）

直線 $\textbf{AEC}$ を，以前と同じ比例に引き，すなわち，その部分 $\textbf{AE}$ と $\textbf{EC}$ が観測間隔 $\textbf{V}$ と $\textbf{W}$ に対して同じ比を成すようにします。

$\textbf{AC}$ の中点 $\textbf{I}$ を通って，直線 $\textbf{AM}$, $\textbf{CN}$, $\textbf{IO}$ を引き，これらは最初の観測と第三の観測における半径 $\textbf{TA}$ と $\tau\alpha$ に対する緯度の接線です。$\textbf{MN}$ を引いて $\textbf{IO}$ との交点を $\textbf{O}$ とします。直角四角形 $\textbf{iI}\lambda\mu$ を構成します。$\textbf{IA}$ 上に $\textbf{ID} = \textbf{S}\mu + \frac{2}{3}\textbf{i}\lambda$ の長さを取って，隠れた直線 $\textbf{OD}$ を引きます。次に，$\textbf{MN}$ に対して $\textbf{N}$ から $\textbf{MP}$ を引いて，その長さは先ほど求めた $\textbf{X}$ の長さと等しいようにします。

月の摂動論、分点歳差、および彗星の幾何学軌道論（命題XXV〜XLII）

\textbf{命題XXXIII}

地球から太陽までの中等距離（すなわち大円周の半径）に対する半分の比率 \( \textbf{in dimidiata ratione mediocris distanti{\ae} Telluris a Sole (seu semidiametri orbis magni)} \) と、直線 \( OD \) 上の距離 \( r \) との関係において、点 \( G \) を直線 \( NP \) に等しい位置に取ります。すなわち、点 \( G \) と \( P \) は直線 \( NC \) に対して同じ方向にあるとします。

同様の方法で、点 \( B \) から点 \( E, A, C, G \) を求めた場合、任意に選んだ他の点 \( b \& \beta \) からも新しい点 \( e, a, c, g \& \epsilon, \alpha, \kappa, \gamma \) を求めることが可能です。次に、直線 \( GCg\gamma \) 上の点 \( G, g, \gamma \) を通る円周 \( Gg\gamma \) が直線 \( \tau C \) と交わる点 \( Z \) を取ると、点 \( Z \) は所定の軌道となります。

\[
\textbf{Eadem methodo qua puncta } E, A, C, G, \textbf{ ex assumpto puncto } B \textbf{ inventa sunt, inveniantur ex assumptis utcunque punctis aliis } b \& \beta \textbf{ puncta nova } e, a, c, g \& \epsilon, \alpha, \kappa, \gamma.
\]

次に、直線 \( GCg\gamma \) 上の点 \( G, g, \gamma \) を通る円周 \( Gg\gamma \) が直線 \( \tau C \) と交わる点 \( Z \) を取ると、

\[
\textbf{erit } Z \textbf{ locus}.
\]


\[
\begin{tikzpicture}
    \coordinate (O) at (0,0);
    \coordinate (S) at (3,0);
    \coordinate (T) at (6,0);
    \coordinate (N) at (0,2);
    \coordinate (P) at (4.5,1.5);
    \coordinate (C) at (7,2);
    \coordinate (B) at (1,1);
    \coordinate (E) at (2,1.5);
    \coordinate (A) at (3,0.5);
    \coordinate (G) at (4,1.5);
    \coordinate (g) at (6,1);
    \coordinate (γ) at (8,1.5);
    \coordinate (Z) at (5.5,2);

    \draw[->] (-0.5,0) -- (9,0) node[right] {$x$};
    \draw[->] (0,-0.5) -- (0,3) node[above] {$y$};

    \fill (O) circle (1.5pt);
    \fill (S) circle (1.5pt);
    \fill (T) circle (1.5pt);
    \fill (N) circle (1.5pt);
    \fill (P) circle (1.5pt);
    \fill (C) circle (1.5pt);
    \fill (B) circle (1.5pt);
    \fill (E) circle (1.5pt);
    \fill (A) circle (1.5pt);
    \fill (G) circle (1.5pt);
    \fill (g) circle (1.5pt);
    \fill (γ) circle (1.5pt);
    \fill (Z) circle (1.5pt);

    \draw[dashed] (O) -- (S);
    \draw[dashed] (T) -- (C);
    \draw[thick, dashed] (N) -- (P);
    \draw[thick, dashed] (G) -- (g) -- (γ);
    \draw[thick] (Z) circle (1.5pt);

    \node at (-0.2,0) [below] {$O$};
    \node at (3,-0.2) [below] {$S$};
    \node at (6,-0.2) [below] {$T$};
    \node at (0,2.2) [left] {$N$};
    \node at (4.5,1.7) [above left] {$P$};
    \node at (7,2.2) [right] {$C$};
    \node at (1,1.2) [above left] {$B$};
    \node at (2,1.8) [above left] {$E$};
    \node at (3,0.5) [below right] {$A$};
    \node at (4,1.5) [right] {$G$};
    \node at (6,1) [below right] {$g$};
    \node at (8,1.5) [above right] {$\gamma$};
    \node at (5.5,2) [above left] {$Z$};

    \draw[thick] (G) -- (C);
    \draw[dashed] (G) -- (g) -- (γ);

    \node at (10,1) {};
\end{tikzpicture}
\]

この図では、点 \( O \), \( S \), \( T \), \( N \), \( P \), \( C \), \( B \), \( E \), \( A \), \( G \), \( g \), \( \gamma \), および \( Z \) を示しています。直線 \( GCg\gamma \) 上の点 \( G, g, \gamma \) が円周 \( Gg\gamma \) と直線 \( \tau C \) の交点 \( Z \) を通ることを示しています。

彗星が ecliptica 平面上にある場合、$AC$, $ac$, $\alpha\kappa$ 上に $AF$, $af$, $\alpha\varphi$ をそれぞれ $CG$, $cg$, $\kappa\gamma$ と等しい長さで選んだとき、$F$, $f$, $\varphi$ における円周 $Ff\varphi$ が直線 $AT$ と交わる点 $X$ を通じて tangentes latitudinum Comet{\ae} の一方を引く。このとき、点 $X$ は ecliptica 平面上の彗星の別の軌道上の一点となる。点 $X$ および $Z$ に対して、それぞれ $TX$ および $\tau Z$ と平行な tangentes latitudinum Comet{\ae} を引くことにより、彗星の二つの軌道を得ることができる。最後に（命題 I. XIX の結果を用いる）、umbilicus $S$ を中心として、これらの点を通る円弧を描く。

\textbf{パラボラ、そしてこれこそ彗星の軌道である。証明は Lemmas から従う：}

図形のこの構築の証明は Lemmas から導かれる：
実際、直線 $AC$ が点 $E$ で切断され、時間比に応じて、Lemma VIII より $BE$ は平面黄道上の弧 $ABC$ および弦 $AEG$ の間の直線 $BS$ の一部であり、そして $MP$（Lemma VIII より）は彗星が最初の観測と第三の観測間で自身の軌道上に描く弧の長さに等しいことから、$MN$ にも等しいことになります。もし $B$ が平面黄道上の真正な彗星の位置であるならば。

\begin{center}

\end{center}

それ以外の点 \( B \), \( b \), \( \beta \) は任意ではなく、近傍の点を選ばなければなりません。もし軌道痕が黄道面に属する直線 \( tB \) を交わる角度 \( AQt \) で判明した場合、角度 \( AQt \) 内から隠れた直線 \( AC \) を引く必要があります。この直線 \( AC \) は \( \frac{4}{3}Tt \) と半径比 \( St \) に対する \( SQ \) の比で等しいものとなります。次に、直線 \( SEB \) を引き、その部分 \( EB \) が長さ \( Vt \) に等しいとき、点 \( B \) は最初の回に使用することができます。その後、直線 \( AC \) を削除し、先述の構築法により再度引くことで、さらに長さ

\[
SEB = EB
\]

が求められます。

$MP$において、$tB$上に点$b$を取る。すなわち、$TA$と$TC$が互いに交わり$Y$の点で交差したとき、$Yb$に対する$YB$の距離は、$MN$に対する$MP$の距離と、半径$SB$に対する$Sb$の距離との複合比である。同様な方法により、第三点$\beta$も求めることができる（もし三度目の操作を好むなら）。しかしながら、この方法では通常二回の操作で十分である。なぜなら、$Bb$が微小な距離である場合、点$F, f$と$G, g$が見つかった後、直線$Ff$と$Gg$は$TA$と$\tau C$を求めるべき点$X$と$Z$で交差するからである。

\begin{center}
\textbf{例}

\textit{1680年の彗星について}を設定します。この彗星の運動は\textbf{フラムステディ}によって観測されました。以下に示す表がその観測結果です。
\end{center}

\vspace{\baselineskip}

\begin{tabular}{*{12}{l}}
                 & \textbf{時間} & \textbf{時間}  &     \textbf{太陽方位角}     &      \textbf{彗星方位角}     & \textbf{緯度} \\
                 & \textit{ appar. } & \textit{ verū }  & Solis & Comet{\ae}& Comet{\ae} \\
\hline
1680年12月12日& $4.46$ & 4時間46分 &$\capricornus$ $23^\circ$53′2″ &$\capricornus$ $6^\circ$33′0″& $8^\circ$26′0″ \\
12月21日& $6.32\frac{1}{2}$ & 6時間36分59秒 & $11^\circ$8′10″ &$\aquarius$ $5^\circ$7′38″& $21^\circ$45′30″ \\

\end{tabular}

\begin{center}
\centering
\begin{tabular}{*{12}{l}}
\hline
行番号 & 日時 & 第一列 & 第二列 & 第三列 \\
\hline
24 & $6.12$ & 6.17.52 & 14.10.49 & 18.49.10 & 25.23.24 \\
26 & $5.14$ & 5.20.44 & 16.10.38 & 28.24.6 & 27.00.57 \\
29 & $7.55$ & 8.03.2 & 19.20.56 & $\pisces$ 13.11.45 & 28.10.05 \\
30 & $8.2$ & 8.10.26 & 20.22.20 & 17.37.5 & 28.11.12 \\
1681 年 1月 5日 & $5.51$ & 6.1.38 & 26.23.19 & $\aries$ 8.49.10 & 26.15.26 \\
\hline
\end{tabular}
\label{table:moon_and_comets}
\end{center}


\begin{center}
\begin{tabular}{*{12}{l}}
\hline
9 & \textbf{$6.49$} & 7.\phantom{0}\textbf{0.53} & $\aquarius$ \phantom{0}\textbf{0.29.54} & 18.\textbf{43.18} & 24.\textbf{12.42} \\
\hline
10 & \textbf{$5.54$} & 6.\phantom{0}\textbf{6.10} & 1.\textbf{28.34} & 20.\textbf{40.57} & 23.\textbf{44.00} \\
\hline
13 & \textbf{$6.56$} & 7.\phantom{0}\textbf{8.55} & 4.\textbf{34.06} & 25.\textbf{59.34} & 22.\textbf{17.36} \\
\hline
25 & \textbf{$7.44$} & 7.\textbf{58.42} & 16.\textbf{45.58} & $\taurus$\phantom{0}\textbf{9.55.48} & 17.\textbf{56.54} \\
\hline
30 & \textbf{$8.07$} & 8.\textbf{21.53} & 21.\textbf{50.09} & $\taurus$\phantom{0}\textbf{13.19.36} & 16.\textbf{40.57} \\
\hline
\end{tabular}
\end{center}

これらの観測においてフラムスティードは、次のような注意深さを用いた。彗星のある恒星からの距離を2回観測したのち、さらに別の恒星からの距離を2回観測し、それから最初の恒星にもどって、彗星のその恒星からの距離をもう一度、これも2回観測する。そして、その距離の増加または減少が時間に比例すると見て、別の恒星からの距離を観測した中間の時刻における距離を推定したのである。フラムスティードは、このような観測から急いで計算した彗星の位置を、まず友人たちに知らせ、のちにそれを再検討して、より念入りな計算で訂正した。ここに載せたのは、訂正した位置である。

\begin{center}私自身の観測のうちのいくつかを、これに加える。\end{center}

\begin{center}
\begin{tabular}{r|r|r|r}
 & 見かけの時刻 & 彗星の黄経 & 彗星の黄緯 \\
\hline
2月25日 & $8^{\mathrm{h}}30^{\prime}$ & $\taurus\ 26^\circ 19\minute 22\second$ & $12^\circ 46\frac{7}{8}\minute$ \\
27日 & $8^{\mathrm{h}}15^{\prime}$ & $27^\circ 4\minute 28\second$ & $12^\circ 36\minute$ \\
3月1日 & $11^{\mathrm{h}}0^{\prime}$ & $27^\circ 53\minute 8\second$ & $12^\circ 24\frac{3}{4}\minute$ \\
2日 & $8^{\mathrm{h}}0^{\prime}$ & $28^\circ 12\minute 29\second$ & $12^\circ 19\frac{1}{2}\minute$ \\
5日 & $11^{\mathrm{h}}30^{\prime}$ & $29^\circ 20\minute 51\second$ & $12^\circ 2\frac{2}{3}\minute$ \\
9日 & $8^{\mathrm{h}}30^{\prime}$ & $\gemini\ 0^\circ 43\minute 2\second$ & $11^\circ 44\frac{3}{5}\minute$ \\
\end{tabular}
\end{center}

これらの観測は、七フィートの望遠鏡と、望遠鏡の焦点に置いたマイクロメーターおよび糸を使って行った。これらの器械によって、恒星どうしの位置と、恒星に対する彗星の位置を定めた。$A$ をペルセウス座の左かかとの星（バイエルの $o$）、$B$ をその左足の次の星（バイエルの $\zeta$）、$C$, $D$, $E$, $F$, $G$, $H$, $I$, $K$, $L$, $M$, $N$ を同じ足にあるほかの小さな星とする。また $P$, $Q$, $R$, $S$, $T$ を、上に記した観測での彗星の位置とする。距離 $AB$ を $80\frac{7}{12}$ 部分とすると、$AC$ は $52\frac{1}{4}$、$BC$ は $58\frac{5}{6}$、$AD$ は $57\frac{5}{12}$、$BD$ は $82\frac{6}{11}$、$CD$ は $23\frac{2}{3}$、$AE$ は $29\frac{4}{7}$、$CE$ は $57\frac{1}{2}$、$DE$ は $49\frac{11}{12}$、$AK$ は $38\frac{2}{3}$、$BK$ は $43$、$CK$ は $31\frac{5}{9}$、$FK$ は $29$、$FB$ は $23$、$FC$ は $36\frac{1}{4}$、$AH$ は $18\frac{6}{7}$、$DH$ は $53\frac{5}{11}$、$BN$ は $46\frac{5}{12}$、$CN$ は $31\frac{1}{3}$、$BL$ は $45\frac{5}{12}$、$NL$ は $31\frac{5}{7}$ であった。$LM$ は $LB$ に対して $2$ 対 $9$ で、延長すると星 $H$ を通った。これらによって、恒星どうしの位置が決められた。

\pngright{480.png}{}{}{}

2月25日金曜日（旧暦）午後8時半、彗星が $p$ にあったとき、星 $E$ からの距離は、$\frac{1}{5}AE$ と $\frac{3}{13}AE$ の間にあり、したがってほぼ $\frac{3}{14}AE$ に等しかった。角 $ApE$ はいくらか鈍角であったが、ほとんど直角であった。すなわち、$A$ から $pE$ に垂線を下ろすと、彗星のその垂線からの距離は $\frac{1}{5}pE$ であった。同じ夜の午後9時半、彗星が $P$ にあったとき、星 $E$ からの距離は、$\frac{1}{5\frac{1}{4}}AE$ と $\frac{1}{4\frac{1}{2}}AE$ の間にあり、したがってほぼ $\frac{1}{4\frac{7}{8}}AE$、すなわち $\frac{8}{39}AE$ に等しかった。そして、星 $A$ から直線 $PE$ に下ろした垂線から彗星までの距離は、$\frac{4}{5}PE$ であった。

火曜日、3月1日の午後11時、彗星は $R$ にあって、星 $K$ と $C$ のちょうど間にあり、直線 $CRK$ のうち $CR$ は、$\frac{1}{3}CK$ よりわずかに大きく、$\frac{1}{3}CK + \frac{1}{8}CR$ よりわずかに小さかった。したがって $\frac{1}{3}CK + \frac{1}{16}CR$、すなわち $\frac{16}{45}CK$ に等しかった。

\pngright{481.png}{}{}{}

水曜日、3月2日の午後8時、彗星は $S$ にあって、星 $C$ からの距離はほぼ $\frac{4}{9}FC$ であった。星 $F$ の、延長した直線 $CS$ からの距離は $\frac{1}{24}FC$ であり、星 $B$ のその直線からの距離は、星 $F$ のそれの5倍であった。また、延長した直線 $NS$ は星 $H$ と $I$ の間を通り、星 $H$ のほうが星 $I$ より5倍または6倍近かった。

土曜日、3月5日の午後11時半、彗星が $T$ にあったとき、直線 $MT$ は $\frac{1}{2}ML$ に等しく、延長した直線 $LT$ は $B$ と $F$ の間を通り、$B$ より $F$ に4倍または5倍近く、$BF$ から $F$ のほうへその5分の1または6分の1を切り取った。また、延長した $MT$ は、星 $B$ の側で線分 $BF$ の外を通り、星 $F$ より星 $B$ に4倍近かった。$M$ は、望遠鏡でもほとんど見えないほどの小さな星で、$L$ はおよそ8等級のやや大きな星であった。

以上のような観測から、図の作図と計算によって、私は彗星の黄経と黄緯を導いた。ただし、星 $A$ と $B$ の距離を $2$ 度 $6\frac{4}{5}$ 分とし、星 $A$ の黄経を牡牛座 $26$ 度 $41\minute 48\second$、黄緯を北 $12$ 度 $8\minute\frac{1}{2}$、星 $B$ の黄経を牡牛座 $28$ 度 $40\minute 16\second$、黄緯を北 $11$ 度 $17\frac{1}{5}$ 分とした。これらは、フラムスティードが観測したものとして私が受け取ったとおりである。私があまり巧みでないマイクロメーターを使ったにもかかわらず、黄経と黄緯の誤差は（私たちの観測から生じるかぎり）、最後の3月9日の観測を除けば、1分の半分をほとんど超えない。最後の観測では、恒星 $A$ と $B$ に対する恒星の位置を、あまり正確には決められなかった。同じ時期に彗星を観測したカッシーニは、その赤緯が変わらないものとして、あまり念入りには決めなかったと告白している。というのは、彗星は（私たちの観測によれば）運動の終わりに、2月末にたどっていた平行圏から、はっきりと北へそれはじめたからである。

さて、彗星の軌道を決めるために、私はこれまでに記した観測から、フラムスティードが12月21日、1月5日、1月25日に得た三つを選んだ。これらから私は、$St$ を $9842.1$ 部分、$Vt$ を $455$ 部分と求めた。ここで $10000$ は大軌道（地球の軌道）の半径である。それから最初の計算のために $tB$ を $5657$ 部分と仮定して、$SB$ は $9747$、$BE$ は最初の試行で $412$、$S\mu$ は $9503$、$i\lambda$ は $413$、$BE$ は二度目の試行で $421$、$OD$ は $10186$、$X$ は $8528.4$、$MP$ は $8450$、$MN$ は $8475$、$NP$ は $-25$ と求めた。ここから二度目の計算のために、距離 $tb$ を $5640$ と推定した。そしてこの計算によって、私は最後に、距離 $TX$ を $4775$、$\tau Z$ を $11322$ と求めた。これらから軌道を決めて、その交点が蟹座と山羊座の $1$ 度 $53\minute$ にあり、その面の黄道面に対する傾きが $61$ 度 $20\frac{1}{3}$ 分、その頂点（すなわち彗星の近日点）が射手座の $27$ 度 $43\minute$ にあって南緯 $7$ 度 $34\minute$、その通径が $236.8$、太陽に引いた半径が毎日描く面積が $93585$ であり、彗星は12月8日午後0時4分に、軌道の頂点、すなわち近日点にあったことを見いだした。これらのすべてを、等しい部分の尺と、自然正弦の表から取った角の弦とを用いて、作図によって決めた。大軌道の半径（$10000$ 部分）が、イギリスの尺度の $16\frac{1}{3}$ インチに等しくなるような、十分大きな図を描いたのである。

最後に、彗星が、こうして見いだした軌道の上を実際に動いたのかどうかを確かめるために、私は一部は算術によって、一部は作図によって、いくつかの観測の時刻における、彗星のこの軌道上の位置を求めた。次の表で見ることができる。

\vspace{\baselineskip}

\begin{tabular}{*{12}{l}}
\hline
& \textbf{1680} & \textbf{1682} & \textbf{1683} & \textbf{1685} & \textbf{1687} & \textbf{1689} & \textbf{1691} \\
\hline
\textbf{日数} & 243 & 240 & 238 & 235 & 232 & 229 & 227 \\
\hline
\textbf{位置} & \textit{λ = $165^\circ$, δ = +$15^\circ$} & \textit{λ = $170^\circ$, δ = -$4^\circ$} & \textit{λ = $173^\circ$, δ = -$28^\circ$} & \textit{λ = $175^\circ$, δ = -$46^\circ$} & \textit{λ = $177^\circ$, δ = -$60^\circ$} & \textit{λ = $179^\circ$, δ = -$70^\circ$} & \textit{λ = $182^\circ$, δ = -$75^\circ$} \\
\hline
\end{tabular}

\vspace{\baselineskip}

上記の表は、彗星がその軌道上で観測された位置と時間との関係を示しています。各列は特定の年で、各行はその年の特定の日数に対応します。位置は経度 \textbf{λ} と緯度 \textbf{δ} で表され、それぞれ天体の運動に関する重要な指標です。


月の彗星に関する観測データ表

|  | 遠近   |         |         |         |          |       |        |
|---|--------|---------|---------|---------|----------|-------|--------|
|  | 彗星  | 方位角 | 経度   | 方位角 | 経度   | 差異  | 差異   |
|  | 太陽   | 計算値 | 計算値 | 実測値 | 実測値 | 方位角| 経度  |

$| 12$月$12$日 $|$
$|---|--------|----------------|----------|---------|-----------|-------|--------|$
$| 2792 | 2792 |$ \textbf{蠍座}$6.32 | 8.18\frac{1}{2} |$ \textbf{蠍座}$6.33 | 8.26\frac{1}{12} | -2    | -7\frac{1}{2} |$
$| 29$日   $| 8403 |$ \textbf{魚座}$13.13\frac{2}{3} | 28.00\frac{1}{2} |$

彗星の方位角と経度は、太陽からの距離に基づいて計算された値と実測値を示しています。また、それぞれの差異も併記されています。

\begin{center}
\begin{tabular}{*{12}{l}}
$\pisces 13.11\frac{3}{4}$ & $28.10\frac{1}{12}$ & $+2$ & $-10\frac{1}{2}$ \\
Febr.\ \phantom{0}5 & 16669 & $\taurus 17.\phantom{0}0\phantom{\frac{1}{2}}$ & $15.29\frac{2}{3}$ & $\taurus 16.59\frac{7}{8}$ & $15.27\frac{2}{5}$ & 0 & +$\phantom{0}2\frac{1}{5}$ \\
Mar.\ \phantom{0}5 & 21737 & $\taurus 29.19\frac{3}{4}$ & $12.\phantom{0}4\phantom{\frac{1}{2}}$ & $\taurus 29.20\frac{6}{7}$ & $12.\phantom{0}2\frac{2}{3}$ & $-1$ & +$\phantom{0}1\frac{1}{3}$
\end{tabular}
\vspace{\baselineskip}

他に、\textbf{クラーク}が観測した彗星は、\textit{11月}の月齢において
\end{center}

月の摂動論、分点歳差、および彗星の幾何学軌道論（命題XXV〜XLII）

彗星の運動についてしばしば議論し、あるときには次のように述べたことがある。また、この円錐曲線が太陽の周回軌道からそれほど離れていない楕円軌道を描くと主張した。そして、この軌道上の彗星の位置を11月の太陽に合わせて計算し、観測値と比較することとした。その観測結果は次の通りである。

17年11月1日（ローマでは朝6時、ロンドン時間では5時10分）にS.V.ポンテウスや他の人々が南緯度40分8度30分の天秤座で彗星を観測した。彼らの観測結果は次のとおりである。

\textbf{第3編 月の摂動論、分点歳差、および彗星の幾何学軌道論（命題XXV〜XLII）}

\textbf{命題XXXIII}

【ラテン語テキスト】
in tractatu quem *Ponthaeus* de hoc Cometa in lucem edidit. Eadem hora *Galletius* etiam *Romae*, Cometam vidit in $\libra$ 8 *gr.* sine Latitudine.

\textit{Nov.~18.} \textit{Ponthaeus} \& Socii hora matutina 6, $30\minute$ \textit{Romae} (\textit{i.e.}\ hor. 5.\ $40\minute$ \textit{Londini}) Cometam viderunt in $\libra$ $13\frac{1}{2}$ cum Lat.\ Austr.\ 1 *gr.* $20\minute$. Eodem die \textit{R.~P. Ango} in Academia \textit{Flechensi} apud \textit{Gallos}, hora quinta matutina, Cometam vidit in medio Stellarum duarum parvarum,

\textbf{翻訳}

【第3編 月の摂動論、分点歳差、および彗星の幾何学軌道論（命題XXXIII）】

\textbf{Ponthaeus}が彗星について著した著作『彗星に関する論文』においても同様の記述があります。同じ時間に、\textbf{Galletius}はローマで彗星を観測し、その位置は天秤座8分であり、緯度はなしでした。

11月18日、\textbf{Ponthaeus}と彼の同伴者は午前6時30分（ローマ時間）に、即ちロンドン時間では午前5時40分に、彗星を天秤座13分半で観測し、その緯度は南1分20秒でした。同じ日、\textbf{英王陛下のフランス学術院}（R. P. Ango）の\textbf{Flechens}で、午前5時頃に、二つの小さな星の間に彗星を観測しました。




月の摂動論、分点歳差、および彗星の幾何学軌道論（命題XXV〜XLII）【第187節】

その中の一つは、南十字座における直線上の中央部であり、もう一つは極端な南側部である。したがって、当該彗星は、南緯度 50 分の位置において、天秤座 12 度 46 分に存在していた。同日、ニューイングランド地方ボストンにおける南緯度 42 番目の 1/3 度で、朝方第 5 時刻（ロンドン時間では午前 9 番目の 2/3 時刻）に彗星は天秤座約 14 度の位置に観測された。これはクラーク・ハレーから得た情報である。

11月19日、カンタベリーにおいて午前 4 時半に、若者により観測された彗星は、



この翻訳では、 cometae visus est という部分を「観測された」と訳し、具体的な位置や時間の記述を忠実に反映しています。また、図の配置も指示通りに行っています。


ある日、彗星は天秤座のスピカから北東風方向に約2分の距離を離れていた。
同じ日に、マサチューセッツ州ニューポートでの観測によれば、彗星はスピカから一度の弧度を離れ、緯度の差が40分であることを確認し、経度の差は大約44分であった。したがって、そのとき彗星は天秤座18分40秒で南緯1分19秒に位置していた。
同じ日に、マリランド州パトゥンティ川近郊ハンティング・クリーク付近のアーサー・ストーラー博士は、以下の観測を行った。

月曜日の午前3時半（つまりロンドン時間で10時）に、観測者が星囲座の尖角星の上に彗星を観測し、その位置はほぼ重なっていました。同一日中の翌日に同じ観測者が彗星を観測すると、それは尖角星より約2分の弧下にありました。これらの観測結果は、ニューヨーク州で行われた観測結果と一致しますが、彗星の昼間の距離（地球からの距離）をわずかに増加させることで、前日には彗星が星囲座の高度よりも上にあったことを考慮に入れます。

52分の約近似で、翌日午後の同じ星は垂直方向に$2^\circ$40′の高さでした。

11月20日。Paduaの天文学者ドン・モンテナーロ教授は朝6時（ロンドン時間では5時10分）に彗星を見ました。その位置は天秤座23度、南緯1度30分です。同一日、ボストンでは彗星がスピカ西方4度の方向に離れており、そのためには天秤座23度24分近辺でした。

11月21日。ポンテウス及其同行者は朝7時4分に彗星を観測しました。

在第27度45分，赤道坐標系中，当南緯1度16分為最大赤緯時，在第五小時観測。蒙特納魯斯観測為第27度51分。同一天，在牙買加島上観察到接近室女座的起始位置，かつ其赤緯約為織女星的赤緯，即1度59分。

11月22日，蒙特納羅在室女座$2^\circ$33′観測到此星体。而在北美波士ニュートン地区，則大約在室女座3度附近出現，具有相近的赤緯。

ac prius.

次に、\textbf{モンテナロ}は11月24日午後$\scorpio$ 12度52分、11月25日午後$\scorpio$ 17度45分を観測したと報告している。 \textbf{ガレリウス}は既に赤道緯度2度を記述している。

\textbf{ポンテオイゥス}と\textbf{ガレリウス}は同じ結果を得ており、\textbf{モンテナロ}と\textbf{アンゴ}は常に増加していることを確認している。これらの観測は非常に詳細であるが、\textbf{モンテナロ}, \textbf{アンゴ}およびニューイングランドの観測者によるものの方が優れているように思われる。

これら全てを比較検討した結果、午後の位置に到達するまで、その他の観測は考慮されないことにする。

月の摂動論、分点歳差、および彗星の幾何学軌道論（命題XXV〜XLII）

\textbf{第3編}

\textbf{彗星の軌道}

\textit{Londini}, hora mat. 5.\ 10$\minute$ reductis, colligo Cometam hujusmodi cursum quamproxime descripsisse.

\vspace{\baselineskip}

観測日時と座標を示す表は以下の通りである。

\begin{tabular}{*{12}{l}}
               & \textbf{長さ（彗星）} & \textbf{緯度（彗星）} \\
\hline
\textit{Nov.}\ 17 & $\libra$ 8.0 & 0.45 南 \\
            18 & 12.52 & 1.2 \\
            19 & 17.48 & 1.18 \\
            20 & 22.45 & 1.32 \\
            21 & 27.46 & 1.44 \\
            22 & $\scorpio$ 2.48 & 1.55 \\
\end{tabular}

これらの観測データから、彗星はこの近似的な軌道を描いたと結論付けることができる。



表23と表24は、彗星が双曲線軌道上にある同一時間における位置を示しています。

表25

\begin{tabular}{*{12}{l}}
               & \textbf{彗星経度} & \textbf{彗星緯度}     \\
\hline
17日新月   & $\textit{\libra}$ 8.3 & 0.23$\textsuperscript{A}$ \\
21日       & $\textit{\libra}$ 28.0 & 1.22$\textsuperscript{A}$ \\
25日       & $\textit{\scorpio}$ 18.17 & 2.6$\textsuperscript{A}$
\end{tabular}


注：`$\textit{\libra}$` と `$\textit{\scorpio}$` は星座のラテン語表記を日本語に置き換え、`$\textsuperscript{A}$` は脚注の表示方法として使用しました。また、 `\textbf{}` で太字を表現しています。

\begin{center}
\textbf{したがって，11月的観測結果および随后三個月中的観測結果与彗星回太阳在これ一抛物線軌道上的运動相符，因而証明了これ是該彗星的真実軌道。なぜなら観測位置与〜によれば誤差和図形方法定義軌道時計算的位置之間的差异，できる很容易地由誤差引起。}

其他方面，彗星所描描的軌道および它在各個位置投下的真実尾巴，在附図中如下所示。
\end{center}


\par\begin{center}
  \centering
  \includegraphics[width=0.8\textwidth]{486.png}

\end{center}\par

月の光軌道を光学的に描いた図を示す：それに対する観測によって尾の度数を定義する。

\textbf{17年11月17日}\ ポンテオによると、尾はさらに15度以上長く現れた。
\textbf{11月18日} 尾は30度あり、太陽に対して直接対向し、ニューアングルで観察され、また$\mars$星まで伸びていた。そのとき、$\mars$星は$\virgo$座9度54分にあった。
\textbf{11月19日} マリー・ランドで尾が見られた角度は


\par\begin{center}
    \centering
    \includegraphics[width=0.8\textwidth]{tail_angle.png}
    \label{fig:tail_angle}

\end{center}\par

でした。

$15$または$20$日。\textbf{10月}、フラムステディによる観測によれば、彗星の尾は蛇$Ophiuchi$の尾と $Aquila$ 南翼の星 $\delta$ の間の距離を通過し、そして $A, \omega, b$ という星（ $Bayer$ 星表において）に近づき始めた。その終点は $Capricornus 19^\circ30'$ で緯度北 $34^\circ25'$ 左右であった。\textbf{11月}、彗星は矢の頭 $\alpha, \beta$（ $Bayer$ 星表において）に近づき始め、そして $Capricornus 26^\circ43'$ で緯度北 $38^\circ34'$ 左右で終了した。\textbf{12月}、彗星は通過した。

月の運動は矢印座の内側を短時間に進み、$\aquarius$ $4^\circ$まで到達し、その後$\aquarius$ を中心として約 $42\frac{1}{2}$ サイクロン径近傍で終了した。この長さは明るい尾の長さを指すものである。暗い光線と晴れた夜空では、ローマ時間 5 時 40 分（ポンテオ観測）に$\textit{Dec.~12}$ の星座でウルグイの肩から約 $\textit{gr.}\, 10$ の位置まで尾が伸びており、その尾端は北西方向に約 45 分の弧を描いた。この数日間では、尾は約 $\textit{gr.}\, 3$ を中心として伸縮しており、したがってその平均長さは$\textbf{約}$ $\textit{gr.}\, 3$ であった。

\begin{center}

\end{center}

月はその星から南西方向に約2グレゴリウス（15分）ずれている。上端は12月21日、座標が魚座22度61分北緯で、この日にほぼ午後8時に昇る。これはカシオペイア座のβ星とシェディル星から等距離にあり、これら二つの星との距離は互いに等しい。そして12月29日には魚座24度47分30秒北緯でセハット星の左方に接するまで進む。二つの星間の距離は北半球の1フィートに相当する。

\[
\textbf{訳注：}
\]
- グレゴリウス（gr.）は旧暦での角度単位で、1グレゴリウス＝15分を表します。
- 午後8時（cathedram）は古代の時間表記で、現在の午後8時を指すものと解釈されます。
- セハット星（Scheat）は北極座の一部であり、魚座24度47分30秒北緯に位置します。

安デメイアは正確に完了し、長さは54分の54としました。そして、1月5日には、安德梅亜星座の59分25秒で終了し、牡牛座19分35秒で終わることになりました。この星は安德梅亜星座の胸部にあるπ星に接しており、右側にはμ星が存在します。そして（私たの観測によれば）、長さは40分でした。また、これは曲線であり、南向きに凸していました。

太陽と彗星の頭部を通る円周角を形成し、彗星の頭部から4度の角度となりました。しかし、その他の終点では...

月の軌道がその円周に対して角度約$10^\circ$または$11^\circ$に傾斜しており、尾の弦はその円周と角度$8^\circ$を形成していました。$1$月$13$日には、尾の端は十分に明るい光でアルメキウス（$\textbf{Alamech}$）とアルゴル（$\textbf{Algol}$）の間で終了し、最弱の光では$\kappa$星周辺のペルセウス座の領域から消失しました。尾の端点から太陽までの距離は約$3^\circ50′$であり、尾の弦がその円周に対して角度$81/2^\circ$を形成していました。$1$月$25$日と$26$日に、最弱の光で長さ$6^\circ$または$7^\circ$に微かに輝いていました。

天上が極めて晴れ渡っていたとき、最も微細な光でも強く感じられるほど、その軌道は十二分の一度以上に達した。
そしてその軸は正確にオリエント・ヒムノス（ Auriga 星座の東方）の Lucidam （一星）へ向かって進行し、太陽との対立点から北方向へ約十分の一角度の傾きを示していた。ついで、1月10日には武装した目で二分の一度長い尾を見た。その光は前記のものよりもさらに微細であり、ガラスを通じて現れなかった。一方、ポンテウスは1月7日に長さが十二分の一度の尾を観測していると記している。

この彗星の軌道が描き終わり、その他の現象も観察した後で、


Orbem jam descriptum spectanti \& reliqua Comet{\ae} hujus Ph{\ae}nomena.

確かに、彗星の体は固体で、堅固かつ固定されたものであり、惑星と同様に堅牢であることが容易に理解できます。なぜなら、もし彗星がただの地球蒸気や太陽蒸気、惑星蒸気の集合体であるとするならば、彗星は現在の軌道上での太陽近接時に即座に消滅するはずです。太陽の熱は放射線密度と等価であり、これは距離の逆数に比例し、すなわち太陽から離れる距離の平方に反比例します。したがって、彗星が太陽から距離 \( \text{Dec.~8} \) で近日点を通過していたとき、その距離は地球から太陽までの距離に対して約 \( 6 : 1000 \) の関係にあるとすると、その当時の彗星周囲の太陽の熱は以下のようになります。

\[
\text{太陽の熱} = \frac{\text{放射線密度}}{\text{距離}^2}
\]

ここで、彗星から太陽までの距離 \( d_{\text{彗星}} \) が地球から太陽までの距離 \( d_{\text{地球}} \) の約 \( \frac{6}{1000} \) 倍であるとすると、

\[
\left( \frac{d_{\text{彗星}}}{d_{\text{地球}}} \right)^2 = \left( \frac{6}{1000} \right)^2
\]

したがって、その時の太陽の熱は

\[
\text{太陽の熱} = \frac{\text{放射線密度}}{\left( \frac{6}{1000} \right)^2 d_{\text{地球}}^2}
\]

となります。

熱は、太陽の夏季の輻射によるもので、私たちは約 \(1000000\) に対して \(36\)、すなわち約 \(28000\) と推定しています。しかし、沸騰する水の熱は、乾燥した土地が夏季の太陽から吸収する熱よりも約三倍大きいです。これは経験的に確認されています。また、融解中の鉄の熱（正しく予測すれば）は、沸騰する水の熱よりも約三倍か四倍大きいでしょう。従って、乾燥した土地が彗星を通過中に太陽光線から吸収する熱は、融解中の鉄の熱よりも約二万倍大きいと考えられます。この程度の熱により、蒸気や揮発物質、そして全体として

\[
\textbf{すべての} \quad \textit{物質が大量に放出される可能性があります。}
\]

月の摂動論、分点歳差、および彗星の幾何学軌道論（命題XXV〜XLII）

彗星は太陽に近づくとき、即ち近日点において、極めて大きな熱を吸収し、その熱を長時間保持することができる。例えば、鉄の球が発光している場合、指の幅ほどの一方の面積が1時間の間に空気中にほとんど熱を放出しない。しかし、より大きい球は、直径に比例して長い間熱を保持することができる。これは、周囲の空気が接触することで冷やされる表面積（その大きさは、球内の熱い物質の量に対して反比例します）がその比率で小さくなるためである。したがって、鉄の発光する球は地球に比してより長い間熱を保持する。

\[
\textbf{訳注：}
\]

この部分では、ニュートンは彗星が近日点において大量の熱を吸収し、その熱を長時間保持できる理由について説明しています。鉄の球体の例を使い、表面積と内部の熱量の関係性を示しており、これは彗星の場合にも同様に適用されるという考えを述べています。

\textbf{第3編 月の摂動論、分点歳差、および彗星の幾何学軌道論（命題XXV〜XLII）}

{\ae}qualis, id est pedes plus minus 40000000 latus, diebus totidem \& idcirco annis 50000, vix refrigesceret. Suspicor tamen quod duratio Caloris ob causas latentes augeatur in minore ratione quam ea diametri: \& optarim rationem veram per experimenta investigari.

Porrò notandum est quod Cometa Mense \textit{Decembri}, ubi ad Solem modò incaluerat, caudam emittebat longe majorem \& splendidiorem quam antea Mense \textit{Novembri}; ubi perihelium nondum attigerat. Et universaliter caudae omnes maximae \& fulgentissimae e Cometis oriuntur, statim post transitum eorum per regionem Solis.

\begin{itemize}
    \item その長さは約4000万ペデ、すなわち約50000年でほとんど冷めないでしょう。しかし、熱の持続時間は隠れた原因により直径に対する比率よりも小さい比率で増加すると推察します。そして、実験を通じてその真の比率を明らかにしたいと思います。
    \item さらに注意すべきは、12月に太陽に向かって急速に進み始めた彗星が、前月である11月と比較して遥かに長い且つ明るい尾を放ったことです。これはまだ periherium（近日点）に達していないときの現象です。
    \item 一般に、彗星の通過した太陽周辺で最も大きいかつ最も輝く尾は常に生成されます。
\end{itemize}

月の加熱は彗星の尾の大きさに影響を及ぼす。したがって、尾は虚無的な蒸気の一端であると考えられる。そして、彗星の頭部または核からその熱によって放出されるものだと結論付けることができる。

他には、彗星の尾に関する三つの主張がある。一つ目は、尾は太陽光線を通じて彗星の透明な頭部で拡散した髪であるとする主張だ。二つ目は、尾は彗星の頭部から地球に向かって進む際に光の屈折によって生じるという主張だ。そして最後に、尾は彗星の頭部から常に昇り続け、太陽から遠ざかる方向へと去っていく雲または蒸気であるとする主張がある。これらの最初の主張は光学に関する知識がまだ十分でない人々によるものであり、彼らは尾を髪と考えている。



これらの主張はそれぞれ異なる視点から尾の性質を説明しようとするものである。太陽光線を通じて拡散した髪という考え方は、尾が透明な頭部から放出されたものとして理解される。一方で、屈折による光の変化により生じた尾は、彗星の進路と地球との関係性によって形成される。最後に、雲または蒸気の主張は、尾が恒常的に昇り続け、太陽から遠ざかる方向へ去っていくという観察結果に基づくものである。



これらの主張はそれぞれ異なる視点から尾の性質を説明しようとするものであり、それぞれが特定の観察や理論に基づいて形成されている。しかしながら、これら全ての主張は彗星の尾の実際の構造と一致していない可能性がある。

月は暗室の中で直接見ることはなく、その光が空気中の粉塵と煙の粒子によって反射される限りである。したがって、空気が濃い煙で満たされているときにはより明るく、眼に与える感覚も強く感じられる。一方、空気が薄いときには弱く感じられる。しかし、天界においては反射する物質がない限り光は存在しない。光は直接目に入ることなく、反射されるまで見えない。視覚は放射線が目に衝突することによってのみ成立する。したがって、反射を行う何かを含む領域が必要である。

\[
\textbf{訳文}
\]


\[
\textbf{数学的注釈}
\]

ここでニュートンは、月の光が地球の大気中の粒子によって反射されることを説明しています。大気が濃い煙で満たされているときには、その反射により月はより明るく見え、眼に与える感覚も強くなる一方、空気が薄いときには弱く感じられるという現象を述べています。また、天界において光が直接目に入ることなく反射されるまで存在しないことを指摘し、視覚は放射線が目に衝突することによってのみ成立すると結論付けます。

\[
\textbf{図表}
\]



この図は空気中の粒子が月の光を反射する様子を示しています。濃い煙で満たされた空気では、反射により月はより明るく見え、眼に与える感覚も強くなることが示されています。

月々の尾は、太陽光が均一に輝くように宇宙全体を照らすことを妨げません。この見解は多くの人々を押さえ込んでいます。実際、月々の尾は決して色で変化しませんが、それらはしばしば屈折と密接に関連しています。固定星および惑星からの光が明確に私達へ到達することから、宇宙の媒質は屈折作用によって影響を受けないことが示されます。実際、エジプト人によって「彗星を梳いた」と呼ばれた彗星が時々見られたとしても、これは稀な出来事であり、それは偶然の雲の屈折によるものです。固定星からの放射と光沢も、屈折により影響を受けることが示されます。

太い望遠鏡や空気の振動を考慮する必要がある：実際、望遠鏡に取り付けられた目は容易に消えてしまう。また、空気中の昇蒸汽の振動により、光線は狭い瞳孔領域を通じて多次元的に回転し、しかし厚いガラスの物体面では全く影響を受けない。

したがって、前の場合には散光が生じる一方で、後者の場合には散光が消失する。そして後者での散光の消失は、光線が空気中を伝搬する際には完全な屈折を感じさせないことを示す。

誰も主張しないように：彗星の尾はその輝度が十分ではないため、望遠鏡で観察されないことが普通である。なぜなら、二次的な光線はその影響を受けないからである。

非光学的力が十分でなく、目で動かすことができないため、固定星の尾は観測されません。科学的に知るべきは、望遠鏡を介して固定星の光は百倍以上強化される可能性があるにもかかわらず、その尾は見つからないことです。惑星の光は豊富ですが、尾は存在しません。しかし、彗星はしばしば非常に長い尾を持ち、頭部の光が薄く鋭い場合があります。例えば1680年の彗星は、その頭部が第2等星とほぼ同じ明るさだった時期に、長さ$40^\circ$、$50^\circ$、$60^\circ$まで目立つ尾を放ち始めました。その後、1月27日と28日に頭部は次のようになって現れました。

月に比して非常に大きい星でありながら、その尾は光ではあるが、十分な明るさではなく、6または7度の長さがあり、極めて暗く、ほとんど見ることができなかった。この尾は12度目までの角度まで延びており、先述した通りである。しかし、2月9日と10日に裸眼で頭部を確認できなくなったとき、望遠鏡を使って2つの角度の長い尾を見た。もし尾が天体の屈折によって発生し、太陽の対立点から気層の形状に偏曲されるのであれば、その偏曲は同じ領域において同一方向に起こるべきである。



この証拠は尾が屈折によって生じるという仮説を否定するものである。

月の摂動論、分点歳差、および彗星の幾何学軌道論（命題XXV〜XLII）において、以下は1680年12月28日午後8時30分ロンドン時間に彗星が位置していた部分です。そのとき、彗星は$\pisces$ $8^\circ$41' 南緯$28^\circ$6'で太陽は$\capricornus$ $18^\circ$26'に存在していました。また、1577年12月29日に彗星が位置していたときも同様で、彗星は$\pisces$ $8^\circ$41' 南緯$28^\circ$40'で太陽は約$\capricornus$ $18^\circ$26'に存在していました。両者の場合において、地球は同じ場所を向いており、彗星は同じ天体に現れていました。

\[
\textbf{訳注：}
\]

この部分では、ニュートンが1680年と1577年の2つの観測データを比較しています。両者の共通点として、彗星の位置（$\pisces$ $8^\circ$41'）と太陽の位置（$\capricornus$ $18^\circ$26'）が一致していることが指摘されています。また、地球も同じ方向を向いていたことから、彗星は同じ天体に見えたという結論が導かれています。

\[
\textbf{補足：}
\]

これらの観測結果は、ニュートンの理論を検証するための重要なデータとなりました。彗星の位置と太陽との関係性から、地球の運動や天体の軌道について深い洞察を得ることができたのです。

月の尾の摂動、分点歳差、および彗星の幾何学的軌道論（命題XXV〜XLII）【第3編】

パート：しかしながら前記の場合においては、彗星の尾（私の観測と他者の観測に基づくもの）は太陽に対し北西方向に角度 $4\frac{1}{2}$ で傾斜していたが、後期には南方向に角度 21 で傾斜していた。したがって、天体の屈折を否定するとすれば、尾の現象は反射する何かから導かれるものであることが残る。

また、尾は頭から始まり、太陽から離れた地域へ上昇することも観測法によって確認される。すなわち、平面彗星軌道において太陽を通過する直線上に存在し、太陽に対して偏角を描くものについては、

\[
\textbf{命題XXV:} \quad \text{彗星の尾は頭から始まり、太陽から離れた地域へ上昇することを確認される。}
\]

この論証は以下の通りである。

\begin{enumerate}
    \item 平面彗星軌道において、彗星が太陽を通過する直線上に存在し、太陽に対して偏角 $\theta$ を描く。
    \item これらの観測結果から、尾の起源と上昇は頭から始まり、太陽から離れた地域へと進展することを推定できる。
\end{enumerate}

以上のように、彗星尾の現象は反射する何かから導かれるものであることが示唆される。

月の対日位置において、彗星は常にその軌道上の特定の部分を向いている。観測者がこれらの平面に固定されている場合、彗星が太陽から直接反対方向にある部分で見え、観測者がこれらの平面から離れるにつれて、彗星の偏角は徐々に感じられ、最終的にはより大きな角度で見られるようになる。この偏角は他の条件が等しいときでも、彗尾が彗星の軌道に対して傾きが大きいほど小さくなり、また彗頭が太陽に近づくほど小さくなる。特に彗頭近くでの偏角の角度を観測する場合、その値はより小さいものとなる。

さらに、彗尾が偏角を示さない場合は直線のように見え、偏角を示す場合はそれらが異なる方向に見える。

月の軌道はまた曲がる。曲がり方は、偏角が大きい場所で大きく、また尾が他の部分と比べて長ければより敏感に曲がる。短い軌道では特にこの傾向が顕著である。

偏角の角度は彗星の頭部近くでは小さく、尾の端部近くでは大きくなる。したがって、尾は凸側を向いており、その部分からの偏角は彗星の頭部を通る直線（太陽から彗星の頭部まで延びた直線）に沿った直線分である。

また、長く広い尾はより明るく光り、これらの尾は凸側を向いている。

月の尾が輝き、その端はより明瞭な境界で終わるほど、凹面に比べて優れています。したがって、彗星の現象は首の運動によって生じます。天体の位置に関係なく、首が観測される場所からではなく、首が見られる場所までです。そのため、これらの現象は屈折によるものではなく、首が物質を供給することで生じます。

実際、私たちの空気中で燃えた物体から発生する煙や蒸気が上昇し、物体が静止している場合に垂直方向に、運動している場合は斜め前方に向かうようにするのと同様に、太陽を中心に重力によって回転している天体からは、煙や蒸気が太陽から上昇すべきであり（既述）、そして上方向へ直線的に進むべきです。

静止するか、斜めに動くかは、物体が上部の蒸気を通過しながら進む場合によって異なります。すなわち、上部の蒸気が昇った後で物体がその場所を離れるなら、静止します。また、蒸気が斜め方向に昇る場合は、物体も斜めに動きます。そして、蒸気が昇る速度が速いほど、その傾斜は小さくなります。特に太陽近傍や煙る物体の近くでは、この傾斜は小さいです。

さらに、この傾斜により蒸気の柱は曲がります。また、前の列の蒸気が少し新鮮であるため、その列はそこよりも少し濃く、光をより豊かに反射し、限界もより明瞭になります。突然で予測不能な尾の揺れや、その形状の不規則性について、一部の人々は時折述べていますが、ここでは詳しく説明しません。


月の軌道や彗星の尾などから生じる雲々について述べます。ここでは何も追加しません；これは、空気の変動や雲の尾が一部で見えにくくなること、あるいは乳暈の一部と混ざって見えることによるものであると考えられます。

しかし、これほど広大な空間を満たすのに足りるだけの蒸気が、彗星の大気から生じうることは、私たちの空気の希薄さから理解されるだろう。地表近くの空気は、同じ重さの水のおよそ850倍の空間を占める。したがって、高さ850フィートの円柱状の空気の柱は、同じ底面をもつ高さ1フィートの水の柱と同じ重さである。ところで、大気の頂上にまで達する空気の柱は、その重さが、高さおよそ33フィートの水の柱に等しい。ゆえに、空気の柱全体から、下の高さ850フィートの部分を取り除けば、残りの上の部分の重さは、高さ32フィートの水の柱に等しい。ここから、（空気の圧縮は上にのしかかる大気の重さに比例し、重力は場所の地球の中心からの距離の平方に反比例するという、多くの実験で確かめられた仮説にもとづいて）第2編 命題XXIIの系によって計算すると、私は次のことを見いだした。空気は、地球の表面から地球の半径1つ分の高さまで上ったところでは、私たちのところよりも希薄であって、その比は、土星の軌道より内側の空間全体の、直径1インチの球に対する比よりも、はるかに大きい。したがって、地球の半径1つ分の高さでの希薄さをもつ私たちの空気の、一辺1インチの球は、土星の天球に至るまで、さらにそのはるか先まで、すべての惑星の領域を満たすだろう。ゆえに、空気はそれより高くなると限りなく希薄になり、彗星のコマ、すなわち大気は、その中心から上って、核の表面のおよそ10倍も高く、さらに尾はもっと高く上るのだから、尾はきわめて希薄でなければならない。もっとも、彗星の大気がずっと厚いこと、また物体の太陽への重力、および空気や蒸気の粒子が互いに及ぼす重力のために、天空の空間や彗星の尾の中では、空気がそれほど希薄にならない、ということも起こりうるかもしれない。しかし、ごくわずかな量の空気と蒸気で、尾のあらゆる現象を説明するのに十分であることは、この計算から明らかである。

また、尾の著しい希薄さは、尾をとおして透けて見える星からも推定される。太陽の光に輝く地球の大気は、わずか数マイルの厚さで、すべての星も月自身もすっかり覆い隠して消してしまう。ところが、同じように太陽の光に照らされた、尾の巨大な厚さをとおして、きわめて小さな星でも、明るさを損なわずに透けて見えることが知られている。そして多くの尾の輝きは、暗い部屋の中で、太陽の光を光の筋として反射する私たちの空気の、1インチか2インチの幅の輝きより大きいことは、ふつうない。
水蒸気が頭から尾端まで上昇した時間は、尾端から太陽へと引いた直線を描き、その点を標識することで推定できます。

月の摂動論、分点歳差、および彗星の幾何学軌道論（命題XXV〜XLII）【第3編】

そこで、直線がその軌道と交わる地点を特定します。すなわち、蒸気は太陽から垂直に上昇すると、彗星がその上昇前の運動を保持しながら、頭部が交差点の位置にあるときに頭部から垂直に上昇を開始します。しかし、蒸気は太陽から垂直に上昇せず、彗星の上昇前の運動を保持しながら斜めに上昇します。したがって、問題の解釈においてより正確なものは、その直線が尾の長さと平行であるか、あるいは（彗星の曲線運動により）尾の線から離れる方向にあることを考慮に入れるべきです。

\textbf{証明}

まず、彗星が太陽から垂直に上昇する場合を考えてみましょう。このとき、蒸気は彗星の頭部が交差点の位置にあるときに垂直に上昇を開始します。しかし、実際には彗星の運動は曲線であり、その上昇前の運動を保持しながら斜めに上昇します。したがって、直線が尾の長さと平行であるか、あるいは尾の線から離れる方向にあることがより正確な解釈となります。

\textbf{命題}

以上の理由により、直線が尾の長さと平行であるか、または尾の線から離れる方向にあることを証明します。

これにより、私は1月25日に太陽軌道の尾が蒸気を放出し始めたことを発見しました。その蒸気は12月11日以前の頭部から昇り始めましたので、その昇り時間は約45日余り消費しました。その後、12月10日に完全に現れた尾は、彗星が近日点を通過してからそれまでの2日間分の距離を昇天し続けました。したがって、蒸気は太陽近傍で最も速く昇り始めましたが、重力による運動と共にその昇りは遅れ続け、さらに昇るにつれて尾の長さが増しました。そして尾が現れた間、ほぼ全般的に蒸気が放出されていました。

\[
\textbf{注記} \ref{wasp504}：尾の長さは蒸気の昇りとともに増加しました。
\]

次に、彗星の尾について述べます。彗星が太陽から昇り始めたとき、その尾も最初に昇り始め、そして尾の終端は最初に昇った蒸気によって形成されます。この蒸気が完全に消失するまで、尾自体は見え続けることになります。特に短い尾を持つ彗星についても同様です：これらの尾は急速で恒久的な運動を伴って頭から昇り始めますが、すぐに消えてしまいます。それらは実際には非常にゆっくりとした運動によって拡散された蒸気の柱であり、その運動は彗星の頭が太陽に近づく初期に獲得したものを共有します。これらの尾は頭と共に空気中を移動し続けます。そして結論として：

\textbf{彗星の尾は、頭からゆっくりとした運動によって拡散された蒸気の柱であり、その運動は彗星が太陽に近づく初期に獲得したものを共有します。これらの尾は頭と共に空気中を移動し続けます。}

spatia caelestia vi resistendi destitui; utpote in quibus non solum solida Planetarum \& Cometarum corpora, sed etiam rarissimi caudarum vapores motus suos velocissimos liberrime peragunt ac diutissime conservant.

Ascensum caudarum ex Atmosphaeris capitum \& progressum in partes a Sole aversas \textit{Keplerus} ascribit actioni radiorum lucis materiam caudae secum rapientium. Et auram longe tenuissimam in spatiis liberrimis actioni radiorum cedere, non est a ratione prorsus alienum, non obstante quod substantiae crassae, impeditissimis in regionibus,



\textbf{訳文}

天体空間に抵抗力は存在しない；つまり、その中では不規則な形状の惑星や彗星の固体物質だけでなく、極めて希薄な尾の蒸気が非常に速く自由に移動し、長期間保存される。

ケプラーは、尾の上昇は大気層の頂点から始められ、太陽を向かない方向へ進展することを、光線の放射が尾と共に引きずる作用により説明している。そして、極めて希薄な媒質が自由空間において光線の作用を受け入れることは、根本的には不合理とは言えないと考えられる。しかし、物質が非常に濃密で、特にその領域は阻害されることを考慮すると、



注記：図494は尾の上昇と太陽からの遠ざかり方を示している。

月は、太陽の光から微細に推進されない。一方で、彼は軽量粒子と重い粒子が存在すると考え、また矢継ぎ早に彗星の尾を浮遊させるとも主張する。しかし地球の物体の重力は物質量に比例し、その量は変化しないことから、私は尾の上昇は物質の希薄化によるものだと推測する。煙は風の衝撃により通路を上昇する。空気は熱によって希薄になり、その特定の重力を減少させ、煙も一緒に引き上げられる。彗星の尾も同様に上昇する理屈は同じではないかと推察される。

月が太陽から同じ方法で上昇したか？すなわち、太陽の放射は通過する媒質を動かさない（反射や屈折以外では）か？反射する粒子はその作用により熱されることで太陽周囲の気体を暖め、それらが包囲されている。これらの粒子は熱によって自己に伝わり、密度が減少し、以前は太陽に向かって特定の密度の重力を向いていたが、現在は上昇し、彗星の尾を構成する反射粒子も一緒に引き上げる。また、それらが太陽を中心に回転していることにより、その作用によって太陽から逃れようとすることが、気体と雲の材料である太陽の大気は静止しているか、または運動しているかに関係なく、蒸発の上昇に寄与する。

\begin{center}

\end{center}

月に近い場所では、より曲がった軌道を描くため、彗星の尾は遅く回転します。これは彗星の尾が太陽近くで上昇しやすくする要因であり、その理由は、太陽周囲の濃密な太陽気体層により、彗星の軌道がより曲がりやすく、それに伴い尾が非常に長いものとなるからです。実際には、当該尾はその運動を保ちつつ、重力に従って太陽に向かって移動し、彗頭と共に楕円軌道上を回ります。そしてこの運動により、彗頭は常に自由に尾の追従を許します。気体の重力は彗頭から尾が後ろに落ちる太陽側へと影響しないためです。

\[
\textbf{訳文：}
\]

月に近い場所では、太陽の回転から受けた影響を受けた彗星は遅く回転します。これは彗星の尾が太陽近くで上昇しやすくする要因であり、その理由は、太陽周囲の濃密な太陽気体層により、彗星の軌道がより曲がりやすくなるからです。それに伴い、彗星は太陽に重力によって引き寄せられ、尾を非常に長いものとして放出します。実際には、当該尾はその運動を保ちつつ、重力に従って太陽に向かって移動し、彗頭と共に楕円軌道上を回ります。そしてこの運動により、彗頭は常に自由に尾の追従を許します。気体の重力は彗頭から尾が後ろに落ちる太陽側へと影響しないためです。

\[
\textbf{数学的記述：}
\]

\[
F \propto 1/r^2
\]

\textbf{第三編 月の摂動論、分点歳差、および彗星の幾何学軌道論}

彗星尾部的重力作用によりその軌道をどのように決めるかについて述べる。

彗星尾部と頭部は共同の重力によってまたは同時に太陽に落下し、または同時に昇天時に遅滞するので、その重力は彗星尾部と頭部が既述の原因からどのような位置関係にも自由に受け入れ、その後自由に保持されることを妨げない。

\textbf{彗星尾部は彗星のperiastron（近地点）から生じ、遠方へ移動し、その頭部と共に遠くへ去る。そして、これらはしばしば長い年月を経て再び私達の側に戻り、またはそれらが}


\par\begin{center}
    \centering
    \includegraphics[width=0.8\textwidth]{495.png}
    \label{fig:comet_orbit}

\end{center}\par

稀薄化は徐々に消失する。その後、彗星の尾が太陽に向かって下降し、新しい短い尾が首からゆっくりと伝播する必要がある。そして、その彗星が太陽の大気層に降り立つ近接点（Perihelion）においては、尾は大幅に増大する。実際、これらの空間は非常に自由で、蒸気が常に希薄化し、拡張されるからである。この理由により、全ての尾は頂点側が広く、首の近くでは狭くなる。しかし、その蒸気は永久的に拡散し、宇宙全体に散布され、その後重力によって惑星へ徐々に引き寄せられ、それらの大気層と相互作用する。


\par\begin{center}
  \centering
  \includegraphics[width=0.8\textwidth]{彗星尾の拡張.png}
  \label{fig:彗星尾}

\end{center}\par

この過程は、尾が太陽に近づくにつれて急速に増大し、遠ざかるにつれて徐々に減少する特徴を持つ。また、尾の形状はその位置によって異なるため、首の近くでは狭く、頂点側では広くなる傾向がある。


\par\begin{center}
  \centering
  \includegraphics[width=0.8\textwidth]{彗星尾の形状.png}
  \label{fig:彗星尾形状}

\end{center}\par

これらの論証は、彗星尾の運動とその大気環境との相互作用を理解する上で重要であり、彗星の軌道や太陽系の大気構造についても深い洞察を提供する。

misceri rationi consentaneum videtur. Nam quemadmodum Maria ad constitutionem Terrae hujus omnino requiruntur, idque ut ex iis per calorem Solis vapores copiose satis excitentur, qui vel in nubes coacti decidant in pluviis, \& terram omnem ad procreationem vegetabilium irrigent \& nutriant; vel in frigidis montium verticibus condensati (ut aliqui cum ratione philosophantur) decurrant in fontes \& flumina: sic ad conservationem marium \& humorum in Planetis Cometae requiri videntur; ex quorum exhalationibus \& vaporibus condensatis, quicquid liquoris per vegetationem \& putrefactionem

\textbf{生成}されるべきである。すなわち、地球の形成に必要なマリアが要求されるように、太陽熱による蒸気の大量発生が求められ、その蒸気が雲となって降雨となり、すべての土地を植物の生育と養分供給のために潤すか、または寒冷な山頂で凝結して泉や川へ流れ込むかのように、月や彗星の水や液状物質の保存にも似たような過程が必要である。これらの発散した蒸気や凝結された蒸気から、植物の生育と腐敗によって生成されるすべての液体が得られるべきである。

月が地表に衝突し、乾土となると、その量は継続的に補給され、再生成されることも可能である。実際、すべての植物は完全に液体から成長し、その後大部分は腐敗によって乾土に戻り、常に腐敗した液体から泥が降り注ぐ。

これにより、乾土層は地殻内側で増加し、もし他所からの補給がない限り、液体は永遠に減少し、ついには枯渇するだろう。しかしながら、その気体は私たちの空気の最小限の部分でありながら極めて微細で優れたもので、万物の生命を必要とすると推察される。これらは主に彗星から来ると考えられる。

\[
\textbf{訳文}
\]



\textbf{数学的注釈}
 

この文脈では、ニュートンは自然現象の説明において気体の役割を強調しています。特に、乾土層の形成と維持について述べていますが、その補充源として彗星の存在を指摘しています。

月の周囲の大気は、彗星が太陽に近づくときに尾を短縮し、（特に彗星が太陽を直視する部分では）狭くなりますが、離れるときには尾が長くなり、それほど頻繁に短時間にその様子を見たのはヘベリウスだけです。また、尾の最も明るい部分は、彗星の頭部がまだ太陽に近づいているときに最大となり、その先端は大気の最も下層で濃い煙によって覆われていることが多いと思われます。実際、すべての大気中の煙は非常に広範囲にわたって拡散しています。

\[
\textbf{訳文}
\]

彗星の大気は、太陽に向かって下降するとき尾が減少し、（特に彗星が太陽を直視する部分では）狭くなりますが、離れるときには尾が広がり、ヘベリウスが正確に記録したように、その様子は頻繁に見られていません。また、尾の最も明るい部分は、彗星の頭部がまだ太陽に近づいているときに最大となり、その先端は大気の最も下層で濃い煙によって覆われていることが多いと思われます。実際、すべての大気中の煙は非常に広範囲にわたって拡散しています。

\[
\textbf{注}
\]

-  ここで説明した現象の一部を示す図が添付されています。
-  また、尾の広がりと縮小の様子をさらに詳しく示す別の図も提供されています。

月の摂動論、分点歳差、および彗星の幾何学軌道論（命題XXV〜XLII）【第236節】

熱を帯びた彗星は通常より太く且つ黒くなる。したがって、我々が扱った彗星については、日と地球から等しい距離にあるときよりも、接近日点後には暗く見えます。10月の月に第三等星程度の明るさを示す一方で、9月の月には第一等星や第二等星ほど明るく見せました。どちらを見たかによっては、前者がより大きな彗星であると報告します。例えば、カンタブリからあるジューヴェニスが11月19日に彗星を観測し、その光は銀色で鈍いものの、室女座の穂先よりも明るく輝き、その後ではより暗く見えたと報告しています。

D. ストーラーは、彼の記述が12月に到着したと述べています。この時期に彗星が最大かつ最も輝く尾を放ち、11月に太陽昇る前に現れた小さな尾を持つ比較的小さな彗星が見られたことがありました。これらの現象の理由は、彗星の頭部の材料が最初には豊富だったが、徐々に消費されたことだと推測されるところです。

また、他の彗星の頭部も同様に、最大かつ最も輝く尾を放った後、暗く小さなものとして描かれる傾向があるようです。実際、1668年3月5日聖マルティヌスの日に7時頃には、彗星の頭部が不明瞭で小さなものとして観測されました。

【第3編 月の摂動論、分点歳差、および彗星の幾何学軌道論】

夜明け前 \textbf{バルセロナのValentinus Estancius}が観測した彗星は、冬至日夕方近い地平線に見られ、頭部は僅かに見え、しかし尾は異常に輝いていた。立っている人々は、海辺から反射されたその姿を容易に確認した。彗星の形状は明るさ23度の帯状で、西方から南方へと傾き、地平線に近い平行線であった。しかしながらその輝度は僅か3日間持続し、その後著しく減少し始めた一方で、尾の長さは輝度が減少するにつれて増大した。そして彗星はさらに赤道より南緯45度付近（すなわち天球の約四分の一）を占領していた。



月は西方から東方へと輝きを拡大し、その光は非常に強く放たれます；
\begin{center}

\end{center}
しかし、完全には現れませんでした。これらの地域では、常に地平線以下で輝度が衰えました。光の増加と彗尾の減少から明らかに、彗頭は太陽から遠ざかり、最初の段階では太陽に最も近かったことがわかります。これは1680年彗星の通常の傾向と一致します。また、1101年または1106年の彗星についても類似の記述があります。その彗星は小さな暗い星でしたが（1680年に見られたように）、そこから放たれた光は非常に明るく、そして東方と北に向かって広がる巨大な帯のように見えました。これは\textit{ヘベリウス}によって報告されています。

月の摂動論、分点歳差、および彗星の幾何学軌道論（命題XXV〜XLII）

\textbf{Simeone Dunelmensi} の修道士は、\textit{2月}の初め、夕暮れ頃に出現し、太陽沈む直前に観測されたと記録している。その後、パリスのマタイウスによれば、その彗星は太陽から約一肘離れており、午後3時（正確には6時）から午後9時まで長さを放っているという。

この彗星はアリストテレスによって記述されたものと同様に熱々であり、それは『 Meteor.\ 1, c. 6 \textit{Lib.\ 1}\ Meteor.\ 6』で述べられている。その頭は最初の日に見えず、太陽よりも前か、少なくともその放射線の内側にあるためである。



 
$\textbf{Simeone Dunelmensi}$ の修道士によると、\textbf{2月}の初めに夕暮れ頃に出現し、太陽沈む直前に観測された。その後、パリスのマタイウスは次のように述べている：「その彗星は太陽から約一肘離れており、午後$3$時（正確には$6$時）から午後$9$時まで長さを放っている」。この彗星はアリストテレスによって記述されたものと同様に熱々であり、それは『 $Meteor.\ 1, c. 6 \textit{Lib.\ 1}\ Meteor.\ 6$』で述べられている。その頭は最初の日に見えず、太陽よりも前か、少なくともその放射線の内側にあるためである。

\begin{align*}
& \text{2月の初めに夕暮れ頃に出現し、}\\
& \text{太陽沈む直前に観測された。}\\
& \\
\textbf{パリスのマタイウス}
& \text{「その彗星は太陽から約一肘離れており、}\\
& \text{午後3時（正確には6時）から午後9時まで長さを放っている」}\\
& \\
& \textbf{アリストテレス}
& \text{『 Meteor.\ 1, c. 6 Lib.\ 1 Meteor.\ 6』で述べられている。}\\
& \text{その頭は最初の日に見えず、太陽よりも前か、}\\
& \text{少なくともその放射線の内側にあるためである}
\end{align*}

月が太陽に沈むと、次の日には可能な限り早く消えました。すなわち、彗星は太陽から最も近い距離で残存し、その後すぐに消失しました。アリストテレスによれば、その原因は彗星の尾の過度の熱さによるものでした。彗星の頭部に散在する火がまだ見えなかったとき、尾の炎が弱まり始めると、彗星の顔色が回復し、その輝きは太陽から約三分の一（すなわち60度）まで拡大しました。彗星は冬の間出現し、オリオン帯に上昇した後で消失しました。

\begin{center}

\end{center}

この彗星は太陽から約60度まで輝きを拡大したので、その軌道は図に示すように推定されました。

月の摂動論、分点歳差、および彗星の幾何学軌道論（命題XXV〜XLII）

その1618年の年には、最も太陽光線に照らされた彗星が現れ、
大小の恒星とほぼ等しく見えたり、それ以上の大きさを示したが、
その後、いくつかの彗星が短い尾を持ちながらもより大きなものとなった。

これらの中には、誰かはその彗星を木星と等しいと主張し、
他の人々は金星や月と同程度であると述べた。

我々は彗星を太陽を中心に非常に偏心率の大きい軌道で回る惑星の一種だと説明した。
そして、非尾を持つ惑星では、より小さなものは軌道が小さく太陽に近い所で回ることから推測され、
同様に、彗星もその近日点で太陽に最も接近するものほど小さく、また軌道が小さいとされる。

\textbf{証明}：彗星は太陽を中心に非常に偏心率の大きい軌道を描き、
その軌道は惑星の軌道よりも大きく、そしてより近接点で太陽に接近する。
したがって、彗星は惑星ほど大きなものであることが推測される。

\textbf{命題XXVI}：彗星は太陽を中心に非常に偏心率の大きい軌道を描く。また、その軌道は惑星の軌道よりも大きく、そしてより近接点で太陽に接近する。したがって、彗星は惑星ほど大きなものであることが推測される。

\textbf{命題XXVII}：彗星の尾は太陽光線によって形成される。彗星が太陽に向かって進むとき、その尾は太陽から遠ざかる方向に伸びる。また、彗星が太陽から離れるとき、尾は太陽に向かう方向に伸びる。

\textbf{命題XXVIII}：彗星の尾の長さは彗星が太陽に接近するほど増加し、彗星が太陽から遠ざかるほど減少する。また、彗星が太陽に最も近い点（近日点）で最大の長さを示す。

\textbf{命題XXIX}：彗星の尾は太陽光線によって推進される。彗星が太陽に向かって進むとき、尾は太陽から遠ざかる方向に伸びる。また、彗星が太陽から離れるとき、尾は太陽に向かう方向に伸びる。

\textbf{命題XXX}：彗星の軌道は太陽を中心に描かれ、その軌道は非常に偏心率が高い。彗星はその軌道で太陽に最も接近する点（近日点）と最も遠ざかる点（遠日点）がある。

\textbf{命題XXXI}：彗星の尾の形状は彗星が太陽に向かって進むとき、尾は太陽から遠ざかる方向に伸びる。また、彗星が太陽から離れるとき、尾は太陽に向かう方向に伸びる。

\textbf{命題XXXII}：彗星の尾の長さは彗星が太陽に最も接近する点（近日点）で最大となる。彗星が太陽から遠ざかる点（遠日点）では、尾は短くなる。

\textbf{命題XXXIII}：彗星の尾は太陽光線によって推進される。彗星が太陽に向かって進むとき、尾は太陽から遠ざかる方向に伸びる。また、彗星が太陽から離れるとき、尾は太陽に向かう方向に伸びる。

\textbf{命題XXXIV}：彗星の軌道は太陽を中心に描かれ、その軌道は非常に偏心率が高い。彗星はその軌道で太陽に最も接近する点（近日点）と最も遠ざかる点（遠日点）がある。

\textbf{命題XXXV}：彗星の尾の形状は彗星が太陽に向かって進むとき、尾は太陽から遠ざかる方向に伸びる。また、彗星が太陽から離れるとき、尾は太陽に向かう方向に伸びる。

\textbf{命題XXXVI}：彗星の尾の長さは彗星が太陽に最も接近する点（近日点）で最大となる。彗星が太陽から遠ざかる点（遠日点）では、尾は短くなる。

\textbf{命題XXXVII}：彗星の尾は太陽光線によって推進される。彗星が太陽に向かって進むとき、尾は太陽から遠ざかる方向に伸びる。また、彗星が太陽から離れるとき、尾は太陽に向かう方向に伸びる。

\textbf{命題XXXVIII}：彗星の軌道は太陽を中心に描かれ、その軌道は非常に偏心率が高い。彗星はその軌道で太陽に最も接近する点（近日点）と最も遠ざかる点（遠日点）がある。

\textbf{命題XXXIX}：彗星の尾の形状は彗星が太陽に向かって進むとき、尾は太陽から遠ざかる方向に伸びる。また、彗星が太陽から離れるとき、尾は太陽に向かう方向に伸びる。

\textbf{命題XL}：彗星の尾の長さは彗星が太陽に最も接近する点（近日点）で最大となる。彗星が太陽から遠ざかる点（遠日点）では、尾は短くなる。

\textbf{命題XLI}：彗星の尾は太陽光線によって推進される。彗星が太陽に向かって進むとき、尾は太陽から遠ざかる方向に伸びる。また、彗星が太陽から離れるとき、尾は太陽に向かう方向に伸びる。

\textbf{命題XLII}：彗星の軌道は太陽を中心に描かれ、その軌道は非常に偏心率が高い。彗星はその軌道で太陽に最も接近する点（近日点）と最も遠ざかる点（遠日点）がある。

\textbf{証明}：彗星の尾の形状や長さは彗星が太陽に向かって進むとき、尾は太陽から遠ざかる方向に伸びる。また、彗星が太陽から


月の軌道を幾何学的に求めた後、その軌道が周期的な横直径と回転時間から推定される天体の位置を決定するために必要な作業は、長期間隔で再帰する天体の軌道を特定することを残します。しかしながら、この問題に対処するためには次の命題を用いることができます。

{}

\begin{center}
\textbf{命題 XLII. 証明 XXI.}

\textit{彗星の軌道を幾何学的に求めた後、それを修正する。}

\textit{作業 1.} 前述した方法で幾何学的に求めた彗星軌道の位置を仮定し、観測によって三つの場所を選択します。
\end{center}

accuratissimis definita, \& ab invicem quam maxime distantia; sitque \( A \) tempus inter primam \& secundam, ac \( B \) tempus inter secundam ac tertiam. Cometam autem in eorum aliquo in Perigaeo versari convenit, vel saltem non longe a Perigaeo abesse. Ex his locis apparentibus inveniantur per operationes Trigonometricas loca tria vera Cometae in assumpto illo plano Trajectori.

Deinde per loca illa inventa, circa centrum Solis ceu umbilicum, per operationes Arithmeticas, ope Prop.\ XXI. Lib.\ I. institutas, describatur Sectio Conica: \& ejus are\ae, radiis a Sole ad loca inventa ductis

\textbf{終焉して，$D$ と $E$ を仮定する。すなわち，$D$ は最初の観測点と第二の観測点間の面積であり，$E$ は第二の観測点と第三の観測点間の面積である。} \textbf{次に，$T$ を全時間とする。この時間において，彗星の速度が第I書の第XVI命題により求められた値で，全体として $D + E$ の面積を描くべきである。}

\textit{\textbf{作図2.}} \textbf{ノドゥム軌道の長さを増加させる。} その長さに20分または30分を加えて，それを $P$ とする。そして，その平面の傾きは日食面に対するままである。次に，前記の三つの観測点から彗星がこの新しい平面上にある位置を求めよ。

三つの真実（先に述べた通り）：その後、その軌道を通過するときの面積 $d$ と $e$、そして全体的な時間 $t$ を描写すべき合計面積 $d + e$ の全期間が与えられる。

\textit{作業3.} \wsp{}最初の操作においてノドゥムの経度を保持し、平面軌道の傾きを平面日食平面上に加える。20分または30分を追加し、それらは $Q$ と呼ばれる。次に、観測された三つの彗星の位置から、以下の場所で求められる：

\[
\textbf{作業3.} \quad \text{最初の操作においてノドゥムの経度を保持し、平面軌道の傾きを平面日食平面上に加える。20分または30分を追加し、それらは $Q$ と呼ばれる。}
\]

\[
\textbf{作業3.} \quad \text{観測された三つの彗星の位置から、以下の場所で求められる：}
\]



この部分は、観測データに基づいて彗星の軌道を計算するための手順を述べています。ノドゥムの経度と平面軌道の傾きが調整され、それらの値が与えられることが示されています。

この新たな平面において、三つの真の位置点が存在し、その軌道はこれらの位置点を通る。そして、これら位置点における二つの面積 $\delta$ と $\epsilon$ を観測によって描き出したことから、全体的な時間 $\tau$ が求められる。この時間 $\tau$ は、指定された位置点における総面積 $\delta + \epsilon$ の形成に必要な時間である。

ここで、$C : A = 1 : B$, $G : D = 1 : E$, $g : d = 1 : e$, $\gamma : \delta = 1 : \epsilon$ と置く。また、最初の観測と第三観測間の真実の時間を $S$ とする。そして、正負符号が観察によって確認された条件下で、数値 $m$ と $n$ を求めることを考える。この条件は $G - C = mG - mg + nG - n\gamma$ である。

\& \( T - S \)は\( mT - mt + nT - n\tau \)と等しい。そして、初めの操作において、\( I \)が軌道平面に対する黄道面の傾きを表し、\( K \)が二次節点間の経度を示すとき、真の軌道平面に対する黄道面の傾きは\( I + nQ \)、真の二次節点の経度は\( K + mP \)となる。そして最後に、初めの操作、二番目の操作、三番目の操作において、\( R \), \( r \), \( \rho \)が軌道上の直線距離を表し、それぞれの直線距離に対応する横断距離は\( \frac{1}{L} \), \( \frac{1}{l} \), \( \frac{1}{\lambda} \)とするとき、真の直線距離は\( R + mr - mR + n\rho - nR \)となる。

\& $\frac{1}{L + ml - mL + n\lambda - nL}$ は、彗星が描く軌道の横径を表す。

しかしながら、横径が与えられた場合、彗星の周期的時間も与えられる。

\textbf{証明終了}

\begin{center}\large \textit{FINIS.}\end{center}

\end{document}
